\documentclass[10pt,a4paper,twoside,bg=print,twocolumn,openany,nodeprecatedcode]{dndbook}
\usepackage[utf8]{inputenc}
\usepackage{tabulary}
\usepackage[normalem]{ulem}
\usepackage{color,soul}
\usepackage{float}
\usepackage{caption}
\usepackage{wrapfig}
\usepackage{tocloft}
\usepackage[toc]{multitoc}
\usepackage[hidelinks]{hyperref}
\raggedbottom
\extrafloats{2000}
\cftsetindents{section}{0em}{0em}
\cftsetindents{subsection}{0em}{0em}
\cftsetindents{subsubsection}{0.2em}{0em}
\setcounter{tocdepth}{1}
\renewcommand*{\multicolumntoc}{3}
\setlist[description]{nolistsep,listparindent=3pt,topsep=6pt}
\date{}
\begin{document}
\frontmatter
\tableofcontents
\mainmatter
\addcontentsline{toc}{chapter}{Introduction: Danger to the Multiverse}
\markboth{INTRODUCTION: DANGER TO THE MULTIVERSE}{INTRODUCTION: DANGER TO THE MULTIVERSE}
\chapter*{Introduction: Danger to the Multiverse}
This adventure celebrates fifty years of Dungeons \& Dragons history. The story spans many beloved settings and wondrous planes of existence. Its cast includes characters iconic to longtime fans. This adventure's stakes involve the fate of all worlds—in other words, the multiverse. If the player characters finish this adventure successfully, they'll reach 20th level and will have thwarted one of the most notorious villains in D\&D's history.

The information in this book is intended for the Dungeon Master only. If you're planning to play through the adventure with someone else as your DM, stop reading now!

\textit{Vecna: Eve of Ruin} is a Dungeons \& Dragons adventure optimized for four to six characters. Characters begin at 10th level and gain a level at the end of each of this book's chapters.

The characters are the heroes of the story. This book describes the locations the characters explore and the challenges they must overcome to successfully complete the adventure. All pertinent details about the adventures' settings and locations are covered in this book.

\section{Existence in Peril}
For many years, the lich Vecna has imagined remaking existence. When this adventure starts, Vecna is already a god. Never truly satisfied, he yearns to become the most powerful being in the multiverse. Unfortunately for his enemies, Vecna recently devised the Ritual of Remaking to turn his dreams into reality.

\subsection{The Ritual of Remaking}
Vecna's Ritual of Remaking involves gathering powerful secrets, extracting energy from them, and interweaving this energy with the lich's essence and magic. After Vecna and his cultists extract enough secrets, the lich plans to weave the ritual alone in the Cave of Shattered Reflection on the plane of Pandemonium.

The cave holds the power to harness energy and reveal fundamental truths, making it the perfect location for the ritual. The combined power of Vecna's magic and the collected secrets reaches critical mass at the end of the ritual, triggering an explosion of sufficient magnitude to destroy Pandemonium and unravel the multiverse. Vecna, with his consciousness preserved inside a magical singularity, can then reknit the multiverse exactly as he wants, sealing the multiverse's terrible fate.

Vecna plans to elevate himself above all others and to rewrite the histories and details of the entire multiverse. In Vecna's new reality, all people serve him, he has never known defeat, and all who oppose him live in torment.

By the time this adventure begins, the lich has gathered a hoard of secrets and is nearly ready to begin the Ritual of Remaking. Only the characters can stop Vecna's evil plan.

\section{Vecna's History}
Throughout the adventure, Vecna proceeds with his plans as if nothing can stop him. However, the waves of magic the lich and his cult unleash in preparing for the Ritual of Remaking garner the attention of powerful individuals.

One is Kas, a vampire warrior who once served Vecna but is now Vecna's eternal enemy. Kas learns the details of Vecna's plans and plots to co-opt the ritual for himself.

\subsection{Kas and Vecna}
Hundreds of years ago, Vecna and Kas were close associates. They were driven men who shared similar evil outlooks. Kas admired Vecna's sadism and thirst for power, while Vecna valued Kas's ferocity and cruelty.

By the time Vecna became a lich and the despotic ruler of an empire on the world of Oerth, Kas was the leader of Vecna's armies and the lich's most trusted lieutenant. Kas pledged his life to Vecna, and Vecna gave Kas generous gifts. The lich created the \textit{Sword of Kas} (see the \textit{Dungeon Master's Guide}) for his lieutenant. This intelligent artifact of pure evil was meant to serve Kas unquestioningly, but the sword soon developed other ideas.

Over time, Kas and Vecna began disagreeing about strategies to expand the lich's empire. Convinced Kas was more gullible than Vecna, the sword urged Kas to kill and supplant Vecna. The sword wanted nothing less than full control of Vecna's empire through Kas.

Kas finally betrayed his liege when he confronted Vecna in the lich's tower. Kas killed Vecna, but before Vecna died, Vecna flung Kas across the multiverse. Kas lost the \textit{Sword of Kas} in his flight. All that was left of Vecna after he died was one hand and one eye. These eventually became artifacts known as the \textit{Eye and Hand of Vecna} (see the \textit{Dungeon Master's Guide}).

In the aftermath of this battle, Kas transformed into a vampire. He became the ruler and prisoner of a Domain of Dread (a mist-bordered realm in the Shadowfell) called Tovag. Eventually, the Dark Powers whispered to Kas that Vecna had risen again, becoming an evil god of secrets and magic on Oerth.

Vecna's defeat of Kas grates on the warlord's ego. Kas aches to annihilate Vecna. Meanwhile, Vecna has been building his power, though the lich yearns to finally destroy his former lieutenant.

\section{Kas's Plan}
Shortly after Vecna began traversing the multiverse to gather secrets, Kas learned of the lich-god's plan from the Dark Powers. Upon working out a bargain with those powers, Kas devised a plot to usurp the power Vecna was gathering.

The moment before the lich unravels existence, Kas plans to slay Vecna and step into the lich's place, reshaping the multiverse to his own whims.

\subsection{Domain of Dread Parolee}
Kas uses an artifact called the \textit{Crown of Lies} to impersonate a powerful wizard and manipulate others into retrieving an item that will allow him to free a demon lord ally named Miska the Wolf-Spider. Together, the two can defeat Vecna. Kas is attuned to the crown throughout this adventure.

\subsection{The Wizards Three}
When the Dark Powers divulged Vecna's plan to Kas, the powers also whispered that Lady Alustriel Silverhand was marshaling allies to stop the lich. Kas didn't relish the thought of facing the might of one of the multiverse's most powerful good-aligned spellcasters, so the warlord devised a plot to trick Alustriel into helping him defeat Vecna.

With the help of spies, Kas learned that Alustriel intended to summon two powerful friends to oppose Vecna: the wizards Tasha and Mordenkainen. (For more information about Alustriel, Tasha, and Mordenkainen, see appendix B.)

Kas knows Tasha only by reputation, and he holds a grudging respect for the daughter of Baba Yaga. Kas is more familiar with Mordenkainen, as the two have clashed on Oerth in the past.

Kas disguised himself as Mordenkainen using the \textit{Crown of Lies} and received Alustriel's summons to her sanctum in Sigil, the City of Doors. Throughout this adventure, the real Mordenkainen is traveling the multiverse, unaware of Kas's impersonation.

When Kas arrives in the sanctum, he realizes that Alustriel plans to combine her magic with Tasha's and Mordenkainen's to craft a \textit{Wish} spell to thwart Vecna. Alustriel intends to dispel the collective power of the secrets Vecna gathered, rendering the lich's Ritual of Remaking useless.

Despite looking identical to Mordenkainen, Kas has no significant magical power. This lack sabotages the \textit{Wish} spell exactly as the warlord hoped, allowing him to suggest an alternate course of action: reassemble the artifact called the \textit{Rod of Seven Parts} and use it to defeat Vecna instead. Kas is lying; he wants the rod only for his own selfish ends.

\begin{DndSidebar}[float=!b]{Kas's Sword}
Kas isn't in possession of the legendary artifact called the \textit{Sword of Kas} during this adventure. However, if the characters wish to find it and use it against the warlord, you might place the artifact somewhere in this adventure for them to find. See the \textit{Dungeon Master's Guide} for more information about the \textit{Sword of Kas}.



\end{DndSidebar}

\subsection{Confronting Vecna}
When the spellcasters' \textit{Wish} spell goes awry in chapter 2 of this adventure, the player characters are shunted to the sanctum in Sigil. Kas doesn't expect this development, but the characters are the perfect patsies. Kas planned to retrieve the pieces of the \textit{Rod of Seven Parts} himself, but now he sends the characters to do so. Kas plans to use the rod to free the demon lord Miska the Wolf-Spider, his ally against Vecna, though he falsely claims that the rod is the key to stopping Vecna. (Read more about Kas's plan in chapters 2 and 9.)

\subsubsection{The Rod's Pieces}
During this adventure, the characters might recover all the pieces of the \textit{Rod of Seven Parts}. Each piece's location is presented in the Rod Piece Locations table below.

\begin{DndTable}[header=Rod Piece Locations]{XcXXX}

Piece & Chapter & Place & Setting & Location\\
First & 2 & Underdark, Toril & Forgotten Realms & A covert base of Lolth worshipers\\
Second & 3 & Astral Sea & Spelljammer & Near the wreckage of a spelljamming ship\\
Third & 4 & Mournland & Eberron & The control room of a deactivated colossus\\
Fourth & 5 & Death House, Barovia & Ravenloft & The house's dungeon\\
Fifth & 6 & Northern Dargaard Mountains, Krynn & Dragonlance & The Three Moons Vault's upper level\\
Sixth & 7 & Isle of Serpents, Oerth & Greyhawk & The Tomb of Wayward Souls\\
Seventh & 8 & Avernus, Nine Hells & Planescape & The Red Belvedere casino\\
\end{DndTable}

\section{Adventure Summary}
The adventure is split into three narrative components:


\begin{description}
\item[Discovering Vecna's Activity.]
In chapter 1, the characters stumble upon a cult of Vecna performing a ritual to extract secrets. When the characters clash with the cult, they become linked to the lich, which sets them on the path to discover Vecna's plot to remake the multiverse.
\item[Seeking the Rod.]
In chapter 2, the characters' link to Vecna shunts them to Sigil when powerful spellcasters botch a \textit{Wish} spell to stop the lich. The characters are asked to retrieve the pieces of the \textit{Rod of Seven Parts}, as detailed in chapters 2 through 8.
\item[Stopping the Ritual of Remaking.]
In chapter 9, "Mordenkainen" reveals himself as Kas in disguise. Kas steals the reassembled rod and heads to the plane of Pandemonium to co-opt Vecna's ritual. The characters can chase and confront Kas, but they must stop the Ritual of Remaking from being completed.
\end{description}

\subsection{Chapter 1 Summary}
Chapter 1 takes place in Neverwinter, a city on the Sword Coast on the world of Toril. The characters are investigating the disappearance of high-ranking nobles in Neverwinter when they stumble upon a cult of Vecna. The characters stop the cult from performing a ritual to steal a kidnap victim's secrets. In so doing, they are shunted to Evernight, a gloomy version of Neverwinter in the Shadowfell, and receive Vecna's Link, which ties their fate to that of the lich-god.

\subsection{Chapter 2 Summary}
The first part of chapter 2 takes place in a secret sanctum in Sigil, the City of Doors. The second part takes place in the Underdark on Toril. When two powerful spellcasters and an impostor try to cast a \textit{Wish} spell to stop Vecna's plot, the botched spell latches on to Vecna's Link and pulls the characters to the trio's location. Lady Alustriel Silverhand, Tasha, and Mordenkainen (secretly Kas) are baffled by this development, but Mordenkainen pretends to improvise a plan to stop Vecna. He falsely claims that the \textit{Rod of Seven Parts}, a powerful artifact, might be the only way to stop Vecna now. He knows the location of the first piece, which is in a covert base for Lolth operatives, and sends the characters after it. Each rod piece will point the way to the next piece, and all pieces are in different locations in multiverse.

\subsection{Chapter 3 Summary}
Chapter 3 takes place on the Astral Plane. The characters learn that the second rod piece is in the wreckage of a spelljamming ship called the \textit{Lambent Zenith}. The ship was carrying the piece when it crashed into a dying god's body adrift on the Astral Sea. The characters soon discover that a dragonlike creature devoured the piece and retreated into the heart of the god. The characters must confront the creature and retrieve the rod piece.

\subsection{Chapter 4 Summary}
Chapter 4 takes place on the continent of Khorvaire in the world of Eberron. The characters learn that the third rod piece is located in the Mournland, a magical wasteland. Eventually, the characters discover that the piece is located within an enormous, deactivated construct called Landro. The characters must navigate this colossus and retrieve the rod piece from Landro's engine room.

\subsection{Chapter 5 Summary}
Chapter 5 takes place in Barovia, a Domain of Dread. The characters learn that the fourth rod piece was taken to a haunted residence called Death House. Cultists and the infamous vampire lord Strahd von Zarovich hope to obtain the rod piece themselves, and the characters must thwart all threats before they can claim the piece.

\subsection{Chapter 6 Summary}
Chapter 6 takes place on the world of Krynn. The characters believe the fifth rod piece is inside a gigantic tree, but with the help of good-aligned werewolves, they learn that an evil group loyal to the mighty warlord Lord Soth took the piece to a complex called Three Moons Vault. The characters must infiltrate the complex and retrieve the rod piece, perhaps freeing a resistance leader in the process.

\subsection{Chapter 7 Summary}
Chapter 7 takes place on the Isle of Serpents on the world of Oerth. The characters learn that the sixth rod piece is in the crypt vault of the Tomb of Wayward Souls, a dungeon created by the lich Acererak. The characters must defeat a false lich that calls himself Rerak—a less-powerful version of the Acererak—before they can claim the rod piece.

\subsection{Chapter 8 Summary}
Chapter 8 takes place near the dragon god Tiamat's lair in Avernus, the first layer of the Nine Hells. The characters learn Tiamat obtained the seventh rod piece and stashed it somewhere near her lair. The characters discover that the rod piece is inside a casino called the Red Belvedere. They must infiltrate a members-only section and either defeat a champion of Tiamat or convince the champion to surrender the rod piece.

\subsection{Chapter 9 Summary}
Chapter 9 starts in Sigil and continues in Pandesmos, the first layer of the chaotic plane of Pandemonium. The rest of the adventure takes place on that plane. With all pieces of the \textit{Rod of Seven Parts} retrieved, the characters must allow the spellcasters to examine the rod—or so "Mordenkainen" claims. The archmage reveals himself as Kas and uses the rod to subdue Alustriel and Tasha. Kas flees to Pandemonium to free Miska the Wolf-Spider and usurp Vecna's ritual. The characters must race to Pandemonium in pursuit of Kas, who has discovered where Vecna's ritual is taking place.

\subsection{Chapter 10 Summary}
In chapter 10, the characters follow Kas through Pandemonium toward the Ruinous Sea, an ocean of swirling, impassable, chaotic magic. On the coast, a battle rages between the demonic forces of Kas, who is attempting to free Miska the Wolf-Spider from his prison in the nearby Ruinous Citadel, and the demon-god Lolth, who is allied with Vecna. The characters can manipulate the battle, or they can race straight to Kas, who eventually reveals the location of Vecna's ritual.

\subsection{Chapter 11 Summary}
In chapter 11, the characters must descend into the Cave of Shattered Reflection, where Vecna weaves his Ritual of Remaking. The ritual is nearing its end, and the lich-god has created several demiplanes that offer glimpses of the multiverse he is creating. The characters must navigate these demiplanes to find the key to entering Vecna's ritual chamber. Once inside, the characters must stop the ritual, which has left the lich-god in a weakened state.

\section{Running the Adventure}
To run this adventure, you need the fifth edition core rulebooks (\textit{Player's Handbook}, \textit{Dungeon Master's Guide}, and \textit{Monster Manual}).

\begin{DndReadAloud}{}
Text that appears in a box like this is meant to be read aloud or paraphrased for the players when their characters first arrive at a location or under specific circumstances, as described in the text.

\end{DndReadAloud}

The \textit{Monster Manual} contains stat blocks for most of the creatures encountered in this adventure. The remaining stat blocks can be found in appendix A or B, as indicated in the text, or in the encounters in which they appear.

When a creature's name appears in \textbf{bold} type, that's a visual clue pointing you to its stat block as a way of saying, "Hey, DM, you better get this creature's stat block ready. You're going to need it."

Spells and equipment mentioned in the adventure are described in the \textit{Player's Handbook}. Magic items not described in the adventure's text are described in the \textit{Dungeon Master's Guide}.

\subsection{Using the Maps}
This book contains a number of interior maps.

\subsubsection{Interior Maps}
Maps that appear in this book are for the DM only. As the characters explore locations on a given map, you can redraw portions of the map on graph paper, an erasable mat, or another surface to help your players visualize locations with unusual shapes or features. Your hand-drawn maps needn't be faithful to the originals, and you can alter a map's features as you see fit. Nor do your maps need to be painstakingly rendered. Omit details that aren't readily visible (such as secret doors and other hidden features) until the characters detect and interact with them. For example, locked doors are indicated on the maps with dots, but you need not include this detail in your hand-drawn maps.

One of the maps within shows the sanctum in Sigil that serves as the characters' home base for most of this adventure. This map is further described in chapter 2, but don't share information with the players that their characters wouldn't know. For example, the characters wouldn't know which magic items Alustriel keeps in area S3 of the sanctum unless they examine that area carefully. The map shows the sanctum as it is before the characters arrive.

\subsection{Nondeadly Resolutions}
This adventure sets up a number of encounters for the characters to fight their foes. However, other nondeadly resolutions are equally valid ways to resolve enemy encounters. The characters might knock out enemies, intimidate them into running away, bribe them for information, or otherwise find creative ways to resolve conflicts. Use your discretion, and if the characters attempt to resolve encounters without violence, go with it if the story allows.

\section{Character Creation}
Before starting this adventure, consider spending your first game session helping your players create characters. This adventure recommends 10th-level characters to start, so make sure your players have appropriately leveled characters before play begins. Chapter 1 provides instructions for starting this adventure with 7th-, 8th-, 9th-, or 11th-level characters.

\subsection{Existing Characters}
One or more of your players might want to continue playing characters who have successfully completed previous D\&D adventures, and that's fine! Any character from any campaign setting is appropriate for this adventure, as long as they aren't higher than 11th level.

The beginning of this adventure takes place in the city of Neverwinter on the world of Toril (in the Forgotten Realms campaign setting). If a character is from a different world or is from somewhere else in that setting, work with the player to devise a reason why that character is in Neverwinter and answers Lord Neverember's call for help in solving the disappearance of several nobles. The reason might range from the magical (an unexpected magical backlash transports them to the city) to the mundane (a family connection brings them there).

If a character takes a break after their previous adventure, consider rolling on or picking an option from the Purpose in Neverwinter table, expanding on it as necessary, to provide a reason the character has come to Neverwinter.

\subsection{New Characters}
Players might create a 10th-level character from scratch for this adventure. For ease of introduction, the character might be from Neverwinter or somewhere else nearby.

If one or more of your players want to create characters from more distant lands, use the suggestions provided in the previous section for why the characters are now in Neverwinter. If they don't already have one, the characters need a reason for Lord Neverember to request their help in chapter 1, even if it's simply because their heroics are widely known.

Regardless, 10th-level characters have already had long careers and earned their abilities through experience. These characters have likely accomplished impressive deeds, so encourage your players to describe how the characters reached 10th level.

The 10th-Level Backstories table contains sample backstories for new characters. Roll on the table or pick your favorite option when working with your players to determine their characters' backstories.

\begin{DndTable}[header=10th-Level Backstories]{cX}

d6 & Deeds Accomplished\\
1 & The character began their adventuring career investigating zombies that rose from a once-peaceful graveyard. Their investigations revealed a cult. The character fought the cultists and defeated the cult's archmage leader.\\
2 & As a bodyguard for a high-ranking political or business figure, the character has staved off attacks from assassins and demons.\\
3 & The character hunts evil dragons or another dangerous kind of creature.\\
4 & The character is a private investigator who, for the right price, retrieves kidnapped loved ones and recovers stolen documents and items from crafty villains.\\
5 & The character studies and preserves nonmagical historical artifacts and has traveled widely in pursuit of rare items of deep significance. While securing items for local museums, the character has thwarted threats ranging from ettins to glabrezus.\\
6 & The character is in the service of a god and on a mission to protect the helpless. The work doesn't pay well, but the character has saved children from hungry monsters, thwarted rampaging monsters in villages, and protected the poor from greedy overlords.\\
\end{DndTable}

\begin{DndTable}[header=Purpose in Neverwinter]{cX}

d6 & Reason\\
1 & The character is in the employ of a powerful merchant, mage, or monarch who sent the character to Lord Neverember as a favor. Neverember needed the character's help to handle growing problems with undead in Neverwinter. The emergency involving the missing nobles is the character's most recent assignment.\\
2 & After a long struggle, the character defeated a powerful sorcerer. In the throes of death, the sorcerer's magic backfired, opening a portal and shunting the character to Neverwinter. The character has been working for Lord Neverember as a mercenary while figuring out what to do next.\\
3 & The character fought an adult black dragon. The fight wasn't going well, and nothingness enveloped the character as the character was about to die. Instead of dying, the character woke up in Neverwinter outside Lord Neverember's villa. It's up to you to decide why; perhaps a god's or spellcaster's intervention is involved.\\
4 & The character defeated an evil leader and was given the pick of the spoils from the villain's lair. While examining the items, the character stumbled upon a \textit{Cubic Gate} that sent them to Neverwinter. The character agreed to work for Lord Neverember in exchange for the spellcasting services they need to return home.\\
5 & The character is an expert at retrieving kidnapped individuals, even those on other planes of existence. The family of one of the kidnapped nobles in chapter 1 reached out to the Harpers, who recommended the character's services. The family promised a high price for the character to come to Neverwinter and look for the missing loved one.\\
6 & The character is an explorer charting the multiverse. During their travels, they've come to recognize the foreboding sense of danger they feel whenever they're about to encounter evil. The character can't shake a sense of dread about Neverwinter, so they've come to the city to ask Lord Neverember if there's a threat.\\
\end{DndTable}

\section{The Power of Secrets}
Vecna uses stolen secrets to power his ritual to remake the multiverse. When the characters stumble on the cult of Vecna, whose members are trying to extract secrets from a captive, they gain access to magic fueled by powerful secrets. The characters can spend secrets like currency once they receive Vecna's Link in chapter 1.

Throughout the adventure, the characters can learn many secrets. Each chapter's beginning includes a "Power of Secrets" section that lists each secret that can be used with these rules. The Secrets Tracker in appendix C helps you keep track of secrets the characters have learned. The Secrets Tracker includes spoilers, so keep it hidden from the players.

\subsection{Learning Secrets}
Every time the characters learn a powerful secret, note it on the Secrets Tracker. Once the characters learn a secret, they can't learn it again. When one character learns a secret, it counts as a secret learned for the whole party.

\subsubsection{Revealing Secrets}
A character can magically spend a secret like currency, revealing it to the multiverse to gain a momentary boon. To do this, the character must use an action to whisper the secret into the wind. The secret is then gone from the minds of every character in the party. There is one exception; if the characters spend the secret they learn from Kas in chapter 10 about Vecna's location, they still know where the lich-god weaves his ritual.

When a character spends a secret, every character in the party gains advantage on d20 rolls for 1 minute.

The characters might hear the information from the secret again, but they can't spend a secret more than once. If the characters share a secret with a nonplayer character, that secret immediately loses its power and counts as revealed. It remains in the minds of every character in the party, but the characters gain no benefit from revealing the secret.

\subsubsection{Keeping Secrets}
When the characters confront Vecna in the Cave of Shattered Reflection in chapter 11, they can use any number of secrets they've kept to help thwart the lich-god's Ritual of Remaking. See chapter 11 for more details about how secrets the characters kept can affect their confrontation with Vecna.

\subsection{Linked to Vecna}
Once the characters are metaphysically linked to Vecna and can spend powerful secrets, their connection to the lich-god might manifest in additional ways, at your discretion.

For instance, the characters might periodically see mental images of Vecna weaving his ritual in a mysterious, crystal-filled cave. Or the characters might dream about foreboding Vecnan images, including the lich-god's unholy symbol or robed cultists worshiping Vecna. Reminding the characters periodically about the lich-god's evil plan adds a sense of urgency to the adventure.

\section{Challenge Ratings}
The Stat Blocks by Challenge Rating table sorts the creatures in this book by challenge rating. It also lists their creature type and where in the book they appear.

\begin{DndTable}[header=Stat Blocks by Challenge Rating]{cXcc}

CR & Stat Block & Creature Type & Chapter\\
1 & \textbf{Warforged warrior} & Construct & A\\
5 & \textbf{Night scavver} & Monstrosity & A\\
5 & \textbf{Kakkuu spyder-fiend} & Fiend & A\\
6 & \textbf{Black rose bearer} & Undead & A\\
6 & \textbf{Moonlight guardian} & Construct & A\\
6 & \textbf{Priest of Osybus} & Humanoid & A\\
7 & \textbf{Blade scout} & Construct & A\\
7 & \textbf{Lost sorrowsworn} & Monstrosity & A\\
8 & \textbf{Bone roc} & Undead & A\\
8 & \textbf{Inquisitor of the Tome} & Humanoid & A\\
8 & \textbf{Star angler} & Monstrosity & A\\
8 & \textbf{Whirling chandelier} & Construct & A\\
9 & \textbf{Blade lieutenant} & Construct & A\\
9 & \textbf{Lonely sorrowsworn} & Monstrosity & A\\
9 & \textbf{Necromancer wizard} & Humanoid & A\\
10 & \textbf{Eye monger} & Aberration & A\\
10 & \textbf{Mirror shade} & Undead & A\\
11 & \textbf{Degloth} & Fiend & A\\
11 & \textbf{Glaive} & Humanoid & 4\\
11 & \textbf{Spiderdragon} & Monstrosity & A\\
11 & \textbf{Vlazok} & Fiend & A\\
12 & \textbf{Blazebear} & Monstrosity & A\\
12 & \textbf{Granite juggernaut} & Construct & A\\
13 & \textbf{Deadbark dryad} & Fey & A\\
13 & \textbf{Adult lunar dragon} & Dragon & A\\
13 & \textbf{Phisarazu spyder-fiend} & Fiend & A\\
14 & \textbf{Cadaver collector} & Construct & A\\
14 & \textbf{Hazvongel} & Fiend & A\\
15 & \textbf{Green abishai} & Fiend & A\\
15 & \textbf{Borthak} & Monstrosity & A\\
15 & \textbf{Deathwolf} & Undead & A\\
15 & \textbf{Relentless impaler} & Fiend & A\\
15 & \textbf{Strahd, Master of Death House} & Undead & B\\
17 & \textbf{Blue abishai} & Fiend & A\\
17 & \textbf{Hertilod} & Monstrosity & A\\
17 & \textbf{Quavilithku spyder-fiend} & Fiend & A\\
18 & \textbf{Citadel spider} & Monstrosity & A\\
18 & \textbf{Cosmic horror} & Aberration & A\\
19 & \textbf{Red abishai} & Fiend & A\\
19 & \textbf{Raklupis spyder-fiend} & Fiend & A\\
19 & \textbf{Lord Soth} & Undead & B\\
19 & \textbf{Tasha the Witch} & Humanoid & B\\
21 & \textbf{Astral dreadnought} & Monstrosity & A\\
21 & \textbf{False lich} & Undead & A\\
21 & \textbf{Alustriel Silverhand} & Humanoid & B\\
22 & \textbf{Camlash} & Fiend & 10\\
23 & \textbf{Kas the Betrayer} & Undead & B\\
23 & \textbf{Windfall} & Humanoid & 8\\
24 & \textbf{Miska the Wolf-Spider} & Fiend & B\\
26 & \textbf{Vecna the Archlich} & Undead & B\\
\end{DndTable}

\chapter{Return from Neverdeath Graveyard}
As the adventure begins, the characters are established heroes currently in the city of Neverwinter on the Sword Coast. Several calamities have battered Neverwinter in the recent past. The greatest was the eruption of nearby Mount Hotenow, which nearly destroyed the city forty years ago, though most of the damage has since been repaired.

Neverwinter is ruled by its Lord Protector, Dagult Neverember. Rising to power through a tenuous claim of descendance from one of Neverwinter's past heroes, Lord Neverember has nevertheless provided stable leadership.

Unknown to the authorities, a cult of Vecna operates in the catacombs beneath Neverwinter's sprawling Neverdeath Graveyard. Cult members have been kidnapping city residents who carry significant secrets, draining their knowledge and their souls in a fell ritual and passing the collected secrets to Vecna as he gathers power for his Ritual of Remaking. (See the introduction for more information about Vecna's plot.) In the process, the kidnap victims become creatures robbed of their knowledge and volition.

\section{Running This Chapter}
In this chapter, the characters discover a cult of Vecna preparing four kidnap victims for a ritual in the catacombs beneath Hallix Mausoleum. Disrupting this ritual hurls the characters and an elf scholar named Eldon Keyward into Evernight, Neverwinter's sinister reflection in the Shadowfell. To return home, the characters must confront the lonely legacy of the Dolindar family and find a rift that leads back to Neverwinter.

\subsection{Character Advancement}
The characters should be 10th level when this chapter begins; see the "Lower-Level and Higher-Level Characters" sidebar for accommodating characters of other levels. If the characters are 10th level or below, they gain a level after returning to Neverwinter from Evernight.

\subsection{Power of Secrets}
The characters can learn three secrets in this chapter applicable to the Power of Secrets rules found in this book's introduction. These secrets are tied to three NPCs whom the characters encounter in the Neverdeath Catacombs:


\begin{description}
\item[Indrina's Secret.]
A prisoner named Indrina knows that Lord Neverember doesn't have a legitimate claim to his title. The characters can learn this secret in area C20.
\item[Sarcelle's Secret.]
A prisoner named Sarcelle recently received a disturbing vision. The characters can learn about Sarcelle's recent vision in area C5.
\item[Umberto's Secret.]
The characters can learn about Umberto's secret role as a historian of Vecna in area C11.
\end{description}

The first time the characters learn one of these secrets, they feel a sense that the information they've discovered is important. Describe the Power of Secrets rules to the players at this time, but don't let them spend any secrets yet. When the characters receive Vecna's Link, they can spend secrets using the Power of Secrets rules, as described in the introduction.

\subsection{Lower-Level and Higher-Level Characters}
This chapter is a preamble to the adventure's primary plot. You can run this chapter for lower-level characters, adapting it as described below. If your characters are 7th, 8th, or 9th level, remove creatures as noted in the Creatures to Remove table.

Additionally, have the \textbf{marid} in area C10 surrender if reduced to fewer than 150 hit points (rather than 100 hit points).

\begin{DndTable}[header=Creatures to Remove]{cX}

Area & Remove\\
C2 & Two wights\\
C7 & Two water weirds\\
C14 & Two cult fanatics\\
C16 & One mage\\
C17 & Two cult fanatics\\
C25 & One nothic\\
C26 & Three nothics\\
Evernight Awakening & Three ghouls\\
B1 & One vampire spawn\\
\end{DndTable}

If the characters are 11th level at the start of this chapter, they don't gain a level for completing it.

\section{Getting Started}
This adventure begins when Lord Dagult Neverember summons the characters to his modest villa in Neverwinter. Several local guards are present, as are three priests of Oghma—a god of inspiration, invention, and knowledge—from Neverwinter's House of Knowledge.

When the characters enter his audience chamber, Lord Neverember is busy discussing politics with his advisers. He breaks off the discussion, gives the characters a smile of recognition, and says the following:

\begin{DndReadAloud}{}
"Greetings, my heroic friends! I'm so glad you came. I daresay, terrible events are afoot. Specifically, four prominent citizens have been kidnapped in the past several days. May I count on your help in rescuing them?"

\end{DndReadAloud}

Lord Neverember describes the kidnap victims as follows:

\textbf{Eldon Keyward} is a highly knowledgeable scholar who specializes in the Outer Planes.

\textbf{Indrina Lamsensettle} is a human actor who moves in Neverwinter's highest social circles.

\textbf{Sarcelle Malinosh} is a human wild-magic sorcerer who plumbs the mysteries of the Outer Planes.

\textbf{Umberto Noblin} is a gnome historian who has written books on various deities.

Lord Neverember confirms that each victim was kidnapped at night. The victims don't know each other, and there appears to be no connection between them.

Lord Neverember funded divinations from the House of Knowledge, hoping to find the victims. The priests reported that the mystical trail of the victims ends at a specific place: Hallix Mausoleum in Neverdeath Graveyard. The priests worry that their inability to see inside this mausoleum means that an unknown opponent is blocking their divinations.

Lord Neverember asks the characters to investigate the disappearances at Hallix Mausoleum. He promises each character a fine house in Neverwinter if they can recover the four missing townspeople and bring the kidnappers to justice.

\subsection{Neverdeath Graveyard}
Lord Neverember and the priests provide an overview of Neverdeath Graveyard, which contains two sprawling, connected cemeteries: the Main Graveyard and the Pauper's Graveyard. A thick stone wall separates the crypts of the wealthy from the graves of the poor.

The Main Graveyard holds several mausoleums, some with expansive underground chambers. The Pauper's Graveyard features numerous simple headstones, but a few civic-minded citizens funded communal catacombs when the graveyard was first built.

\section{In the Graveyard}
The wandering zombies and skeletons of Neverdeath Graveyard don't pose a challenge to a higher-level party, so the characters can reach Hallix Mausoleum without trouble. The characters don't yet realize that the mausoleum leads to a network of catacombs that extends beneath both the Main Graveyard and the Pauper's Graveyard.

\subsection{General Features}
The following features are common throughout the catacombs and chambers.

\subsubsection{Broken Stone}
The old stones used to build the subterranean areas don't fit together well, leaving space for mold or tangled roots. Vecna's cultists have scribbled symbols on the walls, depicting staring eyes and left hands.

\subsubsection{Ceilings}
Ceilings are 10 feet high in passages and 15 feet high in rooms unless otherwise noted.

\subsubsection{Doors}
Doors throughout the area are made of heavy stone with metal hinges. The cult keeps the doors well-oiled, so they don't make noise when opened. All doors are unlocked unless otherwise noted.

\subsubsection{Lighting}
Areas C1–C12 are dark, and characters must have darkvision or a light source to see. Cult members frequent areas C13–C26, so lanterns hung on wall hooks create bright light there.

\subsubsection{Vestiges of the Waterclock Guild}
Part of the catacombs once belonged to an occult organization called the Waterclock Guild. The guild members are gone, though their bound elementals and clockwork mechanisms remain. A network of pipes runs through areas C7–C12 and C14–C15, indicated on map 1.1 by solid lines. The sound of dripping water echoes throughout these areas.

\subsection{Neverdeath Catacomb Locations}
The following locations are keyed to map 1.1.

\subsubsection{C1: Hallix Mausoleum}
\begin{DndReadAloud}{}
The towering stone mausoleums in Neverdeath Graveyard cluster near the wall separating the Main Graveyard from the Pauper's Graveyard to the west. Hallix Mausoleum is a squat, unassuming granite block in the shadow of larger monuments to the west and south. Its metal double door bears a rusty broken chain and a padlock that hangs off the door.

\end{DndReadAloud}

The well-oiled door opens noiselessly. The crypt interior is dusty, with numerous tracks leading to a descending staircase at the rear of the room.

Against the walls rest six stone coffins, three on each side. A stone slab engraved with a name, birth year, and death year covers each coffin. These members of the Hallix family died forty years ago, after Mount Hotenow erupted. The coffins are empty except for scraps of cloth and bits of bone; the cult's ghouls ate the former occupants.

\subparagraph{Noisy Investigations}
If the characters make a lot of noise here, the wights in area C2 investigate.

\subsubsection{C2: Lower Mausoleum}
If the wights described below moved to investigate area C1, omit the last sentence when reading aloud:

\begin{DndReadAloud}{}
Stone stairs descend from Hallix Mausoleum to a large subterranean chamber with stone coffins sitting on sturdy shelves. Part of the west wall has collapsed, creating an opening into another chamber. Marching around the chamber are five pale, desiccated warriors wearing wicked-looking armor.

\end{DndReadAloud}

The cult pressed five wights into service as guards. The wights know the cultists by appearance and don't attack them, but they attack anyone else.

The older Hallix corpses once buried here were fed to the cult's ghouls.

\subparagraph{Secret Door}
One of the empty coffins contains no evidence of a former occupant. The back of this coffin hides a panel with a latch that causes the wall behind it to swing aside. A character who searches the coffin or wall and succeeds on a DC 18 Intelligence (Investigation) or Wisdom (Perception) check finds this secret door. Neither the cultists nor the wights are aware of it. The passageway beyond ends at another secret door that is easily spotted and opened from inside the tunnel. It leads to area C9.

\subparagraph{Treasure}
One open coffin contains four wool cloaks worth 10 gp each and two wide hats worth 5 gp each, which the cultists use to travel inconspicuously aboveground. There is also a \textit{Potion of Invisibility} the cultists were saving for an emergency.

\subsubsection{C3: Uneven Chamber}
\begin{DndReadAloud}{}
Roots protrude through cracks in the ceiling here. A stone stairway in the southeast corner has collapsed, and the nearby walls have crumbled. Three doors in the north wall are shut, and the middle door bears a new padlock. To the west, stairs lead to a small balcony that overlooks the room from five feet above, with just enough room for a door painted with an eye.

\end{DndReadAloud}

This room is inaccessible from the surface since the stairway leading upward has collapsed. See area C5 for more about the padlocked door.

\subsubsection{C4: Trapped Grate}
\begin{DndReadAloud}{}
A metal grate in the floor of this ten-foot-square room blocks access to a shallow stone pit holding a small gold harp, a handful of loose papers, and a piece of bloody cloth.

\end{DndReadAloud}

The iron grate covers a pit that's 5 feet square and 3 feet deep. The openings in the iron grate are 5 inches square. A character who can reach the harp can carefully tilt it and slide it through the grate with a successful DC 18 Dexterity (Sleight of Hand) check. If this check fails by 5 or more, the harp falls back into the pit.

A cultist of Vecna tried to lever the harp from the pit but fell victim to the trap on the grate. The trap tore away half of the cultist's jacket—the bloody cloth now at the pit's bottom—and papers tumbled from the cultist's pocket and through the grate. After that mishap, the cultists decided not to press their luck and left the treasure alone.

\subparagraph{Grate Trap}
A character can detect the grate's trap by examining the grate and succeeding on a DC 14 Intelligence (Investigation) check. The trap activates when more than 10 pounds of pressure is placed on the grate. Poisoned blades extend from grooves in the grate, dealing 14 (4d6) slashing damage to whatever triggered the trap, and if the target of the trap is a creature, it must succeed on a DC 14 Constitution saving throw or take 14 (4d6) poison damage. The trap resets after 1 minute.

\subparagraph{Notes}
The papers detail plans to kidnap a Neverwinter aristocrat named Indrina Lamsensettle. The notes include a map of her estate, schedules of her movements, and suggestions that she knows an important secret about Lord Neverember. A scrawl in the margin of a note claims that "her secrets will make a worthy sacrifice." (The characters can learn more by examining Jerot's papers in area C25.)

\subparagraph{Treasure}
The harp is worth 2,500 gp.

\subsubsection{C5: Sarcelle's Cell}
This room's only door is padlocked from the outside. As an action, a character with thieves' tools can use them to try to open the lock, doing so with a successful DC 18 Dexterity (Sleight of Hand) check. The cult's four mages (in areas C14, C16, and C26) each carry a key to the lock. Because the stone is unevenly set around this door, a character could use an action to try to pull the door aside, doing so with a successful DC 17 Strength (Athletics) check.

\begin{DndReadAloud}{}
This old crypt holds a single open coffin containing a few tattered blankets. A pouf of wild black hair sprouts from the end of one of the blankets.

\end{DndReadAloud}

The cultists repurposed this crypt into a cell for one of their intended ritual victims, Sarcelle Malinosh. Sarcelle is dozing inside the coffin, wrapped in the blankets so only her hair is visible.

Sarcelle is a human wild-magic sorcerer whose spellcasting power was stripped during a recent excursion to a distant plane. Until her magic naturally returns, Sarcelle has the game statistics of a neutral \textbf{mage} without Spellcasting. She responds to some questions with cryptic-sounding predictions, but she tries to keep this irritating habit in check.

Sarcelle wants help freeing herself; she explains that without her magic, she feels uneasy and would appreciate being escorted from Neverdeath Graveyard. She can make her way home from there.

\subparagraph{Sarcelle's Secret}
Sarcelle's psychic explorations showed her a glimpse of a dreadful future. She saw the desiccated figure of a man levitating off the ground, gathering evil energy around himself in glowing wisps. The desiccated man then screamed and the energy exploded, causing something terrible to happen. This vision terrified Sarcelle. A character who interacts with Sarcelle and succeeds on a DC 14 Wisdom (Insight) check realizes that something is bothering the sorcerer. If asked about what's upsetting her, Sarcelle shares her vision.

Learning of Sarcelle's vision counts as a secret for the purposes of the Power of Secrets rules found in this book's introduction.

\subsubsection{C6: Supply Room}
\begin{DndReadAloud}{}
Stone shelves in this room contain boxes and bags. A few crates are stacked against the wall.

\end{DndReadAloud}

The cultists emptied this servants' crypt to store supplies such as lantern oil, chains, and manacles.

\subparagraph{Treasure}
Among the supplies are two Potion of Poison labeled "Healing Use Only."

\subsubsection{C7: Fountain Room}
If the characters minimized the water pressure in area C8, omit the first sentence when reading aloud:

\begin{DndReadAloud}{}
Rusted pipes run along the walls and ceiling, and water flows from nozzles in the ceiling pipes. In the center of the room is a deep stone basin that's set into the floor and filled to the brim. The surface ripples, revealing several watery creatures inside. To the south, a closed door is padlocked.

\end{DndReadAloud}

See area C11 for more information about the padlocked door.

A \textbf{water elemental} and two water weirds live in the 25-foot-deep fountain. These creatures are indifferent toward intruders and attack only in self-defense. Once bound to serve the Waterclock Guild, they're now free but enjoy the perpetual "rain" here. The cult bullies these Elementals, so they remain sulking under the water's surface. Determined not to stand for further intrusion, the Elementals rise to attack anyone other than cultists. The water weirds consider the water elemental their leader. If reduced to fewer than 50 hit points, the water elemental retreats to the basin's bottom with any surviving water weirds.

The water elemental enjoys conversation but speaks Aquan only. Characters who are able to communicate with the water elemental can learn the following from it:


\begin{description}
\item[Cult Activity.]
Cultists who worship a god whose symbol is a hand and an eye recently moved into nearby rooms.
\item[Neighbor.]
A fish-headed creature named Shanzezim lives in area C10.
\item[Prisoner.]
The cultists locked a small creature in an adjacent room (area C11).
\end{description}

\subparagraph{Treasure}
A silver bracelet set with seven small diamonds fell to the bottom of the basin. It's worth 150 gp.

\subsubsection{C8: East Pressure Room}
\begin{DndReadAloud}{}
Pipes along the south wall of this room disappear into the walls near the ceiling. A complicated series of cogs and four hand-turned wheels connect to the pipes.

\end{DndReadAloud}

The wheels control the water pressure through the pipes, but they lack gauges to show how turning the wheels affects the pressure. A character can determine that the water flows west, as well as how to maximize or minimize the water pressure, with an hour of trial and error. A character who succeeds on a DC 16 Intelligence (Investigation) check discovers this information in only 10 minutes. Alternatively, the \textbf{marid} Shanzezim in area C10 can describe how to work the wheels.

\subparagraph{Minimum Pressure}
If the characters minimize the water pressure, the nozzles in area C7 stop flowing.

\subparagraph{Maximum Pressure}
If the characters maximize the water pressure here and in area C12, the basins in areas C14 and C15 start to overflow. After 10 minutes, those areas become \textit{difficult terrain} due to flooding. Five minutes after that, the denizens of area C14 come to investigate areas C8 and C12, bringing along both ghouls from area C17. The cultists shout about "teaching those meddling elementals a lesson" as they arrive, allowing the characters time to set up an ambush or another ploy.

\subsubsection{C9: Clockwork Alcove}
\begin{DndReadAloud}{}
Rusty standpipes and interlocking cogs cover the walls of this small alcove.

\end{DndReadAloud}

The cogs here are jammed together and don't move. Whatever mechanism they connect to is inoperable.

\subparagraph{Secret Door}
One cog on the north wall doesn't connect to anything else on the wall around it. A creature must succeed on a DC 14 Intelligence (Investigation) check to find this loose cog. When turned, the cog causes part of the wall to slide away as a secret door. The cultists aren't aware of this door. The small passageway beyond ends at another secret door that is easily spotted and opened from inside the tunnel. It leads to area C2.

\subsubsection{C10: Improvised Workshop}
\begin{DndReadAloud}{}
Rubble chokes the southeast corner of the room, leaving only a small gap near the uneven ceiling. Some of the rubble has been reassembled into a low table, which bears small clockwork components. A hulking, fish-headed creature wearing exquisite silk finery carefully examines the tiny parts.

\end{DndReadAloud}

The creature is a \textbf{marid} named Shanzezim. The marid was bound by the Waterclock Guild and can't leave the crypts belonging to that organization, even though Shanzezim believes the Waterclock Guild has been defunct for years. Not quite ready to test the binding to make an escape, the marid spends time here trying to reassemble one of the Waterclock Guild's most intricate clocks.

If the \textbf{water elemental} in area C7 fled here, it informed Shanzezim about intruders in the area, so the marid attacks right away to drive off the characters. Otherwise, the marid asks the characters what they want. If a fight breaks out, the marid surrenders if reduced to fewer than 100 hit points or if the characters insist that they aren't with the cult.

\subparagraph{Shanzezim's Lore}
If the characters talk to the marid and reassure Shanzezim they're not part of the cult, Shanzezim offers the characters a gold-colored flywheel from the disassembled clock. The marid is chatty and can share the following pieces of information:


\begin{description}
\item[Clock Assembly.]
Shanzezim decided to reassemble a mechanical clock the Waterclock Guild left behind. The marid believes the thousands of parts are all here and reassembly shouldn't take more than another few years. The length of time doesn't bother Shanzezim, since the marid enjoys the work.
\item[Cult Activity.]
Cruel cultists bully the Elementals in the next room and are keeping a prisoner nearby. The prisoner is a gnome whom Shanzezim hasn't yet met.
\item[Pipe System.]
The pipes running through this area lead to a part of the crypt where Shanzezim can't go, but to which the pipes' water flows. Shanzezim describes how to maximize and minimize the water pressure in areas C8 and C12, opining with delight that turning the flow to maximum pressure in both areas should flood out the recently arrived cultists in short order.
\item[Waterclock Guild.]
Shanzezim is bound to a portion of the catacombs once controlled by a vanished organization of artificers and geomancers called the Waterclock Guild. Additional guild catacombs lie past the collapsed portion of this room, but they hold nothing of interest.
\end{description}

\subparagraph{Beyond the Rubble}
It takes several days of labor to clear the rubble so creatures can pass through it, but Shanzezim is right about there being nothing relevant beyond it. If the characters are determined to explore the other Waterclock Guild chambers, you can invent water-themed or clockwork-based denizens and traps for them to encounter.

\subparagraph{Treasure}
The clock parts include a gold-colored flywheel that isn't a part of the clock Shanzezim is trying to assemble. The flywheel thus doesn't interest the marid, who gives it to the characters. The gold-colored flywheel is magical and has the properties of a \textit{Stone of Good Luck}.

\subsubsection{C11: Umberto's Cell}
This room's only door is padlocked from the outside with a new, sturdy lock. As an action, a character with thieves' tools can try to use them to open the lock, doing so with a successful DC 18 Dexterity (Sleight of Hand) check. The cult's four mages (in areas C14, C16, and C26) each carry a key that unlocks it. If the characters are friendly with the \textbf{water elemental} (see area C7) or Shanzezim (see area C10), either is happy to flow into the crack around the door and burst it open from the inside, much to the surprise of this room's occupant, Umberto Noblin.

\begin{DndReadAloud}{}
Water leaks down the walls of this cell and pools on the floor near a rusty drain. A dejected gnome sits on a sodden mattress in one corner.

\end{DndReadAloud}

Umberto Noblin is a gnome historian who is eager to escape Neverdeath Graveyard. Umberto has the game statistics of a lawful neutral \textbf{mage} without Spellcasting.

\subparagraph{Umberto's Secret}
Umberto knows that cultists of Vecna are his kidnappers, as he's one of Neverwinter's preeminent experts on Vecna's history. He initially keeps his expertise from the characters lest they think he's in league with the cult.

To keep his mind off of the nightmare of his capture and imprisonment, Umberto focuses on complaining about the poor cuisine. If the characters free Umberto and share some tasty food with him, he reveals his expertise in Vecna's history. Umberto especially likes food created with or by magic, such as berries from the \textit{Goodberry} spell.

If Umberto reveals his role as a historian of Vecna, it's all he can talk about. He discusses his latest clandestine research project: the ancient rivalry between Vecna and his treacherous lieutenant, the vampire Kas. The historian has kept this research to himself so other scholars don't beat him to publication; Umberto knows he's a slow writer. The gnome's chattering should be endearing rather than irritating, and you can use Umberto to impart basic history about Vecna and Kas as described in the introduction.

Learning about Umberto's secret research topic counts as a secret for the purposes of the Power of Secrets rules in this book's introduction.

\subsubsection{C12: West Pressure Room}
\begin{DndReadAloud}{}
Pipes climb the north wall, disappearing near the ceiling. A complicated series of cogs and three hand-turned wheels connect to the pipes.

\end{DndReadAloud}

The wheels control the pipes' water pressure. As in area C8, a character can determine how to maximize or minimize the water pressure with an hour of trial and error, with 10 minutes of trial and error if a character succeeds on a DC 14 Intelligence (Investigation check), or with Shanzezim's instructions.

\subparagraph{Minimum Pressure}
Minimizing the pressure here doesn't affect the nozzles in area C7.

\subparagraph{Maximum Pressure}
If the characters maximize the water pressure both here and in area C12, the basins in areas C14 and C15 start to overflow (see area C8 for more information).

\subsubsection{C13: Wall Crossing}
\begin{DndReadAloud}{}
This small room appears to be a crossroads between different parts of the graveyard. Steep stairs descend from the east and west sides of this room. Four bells of different sizes hang from leather cords affixed to the ceiling.

\end{DndReadAloud}

This room is set into the wall separating Neverdeath's Main Graveyard and Pauper's Graveyard. This room lets the cultists pass between the two graveyards without venturing aboveground.

Arriving cultists ring specific bells in a predetermined pattern, based on their rank, so cultists in the common room (area C14) can prepare an appropriate welcome. If the characters ring the bells without knowing their significance and patterns, the cultists in area C14 are alert to trouble.

\subsubsection{C14: Common Room}
\begin{DndReadAloud}{}
Tables and chairs in this crypt are arranged to create a meeting room or mess hall. Water drips from a pipe into a basin in the southeast corner beneath a detailed image of a staring eye gripped in a withered hand. Five robed cultists are in this room, with one bullying the rest.

\end{DndReadAloud}

A neutral evil \textbf{mage} and four neutral evil cult fanatics occupy this room. The mage, a sneering, human bully named Oxtu, insists the cult fanatics call him by his formal title of "Teeth of Vecna." In turn, he refers to them as "Memories of Vecna," their rank. Oxtu likes to describe violent methods of coercing secrets from people, and the cult fanatics hang on his words. Oxtu carries keys that unlock all the prisoner cells (areas C5, C11, C18, and C20).

The cultists are quick to fight intruders. The fanatics try to stay out of the way of Oxtu's spells, but Oxtu makes no effort to exclude them. The cultists all fight to the death.

\subparagraph{A Noisy Fight}
Loud noise here rouses the two mages dozing in area C16.

\subparagraph{Unholy Basin}
The cultists desecrated the basin by placing iconography of Vecna above it. Creatures that aren't devotees of Vecna that drink from the basin must succeed on a DC 17 Constitution saving throw or have the poisoned condition for 1 hour. Cultists who suspect a lack of true devotion from their compatriots challenge each other to drink from the basin to prove their faith. If the characters come in disguise as devotees of Vecna, the cultists demand that they prove their faith by taking a drink.

\subsubsection{C15: Kitchen}
\begin{DndReadAloud}{}
This small crypt has been converted into a kitchen. Dirty utensils soak in a large basin to the east. A stone coffin serves as a firepit; the coffin's lid has been repurposed as a table, which bears platters of dried fruit, nuts, and meat.

\end{DndReadAloud}

This kitchen remains empty except during mealtimes, and no one bothers to keep it clean. Removing the utensils from the basin reveals a wide drain.

\subsubsection{C16: Subleader Quarters}
\begin{DndReadAloud}{}
Bones in nooks along this wall were pushed aside to make room for folded robes and other personal effects. Four narrow cots lie against the north wall. A robed human and a robed elf each rest on a cot.

\end{DndReadAloud}

Two neutral evil mages, an elf man named Hannel and a human woman named Algra, rest here. They are surly, taciturn zealots who venerated Vecna in secret for decades before joining the cult. They love exercising their authority over junior cultists. Each wears a necklace of human teeth in honor of their titles within the cult hierarchy as "Teeth of Vecna." They each carry a key to the prisoner cells (areas C5, C11, C18, and C20).

The mages are quick to fight if they spot intruders, since they don't want the cult exposed. They are determined to vanquish intruders and prove their worth to the cult, even if it means fighting to the death.

\subparagraph{Secret Door}
An urn in a nook on the south wall rotates, sliding aside a wall panel that leads to a short tunnel between this room and area C25. A character who searches the room and succeeds on a DC 15 Wisdom (Perception) check finds the secret door and the means to open it. The mages know about the secret passage, but none of the cult fanatics do. At the far end of the secret passage is another secret door easily spotted and opened from inside the tunnel.

\subsubsection{C17: Library}
\begin{DndReadAloud}{}
The long, low shelves of this room are canted at irregular angles due to the uneven stones in the floor. The shelves are crammed with books, scrolls, and folios. Four robed cultists stand near the south shelves instructing two ghouls to tidy up the books.

A stone coffin rests against the north wall, its top carved to look like pages of an open book. Engraved on the book's pages is a name: Ayren Griffynstone.

\end{DndReadAloud}

Four chaotic evil cult fanatics are supervising two ghouls trying to reorganize the jumble of books the cultists brought to this tomb. The library belongs to Ayren Griffynstone, a human Neverwinter historian. The room hasn't fared well in the graveyard's upheavals, and the uneven floor makes this room \textit{difficult terrain}.

As Vecna is a god of knowledge as well as secrets, the cultists all contributed their personal libraries to this collection. Each of the cult fanatics has their own ideas about how this hodgepodge of eclectic works should be organized, so the ghouls labor under constant streams of conflicting directions.

Everyone here is on edge and grateful for the distraction of a fight. If the fight turns against the fanatics, one tries to escape through the secret door to fetch the demons in area C19.

\subparagraph{Secret Door}
One shelf swings back to reveal a secret passage to area C19. A character who searches the shelves and succeeds on a DC 15 Intelligence (Investigation) check finds the secret door and the means to open it. All the cultists know about this secret door. At the far end of the secret passage is another secret door easily spotted and opened from inside the tunnel.

\subparagraph{Treasure}
The collection of sinister books, many of which are duplicates, includes a few valuable tomes. A book describing the \textit{Eye of Vecna} and \textit{Hand of Vecna} is a masterpiece of writing and artistic illumination worth 450 gp. A book of nonsensical poetry titled Quite Good Verse has a gold-plated cover and is worth 200 gp. A book about Neverwinter's history contains a \textit{Spell Scroll} of \textit{Greater Invisibility} and a \textit{Spell Scroll} of \textit{Major Image} folded in its pages. The characters can find these treasures with 10 minutes of dedicated searching.

\subsubsection{C18: Vacant Cell}
This room has a padlock outside like the other prisoner cells, but the lock hangs open.

\begin{DndReadAloud}{}
This squalid cell contains nothing but a bucket and a small heap of filthy blankets.

\end{DndReadAloud}

The planar scholar Eldon Keyward occupied this cell for many miserable days. He was taken to the ritual cage in area C26, so his cell isn't locked.

\subparagraph{Eldon's Notebook}
Anyone searching the blankets finds Eldon's prize possession: a small notebook filled with his cramped writing about extraplanar intersections, planar conjunctions, and similar esoterica.

\subsubsection{C19: Demon Lair}
\begin{DndReadAloud}{}
Hundreds of names are etched into metal plates set into the walls of this room, many scratched over and unreadable. Two hulking, red-furred, apelike creatures stalk around the room.

\end{DndReadAloud}

The two barlguras here work as the cult's kidnappers. With little to do until the next kidnapping spree other than guard the imprisoned aristocrat in area C20, the demons spend their time scratching out the names on the memorial plates with their claws.

\subparagraph{Secret Door}
One of the nameplates pivots to reveal a secret passage to area C17. A character who searches the room and succeeds on a DC 17 Wisdom (Perception) check finds the secret door and the means to open it. The demons don't know about the secret door, but the cultists do. At the far end of the secret passage is another secret door easily spotted and opened from inside the tunnel.

\subsubsection{C20: Indrina's Cell}
This room's only obvious door is padlocked from the outside with a sturdy, new lock. As an action, a character with thieves' tools can use them to try to open the lock, doing so with a successful DC 18 Dexterity (Sleight of Hand) check. Each of the cult's four mages (in areas C14, C16, and C26) carries a key to it. Once the characters open the door, read the following:

\begin{DndReadAloud}{}
This crypt smells like a sewer. A woman sits on a mattress atop a low shelf, her once-fine clothing in tatters and a silk scarf wrapped around her face.

\end{DndReadAloud}

The prisoner is a human actor named Indrina Lamsensettle. Indrina's normally haughty demeanor has diminished in her imprisonment, though she's determined to make the cult pay once she escapes. Indrina dreams of returning to her estate, cleaning up, and dousing herself in perfumes. She doesn't know anything about Vecna or what the cult has in store for her; she believes that Lord Neverember is behind her imprisonment. Indrina has the game statistics of a lawful neutral \textbf{noble} but is unarmed and unarmored.

\subparagraph{Secret Door}
Part of the south wall swings aside when shoved. The short hall beyond leads to the latrine and smells even worse than Indrina's cell. A character who searches this room and succeeds on a DC 18 Wisdom (Perception) check finds the secret door. Indrina doesn't know the door is there.

\subparagraph{Indrina's Secret}
Indrina collected information from skilled genealogists and assembled proof that Lord Neverember isn't descended from Neverwinter's great hero, Lord Nasher Alagondar, as he claims. Indrina assumes Lord Neverember wants to silence her for what she's discovered.

If the characters admit to working for Lord Neverember, Indrina doesn't reveal what she knows. A character who asks why Indrina is here without revealing who hired the characters can attempt a DC 18 Charisma (Persuasion) check, with advantage if Indrina is given perfume or otherwise removed from the cell's offensive smell. On a success, Indrina reveals her secret knowledge.

Learning Indrina's discovery counts as a secret for the purposes of the Power of Secrets rules in this book's introduction. Lord Neverember casually dismisses Indrina's accusation if it's later brought to his attention, insisting that the woman can't prove anything.

\subsubsection{C21: Haunted Room}
The outside of this door bears a large "X" painted on it in red.

\begin{DndReadAloud}{}
Several urns lie shattered across the floor of this room amid heaps of ash and bone dust.

\end{DndReadAloud}

The cultists disturbed two wraiths bound to the urns while ransacking the room. The cultists quickly retreated and haven't been back. The wraiths emerge from the dust and ash when anyone opens the door, shrieking, "Vecnans, die!" They vent their rage on nearby creatures, preferring to attack cultists. The wraiths fight until destroyed.

\subsubsection{C22: Latrine}
\begin{DndReadAloud}{}
This filthy latrine is merely a deep pit with a few boards across it. The stink is overpowering.

\end{DndReadAloud}

\subparagraph{Secret Door}
The north wall of this disgusting room swings aside when a particular stone is pressed. A character who searches the room and succeeds on a DC 15 Intelligence (Investigation) check finds the secret door. All the cultists know about this secret door and open it from time to time to waft the smell into Indrina's cell (area C20).

\subsubsection{C23: Corridor}
\begin{DndReadAloud}{}
A shuffling cultist bearing a vacant expression moves through the corridor. Several narrow doorways lead off this long hall.

\end{DndReadAloud}

The Vecnan cult leader, Jerot Galgin, used a dreadful ritual to drain the cultist of her knowledge and vitality. This cultist, as well as other people whose secrets the cult have sacrificed to Vecna, has the game statistics of a \textbf{zombie}. The cultist is a Humanoid rather than Undead and isn't immune to poison damage or the poisoned condition. She is dressed like the other cult members and doesn't attack anyone dressed like cultists.

\subparagraph{Made an Example}
Raina Kairls was caught planning to betray the cult to Neverwinter's guards. Jerot first tested the sacrificial ritual on Raina and thinks she serves as a useful reminder of the price of betrayal.

\subsubsection{C24: Empty Crypts}
\begin{DndReadAloud}{}
This narrow doorway leads into a small, empty crypt.

\end{DndReadAloud}

These four rooms are empty.

\subsubsection{C25: Leader's Room}
\begin{DndReadAloud}{}
This large room has been furnished to resemble a cozy bedroom and study. Stooped over a desk and scribbling furiously on a parchment is a gaunt, robed human man. Standing next to the desk is a one-eyed, bipedal horror with spikes growing out of its back. Tapestries depicting feasting undead creatures hang on the wall.

\end{DndReadAloud}

The cult's leader is a neutral evil, human \textbf{necromancer wizard} (see appendix A) named Jerot Galgin. A loyal \textbf{nothic} assistant named Maszundrin never leaves Jerot's side. A devotee of Vecna, Maszundrin learned Common while lurking in the catacombs for decades and considers the cultists vital servants of the lich-god.

Jerot is an aristocrat who has lived his entire life in Neverwinter. He has built his deep faith in Vecna and vast necromantic knowledge over many years, right under the noses of his peers. He's engaging in this current research while his friends and family believe him to be on an extended trip to Waterdeep. Jerot considers his secret life as a cult leader, or the "Thought of Vecna," to be yet another way to honor his evil patron.

Jerot is refining the ritual occurring in the Sacrifice Gallery (area C26) and is too focused on his work to be distracted by combat elsewhere in the catacombs. He trusts his minions to handle any trouble. If intruders reach Jerot's personal chamber, he commands the \textbf{nothic} to enter melee while he fights from a distance, summoning Undead defenders if the nothic falls. Jerot fears exposure more than anything else and thus fights to the death.

\subparagraph{Jerot's Papers}
Jerot's notes on the ritual describe draining and sacrificing a victim's secrets and knowledge to Vecna. His notes illustrate the first test of the ritual, which used a disloyal cultist as the victim. For his future attempts, Jerot has chosen townspeople from Neverwinter whom he believes have particularly meaningful secrets. Their secrets are the cultists' best offerings to Vecna.

A character who examines Jerot's notes and succeeds on a DC 14 Intelligence (Investigation) check finds mention of magical phenomena called "Crevices of Dusk" that sometimes appear in Neverwinter. The notes indicate that these magical gateways connect to a plane populated by Undead, but it's clear Jerot doesn't know much more than that. His notes indicate his resolve to learn more after he finishes his current experiments in stealing and offering secrets to Vecna.

\subparagraph{Secret Door}
A tapestry depicting a feasting ghoul conceals a secret, sliding door leading to a short tunnel between this room and area C16. A character who searches behind the tapestry and succeeds on a DC 10 Wisdom (Perception) check finds the secret door. Once inside the secret passage, a character can easily spot and open the secret door at the opposite end.

\subsubsection{C26: Sacrifice Gallery}
\begin{DndReadAloud}{}
This enormous room features raised galleries at the east and west ends. Six chanting figures ring the east gallery, their hands raised toward a spherical cage hanging from the 30-foot-tall ceiling. The ritual's leader chants from the west balcony, surrounded by hunched, one-eyed creatures with knobby hides. A terrified elf struggles in the dangling cage.

\end{DndReadAloud}

This room's floor is 10 feet lower than the raised galleries.

The room is filled with cultists engaged in an extensive, hours-long ritual to sacrifice the elf Eldon Keyward's secrets to Vecna. The ritual leader is a tall, proud, neutral evil \textbf{mage} named Kendri Nex. The five nothics around her attack intruders on sight. Kendri uses her magic defensively, retreating to the room's floor if pressed. Kendri carries keys that unlock all the prisoner cells (areas C5, C11, C18, and C20), but if the characters haven't already rescued the other prisoners, they might not have the chance to do so, as the encounter likely ends with them being shunted through a planar rift.

The six neutral evil cult fanatics on the raised east gallery don't fight, since they fear interrupting the complicated ritual. They maintain their chanting and wild gesticulations.

\subparagraph{The Cage}
Eldon is a lawful good elf \textbf{priest} who follows Deneir, a god of writing and knowledge. He can't cast spells while he's in the cage. Eldon's cage hangs from a sturdy chain that ends 20 feet above the ground. The door on the cage's side is latched but not locked. Any character who can reach the cage can open its door as an action.

\subsubsection{Fight's End}
This fight ends immediately if the characters kill Kendri, silence any of the cult fanatics, or attempt to free Eldon. Energy from the interrupted ritual opens a latent planar rift that shunts Eldon and the characters into the Shadowfell:

\begin{DndReadAloud}{}
A riot of silvery-purple energy fills the room. You feel a sense of space tearing open—then you're falling, and everything goes dark.

\end{DndReadAloud}

\section{Escaping Evernight}
The cult's disrupted ritual thrusts the characters (along with Eldon) through a Crevice of Dusk, a gap between the Material Plane and the Shadowfell city of Evernight, a gloomy reflection of Neverwinter. The Crevice of Dusk closes immediately after shunting the characters.

The characters experience the following vision:

\begin{DndReadAloud}{}
Around the world and across the planes, you perceive innumerable cults of Vecna. They snatch away people and strip their secrets in rituals like the one you stopped. Behind them, the withered form of Vecna gathers the secrets like threads, adding them to a glowing sphere of hidden knowledge in some impossibly distant place. The vision fades into darkness, leaving only Vecna's glaring left eye.

\end{DndReadAloud}

\subsection{Link to Vecna}
The characters each gain a metaphysical link to Vecna, which follows the rules for \textit{blessings} presented in the \textit{Dungeon Master's Guide}. Vecna's Link is the result of feedback from the interrupted ritual. Vecna is unaware the characters—or anyone, for that matter—are linked with him, so the god has no reason to sever the tie. The link can manifest as subtly or as obviously as each player wishes, from the sensation of a specific smell when the character thinks of Vecna to a loud noise only they hear when the lich's name is uttered.

\begin{DndReadAloud}{}
\subparagraph{Vecna's Link}
You gain a special intuition for secrets. You have advantage on Wisdom (Insight) checks. In addition, you can use an action to cast \textit{See Invisibility} without expending a spell slot. Once you cast that spell in this way, you can't do so again until you finish a long rest.

\end{DndReadAloud}

When the characters acquire this link, remind them about the Power of Secrets rules. Allow them to spend any secrets they've gained so far as usual.

\subsection{Evernight}
Evernight is a forlorn metropolis in the Shadowfell. It has geography similar to Neverwinter's, but it presents as Neverwinter's dismal opposite. The sun never shines on Evernight, and ash-laden fog rises from lava flowing through the city in place of Neverwinter River, choking the city.

While Neverwinter is filled with living creatures trying to build a better future, Evernight is populated by Undead—primarily vampires and ghouls—who prey on each other and on travelers.

Evernight is a crossroads of trade in the Shadowfell and hosts numerous markets, including the lively Corpse Market. There, undead merchants trade in the bodies and blood of the dead—sometimes, the very recently dead.

\subsection{Evernight Awakening}
Since the characters were underground in Neverwinter's graveyard when shunted to Evernight, they're similarly underground in Evernight's graveyard.

Each character appears within an open coffin. The coffins are jumbled near each other in a large, 10-foot-deep grave pit.

Twelve ghouls prowl the area around the party's grave pit. Characters in the grave pit hear the hungry shouts and slavering of approaching ghouls, and you can further ramp up the tension by concealing the total number of ghouls until a character emerges from the pit to look around. The coffin-filled bottom of the grave pit is \textit{difficult terrain}. The muddy, sloping sides require a successful DC 10 Strength (Athletics) check to ascend, but moving down them doesn't require a check.

\subparagraph{Eldon's Latest Imprisonment}
In addition to the open coffins for each character, the grave pit contains another coffin wedged into the dirt and nailed shut. Eldon Keyward is confined within it.

Eldon grunts for aid and pounds at the wood. Unless a character uses an action to open Eldon's coffin, he kicks a side panel free and squirms out after 1 minute of effort. If the fight with the ghouls is still going on, Eldon helps as best he can.

Once the fight is over, Eldon shares the following information, either all at once or in fragments:

\begin{DndReadAloud}{}
"Very bad, but not surprising. I don't know if the cultists planned to send us here, but here we are. Neverwinter has cracks between our world and the Shadowfell—Crevices of Dusk, they're called. Planar travelers sometimes slip through.

"We're in a nasty city called Evernight. It's an evil echo of Neverwinter, populated by undead. We arrived in Evernight's graveyard because we left from Neverwinter's graveyard.

"There's no evidence of the crevice we came through. It doesn't surprise me that it's gone; stable crevices are rarer than spontaneous ones. But we need a stable one to get back to Neverwinter.

"We shouldn't spend a lot of time traipsing through an undead-infested metropolis, hoping to stumble across a gateway to Neverwinter. That feels unsafe. Maybe we should ask around inconspicuously. There's a market east of the Neverwinter Graveyard, which means it's likely there's a market east of Evernight's graveyard, too. But we should be careful."

\end{DndReadAloud}

The characters have little to go on other than Eldon's suggestion. Other explorations of Evernight are both dangerous and fruitless, so you should eventually steer the party back to the market even if they venture elsewhere.

The characters don't have other encounters as they make their way through Evernight's graveyard, though they hear howls and cries through the fog that let them know the graveyard isn't a safe place to linger.

\subparagraph{Eldon's Notebook}
If the characters found Eldon's notebook and return it to him, he's grateful and immediately starts adding notes about the Shadowfell and Evernight.

\subsection{The Corpse Market}
\begin{DndReadAloud}{}
The wall around Evernight's graveyard is riddled with gaps, and the ground near it is covered in fallen rubble, so leaving the graveyard is easy. Directly east of the graveyard, a large market stretches for blocks in every direction. Tattered canvas and shrouds separate the numerous stalls. Feeble moonlight and flickering torches illuminate the city, regardless of the time of day.

\end{DndReadAloud}

Everything in the market exhibits pale, subdued colors, but the atmosphere is lively as ghouls, skeletons, and vampires meander from stall to stall.

\subsubsection{The Market's Goods}
Goods for sale in the Corpse Market include ghoulish remains intended for consumption, such as fingers pickled in brine, jars of blood, and wrapped organs. Shops sell bouquets of dead flowers, frayed burial finery, jewelry displayed on severed hands, elegant canopic jars, and the like.

\subsubsection{Market Clientele}
The Corpse Market occasionally sees living visitors, though the characters are the only ones present now. If the characters don't attempt to hide or disguise themselves, they receive sidelong stares from merchants and customers. Everyone assumes the characters wouldn't be here unless they were under the protection of an influential figure in the city or were powerful travelers in their own right.

\subsubsection{Meeting Sangora}
Shortly after the characters enter the Corpse Market, they draw the attention of a vampire merchant named Sangora. When she sees the characters, she spreads her cloak wide and shows sharp fangs in her smile. She says:

\begin{DndReadAloud}{}
"I am Sangora, proprietor of Sangora Sanguinaries. You're not likely interested in a cup of warm blood, but perhaps you need something else? Something more ephemeral, like directions? Or knowledge? I've been in this city a long, long time."

\end{DndReadAloud}

Sangora is a centuries-old \textbf{vampire} with sunken eyes and a shock of long, white hair. She is inquisitive and happy to gossip.

If the characters don't seem inclined to speak with Sangora, Eldon blurts out a question about finding a Crevice of Dusk.

Sangora sells information at a higher profit than she sells blood, and she's full of useful tidbits about the city. She charges 20 gp for each question she answers, but she also accepts an answer to a probing question instead of payment (primarily about where the characters came from, how they got here, and what they're looking for in Evernight). Sangora isn't looking for a fight.

Sangora can share the following points:


\begin{description}
\item[Evernight.]
Sangora provides background about the city of Evernight, as presented above.
\item[Neverwinter.]
Sangora hasn't been to Neverwinter in more than a decade and remembers the place as a lawless ruin. At that time, the populace viewed Lord Neverember as a manipulative tyrant. Sangora knows that gateways called Crevices of Dusk occasionally appear in Evernight and allow passage to and from Neverwinter.
\item[Crevices of Dusk.]
Sangora explains that residents of both Evernight and Neverwinter dislike these portals being used indiscriminately. Those who know the location of stable Crevices of Dusk either guard or hide them.
\item[Stable Crevice of Dusk.]
Sangora tilts her head with a thoughtful look before revealing that she knows of a stable Crevice of Dusk in a tomb of one of Evernight's former living families, the Dolindars. The Dolindar tomb is in Evernight's graveyard, and Sangora provides true directions to it.
\item[Dolindars.]
Sangora explains that the Dolindar family was exiled to the Shadowfell for reasons they never shared. The family studied wizardry and knew much about planar gates despite being unable to use them due to a family curse. They were miserable and lonely in life.
\end{description}

\section{Dolindar Tomb}
The party's best next stop is the Dolindar tomb, in a part of Evernight's graveyard the characters haven't yet seen. When the characters are ready to explore the Dolindar tomb, Sangora can point the way.

\subsection{General Features}
The following features are common throughout the Dolindar tomb.

\subsubsection{Unending Isolation}
The tomb isn't precisely haunted, but the isolation the Dolindars felt living in a city of the dead suffuses their tomb. Creatures in the tomb or the portico outside it can't muster the will to support others and thus can't take the Help action.

\subsubsection{Stone Construction}
The tomb is made of old, durable stone.

\subsubsection{Lighting}
The tomb is dark. Area descriptions assume the characters have a light source or some other means of seeing in the dark.

\subsubsection{Doors}
The heavy doors throughout the tomb are made of stone and grind noisily when open and shut. No doors are locked except the puzzle door in area B5.

\subsubsection{Ceilings}
Ceilings are 10 feet high throughout the tomb.

\subsection{Dolindar Tomb Locations}
The following locations are keyed to map 1.2.

\subsubsection{B1: Portico}
\begin{DndReadAloud}{}
A roof supported by stone pillars extends from the Dolindar tomb into the weedy yard of uneven earth. The stone door to the tomb is engraved with the word "DOLINDAR" above it.

\end{DndReadAloud}

Four \textbf{vampire spawn} catch up to the party when the characters reach the door to the tomb. After the characters left her, Sangora detailed her conversation to her vampire spawn assistants. The spawn decided to make a meal of the characters, assuming that Sangora would never discover their perfidy.

The ground in the portico is swept and free of weeds, thanks to the efforts of the ghost Newmy (see area B2). The entrance to the tomb isn't locked. It opens onto a small room containing a steep spiral staircase leading 20 feet downward.

\subsubsection{B2: Hall}
\begin{DndReadAloud}{}
Upright slabs are set into the walls of this large burial chamber. Each slab is carved with the faded likeness of a different robed human above indecipherable writing. One of these slabs is blank with a piece of paper stuck to it.

\end{DndReadAloud}

The lesser-known Dolindars are interred here, buried in the alcoves behind the slabs. The images of the dead and their names have worn away. If the characters open these slabs, they find only dust and bones inside.

The piece of paper is on a slab in front of an open nook that never held a dead Dolindar. Instead, a ghost custodian named Newmy lives inside. She's affixed a piece of paper with "Newmy's Room" written on it in Common. If anyone approaches the slab, Newmy pops out, sputtering apologies.

Newmy is a lawful neutral \textbf{ghost} who was once a moon elf. She can cast \textit{Prestidigitation} at will. Newmy isn't quite five feet tall, and she has frizzy blue hair and pale skin shot through with blue veins. She's not interested in fighting, since fights make messes.

\subparagraph{Talking with Newmy}
Newmy would rather talk than fight. She shares the following points:


\begin{description}
\item[Crypt Cleaner.]
The Dolindars hired Newmy and paid her several lifetimes of wages up-front to keep their tomb clean shortly after the tomb was built—more years ago than Newmy can remember. When Newmy died, she returned as a ghost to continue her duties for as long as she was contracted.
\item[Lonely Family.]
The Dolindars were all lonely, and there weren't many of them. They didn't like Evernight but couldn't leave it for some reason.
\item[Puzzle Buttons.]
Newmy doesn't know anything about a rift or a Crevice of Dusk, but she remembers some "puzzle buttons" deeper in the tomb that she doesn't know how to work. Maybe they lead somewhere special.
\item[Restless Dead.]
Newmy doesn't go through the doors at the end of the hall anymore. The Dolindars buried down there "aren't resting right," and Newmy is afraid of them. She hasn't been that way in years and doesn't remember what's there.
\end{description}

\subparagraph{Newmy's Room}
Newmy rests in a burial nook large enough to hold the corpse of a Medium creature. The nook contains old rags and a decrepit broom. Newmy considers it her personal space and grumbles if anyone seems intent on disturbing it.

\subsubsection{B3: Lost Dolindars}
\begin{DndReadAloud}{}
Two stone coffins in this room have been broken open, littering the floor with rubble and dust. A creature with too many arms and spikes in place of hands taps at the room's walls.

\end{DndReadAloud}

The Dolindar siblings buried here were warped into two \textbf{lost sorrowsworn} (see appendix A). A visible sorrowsworn shrieks in rage and attacks anyone she can see. A second sorrowsworn is resting inside one of the broken coffins, initially out of sight but quick to join his sibling in a fight. The sorrowsworn fight until destroyed.

\subparagraph{Treasure}
The rubble includes the nameplates that once adorned the coffins. One reads, "Nolan Dolindar, Beloved Brother" and the other reads, "Evisha Dolindar, Beloved Sister." Each silver nameplate is worth 75 gp.

\subsubsection{B4: Vault}
\begin{DndReadAloud}{}
This vault contains six pedestals, each bearing treasure.

\end{DndReadAloud}

Carvings encircling the base of each pedestal read, "What good are treasures when home is denied?"

\subparagraph{Treasure}
The following Dolindar family treasures sit atop the pedestals:


\begin{itemize}
\item 2,200 gp in neatly stacked piles
\item A golden helmet worth 280 gp
\item A \textit{Driftglobe} with a map of the gate towns of the Outlands carved on it
\item A \textit{Necklace of Adaptation} bearing the inscription "My breath is yours, Kevetta—take it"
\item A book titled \textit{Out of the Endless Prison}, outlining methods of escaping the prison-plane of Carceri, worth 500 gp to a planar scholar (such as Eldon, who promises to buy the book once returned to Neverwinter)
\item A snow globe containing a miniature replica of the city of Neverwinter, worth 350 gp
\end{itemize}

\subsubsection{B5: Puzzle Door}
\begin{DndReadAloud}{}
The door in the east wall of this otherwise empty room bears the inscription "DOLINDAR" above "NO WORLD TO RETURN." Every letter of each word is set into the wall on a separate tile.

\end{DndReadAloud}

The door is sealed with a puzzle that requires pushing the right letter tiles in sequence. A pushed letter makes its whole word sink into the door with a click. The door unlocks when the correct combination is input. (Pushing a second letter in the same word doesn't do anything.) The correct combination, which causes the wall to slide aside, is to push the letters spelling ALONE in this order:

\textbf{A} in DOLINDAR

\textbf{L} in WORLD

\textbf{O} in TO

\textbf{N} in NO

\textbf{E} in RETURN

If the five words sink into the wall in the wrong order, or if the wrong letters are used to push them in, all five words reset with a wave of painful loneliness. Creatures in the room must make a DC 16 Charisma saving throw, taking 22 (4d10) psychic damage on a failed save or half as much damage on a successful one. If the characters have a hard time figuring out this puzzle, Eldon gives them hints. Eldon stays at a safe distance from the trap and doesn't take damage if it's triggered.

The first time a character pushes the tiles incorrectly, they see a small mechanism below the phrase "NO WORLD TO RETURN." This mechanism is a lock. As an action, a character with thieves' tools can use them to try to pick the lock, doing so with a successful DC 16 Dexterity (Sleight of Hand) check. Picking the lock has the same effect as solving the puzzle, granting access to area B6.

\subsubsection{B6: Matriarch's Chamber}
\begin{DndReadAloud}{}
On the far side of this room rests a stone coffin. Between the door and the coffin, the floor is studded with a sharp metal blades. A person-shaped figure with elongated arms lurks near the coffin.

\end{DndReadAloud}

Isolation warped the matriarch of the Dolindar family, Kevetta, into a \textbf{lonely sorrowsworn} (see appendix A). She remains near her coffin and uses her harpoon arms to attack anyone who enters the room. She fights until destroyed but doesn't pursue foes who flee.

\subparagraph{Blades}
The blades on the floor are \textit{difficult terrain}. When a creature moves into or through the blades on its turn, it takes 5 (2d4) slashing damage for every 5 feet it travels. A character can use an action to pull a lever behind the stone coffin to cause the blades to retract into the floor or to raise them again.

\subparagraph{In the Coffin}
Kevetta's coffin still bears her name, but not her body. Instead, the coffin's bottom is a roiling swirl of silvery-purple energy. Eldon confirms what the characters might guess: this is a Crevice of Dusk leading to Neverwinter.

Creatures and objects placed into Kevetta's coffin appear in a dusty, nondescript tomb in Neverwinter's Pauper's Graveyard. The name on the tomb is illegible due to advanced age, but the carved phrase "Home Again, To Rest Forever" is barely legible above the tomb's doorway.

\subsection{The Cult Scatters}
Though any remaining cultists believe they've successfully stolen Eldon's and the characters' secrets, it doesn't take them long to realize the significance of the characters' presence, which implies that important figures in Neverwinter know about the cult.

The cultists flee Neverwinter shortly after the ritual's abrupt end. By the time the characters return to the catacombs, the cultists are gone, along with any loose treasure. Other monsters such as the demons and nothics might still lurk in the graveyard, but they don't know anything about what the cult is doing or why the cultists departed so suddenly.

Before scattering, the cultists murder any kidnapped townspeople the characters didn't free from their cells. After all, the cult needs to keep its secrets.

\section{What's Next?}
Apart from a few loose ends, the characters' adventure is over when they return to Neverwinter. Grateful that the characters have ended the kidnapping threat, Lord Neverember rewards each character with a large house in Neverwinter. These houses can be adjacent to each other or spread around the city, as the characters prefer—Neverwinter contains a considerable number of vacant residences. Lord Neverember also pays for a small army of construction workers and renovators to restore the houses.

The rescued nobles and their families are grateful for the characters' efforts and give them 9,000 gp in total as a reward.

Although this chapter is concluded, each character retains their Vecna's Link. Unknown to the characters, they are about to become embroiled in the rivalry between Vecna and Kas.

If the characters aren't yet 11th level, more adventure awaits; there's no end of work in a frontier city such as Neverwinter. Alternatively, you can award the characters a milestone level-up and move ahead with the action of the next chapter.

\chapter{The Wizards Three}
Time has passed since the characters' experience in Neverdeath Graveyard. That ordeal might seem firmly in the past, but the characters receiving Vecna's Link is the harbinger of events none could predict. In this chapter, the characters become involved in three powerful wizards' desperate bid to stop the remaking of existence at the hands of the lich-god Vecna. The characters are soon led to believe that retrieving and reconstructing a legendary artifact is the only way to avoid a bleak future for the entire multiverse.

\section{Running This Chapter}
This chapter begins sometime after the characters' adventure in Neverwinter—days, months, or years, at your discretion. After you and the players determine what the characters have been doing following their Neverwinter escapades, the heroes are abruptly shunted to a mysterious sanctum in Sigil, the city at the center of the multiverse (see the "Surprise Development" section later in this chapter).

The characters learn a sobering fact from the renowned wizards Alustriel Silverhand, Mordenkainen, and Tasha: the lich-god Vecna is planning to remake the multiverse and emerge as its most powerful being, subjugating all existence. Mordenkainen, who is actually the vampire Kas in disguise, believes the characters are the only ones who can stop Vecna's plans using the fabled \textit{Rod of Seven Parts}, pieces of which are scattered throughout the multiverse.

Alustriel, Mordenkainen, and Tasha aid the characters in their task. This quest leads the characters deep into the Underdark in Faerûn, where the first piece of the rod is hidden in Web's Edge, a secret haven for cultists of the demon-god Lolth.

\subsection{Character Advancement}
The characters should be 11th level when this chapter begins. The characters gain a level after they retrieve the first piece of the \textit{Rod of Seven Parts} from Web's Edge in the Underdark.

\subsection{Power of Secrets}
The characters can learn two secrets in this chapter or later that are applicable to the rules in "The Power of Secrets" section in this book's introduction:


\begin{description}
\item[Gertrude's Secret.]
The characters can learn that Gertrude, the lone survivor of an ambushed caravan, and her friend Rockzanna knew about an imminent attack by Lolth's cultists but said nothing (see area W6a in Web's Edge).
\item[Mordenkainen's Secret.]
Mordenkainen is duping Alustriel, Tasha, and the characters and is actually Kas in disguise. (The characters aren't likely to learn this until later in the adventure.)
\end{description}

\subsection{First Rod Piece}
The Rod of Seven Parts of the \textit{Rod of Seven Parts} is in area W12 of Web's Edge later in this chapter. For more information about the rod and the spell this piece allows its wielder to cast, see this book's introduction.

\section{After Neverdeath}
It's up to you whether the events of this chapter happen after some time passes or immediately following the characters' experience in Neverdeath Graveyard. If you wish to emphasize urgency, this chapter's events happen a few days after the characters report their success to Lord Neverember in Neverwinter. If you'd rather emphasize the methodical nature of an eternal lich-god planning his dominance over the multiverse, months or years could pass before the characters arrive in Sigil.

\subsection{Vecna's Link}
Before resolving how the characters spend the time between their Neverwinter adventure and the events of this chapter, reiterate that each character retains their Vecna's Link. At various times throughout the interim period, the characters are reminded that they're metaphysically linked to Vecna. However you narrate this, make it clear that there's something supernatural about this connection. The link can't be removed, and attempts to research the link are unsuccessful. The characters don't know it yet, but they'll soon learn that the link is a harbinger of events to come.

\subsection{During the Interim}
If months or years pass before the characters meet the Wizards Three (see the "Surprise Development" section later in this chapter), ask the players what their characters did during this time. The party might have stayed together and gone on further adventures, or each character could have gone their own way.

If the players aren't sure or you'd like to offer them suggestions, consult the Interim Events table below. Choose your favorite result or roll randomly as you see fit.

\begin{DndTable}[header=Interim Events]{cX}

d6 & Event\\
1 & Determined to find answers about what happened to them, the characters worked as protection for gravediggers in Neverdeath Graveyard. They fought undead, grave robbers, and sometimes even devils who served evil cultists.\\
2 & Hoping to put their Neverwinter experiences behind them, the characters hopped on a ship to the Moonshae Isles and escorted traders around the Sword Coast area. They fought privateers and oceanic monsters in their bid for high adventure.\\
3 & The characters became obsessed with the cult of Vecna, researching it in prominent libraries and memorizing historical accounts of its activities. As a result, they know everything the Wizards Three tell the party about Vecna except his attempts to re-create the multiverse.\\
4 & Sobered by their experiences, the characters passed the time peacefully in the houses Lord Neverember granted them in Neverwinter. Adventure became less of a priority as they settled into domestic life.\\
5 & Obsessed with the concept of the multiverse, the characters studied the realms that make up existence, traveling between them to learn more. As a result, they know everything the Wizards Three tell the party about the planes where the pieces of the rod are located.\\
6 & Hoping to put their experiences behind them, the characters accepted a call for heroes to confront a murderous young red dragon named Rhagaermati in the wilds outside Myth Drannor. They are fighting the dragon when they experience the events described in the "Surprise Development" section.\\
\end{DndTable}

\section{Rise of the Lich-God}
Regardless of the characters' activities after their Neverwinter adventure, key plots elements of this adventure progress in the background. While the characters live their lives, a powerful wizard and former ruler of the city of Silverymoon, \textbf{Alustriel Silverhand}, detects a sinister wave of magic rippling through the multiverse.

Using her divining abilities, Alustriel eventually traces the magical pulse to Vecna. She learns that Vecna's activities span multiple realms. After a period of investigation, Alustriel discovers the sobering truth: using stolen secrets, Vecna and his cults have siphoned incredible amounts of power from individuals throughout the multiverse. Worse, Alustriel eventually stops detecting these activities, leading the archmage to believe that Vecna plans to unleash this power for a heinous purpose. Alustriel realizes this could unravel the entire multiverse, elevating Vecna and cowing all others to his will.

\subsection{The Archmages' Desperation}
Determined to oppose Vecna but unsure how the lich-god intends to unleash his amassed magic, Alustriel contacts her most powerful allies. Answering her call are the archmages Mordenkainen and a version of Tasha from Oerth, though Alustriel and Tasha don't yet know that Mordenkainen is Kas the Destroyer in disguise. These Wizards Three retreat to a sanctum Alustriel keeps in Sigil. Using their combined magic, the archmages weave a \textit{Wish} spell in hopes of sabotaging Vecna's accumulated power and defusing his ritual.

Instead of any expected effect, the \textit{Wish} spell shunts the characters to Alustriel Silverhand's Sigil sanctum, as explained in the "Surprise Development" section later in this chapter. With time of the essence and the archmages weakened, Mordenkainen suggests a desperate contingency plan. The characters could use the fabled \textit{Rod of Seven Parts} to stop Vecna. The rod's seven pieces are scattered throughout the multiverse, but Mordenkainen knows where the first piece is located.

\section{The Sigil Sanctum}
Hundreds of years ago, Lady Alustriel of Silverymoon needed a secure place to conduct sensitive magic work and research. After discovering a portal from Silverymoon to Sigil, Alustriel traveled to Sigil. There she spent decades creating a private sanctum around the portal.

Alustriel magically fused the portal in her sanctum with a \textit{Well of Many Worlds}, destroying the magic item in the process. As a result, the portal now leads to numerous locations throughout the multiverse. Further, the portal reacts to magic items and artifacts, allowing those who carry such objects to step through the portal to the exact places they wish to travel. Alustriel carefully maintains the sanctum's secrecy, bringing guests here only after first meeting them in a neutral location. So far, either her portal's strange teleportation properties have gone unnoticed by Sigil's Lady of Pain and her agents, or the Lady of Pain is allowing the portal to function this way for reasons of her own.

\subsection{Magic in the Sanctum}
Planar magic functions differently in Sigil. For instance, the only way to reach Sigil is through portals; teleporting into or out of the city doesn't work. Since the characters might spend significant time in Sigil, take note of how the following types of magic work in the city:


\begin{description}
\item[Banishment.]
Effects that banish a target from Sigil treat the target as if Sigil were its home plane.
\item[Extradimensional Space.]
Extradimensional spaces, demiplanes, and pocket dimensions—such as those created by a \textit{Bag of Holding} or the \textit{Rope Trick} spell—function within Sigil, but those spaces follow these restrictions as if they were part of the city.
\item[Planar Travel.]
Effects that allow interplanar travel, such as the \textit{Astral Projection} and \textit{Plane Shift} spells, fail if used to try to enter or leave Sigil, with one exception (see "Teleportation Circles" below).
\item[Summoning.]
Spells, magic items, and effects that summon creatures or objects from other planes, such as a \textit{Ring of Djinni Summoning}, instead summon targets from within Sigil if possible or otherwise fail. Effects that summon a specific target from outside Sigil, such as the \textit{Drawmij's Instant Summons} and \textit{Leomund's Secret Chest} spells, automatically fail.
\item[Teleportation.]
Attempts to teleport into or out of Sigil fail, but such magic functions normally when teleporting within the city.
\item[Teleportation Circles.]
Permanent teleportation circles exist within Sigil, but attempts to create new ones fail. If the Lady of Pain permits, they can be used to enter the city via the \textit{Plane Shift} spell but not to leave.
\end{description}

\subsection{General Features}
The features of Alustriel's sanctum are described in the following sections.

\subsubsection{Ceilings}
The ceilings on the first floor of the sanctum are 20 feet high. The ceilings on the upper floor are 10 feet high.

\subsubsection{Doors}
No doors separate areas S1–S3 of the sanctum, though open archways set off each area. The doors in areas S5–S7 of the sanctum are unlocked. The door to area S8, Mordenkainen's room, is sealed with an \textit{Arcane Lock} spell that requires a successful DC 30 Dexterity (Sleight of Hand) check using thieves' tools to open.

\subsubsection{Lighting}
\textit{Continual Flame} spells cast on sconces bathe most areas in bright light. Only area S8 is unlit.

\subsubsection{Safe from Prying Eyes}
Everyone in the sanctum is under the effect of a \textit{Nondetection} spell while they remain inside the structure.

\subsubsection{Walls}
Alustriel magically reinforced the walls of the sanctum, which are made of a combination of brick and layers of cement. Only exceedingly powerful destructive magic could harm the sanctum's structural integrity.

\subsection{Sanctum Locations}
The following locations are keyed to map 2.1.

\subsubsection{S1: Library}
\begin{DndReadAloud}{}
Gleaming, floor-to-ceiling bookshelves filled with leather-bound tomes wrap around the sanctum's west end, while marble pillars hold up vaulted ceilings to the east. Spiral staircases in the corners lead upward.

\end{DndReadAloud}

This section of the sanctum holds Alustriel's collection of rare books. The books aren't magical, but the characters can use them to research any esoteric topic.

If a character researches a specific topic, roll a d4. On a roll of 1 or 2, the character finds information about the topic. At your discretion, this might grant the character advantage for the rest of the day on any ability checks made to recall or use the information they've learned.

\subsubsection{S2: Parlor}
\begin{DndReadAloud}{}
Plush, antique furniture, including sofas and footstools, is arranged here. To the north is a grand staircase. To the south, a garden of ferns and flowering plants grows around a raised marble dais.

\end{DndReadAloud}

When Alustriel invites dignitaries into her sanctum, she holds business and political discussions in this parlor. The grand staircase ascends to area S4.

\subparagraph{Dais}
The dais in the center of the garden is the site of the portal Alustriel has tweaked to lead to numerous multiplanar locations.

\subparagraph{Garden}
The greenery is the result of Alustriel's dabbling in magical gardening. The plants here need no tending.

\subsubsection{S3: Workspace}
\begin{DndReadAloud}{}
A desk covered in paperwork and books sits in the middle of this work area. Nearby are curio cabinets filled with magical trinkets—everything from necklaces and brooches to sparkling stones and shoes.

\end{DndReadAloud}

Characters transported to the sanctum by the failed \textit{Wish} spell appear here (see the "Surprise Development" section later in this chapter).

\subparagraph{Curio Cabinets}
Alustriel keeps an impressive collection of magic items here. If the characters ask, Alustriel allows each character to borrow their choice of any single magic item of rare or lower rarity from the \textit{Dungeon Master's Guide}. Later in this chapter, Kas, disguised as Mordenkainen, retrieves the \textit{Chime of Exile} from one of these cabinets.

\subsubsection{S4: Lounge}
\begin{DndReadAloud}{}
Roaring fireplaces are built into the east and west walls of this lounge. Comfortable couches and chairs are artfully arranged before a grand staircase that leads downward.

\end{DndReadAloud}

Alustriel and her friends socialize, dine, and relax in this lounge. Alustriel is known for her sumptuous dinner parties. Discreet drawers in the bottoms of the couches store fine china and cutlery for when Alustriel transforms the center table into a dining table. Additional storage is in the northwest corner.

When the characters are shunted to the sanctum, Alustriel offers them use of the lounge.

\subsubsection{S5–S8: Bedrooms}
These small but comfortable bedrooms are where Alustriel and her guests stay when they need personal quarters, whether it's to sleep or be alone.

Area S5 is the bedroom of Alustriel and her wife, Malaina van Talstiv. A retired adventurer, Malaina is a neutral good human who uses the \textbf{assassin} stat block. Malaina is meeting with associates to thwart a scheme against Waterdeep when the characters arrive. When the characters return, Malaina speaks with them (see the "Next Steps" section at the end of this chapter).

Area S6 is Tasha's room. Area S7 is an empty room, which Alustriel offers to the characters for their use. Area S8 is Mordenkainen's bedroom, which contains clues hinting at subterfuge (see the "Evidence of Deception" section later in this chapter).

\section{The Wizards Three}
Three legendary archmages serve as the characters' allies throughout the rest of this adventure, though one of these individuals is an impostor. (See the "Kas the Destroyer" section.) More information about each spellcaster follows.

\subsection{Alustriel Silverhand}
More than seven centuries old, Lady \textbf{Alustriel Silverhand} is beloved throughout Faerûn for her kindliness and for using her incredible spellcasting power to protect the innocent. Alustriel is an immortal daughter of Mystra, a god of magic, and is a chaotic good, human wizard. Alustriel served as High Mage of Silverymoon, battled demons on the Outer Planes, and prevented the lich Larloch from ascending to godhood. She has seen the rise and fall of evil powers, though Vecna's plot troubles her greatly. More about Alustriel, including her stat block, appears in appendix B.

\subsection{Mordenkainen}
Renowned for his bravery but not his judgment, Mordenkainen is a powerful spellcaster from Oerth. The chaotic neutral, human wizard led a council of famous archmages known as the Circle of Eight. He later became trapped in the dread realm of Barovia, where he lost his spellbook and staff as well as his grip on reality and wandered adrift for some time. Alustriel trusts Mordenkainen and respects his gumption, even though he occasionally embarks on tasks that outstrip his capabilities. The real Mordenkainen doesn't appear in this adventure, though the characters likely believe they're interacting with him.

\subsection{Tasha}
Tasha is a chaotic neutral archmage from Oerth, a renowned demonologist, and the adopted daughter of Baba Yaga. Tasha is widely considered the multiverse's foremost authority on the Abyss, having authored the fabled \textit{Demonomicon of Iggwilv}. Although Tasha's motives are fickle, Alustriel has found Tasha an unparalleled ally when their interests align, as Tasha can access and influence places where even the daughter of Mystra dares not tread. The version of Tasha who answers Alustriel's summons is from the past, before Tasha became Iggwilv the Witch Queen. More about Tasha, including her stat block, can be found in appendix B.

\section{Kas the Destroyer}
The vampire warlord Kas maintains his masquerade as Mordenkainen, likely fooling the characters until the events of chapter 9. The blood feud between Kas and Vecna stretches back centuries to a time when Kas served as Vecna's first lieutenant. More about Kas, including his stat block, appears in appendix B.

\subsection{Manipulation in Motion}
After the Dark Powers released him, Kas tracked down a cell of Vecna's cultists, who told the vampire more about the lich-god's planned ritual. The vampire used the \textit{Crown of Lies} to intercept a message from Alustriel to Mordenkainen, then used the crown to mimic Mordenkainen. Kas then met with Alustriel and gained her confidence.

While posing as Mordenkainen, Kas sabotages the archmages' \textit{Wish} spell. As a consequence, the spell shunts the characters to Alustriel's sanctum. Suspecting that the characters are somehow linked to Vecna, Kas urges them to retrieve the pieces of the \textit{Rod of Seven Parts}. As Mordenkainen, Kas claims that the rod is the only way to weaken Vecna to the point where the characters might thwart his ritual and banish the lich-god to Oerth.

In truth, Kas intends to steal the \textit{Rod of Seven Parts} once he has all the pieces. He plans to whisk the artifact to the plane of Pandemonium, use it to free the demon known as Miska the Wolf-Spider, and kill Vecna. Kas doesn't yet know the site of Vecna's ritual is also on Pandemonium, the same plane where Miska is imprisoned, but the vampire finds out later in this adventure.

\subsubsection{Kas in the Sanctum}
While disguised as Mordenkainen, Kas never lets his guard down or removes the \textit{Crown of Lies}. Occasionally during this adventure, Mordenkainen is absent from the sanctum. During these interludes, the vampire warlord travels to other realms and eventually deduces where Vecna plans to enact the multiverse-unraveling ritual. This sets into motion the betrayal that takes place in chapter 9.

\subsubsection{Evidence of Deception}
While in the sanctum, Kas is careful to leave no direct evidence that he's not Mordenkainen. In Mordenkainen's quarters (area S8), Kas leaves spellbooks and collections of notes spread across the desk to reinforce his fake identity.

However, a character who searches Mordenkainen's quarters and succeeds on a DC 30 Intelligence (Investigation) check notices that none of the famed spells Mordenkainen authored are mentioned in these spellbooks or notes. The jumbled notes are written awkwardly, while Mordenkainen is known for his academic, methodical style.

Any thorough search of Mordenkainen's quarters yields a black \textit{Bag of Holding} tucked underneath the mattress. Inside the bag is a mundane silver chain bearing a sword-shaped pendant inlaid with black diamonds, worth 5,000 gp total. Also inside is a magically preserved, leather-bound journal.

The journal's contents are written in a cipher. A character who spends at least 1 hour examining the text can make a successful DC 30 Intelligence (Investigation) check, decoding it on a success. The journal tells the story of two people called "K." and "V." who once crafted murderous battle plans together, but then parted bitterly and became enemies. A successful DC 20 Intelligence (History) check reveals that this story refers to Kas and Vecna.

\subparagraph{Mordenkainen's Secret}
If the characters find out at any point in this adventure that Mordenkainen is Kas in disguise, it counts as a secret for the purposes of the Power of Secrets rules in this book's introduction.

\subsubsection{If They Find Out?}
Although it's unlikely, the characters might become suspicious of Mordenkainen and question him. The characters could investigate the wizard's activities outside the Sigil sanctum, or they might find the strange personal items in Mordenkainen's \textit{Bag of Holding}. If the characters put sufficient effort into their investigations, consider allowing them to discover before chapter 9 that Mordenkainen is an impostor.

\subparagraph{Kas Defeats the Characters}
If the characters confront Kas before chapter 9, allow the battle to play out using the stat block for Kas in appendix B. If Kas defeats the characters, he kills them and absconds to Pandemonium. Alustriel has the clerics of Silverymoon bring the characters back to life, apologetic for not discovering Kas's ruse and angry she was duped. Alustriel suspects Kas knows where Vecna weaves his ritual and encourages the characters to follow the vampire to Pandemonium when they can.

In this case, Kas lacks the completed \textit{Rod of Seven Parts}. Adjust chapter 10 to reflect that Miska never begins to escape his prison. The characters' priority remains stopping Vecna, but if they don't permanently destroy Kas, at some point the vampire hunts down the characters. Once the characters eventually defeat Kas, they learn where Vecna weaves his ritual and can proceed to chapter 11.

\subparagraph{The Characters Defeat Kas}
If the characters defeat Kas, allow them to decide how to proceed. They might still want to reassemble the \textit{Rod of Seven Parts} to use its power against the lich-god. However, to reinforce the story's sense of urgency, feel free to skip the characters ahead to chapter 10. If you do, award the characters the appropriate number of milestone level-ups so they have a fair chance of defeating the threats in the adventure's final chapters.

\section{Surprise Development}
When the failed \textit{Wish} spell transports the characters to Alustriel's sanctum, read or paraphrase the following:

\begin{DndReadAloud}{}
With a flash of multicolored light, the world winks out of existence. Nothingness envelops your senses, though you feel the thrum of magic rushing like an electric spark through your veins.

When you regain focus, you stand in a plush, candlelit parlor. A stately woman in blue robes leans toward you, her brow furrowed in confusion and concern. In the blurry background stand a frowning woman in a flowing, black dress and a man tugging in confusion at his high, crimson collar.

\end{DndReadAloud}

The characters appear in area S3 of the Sigil sanctum. The woman leaning toward the characters is Alustriel, and the two individuals behind her are Tasha and Mordenkainen. "You can't be the answer to our \textit{Wish}," Tasha sneers, while Alustriel shushes her and ushers the characters toward the sanctum's comfortable lounge furniture.

\subsection{Some Answers}
At this point, Alustriel explains where the characters are and how they arrived, though she doesn't understand why they teleported to the sanctum as a result of the \textit{Wish} spell. She then introduces herself and her wizardly companions.

As the characters try to figure out what's going on, Alustriel asks whether they have any ties to Vecna. If the characters mention Vecna's Link, Alustriel reasons that the characters' fate must somehow be interwoven with Vecna's. She surmises that since the \textit{Wish} was unable to fulfill its parameters, it instead found beings tied to Vecna and brought them to Sigil. She then explains everything she knows as described in "Rise of the Lich-God" earlier in this chapter.

\subsubsection{Aftermath of a Wish}
As the characters interact with Alustriel, Tasha stands near them, aloof and occasionally interjecting acerbic comments. Mordenkainen remains silent, focusing intensely on everything the characters say.

After a few minutes of conversation, the characters notice that the Wizards Three seem tired and weakened. Any character proficient in the Medicine skill understands that the wizards are extremely fatigued. Any character proficient in the Arcana skill realizes that the wizards expended a great deal of arcane energy casting the \textit{Wish} spell. (Mordenkainen is faking his fatigue, but in the context of Alustriel and Tasha's real physical weakness, it seems legitimate.)

\subsection{A New Plan}
After the characters absorb their situation, Alustriel expresses concern that the wizards can't stop Vecna. Mordenkainen says, "I have an alternative plan."

Mordenkainen strides over to a cabinet full of magic items and picks up a silver chime. He then approaches the characters, displays the chime, and says the following:

\begin{DndReadAloud}{}
"I worried that perhaps our \textit{Wish} wouldn't stop one as powerful as Vecna. Alustriel, please forgive me, but I did plan a contingency.

"Vecna is too powerful to stop directly at this point. However, no power is absolute. I suspect Vecna will be vulnerable while he weaves his ritual. If you confront and weaken him, this \textit{Chime of Exile} can banish him to Oerth. He would no longer be able to affect multiple realms at once. His plan to remake the multiverse would fail."

\end{DndReadAloud}

Mordenkainen's eyes gleam with excitement as he continues:

\begin{DndReadAloud}{}
"There is an artifact that is, I believe, the only way to weaken Vecna so the chime can be used against him. It's called the \textit{Rod of Seven Parts}. Unfortunately, its component pieces are scattered across existence.

"While my esteemed colleagues researched the magic needed to cast a \textit{Wish} spell powerful enough to stop Vecna, I looked into finding the rod's pieces should we need them. The first piece is deep in the Underdark, somewhere inside a hidden safe house for worshipers of the demon-god Lolth. Web's Edge, it's called.

"Find that first piece of the rod, my friends. The portal on the dais will take you to it. Once you locate the first piece, the magic inherent in the artifact will point you to the second piece, and then to the third, and so on. This sanctum's portal reacts to powerful magic such as this artifact, so as you find each piece, this portal will lead to the next. No need to use another key; the most recent piece of the rod will suffice.

"There's no time to waste. Once you've reassembled the rod, we can all confront Vecna and stop his heinous plan. You'll be the heroes of the multiverse! How does that sound?"

\end{DndReadAloud}

Although taken aback at how much work Mordenkainen has done on this contingency, Alustriel and Tasha admit that the plan Mordenkainen has laid out is the best hope to stop Vecna. Alustriel compliments Mordenkainen on his foresight.

The wizards answer the characters' questions to the best of their ability. Like Mordenkainen, Alustriel and Tasha urge the characters to make haste, though they help the characters prepare for their journey any way they can. Mordenkainen, Alustriel, and Tasha explain that while the characters reassemble the rod, the archmages must continue to research Vecna's plans and potential weaknesses.

Before the characters depart for Web's Edge, Alustriel relays any information they haven't yet learned, including the details found in "Magic in the Sanctum" earlier in this chapter.

\subsubsection{A Significant Lie}
In the guise of Mordenkainen, Kas has lied to the characters about the \textit{Rod of Seven Parts}. While the rod is useful against Vecna, Kas simply wants the characters to retrieve the rod so he can use it to free Miska.

Kas doesn't know that Vecna is weakened to pre-god status while weaving the Ritual of Remaking. The characters will learn this when they confront Vecna in chapter 11.

\subsection{Sending Vecna Back}
The \textit{Chime of Exile} is Alustriel's property and ultimately can send Vecna back to Oerth, interrupting his ritual. The interruption would dissipate the secret-based magic Vecna is weaving, significantly setting back his plot and leaving the lich-god in a weakened state. If the characters wish to take the chime with them, Alustriel assents.

\subsubsection{Chime of Exile}
\textit{Wondrous Item, Very Rare}

This silver chime is engraved with delicate magic sigils. While holding the chime, you can use an action to cast the \textit{Banishment} spell (spell save DC 20). If the target of the spell has 50 hit points or fewer, it automatically fails its saving throw. Once the chime has been used to cast the spell, it can't be used this way again until the next dawn.

\section{Toward Web's Edge}
After stepping through the portal, the characters emerge deep in the Underdark in Faerûn. The rod piece is hidden in Web's Edge, a nearby safe house for agents of the demon-god Lolth.

The characters initially emerge through a door-shaped fissure in one wall of a claustrophobic, seemingly abandoned corridor outside the hidden entrance to Web's Edge. There is no light source, and the invisible portal back to Sigil remains open.

\subsection{Web's Edge Entrance}
The entrance to Web's Edge is hidden mere steps from where the characters arrive in the Underdark. The 10-foot-wide doorway is closed and cloaked with a permanent invisibility effect. Those who can perceive invisible objects see a nondescript iron door. The door is sealed with an \textit{Arcane Lock} spell and requires a successful DC 22 Dexterity (Sleight of Hand) check using thieves' tools to open.

Visible to characters who can see invisible objects are eight symbols subtly carved into the rock above the doorway. Tiny Undercommon characters are embedded in each symbol. When touched in the order that spells out "Web's Edge," the doorway opens.

\section{Web's Edge}
Web's Edge is a hidden meeting place for agents of Lolth who operate throughout the Underdark beneath the Sword Coast, pursuing missions for the glory of their demon-god. These agents use Web's Edge to plot infiltration missions targeting virtuous clerical orders; to sabotage efforts to quash Lolth's worship; to orchestrate the upheaval of good-aligned governments; and to plan large-scale conversion efforts in Underdark cities, including Blingdenstone, Gracklstugh, and Mithral Hall.

Agents might spend a few hours in Web's Edge, attending tactical meetings before dispersing. Others stay for 14 days or more while they await the arrival of colleagues from all reaches of the Underdark. A few spend months at a time in Web's Edge, using its barrack to lie low while authorities elsewhere hunt for them.

Owing to the highly sensitive and wicked work that takes place here, Web's Edge is a carefully guarded secret. Only the highest-ranking covert operatives and members of Lolth's clergy know about it, though powerful demons and devils allied with the Spider Queen are also aware of Web's Edge. A few cultists live in Web's Edge to maintain its facilities and serve as support staff to the operatives. These cultists are confined to the facility for life—a sacrifice they make willingly to show their devotion to the Spider Queen.

\subsection{General Features}
Recurring features of Web's Edge are described in the following sections.

\subsubsection{Agents and Cultists}
No matter why they're in Web's Edge, all Lolth worshipers here are fanatics devoted to the demon-god. In a battle against the characters, these fanatics always fight to the death.

\subsubsection{Ceilings}
The ceilings in most areas of Web's Edge are 20 feet high, with stalactites occasionally jutting down a few feet. The ceilings in the fodder chamber (area W9) and Sacred Web Hall (area W12) are 50 feet tall and relatively smooth.

\subsubsection{Doors}
The doors in Web's Edge are made of iron. The doors to areas W3, W4, and W5 are sealed with \textit{Arcane Lock} spells, requiring a successful DC 20 Dexterity (Sleight of Hand) check using thieves' tools to unlock. Additionally, areas W6a–W6d are locked and each require a successful DC 15 Dexterity (Sleight of Hand) check using thieves' tools to unlock. The \textbf{nalfeshnee} in area W5 carries keys to all the cells in area W6. All other doors in Web's Edge are unlocked.

\subsubsection{Lighting}
There are no natural or magical light sources in the complex. Area descriptions assume the characters have a light source or some other means of seeing in the dark. Light sources bring the attention of the complex's denizens, although the characters might not immediately be recognized as intruders (see "Infiltrating Web's Edge" below).

\subsubsection{Safe from Prying Eyes}
Everyone in Web's Edge is under the effect of a \textit{Nondetection} spell while inside the structure.

\subsubsection{Walls}
The walls in Web's Edge are rocky, jagged, and uneven. Decades ago, Lolth's faithful enlarged the chamber that contains Sacred Web Hall (area W12) using the \textit{Stone Shape} spell. As such, the walls in area W12 are smooth.

\subsection{Infiltrating Web's Edge}
Web's Edge is a dangerous place. If the characters burst into the complex without a plan, they're unlikely to survive the experience, let alone retrieve the first piece of the \textit{Rod of Seven Parts}. See the advice below for handling strategies the characters might adopt.

\subsubsection{Being Sneaky}
If the characters are careful, they might be able to sneak through Web's Edge and steal the rod piece without anyone catching them.

If any devotees in Web's Edge spot the characters infiltrating the safe house, they scream for help to any denizen within earshot and attack unless otherwise noted.

\subsubsection{Impersonating Worshipers}
Masquerading as devotees of Lolth could be an effective strategy. Creatures of all backgrounds and origins use Web's Edge as a meeting place, a safe house, and a place of worship. Web's Edge has only a handful of permanent attendants, none of whom could possibly be familiar with every Lolth agent in the Underdark.

If the characters want to blend in with the Lolth worshipers, any disguises they use (magical or otherwise) should include spider-shaped paraphernalia. A successful DC 14 Wisdom (Religion) or Intelligence (History) check reveals that "The Spider Queen smiles on you" is a typical religious greeting that those in Web's Edge expect to hear from fellow devotees. Any creature in Web's Edge who sees a character wearing something spider-shaped or hears the character speak this greeting has disadvantage on ability checks to see through the characters' disguises.

Devotees in Web's Edge who identify the characters as impostors scream for help and attack.

\subsection{Web's Edge Locations}
The following locations are keyed to map 2.2.

\subsubsection{W1: False Front}
\begin{DndReadAloud}{}
The remains of mold-covered broken wagons and barrels languish in this large foyer. Humanoid skeletal remains lie strewn about, their swords and armor bent and rusted. Something gleams next to the bodies propped along the northern wall. Double doors on the east, northeast, and southeast walls are rusted shut, and a semicircular chamber that opens in the southwest wall contains a shrine.

\end{DndReadAloud}

To convince unwelcome visitors that this cave is unused, the facility's cultists created a false front in this foyer. The skeletal remains and rusted gear belong to five long-dead adventurers whom the characters can identify as two dwarves, an elf, a gnome, and a tiefling.

\subparagraph{Traps}
The gleams in the northern part of the chamber are magical traps placed on fist-size, fake rubies tucked conspicuously next to two skeletons. The fake rubies are fixed to the floor and can't be removed. Any character who touches one of the rubies must make a DC 20 Dexterity saving throw, taking 54 (12d8) lightning damage on a failed save or half as much damage on a successful one.

The traps can't be disabled, but a character who comes within 5 feet of one of the rubies and examines it can make a DC 18 Intelligence (Investigation) check. On a success, the character deduces that the ruby is fake, worthless, and magically trapped to release a violent electric shock.

\subsubsection{W2: Lolth Shrine}
\begin{DndReadAloud}{}
This semicircular chamber forms a natural alcove. Shelves carved into the wall from floor to ceiling hold small, repulsive items, including bloody baubles, shriveled fingers, and idols carved from bone. Two figures crouch before the shelves.

\end{DndReadAloud}

Two agents of Lolth have come to pray at this shrine. They are Makubli Khee, a chaotic evil, hobgoblin \textbf{assassin}, and Torkner Ironteeth, a chaotic evil, duergar \textbf{mage}. Because they're engrossed in their prayers, Makubli and Torkner don't notice the characters unless a member of the party deliberately hails them or the characters have a light source.

The items on the shelves are nonmagical. Any character who looks at the ceiling and succeeds on a DC 15 Wisdom (Perception) check spots a small, spider-shaped carving in the rock—a hint of this shrine's purpose.

\subparagraph{Makubli and Torkner}
Both Makubli and Torkner operate along the well-traveled route between Menzoberranzan and Blingdenstone, sabotaging supply caravans and generally causing mayhem. They usually kill or convert victims who survive their initial attacks, although Makubli recently turned over custody of a cyclops to the Web's Edge prison (see area W6a). Makubli and Torkner dislike each other intensely. Neither knows many other covert Lolth operatives, so clever characters could masquerade as allies (see "Impersonating Worshipers" earlier in this chapter).

The cultists who maintain Web's Edge are friendly with both agents, making Makubli and Torkner good candidates for the characters to impersonate. If either agent suspects that the characters aren't fellow Lolth operatives, they attack.

\subsubsection{W3: Summoning Chamber}
\begin{DndReadAloud}{}
An enormous pentagram drawn in chalk, with stubby unlit candles placed at each of its five points, covers the floor of this open chamber. A hooded figure hunches over a cluttered table in the southwest corner, mumbling profane phrases.

\end{DndReadAloud}

Cultists and agents gather here to commune with Lolth's Abyssal servants and occasionally summon Fiends to assist with missions. The figure at the table is Grottenelle Stonecutter, a chaotic evil, svirfneblin \textbf{mage} who serves as the facility's high summoner.

\subparagraph{Summoning in Progress}
Grottenelle is in the middle of a summoning ritual. Have Grottenelle and the characters roll initiative. If the characters attack or otherwise interrupt Grottenelle before the start of her first turn, the summoning fails. Otherwise, a \textbf{glabrezu} appears in the center of the pentagram at the start of Grottenelle's turn. Grottenelle attacks intruders and commands the glabrezu to do the same if the summoning is successful.

Grottenelle needs the glabrezu to assist in the operation being planned in area W7. At your discretion, the noise from a battle might alert the agents in that area, prompting them to investigate.

\subparagraph{Treasure}
Grottenelle's table holds a \textit{Spell Scroll} of \textit{Circle of Death}, an ornate \textit{+2 Dagger}, and an assortment of unguents and oils worth 500 gp.

\subsubsection{W4: Corridor}
\begin{DndReadAloud}{}
This corridor is empty. Heavy iron double doors bookend the western and eastern access points.

\end{DndReadAloud}

Additionally, if the characters haven't yet explored area W3, read the following aloud:

\begin{DndReadAloud}{}
From the north, you hear a faint voice mumbling profane phrases.

\end{DndReadAloud}

The cultists keep the door to the fodder chamber (area W9) closed in case the food intended for Ker-arach in the Sacred Web Hall (area W12) tries to escape.

\subsubsection{W5: Guardian Chamber}
\begin{DndReadAloud}{}
An enormous, winged biped with a boar's head paces in front of a barred prison door to the southwest. A ring of keys hangs from a hook on the creature's trident.

\end{DndReadAloud}

The creature is a \textbf{nalfeshnee} named Maaltok, the guardian and keeper of the prison cells (areas W6a–W6d). Sometime in the last few decades, Maaltok switched allegiance from Graz'zt to Lolth. Since he is a recent recruit, Maaltok's yochlol handlers assigned him guard duty, which he performs grudgingly.

Maaltok knows each of the prison cells' inhabitants and doesn't expect infiltrators. If the characters are impersonating devotees of Lolth, Maaltok tells them who's imprisoned in each cell.

\subparagraph{Cell Keys}
Maaltok carries keys to each of the cells in the prison (areas W6a–W6d). Stealing Maaltok's key ring without his knowledge requires a successful DC 16 Dexterity (Sleight of Hand) check.

\subsubsection{W6: Holding Cells}
\begin{DndReadAloud}{}
Four jail cells are arrayed along the south wall of this cavern. The first holds a cyclops; the second holds a dead elf; the third holds a bugbear slumped against a bench; and the fourth holds a hunched, gray-skinned figure in robes. A worn chest sits along the room's northwest wall. Leaning next to it is a giant-size club.

\end{DndReadAloud}

The chest contains confiscated weapons and equipment from prisoners kept in the nearby cells. It's locked and requires a DC 20 Dexterity (Sleight of Hand) check using thieves' tools to open.

\subparagraph{Treasure}
Inside the chest are a \textit{+1 Dagger} and a \textit{+1 Longsword} confiscated from the recently deceased elf prisoner, Fernil, whose body is in area W6b. The giant-size club next to the chest is a \textit{+2 Greatclub} that belongs to Gertrude, the cyclops imprisoned in area W6a. The greatclub resizes to serve its wielder.

\subsubsection{W6a: Holding Cell A}
\begin{DndReadAloud}{}
A burly cyclops sits on a bench in this jail cell, shackled to the walls by her hands and legs. She holds her head in her hands, looking defeated.

\end{DndReadAloud}

This jail cell contains a bench and a bucket. The inhabitant is Gertrude, a chaotic neutral \textbf{cyclops} who was the lone survivor of an attack on a supply caravan near the ruins of the city of Ched Nasad. Makubli (see area W2) captured Gertrude weeks ago and brought her here as a prisoner, hoping she would reveal information about Blingdenstone's interest in Ched Nasad. Gertrude is a caravan guard and has no information the cultists can use; as soon as they realize this, they'll kill Gertrude, and she knows it. Her prized \textit{+2 Greatclub} is stashed in the chest in area W6.

\subparagraph{Recruiting Gertrude}
Gertrude hates Lolth and is eager to escape Web's Edge. If she realizes the characters are infiltrators, she begs them to free her. In exchange for her freedom, Gertrude offers to help the characters, including giving them her greatclub from the chest in area W6.

\subparagraph{Gertrude's Secret}
Even if the characters free her, Gertrude remains despondent. If the characters ask her what's wrong, Gertrude reveals that her friend, a svirfneblin named Rockzanna, was involved in a nearby mining operation and discovered that two of the operation's leaders were secret Lolth cultists planning an attack on the operation. Rockzanna was too scared to tell anyone besides Gertrude, and the cyclops fears a terrible fate for the miners.

Regardless of the characters' reaction to this revelation, learning it counts for the purposes of the Power of Secrets rules in this book's introduction. The mining operation is the same one the cultists in area W7 are discussing.

\subsubsection{W6b: Holding Cell B}
\begin{DndReadAloud}{}
This jail cell contains a bench and a tattered blanket. The desiccated corpse of an elf lies on the floor.

\end{DndReadAloud}

The body was once an elf named Fernil Orellian, an adventurer and priest of Corellon Larethian. A cloaker in the Underdark injured Fernil and killed his companions. An agent of Lolth came upon Fernil and marched him to Web's Edge, where he died shortly thereafter.

\subparagraph{Treasure}
Any character who examines the body and succeeds on a DC 16 Intelligence (Investigation) check notices a vial tucked into the corpse's belt. This is a \textit{Potion of Fire Resistance}—the agent of Lolth never noticed it.

\subsubsection{W6c: Holding Cell C}
\begin{DndReadAloud}{}
A bugbear with thick eyebrows is slumped on the bench against the wall in this jail cell.

\end{DndReadAloud}

Although she looks like she's sleeping, the bugbear in this cell, Rothgral, is dead. She deserted a band of mercenaries based near Mithral Hall after meeting a persuasive duergar agent of Lolth named Vundren. Vundren brought Rothgral here, but the cultists determined the bugbear didn't have the skills necessary to become a covert Lolth operative. The cultists imprisoned her until they could determine what to do with her, and she died about a day ago due to neglect.

\subsubsection{W6d: Holding Cell D}
\begin{DndReadAloud}{}
A gray-skinned creature with stringy, black hair and cultist's robes sits hunched in a corner of this cell.

\end{DndReadAloud}

This prisoner is Sril Brayspoke, a chaotic evil \textbf{grimlock}. Until a week ago, Sril was a cultist in residence at Web's Edge devoted to serving the Spider Queen. However, the bumbling Sril accidentally insulted the yochlol in area W12. The powerful Fiend had Sril thrown into jail, and the cultists are considering feeding him to Ker-arach to appease Lolth.

\subparagraph{Recruiting Sril}
Whether the characters are masquerading as Lolth worshipers or not, Sril begs to join the characters. A character who succeeds on a DC 14 Wisdom (Insight) check knows that Sril is still devoted to Lolth and will betray the characters if necessary to return to the cultists' good graces. If Sril is freed and knows the characters are impostors, he reveals their charade as soon as he sees another Lolth devotee.

\subsubsection{W7: Meeting Room}
\begin{DndReadAloud}{}
Relief carvings of spiders in webs decorate this chamber's walls, which are papered over with tactical maps and schematics. Gathered around a paperwork-covered table in the room's center are several cloaked figures as well as a horned, winged devil with a whip on her belt.

\end{DndReadAloud}

Seven Lolth devotees are meeting in this room to plan an assault on a svirfneblin mining operation located about 20 miles from Web's Edge. The devotees have coordinated with two Lolth operatives embedded in the mining operation. The devotees plan to kidnap the most powerful miners and feed them to Ker-arach in area W12. They'll kill any remaining miners who don't worship Lolth.

The devotees include a chaotic evil, elf \textbf{assassin} named Jolera Hartoph; two mages named Bromtok and Shiroktu, who are chaotic evil orcs; two grimlocks named Roltharni and Sharlotte, who are chaotic evil; and an \textbf{erinyes} named Fernitha.

Fernitha, Bromtok, and Shiroktu plan to attack the svirfneblin while the deep gnomes rest. If the characters eavesdrop on the planning before they investigate the summoning chamber (area W3), they learn that the high summoner, Grottenelle Stonecutter, is summoning a glabrezu to help with the operation.

The Lolth devotees here attack any apparent interlopers. The \textbf{glabrezu} summoned in area W3 joins the fight if the characters haven't already defeated the demon (or prevented its summoning).

\subsubsection{W8: Mission Hall}
The east end of this hallway holds a secret door that leads into the Sacred Web Hall (area W12). A character who examines the wall and succeeds on a DC 16 Intelligence (Investigation) check finds the door, which serves as an escape for the cultists who attend to Ker-arach when the creature is in a particularly foul mood.

\subsubsection{W9: Fodder Chamber}
\begin{DndReadAloud}{}
Peering from this rocky room's center are four large lizards. The pungent smell of raw meat hangs in the air.

\end{DndReadAloud}

The cultists keep a supply of prey creatures here to feed Ker-arach. Right now, the chamber is occupied by four giant lizards. The cultists release a lizard into the Sacred Web Hall (area W12) whenever Ker-arach is grumpy, providing her with recreation and a meal.

\subparagraph{Releasing the Lizards}
The cultists recently fed the lizards raw meat from the pantry in area W10, so the creatures are complacent. A character who holds raw meat near a giant lizard can successfully give that lizard simple commands (such as "follow me," "go where I point," or "bite her") for 10 minutes without giving the lizard the meat. If the lizard doesn't get the meat after 10 minutes, it attacks that character.

If either exit is left open, the lizards meander into the complex. If a lizard wanders into area W12, Ker-arach has disadvantage on Wisdom (Perception) checks and initiative rolls made in relation to the characters, since she is distracted while trying to capture the lizard and cocoon it in the chamber's central web.

\subsubsection{W10: Barrack of the Faithful}
\begin{DndReadAloud}{}
Five bedrolls, each with a small chest beside it, are tucked into this room. A large desk stands near the east wall. A small room to the south holds a long table with chairs around it.

\end{DndReadAloud}

This meager barrack houses the cultists who live in Web's Edge.

\subparagraph{Chests}
The chests contain the personal belongings of each cultist. Each chest is locked and requires a successful DC 18 Dexterity (Sleight of Hand) check using thieves' tools to open.

Each chest contains a ceremonial holy symbol of Lolth worth 50 gp. The chest near the bedroll closest to the entrance also contains two Potion of Superior Healing. The chest near the southernmost bedroll contains a \textit{Gem of Seeing}.

\subparagraph{Desk}
A \textit{Detect Magic} spell reveals an aura of conjuration magic around the desk, which bears a magical trap. Any character who touches the desk releases a swarm of spiders and must make a DC 18 Dexterity saving throw, taking 20 (8d4) poison damage on a failed save or half as much damage on a successful one. The trap triggers once, after which the spiders disappear and the desk becomes nonmagical. Casting \textit{Dispel Magic} on the desk removes the trap.

The desk contains notes about covert missions the Lolth worshipers are planning in the next few months. Additionally, it contains gold- and gem-encrusted, spider-shaped knickknacks worth 500 gp total.

\subparagraph{Dining Area}
The resident cultists eat and store their food in the room to the south. Characters who search the room find boxes of dried meat and other shelf-stable provisions. An ice box in the southeast corner is full of raw meat, which the characters can use to command the giant lizards in area W9.

\subsubsection{W11: Passageway}
\begin{DndReadAloud}{}
At the east end of this hallway is an enormous pile of bones.

\end{DndReadAloud}

The bones are remnants of meals eaten by Ker-arach, the spiderdragon in area W12.

\subsubsection{W12: Sacred Web Hall}
\begin{DndReadAloud}{}
An enormous web stretches across this open cavern's center, its strands plastered over stalagmites and stalactites. Crawling on the web is a gigantic reptilian creature with eight legs. Near that creature stands a smaller, one-eyed creature with a body that resembles melting wax.

\end{DndReadAloud}

This cavern's ceilings are 50 feet tall. The cavern is the lair of Ker-arach, the \textbf{spiderdragon} (see appendix A). Ker-arach crawled into this cavern through a temporary rift to the Abyss that opened during a ritual conducted here about half a year ago. Ker-arach brought with her a piece of the \textit{Rod of Seven Parts}, which she uses to answer questions posed directly to the Spider Queen. Ker-arach is essentially a tourist attraction for the few Lolth devotees who know she exists. The cultists consider her a sign of Lolth's favor, even though she eats hundreds of pounds of meat every few days and doesn't otherwise contribute to the complex.

Standing near Ker-arach is a \textbf{yochlol} named Ylellith, which the cultists recently summoned to help with a mission. When the creatures see the characters and realize the characters are enemies, they attack and fight to the death.

\subparagraph{Ker-arach's Web}
The web in the center of the cavern is made of ultra-strong, ultra-sticky strands of Ker-arach's silk. Additionally, Ker-arach has spun other, smaller webs throughout the area. The webs are \textit{difficult terrain}. Any creature that enters the webbing for the first time on a turn or ends its turn there must succeed on a DC 12 Strength saving throw or it becomes stuck and has the restrained condition. As an action, a creature can try to pull itself or another creature within its reach from the webbing, doing so with a successful DC 15 Strength (Athletics) check. A creature freed in this way is no longer restrained by the webbing.

The webs are flammable. Any 5-foot cube of webs exposed to fire burns away, dealing 5 (2d4) fire damage to any creature in that area.

\subparagraph{Retrieving the Rod Piece}
The Rod of Seven Parts of the \textit{Rod of Seven Parts} is wrapped tightly in silk at the center of this cavern's web. Characters who have darkvision or a light source can see the rod piece from up to 30 feet away. A character within reach of the rod piece can use a sharp tool to cut it free of the web as an action. For more about the \textit{Rod of Seven Parts}, see this book's introduction.

\section{Next Steps}
Once the characters have acquired the first piece of the \textit{Rod of Seven Parts}, they can return to Sigil through the portal that remains open outside Web's Edge. At this point, the characters will likely want to rest, confer with their allies in Sigil, and see where this first rod piece points them to go next.

When the characters return to the sanctum, Malaina is there (see "The Sigil Sanctum" earlier in this chapter for more details). Malaina offers to help however she can, including seeking information outside of Sigil or retrieving specific magic items the characters might want.

Once the characters have concluded their business in the sanctum, they can begin their quest for the location of the second rod piece, as described in the next chapter.

\chapter{The Lambent Zenith's Last Voyage}
The hunt for the second piece of the Rod of Seven Parts brings the characters to the starry void of the Astral Plane. Within its silvery depths, alien predators lurk in silence and fallen gods lie in stasis. Adventurers known across the multiverse as spelljammers gallivant through space in ships powered by magic.

In this chapter, the characters search the Astral Sea for the second rod piece. The characters need to explore the wreckage of a spelljamming ship called the \textit{Lambent Zenith}, then retrieve the rod piece from a dragon-like creature that guards it inside the heart of a fallen god.

\section{Running the Adventure}
This chapter begins after the characters retrieve the first piece of the \textit{Rod of Seven Parts}. A character who holds that piece knows instinctively that the next piece is in the part of the Astral Plane called the Astral Sea. As Mordenkainen previously explained, the portal in the Sigil sanctum leads to the general area of the rod piece the characters seek.

\subsection{Features of the Astral Sea}
In the Astral Sea, time is meaningless, and creatures can survive there indefinitely without food or drink.

The locations explored within this chapter are a small fraction of what can be found in the Astral Sea. As such, the guidance here focuses on information relevant to the chapter's contents.

\subsubsection{Air}
The Astral Plane contains breathable, comfortable air. Unless otherwise stated, creatures can breathe normally.

\subsubsection{Gravity}
The characters are in an area of normal gravity during this chapter.

\subsubsection{Movement}
Though many creatures use spelljamming ships or other vessels to traverse the Astral Sea, a vehicle isn't required. In the Astral Sea, a traveler can propel themself by thought alone. A creature can move in any direction at a flying speed in feet equal to 5 × its Intelligence score.

\subsection{Character Advancement}
The characters should be 12th level when this chapter begins. The characters gain a level after they retrieve the Rod of Seven Parts of the \textit{Rod of Seven Parts} from the hertilod.

\subsection{Power of Secrets}
The characters can learn two secrets in this chapter that are applicable to the rules in "The Power of Secrets" section in this book's introduction:


\begin{description}
\item[Figaro's Secret.]
Figaro, the tiefling first mate of the Lambent Zenith, knew about the dangers of the portion of the Astral Sea the ship was passing through but deliberately hid this information from the captain. The characters can learn his secret in area Z8 of the ship's wreckage.
\item[Ikasa's Secret.]
The blink dog Ikasa knows about another survivor of the pirate attack that stranded him and his best friend, the elf Daveras. The characters can learn this secret in area Z12 of the ship's wreckage.
\end{description}

\subsection{Second Rod Piece}
The second piece of the \textit{Rod of Seven Parts} is inside the hertilod in area A2 in the "Heart of Havock" section later in this chapter. For more information about the rod and the spell this piece allows its wielder to cast, see this book's introduction.

\section{A Doorway to Space}
When the characters are ready to continue their adventure, they can step through the portal in the Sigil sanctum and emerge on the Astral Plane in the Astral Sea. Like with the first rod piece, the portal opens near where the second piece is located, but it's up to the characters to find exactly where the piece is ensconced.

\subsection{The Astral Plane}
By researching in the Sigil sanctum or talking with the wizards, the characters can learn the following information about the Astral Plane and the Astral Sea:


\begin{description}
\item[Planar Nature.]
The Astral Sea is in the Astral Plane, which is colloquially known as the realm of thought and dream. Creatures can propel themselves through the plane by merely thinking about moving in a specific direction. Creatures can also traverse the Astral Plane using vessels called spelljamming ships.
\item[Dead Gods.]
Scattered throughout the Astral Sea are the remains of dead and dying gods who are here either because they were forgotten by their worshipers or slain at the hands of more powerful entities.
\end{description}

If the characters ask about the location the rod piece points to, Alustriel determines that the dying god Havock is near the piece.


\begin{description}
\item[Havock.]
Havock's petrified form is hundreds of miles long and weighs thousands of tons. It once had eight legs and two heads that each held a single, unblinking eye, but as Havock lost worshipers and power over millennia, its legs and heads snapped off, destined to drift forever in the Astral Sea.
\item[Timelessness.]
Time has no meaning on the Astral Plane. Creatures on the Astral Plane don't age or experience hunger or thirst.
\end{description}

\section{Into the Astral Sea}
When the characters step through the portal to the Astral Sea to pursue the second rod piece, read or paraphrase the following:

\begin{DndReadAloud}{}
You step through the doorway and enter a silver-clouded void. Lucent wisps of white and gray fog swirl in the distance among pinpricks of starlight. For a split second, you have no sense of direction. Then, you start plummeting. The silver clouds shift, and you see that you're falling toward a colossal, misshapen mass.

\end{DndReadAloud}

The characters emerge from a free-floating doorway on the Astral Plane. On passing through, the characters are subject to the gravity field that extends from the stony mass they're falling toward: the petrified body of a dying god called Havock. The characters can immediately propel themselves using their thoughts as described in the "Movement" section earlier in this chapter.

The character holding the rod piece divines that the next rod piece is located within the stony mass. The doorway is anchored approximately 1 mile from the mass's surface.

\subsection{Stalkers in the Stars}
As the characters move toward the mass, something follows them. Read or paraphrase the following:

\begin{DndReadAloud}{}
Amid the clouds, two large globes of light bob toward you.

\end{DndReadAloud}

Characters can make a DC 18 Wisdom (Perception) check. On a successful check, a character notices a large, anglerfish-shaped outline following each bobbing light—two hungry star anglers (see appendix A) are stalking the party. The star anglers attack immediately.

Once the star anglers are dispatched, the characters can continue on to the surface of the dying god, where the wreckage of the \textit{Lambent Zenith} lies.

\section{Wreck of the Lambent Zenith}
With its navy-blue hull and golden gossamer sails, the spelljamming galleon \textit{Lambent Zenith} once cut an elegant silhouette through the Astral Sea's silver clouds. Its captain, the deva Inda Malayuri, was an arcanist and emissary tasked with guiding the lost and bringing peace to tumultuous worlds. During a voyage, Inda uncovered a piece of the \textit{Rod of Seven Parts}.

Using her extensive arcane knowledge, Inda harnessed the magic within the rod piece to augment the \textit{Lambent Zenith}, allowing the ship to travel to the farthest corners of the multiverse in the blink of an eye.

But the \textit{Lambent Zenith}'s expedition screeched to a halt when it encountered a dying god of chaos named Havock drifting through the Astral Sea. The latent chaos magic around Havock twisted the ship's capabilities; the \textit{Lambent Zenith} was ripped asunder, and its prow plunged into Havock's heart.

When the characters approach Havock and the wreckage (shown on map 3.1), read or paraphrase the following:

\begin{DndReadAloud}{}
The stony mass isn't a planet or an asteroid, but a colossal creature that appears lifeless. Shattered ribs arch over the creature's mossy spine, and the air crackles with decaying magic.

Among the bones is the shipwreck of a large galleon broken into three large chunks: the sterncastle, nestled in the corpse's hip bones; the starboard section, embedded in the ribcage; and the prow, stabbed into the creature's heart.

\end{DndReadAloud}

The \textit{Lambent Zenith} broke into three distinct pieces: the stern segment, the starboard segment, and the prow segment. On their initial approach, the characters are closest to the stern segment and area Z1, but they can approach the segments in whichever order they like.

The character holding the rod piece senses that the next piece is located somewhere in the ship's wreckage.

\subsection{Accessing the Rod Piece}
The rod piece was stored in a safe room in the \textit{Lambent Zenith}'s prow (area Z19). When the ship crashed, the safe room's security wards activated, and the room's doors sealed magically. To access the safe room and uncover the rod piece, the characters must first deactivate the security wards.

\subsubsection{Ward Runes}
The security wards are powered by two magical runes: one in area Z8 and one in area Z13. The runes are invisible. If a creature can see invisible objects, each rune looks like a stylized carving of a crescent moon pulsing with silver light. Casting \textit{Dispel Magic} (DC 17) on a rune destroys it. A rune can also be deactivated if a creature within 5 feet of it says, "The moon sings a song for the lost."

Both runes must be either rendered inactive or destroyed for the safe room's door to open. This allows access to the heart of Havock, where the rod piece is located.

\subsection{Arcane Portals}
The crew of the \textit{Lambent Zenith} used the magic contained within the rod piece to create portals through which the ship could travel across the multiverse with ease. However, when the \textit{Lambent Zenith} approached Havock, the rod's magic mixed with the latent power in Havock's dying body and went awry. This volatile combination now causes teleportation spells and effects to function strangely in the segments of the shipwreck.

When a creature within a wreck segment casts a spell (or uses a similar magical effect) that would teleport it or another creature, instead of the normal effect, a 5-foot-diameter circular portal appears in an unoccupied space within 30 feet of the casting creature. The portal is a glowing ring filled with opaque mist and remains open for 1 minute.

When a portal appears, roll on the Portal Exit table to determine where the portal leads. If the result is where the characters already are, roll again. Any creature or object that passes through an open portal appears in a random unoccupied space in the exit location.

\begin{DndTable}[header=Portal Exit]{cX}

d6 & Exit Location\\
1 & Sterncastle deck (stern segment; area Z1a)\\
2 & Companionway (stern segment; area Z5)\\
3 & Starboard top deck (starboard segment; area Z10)\\
4 & Grell nest (starboard segment; area Z11)\\
5 & Study (starboard segment; area Z13)\\
6 & Bridge (prow segment; area Z17)\\
\end{DndTable}

\subsection{Traversing Wreck Segments}
Characters can forgo the portals and travel between the wreck segments on their own, either by walking or by flying. The segments are 300 feet from each other.

The areas between the segments contain a multitude of hungry astral predators. The creatures trapped in the segments don't traverse the wreck for this reason. When the characters enter one of these areas, roll on the Random Wreck Encounters table to see what confronts the characters.

\begin{DndTable}[header=Random Wreck Encounters]{cX}

d10 & Creature\\
1–2 & Two night scavvers (see appendix A)\\
3–4 & One \textbf{star angler} (see appendix A)\\
5–6 & One \textbf{cloaker}\\
7–10 & No encounter\\
\end{DndTable}

\subsubsection{Stern Segment}
The largest piece of the wreck, the stern segment, is nestled in Havock's hip bones. A handful of stranded shipwreck survivors led by the \textit{Lambent Zenith}'s first mate, Figaro, are encamped on this segment.

A \textbf{death slaad} recently infiltrated the camp, sowing chaos and magically manipulating Figaro's mind. This drove Figaro into a paranoid state, causing him to lock himself in his quarters, where one of the safe room's ward runes is located.

Characters approaching this wreck segment can land on either the sterncastle deck (area Z1a) or in the companionway (area Z5).

\subsubsection{Starboard Segment}
The starboard piece of the ship crashed into Havock's ribcage. A small lifeboat also crashed into this area, carrying refugees from a space pirate attack: Daveras, his blink dog companion Ikasa, and a treant named Redbud. Daveras escaped the wreckage and joined the survivor camp on the stern segment. Unknown to Daveras, Ikasa and Redbud also survived and are on the starboard segment.

Wreckage blocks off this segment's lower decks from the outside. Characters approaching this segment must land on the starboard segment's section of the top deck (area Z10).

\subsubsection{Prow Segment}
The wreck of the \textit{Lambent Zenith}'s prow impaled the dormant god's heart. The ship's captain, Inda, is here, as is the entrance to the safe room.

Characters approaching this segment can land on either the prow segment's top deck (area Z14a) or the lower deck (area Z17).

\subsection{General Features}
The areas of the \textit{Lambent Zenith}'s wreckage have the following features:


\begin{description}
\item[Gravity.]
The wreck and the areas around it have normal gravity.
\item[Lighting.]
Unless otherwise stated, the ambient silvery glow of the Astral Sea fills the wreck, rendering its areas brightly lit.
\item[Walls, Ceilings, and Floors.]
The floors and walls of the ship are made from faded blue wood. The ceilings of the ship's lower decks are 10 feet high.
\end{description}

\subsection{Lambent Zenith Locations}
The areas of the wreck are keyed to map 3.1.

\subsubsection{Z1a–Z1b: Sterncastle Deck and Below}
\begin{DndReadAloud}{}
Scraps of golden sails dangle from the sterncastle's ruined mast. Two armored githyanki patrol the deck and below, silver swords at their hips.

\end{DndReadAloud}

Two lawful neutral githyanki knights named Lysan and Zastra, who served as bosuns on the \textit{Lambent Zenith}, keep watch here. Characters who want to avoid the githyanki must succeed on a DC 15 Dexterity (Stealth) check. The knights notice characters who fail the check, beckoning them to approach and state their business.

Lysan and Zastra are cautious but friendly. They have a sardonic sense of humor developed after being stranded for so long. If the characters express no hostile intent, the bosuns welcome them to the \textit{Lambent Zenith} and call for Ilren, who stepped in as the camp's interim leader when the ship's first mate fell ill. The githyanki explain that they haven't heard from the ship's captain, Inda, who Lysan and Zastra assume either died during the crash or is stranded elsewhere in the wreckage. Lysan and Zastra introduce Ilren to the party.

If the characters are hostile, Lysan and Zastra shout for help and attack the characters. After all combatants have taken a turn, all five sailors in the upper crew quarters (area Z2) join the fight.

\subparagraph{Talking with Lysan and Zastra}
Lysan and Zastra know the following information:


\begin{description}
\item[Last Voyage.]
The Lambent Zenith was a sailing ship that sought to herald hope, community, and order across the multiverse. But this mission came to an end a few months ago when the ship collided with the body of a dying god.
\item[Strange Magic.]
The ship's prow was outfitted with "fancy magic stuff" beyond the githyanki's understanding. This magic malfunctioned when the ship crashed, and now teleportation magic used within the wreck is warped. Instead of teleporting creatures, the magic instead spawns a portal that leads somewhere else in the wreckage.
\item[Terrors of Space.]
The areas beyond the wreck are teeming with hungry astral predators. The survivors are too afraid of these predators to risk traveling between the segments. The camp has also been plagued recently by aberrant monsters such as cloakers, though the survivors don't know why.
\end{description}

If asked for specifics, Lysan and Zastra share a look before directing the characters to speak with the second mate, Kycera, in the galley (area Z3). A character who succeeds on a DC 14 Wisdom (Insight) check intuits that the two are uneasy about some of their crewmates. If the characters press the issue, the githyanki pull them aside and admit that the camp's interim leader, Ilren, and the \textit{Lambent Zenith}'s first mate, Figaro, would have more information. However, Figaro has been acting paranoid of late and refuses to leave his quarters (area Z8).

\subparagraph{Talking with Ilren}
Ilren appears as a boisterous giff—a bulky individual with a head similar to that of a hippopotamus—wearing a sleek red coat. But Ilren is a disguised \textbf{death slaad} who can cast the spell \textit{Modify Memory} (spell save DC 15) once per day. Ilren is the reason the camp suffers Aberration attacks; the death slaad sneakily directs monsters toward the wreck's survivors and revels in the victims' suffering. Ilren also magically modified Figaro's memories and usurped power in the camp.

On meeting the characters, Ilren claims to come from a similar background: the giff is an adventurer with heroic dreams who became stranded on this wreck when his skiff crashed into Havock. Ilren claims that nothing of his skiff remains.

Ilren welcomes the characters to explore the \textit{Lambent Zenith}'s wreck and make themselves at home, but a character who succeeds on a DC 17 Wisdom (Insight) check intuits that Ilren's jolly demeanor belies a suspicious and cold attitude. If asked about Figaro, Ilren breezily explains that Figaro is recovering from an illness but doesn't elaborate. Ilren adds that he is happy to serve as leader while Figaro is indisposed.

Ilren doesn't reveal its true form unless attacked.

\subsubsection{Z2: Upper Crew Quarters}
\begin{DndReadAloud}{}
This room is plain yet inviting. Five sailors lounge in cramped bunks and colorful hammocks strung from the ceiling.

\end{DndReadAloud}

Five lawful good sailors from the Upper Planes (use the \textbf{veteran} stat block, but the sailors are Celestials instead of Humanoids) are lounging here. The sailors are indifferent to the characters but jump into combat to protect themselves and their fellows.

These sailors were lackeys aboard the \textit{Lambent Zenith}. Though they can recount the ship's mission and subsequent crash, these sailors know nothing about the rod piece.

\subsubsection{Z3: Galley and Pantry}
\begin{DndReadAloud}{}
The sound of sizzling and the smell of spices emanate from the galley. A tall orc woman in black leather armor works busily at the stove. A gold tattoo of a crescent moon surrounds her left eye.

\end{DndReadAloud}

Kycera Duskstride is the \textit{Lambent Zenith}'s second mate and is a chaotic good, orc \textbf{assassin}. A former pirate, Kycera saw the error of her ways and joined the \textit{Lambent Zenith}'s crew to find redemption. Now she spends her time in the wreck's galley. Though the survivors don't require food in the timeless void of the Astral Sea, Kycera finds that cooking and eating bring the restless survivors comfort.

\subparagraph{Talking with Kycera}
Kycera gladly discusses the general history of the Lambent Zenith. However, if asked about the rod piece or the security system, Kycera clams up, suspicious of the characters' motives. A character who succeeds on a DC 13 Charisma (Persuasion) check assures her of their heroic purpose, at which point Kycera reveals the following information:


\begin{description}
\item[Precious Cargo.]
The Lambent Zenith carried an extremely rare and powerful magic artifact. As a former pirate, Kycera helped Captain Inda design the security system to protect the artifact.
\item[Safe Room.]
The artifact was stored in a safe room at the front of the ship. During the Lambent Zenith's crash, arcane wards should have sealed the safe room's doors. The safe room is inaccessible while the wards are active.
\item[Ward Runes.]
Two runes hidden on the ship power the wards. Kycera knows one rune was placed in the ship's stateroom, but she doesn't remember the second rune's location. She believes that magic or a pass phrase can deactivate the runes. She doesn't know the pass phrase, but the captain would know.
\end{description}

If asked about Figaro, Kycera expresses annoyance at his seclusion; her tone belies a sense of parental concern. If the characters help Figaro and defeat Ilren, Kycera is grateful and provides all the information above willingly.

\subparagraph{Treasure}
Within the pantry are 2d4 unopened casks of aged wine, each worth 500 gp. Kycera happily allows the characters to take them, as she's concerned about the crew getting too drunk to fend off threats.

\subsubsection{Z4: Mess Hall}
\begin{DndReadAloud}{}
Empty wooden tables and chairs fill the mess hall. In the corner, an elf man plays a game of cards.

\end{DndReadAloud}

The elf is Daveras, a neutral \textbf{druid}. Daveras joined the camp after his lifeboat crashed into Havock. The remains of his lifeboat rest on the starboard segment.

If the characters haven't met Ilren already, Ilren is here playing cards with Daveras.

\subparagraph{Talking with Daveras}
Though gruff, Daveras is friendly. He shares the following information:


\begin{description}
\item[Missing Camp Member.]
Daveras's usual card partner, a halfling woman named Cirit, hasn't shown up for their games lately. Camp members said Cirit left to scout the other wreck segments, but Daveras seems doubtful. Unknown to Daveras, Cirit is actually in the Lambent Zenith's brig (area Z9).
\item[Paranoid Leader.]
Figaro, the camp's first leader, locked himself in his quarters and refuses to communicate. Daveras tried to talk to him, but Figaro used a magical device to erect an impenetrable barrier around the room.
\item[Personal History.]
Daveras wasn't part of the Lambent Zenith's crew. When his spelljamming ship, the Verdant Branch, was attacked by space pirates, he escaped on a lifeboat with his best friend, a \textbf{blink dog} named Ikasa, as well as a treant named Redbud, who had grown into the lifeboat's hull. The lifeboat crashed into a different segment of the wreck, and Daveras outran the raiders to get to the camp on the stern segment. He believes no one else survived.
\end{description}

If Ilren isn't present, Daveras confides that he noticed Figaro's odd paranoia started shortly after Ilren arrived.

If Daveras is reunited with Ikasa (found in area Z12), he is elated and happily vouches for the characters when talking to Figaro in area Z5.

\subsubsection{Z5: Companionway}
\begin{DndReadAloud}{}
The forward end of this hallway is buried in Havock's flesh. The stern end holds two doors and two stairways leading up, and another door on the port side bears a gilded plaque that reads "Stateroom."

\end{DndReadAloud}

The port door is locked and leads to a guest stateroom that now serves as Figaro's living quarters. Due to his paranoid state, Figaro uses a \textit{Cube of Force} to deter any living matter from entering his room.

A character trying to sneak up to Figaro's locked door must succeed on a DC 18 Dexterity (Stealth) check to avoid being noticed. An unnoticed character can, as an action, use thieves' tools to try to pick the door's lock, doing so with a successful DC 20 Dexterity (Sleight of Hand) check.

\subparagraph{Gaining Figaro's Trust}
Figaro deactivates the cube's barrier if he feels he can trust the characters. The barrier allows sound to pass through, meaning that the characters can talk with Figaro. A character talking with Figaro can make a DC 25 Charisma (Persuasion) check. On a successful check, Figaro trusts the characters, unlocks the door, and permits the characters to enter.

By helping the people found throughout the wreck, the party can find individuals who will vouch for them. If two or more individuals vouch for the party, Figaro deactivates the barrier without the characters needing to make a check. If only one individual vouches for the party, a character makes the check with advantage.

\subsubsection{Z6: Lower Crew Quarters}
\begin{DndReadAloud}{}
The bunks and hammocks here are empty, and the walls are covered with colorful pieces of wood, each with something written on it.

\end{DndReadAloud}

The painted pieces of wood are makeshift remembrance markers, honoring those who died in the \textit{Lambent Zenith}'s crash. Each piece of wood has a name scribbled on it with smudged graphite. The survivors long ago burned the deceased's bodies on the ship's decks.

\subsubsection{Z7: Storage}
\begin{DndReadAloud}{}
Unused rigging and dusty old sails fill this old storage closet.

\end{DndReadAloud}

Characters who have a passive Wisdom (Perception) score of 15 or higher hear muffled shouting from below.

\subparagraph{Secret Entrance}
Loose floorboards on the room's port side can be pried up, allowing access to the brig (area Z9). Ilren recently used this passage, then attempted to hide the loose floorboards beneath a canvas tarp. A character who searches the floor and succeeds on a DC 12 Intelligence (Investigation) check notices the edge of the loose floorboards jutting upward beneath the canvas.

\subsubsection{Z8: Stateroom}
\begin{DndReadAloud}{}
The stateroom is simple but elegant. The warm light of the wall sconces bathes the mahogany bookshelves and crimson bedsheets. Sitting on the edge of the bed is a purple-skinned tiefling. He stares blankly at the wall, but his posture is as tense as a wound spring.

\end{DndReadAloud}

The tiefling is Figaro, a lawful good \textbf{mage}. If Figaro allowed the characters into his room willingly or if enough individuals vouch for them, Figaro shakily introduces himself and is willing to talk to the characters. Otherwise, he screams and attacks. Figaro surrenders after he is reduced to half of his hit point maximum, but he still refuses to talk with the characters.

\subparagraph{Talking with Figaro}
If the characters gained Figaro's trust, he answers questions about the Lambent Zenith. He knows the following:


\begin{description}
\item[Experiments with the Rod.]
The Lambent Zenith harnessed the conjuration magic within a piece of the \textit{Rod of Seven Parts} to create portals that allow the ship to quickly traverse the multiverse. The rod piece was stored below the bridge on the ship's prow.
\item[Emergency Security Protocol.]
Knowing that the rod piece shouldn't fall into the wrong hands, the Lambent Zenith's captain installed a security protocol that sealed the rod piece's location in the event of an emergency—such as a crash. As far as Figaro knows, the rod remains in the safe room (area Z19), though that area is inaccessible due to the activated safety protocol.
\item[Ward Runes.]
The security protocol is fueled by two runes in different locations aboard the ship. Both runes must be dispelled to access the room that holds the rod piece. Figaro knows the runes are invisible. He also knows one rune is located in this stateroom, on the bed's headboard, while the other is in the ship's study (area Z13) on the wall above the desk. He suspects casting \textit{Dispel Magic} on a rune would deactivate it, but the caster needs to see the rune. He also suspects the captain knew another way to deactivate it. If she survived on another segment, the characters could try to reach her to ask; if not, it might be recorded it in her notes.
\end{description}

\subparagraph{Mending Figaro's Memories}
Ilren used repeated \textit{Modify Memory} spells to manipulate Figaro's mind. A character who studies Figaro's symptoms and succeeds on a DC 20 Intelligence (Arcana) check recognizes Figaro's glassy gaze as a telltale sign of memory-tampering magic. Restoring Figaro's true memories causes Figaro to immediately recall Ilren's deception. If the characters haven't confronted Ilren, Figaro rushes off to face the death slaad.

\subparagraph{Figaro's Secret}
After recovering his memories, Figaro remains agitated even if he confronts Ilren. If the characters ask what's bothering him, Figaro admits he knew about the dangers of this stretch of Astral Sea but didn't warn the captain, since this route was economical and Figaro believed it would be easy for the crew to traverse.

Regardless of the characters' reaction to this revelation, learning it counts as a secret for the purposes of the Power of Secrets rules in this book's introduction.

\subparagraph{Ward Rune}
The ward rune is located on the bed's headboard.

\subsubsection{Z9: Brig}
\begin{DndReadAloud}{}
Small cells with thin iron bars line both sides of the center hallway. You hear shouts for help from one of the cells.

\end{DndReadAloud}

The shouting is coming from a \textbf{couatl} named Cirit, currently in the form of a halfling \textbf{priest}. Cirit was one of the \textit{Lambent Zenith}'s crew members and was present when Ilren infiltrated the survivors' camp.

Cirit saw Ilren's true aberrant form. But when Cirit confronted Ilren about the deception, the death slaad overpowered her. Believing Cirit to be an ordinary halfling spellcaster, Ilren hoped to use her to create a \textbf{green slaad} and locked her in the brig instead of killing her. Cirit has been trying to escape since.

The door to Cirit's cell is locked. A character can pick the lock with a successful DC 15 Dexterity (Sleight of Hand) check using thieves' tools, or a character can wrest open the door with a successful DC 15 Strength (Athletics) check.

\subparagraph{Talking with Cirit}
On seeing the characters, Cirit is relieved and immediately asks for their assistance. She provides the following information:


\begin{description}
\item[Ilren's Aberrant Nature.]
Cirit saw through Ilren's disguise. She identifies Ilren as a \textbf{death slaad}, an Aberration that relishes the suffering of others.
\item[What Happened.]
Cirit tried to speak with Figaro about Ilren, but Figaro did nothing; Cirit suspects Ilren magically manipulated Figaro's mind. She then confronted Ilren, but the death slaad overpowered her. She doesn't know why Ilren kept her alive, and she's not keen on finding out why.
\end{description}

Cirit implores the characters to help her thwart Ilren. If the characters defeat Ilren, Cirit vouches for them when talking with Figaro in area Z5.

\subsubsection{Z10: Starboard Top Deck}
\begin{DndReadAloud}{}
The deck's flooring here is uneven, with roots woven around the planks. Jutting from the segment's fore is the wreckage of a lifeboat. Atop this smaller wreck is a stout tree with vibrant pink and red blossoms.

\end{DndReadAloud}

The tree atop the crashed lifeboat is a treant named Redbud (use the \textbf{treant} stat block, but Redbud's movement speed is 0 feet) that joined the crew and grew into the lifeboat to strengthen its hull. The treant nearly died when the boat crashed, but Redbud survived by sending roots into the \textit{Lambent Zenith}'s wreck to feed off Havock. Redbud's main purpose now is protecting the blink dog Ikasa (found in area Z11), who remains trapped aboard the lifeboat.

\subparagraph{Interacting with Redbud}
Redbud remains silent and motionless until a creature attempts to enter the segment's lower deck (areas Z11–Z13). Initially, Redbud affably discourages characters from proceeding deeper into the wreck; when pressed, Redbud claims this is to protect them from hungry parasites that have nested below. However, a character who makes a successful DC 13 Wisdom (Insight) check intuits that while Redbud is telling the truth about monsters below, the treant has ulterior reasons for keeping the characters away from the lower deck.

If the characters question the treant's sincerity, Redbud admits to the deception. Redbud explains that a dear friend is trapped below, and the treant's roots are protecting this friend from being eaten by the monsters nesting on the lower deck. Redbud implores the characters to clear out the nest, allowing the characters to proceed further only if they agree.

\subsubsection{Z11: Grell Nest}
\begin{DndReadAloud}{}
Thick, pale roots here form two dense walls on either side of the room. Six bulbous, brain-shaped creatures with snapping beaks and barbed tentacles bob through the air.

\end{DndReadAloud}

Six grells have made their home here, feeding off hapless raiders looking to pillage this wreck segment. The grells are hostile toward all other creatures.

\subparagraph{Redbud's Roots}
The treant's roots have warped and splintered the floorboards, rendering squares with roots \textit{difficult terrain}.

Redbud's roots also block the entrances to the lifeboat wreck (area Z12) and the study (area Z13), the latter of which contains one of the safe room's ward runes. Redbud moves these roots if the characters eradicate the monsters nesting below. The roots also retract if Redbud is killed or forced to relocate.

\subsubsection{Z12: Lifeboat Wreck}
\begin{DndReadAloud}{}
A cot is shoved against one wall here opposite a desk. Sitting in the room's center is a reddish-brown dog with a leather collar studded with glowing crystals.

\end{DndReadAloud}

Ikasa the \textbf{blink dog} is outfitted with a magical collar that allows him to speak Common. Afraid to use his teleportation ability due to the wreck's warped magic, Ikasa remained in the lifeboat, protected by Redbud.

On seeing the characters, Ikasa excitedly bounds over and introduces himself. If the characters mention Daveras, Ikasa asks to be reunited with his old companion.

\subparagraph{Ikasa's Secret}
If the characters speak with Ikasa about Daveras, the blink dog becomes uncharacteristically morose. On further inquiry, Ikasa reveals he saw Daveras's friend, the halfling Palenna Tindertoe, escape the wreckage of the Verdant Branch after pirates attacked it. Palenna floated away from the ship using the power of her mind and likely still floats somewhere in the Astral Sea. Ikasa didn't tell Daveras, who believes he, Redbud, and the blink dog were the only survivors of the attack. Ikasa hasn't had the opportunity to find Palenna due to being trapped on the lifeboat, and he is afraid Daveras will become despondent thinking about his friend trying to survive in the Astral Sea alone.

Regardless of what the characters do on learning this, it counts as a secret for the purposes of the Power of Secrets rules in this book's introduction.

\subparagraph{Treasure}
The lifeboat's interior contains survival necessities. A character who searches the lifeboat and succeeds on a DC 18 Intelligence (Investigation) check also finds an emergency fund of 100 gp stashed in one of the desk's drawers.

\subsubsection{Z13: Study}
\begin{DndReadAloud}{}
Books and scrolls concerning arcane subjects pack the bookshelves here. Pressed against the wall is a redwood desk, atop which are clean sheets of parchment and two inkwells.

\end{DndReadAloud}

\subparagraph{Secret Strongbox}
A character who inspects the bookshelves and succeeds on a DC 15 Intelligence (Investigation) check discovers that one of the books—a volume titled Dissertations on the Abstruse Mind—is a disguised strongbox. The box is locked, but the lock can be picked with a successful DC 15 Dexterity (Sleight of Hand) check using thieves' tools or forced open with a successful DC 11 Strength (Athletics) check. Inda carries the strongbox's key.

The box is trapped. When a creature attempts to open the box by any means other than the key, poisonous gas puffs out the keyhole. The creature must succeed on a DC 20 Constitution saving throw or have the poisoned condition for 1 hour.

The box contains a \textit{Potion of Mind Reading} and a journal. The journal belongs to Inda and contains writings in Common and Celestial regarding her experiments with the rod piece. A character who spends 10 minutes reading the journal learns the pass phrase needed to deactivate the ward runes protecting the safe room: "The moon sings a song for the lost."

\subparagraph{Ward Rune}
The ward rune is on the wall above the desk.

\subsubsection{Z14a–Z14b: Forecastle and Top Deck}
\begin{DndReadAloud}{}
A colossal heart—part flesh and part stone—looms over this segment's top deck. What remains of the ship's prow is plunged deep into the heart, and the deck rumbles as the heart pulses. Scraps of gossamer sails dangle off the ship's broken-off mast, which is jammed through one of the doors leading into the forecastle.

\end{DndReadAloud}

The mast blocks the entrance to the navigation room (area Z16). The door to the captain's quarters (area Z15) is unobstructed but locked; the lock can be opened with a successful DC 15 Dexterity (Sleight of Hand) check using thieves' tools. As an action, a character can attempt a DC 15 Strength (Athletics) check to force the door open. A character who has a passive Wisdom (Perception) score of 13 or higher and who is standing within 5 feet of the door hears someone inside.

\subsubsection{Z15: Captain's Quarters}
\begin{DndReadAloud}{}
In this cabin, clothes are tossed over chair backs, and open books are strewn across tables. At the back of the room stands a deva. Her right leg below her knee is a wood-and-metal prosthetic, and from her back extends a massive, white-feathered wing.

\end{DndReadAloud}

The one-winged woman is the captain of the \textit{Lambent Zenith}, Inda Malayuri. Inda uses the \textbf{deva} stat block, except she has a flying speed of 0 feet.

Inda's right leg and wing were severed in an attack years ago. She now uses prosthetics, but her prosthetic wing was damaged in the crash. She has been trying to repair it in the navigation room (area Z16). She is determined to escape the wreck and reunite with whatever remains of her crew, though she worries the rod piece will be stolen by raiders if she leaves.

\subparagraph{Talking with Inda}
Having encountered raiders attempting to pillage the wreck, Inda is wary of the characters and pointedly asks about their motives. She is hesitant to provide any information about the Lambent Zenith's cargo—especially the rod piece—and isn't swayed by bribes or threats.

A character talking with Inda can make a DC 15 Charisma (Persuasion) check. If a character offers to aid her with repairing her prosthetic wing or to reunite her with her crew in the stern segment, this check is made with advantage. On a successful check, Inda agrees to help the characters. In addition to the general history of the \textit{Lambent Zenith}, Inda knows the following points of information:


\begin{description}
\item[Location of the Rod Piece.]
The rod piece was stored in a safe room directly beneath the ship's bridge. To protect the rod piece from raiders, the safe room was equipped with magical security wards that triggered when the ship crashed.
\item[Deactivating the Ward Runes.]
The safe room's wards are powered by two runes. One is located on the headboard of the bed in the ship's stateroom, and the other is in the study, on the wall above the desk. The runes are invisible, and both runes need to be deactivated to access the safe room. Inda knows the pass phrase for deactivating both runes: "The moon sings a song for the lost."
\item[What's in the Heart.]
Inda has avoided exploring Havock. However, she's seen a serpentine monster emerge from Havock's heart.
\end{description}

If the characters have no way to see the invisible ward runes, Inda lends the party a \textit{Lantern of Revealing}.

\subparagraph{Treasure}
Inda has a \textit{Flame Tongue shortsword} and a \textit{Ring of Evasion} on her bedside table next to the \textit{Lantern of Revealing}. If asked, Inda is reluctant to give away the shortsword or ring but can be convinced with a successful DC 22 Charisma (Persuasion) check. Inda has a key to the strongbox in area Z13 and knows where the secret compartment holding a \textit{Spell Scroll} of \textit{Flame Strike} is in area Z17; she mentions these to the characters if they help her with her prosthetic wing in area Z16 or help her reach her crew in the stern segment.

As an action, a character who can reach the bedside table can try to steal one of the magic items atop it without being seen by Inda, doing so with a successful DC 19 Dexterity (Sleight of Hand) check. On a failed check, Inda sees the attempt and immediately becomes hostile.

\subsubsection{Z16: Navigation Room}
\begin{DndReadAloud}{}
A chunk of mast has pierced the doors to the corridor and fills one corner of the room. A large bronze sphere hovers in the center of the room, spinning idly. Splayed on the table below is a wing-shaped contraption.

This room formerly served as the Lambent Zenith's navigation room. It now functions as Inda's workshop.

\end{DndReadAloud}

\subparagraph{Inda's Prosthetic Wing}
Inda's prosthetic wing sits on the center table. Inda has been trying to repair it, but while she excels at theoretical studies, she is inexperienced with practical engineering.

Characters can attempt to help Inda repair the prosthetic wing. A character must first succeed on a DC 15 Intelligence (Investigation) or Wisdom (Survival) check to know which joints must be reconnected. Then, the character must succeed on a DC 17 Dexterity (Sleight of Hand) check to reconnect the joints; a character who has proficiency with tinker's tools makes this check with advantage.

Once a successful Dexterity (Sleight of Hand) check has been made, Inda knows she can tinker with the wing enough that it will eventually work.

Inda is grateful to the characters for their help fixing the broken wing, even if their efforts are unsuccessful. If the characters escort her to the stern segment of the wreck, she vouches for the characters to Figaro in area Z5.

\subsubsection{Z17: Bridge}
\begin{DndReadAloud}{}
This room ends in a wall of Havock's flesh, and glass shards from the shattered window litter the floor. An armchair sits on a platform toward the back of the room; below it lies a silver-wrought chair, toppled over and broken in half.

\end{DndReadAloud}

The armchair served as the captain's chair. The broken chair on the ground was the ship's helm but is now irrevocably destroyed.

\subparagraph{Treasure}
A character who examines the captain's chair and succeeds on a DC 15 Wisdom (Perception) check finds a secret compartment in the right armrest. The compartment holds a bronze tube containing a \textit{Spell Scroll} of \textit{Flame Strike}.

\subsubsection{Z18: Forward Cargo Hold}
\begin{DndReadAloud}{}
Empty crates and barrels litter this deck, and a double door stands at the fore. The doors' surfaces shimmer as if trapped behind a wall of translucent silver light.

\end{DndReadAloud}

The door to the safe room is sealed by the ship's emergency security wards (see the "Accessing the Rod Piece" section earlier in this chapter). Inspecting the safe room via a \textit{Detect Magic} spell or a similar effect reveals a strong aura of abjuration magic, as well as two tethers stretching toward the stern segment and the starboard segment, respectively. These tethers indicate the links between the wards and the runes in these respective locations.

\subsubsection{Z19: Safe Room}
\begin{DndReadAloud}{}
An overwhelmingly foul stench fills this destroyed and empty room. Where the prow should be is a gaping hole leading into the putrid core of Havock's heart.

\end{DndReadAloud}

The rod piece was consumed by the \textbf{hertilod} lurking within the heart.

\subparagraph{Heart Entryway}
The hole in this area leads into the heart of Havock, as shown on map 3.2 and described in the following section. To retrieve the rod piece, the characters must venture into the heart and confront the monster inside.

\section{Heart of Havock}
Though much of Havock is petrified and crumbling, the god's heart still beats. Over time, the heart spawned a parasitic monstrosity known as a hertilod, which feasts on the god's residual divine power and terrorizes hapless astral travelers who stumble on Havock's corpse. When the \textit{Lambent Zenith} crashed into Havock, the god's heart began absorbing pieces of the ship's prow, and the hertilod gorged itself on the spelljamming vessel's detritus—including the piece of the \textit{Rod of Seven Parts}.

\subsection{General Features}
The areas of the heart of Havock have the following notable features:


\begin{description}
\item[Foul Air.]
The air throughout the heart is foul. Any creature that breathes the foul air has the poisoned condition until it breathes fresh air again.
\item[Lighting.]
All areas inside the heart are dimly lit by floating orbs of pinkish bioluminescence.
\item[Walls, Ceiling, and Floor.]
The walls, ceiling, and floor of the heart are made of purple flesh. Parts of this flesh have been petrified, hardening into gray crystal, while other parts remain soft and spongy. The heart beats erratically, sending a dull pulse echoing through its chambers every few minutes. The ceilings in the heart are all 50 feet high.
\end{description}

\subsection{Areas of the Heart}
The areas of the heart are keyed to map 3.2.

\subsubsection{A1: Entry Atrium}
\begin{DndReadAloud}{}
Beyond the shattered safe room is an oval chamber with walls of purple flesh dimly lit by pinkish, bioluminescent orbs. The air reeks of death. An opening in one wall leads to another chamber.

\end{DndReadAloud}

A character holding the first rod piece divines that the next piece is deeper in the heart.

\subparagraph{Treasure}
The heart absorbed detritus from the wreckage of the Lambent Zenith's safe room and bridge. A character who searches through the detritus and succeeds on a DC 15 Wisdom (Perception) check uncovers a gilded spyglass worth 1,000 gp.

\subsubsection{A2: Ventricle Chamber}
\begin{DndReadAloud}{}
The walls and floor of this massive chamber alternate between gray, petrified stone and flexing muscle and flesh. Clinging to the ceiling is a serpentine monstrosity, like a skinless snake. Its long, draconic snout drips with venom as it slumbers.

\end{DndReadAloud}

The creature sleeping on the ceiling is the \textbf{hertilod} (see appendix A) that swallowed the next rod piece. A character can move through the area without waking the hertilod by succeeding on a DC 18 Dexterity (Stealth) check. Otherwise, the hertilod wakes when it detects a creature within 30 feet of itself and attacks.

\subparagraph{Retrieving the Rod Piece}
The rod piece is deep within the hertilod's gullet. If the hertilod is forced to regurgitate creatures in its gullet, roll a d6. On a roll of 4 or higher, the Rod of Seven Parts is also regurgitated and lands in an unoccupied space within 10 feet of the hertilod. Alternatively, a creature within the hertilod's gullet can use its action to make a DC 15 Intelligence (Investigation) check, finding the rod piece on a successful check. Once the hertilod is dispatched, a character can easily recover the rod piece from the hertilod's gullet. For more about the \textit{Rod of Seven Parts}, see this book's introduction.

\subparagraph{Caved-In Artery}
A 10-foot-wide petrified artery leads from this chamber. The artery has caved in, blocking that exit from the chamber.

A character who uses an action to dig through the cave-in can make a DC 15 Strength (Athletics) check. On a successful check, the character clears a passage through the cave-in large enough for a Medium creature to squeeze through. Exiting the heart through this passage leads to the area above the \textit{Lambent Zenith}'s prow (area Z19).

\subsubsection{A3: Flooded Atrium}
\begin{DndReadAloud}{}
The air here is humid. Pooled on the floor is black sludge that shimmers like an oil slick.

\end{DndReadAloud}

A character who examines the sludge and succeeds on a DC 13 Intelligence (Arcana) check identifies it as Havock's blood. A creature that touches the blood immediately takes 10 (3d6) force damage. For each minute a creature spends in contact with the blood, the creature takes this damage again.

\subparagraph{Caved-In Artery}
A 10-foot-wide petrified artery leads from this chamber. The artery has caved in, blocking that exit from the chamber. Characters can clear that exit as described in area A2.

\section{Next Steps}
Once the characters have successfully retrieved the second piece of the \textit{Rod of Seven Parts} from the hertilod, they can return to the portal to Sigil that floats above Havock. After the rod piece is removed from Havock's vicinity, teleportation magic within the wreck segments functions normally.

\chapter{The Ruined Colossus}
Searching for the Rod of Seven Parts of the \textit{Rod of Seven Parts} takes the characters to the continent of Khorvaire on Eberron. The scars of the Last War are still fresh on Khorvaire. Nowhere is this truer than in the Mournland, a once-glorious nation destroyed during the apocalyptic Day of Mourning. Here, the characters must search for the next piece of the rod.

Curtains of thick, gray mist blanket the Mournland. Sentient constructs, desperate adventurers, and ravenous mutated monsters roam the blasted landscape. The characters must navigate this wasteland in search of the rod piece, ultimately discovering that the piece is located inside the remains of an enormous, highly advanced, bipedal war machine called a colossus.

\section{Running This Chapter}
This chapter begins after the characters retrieve the second piece of the \textit{Rod of Seven Parts}. The characters can knit the pieces together or keep them separate, as described in the introduction. Regardless, when a character holds the second piece of the rod, they intuitively know that the next piece is located somewhere on the continent of Khorvaire on the world of Eberron.

Specifically, the piece lies in the Mournland, on the slopes of Mount Ironrot. If the characters ask the Wizards Three in the Sigil sanctum why it's unclear exactly where the third piece is located, Alustriel shares the information in the "Exploring Mount Ironrot" section later in this chapter.

This chapter first describes Mount Ironrot, including locations and encounters awaiting the characters. The chapter then details the ruins of a colossus called Landro, whose dangers the characters must survive to claim the third piece of the \textit{Rod of Seven Parts}.

\subsection{Character Advancement}
The characters should be 13th level when this chapter begins. The characters gain a level after they retrieve the third piece of the \textit{Rod of Seven Parts} from the graymatter engine of Landro.

\subsection{Power of Secrets}
The characters can learn two secrets in this chapter that are applicable to the rules in "The Power of Secrets" section in this book's introduction:


\begin{description}
\item[Mercy's Secret.]
Mercy is the leader of a band of warforged pilgrims. Mercy was separated from their best friend, the warforged Filch, after the two were separated after the Day of Mourning. They feel guilty that they haven't searched for Filch. Characters can learn Mercy's secret in the "Warforged Pilgrims" section later in this chapter.
\item[Kalyth's Secret.]
Kalyth, the leader of a band of veterans, lost valuable magic items and money that could have prevented the current financial strain she and her allies are under. Characters can learn Kalyth's secret in the "Veterans' Camp" section later in this chapter.
\end{description}

\subsection{Third Rod Piece}
The third piece of the \textit{Rod of Seven Parts} is found in area L28 within Landro. For more information about the rod and the spell this piece allows its wielder to cast, see this book's introduction.

\section{Doorway to Eberron}
Before seeking the third rod piece, the characters can rest and prepare in the sanctum in Sigil. By doing some research in Sigil or conversing with the Wizards Three, the characters can learn the following about their next destination:


\begin{description}
\item[Ruins of War.]
The Mournland is a ruined wasteland. Before the Last War—a global conflict waged across the continent of Khorvaire for over a century—the Mournland was the nation of Cyre. An apocalyptic event called the Day of Mourning destroyed Cyre and transformed the region. Remnants of the Last War's battlefields, such as the fire-spewing war machines known as colossi, lie strewn about the Mournland, and warforged wander the land.
\item[Warforged.]
Warforged are common on Khorvaire. These Constructs are formed from wood and steel, then magically imbued with life and sentience. Warforged were originally created to fight in the Last War, though many survived that conflict and now try to understand their place in the world.
\item[Unreliable Magic.]
In the Mournland, magic doesn't always function as intended, particularly teleportation and divination magic. See the "Regional Effects" section later in this chapter for more.
\end{description}

\section{An Important Detail}
Before the characters leave Sigil for Mount Ironrot, the Wizards Three share their concerns that this rod piece will be particularly challenging to find since the Mournland is difficult to navigate. The wizards suspect that the rod piece is hidden in one of the dozens of ruined colossi scattered across Mount Ironrot (as shown on map 4.1).

Instead of searching each colossus individually, the wizards suggest augmenting the rod's divining power with local magic that can penetrate the Mournland's impeded navigation. The characters can do this by tuning their second rod piece to a working \textit{Docent} from inside one of the fallen colossi. A \textit{Docent} is a small, sentient metal sphere that employs its sentience and magic on behalf of an attuned warforged; the full description of a \textit{Docent} is found in the "Finding a Docent" section later in this chapter. The wizards correctly postulate that a \textit{Docent}'s magic will stabilize and augment the rod's divining powers so it will directly point the way to the third piece.

\subsection{About Docents}
Long ago, colossi were powered and controlled by Docent. The techniques and tools to create \textit{Docents} were lost, but working \textit{Docents} can sometimes still be found in colossi. The Wizards Three suggest the characters find a working \textit{Docent} and tune it to the rod piece.

\section{Mount Ironrot}
When the characters are ready, they can step through the doorway in Sigil. They arrive outside a jet-black, glass oval portal near the slopes of Mount Ironrot.

\subsection{Running Mount Ironrot}
As the characters search for a \textit{Docent}, they encounter the Mournland denizens described in the subsequent "Mount Ironrot Encounters" section. Run an encounter each time the characters travel from one location to another. All three encounters should occur before the characters proceed to the "Landro" section.

From their expeditions and these encounters, the characters learn of a working \textit{Docent} in a ruined village called Ialos. They also learn that this \textit{Docent} is central to a conflict between rival groups of Mournland scavengers: a band of veterans and a community of warforged pilgrims.

\subsection{Exploring Mount Ironrot}
Use the following rules when the characters travel around Mount Ironrot.

\subsubsection{Regional Effects}
Mount Ironrot is affected by these environmental phenomena:


\begin{description}
\item[Impeded Navigation.]
In outdoor areas, creatures have disadvantage on Wisdom (Survival) checks made to navigate.
\item[Lightly Obscured.]
Mist perpetually shrouds outdoor areas, rendering them \textit{lightly obscured}.
\item[Muddled Magic.]
Creatures and objects within 1 mile of Mount Ironrot can't be perceived through magical scrying sensors. Spells or magical effects that would reveal a creature's or object's location fail while that creature or object is within 1 mile of the mountain. The third piece of the \textit{Rod of Seven Parts} points only to Mount Ironrot without precision.
\end{description}

\subsubsection{Traveling}
Distances between locations anywhere in the Mournland are indeterminate and ever-changing. Map 4.1 shows the approximate positions of locations around Mount Ironrot relative to one another. The map doesn't include a scale due to this effect.

Whenever the characters travel from one location to another, roll 2d4 to determine the number of hours it takes the party to reach their destination. Have the players designate one party member to make a DC 15 Wisdom (Survival) check on behalf of the group each time the characters travel. On a failed check, double the number of hours the party must travel to reach their destination.

\subsubsection{Dangers}
There is no safe food to forage in the Mournland, and monsters are everywhere. If the party lingers too long between locations, introduce random encounters with creatures such as blazebears (see appendix A) or blade scouting parties (see the "Blade Scouting Party" section later in this chapter).

\subsection{Mount Ironrot Locations}
The following locations are keyed to map 4.1.

\subsubsection{Oval Portal}
\begin{DndReadAloud}{}
A twenty-foot-tall, vertical, black oval made of glass sits on the muck atop a low hill. The oval reflects the surrounding mist-shrouded landscape with supernatural clarity. The oval looks like a giant mirror, but it functions as a door between here and Sigil.

\end{DndReadAloud}

A \textit{Detect Magic} spell reveals an aura of conjuration magic emanating from the oval. This oval served as a conduit for powerful teleportation magic in Cyre, but the Day of Mourning rendered its magic unstable. The oval serves as a portal to Sigil while the characters carry the second piece of the rod.

If a creature tries to teleport to the Mournland using a \textit{Teleport} spell, a \textit{Plane Shift} spell, or similar magic, the spell's caster must make a DC 15 Constitution saving throw. On a failed save, the spell fails and has no effect. On a successful save, the spell works, but teleported creatures arrive in an empty space near this oval, regardless of their intended destination in the Mournland.

A creature that studies the reflection of the Mournland in the oval clearly sees the ruined colossi, the veterans' camp, and Ialos (see map 4.1 and the sections below). Otherwise, the view from this hill is limited, offering a glimpse of only the nearest ruined colossus.

\subsubsection{Ruined Colossi}
Each time the characters investigate one of the colossi shown on map 4.1, read or paraphrase the following:

\begin{DndReadAloud}{}
What was once a massive, bipedal war machine made of stone, metal, and wood now lies in ruin. Fragments of the mechanical titan lie scattered about, each piece as big as a barn.

\end{DndReadAloud}

These colossi are smaller than Landro and are easily accessible, making the task of searching them relatively quick. For this reason, no maps of these ruined colossi are provided. You can simply narrate the characters' foray into each colossus.

It takes 1d4 hours for the characters to search a colossus's wreckage for a \textit{Docent}. At the end of this time, have one character make a DC 17 Intelligence (Investigation) check for the group. On a failed check, the characters don't find a \textit{Docent}, but they can't be sure they didn't simply miss it; the characters can spend another 1d4 hours to reattempt the check. On a successful check, the characters not only determine that the colossus's \textit{Docent} is definitely gone, but also uncover one of the following clues (determined by you):


\begin{description}
\item[Deep Tracks.]
Deep, angular boot prints belonging to unusually heavy, bipedal creatures lead to and from the wreckage. Examining multiple colossi and triangulating the boot prints' directions (no check required) points to the village of Ialos.
\item[Second Party.]
Smaller, fresher boot prints made by lighter bipeds meander aimlessly about the wreckage. It's impossible to discern the aim of these lighter bipeds from these tracks alone.
\item[Skilled Extractor.]
The Docent was removed with surgical precision, likely by someone who knew exactly what they were looking for.
\end{description}

\subsubsection{Landro}
This unusually large colossus contains the third piece of the \textit{Rod of Seven Parts}. Landro is described later in this chapter.

\subsubsection{Veterans' Camp}
\begin{DndReadAloud}{}
A crumbling stone bridge crosses a river of indigo water. Beneath one side of the bridge, soiled leather tents and bedrolls are arranged around a firepit.

\end{DndReadAloud}

This campsite belongs to the ex-soldier Kalyth and her two companions (see the "Cyran Veterans" section later in this chapter). The veterans rest here between expeditions around Mount Ironrot. If the characters meet the veterans here and aren't hostile toward them, Kalyth offers to share her group's meager food supplies with the party.

\subparagraph{Kalyth's Secret}
If the characters are friendly with Kalyth, she asks to speak with them privately. In a quiet spot, Kalyth admits that in the aftermath of the Day of Mourning, she lost several valuable magic items and a purse full of gold that she could've used to stave off the veterans' current financial plight. The items are long gone, and Kalyth fears that the other veterans would reject her—or worse—if they found out she once had the money they desperately need and lost it.

Regardless of the characters' reaction to this revelation, learning it counts as a secret for the purposes of the Power of Secrets rules in this book's introduction.

\subparagraph{Treasure}
A character who searches the area and succeeds on a DC 14 Wisdom (Perception) check spots a loose brick among the bridge's remains. Behind the brick are a Potion of Greater Healing and a moldy hunk of bread—part of the veterans' provisions.

\subsubsection{Ialos}
\begin{DndReadAloud}{}
A windmill towers above this ruined village. The windmill's sides and blades are reinforced with badly rusted steel plates, and the attached wooden outbuilding is in shambles. Muddy roads separate the mill from the ruined buildings. In the center of a nearby intersection stands an eroded stone statue of a humanoid draped in colorful scraps of cloth.

\end{DndReadAloud}

Mournland-roaming pilgrims (see the "Warforged Pilgrims" section later in this chapter) turned this ruined mill into an ossuary for fallen warforged.

From this base, pilgrims venture out to find relics or lost warforged. At such times, 2d4 pilgrims (use the \textbf{warforged warrior} stat block in appendix A) remain in Ialos to guard the mill or patrol its surroundings. When no expedition is underway, thirteen pilgrims (including their leader, Mercy, who also uses the \textbf{warforged warrior} stat block) occupy Ialos. The pilgrims defend their home against attackers, but they flee when they're outmatched.

The following locations are keyed to map 4.2:


\begin{description}
\item[Founder's Statue.]
Colorful scraps of fabric drape over the shoulders and outstretched arm of this statue. The warforged pilgrims decorated the statue of Ialos's founder as a gesture of respect; they are grateful to the village's original builders for creating this space they can use to mourn.
\item[Mill.]
The millstones have been removed from this mill, leaving an empty space. Warforged pilgrims use the mill to plan expeditions.
\item[Ossuary.]
Bodies of deceased warforged line this cleared-out basement. Each body is carefully secured to the wall in a dignified pose: back straight, arms crossed, and head slightly lowered.
\item[Ruins.]
Except for the mill, every building in Ialos was destroyed on the Day of Mourning. Broken stone and rotten timber lie strewn about the village.
\item[Storeroom.]
Objects the warforged pilgrims deem sacred line the shelves of this storeroom (see "Treasure" below).
\end{description}

\subparagraph{Treasure}
The storeroom contains curiosities and trinkets salvaged from Mournland ruins. Some of these items had magical properties until the Day of Mourning destroyed their magic. The storeroom's notable treasures include the following:


\begin{description}
\item[Art Objects.]
Clean tapestries, intact sculptures, and leather-wrapped oil paintings are carefully piled in a corner of the storeroom. Altogether, these relics of old Cyre are worth 2,000 gp.
\item[Docents.]
Half a dozen metal orbs line one shelf, but only one is a functional \textit{Docent} (see the "Docent" section later in this chapter). These items were taken from the area's ruined colossi. A character can touch a piece of the \textit{Rod of Seven Parts} to the working Docent to determine the location of the third piece (see the "Finding a Docent" section later in this chapter).
\item[Magic Ring.]
An iron ring etched with angelic feathers hangs from a silver chain on the wall. This is a \textit{Ring of Feather Falling}.
\end{description}

\subsection{Mount Ironrot Encounters}
These encounters should occur, in the order presented, as the characters travel from one location on map 4.1 to another. Each encounter occurs only once.

\subsubsection{Warforged Pilgrims}
The characters spot a band of five bipeds walking single file across a distant ridge. These are warforged pilgrims. Each uses the \textbf{warforged warrior} stat block (see appendix A).

The pilgrims are peaceful folk who wander the Mournland in search of friendly travelers and artifacts important to warforged. A purple-hued ex-soldier named Mercy is the group's leader. Mercy is initially indifferent to the characters, but the leader responds positively to any aid the characters offer, such as advice or adventuring supplies.

\subsubsection{Pilgrims' Problems}
At first, Mercy speaks vaguely about the pilgrims' goals. If the characters offer aid or helpful advice, Mercy talks about the pilgrims' expeditions to ruined colossi. Mercy then tells the characters that the pilgrims have been extracting Docent from these ruins, and that the pilgrims consider the items sacred. Mercy also talks about the pilgrims' recent run-ins with rival salvagers. Mercy mentions two other groups in particular: a band of Cyran veterans that watches Mercy's group from the shadows, and a gang of openly hostile warforged that the pilgrims avoid at all costs. This latter group, Mercy explains, is loyal to the Lord of Blades, a local warmonger who conscripts warforged to his bloody cause of wiping out anyone who opposes him.

If the characters offer to help the pilgrims, Mercy suggests they meet at the pilgrims' base in Ialos (see "Ialos" earlier in this chapter). Mercy and the pilgrims happily travel with the characters, though they hide if violence breaks out, emerging only when it's safe. If the characters wish to attend to other business before meeting the pilgrims at Ialos, Mercy draws a rough map in the mud to show the way. The pilgrims head straight to Ialos and are there when the characters arrive.

\subparagraph{Mercy's Secret}
If the characters befriend Mercy, they sense during their conversations with the warforged that something weighs heavily on Mercy. If the characters ask the warforged about their melancholy, Mercy asks to speak with the characters privately. In private, Mercy admits that their best friend has been missing in the Mournland since the Day of Mourning. Weighted with the responsibility of leadership, Mercy has never searched for this friend, a blue-and-red warforged named Filch. If the characters offer to search for Filch, they find the ex-soldier in area L4 of Landro.

Regardless of the characters' reaction to this revelation or whether they offer to look for Filch, learning it counts as a secret for the purposes of the Power of Secrets rules in this book's introduction.

\subsubsection{Cyran Veterans}
The characters encounter three lawful neutral veterans wandering through the mist. Once it's clear the characters aren't monsters or violent marauders, the wanderers introduce themselves as the Turquoise Spear, a small platoon of Cyran ex-soldiers that scours the Mournland for relics.

An orc artillerist named Kalyth speaks for her companions, two halfling soldiers named Dortle-Lynn and Grezan. Kalyth explains that the Turquoise Spear recently discovered a stockpile of Docent taken from colossi. She wants to claim the \textit{Docents} to sell in markets outside the Mournland, but she'd happily let the characters borrow a \textit{Docent} beforehand. Unfortunately, Kalyth says, the \textit{Docent} stockpile is in Ialos, a ruined village currently occupied by "rogue warforged." The Turquoise Spear attempted to infiltrate the ruins, but the warforged made it clear they wouldn't give up their stockpile without a fight.

Kalyth doesn't know that the warforged are peaceful pilgrims, nor does she know that the pilgrims regard the \textit{Docents} as sacred. Regardless of these facts, the veteran insists that these treasures rightfully belong to "living, breathing Cyrans."

If the characters agree to help the Turquoise Spear, Kalyth gives the party directions to Ialos (see "Ialos" earlier in this chapter). She adds that the town's most prominent structure is a large metal windmill that glints faintly in the mist.

\subsubsection{Blade Scouting Party}
As they pass between scorched thickets, the characters are ambushed by violent followers of a bloodthirsty warmonger called the Lord of Blades. These scouts have orders to search the region for a suitable colossus to turn into a new outpost, but they'll happily make time to slay a few people in their way. The scouting party includes two blade scouts and one \textbf{blade lieutenant} (see appendix A for both stat blocks).

The blades fight until destroyed. They remain tight-lipped if captured, but if a character demands information and succeeds on a DC 15 Charisma (Intimidation) check, the lieutenant says his party's commander is a powerful warforged headhunter named Glaive. "She'll make mincemeat out of you, interlopers!" the lieutenant scoffs. The blades can point the party to ruined colossi in the area but otherwise have no useful information.

\section{Finding a Docent}
The party's first task in the Mournland is to find a working \textit{Docent} to which they can tune the second piece of the \textit{Rod of Seven Parts}. In the course of their exploration, they learn that the only operational \textit{Docent} in the region is in a storeroom in Ialos.

\subsection{The Dispute}
By the time the characters reach Ialos, they should have met the pilgrims and the Cyran veterans. Both sides might ask the characters to intervene and resolve their dispute. Here is a breakdown of the situation:


\begin{description}
\item[What Kalyth Wants.]
The Cyran veterans want to sell treasures from their ruined homeland in markets outside the Mournland. With the money, Kalyth says, she and her comrades can fund new, better lives for themselves.
\item[What Mercy Wants.]
Mercy regards Docents (and any other items intended for warforged use) as sacred. The pilgrims believe that such objects belong in the hands of warforged who fought in the Last War.
\item[What They'll Settle For.]
Ultimately, Kalyth's group wants money, and Mercy's group wants respect. If both parties get what they want, the dispute ends.
\end{description}

\subsubsection{Settling the Dispute}
Mercy doesn't understand that the Cyran veterans' true motivation is money, nor does Mercy recognize the value of the items in the pilgrims' possession. Though reluctant to bargain with the veterans, whom the pilgrims consider warmongers, Mercy can be persuaded to negotiate if the characters act as intermediaries.

If the characters explain the veterans' financial need and convince Mercy that the veterans mean no harm, Mercy offers to give the characters the art objects kept in the pilgrims' storeroom. If a character succeeds on a DC 14 Charisma (Persuasion) check during this discussion, Mercy also offers the characters the \textit{Ring of Feather Falling} in the storeroom.

Kalyth doesn't realize that the pilgrims' have valuable art objects or a magic ring. If told about these treasures, Kalyth says the Cyran veterans would gladly take either option instead of the Docent. If Mercy offers the ring as well, Kalyth insists the characters keep either the art or the ring, whichever they'd prefer, in thanks for brokering a peaceful resolution.

\subsubsection{Helpful Information}
If the resolution allows the warforged pilgrims to keep their \textit{Docents} and no one is harmed in the dispute, once it's clear the characters are looking for a piece of the \textit{Rod of Seven Parts}, Mercy tells the characters the following about Landro's unique power source:

\begin{DndReadAloud}{}
"I've heard tales of an unusual colossus called Landro powered not by a \textit{Docent}, but by a device called a graymatter engine that incorporated part of an artifact—perhaps that's what you seek? The graymatter engine was said to generate a magical barrier around Landro, making it impossible to enter Landro except where the barrier has already been broken."

\end{DndReadAloud}

If the characters negotiate a deal that the Turquoise Spear finds agreeable, Kalyth tells the party the following about the third rod piece:

\begin{DndReadAloud}{}
"Rumor has it there's a colossus powered by part of an artifact. Supposedly, if you remove this piece, the colossus will begin a self-destruct sequence, and there's no way to stop the colossus from detonating in a terrible conflagration."

\end{DndReadAloud}

\subsection{Tuning the Rod}
Regardless of which side, if any, the characters take in this conflict, the characters must pair the second piece of the \textit{Rod of Seven Parts} to the working \textit{Docent} in Ialos to continue the adventure. (A character need not be attuned to the \textit{Docent} to pair the \textit{Docent} to the rod piece.)

When a character touches the rod pieceto the working \textit{Docent}, the piece emits a loud, metallic ping. Thereafter, the piece points in the direction of the third rod piece inside Landro. If the party travels in that direction, the characters eventually reach the ruins of Landro (see the "Landro" section below).

\subsubsection{Docent}
\textit{Wondrous Item, Rare (Requires Attunement by a Warforged})

A \textit{Docent} is a 2-inch-diameter metal sphere studded with dragonshards. To attune to a \textit{Docent}, you must embed the item somewhere on your body, such as your chest or your eye socket.

\subparagraph{Sentience}
A Docent is a sentient item of any alignment with an Intelligence of 16, a Wisdom of 14, and a Charisma of 14. It perceives the world through your senses. It communicates telepathically with you and can speak, read, and understand any language it knows (see "Random Properties" below).

\subparagraph{Life Support}
Whenever you end your turn with 0 hit points, the Docent can make a Wisdom (Medicine) check with a +6 bonus. If this check succeeds, the Docent stabilizes you.

\subparagraph{Random Properties}
A Docent has the following properties:


\begin{description}
\item[Languages.]
The Docent knows Common, Giant, and 1d4 additional languages chosen by the DM. If a Docent knows fewer than six languages, it can learn a new language after it hears or reads the language through your senses.
\item[Skills.]
The Docent has a +7 bonus to ability checks using one of the following skills (roll a d4): (1) Arcana, (2) History, (3) Investigation, or (4) Nature.
\item[Spells.]
The Docent knows one of the following spells and can cast it at will, requiring no spell components (roll a d6): (1–2) \textit{Detect Evil and Good} or (3–6) \textit{Detect Magic}. The Docent decides when to cast the spell.
\end{description}

\subparagraph{Personality}
A Docent is designed to advise and assist the warforged it's attached to, including acting as a translator. The Docent's properties are under its control, and if you have a bad relationship with your Docent, it might refuse to assist you.

\section{Landro}
In the Last War's climactic final years, Cyran artificers crafted massive war machines called colossi. One of the nation's greatest specimens was deployed to the battlefield just as the Day of Mourning swept over Cyre and destroyed the nation. Instead of teleporting to the front lines, this colossus, called Landro, appeared miles away, half-buried in the face of a soaring mountain. Landro's crew was killed instantly.

Landro now looms empty and foreboding over Mount Ironrot's eastern flank. The experimental magic that powers Landro—an eldritch machine called the graymatter engine, fueled by a piece of the \textit{Rod of Seven Parts}—continues to function, creating a magical force field around the colossus and altering the gravity within it. In addition to the aberrant monsters and restless Undead that haunt Landro's remains, another faction is determined to plumb the ruins' depths: violent, zealous warriors loyal to the Lord of Blades.

\subsection{Lord of Blades' Strike Squad}
While the characters explore Landro, they're pursued by a warforged named \textbf{Glaive} (see the accompanying stat block). Glaive is a high-ranking commander loyal to the Lord of Blades. She has orders to find a ruined colossus suitable to transform into an outpost for the blades, and Landro fits this description perfectly. When she becomes aware of the party's presence in Landro, Glaive relishes the opportunity to destroy these "meddlesome interlopers," as she calls them.

Glaive is accompanied by two blade scouts (see appendix A) named Rack and Crunch. The three blades prefer to split up and sneak around Landro and the adjoining caves. They take turns using hit-and-run tactics to harry the characters, creating a sense that more blades lurk around every corner. As long as Glaive is around, the characters can't finish an uninterrupted long rest inside Landro or in the caves of Mount Ironrot.

It's up to you where the characters encounter this strike squad, though Landro's location descriptions provide suggestions. If you don't want to ambush the characters with this encounter, it should occur in area L27. Glaive is devoted to her mission and fights the characters until destroyed, but her subordinates are less motivated. If Glaive is defeated, Rack and Crunch flee.

\subsubsection{Glaive}
Since the Mournland's earliest days, Glaive has wielded her namesake weapon in service to the Lord of Blades' bloody conquest of the blighted region. Among the blades, Glaive is best known for her talent at self-modification. "Glaive" is etched onto the back of her neck-plate. It is the only name Glaive has ever known, and fellow blades don't dare call her anything else. Mournland adventurers, however, refer to the terrifying commander by another name: Kill Switch.

\subsection{Entering Landro}
The invisible barrier described in the "Landro Features" section below makes it impossible to enter the colossus via typical means. The barrier is damaged in three places, allowing ingress via paths along Mount Ironrot into areas L5, L11, and L16 of the colossus. Map 4.3 includes a side view of Landro and the paths that lead into the colossus's interior.

The easiest way to enter the colossus is through the mountain caves that lead to area L5. The characters can also search higher up the mountain along outdoor paths for the entrances to caverns that lead to areas L11 and L16. Searching for area L11 takes 2d4 hours, and during the characters' search, a hungry \textbf{roc} ambushes them. Searching for area L16 takes 3d4 hours, and during the characters' search, they stumble on a nest with two rocs and one roc egg the height of a human. The angry roc parents attack the characters, though the rocs don't pursue if the characters flee.

\subsection{Landro Features}
The following features are common to all areas of Landro unless otherwise noted.

\subsubsection{Antigravity Magic}
The piece of the \textit{Rod of Seven Parts} linked to the graymatter engine creates special gravity-altering effects in parts of Landro. The areas affected by these effects are denoted on map 4.3 as the swirled features in areas L5, L6, L7, L8, and L9. These have the following effects:


\begin{description}
\item[Antigravity Wells.]
Landro's ankles, legs, knees, hips, spine, and elbow are filled with magical energy meant to allow the colossus's body to bend and rotate freely. Creatures and objects can enter and exit these permeable "antigravity wells" like any normal space. While inside an antigravity well, a creature is affected by a \textit{Levitate} spell (no saving throw). Regular notches in the walls make it relatively easy for creatures to move up or down a well to reach other areas of the colossus.
\item[Gravity Tiles.]
Landro's upper levels are enchanted with gravity magic that keeps passengers safely rooted to the floor. Creatures can move across the floor in areas L20–L28 normally, no matter the orientation of Landro's upper body.
\end{description}

The effects of the antigravity wells and the gravity tiles end if Landro's shutdown sequence is activated (see the "Shutdown Sequence" section later in this chapter).

\subsubsection{Ceilings}
The height of the ceilings varies considerably in Landro; see individual area descriptions for details. In areas L1–L4 of the cave network, the ceilings are 15 feet high unless otherwise noted.

\subsubsection{Graymatter Fluid}
A steady stream of thin, oily liquid leaks from the graymatter engine and runs into other parts of Landro and the connected caverns. This graymatter fluid is the engine's sensory appendage. The graymatter engine can't move or control this liquid, but it has blindsight within 60 feet of anywhere the graymatter fluid touches.

When a creature looks at a puddle of graymatter fluid, the creature sees a twisted, monstrous version of itself with an inscrutable expression rather than its actual reflection. A creature that ingests any amount of graymatter fluid must succeed on a DC 14 Constitution saving throw or experience the following effects:


\begin{description}
\item[Altered Speech.]
The creature gains telepathy to a range of 30 feet but loses the ability to speak. Any attempts to speak generate only incoherent babbling. This effect ends after the creature finishes a long rest.
\item[Psychic Damage.]
For each ounce of graymatter fluid consumed, the creature takes 11 (2d10) psychic damage. A creature killed by this damage rises as a \textbf{zombie} 1d4 hours after dying.
\end{description}

\subsubsection{Invisible Barrier}
The graymatter engine creates an invisible barrier around the intact parts of Landro. Nothing can penetrate the barrier, which is similar to the effect of a \textit{Wall of Force} spell, though a \textit{Disintegrate} spell can't destroy it. The barrier is broken in compromised sections of Landro (areas L5, L11, and L16).

\subsubsection{Lighting}
Landro's interior areas and the adjoining caves are dark. Area descriptions assume the characters have a light source or some other means of seeing in the dark. During the day, the parts of Landro exposed to the outdoors are shrouded in the same \textit{lightly obscuring} mist as the rest of Mount Ironrot (see "Exploring Mount Ironrot" earlier in this chapter).

\subsubsection{Walls}
As long as the graymatter engine is active, Landro's body—which is made from steel-plated stone blocks—is immune to damage. The adjacent caves are made of naturally formed stone.

\subsection{Landro Locations}
The following locations are keyed to map 4.3.

\subsubsection{L1: Cave Entrance}
\begin{DndReadAloud}{}
The jagged cave entrance in the side of the mountain resembles a yawning maw. Warped stalactites droop at odd angles like monstrous fangs. A shallow stream of thin, gray liquid dribbles from the hole. Inside the cave, darkness awaits.

\end{DndReadAloud}

The shallow stream of graymatter fluid originates from the upper levels of Landro.

\subsubsection{L2: Cave Graveyard}
\begin{DndReadAloud}{}
The corpses of recently fallen human soldiers bearing fine weapons and well-oiled armor lie scattered about this musty cavern.

\end{DndReadAloud}

A character who examines the area and succeeds on a DC 12 Intelligence (Investigation) check notices a thick layer of powdery residue has settled across the scene, suggesting the remains have possibly been here for years. Anyone who examined the thin, gray liquid in the cave entrance realizes that the white powder is residue from this liquid.

The corpses are Cyran soldiers who were instantly killed by mysterious magic on the Day of Mourning. The Mournland's magic has preserved their bodies, which a character can determine by examining the corpses and succeeding on a DC 15 Intelligence (Arcana) or Wisdom (Medicine) check.

If a creature picks up the armor or weapons, two flying swords and two suits of \textbf{animated armor} rise from the scattered gear to reclaim the stolen treasure. They attack the thief and the thief's allies, but they don't pursue foes outside this area.

\subsubsection{L3: Cave Shack}
\begin{DndReadAloud}{}
Rickety metal sheets separate the east end of this cave from the rest of the complex. Aggressive growling emanates from behind the sheets.

\end{DndReadAloud}

This area is separated from the rest of the caves by rickety metal sheets scavenged from the debris of Landro's right foot (area L5). The area serves as a shelter for three Brelish soldiers who were transformed by the strange magic of the Day of Mourning. Use the \textbf{fomorian} stat block for each soldier, except they are Large instead of Huge.

The cowardly soldiers make occasional forays into the caves to smash rats or other easy prey. They attack only if they can surprise the characters, perhaps at the narrow tunnel to the north (area L4), or if the characters are already in a fight with Glaive's squad. The soldiers know nothing of Landro other than to avoid drinking the graymatter fluid.

\subsubsection{L4: Bottleneck}
If the characters haven't run into significant challenges yet, Glaive and one of her blade scout companions attack the characters from both ends of this narrow tunnel. After all combatants have taken a turn, the attackers flee into Landro and hide for another surprise attack.

Regardless of whether the characters fight the strike squad, if they search this area, they find a comatose \textbf{warforged warrior} (see appendix A) in the northeast corner of this area. A character who tries to rouse the warrior and succeeds at a DC 18 Charisma (Persuasion) check wakes the warrior, who identifies themself as Filch. Filch wandered in here shortly after the Day of Mourning looking for their best friend, the warforged Mercy (see the "Warforged Pilgrims" section earlier in this chapter). The stress of the situation caused Filch's comatose state. Mercy and Filch would be happy to be reunited.

\subsubsection{L5: Right Foot}
\begin{DndReadAloud}{}
Aside from a narrow passage, the 10-foot-tall toe section of the colossus's foot has collapsed. Oversize, skull-shaped mushrooms pockmark the heap of debris that marks the rough division between the colossus's interior and the adjacent cave. At the chamber's far end, motes of dust and debris float in a cylinder of glowing, green light that stretches from floor to ceiling. The cylinder of light continues upward into Landro's leg.

\end{DndReadAloud}

Within the mushroom patch are two shriekers. The shriekers' high-pitched screams might attract the Brelish soldiers from area L3 or alert one of Glaive's scouts to the characters' location.

The 20-foot-diameter cylinder of light extending up at the back of this foot is an antigravity well (see the "Antigravity Wells" section earlier in this chapter). The well connects to Landro's right hip (see area L8).

\subsubsection{L6: Left Foot}
\begin{DndReadAloud}{}
The stone door at the toe end of this hollow foot has broken in several parts and is jammed shut. Through the cracks in the door, you can see the solid stone of the surrounding mountain.

\end{DndReadAloud}

A character who searches the area and succeeds on a DC 11 Wisdom (Perception) check spots a dusty, lusterless gemstone lying in a corner. This is an id crystal, which the graymatter engine's consciousness can be transferred into (see the "Crystal Companion" section later in this chapter).

\subsubsection{L7: Leg Shafts}
Each of Landro's 130-foot-tall legs contains two 20-foot-diameter, cylindrical antigravity wells, one in its calf and one in its thigh, connected at the knee by a 20-foot-diameter antigravity sphere.

\subsubsection{L8: Armory}
The two 20-foot-diameter antigravity spheres in this 30-foot-tall chamber serve as Landro's hip sockets. Gaps in the floor and ceiling of each sphere grant access to the colossus's legs (area L7) and abdomen (areas L9–L11), respectively. A 10-foot-wide antigravity well cylinder serves as a spine, connecting to area L9.

\subparagraph{Treasure}
A weapon rack on the wall opposite Landro's spine contains two Halberd and one \textit{hand crossbow} with a case of 20 Crossbow Bolts (20).

\subsubsection{L9: Blocked Antigravity Well}
Rubble in area L18 blocks travel through Landro's spinal antigravity well above this floor. Otherwise, this area is empty.

\subsubsection{L10: Artificer Quarters}
\begin{DndReadAloud}{}
Absolute silence fills this tranquil sleeping area. The walls are lined with wooden bunk beds draped in green and purple quilts. Motes of dust drift with a peculiar languidness. A woman sleeps in one of the beds.

\end{DndReadAloud}

This 20-foot-high chamber is affected by a persistent \textit{Silence} spell. A \textit{Dispel Magic} spell ends the effect.

\subparagraph{Aura of Slumber}
Lingering in this chamber causes creatures to feel drowsy. A creature that ends its turn in this chamber must succeed on a DC 13 Constitution saving throw or have the unconscious condition for 1 minute. If the saving throw fails by 5 or more, the creature is unconscious for 1 hour instead. The creature wakes up if it takes damage or if another creature uses an action to wake it.

\subparagraph{Sleeping Ghost}
The figure dozing in one of the beds is the spirit of Alamar-Vatashi, a soldier who overslept on the Day of Mourning. In life, Alamar-Vatashi was a kalashtar, a Humanoid bound to a dream-spirit called a quori. The Day of Mourning caused her to remain asleep in this chamber. She is a neutral good \textbf{ghost} who can speak Common; she can also communicate telepathically within a range of 30 feet.

If Alamar-Vatashi awakens and becomes aware of the characters, she telepathically asks why they aren't sleeping and yawns deeply. If questioned, she says she's too tired to talk, but she offers to share her dreams with the characters if they wish. A character who touches Alamar-Vatashi's incorporeal palm experiences a vivid vision of the day Landro came to this mountain on the Day of Mourning. This touch doesn't damage the character. Read or paraphrase the following to describe the vision:

\begin{DndReadAloud}{}
As if through the eyes of a soaring bird, you see a titanic stone statue standing amid a circle of glowing sigils. A dozen robed figures, as small as ants, stand around the circle and murmur arcane words, their heads bowed in concentration. Impossibly, the stone colossus begins to levitate.

As silvery light swirls around the colossus, the ground quakes, breaking the mages' focus. The colossus is buried in crumbling stone, surrounded by shifting mountains and sickly gray clouds. Years pass in rapid succession while the colossus stands like a stone sentinel, stuck inside the craggy mountainside.

\end{DndReadAloud}

The vision lasts only a moment but feels like it lasts hours to the creature experiencing the vision. A character who receives this vision gains the benefits of a short rest. No character can receive the vision more than once.

\subsubsection{L11: Workshop}
\begin{DndReadAloud}{}
Fine tinkering tools, woodcarving equipment, and metal plates hang from hooks above a stone table and counters in this crowded workshop. In one corner, a faucet juts from a tall steel vat. The opposite corner has collapsed, allowing access to a tunnel beyond. A rivulet of graymatter fluid oozes across the floor and down the tunnel.

\end{DndReadAloud}

The ceiling here is 20 feet high. Any characters proficient in \textit{tinker's tools} or \textit{woodcarver's tools} can assemble a set of either tools from the items in this room with 1 minute of work.

\subparagraph{Repair Paste}
The steel vat contains magical paste specially formulated for strengthening Constructs. Turning the knob on the attached faucet releases a palm-size dose of paste. A creature can take 1 minute to apply a dose of the paste to a Construct to grant it 10 \textit{temporary hit points}. Once dispensed, a dose of paste loses its magic if not applied within 1 minute. The vat contains three doses of paste.

\subsubsection{L12: Lost Soldier}
\begin{DndReadAloud}{}
The ghostly image of a human Cyran soldier paces nervously at the bottom of this sloping cavern. On seeing you, the figure smiles with relief and waves, shouting, "Please, help me find my comrades!"

\end{DndReadAloud}

The lawful good \textbf{ghost}, who says its name is Chandry, pleads with the characters to help it reunite with its comrades inside Landro.

Suddenly, Chandry cries out with joy and sprints toward the colossus. When the ghost comes within arm's reach of Landro, it blinks out of existence. Chandry then reappears at the bottom of the cave with no memory of the last few moments' events and asks the characters for their help again. Every time Chandry approaches Landro, the ghost disappears and reappears in the same way. If the characters tell Chandry about the fallen soldiers in area L2, the ghost is put to rest and doesn't reappear.

\subsubsection{L13: Dragonshard Pool}
\begin{DndReadAloud}{}
The deep pool of gray liquid in the corner of this chamber glitters oddly. Around the pool, bones and battered armaments litter the cave floor.

\end{DndReadAloud}

Pulverized dragonshards—special gemstones native to Khorvaire—mixed with the graymatter fluid here, creating a pool with strange magical properties.

\subparagraph{Enchanted Pool}
A \textit{Detect Magic} spell reveals an aura of enchantment magic emanating from the pool. Drinking from it exposes a creature to graymatter fluid (see the "Graymatter Fluid" section earlier in this chapter). When a nonmagical item is dipped into the pool, roll on the Dragonshard Pool Effects table to determine what happens to the item. An item can be affected by the pool only once; subsequent exposures have no effect.

The pool has 4 charges. It expends 1 charge each time an item is dipped into it. The pool regains all expended charges daily at dawn.

\begin{DndTable}[header=Dragonshard Pool Effects]{cX}

d6 & Effect\\
1–2 & The item is destroyed.\\
3–4 & The item gains a cosmetic defect such as discoloration, mild warping, or mysterious engravings.\\
5–6 & The item craves more graymatter fluid. Every time the item's wielder comes within 30 feet of graymatter fluid, the item emits a high-pitched hum that lasts until the character moves away from the fluid.\\
\end{DndTable}

\subparagraph{Treasure}
A suit of \textit{chain mail}, two Shortsword, and a \textit{mace} lie on the floor. All are twisted and malformed, as if they were affected by a roll of 3 or 4 on the Dragonshard Pool Effects table.

\subsubsection{L14: Blazebear Den}
\begin{DndReadAloud}{}
A pile of leather scraps, broken bones, and sludgy offal fills a corner of this cave. Nesting in this nauseating pile is a monstrous beast that resembles a bear with three long, fleshy tentacles sprouting from its skull. Each tentacle is topped with a glowing, knobby lump of flesh.

\end{DndReadAloud}

The monster dwelling here is a \textbf{blazebear} (see appendix A). It attacks any creature it sees and fights until destroyed.

\subsubsection{L15: Overhang}
\begin{DndReadAloud}{}
Two tunnels branch off opposite sides of this thirty-foot-high cavern. The tunnels reconvene at an overhang fifteen feet above the eastern passage.

\end{DndReadAloud}

Assuming they haven't both been defeated, one of Glaive's blade scouts (see appendix A) hides on the overhang here while the other scout hides down the eastern tunnel, near area L16. When the characters are within range, the scouts fire their crossbows at a random character.

Any creature atop the overhang has \textit{half cover} from creatures in the tunnel below.

\subsubsection{L16: Ruined Chamber}
This 30-foot-tall chamber is choked with debris and is \textit{difficult terrain}.

\subsubsection{L17: Bridge}
\begin{DndReadAloud}{}
A wood-and-steel captain's chair tops the platform in the center of this room. Slumped in the chair is a skeletal corpse wearing Cyran regalia. The skeleton wears a silver helmet with a sizable dent. A view port made from shimmering crystal overlooks the mists of the Mournland.

\end{DndReadAloud}

As long as Landro's magical barrier is active, the view port's crystal is immune to damage. The ceiling is 30 feet high.

\subparagraph{Broken Control Helmet}
The silver helmet once created a magical link between its wearer and the graymatter engine. A creature that puts its ear to the helmet hears a faint, metallic clanking. A creature that dons the helmet must make a DC 13 Constitution saving throw. On a failed save, the creature has the incapacitated condition for 1 minute. On a successful save, the creature experiences a vision of what happened to Landro on the Day of Mourning (use the read-aloud text of the vision in area L10). The helmet loses this effect if removed from this room.

\subsubsection{L18: Collapsed Antigravity Well}
\begin{DndReadAloud}{}
Rubble chokes this shaft. A steady trickle of gray liquid burbles through the rubble.

\end{DndReadAloud}

Damage to Landro's spine has disabled this antigravity well from this point upward. The rubble blocks access to area L9 below, but the path up to area L20 and beyond is clear.

Without the aid of antigravity magic, a creature can climb the notched walls inside the shaft only by succeeding on a DC 12 Strength (Athletics) check. From this area, the shaft stretches up 100 feet before terminating at the roof of Landro's head (area L28).

\subsubsection{L19: Muster Point}
Three Cyran soldiers in full regalia stand diligently at the arrow slits in this 30-foot-tall gathering hall. When they become aware of the characters, the soldiers turn to the party, revealing ghostly, screaming skulls instead of faces. The soldiers are three hostile wights. Three will-o'-wisps, previously invisible, appear at the start of combat and fight alongside the wights.

\subsubsection{L20: War Room}
\begin{DndReadAloud}{}
This thirty-foot-tall chamber centers on a massive, stone table. Dozens of large maps and scroll cases lie scattered across the table.

\end{DndReadAloud}

The table is a magically animated servant (use the \textbf{stone golem} stat block). It is designed to defend Landro from enemy boarders. As soon as a creature disturbs the table's contents, the table reveals its true nature and attacks.

\subparagraph{Treasure}
The vellum maps depict regions of Cyre that bear no resemblance to the current topography. Though useless to navigators, the maps are worth 2,500 gp to historians, collectors, or Cyran expatriates.

\subsubsection{L21: Shoulders}
Ballistae mounted on wheels point out the windows at each of Landro's shoulders. A total of twelve ballista arrows are scattered on the floors in these areas: five in area L21a and seven in L21b. The ceiling in these areas is 40 feet high.

\subsubsection{L22: Walkway}
A 3-foot-tall stone battlement wraps around this open-air walkway.

\subsubsection{L23: Holding Cells}
\begin{DndReadAloud}{}
Iron bars along either side of this hallway form two cells, each with a locked iron door.

\end{DndReadAloud}

The forearm and upper arm both house pairs of holding cells. Designed to hold prisoners of war, these empty cells never saw use. The ceiling in each room is 15 feet high.

\subsubsection{L24: Intake Chamber}
\begin{DndReadAloud}{}
This small, fifteen-foot-tall chamber features a metal hatch in the floor and a lever attached to a nearby wall.

\end{DndReadAloud}

When the lever on the wall is flipped, the hatch in the floor opens, and a 10-foot-diameter, 30-foot-deep cylindrical antigravity well emanates from the colossus's palm. Any Large or smaller unattended object that enters this antigravity well slowly floats into this chamber. The antigravity well deactivates and the hatch closes when the lever is returned to its original position.

\subsubsection{L25: Vault}
\begin{DndReadAloud}{}
The floor of this ruined, thirty-foot-tall chamber is bare except for a large, wooden chest. The ostentatious chest's dark walnut planks shine, and its iron straps are gilded with a thick sheet of gold leaf. An iron padlock crafted in the shape of a grinning demon hangs from the chest's hasp.

\end{DndReadAloud}

The chest is a \textbf{mimic} in disguise. As soon as a creature touches it, the mimic attacks. Once it has grappled a creature, the mimic attempts to drag its victim outside so it can drop the victim off the edge of the walkway (area L22). Creatures that fall take 45 (13d6) bludgeoning damage.

\subsubsection{L26: Throat}
\begin{DndReadAloud}{}
This twenty-foot-tall cylinder is the equivalent of Landro's throat. Above, a tangle of metal dangles like a mechanical uvula.

\end{DndReadAloud}

If she still lives, Glaive has set a trap here for the characters. She tore out the weapons system from area L27 and moved it to the edge of that chamber so that it points down Landro's throat.

As soon as a character is halfway up this area, Glaive activates the weapon, which ejects a jet of water in a 60-foot line. Any creature in this line must succeed on a DC 15 Dexterity saving throw or take 22 (4d10) bludgeoning damage and fall down to area L18, landing with the prone condition.

Glaive compromised the weapons system when she moved it; once fired, it breaks, losing its magic.

\subsubsection{L27: Weapons System}
\begin{DndReadAloud}{}
The mechanical guts of what must have been a magical cannon were torn out of the platform in this room. To the south, a three-foot-tall ledge is all that separates this chamber from the open air. If not for the dense mist of the Mournland, this window would afford a spectacular view.

\end{DndReadAloud}

In this 15-foot-tall room, which is effectively Landro's mouth, \textbf{Glaive} makes her final stand against the characters if she has survived this long. Her blade scouts (see appendix A), Rack and Crunch, fight alongside her if they're still alive.

\subsubsection{L28: Graymatter Engine}
\begin{DndReadAloud}{}
Dripping pipes and rusty chains hang from the fifteen-foot-high ceiling of this chamber. Deep-red light courses rhythmically through the pipes, which converge on a large object atop a circular dais.

The object resembles an oversize brain made of iron-gray ceramic. Its surface is molded with countless grooves that form mesmerizing patterns. A crack along the ceramic brain's frontal lobe leaks a thin, gray liquid that pools around the dais.

Floating above the brain and scattering light across the room is a small, slender object: a piece of the \textit{Rod of Seven Parts}.

\end{DndReadAloud}

Years after its abandonment, the graymatter engine still churns with arcane power. Crafted from experimental designs based on unfamiliar technology, the graymatter engine converts the power of the third \textit{Rod of Seven Parts} piece into magical effects throughout Landro. The colossus's force field, weapons system, and antigravity wells are all powered by this eldritch machine.

When the characters arrive in this area, the graymatter engine manifests a physical entity to converse with the party. Proceed to the "Graymatter Guardian" section below.

\subparagraph{Hidden Sentinels}
This chamber contains false walls made of thin metal plates. A character investigating the walls can tell that they are hollow by succeeding on a DC 17 Wisdom (Perception) check. Behind these walls stand four defense sentinels bound to the graymatter engine (each uses the \textbf{shield guardian} stat block but is Medium instead of Large).

If a sentinel takes damage or the rod piece is removed from the graymatter engine, the walls around all four sentinels slide into the floor and the sentinels activate, attacking all intruders.

\subparagraph{Retrieving the Rod Piece}
As long as the graymatter engine is active, the rod piece floating above the engine's ceramic shell is held in place with powerful magic. As an action, a creature can try to remove the piece. A creature that touches the piece must make a DC 15 Wisdom saving throw. On a failed save, the creature takes 11 (2d10) psychic damage, and the piece doesn't move. On a successful save, the creature takes half as much damage but is able to grab and move the piece freely. Successfully casting \textit{Dispel Magic} (DC 14) on the piece suppresses the magic for 1 minute, during which time the piece can be touched and removed without requiring creatures to make a saving throw.

When the piece is removed, the four defense sentinels attack the characters, and Landro's shutdown sequence begins (see the "Shutdown Sequence" section below). For more about the \textit{Rod of Seven Parts}, see this book's introduction.

\subsection{Graymatter Guardian}
The graymatter engine has—or, more aptly, is—a mind of its own. It sensed the characters when they first entered Landro and has watched them curiously. When the characters arrive in this room, the graymatter engine manifests a swirling mass of animated graymatter fluid.

The animated fluid takes the shape of a tall, lithe humanoid and introduces itself as Landro. This neutral-aligned being uses the \textbf{water elemental} stat block, except it has telepathy out to 30 feet.

Landro isn't interested in fighting. Rather, it wants to question the newcomers who've made it all the way to the colossus's uppermost level.

\subsubsection{Landro the Host}
In its manifestation as Landro, the graymatter engine assumes the role of a gracious host. It is far more interested in learning about its visitors than in talking about itself. If pushed, Landro briefly explains how it was created by Cyran mages to control the colossus of the same name. On the day of its deployment, the colossus teleported halfway into this mountain, where it has remained since. Unable to wrench its body from Mount Ironrot, the graymatter engine fell into a despondent torpor that lasted until the arrival of its recent visitors.

\subsubsection{Sharing Secrets}
Landro understands the world by peering into the thoughts and memories of living beings. During conversation, the entity asks the characters to share a secret with it.

\begin{DndReadAloud}{}
"I already know a great deal," Landro says. "The knowledge I now seek is of the clandestine variety. I would be interested to learn your secrets."

\end{DndReadAloud}

Landro wishes to learn secrets from the characters through a type of osmosis. Landro tells the characters that this process is harmless and takes only a moment, though it does require physical contact with the head of a consenting character, as well as spending a secret the character has learned. If the characters reject the idea, Landro doesn't push the matter.

If the characters consent to the entity's request, Landro extends a ribbonlike tentacle of graymatter fluid. Landro uses this tentacle to briefly touch the head of each consenting character in turn. An affected character hears a deep, sloshing sound and experiences a sensation like being underwater, but when the tentacle is removed, the character is as dry as before.

After this process, Landro nods in quiet appreciation and thanks the party for sharing their secrets.

\subparagraph{Secrets Revealed}
If any character consents to the graymatter engine's secret-learning process, the most recent secret the party learned is spent, as described in "The Power of Secrets" section in the book's introduction. Instead of the usual benefit of spending the secret, the characters all gain \textit{inspiration}. Mark off a secret spent in this way on the Secrets Tracker in appendix C. No matter how many characters consent, only one secret is spent.

\subsubsection{The Engine's Aid}
If the characters befriend Landro, or if two or more of the characters agree to share a secret with it, the entity warns the party of the engine room's defense sentinels and explains how the colossus's shutdown sequence works. Even though it is mentally merged with the colossus of the same name, the entity can't control the colossus's systems.

\subsubsection{Crystal Companion}
At the end of their conversation, Landro asks to travel with the party so it can continue to learn from them. If the characters find a special kind of gem called an id crystal, Landro explains, they can transfer the graymatter engine's consciousness to it. Landro tells the characters that a suitable id crystal can be found in the colossus's left foot (area L6).

When an id crystal touches the graymatter engine, the graymatter engine's consciousness is transferred to the crystal, which becomes an Elemental Gem, Emerald.

As soon as the graymatter engine's consciousness is transferred to the id crystal, the rod piece powering the graymatter engine falls to the floor, activating the hidden sentinels in area L28 and starting Landro's shutdown sequence (see below).

\subsubsection{Shutdown Sequence}
When Landro's shutdown sequence is activated, a magical, soothing voice announces: "One minute until self-destruct." The following things then happen immediately:


\begin{description}
\item[Antigravity Deactivated.]
Landro's antigravity magic—including the antigravity wells and floors—is dispelled.
\item[Barrier Dropped.]
The invisible barrier around Landro disappears.
\item[Self-Destruct Initiated.]
Characters inside Landro must succeed on a DC 14 Dexterity saving throw or have the prone condition as the colossus heaves and shudders. After 1 minute, the colossus explodes in a fiery conflagration. When this happens, creatures inside Landro must make a DC 18 Dexterity saving throw, taking 99 (18d10) fire damage on a failed save or half as much damage on a successful one. Unattended, nonmagical objects inside Landro are destroyed by the explosion.
\end{description}

\section{Next Steps}
Once they've acquired the third piece of the \textit{Rod of Seven Parts}, the characters can return to the oval portal on Mount Ironrot's outskirts. Stepping through the portal returns the characters to the sanctum in Sigil.

\chapter{Death House}
The fourth piece of the \textit{Rod of Seven Parts} is in Barovia, one of the Shadowfell's Domains of Dread. With the aid of an Ulmist inquisitor, the party infiltrates the basement of an eerie locale called Death House. Here, they must stop cults plotting to use the fourth rod piece in vile rituals while also preventing the artifact from falling into of the hands of the infamous vampire Strahd von Zarovich.

\section{Running This Chapter}
This chapter begins after the characters retrieve the third piece of the \textit{Rod of Seven Parts}. When a character holds this piece, they instinctively know that the next closest piece is located in the village of Barovia, in the domain of Strahd von Zarovich. Although it's a small village, Barovia teems with supernatural threats—including the forebodingly named Death House, the location of the fourth piece.

By researching Death House in the Sigil sanctum or asking the Wizards Three about it, the characters learn that Death House is owned by two Barovians named Gustav and Elisabeth Durst, who run a small cult devoted to Barovia's Darklord, Strahd von Zarovich. Given the domain's many greater threats, this minor cult and the house draw scant attention. If the characters ask, Alustriel shares the information in the "Knowledge of Barovia" section. Neither the wizards nor the available research materials in Sigil can reveal further details.

This chapter begins with the characters' approach to Death House. Their journey is interrupted by a mob of scared peasants as well as a potentially helpful inquisitor. Most of the chapter describes the horrors the characters encounter as they explore Death House. Once the characters acquire the rod piece, Strahd arrives at Death House to toy with the characters and tries to block their escape.

\subsection{Character Advancement}
The characters should be 14th level when this chapter begins. They gain a level after they retrieve the fourth piece of the \textit{Rod of Seven Parts} from Death House.

\subsection{Power of Secrets}
The characters can learn one secret in this chapter that is applicable to the rules in "The Power of Secrets" section in this book's introduction:


\begin{description}
\item[Sarusanda's Secret.]
Sarusanda is an Ulmist inquisitor, but her father, Galias, joined the evil priests of Osybus. Sarusanda expected to find and slay him in Death House. The characters can learn this secret during one of their encounters with Sarusanda, as described in the "Meetings with Sarusanda" section.
\end{description}

\subsection{Fourth Rod Piece}
The Rod of Seven Parts of the \textit{Rod of Seven Parts} is in area D38 of Death House. For more information about the rod and the spell this piece allows its wielder to cast, see this book's introduction.

\section{Barovia}
The third rod piece allows the characters to step through the portal in the Sigil sanctum and emerge in the western outskirts of the village of Barovia.

If the characters haven't already, allow them to research Barovia in Sigil or ask the wizards about the place, then convey the points listed in the "Knowledge of Barovia" section below. When you're ready to start the chapter, proceed to the "Arriving in Barovia" section.

\subsection{Knowledge of Barovia}
Characters who research Barovia can learn the following:


\begin{description}
\item[Domain of Dread.]
Barovia is the name of a village as well as the name of the Domain of Dread that encompasses that village. A Domain of Dread is a demiplane hidden in the Plane of Shadow. Every Domain of Dread is separated from the rest of the multiverse by the mysterious Mists.
\item[The Mists.]
The Mists are unfathomable and unpredictable. Once an individual has been taken by the Mists, there is little chance of escape. The only entities with influence over the Mists are the unknowable Dark Powers, which control the Domains of Dread, and the Darklords, who each rule a domain created to torment them.
\item[Strahd the Vampire.]
As Barovia's Darklord, the vampire Strahd von Zarovich wields immense power and usually has the final say over who comes and goes in his domain.
\end{description}

\section{Arriving in Barovia}
When the characters step through the portal in Sigil, they wade through a thick mist before emerging on the western outskirts of Barovia. Read or paraphrase the following:

\begin{DndReadAloud}{}
Thick, gray fog shrouds this small, gloomy village. Dismal houses and outbuildings line the cobblestone streets. The sound of a child weeping echoes preternaturally from a tall house looming in the distance.

\end{DndReadAloud}

If the characters attempt to leave or circumvent the village, the Mists return them to where they started.

The third piece of the \textit{Rod of Seven Parts} points toward the tall, gloomy house on the village's far side, where the sound of weeping is coming from.

\subsection{Unwelcome Party}
While making their way across the village, the characters are accosted by a throng of panicked villagers fearful of the newcomers. Read or paraphrase the following:

\begin{DndReadAloud}{}
As you move through town, a handful of murmuring villagers follow you at a distance. More villagers emerge from houses on all sides, and soon you're surrounded by a small mob. A farmer carrying a pitchfork points at you and calls you interlopers. Other folks brandish brooms, axes, and large stones. They scream at you to leave the village.

\end{DndReadAloud}

The mob consists of twenty hostile commoners who surround the characters. As Sarusanda later explains to the characters, the priests of Osybus foresaw the characters' arrival and sowed foul rumors about them to turn the villagers against them. The characters must diffuse the situation with the villagers before they can proceed.

The characters can disperse the mob in a variety of ways. They can lie about their identities or intentions; they can convince the villagers that they mean no harm or that they hope to help the children crying in the distant house; or they can threaten the villagers' lives. To successfully sway the mob in one of these ways, at least one character must succeed on a DC 18 Charisma (Deception, Intimidation, or Persuasion) check, using whichever skill is appropriate. A character who casts \textit{Calm Emotions} on the mob's area has advantage on this check.

Until the characters disperse or escape the mob, one randomly determined character takes 1 bludgeoning damage from hurled debris at the start of each turn. Each round on initiative count 0, the DC of the check to disperse the mob increases by 1. The villagers immediately disperse if the characters attack or deal damage to any of them.

\subsection{The Ulmist Inquisition}
After the characters are free of the mob, a cloaked adventurer named Sarusanda Allester approaches them. Sarusanda is a lawful neutral, human \textbf{inquisitor of the tome} (see appendix A) who speaks Celestial, Common, Draconic, and Elvish, and can cast \textit{Speak with Dead} at will. She introduces herself and asks to speak with the characters in private. The characters can slip into an alley for this discussion, or they can duck into the Blood of the Vine, the village tavern, for a quiet conversation.

\subsubsection{Sarusanda's Mission}
Sarusanda's initial impression of the characters depends on how the characters handled the angry mob. If the party dispersed the mob without harming any villagers, Sarusanda is friendly toward the characters. Otherwise, Sarusanda is indifferent but wary.

Sarusanda explains that she is a member of the Ulmist Inquisition, an organization dedicated to rooting out evil throughout the multiverse. She asks the characters what business they have at the house on the far side of the village.

However the characters respond, Sarusanda says she is also going to the building. "The locals call it Death House," she says. "Perhaps we can help each other achieve our goals there."

In the course of their conversation, Sarusanda conveys the following information to the party.

\subsubsection{Cultists in Death House}
Death House is owned by Gustav and Elisabeth Durst. It's an open secret that the Dursts use the building to host cultists who venerate the Darklord of Barovia, Strahd von Zarovich. This minor cult occasionally causes trouble, but until now, it hasn't merited intervention by the Ulmist Inquisition.

\subparagraph{Unexpected Discovery}
Recently, the cultists in Death House obtained something important: a fragment of the \textit{Rod of Seven Parts}. Sarusanda believes the cultists don't know much about the rod piece, but she suspects the cultists plan to use it to somehow attract Strahd's attention.

\subsubsection{Priests of Osybus}
The Death House cultists' discovery roused the attention of another wicked group, the priests of Osybus. Unaffiliated with the cultists in Death House, the priests of Osybus are necromancers who steal souls to fuel their evil, life-prolonging magic. They wish to claim the rod piece for use in their necromantic rituals, and they've divined the characters' arrival and know the characters will try to stop them. To slow the characters down, the priests spread dreadful rumors about them throughout town. This stirred the villagers into a frenzy. Sarusanda heard about the priests' activities, discovered the presence of the rod piece, and traveled to the village. She has vowed to stop the priests of Osybus's activities in Death House and any other evil happening there.

Eight priests of Osybus have infiltrated Death House. They are all neutral evil humans who speak Abyssal, Common, and Infernal.

\subsubsection{Strahd Watches}
If Sarusanda's information is correct, Strahd is also aware of the activities at Death House.

\subsubsection{Sarusanda's Plan}
Sarusanda needs to act quickly to stop the priests of Osybus and Strahd from acquiring the rod fragment. She doesn't care if the characters take the artifact; as long as it stays out of her enemies' clutches, she'll consider her work done. To this end, she invites the characters to join her cause. Regardless of whether the characters cooperate with Sarusanda or ignore her, their interaction with her now affect subsequent meetings with her inside Death House.

\subsection{Approaching Death House}
As the characters come within sight of Death House, read or paraphrase the following:

\begin{DndReadAloud}{}
A boy and a girl stand in the middle of the dirt road outside a grim house. The boy is weeping and clutches a stuffed doll. The girl is trying to quiet the boy. She turns to you. "There's a monster in our house!" she says. "Mom and Dad told us to play outside. Please, won't you make it safe?"

\end{DndReadAloud}

The children are Rosavalda "Rose" and Thornboldt "Thorn" Durst. Rose explains that the monster arrived just after the children's parents, Gustav and Elisabeth, ordered Rose and Thorn outside to play. Rose is especially worried for Brigetta, the family's nursemaid, who's always been kind to both children. After Rose and Thorn came outside, they heard screams inside the house, followed by terrible howls coming from the basement. Recently, a group of robed adults (priests of Osybus) arrived and entered the house without a word.

If calmed with a successful DC 10 Charisma (Persuasion) check, Thorn sniffles that his parents often invite weird friends over for parties, but this time seems different.

The children have no other useful information, and neither has any idea how much time has passed since the events they've described. They have no idea that their parents are cult leaders in possession of a powerful artifact. If the characters express hesitation about leaving the children alone, a friendly neighbor arrives and offers to keep the children safe.

\subsubsection{Sarusanda Splits Off}
As soon as the characters enter Death House, Sarusanda suggests she and the party explore the house separately. This way, you can focus on the characters without having to run Sarusanda alongside them. The characters will have several opportunities to run into Sarusanda in the house. See the "Meetings with Sarusanda" section for more details.

If you wish, Sarusanda can stay with the characters, and the events described in that section happen while she accompanies the party.

\section{Death House}
Before the characters enter Death House, make sure the players understand their goals. The characters can't leave until they encounter Strahd von Zarovich, as described in the "Leaving the House" section later in this chapter.

\subsection{Death House Objective}
The characters must take the fourth piece of the \textit{Rod of Seven Parts} from the Strahd-venerating cultists who operate from Death House's basement.

The priests of Osybus have their own plans for the rod piece. As long as the eight priests inside the house are alive, they'll try to keep the characters from taking the artifact.

\subsection{Death House Encounters}
The foul energies coalescing in Death House have spawned numerous horrors. When the characters enter locations with an "Encounter" subsection for the first time, the description will instruct you to roll on the Death House Encounters table below. With the exception of Meetings with Sarusanda, if you roll an encounter the characters have already overcome or an encounter that would be difficult to resolve in the location, choose an unencountered result instead.

\begin{DndTable}[header=Death House Encounters]{cX}

d12 & Encounter\\
1–4 & Meeting with Sarusanda (see below)*\\
5–8 & Three priests of Osybus (see appendix A) lurk nearby, poised to strike†\\
9 & One \textbf{whirling chandelier} (see appendix A) falls on a random character before two invisible stalkers join the surprise attack\\
10–11 & Four \textbf{vampire spawn} attack the party\\
12 & Four helmed horrors, disguised as inanimate suits of armor, attack the characters\\
\end{DndTable}

\subsection{Meetings with Sarusanda}
The characters occasionally reunite with Sarusanda as they explore Death House. Each time you roll this result on the Death House Encounters table, run one of the following encounters, starting with the first and proceeding in order.

\subsubsection{Wounded Pride}
\begin{DndReadAloud}{}
The remnants of a deadly battle lie about this chamber: broken bits of furniture, scattered weapons, and two bodies shrouded in dark robes. Sarusanda kneels in one corner, wrapping a linen bandage around her leg. She startles at your intrusion, then relaxes.

\end{DndReadAloud}

The robed bodies belong to two priests of Osybus (see appendix A), each of whom is stable at 0 hit points. When the characters enter the room, each priest's Tattoo of Osybus trait causes the priest to rise on its next turn, at which point it attacks the characters.

\subparagraph{Development}
Sarusanda has 30 hit points remaining and is obviously in pain. If the characters offer to heal her, she downplays her injuries. A character who succeeds on a DC 15 Wisdom (Insight) check realizes that it's not just Sarusanda's body that's hurt—her pride is badly wounded too.

If unaided, Sarusanda gains 1 level of exhaustion, which she retains until she finishes a long rest outside Death House.

\subsubsection{Séance}
\begin{DndReadAloud}{}
Pungent sticks of incense burn around this room's perimeter. Sarusanda stares intently at a human skull clutched in her outstretched hand.

"Now tell me, servant of evil," Sarusanda demands of the skull, "where do your fellow occultists hide?"

The skull replies, its teeth chattering, "Occultists? How droll! A better question: where are your manners?"

\end{DndReadAloud}

Sarusanda uses \textit{Speak with Dead} to commune with a preserved skull she found in this room. The skull belonged to a cult member named Elya who was sacrificed by his fellow cultists.

\subparagraph{Interrogating the Skull}
Elya's skull offers its words in hair-raising, singsong whispers. The skull doesn't know anything about the \textit{Rod of Seven Parts} or the priests of Osybus. It can, however, give directions to the cultists' ritual chamber (area D38) and describe their general motives: "We wish to impress our glorious lord, Strahd von Zarovich!"

Once Sarusanda is done speaking with the skull, it cackles in grim amusement, then floats into the air. Three will-o'-wisps emerge from corners of the room and join the skull (use the \textbf{flameskull} stat block) in attacking the party and Sarusanda.

\subsubsection{Crisis of Faith}
Run this encounter with Sarusanda after the characters have claimed the rod piece from the cultists (see area D38) but before they encounter Strahd.

\begin{DndReadAloud}{}
Sarusanda stands in the middle of the room. Her shoulders are slumped beneath her grimy armor, and dark circles underline her weary eyes. She doesn't turn to face you, but she holds a finger in the air to indicate silence. A moment later, the room erupts into chaos as three tattooed figures dart from the shadows.

\end{DndReadAloud}

Three priests of Osybus (see appendix A) attack Sarusanda and the party, entering the room from multiple directions if possible.

With or without the characters' help, Sarusanda fights passionately against the priests. She discards her normally staid demeanor for this battle and openly rages at the priests, making careless tactical decisions amid the turmoil.

\subparagraph{Sarusanda's Secret}
After the battle, a character who asks Sarusanda what's wrong is greeted with a long, bitter sigh. The inquisitor then reveals the source of her distress.

Shortly after she joined the Ulmist Inquisition, Sarusanda's father—a man named Galias—joined the priests of Osybus. Sarusanda witnessed him commit a terrible act in the cult's name, though she doesn't give details, and the inquisitor admits she didn't arrest her father. Sarusanda considers this her greatest failing.

Galias was fully inducted into the priests of Osybus, and Sarusanda heard that he would be among the cultists at Death House today. She is enraged that she hasn't found Galias so she can confront him.

Regardless of the characters' reaction to this revelation, learning it counts as a secret for the purposes of the Power of Secrets rules in this book's introduction.

\subsection{Death House Features}
The locations in Death House have the following features.

\subsubsection{Ceilings}
Ceilings vary in height by floor. The first floor has 10-foot-high ceilings, the second floor has 12-foot-high ceilings, the third floor has 8-foot-high ceilings, and the attic has 13-foot-high ceilings.

\subsubsection{Haunted Doors}
When Strahd enters the house later in this chapter, certain doors marked on the map become haunted doors; see the "Haunted Zones" section for details.

\subsubsection{Lighting}
Unless otherwise noted, each room in the house is lit with bright light by oil lamps, a fireplace, or some other light source when the characters arrive. The cultists take oil lamps into the ritual chamber (area D38) when they gather there.

\subsubsection{The Mists}
When the characters enter Death House, the Mists surround the building and prevent them from leaving. A creature that enters the Mists around Death House reemerges seconds later in a random room in Death House. Only Strahd can disperse the Mists (see the "Leaving the House" section).

\subsubsection{Self-Repair}
After 24 hours, any damage to the house repairs itself. The house also repairs itself as soon as Strahd enters the building.

\subsubsection{Walls}
The walls of the house are made of 1-foot-thick brick. The dungeon is carved from earth, clay, and rock. Dungeon tunnels are 5 feet wide by 10 feet high with timber braces at 5-foot intervals.

\subsection{Death House Locations}
These locations are keyed to map 5.1.

\subsubsection{D1a–D1b: Portico and Antechamber}
\begin{DndReadAloud}{}
A hinged, wrought-iron gate fills the archway leading to a stone portico. Oil lamps hang from the portico ceiling by chains, flanking an oaken double door.

The double door opens into an empty antechamber. On the south wall is a shield emblazoned with a stylized golden windmill on a red field.

A closed double door stands a few steps beyond.

\end{DndReadAloud}

The doors are unlocked.

\subsubsection{D2a–D2b: Main Hall and Cloakroom}
\begin{DndReadAloud}{}
This hall runs the width of the house, with a black marble fireplace at one end and a sweeping red marble staircase at the other.

The door to a cloakroom on the east wall is slightly ajar; the black cloaks inside are damp.

\end{DndReadAloud}

The hall is bare of furniture. Faint footprints lead to the cloakroom (area D2b).

\subparagraph{Encounter}
Roll on the Death House Encounters table the first time a character enters this room.

\subsubsection{D3: Den of Wolves}
\begin{DndReadAloud}{}
This oak-paneled room looks like a hunter's den. A stag's head is mounted above the fireplace, and three stuffed wolves are positioned around the outskirts of the room.

\end{DndReadAloud}

\subparagraph{Trapdoor}
A character searching the southwest corner of the room finds a trapdoor with a successful DC 20 Intelligence (Investigation) check. The trapdoor is barred from the other side (area D32). It is a Medium object with AC 13, 15 hit points, and immunity to poison and psychic damage.

\subsubsection{D4a–D4b: Kitchen and Pantry}
\begin{DndReadAloud}{}
The kitchen is tidy, with dishware, cookware, and utensils neatly placed on shelves. A dome-shaped stone oven stands near the east wall. Near the oven is a well-stocked pantry.

\end{DndReadAloud}

The kitchen is area D4a, while the pantry is area D4b.

\subparagraph{Dumbwaiter}
Behind a small door in the southwest corner of the kitchen is a dumbwaiter—a 2-foot-wide stone shaft containing a wooden elevator box attached to a hand-operated rope-and-pulley mechanism. The shaft connects to the servants' quarters (area D7a) and the Durst parents' bedroom (area D12a).

A Small character can squeeze into the elevator box with a successful DC 10 Dexterity (Acrobatics) check. The dumbwaiter's rope-and-pulley mechanism can support 200 pounds of weight before breaking.

\subsubsection{D5: Dining Room}
\begin{DndReadAloud}{}
The centerpiece of this wood-paneled dining room is a carved mahogany table, covered with resplendent silverware and crystalware and surrounded by eight high-backed chairs.

\end{DndReadAloud}

This area is empty except for the immaculate table dressing.

\subparagraph{Treasure}
The silverware and crystalware here are worth 2,500 gp total.

\subsubsection{D6: Upper Hall}
\begin{DndReadAloud}{}
Unlit oil lamps are mounted on the walls of this elegant hall. Hanging above the mantelpiece is a wood-framed portrait of a family: a man and a woman with two smiling children—the same children you saw outside. The marble staircase continues upward.

\end{DndReadAloud}

The red marble staircase that started on the first floor continues its upward spiral to area D11. A cold draft flows down the steps.

\subsubsection{D7a–D7b: Servants' Room and Closet}
\begin{DndReadAloud}{}
An undecorated bedroom contains a pair of beds with straw-stuffed mattresses.

\end{DndReadAloud}

An empty footlocker rests at the foot of each bed. Tidy servants' uniforms hang from hooks in the adjoining closet (area D7b).

\subparagraph{Dumbwaiter}
A small door in the corner of the room opens onto the shaft of the dumbwaiter.

\subsubsection{D8: Library}
\begin{DndReadAloud}{}
A mahogany desk faces the fireplace of this private library, and two overstuffed chairs fill the room's corners. Floor-to-ceiling bookshelves line the south wall. A rolling wooden ladder allows access to the higher shelves.

\end{DndReadAloud}

The bookshelves hold hundreds of tomes on topics such as history, warfare, and alchemy.

\subparagraph{Encounter}
Roll on the Death House Encounters table the first time a character enters this room.

\subparagraph{Secret Door}
A secret door behind a bookshelf in the southeast corner of the room can be opened by pulling on a switch disguised as a red-covered book with a blank spine. A character inspecting the bookshelf spots the fake book with a successful DC 13 Intelligence (Investigation) check. Unless the secret door is propped open, springs in the hinges cause it to close on its own. Beyond the secret door lies area D9.

\subsubsection{D9: Secret Room}
This secret room contains bookshelves packed with tomes describing Fiend-summoning rituals.

\subparagraph{Trapped Chest}
An unlocked treasure chest stands against the room's south wall. When a creature opens this chest's lid, poison-tipped darts shoot from a spring-loaded mechanism attached to the lid's underside. Each creature within 10 feet of the chest that isn't behind \textit{total cover} takes 2 (1d4) piercing damage and must succeed on a DC 15 Constitution saving throw or take 22 (4d10) poison damage and have the poisoned condition for 1 hour.

A character who succeeds on a DC 20 Intelligence (Investigation) check while examining the chest spots the trap mechanism. As an action, a character can use thieves' tools to make a DC 15 Dexterity (Sleight of Hand) check to disarm the trap. Failing this check triggers the trap.

\subparagraph{Treasure}
The chest contains three blank books with black leather covers (worth 25 gp each) and three Spell Scroll (\textit{Bless}, \textit{Protection from Poison}, and \textit{Spiritual Weapon}).

\subsubsection{D10: Conservatory}
\begin{DndReadAloud}{}
A harpsichord and bench fill this elegant hall's northwest corner. On the opposite side of the room, a large standing harp rests near the room's fireplace. Upholstered chairs line the walls.

\end{DndReadAloud}

\subparagraph{Encounter}
Roll on the Death House Encounters table the first time a character enters this room.

\subsubsection{D11: Balcony}
\begin{DndReadAloud}{}
At the top of the red marble staircase is a balcony with a suit of black plate armor standing against one wall.

\end{DndReadAloud}

\subparagraph{Secret Door}
A character who examines the west wall and succeeds on a DC 15 Intelligence (Investigation) check finds a secret door that pushes open to reveal a cobweb-filled wooden staircase leading to the attic.

\subsubsection{D12a–D12c: Master Suite}
\begin{DndReadAloud}{}
Burgundy drapes cover the windows of this large bedroom. The furnishings include a four-poster bed, a matching pair of wardrobes, and a vanity with a jewelry box. A portrait of a man and woman hangs above a fireplace. A parlor in the southwest corner contains a table and two chairs.

\end{DndReadAloud}

The couple in the portrait are the same people from the family portrait in the upper hall (area D6). A door facing the foot of the bed opens to an empty closet (area D12b). A door in the parlor leads to an outside balcony (area D12c).

\subparagraph{Dumbwaiter}
A dumbwaiter in the southwest corner connects to areas D4a and D7a below.

\subparagraph{Treasure}
The jewelry box on the vanity is made of silver with gold filigree (worth 75 gp). It contains three gold rings (worth 25 gp each) and a thin platinum necklace with a topaz pendant (worth 750 gp).

\subsubsection{D13: Bathroom}
\begin{DndReadAloud}{}
This dark room contains a barrel under a spigot in the east wall, a wooden claw-foot tub, and a small iron stove with a kettle resting atop it.

\end{DndReadAloud}

This room is otherwise empty.

\subsubsection{D14: Storage Room}
Dusty shelves line the walls of this room. The shelves hold linens and other household goods.

\subsubsection{D15a–D15c: Nursemaid's Suite}
\begin{DndReadAloud}{}
This bedroom contains a large bed, two nightstands, and a large mirror. A double door fitted with stained-glass windows opens onto a balcony. One of the nightstands has been knocked over, and one of the windows has been smashed.

\end{DndReadAloud}

This elegantly appointed bedroom (area D15a) and an adjoining nursery (area D15b) were occupied by the Dursts' nursemaid, Brigetta, the cult's latest victim.

The double door opens onto a balcony (area D15c) overlooking the front of the house.

\subparagraph{Encounter}
Roll on the Death House Encounters table the first time a character enters this room.

\subparagraph{Signs of Abduction}
A character who succeeds on a DC 15 Intelligence (Investigation) check ascertains that the disarrayed nightstand and broken window are the result of a recent struggle. Rainwater on the floor around the broken window indicates the struggle occurred within the last few hours.

A character who examines the area and succeeds on a DC 15 Wisdom (Survival) check finds scuffs on the floor in area D15a from the victim's shoes. These scuffs create a trail to a secret door that leads up to the attic.

\subparagraph{Secret Door}
A character who examines the mirror and succeeds on a DC 15 Intelligence (Investigation) check finds a secret door behind the mirror. (Characters who saw scuff marks leading to the secret door find it without needing to make a check.) The secret door pushes open to reveal a cobweb-filled wooden staircase leading to the attic.

\subsubsection{D16: Attic Hall}
This bare hallway can be accessed by the staircase behind the secret doors in areas D11 and D15a.

\subsubsection{D17: Spare Bedroom}
\begin{DndReadAloud}{}
This room contains a narrow bed, a nightstand, a writing desk, and a rocking chair.

\end{DndReadAloud}

This room is otherwise empty.

\subsubsection{D18: Storage Room}
\begin{DndReadAloud}{}
This chamber is packed with old furniture draped in white sheets.

\end{DndReadAloud}

\subparagraph{Secret Door}
A character who examines the east wall and succeeds on a DC 15 Intelligence (Investigation) check finds a secret door that opens to reveal a secret stair (area D21) to the dungeon (area D22).

\subsubsection{D19: Spare Bedroom}
\begin{DndReadAloud}{}
This room contains a narrow bed, a nightstand, a rocking chair, an empty wardrobe, and a small iron stove.

\end{DndReadAloud}

This room is otherwise empty.

\subsubsection{D20: Children's Room}
\begin{DndReadAloud}{}
This bedroom contains two small, wood-framed beds, a toy chest, and a dollhouse—a perfect miniature replica of this dreary mansion.

\end{DndReadAloud}

Characters who investigate the dollhouse find all the house's secret doors, including one in the attic leading to a spiral staircase (area D21) that descends below the house.

\subparagraph{Encounter}
Roll on the Death House Encounters table the first time a character enters this room.

\subparagraph{Treasure}
Inside the toy chest is a doll that resembles a tall, pale human with black hair and a pronounced widow's peak. Something papery rattles inside the doll's hollow body. A character who removes the doll's head and looks down its neck finds a \textit{Spell Scroll} of \textit{Detect Evil and Good}.

\subsubsection{D21: Secret Stairs}
This narrow spiral staircase is made of creaky wood contained within a 5-foot-wide shaft of mortared stone. The staircase descends 50 feet from the house's attic to area D22 in the dungeon level.

\subsubsection{D22: Dungeon Level Access}
The wooden spiral staircase from the attic (area D21) ends here. A narrow tunnel stretches southward before branching east and west.

\subparagraph{Chanting}
From the moment they arrive in the dungeon, the characters hear an eerie, incessant chant echoing throughout the chambers. It's impossible to determine which direction the sound is coming from until the characters reach areas D26 or D29. They can't discern its words until they reach area D35.

\subsubsection{D23c–D23f: Family Crypts}
\begin{DndReadAloud}{}
Each crypt is open to the hallway. The stone slabs meant to seal the rooms lean against the walls outside the crypts.

\end{DndReadAloud}

These crypts are intended to hold the remains of the Durst family members on their deaths.

The crypts in areas D23a and D23b are empty, and the slabs for these crypts are unmarked.

The remaining crypts (areas D23c–D23f) each contain an empty coffin on a stone bier.

\subsubsection{D24: Cult Initiates' Quarters}
A wooden table and four chairs stand at the east end of this room. Four alcoves containing straw pallets open off the west half of the room.

\subsubsection{D25: Cultist Quarters and Well}
\begin{DndReadAloud}{}
This area centers on a 4-foot-diameter, 30-foot-deep cistern with a 3-foot-high stone lip. A wooden bucket hangs from a rope-and-pulley mechanism bolted to the crossbeams above the well. Five small rooms with no doors branch off the well room.

\end{DndReadAloud}

The five small bedrooms serve as quarters for senior cultists. Each contains a wood-framed bed with a straw mattress and a wooden chest. Each chest is secured with a rusty iron padlock that can be picked using thieves' tools with a successful DC 15 Dexterity (Sleight of Hand) check.

\subparagraph{Encounter}
Roll on the Death House Encounters table the first time a character enters this room.

\subparagraph{Treasure}
In addition to worthless personal effects, each chest contains one of the following:


\begin{itemize}
\item 110 gp and 60 sp in a pouch made of human skin
\item Three pieces of cut jade (worth 100 gp each) in a folded piece of black cloth
\item A black leather eye patch with a peridot (worth 500 gp) sewn into it
\item A chess set with pieces made of obsidian and chalcedony (worth 250 gp)
\item A \textit{+2 Shortsword}
\end{itemize}

\subsubsection{D26: Hidden Spiked Pit}
\begin{DndReadAloud}{}
As you move farther down this tunnel, the chanting heard throughout the dungeon gets louder to the west.

\end{DndReadAloud}

A character who examines the floor and succeeds on a DC 15 Intelligence (Investigation) check notices a suspicious absence of footprints in this hallway. A character searching the floor for traps finds a 5-foot-long, 10-foot-deep pit hidden under rotted wooden planks. The pit has poisoned wooden spikes at the bottom. The first character to step on the cover falls through, landing prone and taking 3 (1d6) bludgeoning damage from the fall plus 11 (2d10) piercing damage and 11 (2d10) poison damage from the spikes.

\subsubsection{D27: Dining Hall}
\begin{DndReadAloud}{}
This room contains a plain wooden table flanked by long benches. Bones lie scattered on the floor.

\end{DndReadAloud}

These moldy Humanoid bones are the remains of the cult's vile banquets.

\subsubsection{D28: Larder}
Aside from scraps of food, this alcove is empty.

\subsubsection{D29: Intersection}
The chanting heard throughout the dungeon is noticeably louder to the north of this intersection.

\subsubsection{D30: Stairs Down}
Any character standing at the top of this 20-foot-long staircase realizes the chants originate from somewhere below. Characters who descend the stairs and follow the hall beyond arrive in area D35.

\subsubsection{D31: Darklord's Shrine}
\begin{DndReadAloud}{}
This room is full of moldy skeletons that hang from rusty shackles against the walls. A wide alcove in the south wall contains a painted wooden statue carved in the likeness of a gaunt, pale-faced man wearing a black cloak.

\end{DndReadAloud}

The statue depicts Strahd, to whom the cultists make sacrifices in the vain hope that he will reveal his secrets to them.

\subparagraph{Encounter}
Roll on the Death House Encounters table the first time a character enters this room.

\subparagraph{Secret Door}
Characters searching the room find a secret door in the middle of the east wall with a successful DC 10 Intelligence (Investigation) check. The door pulls open to reveal a stone staircase that climbs 10 feet to a landing (area D32).

\subsubsection{D32: Hidden Trapdoor}
The staircase ends at a landing with a 6-foot-high wooden ceiling with a trapdoor set into it. The trapdoor is bolted shut from this side and can be pushed open to reveal the den (area D3) above.

\subsubsection{D33: Cult Leaders' Den}
\begin{DndReadAloud}{}
A chandelier hangs above a table in the middle of this room. Two high-backed chairs flank the table, which has an empty clay jug and two clay flagons atop it.

\end{DndReadAloud}

This is Gustav and Elisabeth's den. Iron candlesticks stand in two corners, their candles lit.

\subsubsection{D34: Cult Leaders' Quarters}
\begin{DndReadAloud}{}
This room contains a wood-framed bed with a feather mattress, a wardrobe containing several robes, a pair of iron candlesticks, and an open crate containing thirty torches and a leather sack with fifteen candles inside it.

\end{DndReadAloud}

At the foot of the bed is an unlocked but empty wooden footlocker. This is Gustav and Elisabeth's bedroom.

\subsubsection{D35: Reliquary}
\begin{DndReadAloud}{}
Loud chants issue from a hallway near the southwest corner. You discern many voices repeating, "She is slain, she is risen."

\end{DndReadAloud}

Lurking in the shadows are two shambling mounds that move to attack as soon as anyone enters this room. The cult keeps "relics" in this chamber. These worthless items—human remains, mundane material components for spellcasting, and so forth—are stored in thirteen niches along the walls.

\subsubsection{D36: Prison}
\begin{DndReadAloud}{}
The alcoves here feature chains ending in shackles attached to the back walls.

\end{DndReadAloud}

Despite the ominous trappings, the alcoves are empty.

\subparagraph{Secret Door}
Characters searching the area find a secret door in the south wall with a successful DC 15 Intelligence (Investigation) check. It pulls open to reveal area D38 on the other side.

\subsubsection{D37: Tunnel to Portcullis}
\begin{DndReadAloud}{}
This tunnel is blocked by a rusty iron portcullis.

\end{DndReadAloud}

This tunnel slopes down at a 20-degree angle into murky water and ends at a rusty portcullis. The floor around the portcullis is submerged under 2 feet of soupy water.

The portcullis can be forcibly lifted with a successful DC 20 Strength (Athletics) check. Otherwise, it can be raised or lowered by turning a wooden wheel in the east wall of area D38. (The wheel is beyond the reach of someone east of the portcullis.)

\subsubsection{D38: Ritual Chamber}
\begin{DndReadAloud}{}
Featureless stone pillars support the high ceiling in this forty-foot-square chamber. Murky water covers the floor, and a breach in the west wall leads to a dark cave heaped with refuse. Stairs lead up to dry stone ledges that hug the walls. Fourteen hooded figures stand atop the ledges, chanting loudly.

In the middle of the room, an octagonal dais rises above the water. Atop the dais is a stone altar drenched in freshly spilled blood. The source of the blood is the body of lifeless woman dangling above the altar from chains mounted to the ceiling. A ceramic stake has been plunged through the woman's heart. Standing next to the altar is a large fiend wielding a gruesome crimson spear.

\end{DndReadAloud}

The cultists chant: "She is slain, she is risen. Come to us, Lord Strahd. Receive our gift!"

The water here is 2 feet deep. The ledges and central dais are 5 feet high (3 feet higher than the water's surface), and the chamber's ceiling is 16 feet high (11 feet above the dais and ledges). The chains dangling from the ceiling are 8 feet long.

Embedded in the east wall is a wooden wheel connected to hidden chains and mechanisms. A character can use an action to turn the wheel, raising or lowering the portcullis in the eastern wall (see area D37). The hole in the west wall contains a heap of stones, rags, and plant matter.

\subparagraph{Cult's Ritual}
The figures along the ledge include Gustav and Elisabeth (cult fanatics), five senior cult members (cult fanatics), and seven cultists. The cult members fight only to defend themselves.

Dangling from the chains above the altar is the corpse of the Dursts' nursemaid, Brigetta. The cult ritually sacrificed Brigetta using a ceremonial stake tipped with the Rod of Seven Parts of the \textit{Rod of Seven Parts}, which is still embedded in Brigetta's chest.

A \textbf{relentless impaler} (see appendix A) has risen from the pool of Brigetta's blood on the altar. When the impaler detects the characters, it attacks. The impaler hunts them throughout the house until it is destroyed or the characters defeat it. For its Bloodheart Stake trait, the impaler is bound to the stake in Brigetta's corpse.

\subparagraph{Retrieving the Rod Piece}
If the relentless impaler is destroyed, the cultists cower in fear and allow the characters to take the rod piece from Brigetta's corpse. For more about the \textit{Rod of Seven Parts}, see this book's introduction.

\subparagraph{Development}
Once the characters claim the rod piece, Gustav and Elisabeth regard them with bitter disdain. "None dare trespass in the territory of Strahd," Gustav says. "You will answer to our master," Elisabeth adds with sinister certainty. They offer no further comments or information. As the characters leave the area, the cultists chant, "He is the Ancient. He is the Land."

The cult leaders aren't bluffing. As soon as the characters take the fourth rod piece, Strahd enters Death House and makes his presence known. See the "Strahd's Presence" section below.

\subsection{Sarusanda's Company}
Once the characters have defeated the Death House cultists and claimed the rod piece, run the "Crisis of Faith" encounter detailed in the "Meetings with Sarusanda" section earlier in the chapter at a time of your choosing to ensure she rejoins the characters before they face Strahd.

Having slain the priests of Osybus, seen the cult's plans thwarted, and made sure the \textit{Rod of Seven Parts} piece is safe with the characters, Sarusanda's work in the house is done. She accompanies the characters for the remainder of their time in Death House (see the "Leaving the House" section).

\section{Strahd's Presence}
Strahd has finally deemed Death House worthy of attention. Once the party retrieves the rod piece in the ritual chamber, read the following to describe Strahd's arrival at Death House. The characters won't yet see Strahd, but they'll feel his presence.

\begin{DndReadAloud}{}
A gust of wind blasts through the halls of Death House, and the temperature drops precipitously. Candles and lamps snuff out and relight with dark-purple flames. The wind takes on a monstrous timbre—like screeching bats or howling wolves—and crescendos to a deafening roar. Suddenly, it stops.

For a long moment, you hear nothing. Then, as clearly as if it were happening right next to you, the house's front door creaks open and slowly closes, latching with a crisp click. Then, once again, nothing.

\end{DndReadAloud}

The characters must now contend with Strahd's nightmare magic in Death House until they find and face the Darklord directly.

\subsection{Haunted Zones}
Strahd can twist his domain to suit his whims. Starting when he makes his presence known, certain doorways in the house no longer connect to the rooms they did before. Instead, creatures that pass through these doorways enter an isolated demiplane called a haunted zone.

\subsubsection{Entering a Haunted Zone}
A character enters a haunted zone whenever they do one of the following:


\begin{itemize}
\item Finish a long rest in Death House
\item Move through a haunted doorway (marked on map 5.1)
\end{itemize}

Only one haunted zone manifests at a time. As long as at least one creature is in a haunted zone, any other creatures that enter a haunted zone appear in that same zone regardless of how they arrive. Once all the characters inside a haunted zone exit it, that zone can't be encountered again. The haunted zones are described below in the "Death House Hauntings" section.

\subparagraph{Haunted Doorways}
The room on the other side of a haunted doorway is filled with magical darkness and silence, giving no hint of whatever is within. A creature that partially passes through a haunted doorway—such as by poking a limb or held item into the darkness—is magically pulled into the haunted zone.

Haunted doorways are one-way passages. Within the demiplane, nothing exists except what is described in each entry. When a character passes through a haunted doorway, that haunted doorway turns back to a normal doorway after 1 minute.

\subsubsection{Exiting a Haunted Zone}
To leave a haunted zone, a character must perform the task specified in the haunted zone's "Correct Exit" section. The "Failed Exit" section describes how creatures that don't perform the required task are eventually forced from the haunted zone.

\subparagraph{Finding the Correct Exit}
A character who casts \textit{Detect Evil and Good} or \textit{Find the Path} or who has truesight sees a faint aura around an object related to the task necessary to exit the room. A character can also locate such an object by looking at the area and succeeding on a DC 15 Intelligence (Investigation) check.

\subparagraph{Leaving a Haunted Zone}
A character who exits a haunted zone appears somewhere in Death House. Characters who leave the same haunted zone emerge in the same part of Death House. The transition is sudden and disorienting, like waking from a dream. Roll on the Emergence from a Haunted Zone table below to determine where the characters appear when they leave a haunted zone, or choose one of the locations listed.

\begin{DndTable}[header=Emergence from a Haunted Zone]{cX}

d6 & Location Emerged\\
1 & The characters sit in the dining room (area D5), sated and covered with sticky crumbs, though there's no food in sight.\\
2 & The characters stand in the storage room (area D18), draped in sheets.\\
3 & The characters kneel at the side of the cultists' well (area D25), desperately parched.\\
4 & The characters are slumped over the furniture in the storage room (area D18).\\
5 & The characters lean against the railing of the nursemaid's balcony (area D15c).\\
6 & The characters are shackled to the wall in the prison (area D36), each character in a different alcove. The characters can easily break their bonds.\\
\end{DndTable}

Once the characters have encountered all four haunted zones, they don't encounter any more. Long rests function as normal, and the haunted doorways become normal doorways.

\subsection{Death House Hauntings}
The characters enter the following haunted zones in the order presented.

\subsubsection{Castle Dining Hall}
\begin{DndReadAloud}{}
You stand at one end of an impossibly long wooden dining table in the hall of a stone keep. Gilded plates, cutlery, and goblets of blood-red wine are set for each member of your party.

At the other end of the table stands a tall, pale man. He wears regal, red garb, and his raven hair comes to a prominent widow's peak. He smiles menacingly at you as he raises his goblet in a toast. "Welcome, my guests. May you find the eternal hospitality of Strahd von Zarovich to your liking. Should at any point you wish to leave, simply find me and ask."

\end{DndReadAloud}

The "Strahd" in this haunted zone is an illusion. The illusion stands still in anticipation of the characters returning its toast, at which point the illusion quaffs its drink; it doesn't respond to anything else. If the illusion is struck with an attack, it dissolves into a cloud of bats (see "Failed Exit" below).

The chamber distorts and elongates as a character attempts to move throughout it. No matter how fast or far a character moves, that character winds up a step or two away from the end of the table where they started.

\subparagraph{Correct Exit}
To escape the dining hall, each character trapped in it must drink from one of the chalices on the table. The wine is thick and metallic tasting. When the character drinks it, the character's vision slowly turns red before returning to normal when the character exits the haunted zone. The only thing the character can hear during this time is the sound of Strahd's deep, spooky laughter.

\subparagraph{Failed Exit}
If after 1 minute no one has drunk from a chalice, Strahd's illusion dissolves into a cloud of bats that quickly fills the chamber. At the start of each turn, the bats deal 7 (2d6) piercing damage to each creature in the haunted zone. After 5 turns, the bats completely obscure the characters' vision, and the characters exit the haunted zone.

\subsubsection{Endless Graveyard}
\begin{DndReadAloud}{}
Countless headstones dot this barren landscape as far as you can see. A pale full moon looms over the horizon.

\end{DndReadAloud}

The headstones continue in every direction no matter how far a character travels. Each headstone is inscribed with the name of a different individual who died within Strahd's domain (the names eventually repeat). Buried beneath each tombstone is a coffin that either is empty or contains an illusion of Strahd's corpse.

\subparagraph{Correct Exit}
A character who reads the headstones finds one bearing their own name. A character can exit this haunted zone by disinterring the empty coffin buried beneath a headstone bearing the character's name, opening the coffin, and lying in it.

\subparagraph{Failed Exit}
If a character fails to leave the graveyard, that character must eventually make camp. When a character finishes a long rest in this haunted zone, the character gains 1 level of exhaustion, doesn't gain the benefits of the long rest, and exits the haunted zone.

\subsubsection{Ghostly Recital}
\begin{DndReadAloud}{}
You're seated in an ornate chair on the edge of a dimly lit parlor. Eight ghostly figures sit in identical chairs to either side of you. At the parlor's far end, a pale man at a harpsichord slowly plinks out a haunting elegy. The ghostly figures listen to the song with rapt attention.

\end{DndReadAloud}

This zone resembles Death House's conservatory (area D10), except the room's doors and windows overlook a black, featureless void.

If a character makes any noise or otherwise disrupts the performance, the musician—yet another illusion of Strahd—ceases playing, scowls at the character, and disappears. Eight specters fly out of the harpsichord and attack the characters in anger.

\subparagraph{Correct Exit}
A character at the harpsichord can make a DC 14 Charisma (Performance) check to play the rest of the interrupted melody. On a successful check, the ghostly figures sigh with contentment, and the characters exit the haunted zone.

\subparagraph{Failed Exit}
If after 1 minute no one has successfully played the rest of the song, the illusion of Strahd reappears at the harpsichord and hammers out a series of chords. Each character takes 11 (2d10) thunder damage and exits the haunted zone.

\subsubsection{The Crying Room}
\begin{DndReadAloud}{}
The floor of this bedroom is littered with creepy dolls, monstrous stuffed animals, and misshapen toys. The two children you saw outside the house stand in the middle of the doorless, windowless room, weeping uncontrollably. Large tears run down their faces, pooling in impossible quantities on the floorboards.

\end{DndReadAloud}

The children are illusions of Rose and Thorn. They are inconsolable; any attempt to interact with the children reveals their illusory nature. Every minute, the water covering the floor of the 10-foot-high bedroom rises an additional 1 foot.

A character who inspects the items in the room discovers that the dolls look like cartoonish versions of the characters, Sarusanda, and Strahd.

\subparagraph{Correct Exit}
To exit the haunted zone, a character must pull the head off the Strahd doll. If the characters struggle to figure this out, allow them to examine the doll and attempt a DC 13 Intelligence (Investigation) check to realize the truth. The doll's neck reveals itself to be a drain hole, and the head is a plug. All the water in the room spirals into the drain, pulling the characters with it and causing them to exit the zone.

\subparagraph{Failed Exit}
If the characters are still in this haunted zone after 10 minutes, the water level reaches the room's ceiling, and each character in this haunted zone gains 1 level of exhaustion and exits the zone.

\section{Leaving the House}
The characters can attempt to leave Death House at any time. If the characters haven't yet faced Strahd, the Mists invariably return them to a random room. Teleportation spells likewise fail.

To leave Death House, the characters must first face Strahd.

\subsection{Reuniting with Sarusanda}
If Sarusanda hasn't yet rejoined the characters, run the "Crisis of Faith" encounter detailed in "Meetings with Sarusanda"; she then rejoins the party and is with them when they face Strahd.

\subsection{Facing Strahd}
After he arrives at Death House, \textbf{Strahd, Master of Death House} (see appendix B) waits for the characters in the house's main hall (area D2a). He is accompanied by two \textbf{vampire spawn}.

If the characters don't intend to return to the main hall, Strahd and his minions wait for the characters in another room of the house, such as the dining room (area D5) or the den of wolves (area D3).

When Strahd encounters the characters, read the following:

\begin{DndReadAloud}{}
"So, intruders into my realm, we finally meet. From the way the villagers spoke about you, I expected you to be more impressive. Ah, well. I am forever disappointed with the insects that think themselves worthy of meeting me. I suppose I could entertain myself with destroying you."

\end{DndReadAloud}

\subsubsection{Sarusanda's Influence}
If Sarusanda is with the party, Strahd regards the Ulmist inquisitor with grave respect and allows her to leave the house. "I will permit you to flee this place," he intones. "But only if you do so immediately, and you never return."

He glances at the characters, then asks Sarusanda, "Do you know them?"

\subparagraph{Sarusanda's Answer}
If Sarusanda is indifferent to the characters, she denies any affiliation with them and takes her leave. Strahd applauds the characters for surviving so long in his realm, then attacks them.

If Sarusanda is friendly with the characters, she says that they are also members of the Ulmist Inquisition. Strahd asks the characters if this is true. A character can then make a DC 20 Charisma (Deception) check to convince Strahd they are members or otherwise important to the group. On a successful check, Strahd allows the party to leave Death House with Sarusanda. On a failure, Strahd sees through the ruse and dismisses Sarusanda before challenging the characters to combat.

\section{Next Steps}
Once the characters have acquired the Rod of Seven Parts of the \textit{Rod of Seven Parts} and escaped Death House, they can travel to the western outskirts of Barovia without trouble. The characters easily find the portal to Sigil in the mist.

\chapter{Night of Blue Fire}
To find the next piece of the \textit{Rod of Seven Parts}, the characters must travel to the world of Krynn where the Blue Fire Wardens, a coalition of benevolent lycanthropes, have clashed with the death knight \textbf{Lord Soth} and his minions. By infiltrating a heavily fortified keep, rescuing the Blue Fire Wardens' leader, and unraveling the schemes of the evil mage Teremini Nightsedge, the characters can claim the fifth rod piece.

\section{Running This Chapter}
This chapter begins after the characters retrieve the fourth piece of the \textit{Rod of Seven Parts}. When a character holds this piece, they instinctively know that the fifth piece is located somewhere on the war-torn world of Krynn. That character senses that the fifth piece is located near a massive, magical tree, in a region called the Northern Dargaard Mountains, but this indication seems inexact. The fifth piece's location is unclear to the characters.

Neither the characters nor the Wizards Three know why the location of the fifth piece is ambiguous. In truth, the fifth piece was housed in the tree until recently, when it was taken to a nearby castle called Three Moons Vault. The crimson moonlight that now surrounds the castle warps the rod's magical properties. To retrieve the fifth rod piece, the characters must dispel the magical moonlight by disrupting a wicked ritual in the vault.

This chapter starts when the characters step through the portal in Sigil, emerge on Krynn, and begin their search for the magical tree. Though initially fruitless, the characters' search puts them in contact with benevolent lycanthropes who know where to find the fifth piece. The characters must then infiltrate the Three Moons Vault. There, the characters must secure the fifth rod piece from Teremini Nightsedge, an ally of Lord Soth.

\subsection{Character Advancement}
The characters should be 15th level when this chapter begins. The characters gain a level after they retrieve the fifth piece of \textit{Rod of Seven Parts} from the upper level of Three Moons Vault.

\subsection{Power of Secrets}
The characters can learn two secrets in this chapter that are applicable to the rules in "The Power of Secrets" section in this book's introduction:


\begin{description}
\item[Gazaia's Secret.]
The dryad Gazaia hid and watched while soldiers attacked the peylon tree where she lived and looted the fifth piece of the \textit{Rod of Seven Parts}, which was sustaining the tree. Gazaia now hides in the grotto (area P4) described in the "Peylon Tree Locations" section later in this chapter.
\item[Valendar's Secret.]
The \textbf{werewolf} Valendar led an assault against his enemies without properly planning the mission. Valendar is the leader of the Blue Fire Wardens and is held captive in area V7 of the Three Moons Vault.
\end{description}

\subsection{Fifth Rod Piece}
The fifth piece of the \textit{Rod of Seven Parts} is in area U5 in the upper level of the Three Moons Vault. For more information about the rod and the spell this piece allows its wielder to cast, see this book's introduction.

\section{Krynn}
Once the characters decide to seek the fifth piece of the \textit{Rod of Seven Parts}, the portal in the Sigil sanctum takes them to the world of Krynn, where evil plots complicate the party's search.

\subsection{Knowledge of Krynn}
The libraries in Sigil contain the following useful information about Krynn:


\begin{description}
\item[Dragons of Krynn.]
Dragons have a strong presence on Krynn. Chromatic dragons swear allegiance to the expansionistic Dragon Armies of Queen Takhisis. Unusual dragons such as sapphire dragons and lunar dragons pursue hidden agendas.
\item[Ruins of the Cataclysm.]
A world-shattering event called the Cataclysm swept across Krynn centuries ago. The apocalypse destroyed nations, unleashed monsters, and cast Krynn into a dark age.
\item[War Torn.]
Recently, war has spread across Krynn. Infamous warlords such as the elf mage Feal-Thas and the death knight \textbf{Lord Soth} (see appendix B) have taken up arms in their bids for greater power.
\end{description}

\section{Arriving on Krynn}
When the party steps through the Sigil portal and arrives on Krynn, they arrive through a doorway in a massive tree. Read the following:

\begin{DndReadAloud}{}
The rising sun limns the rolling hills around you. In the purple sky hang three moons—a red moon, a white moon, and a black moon discernible only because of the stars it blots out. Behind you stands a tall fruit tree in early bloom, its bright-orange buds still flowering. The doorway through which you've just emerged is a deep, dark furrow in the tree's bark.

The fourth piece of the \textit{Rod of Seven Parts} points to another tree on a nearby hill. That massive tree's ashen bark and leafless branches suggest it is dying.

\end{DndReadAloud}

The fourth rod piece points its wielder to a dying peylon tree. This peylon tree grew to a titanic height thanks to the magic of the rod piece that was previously embedded in its trunk. The rod piece was recently removed, resulting in the tree's current state of decay.

The tree still bears traces of the rod's magic. The fourth piece of the \textit{Rod of Seven Parts} points here due to those traces, and the crimson moonlight and latent magic from the Three Moons Vault warps the artifact's divinatory properties. The characters must investigate the tree and talk to its occupants to learn the fifth rod piece's true whereabouts.

\section{The Peylon Tree}
The characters begin their exploration of the peylon tree at area P1. When they arrive at the tree, read or paraphrase the following:

\begin{DndReadAloud}{}
This tree is hundreds of feet tall. The tree's leafless branches and bark are pale and soft-looking, and the hill is covered in the decomposing remains of sticky, gray fruits that hum with swarms of flies. Deep holes in the tree's trunk—like rents in a rusted suit of armor—provide glimpses inside the tree's rotten, hollow interior.

\end{DndReadAloud}

\subsection{Tree Features}
The peylon tree has the following features:


\begin{description}
\item[Light.]
During the day, the peylon tree's interior is dimly lit by light filtering in from outside.
\item[Trunk.]
The tree's rotten bark is soft and easy to climb both inside and out.
\end{description}

\subsection{Peylon Tree Locations}
The following locations are keyed to map 6.1.

\subsubsection{P1: Rotted Roots}
\begin{DndReadAloud}{}
Thick, gnarled roots spread in all directions. A gaping fissure at the base of the tree forms a rough, arched entrance into the trunk.

\end{DndReadAloud}

A character who examines the ground around the tree can make a DC 14 Wisdom (Perception) or Wisdom (Survival) check. On a successful check, the character spots vague humanoid footprints in the soft dirt. The tracks lead into area P2.

\subparagraph{Alternative Entrance}
A character who scales the roots that form a mound along the east side of the tree can enter through the east side of area P3, 60 feet up.

\subsubsection{P2: Hollow}
\begin{DndReadAloud}{}
A ledge of interwoven roots and packed dirt hangs sixty feet above the hollow's eastern half, and a boulder leans against the hollow's northern wall.

\end{DndReadAloud}

The fourth piece of the \textit{Rod of Seven Parts} points its wielder to the boulder along the tree's interior northwest curve.

If the characters enter this area from area P1, they attract the attention of the treant and spiders dwelling on the ledge above (see area P3).

\subparagraph{Grotto Entrance}
Inspecting the boulder reveals a hole dug into the loamy soil beneath it. The hole connects to a small grotto beneath the tree (area P4). With the boulder atop it, the hole is big enough for a Small or smaller creature to fit through.

A creature can spend 1 minute digging to widen the hole so Medium creatures can fit through. Alternatively, the boulder can be pushed aside by any number of creatures with a combined Strength score of 30 or more.

\subsubsection{P3: Ledge}
\begin{DndReadAloud}{}
This thick, tangled mat of roots overlooks a broad hollow within the peylon tree's trunk.

\end{DndReadAloud}

Rosintar, a neutral evil \textbf{treant} who hates intruders, hides on this ledge above area P2. Two giant spiders lurk in the ledge's corner and follow Rosintar's commands.

While motionless, Rosintar is indistinguishable from the rest of the peylon tree. A character looking out for trouble notices the two giant spiders with a successful DC 17 Wisdom (Perception) check. Shortly after the characters arrive in area P2, Rosintar silently signals for the spiders to sneak up on the party. Any character with a passive Wisdom (Perception) score of 15 or higher notices the spiders readying to attack. As soon as the spiders attack, Rosintar hurls a rock at a random character, then fights viciously to drive the intruders away.

If the treant notices characters arriving on the ledge via the alternative entrance (see area P1), it screeches in surprise and attacks the intruders.

Rosintar fights until reduced to 30 or fewer hit points. At that point, Rosintar surrenders.

\subparagraph{What the Treant Knows}
If the characters allow Rosintar to surrender, the treant tells the party that a cruel dryad named Gazaia dwells in the grotto beneath the tree, accessible via a hole under the boulder (see area P4). The treant bemoans that it didn't notice "the previous intruder"—a small person in a dark-blue cloak—until the intruder had already squeezed into the hole under the boulder. Rosintar doesn't go near the hole, both because the treant wouldn't fit and because it is afraid of Gazaia. Of the blue-cloaked intruder's fate, Rosintar says flatly and with certainty: "The blue-cloaked gnat must be dead."

\subsubsection{P4: Grotto}
It's a 50-foot drop from the hole in area P2 to the floor of this underground chamber. Even during the day, the fetid grotto is dark.

\begin{DndReadAloud}{}
Tangled roots anchor the dirt walls of this damp, subterranean chamber. The rinds of large, rotten fruits litter the floor. A thick taproot hangs from the ceiling. In one corner of the grotto, a kender wearing a blue cloak rummages in the dirt on hands and knees.

\end{DndReadAloud}

The rummaging figure is a kender named Riffel who also happens to be a werewolf (use the \textbf{werewolf} stat block, except Riffel's size is Small and his alignment is neutral good). When Riffel realizes he has company, he jumps to his feet and draws his spear, but he doesn't fight unless attacked. He is initially wary of the characters but is willing to hear them out. If a character mentions the rod piece, or if a character tries to befriend Riffel and succeeds on a DC 13 Charisma (Persuasion) check, Riffel relaxes and puts away his weapon.

\subparagraph{Riffel's Quest}
Riffel says he is searching this grotto for a fresh peylon fruit. He explains that he is a member of a group called the Blue Fire Wardens: naturalists who oppose the death knight Lord Soth. The wardens recently attempted to besiege a group of Soth's followers in a keep called the Three Moons Vault, but the attack failed. Riffel and his fellow survivors fled to the nearby wetlands, but then were attacked by a violent monster.

If the characters mention the rod piece, Riffel says:

\begin{DndReadAloud}{}
"The artifact was taken to the Three Moons Vault, but you'd be foolish to simply walk in and try to take it. I can sneak you into the vault—but first, I need your help."

\end{DndReadAloud}

Riffel explains that he needs the peylon fruit to distract a \textbf{borthak}—a bog monster that attacked the Blue Fire Wardens and trapped Riffel's allies.

Riffel continues his explanation:

\begin{DndReadAloud}{}
"Borthaks love peylon fruit. When we were attacked, I ran here to find a peylon fruit so I could return and distract the borthak so the others can escape. It's not fruit season, but this tree used to be magical, so I thought it might still have fruit. Anyway, if you help me find some, I'll help you find your artifact."

\end{DndReadAloud}

\subsection{The Deadbark Dryad}
As Riffel and the characters talk, the peylon tree's guardian emerges from the grotto's wall and stands before her uninvited guests. This is Gazaia, a \textbf{deadbark dryad} (see appendix A). Annoyed at the disturbance, Gazaia tells the party and Riffel to state their business or leave her sanctum at once.

\subparagraph{The Last Peylon Fruit}
Gazaia says that there is one ripe peylon fruit left in this tree, and she's willing to give it to the party—for a price. Gazaia requests any magic item in exchange for the fruit and will take no other payment.

\begin{DndSidebar}[float=!b]{Kender}
The NPC Riffel is a kender \textbf{werewolf}. On the world of Krynn, kender are a race of Humanoid creatures with pointed ears. They are similar to halflings on other worlds and about the same size. For more information about kender, see Dragonlance: Shadow of the Dragon Queen.



\end{DndSidebar}

\subsubsection{Bargaining with Gazaia}
Gazaia can convey the following points to the party:


\begin{description}
\item[Corruption of Soth.]
Not long ago, Gazaia and her tree were verdant and thriving. A powerful artifact buried in this hill—the rod piece—infused the peylon tree with its magic and enticed Gazaia to become the tree's guardian. When Lord Soth's soldiers recently stole the artifact, Gazaia failed to defend her charge. The tree turned fetid without the artifact, and Gazaia became angry, vengeful, and grief-stricken.
\item[Gazaia's Anger.]
If the characters refuse Gazaia's offer and fail to promptly leave, or if Gazaia catches a character attempting to take the peylon fruit by stealth or force, she attacks the party viciously and fights until destroyed.
\item[Gazaia's Secret.]
The \textbf{deadbark dryad} bitterly recounts the tale of Lord Soth's soldiers assaulting the tree. If a character tries to comfort Gazaia, she confesses that she hid while the soldiers were around the tree. Gazaia feels extremely guilty that she didn't defend her charge. Regardless of the characters' reaction to this revelation, learning it counts as a secret for the purposes of the Power of Secrets rules in this book's introduction.
\end{description}

\subsubsection{Gazaia's Death}
If Gazaia dies before she shares information with the party, the characters find the ripe peylon fruit in her possession.

\subsection{Leaving the Peylon Tree}
Once the party has the peylon fruit, they can leave the peylon tree with Riffel and continue the adventure.

\section{Bittergrass Fen}
To help Riffel save his fellow wardens, the characters must follow the kender \textbf{werewolf} to a marshy lowland called Bittergrass Fen.

\subsection{Journey to Bittergrass Fen}
By foot, the trek to Bittergrass Fen takes an hour. Along the way, Riffel conveys the following:


\begin{description}
\item[Teremini's Vault.]
The Three Moons Vault is a heavily fortified keep in the nearby mountains. It's overseen by the \textbf{archmage} Teremini Nightsedge, who serves the dreaded Lord Soth.
\item[The Stolen Shard.]
Teremini's soldiers stole a magical shard at the heart of the peylon tree. The tree and its dryad became corrupted shortly thereafter. (The characters are certain this is the fifth piece of the \textit{Rod of Seven Parts}.)
\item[The Blue Fire Wardens.]
Riffel is a member of the Blue Fire Wardens: benevolent lycanthropes and trackers who oppose Lord Soth and worship the nature god Habbakuk. The wardens' attack on the Three Moons Vault was a catastrophic failure that resulted in the capture of the group's leader, Valendar. Valendar knows more about the vault and the stolen rod piece than any of the other wardens.
\end{description}

\subsection{Approaching the Fen}
Once the party arrives at Bittergrass Fen, show your players map 6.2. Read the following to describe the scene:

\begin{DndReadAloud}{}
Rising from the fen are two rows of standing stones, each one a rough-carved pillar twenty feet high and ten feet thick. Floating inches above each standing stone is a ten-foot-diameter boulder. West of the fen is a swiftly flowing creek with muddy banks. Slopes to the north and east meet to form a rocky bluff.

An enormous, slavering monster attacks the bluff face, where a stone arch marks the entrance to an underground temple. A crumbling stone door separates the monster from the temple's interior.

\end{DndReadAloud}

The Blue Fire Wardens are trapped inside their own temple. If the characters don't intervene, the \textbf{borthak} (see appendix A) will soon break through the temple door.

The \textbf{borthak} is too large to fit through the stone door and into the temple, though smaller creatures can slip through with some effort. Left alone, the borthak uses its action on each of its turns to attack the temple door. The door has AC 20, 225 hit points, and immunity to poison and psychic damage. If the borthak succeeds in breaking down the door, it enters the temple and attacks the seven wardens (use the \textbf{werewolf} stat block, except their alignments are neutral good) who are trapped inside.

The underground temple isn't shown on map 6.2. Hewn out of the rock and earth, it has a central gathering area and six adjoining cells, where the wardens sleep and pray behind thin wooden doors. The temple contains supplies but nothing of value.

\subsection{Distracting the Borthak}
If one or more characters attack the borthak, it stops attacking the temple long enough to defend itself.

As an action, a character can use a ripe peylon fruit to distract the borthak. If the fruit is hurled toward the borthak, the monster moves toward the fruit on its next turn and uses its action to devour it. If a character holding the fruit moves within 20 feet of the borthak, the borthak pursues and attacks that character, eager to obtain the fruit.

While characters contend with the \textbf{borthak}, Riffel moves quickly toward the temple door and crawls through a narrow gap underneath it. The gap is just big enough for a Small character to squeeze through. Once inside, Riffel urges the other wardens to evacuate the temple while the borthak is distracted. On Riffel's next turn, he and the wardens push open the temple door so everyone can escape. Once outside, they skirt along the bluff, heading north. If one or more characters assist with the evacuation, they can keep an eye on the borthak and distract it if necessary while Riffel helps the wardens get to safety.

\subsection{Bittergrass Fen Features}
The fen has features the characters can use to their advantage in a confrontation with the \textbf{borthak}:


\begin{description}
\item[Floating Boulders.]
The boulders floating above the stone pillars are held aloft by ancient druidic magic. Casting \textit{Dispel Magic} on a boulder causes it to fall and tumble to the ground. As an action, a character within reach of a floating boulder can try to push it, doing so with a successful DC 15 Strength (Athletics) check. Pushing the boulder ends the magic on it and causes it to fall in whichever direction the character prefers. Any creature in the path of a falling boulder must succeed on a DC 15 Dexterity saving throw to avoid it, taking 11 (2d10) bludgeoning damage on a failed save.
\item[Muddy Bank.]
The ground within 20 feet of the creek's eastern shore has turned to sticky mud and is \textit{difficult terrain}.
\end{description}

\subsection{Escaping the Fen}
Once all the wardens escape the temple, each warden transforms into a wolf and flees from the borthak. End the characters' encounter with the borthak when they defeat it or after they and all the wardens have escaped, whichever happens first.

\section{The Blue Fire Wardens}
Once the wardens are safe, a human werewolf named Argentia Skywright (use the \textbf{werewolf} stat block, except her alignment is neutral good) thanks the characters on behalf of her group. She expands on what Riffel told the party earlier about the Three Moons Vault, conveying the following points:


\begin{description}
\item[Blue Fire Attack.]
Shortly after Teremini Nightsedge's soldiers seized the rod fragment from the peylon tree, the Blue Fire Wardens assaulted the Three Moons Vault. They refer to their attack as the Night of Blue Fire. Teremini, a renegade from the Mages of High Sorcery's Order of the Red Robes, foresaw the attack and used powerful lunar magic to create a shroud of transformative red moonlight around the vault. Werewolves that step into this moonlight lose control of their powers and succumb to raw, animalistic instinct, making it easier for Teremini to manipulate or trap them.
\item[Mage's Ritual.]
Teremini learned her red moonlight magic from a lunar dragon named Orinix. Orinix also taught Teremini a ritual to make the wall of transformative moonlight permanent and tasked her with completing this lengthy, taxing ritual. While she focuses on the ritual, powerful magical barriers surround Teremini and her ritual components, which include the fifth rod piece.
\item[Stopping the Ritual.]
Argentia insists there must be a way to disrupt Teremini's ritual so the characters can seize the rod piece. She encourages the characters to find and rescue Valendar, the wardens' imprisoned leader, who might know how to stop Teremini's ritual.
\end{description}

\subsection{Wardens' Aid}
Argentia gives the characters a \textit{Spell Scroll} of \textit{Moonbeam}. She also explains that if Valendar is stuck in his wolf or hybrid form, they must bring Valendar close to death before he'll return to his true form, which is a human.

Riffel agrees to take the characters to the Three Moons Vault's secret entrance. Before the party leaves, Argentia leads a quick rite to cast a protective ward over Riffel. For the next 12 hours, Riffel and allies within 10 feet of him are immune to the Forced Transformation effect of Teremini's curtain of red moonlight (see the "Environmental Effects" subsection of the "Three Moons Vault" section). The wardens have enough power to cast this ritual only once, so it's imperative Riffel return to Bittergrass Fen with Valendar as soon as possible.

\section{Three Moons Vault}
The Three Moons Vault is a remote dungeon complex in the mountainous outskirts of Lord Soth's domain.

Millennia ago, a chunk of the red moon Lunitari sheared off, broke into meteorites, and rained down on Krynn. The meteorites' impact formed craters throughout the Northern Dargaard Mountains. Ancient, moon-venerating wardens chose one of these craters as the site for a shrine where they could revere Krynn's moons. The wardens built three magical towers: a spacious white tower for Solinari, a handsome red tower for Lunitari, and a solemn black tower for Nuitari.

The Cataclysm destroyed the white tower, disrupting the balance of magic between the three towers. Lord Soth's forces from nearby Dargaard Keep assumed control of the ruins. They built fortified walls around the site and carved out a subterranean vault to contain prisoners and treasure. The vault is now under the purview of a servant of Lord Soth and operates as a prison for the overlord's enemies and a safe house for his underlings.

\subsection{Teremini Nightsedge}
A lawful evil, elf \textbf{archmage} named Teremini Nightsedge oversees the Three Moons Vault. Though a loyal servant of Lord Soth, Teremini has her own agenda for the ancient moon towers. When she was young, Teremini tried to join the wardens but was rejected after her reverence of Lord Soth became clear. Now, she hopes to complete her revenge by permanently shrouding the vault in crimson moonlight, kidnapping the Blue Fire Wardens she hates so deeply, and bringing them here to be slain, thereby completing her revenge. Thanks to her recent acquisition of a piece of the \textit{Rod of Seven Parts}, Teremini has commenced her plan.

\subsection{The Characters' Goals}
The characters have three goals while inside the Three Moons Vault. As the characters explore the vault, NPCs such as Valendar (area V7) and Casivus (area V15) can help devise plans for achieving these goals.

\subsubsection{Rescue Valendar}
Valendar, the leader of the Blue Fire Wardens, is trapped in the vault and might be stuck in his wolf or hybrid form. The characters need to find him, disrupt the magic forcing him to remain transformed, and restore him to his human form.

\subsubsection{Stop Teremini's Ritual}
If the characters rescue Valendar, he tells them that Teremini needs three special crystals in specific places to complete her ritual. If the characters find and sabotage the special crystals, they'll disrupt Teremini's ritual.

\subsubsection{Get the Rod Piece}
Teremini is using a piece of the \textit{Rod of Seven Parts} as part of a ritual that would permanently shroud the vault and a 1-mile radius around it in magical crimson moonlight. The characters need to disrupt Teremini's ritual before they can retrieve the rod piece. They can then return to Sigil.

\subsection{Environmental Effects}
The magical red moonlight that surrounds the Three Moons Vault is temporary, although Teremini is working to make the effect permanent. The moonlight creates the following effects in a 1-mile radius around the Three Moons Vault. These effects are suppressed in the vault's locations.

\subsubsection{Forced Transformation}
A \textbf{werewolf} that enters the red moonlight changes into its wolf or hybrid form, and it has the poisoned condition as long as it is within 1 mile of Three Moons Vault. As long as it remains in this area, the creature can't willingly change shape unless it is reduced to 10 hit points or fewer; a \textit{Remove Curse} spell suppresses Forced Transformation for 1 hour.

\subsubsection{Reduced Gravity}
Objects and creatures take half as much damage from falls. Creatures can jump twice as far as normal.

\subsection{Entering the Vault}
The Three Moons Vault is built amid rocky crags in the Northern Dargaard Mountains. Its main entrance is in plain sight and heavily guarded, but Riffel knows the way to a second, less-protected passage at the end of an old goat path. He takes the characters to this entrance, then he waits there so he can escort Valendar to Bittergrass Fen. It takes about one hour to traverse the distance between the fen and the vault.

The characters begin their exploration of the vault at area V1.

\subsection{Vault Features}
Maps 6.3 and 6.4 depict the Three Moons Vault. The vault is a complex, castle-like structure with a dungeon level, several towers, a secret entrance through the cliffs, and a drawbridge over a moat. Unless the text says otherwise, areas in and around the Three Moons Vault share a number of features.

\subsubsection{Ceilings}
The vault's ceilings are 10 feet high.

\subsubsection{Doors}
The vault's doors are stone slabs set on brass hinges. Unless otherwise noted, each door is unlocked.

\subsubsection{Lighting}
Unless otherwise noted, interior areas are dark. Area descriptions assume the characters have a light source or some other means of seeing in the dark.

\subsubsection{Moonlight Mirrors}
Several areas in the vault contain objects known as moonlight mirrors, which each reflect a powerful beam of light from one of Krynn's moons. Moonlight mirrors can be used to disrupt Teremini's ritual, as described later in this chapter.

A moonlight mirror is a Medium object with AC 13; 5 hit points; and immunity to poison, psychic, and radiant damage. A \textit{Detect Magic} spell reveals an aura of evocation magic emanating from the mirror. Unless covered, the reflective side of a moonlight mirror casts bright light in a 20-foot hemisphere, in the color of the moon the mirror reflects. The text notes which ability checks, if any, are required to remove a moonlight mirror from where it is found.

\subsection{Vault Locations}
The following locations are keyed to map 6.3.

\subsubsection{V1: Entrance}
Around the corner from where Riffel leaves the party is the vault's entrance. Read the following:

\begin{DndReadAloud}{}
The goat path dead-ends at the convergence of three steep cliff walls. Two austere pillars flank a stone double door set into the wall. Next to each pillar stands a silver statue of a tall, mostly human-shaped individual.

\end{DndReadAloud}

Two moonlight guardians (see appendix A) guard this secret entrance. They attack any intruders on sight and fight until destroyed.

\subparagraph{Trapped Doorway}
A \textit{Detect Magic} spell reveals an aura of evocation magic radiating from the doorway. If a creature pulls on either of the doors' crescent-shaped handles, the doorway emits a 30-foot cone of silvery light for 1 minute, centered on the middle of the double door. A creature that enters that area for the first time on a turn or starts its turn there must make a DC 17 Constitution saving throw, taking 27 (5d10) radiant damage on a failed save or half as much damage on a successful one. A creature not in its true form has disadvantage on this save; if it fails its save, the creature instantly reverts to its true form and can't assume a different form until it leaves the light.

Once this trap is triggered, it can't be triggered again.

\subsubsection{V2: Anteroom}
\begin{DndReadAloud}{}
A stone statue of an elegant elf in flowing robes watches over this chamber.

\end{DndReadAloud}

An inscription at the base of the statue reads, in Elvish: "It is not surrender. Live to fight again another day."

\subparagraph{Treasure}
A weapon rack in the room holds a black steel \textit{+1 Longsword} and a steel shield bearing the stylized floral emblem of the Order of the Rose. An armor stand bears a black breastplate. A small, unlocked chest holds a \textit{Potion of Vitality}.

\subsubsection{V3: Escape Tunnel}
\begin{DndReadAloud}{}
This tunnel stretches one hundred feet before ending at a flat wall.

\end{DndReadAloud}

The secret door in the wall is obvious: emblazoned on the masonry is a wide circle of dimly glowing runes. A character who examines the secret door immediately realizes they can trace the runes to cause that circular section of wall to disappear, revealing the remainder of the tunnel on the other side. The wall re-forms a few seconds after it disappears.

The secret door can be activated from either side of the wall, but the runes are visible only when viewed from the north. From the southern side of the door, a creature that searches the wall and succeeds on a DC 20 Intelligence (Investigation) check can see the faintest traces of the runes' light through the wall's brickwork. The creature can trace those runes to activate the door.

\subsubsection{V4: Halls}
\begin{DndReadAloud}{}
Four ghostly soldiers patrol these hallways. Their eyes are red, and each wears spectral, ridged armor and carries a hollowed-out goat horn.

\end{DndReadAloud}

The creatures are four wraiths. When Akaazi (see area V34) isn't present, these wraiths lead the small army of Undead soldiers stationed in the underground garrison (area V8). A wraith named Guerthel is Akaazi's favorite; he carries an iron key that unlocks the doors to areas V7 and V10.

\subparagraph{Raising the Alarm}
If a wraith spots an intruder or is attacked, it uses its action to blow its horn. Creatures within 100 feet of that wraith can hear the horn, prompting the wraith's allies in area V8 to join the fight. A horn blown in the dungeon level can't be heard from the keep above and vice versa.

\subsubsection{V5: Scrying Chamber}
\begin{DndReadAloud}{}
A fifteen-foot-diameter sphere made of shimmering silver floats above a hexagonal dais on the west side of this thirty-foot-high chamber. The sphere is as reflective as a mirror, and it appears smooth and sleek.

\end{DndReadAloud}

\subparagraph{Scrying Sphere}
The silver sphere is a scrying device that allows Teremini to remotely communicate with Lord Soth, who has a similar sphere in his home, Dargaard Keep.

A creature that touches the seemingly solid sphere discovers that it is made from harmless, viscous silver liquid that is cool to the touch. Each time the sphere is touched, Lord Soth hears a faint pinging noise. Soth can respond by standing before his sphere and speaking. When he does, his helmeted face manifests, but it's composed of the sphere's silvery goo. (For more about \textbf{Lord Soth}, see appendix B.)

Unknown to Teremini, Lord Soth can send a weaker facsimile of himself through this scrying sphere. If the characters respond to Lord Soth's manifested face with anything other than deference, a \textbf{death knight} made from the sphere's silvery metal emerges from the sphere. The death knight inflicts punishment in Soth's stead, following his directives but requiring no action on Soth's part. If the death knight is destroyed, the sphere is also destroyed and disappears.

\subsubsection{V6: Empty Cells}
\begin{DndReadAloud}{}
Iron bars form three empty cells along the walls of this room. A nondescript stone door to the east indicates the room's exit. Lupine howls of pain sound from beyond the iron door to the west.

\end{DndReadAloud}

The howls can be clearly heard coming from the other side of the locked iron door (area V7). The door can be unlocked with Guerthel's key (see area V4), or a character using thieves' tools can use an action to try to pick the lock, doing so with a successful DC 16 Dexterity (Sleight of Hand) check.

\subsubsection{V7: Valendar's Cell}
When the characters enter this room, read the following:

\begin{DndReadAloud}{}
Bright-red light fills this twenty-foot-tall oval chamber. The light emanates from a large circular mirror set into a specially made groove on the ceiling. Kneeling on the floor of the chamber is a werewolf. It howls terribly, clutching and clawing at its body in fury. At the sight of the cell door opening, the werewolf leaps up and sprints toward you, claws extended!

\end{DndReadAloud}

The red moonlight from the magical mirror forced Valendar (use the \textbf{werewolf} stat block, except his alignment is chaotic good) to transform into his hybrid form. Valendar's imprisonment, anger, and fear make it difficult for him to think clearly. He attacks the characters, thinking they are his tormentors.

\subparagraph{Red Moonlight Mirror}
The mirror in the ceiling sheds bright-red light in the whole room. Removing the mirror from the ceiling requires no special effort.

\subparagraph{Saving Valendar}
Valendar can't transform into his true form as long as he is in the area of moonlight emanating from the mirror. Valendar reverts to his true form if he is reduced to 10 hit points or fewer.

\subparagraph{What Valendar Knows}
If the characters help Valendar transform back into his true form, he thanks them wholeheartedly and shares everything he knows about the Three Moons Vault. He conveys the following points:


\begin{description}
\item[General Layout.]
Valendar describes the general layout of the Three Moons Vault. He also explains that each of the vault's three towers is topped with a special room called a lunarium, which is an observatory containing a model.
\item[Lunariums and Crystals.]
Within each lunarium is a "lunar crystal" infused with the moonlight of one of Krynn's three moons. These crystals are essential to Teremini's ritual and are protected by the ritual's magic.
\item[Magic Balance.]
Krynn's moons represent a delicate balance of magical power, and Teremini's ritual relies on this balance. If the lunar crystals' magic is unbalanced, the ritual will be stopped.
\item[Moonlight Mirrors.]
The vault contains several other mirrors like the one in Valendar's cell. Valendar suggests the characters find three moonlight mirrors, one for each moon, then shine their light onto different lunar crystals to unbalance the crystals' magic.
\end{description}

Valendar needs to return to his fellow wardens. After he's spoken to the characters, he departs, sneaking through the vault's secret escape tunnel to meet Riffel outside.

\subparagraph{Valendar's Secret}
As the characters talk with Valendar, the leader adopts a sheepish demeanor, especially when talking about his capture. If the characters ask why, Valendar admits that he didn't properly scout the Three Moons Vault before the wardens' attack on the castle. Valendar underestimated Teremini, assuming she was a powerless lackey of Lord Soth. Valendar believes the wardens' defeat was his fault, and he's glad that no one else was captured in the assault.

Regardless of the characters' reaction to this revelation, learning it counts as a secret for the purposes of the Power of Secrets rules in this book's introduction.

\subsubsection{V8: Garrison}
\begin{DndReadAloud}{}
All manner of undead creatures—mostly skeletons and zombies—impatiently mill around this area. They look ready for a fight.

\end{DndReadAloud}

The vault's Undead soldiers occupy this grim chamber, ready for deployment at a moment's notice. Unless the soldiers have been called elsewhere to defend the vault or perform drudge work, the garrison contains sixteen skeletons, nine zombies, and two ogre zombies.

\subsubsection{V9: Winding Staircase}
A spiral staircase occupies most of this small chamber. Short hallways lead east and west from the staircase, which ascends 50 feet before connecting to area V32.

\subsubsection{V10: Trapped Black Rose Bearer}
Both doors to this room are locked. A character can take the key from Guerthel (see area V4), or as an action, a character can use thieves' tools to try to pick the lock, doing so with a successful DC 10 Dexterity (Sleight of Hand) check. If the characters enter this room, read the following:

\begin{DndReadAloud}{}
Though this tiered chamber was once elegantly appointed, its brass candelabras, velvet runners, and red tapestries were destroyed—likely by the room's occupant, a desiccated Undead currently tearing at the walls. A single doorway is set into the short north wall of the room.

\end{DndReadAloud}

The \textbf{black rose bearer} (see appendix A) trapped here is raging at everything in sight. Akaazi (see area V34) locked the bearer in here to prevent it from wreaking havoc throughout the dungeon. The black rose bearer attacks anyone who enters and fights until destroyed.

The doorway in the north wall opens to a stairway ascending to area V35.

\subsubsection{V11: Vault Access}
\begin{DndReadAloud}{}
Two minotaur skeletons guard this curved hallway. Beyond them, heavy, round metal doors seal five doorways. A brown stain covers the ground in front of the northernmost door.

\end{DndReadAloud}

The two minotaur skeletons attack intruders on sight and fight until destroyed.

\subsubsection{V12: Treasure Vaults}
Each of these vault rooms is sealed behind a locked iron door. As an action, a character can use thieves' tools to try to unlock a door, doing so with a successful DC 18 Dexterity (Sleight of Hand) check, or force open a locked door, which requires a successful DC 18 Strength (Athletics) check.

\subparagraph{V12a: Treasures}
The door to this vault is trapped with a magical glyph. A character who examines the door can find the glyph by succeeding on a DC 17 Intelligence (Investigation) check, and any character can disable the glyph by succeeding on a DC 20 Intelligence (Arcana) check. When a creature other than Teremini opens the door, a 20-foot-radius sphere of fire explodes from the glyph. Each creature in that area must make a DC 17 Dexterity saving throw, taking 22 (5d8) fire damage on a failed save or half as much damage on a successful one. Once this trap is triggered, it can't be triggered again.

This vault contains 1,900 gp, ten gems worth 100 gp each, and six paintings worth 250 gp each.

\subparagraph{V12b: Guardian Chamber}
A \textbf{beholder zombie} waits inside this vault. Akaazi keeps this zombie apart from the other Undead and reserves it for special tasks. It has orders to attack non-Undead creatures other than Akaazi and Teremini.

\subparagraph{V12c: Empty Vault}
This vault is empty.

\subparagraph{V12d: Leader's Gear}
Heaped on this vault's floor are Valendar's belongings, including his leather armor, a Potion of Greater Healing, a \textit{Spell Scroll} of \textit{Stoneskin}, and a brilliant-blue cloak emblazoned with the holy symbol of Habbakuk.

\subsubsection{V13: Ruined Vault}
\begin{DndReadAloud}{}
Fuzzy brown mold covers this area like carpeting. A pile of rubble slopes upward to form a rough stairway.

\end{DndReadAloud}

In addition to brown mold's usual effects, the areas of brown mold in this chamber are \textit{difficult terrain}. A slope made of rubble connects this room to the white tower's basement (area V14) 20 feet above.

\subsubsection{V14: Flooded Basement}
\begin{DndReadAloud}{}
Several inches of standing water cover the floor of these ruined rooms. Rubble in the northeast corner slopes downward to form a rough stairway.

\end{DndReadAloud}

Stone walls once separated this basement level into four rooms, but the Cataclysm destroyed the north and west chambers. The standing water makes the floor difficult terrain in this area. The sloping rubble descends to area V13.

\subsubsection{V15: Lunar Shrine}
\begin{DndReadAloud}{}
A larger-than-life alabaster statue of a humanoid in flowing robes watches over this shrine. The figure bears a moonlike sphere instead of a face. The sphere is illuminated by white light shining from a silver mirror in the crook of the statue's left arm. Coiled around the plinth is an enormous, green, snakelike creature with a human face. It stares shrewdly at you as you enter.

\end{DndReadAloud}

The statue is a depiction of Solinari, the god of Krynn's white moon and good magic.

\subparagraph{White Moonlight Mirror}
The mirror in the statue's arm is a white moonlight mirror. Removing the mirror from the statue requires no special effort.

\subparagraph{Shrine's Guardian}
Wrapped around the statue is a \textbf{guardian naga} named Casivus. Casivus has guarded this shrine since long before Lord Soth's forces arrived, and it chafes at the evil doings of Teremini and her minions.

The naga is initially indifferent to the characters. If a character mentions that the party opposes Teremini's plans or engages the naga in polite conversation and succeeds on a DC 14 Charisma (Persuasion) check, Casivus becomes friendly toward the group. In this case, the naga permits the characters to take the mirror from the statue. Casivus also imparts two useful points of information:


\begin{description}
\item[Position of the Moons.]
To disrupt Teremini's ritual, the characters must shine different colors of moonlight onto the lunariums' crystals in a specific combination. This combination depends on the current phases and positions of Krynn's moons. Teremini has a magical orrery in the red tower that shows the moons' positions. A character can study this orrery to determine the moonlight-crystal combinations necessary to disrupt the ritual.
\item[Stairwell Passphrase.]
Casivus tells the characters the passphrase to open the locked door inside the chamber northeast of here (area V16). The passphrase is "buried is best."
\end{description}

Casivus has no wish to leave this shrine or accompany the characters, but the naga would be pleased to see Teremini and her followers gone.

\subsubsection{V16: Locked Stairwell}
\begin{DndReadAloud}{}
An enclosure with a shuttered door stands in the southernmost part of this otherwise empty room.

\end{DndReadAloud}

The door to this stairwell is locked with an \textit{Arcane Lock} spell. Teremini, Akaazi, and Casivus know the door's passphrase ("buried is best"). As an action, a character can try to force open the door, doing so with a successful DC 24 Strength (Athletics) check.

The circular staircase connects to area V17, which is 30 feet above.

\subsubsection{V17: White Tower Ruins}
\begin{DndReadAloud}{}
The walls and floor here are crumbling. A wooden ladder leans against the edge of the ceiling above.

\end{DndReadAloud}

The white tower's ground floor lies in ruin. A ladder connects this floor to area V19.

\subsubsection{V18: Veteran Camp}
\begin{DndReadAloud}{}
A sloping mound of rubble creates a rough bridge over the moat west of the white tower. Six humans in dented armor pace atop the slope.

\end{DndReadAloud}

The slope immediately leading into this area is guarded by six veterans (lawful evil humans). A seventh veteran is resting in the tent to the north and rushes to help their allies when a fight breaks out. The veterans defend the keep with their lives, knowing that to do otherwise would invite Teremini's wrath.

\subsubsection{V19: White Tower Lookout}
\begin{DndReadAloud}{}
A single human lookout in dented armor stands watch in the southern portion of this area.

\end{DndReadAloud}

The guard is a lawful evil, human \textbf{veteran}. The veteran carries a hollow goat's horn in addition to her usual gear. If she expects attackers, she pulls up the ladder that connects to area V17 and blows her horn to raise the alarm.

A spiral staircase ascends to area V20.

\subsubsection{V20: Ruined Shrine}
\begin{DndReadAloud}{}
Ornate pillars bearing countless impressions of celestial bodies encircle this ruined chamber. The floor tiles are cracked and loose, and part of the wall has caved in.

\end{DndReadAloud}

A staircase in the southern part of the room ascends to the white tower's lunarium (area U1).

\subparagraph{White Moonlight Mirror}
Hidden in a cavity beneath the flagstone tiles is a white moonlight mirror wrapped in old oilcloth. A character can find the mirror by searching the floor and succeeding on a DC 12 Intelligence (Investigation) check. Removing the mirror from the cavity requires no special effort.

\subsubsection{V21: Drawbridge}
\begin{DndReadAloud}{}
The castle's drawbridge is lowered over the murky moat.

\end{DndReadAloud}

Soth's rank-and-file soldiers use this gate to enter the keep.

\subparagraph{Drawbridge}
To raise or lower the wooden drawbridge, both levers inside the keep must be pulled at the same time. An automatic mechanism then raises or lowers the bridge over the course of 1 minute.

\subparagraph{Moat}
The moat surrounding the Three Moons Vault is 50 feet deep. A nearby stream keeps the moat filled with ice-cold snowmelt from the mountains.

\subsubsection{V22: Lookout Turret}
\begin{DndReadAloud}{}
A ladder allows access from the courtyard to this fortified platform. Two guards in dented armor, each carrying hollowed-out goat's horns, keep watch here.

\end{DndReadAloud}

Two veterans (lawful evil humans) keep watch atop the turret. If they see or hear the characters approach, they sound their horns, and allies from area V18 arrive in 5 minutes.

\subsubsection{V23: Wall Walk}
The stone walkway atop the keep's curtain wall passes through the second floor of each moon tower, allowing access. The walkway can be reached via staircases in the courtyard (area V25).

\subsubsection{V24: Rookery}
\begin{DndReadAloud}{}
Two enormous, skeletal, birdlike creatures perch atop the wooden platform here.

\end{DndReadAloud}

These two bone rocs (see appendix A) are trained to ferry passengers to and from a similar rookery near the bottom of the mountain.

A bone roc attacks only if threatened or harmed.

\subsubsection{V25: Courtyard}
\begin{DndReadAloud}{}
The center of the keep features an open-air plaza paved with moon rock tiles.

\end{DndReadAloud}

From this courtyard, individuals can access the ground floors of the three moon towers, most of the keep's other buildings, and staircases to the curtain wall walkways. A hundred feet above looms the moondisk (area U5).

\subsubsection{V26: Barrack}
\begin{DndReadAloud}{}
This sparsely furnished area includes simple beds and empty wash tubs.

\end{DndReadAloud}

When non-Undead individuals stay at the Three Moons Vault, they reside here.

\subsubsection{V27: Armory}
\begin{DndReadAloud}{}
Weapon racks and armor stands line the walls of this cramped storage room.

\end{DndReadAloud}

\subparagraph{Treasure}
The equipment here includes the following:


\begin{itemize}
\item Three Longsword
\item Two Halberd
\item Two Crossbow Bolt Case, each containing 20 Crossbow Bolts (20)
\item Two Shield
\item A suit of Splint Armor mail
\item A bottle of \textit{Oil of Sharpness}
\end{itemize}

\subsubsection{V28: Royal Quarters}
\begin{DndReadAloud}{}
A grim iron throne and an iron statue of an intimidating knight stand in the middle of this room. Both look out a window at the courtyard to the southeast. Fastened to the throne's headrest is a mirror that seems to draw light in rather than reflect it.

\end{DndReadAloud}

On the rare occasions when Lord Soth visits the Three Moons Vault, this chamber is set aside for him to glower over his subjects.

\subparagraph{Black Moonlight Mirror}
The mirror affixed to the throne is a black moonlight mirror. Removing the mirror from the throne requires no special effort.

\subsubsection{V29: Orrery}
\begin{DndReadAloud}{}
A large model of a planet and three moons hangs in the center of this high-ceilinged chamber. The celestial bodies are affixed to bronze rings via delicate rods that move slowly under their own momentum. The device is illuminated by a beam of red moonlight that emanates from a mirror mounted twenty feet up the southwest wall. Twin staircases ascend to a railed walkway overlooking the orrery.

\end{DndReadAloud}

This enchanted orrery accurately depicts the current positions of Krynn's three moons—Solinari, Lunitari, and Nuitari—relative to Krynn. If the moonlight mirror in area V30 is removed, the orrery stops moving.

Staircases on either side of the orrery lead to area V30 above.

\subparagraph{Deciphering the Orrery}
A character can study the orrery to deduce the relative positions of Krynn's moons. The character gleans that, currently, Lunitari prevails over Solinari, Solinari prevails over Nuitari, and Nuitari prevails over Lunitari.

Based on this deduction, the character understands that Teremini's ritual can be disrupted by shining red moonlight on the white crystal, white moonlight on the black crystal, and black moonlight on the red crystal.

\subsubsection{V30: Red Mirror Stairwell}
\begin{DndReadAloud}{}
A stone guardrail protects anyone standing on this elevated walkway from falling to the floor below. One wall of the walkway bears a shiny red mirror. Stairs in the center of the space ascend to another chamber.

\end{DndReadAloud}

\subparagraph{Red Moonlight Mirror}
The mirror affixed to the wall is a red moonlight mirror. Removing the mirror from the wall requires no special effort.

\subsubsection{V31: Teremini's Quarters}
\begin{DndReadAloud}{}
This bedroom includes a four-poster bed, a writing desk, several bookshelves, and a footlocker.

\end{DndReadAloud}

Teremini dwells in this finely decorated bedroom.

\subparagraph{Mimic}
The footlocker is a \textbf{mimic}. Teremini trained the mimic to attack anyone who snoops around her room.

\subparagraph{Logbook}
On the desk is a log of all the treasures and prisoners transported to and from the Three Moons Vault. Of interest to the characters are the first and last pages of the log.

The first page describes the six moonlight mirrors found in the original moon towers during the vault's construction. The log notes the present location of these mirrors:

\textbf{White Mirrors}: The lunar shrine (area V15) and the ruined shrine (area V20)

\textbf{Red Mirrors}: Valendar's cell (area V7) and the red mirror stairwell (area V30)

\textbf{Black Mirrors}: Royal quarters (area V28) and Akaazi's quarters (area V36)

The log's last pages mention "a piece of the legendary \textit{Rod of Seven Parts}, seized from an ungrateful dryad" as well as "three lunar crystals, procured from a powerful dragon." The logbook doesn't specify where the rod piece or the crystals are currently stored.

A staircase at the northern end of the room ascends to the red tower's lunarium (area U2).

\subsubsection{V32: Stairwell}
This stairwell connects to area V9 below.

\subsubsection{V33: Muster Yard}
\begin{DndReadAloud}{}
This empty yard has a floor of thick mud.

\end{DndReadAloud}

Teremini's soldiers conduct training drills in this muddy yard. When the keep's alarm is raised, they rally here and prepare to defend the keep. The mud is \textit{difficult terrain}.

Two earth elementals, summoned and bound by Teremini, dwell beneath the mud. If they detect intruders with their tremorsense, the elementals rise from the mud and attack.

\subsubsection{V34. Circle of Undeath}
\begin{DndReadAloud}{}
An interior circular stone wall creates a ringed hallway inside this tower. Archways in the wall are covered in black curtains. The outer hallway is eerily empty, but you can hear chanting coming from beyond the archways.

\end{DndReadAloud}

Three archways, each shrouded by a black curtain, connect the two spaces in this area. Once the characters walk through a curtain into the inner space, read the following:

\begin{DndReadAloud}{}
Black candles cast dim light throughout this inner room, revealing a circle of unholy runes drawn in chalk on the floor. A black-robed figure crouches, examining the runes, while a zombie-like creature holding a black rose in a glass jar stands nearby.

\end{DndReadAloud}

A neutral evil human \textbf{necromancer wizard} (see appendix A) named Akaazi uses this chamber to conduct necromantic rituals and create Undead soldiers for Lord Soth's army. Akaazi is accompanied by a \textbf{black rose bearer} (see appendix A).

As soon as she's aware of intruders, Akaazi commands the bearer to attack and fights to the death in the name of Lord Soth.

If the characters look up, they can see a black moonlight mirror hanging in area V36.

\subsubsection{V35: Temple}
\begin{DndReadAloud}{}
A stone altar at the back of this temple is its only notable fixture. Dusty windows set in iron frames overlook the courtyard.

\end{DndReadAloud}

Akaazi's eerie black tower leads into this temple, so the chamber is rarely used. The north staircase descends 30 feet to area V10.

\subparagraph{Secret Door}
A section of this chamber's western wall is hollow and set onto hinges that allow it to swing out toward the courtyard. The secret door can be detected with a successful DC 13 Wisdom (Perception) check.

\subparagraph{Treasure}
The back of the altar is hollow and bears a single shelf, on which sits a \textit{Wand of Enemy Detection}.

\subsubsection{V36: Akaazi's Quarters}
\begin{DndReadAloud}{}
This room contains a simple sleeping mat and a bookcase full of grim tomes of little value. Two doors exit the tower and connect to the wall walk outside.

\end{DndReadAloud}

\subparagraph{Black Moonlight Mirror}
A circular opening in the floor allows a creature to view the chamber below (area V34). A black moonlight mirror is suspended from three iron chains connected to the sides of the opening. The mirror can be safely detached from one iron chain with a successful DC 15 Dexterity (Sleight of Hand) check. If a character fails this check by 5 or more, the mirror comes loose from the chains and falls to the floor 20 feet below.

\subsubsection{V37: Aviary}
\begin{DndReadAloud}{}
Black corvids fill the small cages in this bleak, windowless aviary. The birds squawk loudly whenever anyone enters.

\end{DndReadAloud}

Twenty ravens used to deliver messages occupy the flimsy wooden cages.

A staircase on the eastern wall ascends to the black tower's lunarium (area U3).

\subsection{Upper-Level Features}
In addition to the features of the Three Moons Vault listed earlier in this chapter, the vault's upper level has a number of important features.

If the characters don't know how to stop Teremini's ritual yet, one of them realizes that the answer must have to do with Krynn's moons, and that they should examine the orrery in area V29.

\subsubsection{Colored Glass}
Each lunarium (areas U1, U2, and U3) is surrounded by six columns that support a hemispherical rooftop made of enameled glass. Panes of enameled glass are set between each column. Each 10-foot square of glass is a Large object with AC 10, 10 hit points, vulnerability to bludgeoning and thunder damage, and immunity to poison and psychic damage.

The glass is enchanted so that only one color of moonlight shines through it:


\begin{itemize}
\item Solinari's white light shines into the white tower's lunarium (area U1).
\item Lunitari's red light shines into the red tower's lunarium (area U2).
\item Nuitari's black light shines into the black tower's lunarium (area U3).
\end{itemize}

\subsubsection{Lunar Crystals}
Shallow steps ascend to a dais at the center of each lunarium. A specially carved pillar atop the dais bears a lunar crystal. These lunar crystals were specially created by the lunar dragon Orinix.

These crystals are components for Teremini's ritual, but they have no value or special properties. When the characters arrive, these daises are surrounded by barriers of solid moonlight (see below). The crystals can't be smashed while Teremini concentrates on her ritual.

\subsubsection{Solid Moonlight}
The moonbridges and moondisk (areas U4 and U5) are made of solid moonlight that shimmers like opalescent glass. If a creature not in its true form starts its turn touching solid moonlight, the creature must succeed on a DC 15 Constitution saving throw or have the poisoned condition until the start of its next turn.

While Teremini concentrates on her ritual, spherical barriers of solid moonlight also surround each lunar crystal dais.

Nothing can physically pass through solid moonlight or teleport through it. It is immune to all damage and can't be dispelled by the \textit{Dispel Magic} spell. Any solid moonlight caught in the area of a \textit{Sunburst} spell is dispelled for 10 minutes. Solid moonlight also extends into the Ethereal Plane, blocking ethereal travel through it. The ritual is powered by the rod piece and is therefore unaffected by the \textit{Antimagic Field} spell.

\subsection{Upper-Level Locations}
The following locations are keyed to map 6.4.

\subsubsection{U1: White Lunarium}
\begin{DndReadAloud}{}
White moonlight glows around this smooth dais.

\end{DndReadAloud}

The moonlight barrier around this lunarium's dais is trapped. When the barrier is touched, it emits a blinding flash of white light in a 30-foot radius centered on the dais. Any creature in that area must succeed on a DC 17 Constitution saving throw or have the blinded condition for 24 hours.

\subsubsection{U2: Red Lunarium}
\begin{DndReadAloud}{}
A shadowy, skeletal wolf walking upright like a human stalks the area here.

\end{DndReadAloud}

The skeletal wolf is a \textbf{deathwolf} (see appendix A) created from the body of a slain Blue Fire Warden. Akaazi ordered the deathwolf to attack any intruders during Teremini's ritual. If it sees the characters, the deathwolf attacks viciously and stops at nothing to destroy its opponents.

\subsubsection{U3: Black Lunarium}
\begin{DndReadAloud}{}
Two desiccated, undead creatures, each holding a bell jar containing a black rose, stand guard here.

\end{DndReadAloud}

Two black rose bearers (see appendix A) have been ordered to defend this lunarium against intruders.

\subsubsection{U4: Moonbridges}
Three bridges of solid moonlight span the gaps between the lunariums and the moondisk.

\subsubsection{U5: Moondisk}
\begin{DndReadAloud}{}
A robed, hooded woman stands at the center of a circular platform, her arms stretched wide and her eyes closed. A rod piece floats between her outstretched hands.

\end{DndReadAloud}

Teremini (lawful evil, elf \textbf{archmage}) performs her ritual here. The Rod of Seven Parts of the \textit{Rod of Seven Parts} floats between Teremini's outstretched hands.

When Teremini notices the characters approaching, she shouts, "Death to all who oppose Lord Soth! And death to those pitiful wardens! You'll never stop me!"

If the characters disrupt Teremini's ritual, she screams in rage and attacks them. Teremini pockets the rod piece and fights to the death. For more about the \textit{Rod of Seven Parts}, see this book's introduction.

\subsection{Disrupting the Ritual}
Teremini's ritual is disrupted if either of the following occurs:


\begin{itemize}
\item Red moonlight shines on the white crystal, white moonlight shines on the black crystal, and black moonlight shines on the red crystal.
\item A lunar crystal is removed from its pillar. The ensuing chaos after the ritual is disrupted is an ideal time for the characters to seize the rod piece from Teremini, who is distracted as she rallies the vault's forces to fend off the dragon Orinix. If the characters fail to thwart Teremini's ritual, they should fight her elsewhere in the vault so they can retrieve the rod piece.
\end{itemize}

\begin{DndSidebar}[float=!b]{Ritual Duration}
The text is intentionally vague about how much time it takes Teremini to complete her ritual. The players should feel that time is of the essence, but not so much that their characters can't afford to retreat and take a long rest if truly necessary.



\end{DndSidebar}

\subsubsection{Moonlight Shifts}
The curtain of crimson moonlight surrounding the Three Moons Vault disappears. Areas of solid moonlight—including the moondisk, moonbridges, and barriers around the lunarium daises—instantly dissolve. Any creature standing on one of these surfaces falls to the courtyard 100 feet below.

\subsubsection{Orinix Arrives}
A portal to Lunitari opens where the moondisk once hovered. Orinix, an \textbf{adult lunar dragon} (see appendix A) flies from this portal and looms over the courtyard. Displeased at Teremini's failure, the dragon attacks everyone in sight, focusing foremost on Teremini.

\section{Next Steps}
Once the characters have the fifth piece of the \textit{Rod of Seven Parts}, they can return to Sigil via the doorway in the tree through which they arrived on Krynn.

An enormous, magical tree seems to be the location of the next rod piece, though the characters soon learn the piece is elsewhere.

\chapter{Tomb of Wayward Souls}
Retrieving the sixth piece of the \textit{Rod of Seven Parts} leads the characters to a chain of tropical islands on the world of Oerth. There, the party must plumb the depths of a deadly complex called the Tomb of Wayward Souls, which was built to lure in and slay treasure-seekers. This labyrinthine maze of deadly traps was crafted by the diabolical archlich Acererak, who takes perverse pleasure in littering his dungeons with the bones of defeated adventurers.

\section{Running This Chapter}
This chapter begins after the characters retrieve the fifth piece of the \textit{Rod of Seven Parts}. A character who holds that piece instinctively knows that the sixth piece is located on the Isle of Serpents on Oerth, a world known for its legendary dungeons and magical treasures.

The characters arrive on the island in time to witness a battle between a \textbf{kraken} and a crew of archaeologists. The archaeologists are in dire straits, and it's up to the characters to decide whether to help. Eventually, the characters realize that the sixth piece of the \textit{Rod of Seven Parts} is located in the complex the archaeologists were exploring: the Tomb of Wayward Souls.

The archaeologists haven't explored far into the deadly complex, so the characters must face its challenges without much information. The search for the rod piece culminates in a showdown with Rerak, an empowered simulacrum of Acererak.

\subsection{Character Advancement}
The characters should be 16th level when this chapter begins. The characters gain a level after they retrieve the sixth rod piece from the Tomb of Wayward Souls.

\subsection{Power of Secrets}
The characters can learn two secrets in this chapter that are applicable to the rules in "The Power of Secrets" section in this book's introduction:


\begin{description}
\item[Marian's Secret.]
Marian Xavere, one of the archaeologists, holds a lifelong fascination with Acererak and the archlich's evil magic, and she was previously tempted to study necromancy. See the "Archaeologist Camp" section later in this chapter for more information.
\item[Rerak's Secret.]
The false lich resents his imprisonment in the Tomb of Wayward Souls and never wanted to enact Acererak's will. Rerak is the guardian of the rod piece the characters seek and waits in area T26 of the complex.
\end{description}

\subsection{Sixth Rod Piece}
The Rod of Seven Parts of the \textit{Rod of Seven Parts} is in area T27 of the Tomb of Wayward Souls. For more information about the rod and the spell this rod piece allows its wielder to cast, see this book's introduction.

\section{Oerth}
This chapter takes place on Oerth, in the Greyhawk campaign setting. The world is rife with elaborate dungeons, magical treasures, and other fantastical hallmarks of sword and sorcery.

\subsection{Isle of Serpents}
Characters who research the Isle of Serpents in Sigil can learn the following:


\begin{description}
\item[Acererak's Machinations.]
Long ago, the archlich Acererak visited the Isle of Serpents in his search for precious artifacts. Acererak demolished settlements across the island, stealing treasures and killing indiscriminately. He then built one of his infamous tombs near the coast. The descendants of those who survived Acererak's onslaught remain on the island and are still working to rebuild. Many seek to reclaim the treasures that were stolen from them.
\item[Geography.]
The Isle of Serpents is a part of an archipelago known as the Asperdi-Duxchan Chain. The island's temperate climate creates a haven for venomous snakes and poisonous fauna, and its craggy cliffs contain treacherous coastal cave systems.
\item[Inhabitants.]
Popular accounts of the island describe it as unlucky, leading outsiders to avoid it. (The island's inhabitants know this rumor is false, though likely tied to Acererak's long-ago activities, and most families have been here for many generations. The characters learn this when they visit the island.)
\item[Legends.]
Legends speak of Acererak's infamous Tomb of Horrors, a deadly labyrinth filled with monsters and puzzles. Stories state that the archlich created similar tombs to guard precious magical artifacts. One such tomb is located on the isle.
\end{description}

\section{Down to Oerth}
When the party steps through the Sigil portal to the Isle of Serpents, they emerge through a ground-level opening in a cliff. The opening overlooks a lagoon, as shown on map 7.1.

\subsection{Kraken Attack!}
The characters arrive in the midst of chaos on the lagoon's beach. Read or paraphrase the following:

\begin{DndReadAloud}{}
You feel warm sea air as you gaze upon the craggy shore of an island lagoon. Rising from the lagoon's depths are tentacles that lash toward the beach at a group of panicked people dressed in hiking gear.

\end{DndReadAloud}

The tentacles belong to a \textbf{kraken} that took damage in an altercation with a hunting boat earlier and has been reduced to 250 hit points. It was lurking off the lagoon's coast when a group of archaeologists stopped here to fish. Hungry and irate, the kraken ambushed the people, who are desperately trying to escape. The kraken is positioned along the easternmost portion of the 100-foot-deep water shown on map 7.1. The archaeologists are scattered throughout the 2-foot-deep water and along the beach.

There are fifteen human and elf archaeologists on the beach. With the exception of Laysa Matulin, Talo Yieria, Vogren Starcloak, and Marian Xavere (who are described in the following sections), the archaeologists use the \textbf{commoner} stat block.

If the characters join the battle, the \textbf{kraken} focuses on them while the archaeologists retreat. If the characters don't join the battle, the kraken gobbles up eight of the archaeologists before sinking into the lagoon, satisfied. Laysa, Talo, Vogren, and Marian are among the survivors.

\subsubsection{Meeting Laysa Matulin}
When the kraken is no longer a threat, the leader of the archaeologists approaches the characters.

\begin{DndReadAloud}{}
A stout human woman with bronze skin and curly dark hair held back with a colorful scarf runs toward you. Her expression is a mixture of relief and excitement.

\end{DndReadAloud}

Laysa is a chaotic good, human, pirate-turned-archaeologist (use the \textbf{assassin} stat block). Descended from the island's original inhabitants, Laysa seeks to reclaim the treasures stolen from her ancestors. She recruited a crew to help her explore the complex where she believes the treasures are held.

Laysa invites the characters to the archaeologists' camp to rest and chat.

\subsection{Archaeologist Camp}
The archaeologists set their camp on the beach to the lagoon's south, where the cliff gives way and allows passage inland. Describe the camp as follows:

\begin{DndReadAloud}{}
The tents here sit between the shore and the jungle. Armor and digging equipment lie scattered across the sand. To one side, two men—an elf and an orc—tend to wounded colleagues. They each have a symbol of an arc of seven stars hanging from their necks, and their hands glow with divine magic. On the opposite side of camp, a human woman pores over a spellbook.

\end{DndReadAloud}

The archaeologists are investigating the ancient complex in which a piece of the \textit{Rod of Seven Parts} waits. The characters can chat with the archaeologists to learn about the complex and ask to borrow magic items to aid their quest.

\subsubsection{Laysa's Goals}
At camp, Laysa is happy to share what she knows. She volunteers the following information:


\begin{description}
\item[Foreign Complex.]
The complex's architectural style is different than local traditions, so the complex must have been built by an outsider. Based on the complex's defenses, Laysa believes something valuable is held within.
\end{description}


\begin{description}
\item[Some Progress.]
Laysa's crew uncovered the complex's entrance, as well as two false entrances laden with traps. The archaeologists haven't been able to explore much more, but they think they've identified how to delve farther into the complex.
\item[Stolen Treasures.]
Long ago, an evil mage with a skeletal appearance demolished Laysa's ancestral village on this island. The mage absconded with several treasures and was never seen again.
\end{description}

Laysa is unfamiliar with the \textit{Rod of Seven Parts}, but she agrees that such a powerful magic item could be found within the complex. She isn't aware of any other location where the rod piece might be.

Laysa provides the characters with a rough map of the complex (map 7.2) and can lead the party to its entrance when they're ready.

\subparagraph{Stolen Treasures}
Laysa hopes to recover six ancestral treasures stolen from her people. Some of these treasures are magic items, although Laysa describes them only as detailed below. The items (and their locations on map 7.3) are as follows:


\begin{itemize}
\item Ebony wand decorated with bones and feathers (area T9)
\item Ring of golden stars (area T13)
\item Sacred scimitar dedicated to serpent spirits (area T17)
\item Set of copper tablets engraved with incantations (area T19)
\item Crystal orb used in rituals (area T23)
\item Blue silk sash used in ceremonies (area T27)
\end{itemize}

If the characters uncover any of these items from the complex and return them to Laysa, she is ecstatic. For each item returned, Laysa pays the characters five gemstones worth 500 gp apiece. She keeps the gemstones in a pouch that hangs from her belt.

\subsubsection{Talo Yieria and Vogren Starcloak, Priests and Healers}
The elf man is Talo Yieria, and the orc man is Vogren Starcloak. Both are neutral good priests of Celestian, the enigmatic deity of stars and wanderers on Oerth. Though apprehensive about exploring the complex, the couple are dear friends to Laysa and support her efforts.

Talo and Vogren are well versed in the religious history of the island. If asked, they provide the following advice:


\begin{description}
\item[Reciprocity.]
Powerful nature spirits reside on the island and don't take kindly to those who would exploit it. No one should take anything without leaving something in return.
\item[Respect.]
Offerings of money, food, or crafted items show respect to the island's spirits. Acts of kindness can also gain the spirits' favor.
\end{description}

Should the characters return to the camp partway through exploring the complex, Talo and Vogren offer to use magic to heal them.

\subsubsection{Marian Xavere, Resident Mage}
The human woman poring over a spellbook is the camp's resident arcana expert, Marian Xavere (neutral \textbf{mage}). Hailing from the mainland, Marian joined Laysa's crew because of her intense fascination with magical traps and artifacts.

Marian's arcane research led her to study Acererak, and she can tell the characters about the archlich. The characters learn the information about Acererak presented in this chapter and in Acererak's entry in appendix B. Marian doesn't realize exactly how dangerous the complex here is.

\subparagraph{Magic Items}
Marian has a \textit{Lantern of Revealing} sitting near her and a \textit{Gem of Seeing} strapped to her belt. If asked, she lends the characters the lantern but is hesitant to lend the gem. A character can convince her to lend them the \textit{Gem of Seeing} by succeeding on a DC 18 Charisma (Persuasion) check.

\subparagraph{Marian's Secret}
When the characters speak at length with Marian, they notice the mage speaks enthusiastically, but somewhat guiltily, about Acererak. If the characters remark on or ask about this, Marian admits that she once admired the archlich and even considered studying necromancy. When Marian realized the great evil involved, though, she eschewed this grim path.

Regardless of the characters' reaction to this revelation, learning it counts as a secret for the purposes of the Power of Secrets rules in this book's introduction.

\section{Into the Tomb}
After devastating the island's native villages, Acererak constructed one of his tombs here. Though outsiders rarely visited the Isle of Serpents, whispers of an island complex laden with treasures spread across Oerth. Over time, this complex became known as the Tomb of Wayward Souls, since explorers would venture inside only to die.

Laysa and her crew made some progress exploring the complex, demarking the false entrances and traps. The complex's entrance is about 1 mile from the archaeologists' camp.

\subsection{Features of the Tomb}
Unless otherwise stated, the areas of the complex have the following features.

\subsubsection{Ceilings, Doors, Floors, and Walls}
The ceilings, floors, and walls of the complex are constructed from limestone. Chamber ceilings are 20 feet high, while hallway ceilings are 10 feet high. Doors are wooden and unlocked unless otherwise specified; some must be opened by special means. The complex's walls are immune to any spell or magical effect that would change their shape, such as a \textit{Passwall} or \textit{Stone Shape} spell.

\subsubsection{Hidden Doors}
Illusion magic conceals many of the complex's doors, making them look like solid walls. These hidden doors are detectable by touch, and creatures with truesight see the doors as if they weren't hidden. Casting \textit{Dispel Magic} on a hidden door permanently removes the magic concealing it. A hidden door is otherwise identical to other doors in the complex. Locations of hidden doors are marked on map 7.3.

\subsubsection{Lighting}
The interior of the complex is brightly lit by crystal sconces holding eerie green flames. The flames can't be extinguished by any means.

\subsubsection{Planar Trammels}
When a creature within the complex casts a spell or uses an effect that would transport itself or another creature to a different plane (such as casting \textit{Astral Projection}, \textit{Banishment}, \textit{Etherealness}, or \textit{Plane Shift}), the spell slot and components, item charges, or other resources are expended as normal, but nothing happens.

\subsubsection{The Rod Piece and Its Guardian}
The sixth piece of the rod can be found in area T27. The piece is guarded by Rerak, an empowered simulacrum of Acererak designed to harvest souls for the archlich. When an adventurer dies in the complex, Rerak traps their soul using foul magic and shunts it to the gems embedded in the false lich's eyes. The souls remain in the false lich's eyes for 24 hours, after which the souls are transferred to Acererak and are trapped forever. Rerak also resets traps that adventurers have triggered in the complex.

A long time has passed since the archlich created Rerak. The simulacrum is torn between dutifully fulfilling Acererak's will and resenting the archlich for imprisoning him within this complex. Retrieving the rod piece requires confronting Rerak in his crypt, though the characters can recover the piece without fighting the simulacrum. For more information about confronting Rerak, see area T26.

\subparagraph{Corrupted Mirages}
Rerak's resentment, intermingled with the magic of the rod piece, causes strange mirages to appear throughout the complex. These mirages are described where they're encountered (see areas T9, T13, T21, and T24).

\subsection{Tomb Locations}
The areas of the tomb are keyed to map 7.3.

\subsubsection{T1: False Entrance}
\begin{DndReadAloud}{}
This wide corridor ends at two stone doors with metal handles. At the entrance, the archaeologists posted symbols for "danger" and "fire."

\end{DndReadAloud}

The doors at the end of the corridor are fake and can't be opened. Behind them is a solid wall.

\subparagraph{Explosion Trap}
A creature that steps inside the corridor triggers a trap. When the creature touches the corridor floor, the floor tiles depress, and a slab of stone lodged in the ceiling lowers to the ground, sealing the corridor before flooding it with magical fire. Each creature inside the closed-off corridor must make a DC 20 Dexterity saving throw, taking 28 (8d6) fire damage on a failed save or half as much damage on a successful one.

The stone slab seals off the corridor until the next dawn, when the slab rises into its ceiling cavity and the trap resets. In the meantime, a creature can use an action to try to lift the slab, creating a passable opening with a successful DC 18 Strength (Athletics) check.

\subsubsection{T2: False Entrance}
\begin{DndReadAloud}{}
Strange pockmarks line the walls of this corridor, at the end of which is a set of dark wooden doors.

\end{DndReadAloud}

\subparagraph{Arrow Trap}
Spring-loaded arrows are hidden in the pockmarked walls. A creature that enters the trapped area marked on map 7.3 for the first time on a turn or ends its turn there triggers a fusillade of arrows and must make a DC 20 Dexterity saving throw, taking 27 (6d8) piercing damage from the arrows on a failed save or half as much damage on a successful one. The 5-foot-wide area along the east wall is outside the trap's area and safe to stand in.

\subparagraph{Doors' Secret}
The doors at the end of the hallway are fake and can't be opened. Tiny text is scrawled on the doors' front:

\begin{DndReadAloud}{}
"Acererak commends you! Alas, these doors lead nowhere, but your skill shall not go unrewarded. Should you wish to plumb the depths of my tomb, heed these words: the blue one craves your magic."

\end{DndReadAloud}

This is a clue for opening the door in area T4.

\subsubsection{T3: Tomb Entrance}
\begin{DndReadAloud}{}
This hallway has been partially excavated. Lifted stone lids reveal a ten-foot-deep pit filled with spikes in the center of the corridor. A thin plank of wood is laid over the pit to allow for safe crossing.

\end{DndReadAloud}

Characters can use the wooden plank to cross the pit safely. A creature that falls into the pit takes 5 (1d10) piercing damage from the spikes, which turn the floor of the pit into \textit{difficult terrain}.

At the bottom of the pit, situated in the middle of its northeast wall, is a hidden door that leads to area T5.

\subsubsection{T4: Face of the Great Blue Devil}
\begin{DndReadAloud}{}
A relief of a scowling, fiendish face made of azure mosaic tiles spans the hallway's back wall. The face's mouth is open, revealing a black maw, and its eyes are carved from clouded white crystal. To the right of the mosaic is a closed door.

\end{DndReadAloud}

A character who examines the relief with a \textit{Detect Magic} spell or similar magic finds that the relief's mouth radiates conjuration magic and its eyes radiate abjuration magic.

\subparagraph{Mouth Portal}
The relief's mouth is 10 feet wide and several feet tall. Objects and creatures that enter the relief's mouth are teleported to area T14.

\subparagraph{Opening the Devil's Door}
The relief is the key to opening the door to area T6. The door can be opened in no other way.

When a creature standing within 5 feet of the relief casts a spell using a spell slot, the spell slot is expended as normal, but the spell has no effect. Instead, the magic is absorbed by the relief, filling it with a number of charges equal to the level of spell slot expended and causing the relief's eyes to glow green. Once the relief has absorbed 6 or more charges, the door opens. The door closes on its own after 10 minutes unless it is held or wedged open. When the door closes, the relief's eyes stop glowing, and opening the door again from this area requires the expenditure of more spell slots. The door can be opened from the east side like any normal door.

\subsubsection{T5: Cave of Seven Casks}
\begin{DndReadAloud}{}
Crooked wooden planks line the floor of this cave. Seven casks are embedded along one wall. Each cask has a different symbol painted on it, with a spigot beneath it. Cups and tankards line shelves on the opposite wall.

\end{DndReadAloud}

A character who studies the shelves notices a repeating message carved in Celestial along the shelves' edges:

\begin{DndReadAloud}{}
"A septet of libations I present to you,

But only one will help you escape this tomb.

Six drinks are magic; one is mundane.

Those marked with stars are vitality's bane.

Blue is neither blessing nor curse.

Moon cleanses you of illness or worse.

Green's neighbor is never boring—

Drink deep; pass through the flooring."

\end{DndReadAloud}

The message gives clues as to the casks' contents.

Each cask contains 20 pints of liquor. In order from left to right, the casks are marked with the following symbols:


\begin{description}
\item[Green Star.]
The liquor in this cask is infused with necromantic magic. A creature that drinks any amount of the liquor must make a DC 20 Constitution saving throw, taking 22 (5d8) necrotic damage on a failed save or half as much damage on a successful one.
\item[Blue Square.]
A creature that drinks a pint or more of the liquor in this cask gains the benefit of a \textit{Potion of Diminution}. A character reduced to Tiny size by the liquor can't fit through the narrow gaps between the floorboards (see the "Tunnel to Area T8" section for more information).
\item[Green Crescent Moon.]
A creature that drinks a pint or more of the liquor in this cask gains the benefit of an \textit{Elixir of Health}.
\item[Red Square.]
A creature that drinks a pint or more of the liquor in this cask gains the benefit of a \textit{Potion of Gaseous Form}. A character in gaseous form can easily pass through the narrow gaps between the floorboards (see "Tunnel to Area T8" below).
\item[Blue Circle.]
A creature that drinks a pint or more of the liquor in this cask gains the benefit of a \textit{Potion of Growth}.
\item[Green Circle.]
The liquor is this cask is not magical, but it is toxic. A creature that drinks any amount of the liquor must succeed on a DC 25 Constitution saving throw or have the poisoned condition for 1 hour.
\item[Red Star.]
The liquor in this cask acts as a magical paralytic. Upon drinking any amount from this cask, a creature must succeed on a DC 20 Constitution saving throw or have the paralyzed condition for 1 hour.
\end{description}

Liquor removed from this cave immediately loses its magical properties. Removing a cask from this room causes the liquor within the cask to evaporate.

\subparagraph{Tunnel to Area T8}
Characters who examine the floor can see an empty basement through the 1-inch cracks between the floorboards. The basement has a tunnel exiting to the east, toward area T8. The easiest way to reach this tunnel is to drink from the cask marked with the red square to adopt a cloudlike form. The floorboards are otherwise immovable and immune to all damage.

If the characters have trouble solving the puzzle, any character who walks across the floorboards realizes that the only way past this room is through the thin gaps between the floorboards. If the characters still struggle, allow them to make a DC 12 Intelligence (Arcana) check to determine that the magical liquor in one or more of the casks might facilitate travel through the cracks between the floorboards.

\subsubsection{T6: Hall of Gemstones}
\begin{DndReadAloud}{}
Rows of colorful jewels line this corridor's walls. At the east end stands an emerald-colored statue of a crashing wave, the words "This wave is green" carved into its base. A door is set in the corridor's south wall.

\end{DndReadAloud}

Close inspection of the statue uncovers a divot in the statue's base big enough to hold one jewel.

\subparagraph{Gemstone Puzzle}
All the jewels embedded in the walls are made of worthless crystal except for one: an emerald worth 1,000 gp. A character who spends at least 1 minute examining the jewels spots the emerald with a successful DC 20 Intelligence (Investigation) check. Prying the emerald from its wall socket is easily done. Placing the emerald in the statue opens the door sealing off the passage to area T16. The door resists all other attempts to open it. Once opened, however, it ceases to be locked and can thereafter be opened and closed without the emerald.

A character who succeeds on a DC 16 Intelligence (Arcana) check recalls that emeralds are used to contain water elementals.

Whenever a creature touches one of the other jewels embedded in the walls, the statue glows and launches a blast of freezing magic down the hallway in a 30-foot line that is 5 feet wide. Each creature in that area must make a DC 20 Dexterity saving throw, taking 33 (6d10) cold damage on a failed save or half as much damage on a successful one.

\subsubsection{T7: Skeleton Closet}
\begin{DndReadAloud}{}
A moldering skeleton draped in cobwebs sits in one corner of this small room. A gold choker set with a large black stone hangs around the skeleton's neck.

\end{DndReadAloud}

If examined with a \textit{Detect Magic} spell or similar effect, the necklace radiates an aura of necromancy magic.

\subparagraph{Necklace Trap}
The first time a creature touches the necklace, it releases a burst of deathly energy in a 5-foot-radius sphere. Each creature in that area must make a DC 22 Constitution saving throw, taking 36 (8d8) necrotic damage on a failed save or half as much damage on a successful one.

Once the trap has been triggered, the necklace becomes a mundane choker worth 1,200 gp.

\subsubsection{T8: Room of Myriad Archways}
\begin{DndReadAloud}{}
Each wall of this octagonal room contains a stone archway. A veil of white fog fills the room.

\end{DndReadAloud}

The fog is magical and can't be cleared or dispelled.

\subparagraph{Archways}
The archways lead to the following destinations:


\begin{description}
\item[North Archway.]
Any creature or object that enters this one-way portal is teleported to area T3, appearing at the southwest end of the corridor or the nearest unoccupied space.
\item[Northeast Archway.]
This archway connects to area T11. The tunnel leading there is made of smooth, gray stone.
\item[East Archway.]
Any creature or object that enters this one-way portal is teleported to an unoccupied space in area T20a. If there's not enough space in that location for the creature or object, it is unable to pass through the portal and is rebuffed instead.
\item[Southeast Archway.]
This archway connects to area T7 via a hidden doorway. The tunnel beyond the archway is made of obsidian and covered in cobwebs.
\item[South Archway.]
Any creature or object that enters this one-way portal is teleported to a random unoccupied space in area T14.
\item[Southwest Archway.]
This archway connects to a basement under the floorboards of area T5 via a dirt tunnel.
\item[West Archway.]
Any creature or object that enters this one-way portal is teleported to an unoccupied space in area T20b. If there's not enough space in that location for the creature or object, it is unable to pass through the portal and is rebuffed instead.
\item[Northwest Archway.]
This archway connects to area T9. The tunnel leading there is covered with moss.
\end{description}

\subsubsection{T9: Forest of Spirits Mirage}
\begin{DndReadAloud}{}
You arrive at an expansive glade of trees. Bobbing lights flit between the branches, and at the center of the glade stands a towering banyan tree wrapped in sickly looking vines. A path winds through the glade but disappears when you try to focus closely on it.

\end{DndReadAloud}

The banyan tree is a neutral \textbf{treant} named Abalahin. Abalahin is real, while the rest of the forest is a mirage. The mirage is tactile, so creatures can interact with it.

\subparagraph{Piercing the Illusion}
The mirage is linked to Abalahin's existence. If Abalahin is killed, the mirage ends, revealing the room to be a stone chamber containing Abalahin's remains.

\subparagraph{Befriending Abalahin}
Abalahin is covered with strangling vines that render the treant sluggish and distant. A \textit{Remove Curse} spell or similar effect causes the blighted vines around Abalahin to dissolve. Once the vines are gone, Abalahin is grateful to the characters and, in return, suppresses the mirage, revealing the room's exits.

\subparagraph{Bobbing Spirit Lights}
There are 1d8 bobbing spirit lights (each uses the \textbf{will-o'-wisp} stat block and speaks Common and Sylvan) that fill the illusory glade at any given time. The spirit lights are part of the mirage, are indifferent to the characters' presence, and disappear if attacked. Any character with a passive Wisdom (Perception) score of 17 or higher discerns the spirit lights whispering and overhears one of the following remarks:


\begin{itemize}
\item "The next challenge is crushing—move fast!"
\item "The master of this tomb longs to be free. Wouldn't you, if you were he?"
\item "Steal not from the library, for it holds trouble."
\item "Illusions are we, born from a mix of wanderlust and corrupted magic. Many more can you find in this forgotten place."
\item "The jade serpents crave offerings."
\end{itemize}

\subparagraph{Treasure}
Close inspection of Abalahin's roots reveals an ebony wand decorated with bones and feathers (a \textit{+3 Wand of the War Mage}). A character who succeeds on a DC 15 Charisma (Persuasion) check can convince Abalahin to relinquish the wand voluntarily. The wand can also be removed without attracting Abalahin's notice by a character who succeeds on a DC 25 Dexterity (Sleight of Hand) check. The wand is one of the treasures Laysa mentioned earlier in this chapter.

Abalahin is hostile toward creatures it catches trying to steal the wand. Once the treant has been defeated, the wand can easily be retrieved.

\subsubsection{T10: Crushing Room Trap}
This room appears to be empty. When one or more characters enter the room, read the following:

\begin{DndReadAloud}{}
A click sounds from above. The room's ceiling begins descending rapidly.

\end{DndReadAloud}

The 20-foot-high ceiling drops suddenly, stopping 6 inches above the floor before rising just as quickly and plummeting again. This rapid, piston-like movement lasts for 1 minute, after which the ceiling returns to its original height and locks in place until the trap is triggered again.

Each time the trap triggers, have the characters roll initiative. The ceiling drops to its lowest point on initiative count 20 of each round, losing initiative ties, until the trap resets.

Any creature in the room when the ceiling drops must make a DC 20 Dexterity saving throw. On a failed save, the creature takes 55 (10d10) bludgeoning damage and has the prone condition. On a successful save, the creature can rush from the room, taking no damage, or remain in the room and take half as much damage only. A creature that can fit in spaces as narrow as 6 inches has advantage on the save thanks to the narrow gap that remains between the ceiling and floor when the ceiling is at its lowest level.

A Small or larger creature subjected to the trap can choose to make a DC 20 Strength saving throw instead of a Dexterity saving throw, provided the creature doesn't have the prone condition, with the aim of slowing the ceiling's descent for the betterment of others. The creature making the Strength saving throw takes 55 (10d10) bludgeoning damage regardless of whether its save succeeds or fails. If the save succeeds, however, the trap's damage to other creatures is reduced to 0 until initiative count 20 of the next round (losing initiative ties).

The ceiling is too heavy to be stopped by any braces, but an \textit{Immovable Rod} or similar item can prevent it from dropping so far as to crush anyone.

\subsubsection{T11: Strange Altar}
\begin{DndReadAloud}{}
A dark, triangular altar stands in the center of this chamber. Carved into each side of the altar is a niche containing a chest. One chest is plated with gold, one is plated with silver, and the last is plated with lead.

\end{DndReadAloud}

The chests are unlocked, and their contents are as follows:


\begin{description}
\item[Gold Chest.]
This \textit{chest} is worth 100 gp on its own and contains 100 gp.
\item[Silver Chest.]
This \textit{chest} is worth 50 gp on its own and contains a ceremonial \textit{dagger} plated with electrum with an amethyst set in its pommel (worth 750 gp).
\item[Lead Chest.]
This worthless \textit{chest} contains a \textit{Spell Scroll} of \textit{Knock}, which the altar trap values at 200 gp (see below).
\end{description}

\subparagraph{Altar Trap}
Removing a chest or its contents from a niche causes the altar's carvings to glow with pale light. Creatures in the room have 1 minute to either return the chest and its contents to the niche or place one or more items of equal or greater value into the niche. Otherwise, one of the following effects occurs when the time expires, depending on the chest involved:


\begin{description}
\item[Gold Chest.]
Six swarms of poisonous snakes teleport into the room. The snakes are hostile.
\item[Silver Chest.]
A burst of desiccating energy radiates from the altar in a 15-foot-radius sphere. Each creature in the area must make a DC 17 Constitution saving throw, taking 18 (4d8) necrotic damage on a failed save or half as much damage on a successful one.
\item[Lead Chest.]
A \textbf{death knight} teleports into the room, appearing in a random unoccupied space within 15 feet of the altar. The knight is hostile.
\end{description}

\subsubsection{T12: Room of Respite}
\begin{DndReadAloud}{}
A marble gazebo with a fountain at its center is surrounded by flowering bushes.

\end{DndReadAloud}

This room is safe.

\subparagraph{Fountain}
The fountain's waters have healing properties. A creature that drinks from the fountain regains 21 (6d6) hit points. Once a creature has benefited from the fountain's magic, it can't do so again until the next dawn.

\subsubsection{T13: Celestial Sky Mirage}
A hidden door conceals this room. Creatures that step through the doorway are seemingly transported to a location outdoors, thanks to illusion magic:

\begin{DndReadAloud}{}
You're transported to a grassy field with a starry night sky above you.

\end{DndReadAloud}

The grassy field and night sky are illusions concealing a bare, rectangular room. Characters can move only as far as the room's walls and ceiling allow, for though these surfaces are hidden to normal eyes, they can be felt. Creatures with truesight recognize the illusions as mirages and can see the room as it actually is; they can also see an engraved black box resting on a black dais in the middle of the room. Characters who can't see the box and dais can locate them by touch after a careful search of the area.

\subparagraph{Observing the Night Sky}
A character observing the mirage of the night sky can make a DC 14 Intelligence (Nature) check. On a success, a character realizes that, while beautiful, the night sky is wholly inaccurate, as if created by someone who has never seen the sky before. Succeeding on this check also allows the character to see through the room's illusions, revealing not only the walls and ceiling but also the box and the dais on which it sits.

\subparagraph{Box Trap}
The engraved box can't be forced open, but a character can try to unlock it using thieves' tools, doing so with a successful a DC 20 Dexterity (Sleight of Hand) check. A \textit{Knock} spell or similar magic also unlocks the box.

A character who inspects the box before opening it can make a DC 15 Intelligence (Investigation) check. On a successful check, the character notices a spring-loaded compartment beneath the box's lid. A character can use thieves' tools to try to dismantle this compartment and disable the trap, doing so with a successful DC 15 Dexterity (Sleight of Hand) check. If the trap isn't deactivated, poisoned darts burst from the box when it's opened. Each creature within 10 feet of the box when the trap is sprung must succeed on a DC 20 Dexterity saving throw or take 10 (4d4) piercing damage and have the poisoned condition for 1 hour.

\subparagraph{Treasure}
Contained within the box is a delicate gold ring engraved with geometric star designs. This ring isn't magical or of any significant monetary value, but Laysa's ancestors once used it in traditional storytelling ceremonies involving astrology. This is one of the treasures Laysa mentioned earlier in this chapter.

\subsubsection{T14: Juggernaut Arena}
\begin{DndReadAloud}{}
Thin, obsidian pillars are spaced throughout this musty room, the ceiling of which is only ten feet high. In the flickering light, you see two massive, rolling constructs roving the space.

\end{DndReadAloud}

The room's ceiling is 10 feet high—barely large enough to accommodate the two granite juggernauts (see appendix A) that trundle between the room's pillars. The juggernauts attack any intruders they notice, but they won't strike pillars.

The juggernauts fight until they are destroyed or until a creature destroys the triangular stone in the control room (area T15).

\subsubsection{T15: Juggernaut Control Room}
A hidden door conceals this room.

\begin{DndReadAloud}{}
A five-foot-tall, triangular stone covered in runes stands in this bare room's center. At the top of the stone, a plum-sized diamond pulses with red light.

\end{DndReadAloud}

The triangular stone is immovable and powers the juggernauts in area T14. A character who examines the triangular stone and succeeds on a DC 17 Intelligence (Arcana) check can decipher its runes and discern its purpose.

\subparagraph{Destroying the Stone}
The triangular stone has AC 20; 100 hit points; immunity to poison and psychic damage; and immunity to bludgeoning, piercing, and slashing damage from nonmagical attacks. When the stone is destroyed, the juggernauts in area T14 stop moving and become inanimate objects.

\subparagraph{Treasure}
Once the stone has been destroyed or both juggernauts have been defeated, characters can take the diamond, which is worth 5,000 gp. The diamond stops glowing when it leaves the room.

\subsubsection{T16: Acid Pit Crossing}
\begin{DndReadAloud}{}
Filling most of this chamber is a ten-foot-deep pit of green acid, over which stretches a three-foot-wide stone bridge. Nozzles embedded in the north and south walls shoot blasts of air across the bridge.

\end{DndReadAloud}

A creature that falls into the acid takes 21 (6d6) acid damage. Any creature that starts its turn in the acid also takes this damage.

\subparagraph{Gust Nozzles}
Air blasts from nozzles in the walls north and south of the bridge with such force that they can't be easily stoppered. When a creature moves onto the bridge for the first time on a turn or starts its turn there, or when a flying creature enters a space above the bridge or the acid pit for the first time on a turn, the creature must succeed on a DC 20 Dexterity saving throw. On a failed save, the creature is pushed by the gusts 5 feet north or south (determined randomly) and has the prone condition; if the space isn't on a surface that can support the creature, the creature falls as normal.

\subsubsection{T17: Jade Serpent Guardians' Chamber}
\begin{DndReadAloud}{}
A jade statue stands between this room's two exits. The statue depicts an armored, feminine figure from the waist up, with six arms and a gleaming curved sword in each hand; from the waist down, the statue's body is serpentine. Two statues depicting large cobras with jade scales sit in alcoves to the east.

\end{DndReadAloud}

If one or more characters enter the room, read the following:

\begin{DndReadAloud}{}
The center statue's eyes glow with white light, and the cobra statues begin writhing toward you.

\end{DndReadAloud}

The central statue uses the \textbf{marilith} stat block but is a Construct. This central statue controls the two cobra statues; each cobra statue uses the \textbf{spirit naga} stat block but is a Construct. The statues speak Common. The statues are hostile toward any creature that enters the chamber.

\subparagraph{Appeasing the Statues}
A character can use an action to try to appease the statues with a DC 20 Charisma (Persuasion) check; if the character makes an offering of food or treasure worth 50 gp or more, this check is made with advantage. On a successful check, the statues become friendly toward the characters and cease attacking.

The statues served as island guardians for centuries before Acererak abducted them. If informed about Laysa, the statues are thrilled to learn about her efforts.

Acererak's magic confines the statues to this room, but previous adventurers have told them about surrounding areas. The statues can share details about what's in areas T16 and T18.

\subparagraph{Treasure}
One of the central statue's swords functions as a Scimitar of Sharpness in the hands of a Humanoid. The scimitar is one of the treasures Laysa mentioned earlier in this chapter. If told about Laysa, the central statue willingly gives the sword; otherwise, the sword can be retrieved when all the statues have been defeated.

\subsubsection{T18: Hall of Discordance}
\begin{DndReadAloud}{}
Eight niches are built into these walls. Each niche contains a musical instrument: a lute, a lyre, a viol, a flute, a drum, a dulcimer, a shawm, and a horn. At the end of the hallway is a double door emblazoned with two theater masks—one elated and one despairing.

\end{DndReadAloud}

When a creature comes within 5 feet of any of the instruments, they all glow with a purple light and begin playing on their own, filling the hall with an awful din. Each creature in the hall must make a DC 20 Wisdom saving throw, taking 16 (3d10) psychic damage on a failed save or half as much damage on a successful one. After the damage is dealt, the instruments cease to play or glow, becoming normal items of their kind.

\subparagraph{Mask Door}
Upon closer inspection, the mouths of both masks are open slots, each large enough to fit a Medium or Small Humanoid's arm. The doors are locked.

Thin script engraved above the pair of masks reads as follows:

\begin{DndReadAloud}{}
"Pay the toll, in blood or in song."

\end{DndReadAloud}

There are two ways to open the doors. A character can play one of the hallway's instruments and make a DC 20 Charisma (Performance) check; if the character doesn't have proficiency with the chosen instrument, this check is made with disadvantage. On a successful check, the doors open. On a failed check, the character takes 16 (3d10) psychic damage, and the doors remain locked.

Alternatively, a character can place one of their arms into a mask's open mouth. When each mouth has an arm in it, the mouths clamp down. If both arms belong to the same character, that character takes 28 (8d6) piercing damage; if the arms belong to different characters, each character takes 14 (4d6) piercing damage. The mouths then unclamp, and the doors unlock.

\subsubsection{T19: Library}
\begin{DndReadAloud}{}
Shelves filled with tomes and scrolls line the walls of this expansive chamber. Five stone statues depicting skeletal warriors stand sentinel throughout the room.

\end{DndReadAloud}

\subparagraph{Guardians}
Whenever an item is removed from this room's shelves, one of the five skeletal statues here springs to life. The statue uses the \textbf{stone golem} stat block but is Medium and hostile, attacking whichever creature took the item. The statue remains active until the triggering item is returned or until the statue is destroyed. Multiple statues can be active at once.

\subparagraph{Treasure}
A thorough search of the shelves' contents takes 10 minutes and yields the following items:


\begin{itemize}
\item \textit{Spell Scroll} of \textit{Greater Restoration}
\item Set of engraved copper sheets that functions as a \textit{Manual of Bodily Health} (one of the treasures Laysa mentioned earlier in this chapter)
\item Thin, leather-bound journal written in Common that explains that the combination for the crypt doors (area T25) is the word "die" in Draconic
\end{itemize}

\subsubsection{T20a–T20b: Void Closets}
These two storage closets are identical. The door to each room is locked. As an action, a character can use thieves' tools to try to unlock either door, doing so with a successful DC 15 Dexterity (Sleight of Hand) check. Each door can also be forced open from either side with a successful DC 11 Strength (Athletics) check, or with a \textit{Knock} spell or similar magic.

\begin{DndReadAloud}{}
This room is pitch dark.

\end{DndReadAloud}

The interior of each closet is filled with magical darkness that can't be dispelled.

\subparagraph{Spheres of Annihilation}
In the center of each room is an uncontrolled \textit{Sphere of Annihilation}. The magical darkness obscures the sphere; only creatures with truesight can detect its presence.

\subsubsection{T21: Phantasmal Mirage}
\begin{DndReadAloud}{}
Eerie sounds echo in this pitch-dark room.

\end{DndReadAloud}

This room is filled with magical darkness, and neither the darkness nor the eerie sounds can be dispelled. Any creature that enters this room for the first time must succeed on a DC 18 Wisdom saving throw or have the frightened condition, as the darkness and eerie sounds conjure up the creature's deepest fears. The frightened condition ends on the creature when it leaves the room.

While frightened in this way, a creature takes 11 (2d10) psychic damage for every 5 feet it moves. After taking this damage, the creature can repeat the saving throw, ending the frightened condition on a itself on a success.

\subparagraph{Piercing the Illusion}
Creatures with truesight see the room for what it is: an abandoned study. A dusty desk and empty bookcase line the east wall.

\subsubsection{T22: Vortex Bridge}
\begin{DndReadAloud}{}
A dull roar fills this hexagonal chamber. A bridge consisting of engraved, colored limestone tiles stretches between two ledges. A pitch-black vortex of energy churns beneath the bridge.

\end{DndReadAloud}

The tiled bridge stretches over a gravity vortex 60 feet below. The force of the vortex reduces the flying speed of any creature within the room to 0 feet.

\subparagraph{Bridge Puzzle}
Close inspection of the bridge's tiles reveals that each one bears an engraved symbol representing one of the eight schools of magic. A \textit{See Invisibility} spell or similar effect uncovers concealed writing on the ledge before the bridge, which reads as follows:

\begin{DndReadAloud}{}
"The vortex eats all magic, save for that from which it was born."

\end{DndReadAloud}

To cross the bridge safely, creatures must step on the tiles bearing the symbol for evocation magic, which are red, or the tiles bearing the symbol for transmutation magic, which are orange. The tiles are shown on map 7.4.

When observed within the range of a \textit{Detect Magic} spell or similar magic, the vortex radiates an aura of evocation and transmutation magic. A character who observes the vortex or witnesses the flight-inhibiting effect of the room's gravity well can make a DC 16 Intelligence (Arcana) check, recalling the following facts on a success:


\begin{itemize}
\item The spell \textit{Reverse Gravity} is a transmutation spell.
\item The school of transmutation magic deals with manipulating energy and matter.
\item Evocation spells create powerful elemental effects.
\end{itemize}

If the players have trouble with this puzzle, reveal that the characters must step on only evocation or transmutation tiles, and tell them which colors these tiles are. When a creature steps on an incorrect tile, the tile breaks away, and both the tile and the creature are sucked into the vortex below. The creature takes 36 (8d8) force damage and is teleported to a safe, unoccupied space at the south end of the room or in the tunnel leading to area T19.

\subsubsection{T23: Arcane Laboratory}
\begin{DndReadAloud}{}
Three tall, glass cylinders line the back wall of this dimly lit chamber. Each cylinder contains a decaying corpse, half-buried in snow topped with red powder. Against the north wall stands a bare wooden desk.

\end{DndReadAloud}

This chamber bears the remnants of Acererak's endeavors to create empowered simulacra to guard his many tombs. A character who examines the cylinders and succeeds on a DC 20 Intelligence (Arcana) check recognizes the snow and powdered ruby in each glass cylinder as components for creating a simulacrum.

\subparagraph{Secret Desk Compartment}
A character who searches the desk and succeeds on a DC 20 Intelligence (Investigation) check discovers that one of its drawers contains a locked false bottom. Rerak (see area T26) carries the key to this compartment, but a character with thieves' tools can use them to try to unlock the compartment, doing so with a successful DC 20 Dexterity (Sleight of Hand) check. A \textit{Knock} spell or similar magic also unlocks the compartment.

The secret compartment bears a magical trap. If the compartment is opened by any means other than using the key found on Rerak in area T26, magical lightning shoots from the drawer. Each creature in the room must make a DC 15 Dexterity saving throw, taking 26 (4d12) lightning damage on a failed save or half as much damage on a successful one. A \textit{Detect Magic} spell or similar effect unveils the source of this trap: a tiny rune inscribed in the drawer. The trap can be disarmed by dispelling the rune with a \textit{Dispel Magic} spell (DC 17).

\subparagraph{Treasure}
Within the hidden desk compartment is a rose quartz \textit{Crystal Ball} and a stack of notes written in Common. The Crystal Ball is one of the treasures Laysa mentioned earlier in this chapter.

A character who spends 10 minutes studying the notes learns about Acererak's endeavors to create empowered simulacra of himself to guard his tombs. Acererak speaks of his simulacra with disdain, viewing them as tools. He finds the fact that his simulacra can develop sentience a sadistic joke.

The notes also explain that the combination for the crypt doors (area T25) is the word "die" in Draconic.

\subsubsection{T24: Underwater Trench Mirage}
Characters standing outside this room can see that it is empty, but a \textit{Detect Magic} spell reveals the presence of illusion magic throughout. When one or more characters enter the room, read the following:

\begin{DndReadAloud}{}
The walls, floor, and ceiling of the chamber vanish, replaced with a dark ocean in which you are submerged. There is no visible surface. Beneath you swims a large, finned shadow.

\end{DndReadAloud}

This chamber contains a mirage of a deep, underwater trench. Though illusory, the seawater is tactile, and creatures unable to breathe underwater suffocate while in this chamber.

A hostile \textbf{giant shark} swims beneath the characters. The shark is part of the mirage, though it can harm characters who don't pierce the illusion. The shark can't harm anyone outside the room.

\subparagraph{Piercing the Illusion}
A creature that succeeds on a DC 18 Intelligence (Investigation) check can see through the illusion, perceiving the empty room and its exits. A creature with truesight can also see through the illusion.

Upon realizing the illusion, a creature can breathe the water as easily as it breathes air. However, the creature continues to interact with the water as if it were tactile, such as swimming through it.

\subsubsection{T25: Crypt Doors}
\begin{DndReadAloud}{}
This corridor ends at a massive double door with a carved relief depicting a river of screaming souls flowing through a ravaged landscape. Instead of handles, the doors bear a trio of concentric, rune-carved platinum circles that can be rotated. At the circles' center is an engraving of a leering skull.

\end{DndReadAloud}

The platinum circles are the locking mechanisms for the crypt doors. A character who can read Draconic sees the letters \textit{a} through \textit{z} in Draconic on each circle.

When the three concentric rings are aligned so they spell the word "die" in Draconic, the crypt doors unlock. This combination can be found in areas T19 and T23. Alternatively, a character using thieves' tools can spend 1 minute trying to unlock the doors, doing so with a successful DC 25 Dexterity (Sleight of Hand) check. A \textit{Knock} spell or similar magic also unlocks the doors.

\subparagraph{Knockout-Gas Trap}
If the double door is opened via any other means besides the correct combination, a cloud of sweet-smelling gas is released from the skull engraving. The gas fills the widened area immediately south of the double door. Each creature in that area must succeed on a DC 20 Constitution saving throw or have the unconscious condition for 2d4 + 10 minutes. An unconscious creature can repeat this saving throw each time it takes damage, ending the effect on itself on a successful save.

\subsubsection{T26: Crypt}
\begin{DndReadAloud}{}
Six stone pedestals stand near the walls of this expansive room. From left to right, they bear a blue scarf, a golden ring, a glittering scimitar, an ebony wand wrapped with feathers and bone, a set of engraved copper sheets, and a crystal ball made of rose quartz.

At the room's center looms a black throne, upon which sits a skeleton dressed in luxurious robes and headdress. A platinum door stands behind the throne.

\end{DndReadAloud}

The figure sitting on the throne is Rerak, a \textbf{false lich} (see appendix A). Although alert, Rerak remains motionless until he takes damage or a creature moves within 15 feet of him, at which point he becomes hostile and attacks.

\subparagraph{Appeasing Rerak}
Upon creation, Rerak was magically bound to Acererak's piece of the \textit{Rod of Seven Parts}; as long as the rod piece remains within the tomb, so must Rerak. Rerak resents his creator for leaving him to languish within this tomb, and though he remains loyal to Acererak, he longs to experience the world beyond the tomb's walls.

Throughout the fight, Rerak commands the characters to flee. Each time Rerak does this, each character can make a DC 19 Wisdom (Insight) check without taking an action. On a successful check, the character discerns a melancholy tone to Rerak's words, as if he loathed that he must kill them.

During combat, a character can use an action to try to convince Rerak to turn against his creator. Have that character make a DC 23 Charisma (Persuasion) check. If the character mentions the mirages or Acererak's disdain toward his creations (found in the notes in area T23), this check is made with advantage. On a successful check, Rerak's loyalty to Acererak is shaken. After three successful checks, Rerak halts the battle and crushes his eye sockets' gems, severing his connection to Acererak and releasing any souls trapped within.

\subparagraph{Rerak's Secret}
If the characters convince Rerak to turn against Acererak and talk with the false lich about his imprisonment in the tomb or his mistreatment at the archlich's hands, Rerak reveals how unhappy and lonely he has been for centuries. Rerak also reveals that as far back as his earliest memory, the false lich never truly wanted to fulfill Acererak's will.

Regardless of the characters' reaction to this revelation, learning it counts as a secret for the purposes of the Power of Secrets rules in this book's introduction.

\subparagraph{False Treasures}
The items on the pedestals are false re-creations of treasures that can be found throughout the complex. A creature that touches one of these false treasures must immediately succeed on a DC 15 Constitution saving throw or take 10 (3d6) necrotic damage and have the stunned condition until the start of its next turn.

\subparagraph{Keys}
Rerak carries two keys: a platinum key shaped like a skeleton that unlocks the door to the crypt's vault (area T27) and a plain iron key that unlocks the secret compartment in area T23. Characters who defeat Rerak find the keys amid his dusty remains.

If the characters convinced Rerak to turn against his creator, Rerak provides the keys willingly.

\subsubsection{T27: Crypt Vault}
The door to the crypt vault is locked and protected by magical wards. It can be unlocked only with the platinum skeleton key in Rerak's possession. A creature that attempts to open the vault by any other means must make a DC 25 Dexterity saving throw, taking 28 (8d6) force damage on a failed save or half as much damage on a successful one.

The Rod of Seven Parts of the \textit{Rod of Seven Parts} rests on the pedestal.

\begin{DndReadAloud}{}
Coins and precious stones are strewn across this vault's floor. A rod piece rests atop a pedestal draped with a blue silk sash embroidered with gold thread.

\end{DndReadAloud}

For more about the \textit{Rod of Seven Parts}, see this book's introduction.

\subparagraph{Treasure}
If a character shows the silk sash to Laysa, she identifies it as a garment given to her people that functions as a \textit{Robe of Stars}. The robe is one of the treasures Laysa mentioned earlier in this chapter.

The mundane treasure within the vault includes the following:


\begin{itemize}
\item 479 sp, 342 gp, and 179 pp
\item Jade game board with platinum playing pieces, worth 7,500 gp in total
\item Set of ruby earrings worth 500 gp
\item Three fire opals worth 1,000 gp each
\end{itemize}

\subparagraph{Exit}
The door in the middle of the north wall is locked and can be unlocked only by using Rerak's platinum skeleton key. The door opens into a cave tunnel that leads to the beach, about 1 mile south of the archaeologists' camp.

\section{Next Steps}
Once they've left the Tomb of Wayward Souls, the characters can return to the archaeologists' camp without incident. Laysa celebrates the characters' safe return; if the characters brought back any of the ancestral treasures she's been searching for, Laysa pays them as promised.

If Rerak is with the party, he volunteers to lend his knowledge of the tomb to help the archaeologists in their excavation. The archaeologists are apprehensive but nonetheless accept his aid.

While Rerak aids the archaeologists, he doesn't hurt them. Rerak won't return to Sigil with the characters; the simulacrum worries that the magic that sustains him might extinguish if he leaves Oerth. Rerak hopes to make a home on the Isle of Serpents with anyone who welcomes him.

The characters can return to Sigil through the doorway near the lagoon at their leisure.

\chapter{The Dragon Queen's Pride}
The adventurers discover that the final piece of the \textit{Rod of Seven Parts} is in the heart of Avernus, the first layer of the Nine Hells. The characters must infiltrate a diabolical casino dedicated to Tiamat, Queen of Evil Dragons, whose lair is nearby.

\section{Running This Chapter}
This chapter begins after the characters retrieve the sixth piece of the \textit{Rod of Seven Parts}. A character who holds this piece knows instinctively that the seventh piece is located somewhere inside a casino called the Red Belvedere in Avernus. The Wizards Three warn the characters that finding this location might be difficult since the Nine Hells are impossible to map and difficult to traverse.

Once the characters enter Avernus, they must hitch a ride on an infernal war machine and traverse the hellish landscape to reach the Red Belvedere. The characters must manipulate the Fiends in the casino's highest ranks to enter the exclusive Ruby Sanctum, then fight Tiamat's champion—or convince Tiamat to give up the final rod piece.

\subsection{Character Advancement}
The characters should be 17th level when this chapter begins. They gain a level once they retrieve the seventh and final piece of the \textit{Rod of Seven Parts} from the casino's Ruby Sanctum.

\subsection{Power of Secrets}
The individuals the characters meet in Avernus are primarily Fiends and other unpleasant beings. As a rule, these creatures lack the remorse, guilt, and shame that would give their secrets power if revealed to the characters. The characters don't learn any secrets applicable to the rules in "The Power of Secrets" section found in this book's introduction, though they can spend secrets during this chapter normally.

\subsection{Seventh Rod Piece}
The final piece of the \textit{Rod of Seven Parts} is in area N5 in the Ruby Sanctum, located in the Dragon's Pride area of the Red Belvedere casino. For more information about the rod and the spell this rod piece allows its wielder to cast, see this book's introduction.

\section{Avernus}
This chapter takes place in Avernus, the first layer of the Nine Hells of Baator.

\subsection{Knowledge of Avernus}
Characters who research Avernus in Sigil can learn the following:


\begin{description}
\item[Blood War.]
Avernus is a major battlefield for the Blood War, the raging conflict between demons and devils. It would be fruitless to fight every Fiend the characters come across in Avernus, since conflicts invariably escalate as infernal combatants join in against interlopers in the Nine Hells. To accomplish their goals, the characters might need to work with Fiends.
\item[Environment.]
Avernus is a hellscape in the most literal sense: a ravaged battlefield littered with decaying carcasses and wrecked war machines. A foul haze perpetually shrouds the skies. Ambient light swells above the horizon in a grim mockery of a sunset, but the sky is bereft of celestial bodies. Searing gusts of wind blast across the landscape, reeking of ash and brimstone. Due to this perpetual smoky twilight, there are no real days or nights in Avernus. Time passes as usual, though, and can be tracked in hour increments.
\item[History.]
Long ago, Avernus was a tantalizing paradise created by the archdevil Asmodeus to tempt mortals. However, the incursion of the River Styx and hordes of demons ruined that paradise, leaving behind gore and destruction.
\item[Queen of Dragons.]
Tiamat is the five-headed progenitor of the chromatic dragons in some realms and embodies the vices of evil dragonkind. Her lair is located in Avernus, connected to numerous monuments and temples scattered across the hellscape.
\end{description}

\subsection{Spell Modifications}
Drenched in the eldritch energy of the Nine Hells, Avernus warps magic used within its borders. At your discretion, a spell can be cosmetically modified to demonstrate the corruption of Avernus; these modifications don't change the spell's mechanical effects. The following are some examples:

\textit{Hunter's Mark}. A spectral crown of spiked iron encircles the target's head.

\textit{Mage Hand}. The conjured hand takes the shape of a taloned claw.

\textit{Magic Missile}. The missiles emit a haunting wail when fired.

\subsection{Traversing Avernus}
Avernus warps the senses, making distance impossible to gauge. The Wizards Three can't be sure where the characters will emerge in Avernus from the portal in Sigil, so they'll have to find their way to their destination. Residents of Avernus use vehicles known as infernal war machines to make trips more manageable. The Wizards Three suggest that the characters might be able to reach the Red Belvedere by using or commandeering one of these vehicles.

\section{Descending into Avernus}
When the characters step through the portal in Sigil, they emerge along a shelf in a mountainside cliff. Read or paraphrase the following to set the scene:

\begin{DndReadAloud}{}
Shrieking gusts of hot air assault your senses as you step into Avernus. Miasmic clouds roil over a ravaged landscape of broken rocks, bleached bones, and jagged metal. The dusty atmosphere reeks of sulfur and tar.

\end{DndReadAloud}

The characters know they need to travel to the Red Belvedere, but they don't know where the casino is located. The warped nature of the Nine Hells defies all mapping attempts, so the characters must figure out how to get to their destination.

If consulted, the \textit{Rod of Seven Parts} points away from the mountains and down the cliff, where a rough path skirts the cliff's base.

\subsection{Hell of a Ride}
The characters can descend the cliff without trouble. Once they do, they catch the attention of a crew of scavengers. Read or paraphrase the following:

\begin{DndReadAloud}{}
Out of the red dust clouds barrels an angular metal vehicle, ten feet tall, with spiked tires and hellfire roaring in its engine. Tortured screams issue from its furnace. The vehicle slows to a halt, and a winged silhouette jumps out and waves at you.

\end{DndReadAloud}

The infernal war machine Venatrix has room for four Medium or smaller creatures to ride inside the vehicle, including its driver, and for four Medium or smaller creatures to ride in its upper, open-air area. Three \textbf{erinyes} crew the Venatrix: Nykaia, who leaps out to greet the characters, and Kypris and Mykale, who are inside.

Nykaia is the group's glib and calculating leader, Kypris is their unflappable engineer and driver, and Mykale is their overeager weapons specialist. The three devils travel the wasteland of Avernus harvesting scrap metal and other valuable materials. The group is searching for additional crew members to aid them in their latest score.

Seeing that the characters lack a vehicle, Nykaia offers them a deal: if the characters help her with a job, she and her crew will give them a ride to wherever they'd like. If the characters refuse, skip to the "Stop for Directions" section.

The three erinyes are hoping to take down a \textbf{goristro}: a minotaur-like demon that carries demonic legions and rare treasures through the battlefields of the Blood War. They need someone to focus on the goristro while they take out the soldiers on its back.

Nykaia knows about the Red Belvedere but refuses to provide any more information until the characters help with the goristro. If the characters agree to her deal, she welcomes them aboard the Venatrix. Nykaia isn't hostile toward the characters and won't attack them unless the characters threaten her or start a fight.

If the characters attack Nykaia or attempt to take the Venatrix by force, Kypris and Mykale emerge from the war machine and join the fight. If the erinyes are in danger of losing the fight, they surrender and beg the characters to spare them, offering to drive the characters to the Red Belvedere in exchange.

\subsubsection{Driving the Venatrix}
The following rules are an adjusted, condensed version of the infernal war machine rules found in Baldur's Gate: Descent into Avernus.

\subparagraph{Driving}
The helm of an infernal war machine is a chair with a wheel, levers, pedals, and other controls. The helm requires a driver to operate.

While the Venatrix's engine is on, the driver can use an action to propel the vehicle up to its speed (see the "Vehicle Statistics" section below) or bring the vehicle to a dead stop. While the vehicle is moving, the driver can steer it along any course.

If the driver has the incapacitated condition, leaves the helm, or does nothing to alter the Venatrix's course and speed, the vehicle moves in the same direction and at the same speed it did during the driver's last turn until it hits an obstacle big enough to stop it.

As a bonus action, the driver can do one of the following:


\begin{itemize}
\item Start the Venatrix's engine or shut it off.
\item Cause the Venatrix to take the Dash or Disengage action while the vehicle's engine is running
\item Insert a \textit{Soul Coin} into the engine's furnace
\end{itemize}

\subparagraph{Soul Fuel}
The Venatrix's engine has a furnace fueled by Soul Coin (described later in this chapter). Among the vehicle's helm controls is a narrow slot into which Soul Coins can be fed. The vehicle's furnace consumes a Soul Coin instantly, expending all the coin's remaining charges at once and destroying the coin in the process. The soul trapped in the coin becomes trapped in the furnace instead, giving the vehicle 24 hours of fuel for each charge the Soul Coin had when it was consumed (maximum 72 hours). The furnace can hold any number of souls, their screams of anguish audible out to a range of 60 feet.

The Venatrix currently has enough soul fuel to run for 60 hours.

\subparagraph{Vehicle Statistics}
The Venatrix is a Gargantuan vehicle that can hold eight Medium creatures and carry up to 1 ton of cargo. It has a speed of 100 feet; AC 19; a damage threshold of 10; 200 hit points; immunity to fire, poison, and psychic damage; and a +4 bonus to Strength and Dexterity saving throws.

\subparagraph{Weapon Stations}
The Venatrix is equipped with two harpoon guns and an infernal screamer, which is a writhing humanoid torso made of melting wax with a barbed hand crank between its shoulder blades. The weapons are mounted to the vehicle's exterior, so any creature riding on the outside of the vehicle can operate a weapon. Each weapon station requires one creature to operate it.

A creature operating a harpoon gun can use an action to make a ranged weapon attack with it (+6 to hit, range 120 ft., one target). On a hit, the harpoon deals 10 (2d8 + 1) piercing damage.

A creature operating the infernal screamer can use an action to turn its crank, unleashing a telepathic shriek of agony at one target the creature can see within 120 feet of itself. The target must make a DC 15 Wisdom saving throw, taking 26 (4d12) psychic damage on a failed save or half as much damage on a successful one.

\subsection{Goristro Raid}
This encounter occurs only if the characters teamed up with the \textbf{erinyes} in the previous encounter. If the characters instead defeated the erinyes and claimed the Venatrix for themselves or bypassed the erinyes entirely, skip to the "Stop for Directions" section.

It takes about an hour for the Venatrix to travel to the goristro's location. As the Venatrix approaches the goristro, read or paraphrase the following:

\begin{DndReadAloud}{}
Thunderous footsteps shake the ground as you approach a horned, red-furred demon over twenty feet tall. Two winged shapes flank its towering form. On its back is a tented palanquin occupied by demonic soldiers swarming like insects.

\end{DndReadAloud}

The hulking \textbf{goristro} is flanked by two vrocks. Inside the goristro's palanquin are two hezrous and five quasits, as well as the "treasure" that the \textbf{erinyes} are after: a \textbf{unicorn} named Sterling in an iron cage.

Nykaia explains that Kypris will drive the war machine close to the goristro. The three erinyes will then fly to the palanquin to defeat the demons within while the characters slay the goristro.

The goristro defends itself. If the vrocks see that the goristro is losing, they join the fight. The three erinyes easily dispatch the demons in the palanquin.

\subparagraph{True Treasure Revealed}
Once combat finishes, the three erinyes carry down their claimed treasure:

\begin{DndReadAloud}{}
A large cage lands with a thud before the vehicle. Through the iron bars, you see the white fur and gilded mane of a unicorn as it nickers in fear.

\end{DndReadAloud}

Mykale excitedly explains that a unicorn's horn is a highly valuable ritual ingredient that will ingratiate the erinyes with powerful infernal war leaders. The unicorn desperately wishes to be freed and says as much in Celestial and Elvish.

Characters who try to convince the erinyes to release the unicorn must make a DC 30 Charisma (Persuasion) check. If the characters offer a rare or rarer magic item in exchange for Sterling's freedom, this check is made with advantage. On a successful check, the erinyes give the characters custody of Sterling. If the characters attempt to take Sterling by force, the erinyes become hostile and attack.

With Sterling (or a traded magic item) now in their possession, the three erinyes uphold their end of the deal and take the characters to the Red Belvedere. If the characters slay two of the erinyes, the third erinyes surrenders and offers to drive the characters to the Red Belvedere in exchange for her life.

\subsection{Stop for Directions}
If the characters are forced to search the plains and battlefields of Avernus for the Red Belvedere because they have no one to lead them there, have the party's navigator make a DC 20 Wisdom (Survival) check at the end of every hour of travel. After each successful check, the characters find something to steer them in the right direction; determine the source of this help by rolling on the Stop for Directions table. After four successful checks, the characters locate the Red Belvedere.

\begin{DndTable}[header=Stop for Directions]{cX}

d6 & Source of Help\\
1–3 & An iron signpost in the middle of nowhere points in the direction of the Red Belvedere.\\
4 & A friendly \textbf{imp} knows the way to the Red Belvedere. After tagging along for an hour and giving directions, the imp turns invisible and leaves.\\
5 & A \textbf{ghost} doomed to wander the battlefields of Avernus points the way to the Red Belvedere.\\
6 & An indifferent \textbf{bearded devil} hitchhiker offers directions, preferably in exchange for food, drink, or a free ride.\\
\end{DndTable}

\section{The Red Belvedere}
Eventually, the characters reach the Red Belvedere. The casino is a remnant of Avernus's original existence: a bucolic paradise designed to tempt mortal souls. Its edifice was destroyed in the carnage of the Blood War, but a champion of Tiamat named \textbf{Windfall} rebuilt the casino and rededicated it to the Dragon Queen as a celebration of her god's greed and vanity.

The Red Belvedere is more than a monument to the Dragon Queen. The casino is connected via a tunnel to Tiamat's lair, and Tiamat can hear every prayer and proclamation of her power uttered within the buildings' walls.

As the characters approach the Red Belvedere, read or paraphrase the following:

\begin{DndReadAloud}{}
Jagged mountains give way to an unexpectedly beautiful sight: a sprawling, palatial complex glittering with silver and gold. Light dances through the front rotunda's red stained-glass dome, reflecting off the overcast sky to shower the buildings with a rosy glow. A carved stone sign at the front of the complex reads in numerous languages, "The Red Belvedere."

\end{DndReadAloud}

If the crew of the Venatrix is with the party, the erinyes bid the characters goodbye at the front of the casino.

\subsection{Casino Lobby}
The Red Belvedere's layout is shown on map 8.1.

When the characters first enter the casino's rotunda lobby, read or paraphrase the following:

\begin{DndReadAloud}{}
Sweet perfume and glittering lights chase away all evidence of the war-torn landscape outside. A gilded statue of a roaring dragon presides over the casino's lobby. Distantly, you hear clattering coins and patrons laughing in preternaturally deep voices.

A winged tiefling soars above the spectacle. Patches of multicolored scales freckle her skin, and her prim tailcoat shimmers with every hue.

\end{DndReadAloud}

The winged tiefling is \textbf{Windfall}, a champion of Tiamat and the proprietor of the Red Belvedere (see the accompanying stat block). A performer at heart, Windfall is ostentatious and charismatic, making small talk with regular patrons and jovially welcoming new faces to the casino.

\subsection{Meeting Windfall}
As a dedicated champion of Tiamat, \textbf{Windfall} has been granted phenomenal power by her master. Her skin glitters with patches of multicolored scales, and in combat, her blade sings with all five of the chromatic dragons' elements. Windfall's enchanted tailcoat shimmers with the colors of the Dragon Queen, and she uses this coat to dazzle patrons and enemies alike.

When the characters arrive at the Red Belvedere, Windfall introduces herself and answers any questions the characters have about the Red Belvedere. She provides the following information:


\begin{description}
\item[Casino.]
The Red Belvedere is a decadent respite from the Blood War's horrors. Here, adventurers and devils alike can indulge in their vices in five different rooms, each with its own theme.
\item[Currency.]
The Red Belvedere has its own in-house currency (see the "Casino Currency" section). Both gold and Soul Coin can be traded for in-house currency at the exchange desk in the lobby.
\item[Dragon Queen Dedication.]
The casino is dedicated to Tiamat to celebrate the Dragon Queen's avarice and pride. Those who don't worship Tiamat are also welcome to partake in the casino's offerings.
\item[Rod Piece.]
If the characters mention the \textit{Rod of Seven Parts}, Windfall artfully dodges the topic. She claims that while the casino has been visited by several high-profile adventurers and scholars, she hasn't seen anything matching the rod piece's description. In truth, Windfall knows exactly where the final piece of the rod is: her private sanctum, adjacent to Tiamat's lair. A character who succeeds on a DC 32 Wisdom (Insight) check discerns only that Windfall is concealing information.
\end{description}

\subsection{Where's the Rod Piece?}
While inside the Red Belvedere, the character holding the sixth rod piece can determine that the seventh piece is located on the casino's premises one floor down.

The final piece of the \textit{Rod of Seven Parts} is held in Windfall's sanctum within the casino's exclusive club, Dragon's Pride. Only those personally invited by Windfall or one of the casino's pit masters are allowed inside Dragon's Pride.

If the characters ask her for permission to enter Dragon's Pride, Windfall scoffs and says they'll need to ingratiate themselves to a pit master to gain access to the exclusive area. It amuses the tiefling to watch the characters try to gain favor with at least one of the casino's pit masters (see the "Casino Pit Masters" section). Windfall is certain the characters will never convince the pit masters to let them into Dragon's Pride.

\subsection{Casino Currency}
The Red Belvedere uses a special type of in-house currency known as a talon. An exchange desk in the casino's lobby allows patrons to trade in coins for talons and vice versa. One talon is worth 10 gp.

A talon looks like an iridescent metal coin stamped with the silhouette of a dragon's claw on both sides. Since the Red Belvedere is dedicated to Tiamat, each talon is imbued with a mote of the Dragon Queen's power. Carrying and exchanging talons bear unpleasant consequences:


\begin{description}
\item[Curse.]
A creature that carries any number of talons is cursed. While cursed in this way, the creature has disadvantage on Intelligence and Wisdom saving throws as its mind is clouded by jealous and selfish thoughts. The curse ends on an affected creature only when the creature no longer holds any talons. A \textit{Remove Curse} spell or similar magic suppresses the curse's effects for 1 hour. When the curse ends on a creature, that creature must succeed on a DC 18 Constitution saving throw or gain 1 level of exhaustion, as the rampant avarice takes a physical toll.
\end{description}

\subparagraph{Soul Coins}
The devils of the Nine Hells also use a type of currency called Soul Coin. Soul Coins are used among the infernal hierarchy to purchase favors, bribe the unwilling, and reward the faithful for services rendered. A single Soul Coin can be exchanged at the casino for 300 talons and vice versa.

\subsubsection{Soul Coin}
\textit{Wondrous Item, Uncommon}

\textit{Soul Coins} are about 5 inches across and an inch thick, each minted from infernal iron. Each coin weighs one-third of a pound and is inscribed with Infernal writing and a spell that magically binds a single soul to the coin. Because each \textit{Soul Coin} has a unique soul trapped within it, each has a story. A creature might have been imprisoned as a result of defaulting on a deal, while another might be the victim of a night hag's curse.

\subparagraph{Carrying Soul Coins}
Anyone holding a Soul Coin feels the soul bound within it—overcome with rage, fraught with despair, or whatever emotion the soul exudes.

An evil creature can carry as many \textit{Soul Coins} as it wishes (up to its maximum weight allowance). A good or neutral creature can carry a number of \textit{Soul Coins} equal to or less than its Constitution modifier. A non-evil creature carrying a number of \textit{Soul Coins} greater than its Constitution modifier has disadvantage on its attack rolls, ability checks, and saving throws.

\subparagraph{Using a Soul Coin}
A Soul Coin has 3 charges. A creature carrying the coin can use its action to expend 1 charge from a Soul Coin and do one of the following:


\begin{description}
\item[Drain Life.]
You siphon away some of the soul's essence and gain 1d10 \textit{temporary hit points}.
\item[Query.]
You telepathically ask the soul a question and receive a brief telepathic response, which you can always understand. The soul knows only what it knew in life, but it must answer you truthfully and to the best of its ability. The answer is no more than a sentence or two and might be cryptic.
\end{description}

\subparagraph{Freeing a Soul}
Casting \textit{Remove Curse} or another spell that removes a curse on a Soul Coin frees the soul trapped within it, as does expending all of the coin's charges. The coin rusts from within and is destroyed once the soul is released. A freed soul travels to the realm of the god it served or the Outer Plane most closely tied to its alignment (DM's choice).

A soul can also be freed by destroying the coin that contains it. A \textit{Soul Coin} has AC 19, 1 hit point for each charge it has remaining, and immunity to all damage except damage dealt by a \textit{Hellfire Weapon} or an infernal war machine's furnace.

\subsection{Casino Security}
Each room of the casino is supervised by a trio of pit fiends. They are silent observers and intervene only if a patron starts a fight or is caught cheating (see the "Cheating in the Red Belvedere" sidebar).

On approaching a troublesome patron, the pit fiends confiscate the patron's talons and escort the patron from the casino. The patron is then barred from the casino for 24 hours. A barred patron can attempt to reenter by succeeding on either a DC 20 Dexterity (Stealth) check to avoid the pit fiends' notice or a DC 15 Charisma (Persuasion) check to convince the pit fiends to overlook the patron's past transgressions and allow reentry. Patrons barred more than once might be unable to reenter the casino for longer than 24 hours, at your discretion.

The pit fiends have no qualms about killing a patron who attacks them.

\subsection{Casino Pit Masters}
A pit master presides over each room of the Red Belvedere. (The rooms of the casino are displayed on map 8.1.) The characters must gain favor with at least one pit master to access Dragon's Pride. If the characters ask the patrons milling around the casino lobby, they can learn the information presented below about each pit master.

\subsubsection{Kaylan Renaudon}
\textit{Master of the Stygian Maze}

Kaylan is a devious \textbf{vampire} who arrived in Avernus after dodging a group of vampire hunters on the Material Plane. When Windfall refurbished the Red Belvedere, the tiefling employed Kaylan to design the ever-shifting halls of the casino's Stygian Maze. Kaylan now presides over the maze, gleefully watching patrons wander the confounding passages in search of riches.

Kaylan has a special \textit{Sword of Life Stealing} that, in addition to its usual properties, allows the vampire to draw sustenance from a soul trapped within a \textit{Soul Coin}. Feeding off a \textit{Soul Coin} in this way irrevocably destroys both the coin and the soul within, which has left Kaylan perpetually hungry for more.

\subparagraph{Gaining Kaylan's Favor}
Kaylan desires to feed off as many Soul Coins as he can. The characters can give Kaylan three Soul Coins in exchange for access to Dragon's Pride.

\begin{DndSidebar}[float=!b]{Cheating in the Red Belvedere}
Cheating by magical means is nearly impossible within the casino. Magical effects that enhance personal skill, such as Bardic Inspiration or the \textit{Enhance Ability} spell, don't directly tamper with the game and aren't considered cheating; however, casting an illusion on your hand of cards or invisibly moving dice draws the pit fiends' attention.



Physical methods of deceit (hidden cards, weighted dice, and so on) are more likely to pass unnoticed. A character attempting to cheat in this way must make a DC 30 Dexterity (Sleight of Hand) check; if the character is proficient with the gaming set of the game being played, the character can make this check with advantage. On a successful check, the character bends the game to their will. Encourage the player to describe how the character successfully cheats. On a failed check, the pit fiends notice the character's attempt, confiscate the character's talons, and escort the character from the casino. A \textbf{pit fiend} can be placated with a minimum bribe of 250 talons or 2,500 gp.



\end{DndSidebar}

\subsubsection{Khai Kiroth}
\textit{Master of the Scarlet Coliseum}

Khai is a boisterous \textbf{red abishai} (see appendix A) and the master of the casino's fighting pit. He relishes bloodshed, and his laughter booms over the gore-stained arena. Khai fancies himself the strongest of the pit masters, and he chafes under Windfall's leadership. Khai dedicates every death in the coliseum to Tiamat.

\subparagraph{Gaining Khai's Favor}
Khai prizes battle prowess above all else. Should the characters win three matches in the coliseum, Khai challenges them to combat; if the characters reduce Khai to 50 hit points or fewer without killing him, the abishai rewards them with access to Dragon's Pride. See the "Scarlet Coliseum" section later in this chapter for more information about this option.

\subsubsection{Nyssa Otellion}
\textit{Master of the Cerulean Hall}

A stoic and haughty \textbf{blue abishai} (see appendix A), Nyssa serves as the master of the casino's games of strategy and intellect. Nyssa prides herself on being an unbeatable foe in battles of wit. However, Nyssa doesn't hesitate to resort to underhanded methods to achieve victory, citing her opponent's inability to notice her cheating as a failure of intellect.

\subparagraph{Gaining Nyssa's Favor}
Nyssa admires those who are extremely clever and observant. See the "Cerulean Hall" section later in this chapter for more information about gaining favor with her.

\subsubsection{Rezran "Snake Eyes" Agrodro}
\textit{Master of the Viridian Den}

Rezran, a greedy \textbf{green abishai} (see appendix A), is the master of the casino's games of chance. His nickname "Snake Eyes" comes from his penchant for rolling ones in Triple Hydra (see the "Viridian Den" section for an explanation of the game's rules). Rezran is an opportunist to his core, always chasing after riches. He is a devoted follower of Tiamat and covets the Dragon Queen's amassed wealth.

\subparagraph{Gaining Rezran's Favor}
Rezran observes the characters as they partake in the Viridian Den's games. If the characters amass a total of 1,000 talons or more while in the den, Rezran is impressed with their luck and grants the characters access to Dragon's Pride.

\subsubsection{Uvashar}
\textit{Master of the Alabaster Racetrack}

Uvashar is a suave, white-furred \textbf{rakshasa} and the master of the casino's nightmare racing circuit. Uvashar can often be found indulging in fine food and music in a private viewing box.

Duplicitous and manipulative, Uvashar likes to assume the guise of a casino patron and sow discord among the racetrack's audience members. The rakshasa rigs races to add excitement to an event. Uvashar doesn't worship Tiamat like many of the casino's pit masters and patrons, but the rakshasa nonetheless maintains a grudging allegiance to the Dragon Queen.

\subparagraph{Gaining Uvashar's Favor}
Uvashar holds great esteem for those who keep their promises and uphold their end of a deal—especially when it requires a fair amount of skill to successfully pull off. One potential encounter to gain Uvashar's favor is presented later in the "Alabaster Racetrack" section later in the chapter.

\section{Casino Rooms}
The casino's rooms are described in the following sections.

\subsection{Alabaster Racetrack}
\begin{DndReadAloud}{}
Sleek white pillars surround what looks like a racetrack. The stands bustle with patrons, all eager to bet on the impending race.

\end{DndReadAloud}

This arena hosts nightmare racing. Audience members can place bets, earning a payout depending on which steeds place in the top two of a given race. Up to six nightmares race at any given time; those who bet on the first-place steed receive triple their initial bet, and those who bet on the second-place steed receive double their initial bet.

To simulate a race, have each participating character place a bet on a single number from 1 to 6—each number represents a nightmare. Roll a d6 twice to determine which nightmares win first and second place, respectively. If both rolls are the same number, reroll one die to determine which nightmare wins second place.

\subsubsection{To Cheat or Not to Cheat}
A few moments after the characters enter the racetrack's area for the first time, an elf in a stylish white suit approaches them and offers them a deal. Read or paraphrase the following:

\begin{DndReadAloud}{}
"Looking to make some coin, eh? So am I. I've got a little trick planned that'll cinch the race in my favor, winning me a thousand talons at least. You want in? I'll give you a cut of the winnings, and if it all goes swimmingly, I'll take you to the elite floor."

\end{DndReadAloud}

The elf is the racetrack's pit master Uvashar (see the "Casino Pit Masters" section earlier in the chapter), concealed with an illusory disguise. A character who succeeds on a DC 18 Intelligence (Investigation) check pierces the illusion and sees the rakshasa's true appearance.

Uvashar profits regardless of which nightmare wins the race; the rakshasa merely enjoys tempting mortals with offers to cheat and observing their reactions. If the characters accept Uvashar's deal, Uvashar gives the characters a \textit{Potion of Speed} and instructs them to slip the potion into the feed of the nightmare named Sunshroud. Uvashar then points the characters in the direction of the stables.

\subparagraph{Uvashar Unmasked}
If any character confronts Uvashar about their disguise, Uvashar rescinds the deal, bitterly commenting that the character has ruined the fun. The rakshasa nonetheless commends the character's observation skills before leaving to find another group to toy with. The characters must find another way to gain Uvashar's favor or ingratiate themselves to another pit master.

\subsubsection{Stables}
The racetrack's stables are depicted on map 8.2. If the characters enter the stable area, read or paraphrase the following:

\begin{DndReadAloud}{}
Six nightmares snort and paw at the ground in their respective stalls, their flaming manes dancing as devil jockeys attend to the steeds. The reeking stench of offal rises from the feed buckets. In one corner stands a bored-looking pit fiend.

\end{DndReadAloud}

One \textbf{pit fiend} serves as security for the area. The six nightmares (named Chicanery, Far Gambit, Blinx, Cursed Cavalry, Sunshroud, and Retribution) stamp impatiently in their stalls, and a \textbf{barbed devil} jockey attends to each one.

Characters must succeed on a DC 22 Dexterity (Stealth) check each time they enter or exit the stables to avoid being noticed and attacked by the pit fiend.

\subparagraph{Rigging the Race}
Successfully slipping the \textit{Potion of Speed} into Sunshroud's feed bucket requires multiple steps:


\begin{description}
\item[Step 1.]
One or more characters must successfully sneak into the stables.
\item[Step 2.]
After locating Sunshroud's stall, a character must succeed on a DC 25 Dexterity (Sleight of Hand) check to pour the potion into the feed bucket. On a failed check, Sunshroud and its jockey notice the character and make a ruckus, attracting the attention of the pit fiend.
\item[Step 3.]
The potion reacts with the offal in the nightmare's feed bucket, causing a cloud of toxic gas to erupt. Characters within 5 feet of the feed bucket must succeed on a DC 20 Constitution saving throw or have the poisoned condition for 1 hour. On a failed saving throw, the character's coughing also alerts the pit fiend.
\item[Step 4.]
The characters must successfully sneak out of the stables.
\end{description}

If a character executes all four steps without attracting the attention of the pit fiend, Sunshroud consumes the \textit{Potion of Speed} and wins handily.

\subsubsection{Pit Master's Confrontation}
If the characters successfully rig the race in Sunshroud's favor, a pit fiend approaches the party after the race. The pit fiend explains that the racetrack's pit master wants to meet with the characters, and it escorts the party to Uvashar's private viewing box. Read or paraphrase the following if the characters accompany the pit fiend to the viewing box:

\begin{DndReadAloud}{}
This box overlooks the racetrack, far from the crowds in the stands. Charcuterie plates line the tables. Reclining in a chair is a bipedal fiend with fur like a white tiger, dressed in a sharp suit and watching you slyly.

\end{DndReadAloud}

Uvashar cuts straight to the point and asks the characters if they tampered with the race.

\subparagraph{Telling the Truth}
If the characters tell the truth, Uvashar tsks before magically appearing as the suited elf from earlier. The rakshasa admonishes the characters and dismisses them without reward. The characters must find an alternate way into Dragon's Pride.

\subparagraph{Lying}
A character who attempts to lie about the party's involvement must make a DC 18 Charisma (Deception) check. On a successful check, Uvashar is impressed. On a failed check, Uvashar critiques the character's skill before laughing and commending the party's commitment. In either case, Uvashar honors the deal, offering the party access into Dragon's Pride (see the "Dragon's Pride" section later in this chapter).

\subsection{Cerulean Hall}
\begin{DndReadAloud}{}
Blue crystal pillars reach toward this gaming hall's ceiling. All is quiet except for the clinking of coins and shuffling of cards. Walking between tables is a lean, blue-scaled devil with dragon wings.

\end{DndReadAloud}

The blue devil is the room's pit master, Nyssa Otellion (see the "Casino Pit Masters" section). Nyssa stoically moves between game tables, searching for a worthy opponent.

\subsubsection{Games in the Cerulean Hall}
The Cerulean Hall is home to games of strategy and intellect, such as dragonchess and three-dragon ante. The room is divided into three sections based on skill: an amateur players' section, an intermediate players' section, and an expert players' section.

The minimum bet to participate in a game is 15 talons, regardless of section. At the start of a game, all participants place their starting bets in a central pot; games like dragonchess are typically limited to two participants, but other games can have five or more participants. For each turn of the game, all participants make Intelligence checks using the relevant gaming set (a dragonchess set, a playing card set, or another suitable set). The participant with the highest total roll wins the turn. In the event of a tie, no one wins the turn. The first participant who wins three turns wins the game and collects the pot.

For NPC participants, the following table shows the average bet and skill modifier to use depending on the section the NPC is playing in.

\begin{DndTable}[header=Cerulean Hall Betting]{XXc}

Section Level & Average Bet & Modifier\\
Amateur & 17 (1d4 + 15) talons & +4\\
Intermediate & 20 (2d4 + 15) talons & +8\\
Expert & 25 (4d4 + 15) talons & +12\\
\end{DndTable}

\subsubsection{A Contest of Wits}
Any character who wins a game in the experts' section gains Nyssa's attention. The pit master approaches the character after the game and offers a challenge of wits. Read or paraphrase the following:

\begin{DndReadAloud}{}
"Not bad. You are clever—I acknowledge that much. But I wonder how you would fare against the master of this chamber. I hereby challenge you to a dragonchess match: winner takes all."

\end{DndReadAloud}

If the character accepts her challenge, Nyssa insists on playing the match immediately. Nyssa wagers 50 talons on this match of dragonchess, and her modifier is +12.

\subparagraph{Nyssa's Cheating}
Nyssa is wearing a \textit{Medallion of Thoughts}, which she uses to scan her opponent's surface thoughts and stay one step ahead in the match. Nyssa uses a \textit{Disguise Self} spell to hide the medallion from view. As long as Nyssa's opponent can be affected by divination magic, the medallion allows Nyssa to treat any roll of a 9 or below as a 10 when she makes Intelligence checks using a dragonchess set.

A character that succeeds on a DC 20 Intelligence (Investigation) check can pierce the illusion and see the medallion, which can be identified by any character who succeeds on a DC 17 Intelligence (Arcana) check. A character that casts \textit{Detect Magic} or similar magic notices the aura of illusion magic concealing Nyssa and the medallion.

If a character accuses Nyssa of cheating, she asks the character for proof. If the character correctly identifies the medallion, Nyssa smiles widely and commends the character's ability to see through her ruse. She then concedes the match.

\subparagraph{End of the Match}
If the character wins the dragonchess match or Nyssa concedes, Nyssa compliments the character's intellect and thanks the character for the most thrilling match she's had in a while. She then offers the character and their companions access to Dragon's Pride (see the "Dragon's Pride" section).

If Nyssa wins, she scoffs at the character and haughtily collects her winnings.

\subsection{Scarlet Coliseum}
\begin{DndReadAloud}{}
Jeers and shouts rise from this open-air coliseum made of deep-red stone. Over the sound of roars and clattering metal, a booming voice excitedly commentates on the fight.

\end{DndReadAloud}

The booming voice is the magically amplified voice of Khai Kiroth, the coliseum's pit master and commentator (see the "Casino Pit Masters" section earlier in this chapter).

Watching a coliseum match is free. Matches pit two teams against each other: one consisting of creatures enlisted by the casino, the other consisting of casino patrons looking to prove their mettle in combat. Before each match, audience members can bet on which team will emerge victorious. At the conclusion of a match, audience members who bet on the winning team receive twice their initial bet.

\subsubsection{Entering the Coliseum}
It costs 15 talons to enter the coliseum as a combatant. A maximum of three combatants are allowed on the patrons' team. If all three combatant slots aren't filled by characters, the remaining slots are filled by veterans.

The arena matches are divided into three tiers, with victory in a tier granting access to the next one. Victory in a match goes to the side that reduces all combatants on the other side to 0 hit points. If the patrons' team wins a match, each combatant on the team receives 100 talons. The matches have one-hour breaks between them, allowing participants to take a short rest in the coliseum's wings.

Map 8.3 depicts the coliseum's arena. The three tiers of matches are as follows:


\begin{description}
\item[Tier 1.]
Barlguras are released into the arena—one for each member of the opposing team.
\item[Tier 2.]
A \textbf{purple worm} erupts through the floor of the arena.
\item[Tier 3.]
Two iron statues built to resemble pit fiends lumber into the arena. They use the \textbf{iron golem} stat block but have iron wings that give them a flying speed of 30 feet.
\end{description}

If a participant wins a match but doesn't wish to continue to the next tier, the participant can pay an exit fee of 15 talons and swap with another willing participant. A participant's companions can join them in the coliseum's wings between matches.

\subsubsection{Pit Master's Challenge}
Characters who win three matches gain Khai's attention. Khai meets them in the coliseum's wings after the match and speaks with them. Read or paraphrase the following:

\begin{DndReadAloud}{}
"This arena hasn't seen warriors like you for quite some time. Would you indulge in one final bout—against me? Truly, a special event! And, of course, I offer splendid payment for your victory. How does one thousand talons sound, along with access to our casino's exclusive club floor, Dragon's Pride?"

\end{DndReadAloud}

Khai challenges the characters to one final arena match—against him. Since this is a special match, he allows all the characters to fight. If the characters agree, Khai declares that they shall meet in the arena in one hour.

\subparagraph{Final Showdown}
An excited crowd has gathered in the coliseum. Khai has two horned devils on his team. Khai fights ruthlessly but offers cheerful commentary during the combat, lauding the characters for their skill and mocking their failings.

On being reduced to fewer than 25 hit points, Khai concedes the match, declaring the characters the victors. As promised, Khai hands the characters 1,000 talons and offers them access to Dragon's Pride (see the "Dragon's Pride" section).

\subsection{Stygian Maze}
\begin{DndReadAloud}{}
This area's foyer is carved from glistening obsidian. An austere, fanged man in a high-collared cloak sits at an intricately carved desk, next to which is a simple stone doorway. Behind the desk, a floor-to-ceiling window overlooks a stone maze shrouded in fog.

\end{DndReadAloud}

The man behind the desk is the maze's pit master, the vampire Kaylan Renaudon (see the "Casino Pit Masters" section earlier in this chapter). Kaylan provides the following information:


\begin{description}
\item[Maze.]
The maze challenges participants' navigational skills and adventuring prowess. Participants traverse magically shifting hallways in search of treasure, avoiding perils along the way.
\item[Participation.]
Entering the maze costs 20 talons per participant. Participants can explore the maze alone or in groups.
\item[Treasure.]
Treasure stashes found throughout the maze typically contain talons, though some also include Soul Coin.
\end{description}

A character who succeeds on a DC 15 Wisdom (Insight) check notices a flicker of hunger appear in Kaylan's eyes when he mentions Soul Coin. If pressed, Kaylan admits that he feeds off \textit{Soul Coins} using a special \textit{Sword of Life Stealing}, which he always keeps with him. If the characters bring him \textit{Soul Coins}, he'll give the characters 500 talons for each coin.

Bringing three or more \textit{Soul Coins} to Kaylan earns his admiration, and he offers the characters access to Dragon's Pride in return (see the "Dragon's Pride" section).

\subsubsection{Maze Features}
The maze's passageways constantly move, making the maze impossible to map. All areas of the maze have the following features:


\begin{description}
\item[Ceilings.]
Ceilings are solid but appear to be made of thick fog. Ceilings are 15 feet high in rooms and 10 feet high in passageways.
\item[Floor and Walls.]
The floor and walls are constructed of smooth black stone.
\item[Lighting.]
The maze is dimly lit by hovering globes of magical light.
\item[Room Size.]
Unless otherwise stated, all chambers within the maze are circular rooms 30 feet in diameter.
\end{description}

\subsubsection{Navigating the Maze}
Before the characters explore the maze, have them designate a leader. If a character is exploring the maze alone, that character is the leader. The leader then ventures through the maze, making a DC 20 Intelligence (Investigation) or Wisdom (Survival) check to discern the best pathway. On a successful check, the leader correctly navigates to a treasure chamber; roll randomly or choose an option from the Treasure Chambers table to determine what room is encountered. On a failed check, the leader stumbles into a challenge chamber; roll randomly or choose an option from the Challenge Chambers table to determine which room is encountered. Each room should be encountered only once. If you roll a result the characters have already encountered, or if the characters would leave the maze before finding a chest with one or more \textit{Soul Coins}, reroll or choose another option.

\begin{DndTable}[header=Treasure Chambers]{cX}

d10 & Chamber\\
1–3 & A simple table sits at the center of this chamber. On the table are two Potion of Greater Healing and a bag of 30 talons.\\
4–6 & This chamber contains a gurgling pool of turquoise water. A creature that spends 1 minute bathing in the pool gains 2d10 temporary hit points.\\
7–8 & Planters filled with sweet-smelling flowers line this room. At the center of the room is an unlocked wooden chest containing 75 talons.\\
9 & This chamber contains a locked iron chest. A character can open the chest with a successful DC 15 Dexterity (Sleight of Hand) check using thieves' tools, or it can be forced open with a successful DC 17 Strength (Athletics) check. A \textit{Knock} spell or similar magic also opens the chest. Inside the chest are 100 talons and three Soul Coin.\\
10 & One of this room's walls has a simple stone doorway, which leads out of the maze and returns the characters to the maze lobby.\\
\end{DndTable}

\begin{DndTable}[header=Challenge Chambers]{cX}

d10 & Chamber\\
1–3 & Five mimics, all disguised as treasure chests, occupy this room. Disturbing one mimic causes them all to attack.\\
4–6 & This room is filled with 1-foot-deep water. A thin wooden plank spans the water, connecting the room's entrance and exit. A creature crossing the plank must succeed on a DC 17 Dexterity (Acrobatics) check or fall into the water. The first time a creature touches the water, three water elementals appear and attack.\\
7–8 & A hostile but sleeping \textbf{behir} lies curled around a chest. A creature that approaches the behir must succeed on a DC 16 Dexterity (Stealth) check to avoid being noticed. The chest is unlocked and contains 50 talons and one \textit{Soul Coin}.\\
9 & A large vase filled with 100 talons sits in the center of the room. The vase is guarded by three invisible stalkers, which attack any creature that touches the vessel.\\
10 & A hostile \textbf{iron golem} patrols this room and immediately attacks intruders. If the golem is destroyed, a character who searches the golem's remains finds three Soul Coin.\\
\end{DndTable}

\subsection{Viridian Den}
\begin{DndReadAloud}{}
The walls and ceilings of this cozy den are made of luminescent jade. Rowdy patrons gather around game tables, scowling at dice and cards as they test their luck. A green-scaled draconic devil moves between games, idly chatting with players.

\end{DndReadAloud}

The green devil is the room's pit master, Rezran Agrodro (see the "Casino Pit Masters" section). Rezran observes the room, occasionally swooping in to speak with a patron who's on a winning streak.

\subparagraph{Pit Master's Favor}
If the characters have 1,000 talons or more and are in the Viridian Den, Rezran is impressed and invites them to Dragon's Pride.

\subsubsection{Games in the Viridian Den}
The characters can partake in two games while visiting the Viridian Den.

\subparagraph{Drake's Auction}
This card game has a minimum bet of 20 talons, which goes into a pot at the center of the table. The goal of the game is to amass the highest-ranking hand of cards.

At the start of the game, the croupier deals each participant a pair of cards. The croupier draws three cards to lay face up next to the pot; these cards are up for auction. Participants bid on the auction cards, and the participant with the highest bid wins the card and adds their bid to the pot. Once all three auction cards have been sold, all participants reveal their hands, and the highest-ranking hand wins the whole pot.

To simulate play of this game, have each participant roll 2d12 in secret. Then, roll 3d12 and show the numbers rolled—these represent the auction cards. On winning a bid, the participant can add the number rolled on the die won to their total roll. At the end of the game, the participant with the highest total wins. If the highest totals are tied, no one wins.

\subparagraph{Triple Hydra}
The minimum bet for this game is 25 talons. In this dice game, each participant rolls 3d6. If the participant rolls none of the same number, the participant wins nothing. If the participant rolls two of the same number, the participant wins an amount equal to twice their starting bet. If the participant rolls three of the same number, the participant wins an amount equal to triple their starting bet.

If a participant rolls two or more 1s, this roll is called snake eyes. In addition to the normal payout, a participant who rolls snake eyes also gets an additional payout equal to their starting bet.

\section{Dragon's Pride}
Occupying the lower floor of the Red Belvedere, Dragon's Pride is an exclusive club reserved for the casino's most valued patrons. Here, club members can spend talons on opulent jewelry, refreshing cocktails, and rejuvenating spa treatments.

This floor is also home to the casino proprietor's office, which connects to a hidden sanctum used to communicate directly with Tiamat. Within this sanctum, Windfall keeps the final piece of the \textit{Rod of Seven Parts}, hoping to amplify the rod's ability to cast \textit{Simulacrum} with Tiamat's power and re-create her master's beloved lost creation: Sardior, the ruby dragon.

\subsection{Entering Dragon's Pride}
Dragon's Pride occupies the lower floor of the casino and is accessed via an elevator. This elevator is magically concealed within the back of the large dragon statue in the casino's lobby. To view the elevator's entrance and pass through, a creature must bear the club's mark: an enchanted tattoo of a dragon claw. Only Windfall and the casino's pit masters can bestow or remove this tattoo. As soon as the tiefling or a pit master grants permission to enter Dragon's Pride, the mark magically appears on a creature's hand. To creatures lacking the mark, the club's elevator is undetectable and impassable by any other means.

Once the characters have gained favor with at least one of the casino's pit masters, that pit master marks each of the characters with the club's tattoo and instructs the characters to head through the back of the lobby's statue. Now bearing the club's mark, the characters can see the entrance to Dragon's Pride:

\begin{DndReadAloud}{}
What was once a solid statue now appears translucent. Within, a stone platform slowly levitates between this floor and the one below.

\end{DndReadAloud}

As long as they bear the mark, the characters can enter and exit Dragon's Pride at their leisure.

\subsection{Features of Dragon's Pride}
Unless otherwise stated, the areas of Dragon's Pride have the following features.

\subsubsection{Ceilings, Floors, and Walls}
The ceilings, floors, and walls are carved from smooth red stone that's slightly warm to the touch. The ceilings are 15 feet high unless otherwise noted.

\subsubsection{Lighting}
Rooms are brightly lit by multicolored magical sconces. Tapping a sconce turns it off.

\subsubsection{Security}
A \textbf{pit fiend} stands guard in each room. The pit fiends are silent observers who intervene if a fight breaks out.

\subsubsection{Wards}
Magical wards surround Dragon's Pride. If a creature without the club's mark attempts to teleport into Dragon's Pride, the creature takes 33 (6d10) psychic damage, and the attempt fails.

\subsection{Dragon's Pride Locations}
The areas of Dragon's Pride are keyed to map 8.4.

\subsubsection{N1: Wyrmsong Cantina}
\begin{DndReadAloud}{}
Lavish couches surround low tables in this dimly lit lounge. At the center of the room is a circular bar. Liquor bottles on shelves flash colors in time with faintly playing jazz, while a frost-blue devil pours drinks behind the bar.

\end{DndReadAloud}

The bartender is an \textbf{ice devil} named Oganath, who uses its frigid claws to instantly chill the drinks it serves. At any given time, 2d4 horned devils relax in the lounge.

\subparagraph{Oganath's Cocktails}
An expert in arcane mixology, Oganath serves a myriad of magic drinks at the bar:


\begin{description}
\item[Blazing Bloody Jack (Costs 100 Talons).]
Motes of coagulated, red juice are suspended throughout this amber-colored cocktail, which tastes both sweet and spicy. When consumed, this drink functions as a \textit{Potion of Heroism}.
\item[Joker's Sky (Costs 125 Talons).]
This lustrous blue drink is topped with fluffy, white cream. When consumed, this drink functions as a \textit{Potion of Flying}.
\item[Salubra Slinger (Costs 50 Talons).]
This fizzy, berry-flavored drink causes the imbiber's mouth to feel slightly numb. When consumed, this drink functions as an \textit{Elixir of Health}.
\end{description}

A drink loses its magic if it is removed from the Red Belvedere. Oganath also offers a variety of mundane cocktails, which cost 3 talons each.

\subsubsection{N2: Enchanting Arcana}
\begin{DndReadAloud}{}
A sleek rakshasa in shimmering teal robes runs this arcane jewelry store. Elegant, bejeweled accessories are displayed for sale on the counter.

\end{DndReadAloud}

The store is run by a vainglorious \textbf{rakshasa} lapidary named Krysocol, who has several pieces for sale:


\begin{itemize}
\item Two black jade pendants (each an \textit{Amulet of Proof against Detection and Location}) for 300 talons each
\item A brass necklace hung with citrine spheres (a \textit{Necklace of Fireballs} with five beads) for 1,250 talons
\item A purple ring with white diamond stars (a \textit{Ring of Shooting Stars}) for 2,000 talons
\end{itemize}

\subparagraph{Forged in Hell}
Any magic item created by Krysocol bears the following additional effect, which the rakshasa is loath to disclose:


\begin{description}
\item[Curse.]
As long as you wear this item, you are cursed. The item can't be removed until you are targeted by a \textit{Remove Curse} spell or similar magic. While cursed in this way, every time you finish a long rest, you must make a DC 10 Charisma saving throw. On a failed save, you immediately drop everything you are carrying or wearing and are transformed into a \textbf{lemure}. Only a \textit{Wish} spell can reverse this transformation.
\end{description}

\subsubsection{N3: Virtuous Vices}
\begin{DndReadAloud}{}
Fluffy clouds of steam billow from this spa, redolent of ginger and citrus. Two erinyes—one in a white robe, one in black—stand behind the counter. Behind them are several hot springs and massage stations.

\end{DndReadAloud}

The spa is run by two \textbf{erinyes} named Vitia and Vertu. They offer two services:


\begin{description}
\item[Brimstone Springs (Costs 100 Talons).]
This service takes 1 hour and includes a relaxing hot spring bath and massage. A creature who partakes in this service gains 4d10 + 4 \textit{temporary hit points} in addition to the benefits of a short rest.
\item[Luxe Reawakening (Costs 450 Talons).]
This rejuvenating service takes 3 hours. A creature who partakes in this service gains 6d10 + 6 \textit{temporary hit points} in addition to the benefits of a long rest.
\end{description}

\subsubsection{N4: Windfall's Office}
The door to the office is locked and guarded by two pit fiends. Each pit fiend standing guard also carries a key to the office door for emergency purposes. Opening the door without the proper key requires thieves' tools and a successful DC 20 Dexterity (Sleight of Hand) check, or a \textit{Knock} spell or similar magic.

If the office door is opened by any means other than the key, its trap activates, spraying colorless poisonous gas in a 15-foot cone. Each creature in that area must succeed on a DC 20 Constitution saving throw or take 18 (4d8) poison damage. This trap can't be disarmed.

The office's interior is devoid of occupants. When the characters venture inside, read the following:

\begin{DndReadAloud}{}
Bookshelves line the walls of this room. A bed draped in plush linens is built into an alcove, across from which is a simple work desk covered in papers and books. Mounted on the back wall is the sculpted bust of a smiling red dragon.

\end{DndReadAloud}

The desk is covered with notes about magical constructs, as well as religious and historical texts written in Draconic. A character who spends 10 minutes studying these items learns the following:


\begin{description}
\item[The First World.]
Tiamat worked alongside the platinum dragon Bahamut to shape the Material Plane in the form of a single First World. An unknown cataclysm sundered this First World, resulting in the numerous universes now occupying the Material Plane.
\item[The Ruby Dragon.]
While creating the First World, Bahamut and Tiamat also created Sardior, a dragon with ruby scales. Sardior was destroyed in the sundering of the First World.
\item[The Simulacrum Project.]
Windfall acquired the Rod of Seven Parts of the \textit{Rod of Seven Parts} some time ago. She used it to create an image of Sardior, who she believes resembled a red dragon. Windfall hopes to find a way to use the rod piece to bring the image to life.
\end{description}

\subparagraph{Secret Entrance}
The bust on the wall projects an illusion that conceals the entrance to the Ruby Sanctum (area N5). A \textit{See Invisibility} spell or similar magic reveals the entryway's outline in the wall. A character searching the room for secret doors finds this hidden entrance with a successful DC 18 Intelligence (Investigation) check.

\subparagraph{Treasure}
Most of the items in Windfall's office are lavish but mundane and too big to transport. However, a character who searches the room finds a \textit{Tome of Leadership and Influence} tucked into a bookcase.

\subsubsection{N5: Ruby Sanctum}
\begin{DndReadAloud}{}
The walls of this cavernous space are made of faceted, ruby-red crystal. A thin tunnel juts from the back of the cavern and leads into darkness.

A stalagmite rises from the floor in the center of the room. A similarly sized stalactite hangs from the ceiling directly above it. The space between the two is occupied by a crystalline rod piece pulsing with energy. The piece projects the holographic image of a sleeping red dragon. Windfall stands near the projection, watching it intently.

\end{DndReadAloud}

The ceilings in this chamber are 45 feet high. \textbf{Windfall} is hostile to trespassers in the sanctum.

The crystalline object floating between the stalagmite and stalactite is the Rod of Seven Parts of the \textit{Rod of Seven Parts}. Attempting to remove the rod piece from its position immediately causes the image of the dragon to disappear, prompting Windfall to attack.

\subparagraph{Draconic Intervention}
When Windfall is first reduced to 0 hit points, multicolored light surrounds her as five hissing voices echo throughout the chamber:

\begin{DndReadAloud}{}
"Upstart insects! Who dares make a racket outside my lair? Who dares interrupt the musings of the Dragon Queen?"

\end{DndReadAloud}

The facets of the sanctum's walls reflect images of Tiamat's eyes as she talks. Windfall regains 300 hit points but has the paralyzed condition due to Tiamat's magic as the Dragon Queen speaks to the characters directly.

\subparagraph{Retrieving the Rod Piece}
If the characters inform Tiamat of Vecna's plot, Tiamat growls in disdain before grudgingly allowing the characters to leave with the final piece of the rod. She explains that the Material Plane is partly her creation—and she refuses to let an overzealous archlich unmake it. The Dragon Queen then commands the characters to leave the casino at once.

Otherwise, Tiamat releases Windfall from her magical restraints, instructs the tiefling to eradicate the party, and leaves to commune with her other champions. Only after Windfall is defeated can the characters acquire the final piece of the rod. For more about the \textit{Rod of Seven Parts}, see this book's introduction.

\subparagraph{To Tiamat's Lair}
The narrow tunnel to the south eventually leads to Tiamat's lair. This entrance is currently guarded by Drekarvynix, an \textbf{ancient red dragon}, who demands the characters leave unless they can prove their legitimate business in Tiamat's lair. Unless the characters provide a compelling reason to enter, Drekarvynix attacks them. The red dragon doesn't pursue the characters back through the tunnel to Dragon's Pride.

\section{Next Steps}
Once the characters have retrieved the last piece of the rod, they can go back to the doorway through the mountainside cliff in Avernus and return to Sigil without further complications.

\chapter{The Betrayer Revealed}
Once the characters have all of the \textit{Rod of Seven Parts}, all eyes turn to Alustriel, Tasha, and the impostor Mordenkainen. The plan was for the characters to work with the Wizards Three to pin down Vecna's location and stop the archlich's Ritual of Remaking. However, this chapter reveals that the characters—as well as Alustriel and Tasha—have been deceived.

The fake Mordenkainen reveals himself to be \textbf{Kas the Betrayer}, Vecna's former ally and current archenemy. The characters must chase Kas across the multiverse to Pandesmos, the first layer of the plane of Pandemonium, where the vampire plans to co-opt Vecna's ritual and become the most powerful being in existence.

\section{Running This Chapter}
This chapter begins after the characters retrieve the final piece of the \textit{Rod of Seven Parts} from Avernus. Whether the characters are carrying the previous six rod pieces or those pieces remained in Sigil, the rod knits itself into one artifact when all the pieces are within a few inches of each other.

The rod becomes complete, but no character is likely yet attuned to it, and the Wizards Three in Sigil asked the characters to return to the sanctum once all the rod pieces were found. If the characters hesitate after retrieving the final rod piece, encourage them to quickly return to their allies at the sanctum in Sigil. The characters' return to the sanctum triggers the next part of the story.

This chapter's early scenes include the chaos caused when Kas steals the \textit{Rod of Seven Parts} and gravely wounds Alustriel and Tasha. The chapter progresses to Pandesmos, where the characters find a war erupting between Kas's and Vecna's demonic forces. The characters must stay on Kas's trail if they hope to find Vecna, potentially neutralize Kas, and end the threat to the multiverse.

\subsection{Character Advancement}
The characters should be 18th level when this chapter begins. The characters gain a level after they navigate Pandesmos Outlook and arrive at Carapace Ridge, where the war between the forces of Kas and Vecna rages.

\subsection{Power of Secrets}
The characters can learn one secret in this chapter that is applicable to the rules in "The Power of Secrets" section in this book's introduction:


\begin{description}
\item[Malaina's Secret.]
Alustriel's wife was suspicious of Mordenkainen since he arrived in the sanctum, but she said nothing to Alustriel, Tasha, or the characters. The characters can learn this secret in the "Malaina's Mistrust" section later in this chapter.
\end{description}

\section{Returning from Avernus}
When the characters return from Avernus with the final piece of the \textit{Rod of Seven Parts}, Kas—disguised as Mordenkainen—greets them inside the sanctum. He appears giddy yet terrified, unable to hide his honest emotions.

As soon as the disguised Kas sees the characters step through the doorway with the final rod piece, he exclaims the following:

\begin{DndReadAloud}{}
"My friends! Have you got it? The final piece? I can tell from the spring in your step that you do! How exciting!

"While you've been away, we've deduced that Vecna needs to conduct his ritual in a place of great import where he can generate power from the secrets his cult has stolen.

"It's likely to be a far-flung location, and on a plane we'd not expect. I'd say it's like finding a needle in a haystack, but I have colleagues who might help us determine the location. I'll speak with them today, and hopefully we'll have our answer.

"You've all been through a lot. You should rest in your quarters. In the meantime, I'll make sure the rod is properly reassembled and working to its full potential. You're going to need your energy for when we face Vecna!"

\end{DndReadAloud}

Kas strongly encourages the characters to rest now and leave the \textit{Rod of Seven Parts} in the wizards' care. Kas wants to keep the characters away from the rod and the wizards for the next few hours. Ever the villain, Kas plans to take Alustriel and Tasha by surprise, steal the rod, and flee to Pandemonium to free \textbf{Miska the Wolf-Spider} before confronting Vecna. Even if a character is already attuned to the rod, Kas plans to steal it.

Despite what he has said while posing as Mordenkainen, Kas knows Vecna is conducting his Ritual of Remaking on the plane of Pandemonium. After receiving a message in a dream from the Dark Powers, Kas realized that the magic inherent in his \textit{Crown of Lies} would transport him through Alustriel's portal to roughly the location of Vecna's ritual. Now that the characters have retrieved the \textit{Rod of Seven Parts}, the characters are mere inconveniences to Kas, but he doesn't want to fight them, as the fight might significantly weaken him, if not worse.

\subsection{If the Characters Decline}
If the characters are suspicious, or if they are being particularly cautious, they might not want to hand over the \textit{Rod of Seven Parts} or rest in their quarters. Below are options for handling characters who don't take Kas's suggestions.

\subsubsection{The Other Wizards}
Kas told Alustriel and Tasha that he wishes to tinker with the complete \textit{Rod of Seven Parts}. He has convinced the others that he can fortify the artifact to prevent it from shattering again when the characters use it to stop Vecna's ritual. If the characters insist on keeping the rod, Alustriel and Tasha approach them and explain this. They ask the characters to relinquish the rod, acknowledging that Mordenkainen might be a bit eccentric, but his magical theories are sound. Alustriel and Tasha extol the virtues of ensuring that the rod won't sunder again.

\subsubsection{Missing Malaina}
If the characters resist resting in their quarters, Alustriel approaches them privately. She's concerned that her wife, Malaina, left the sanctum several hours ago and hasn't returned. The lengthy absence isn't like Malaina, Alustriel explains, and the wizard is worried for her wife's safety. Malaina is resisting any magical attempts to contact her, Alustriel tells the characters.

When she left, Malaina told Alustriel that she was going to a shop called Velgar's Vittles in Sigil's Market Ward. The Wizards Three were busy, so Alustriel didn't ask for more information. Now, Alustriel expects to imminently be occupied helping Mordenkainen and Tasha locate Vecna. Alustriel can't leave the sanctum, so she asks the characters to find her wife before the group leaves to confront the archlich. Alustriel provides directions to Velgar's Vittles.

In reality, Malaina was suspicious of Mordenkainen, but she didn't want to make accusations without proof. Malaina was seeking a scholar named Velgar Vera (neutral good, orc \textbf{mage}). Velgar is a friend to Mordenkainen and has visited Mordenkainen's home in Avernus. Velgar wasn't in his shop earlier today, so Malaina lingered in the Market Ward, waiting for the orc to return.

\subparagraph{Following the Lead}
If the characters go to Velgar's Vittles, Velgar reports that Malaina met with him briefly a few hours ago to ask about Mordenkainen. Velgar tells the characters the same thing he told Malaina: the shopkeeper spoke with Mordenkainen last week, and the famous wizard mentioned nothing about Alustriel, Tasha, or Vecna, or anything about recently being in Sigil.

By the time the characters return to the sanctum, it's too late to stop Kas. Proceed to the "Kas Reveals Himself" section.

\subsubsection{Fighting Kas}
It's unlikely, but one or more characters might be present when Kas reveals his true identity. In this case, use the stat block presented in the "\textbf{Kas the Betrayer}" section in appendix B. Kas focuses on stealing the rod and using the portal to travel to Pandesmos. If the characters defeat Kas, he reveals the location of Vecna's ritual: the Cave of Shattered Reflection in Pandesmos on the plane of Pandemonium. Kas also reveals he planned to use the rod to free \textbf{Miska the Wolf-Spider} and ally with the demon lord to overtake Vecna.

It's up to the characters what happens to Kas next. Alustriel and Tasha are likely injured in the fight and therefore disinclined to join the characters in their pursuit of Vecna. However, Alustriel and Tasha urge the characters to take the \textit{Rod of Seven Parts} to aid their fight against the lich-god. At your discretion, Alustriel and Tasha might join the characters on the journey to Pandesmos, but they focus on neutralizing Kas's powerful remaining forces instead of joining the fight against Vecna. Alustriel or Tasha might join particularly persuasive characters in their pursuit of Vecna, but this will require significant effort on your part and might make the ensuing chapters significantly easier for the characters.

When the characters are ready, Alustriel can direct them to a portal in Sigil that leads to Pandesmos. It's up to you whether the door opens to Pandesmos Outlook (shown on map 9.1) or directly at the cave entrance. In the former case, Alustriel can provide directions to the cave from the outlook. If this happens, skip parts of the next chapter as appropriate. In the latter case, the characters gain an additional level and skip to chapter 11.

\section{Kas Reveals Himself}
If the characters return to their quarters to rest, Kas spends the next hour going through the motions of pretending to examine the \textit{Rod of Seven Parts} with Alustriel and Tasha. In reality, Kas is waiting for the characters to fall asleep or become distracted in their quarters. Once he assumes the characters are sleeping or oblivious, he reveals his true form.

Kas takes Alustriel and Tasha by surprise, severely wounding them, then steals the rod and steps through the sanctum's portal to Pandemonium, where the characters face him in the next chapter.

Once Kas escapes the sanctum, Malaina returns to find the immediate aftermath of Kas's betrayal. She rushes to alert the characters.

\subsection{Reality Dawns}
When Kas reveals himself, a cacophony erupts downstairs. Sleeping characters are awakened, and otherwise occupied characters hear the following:

\begin{DndReadAloud}{}
Violent slamming and terrible crashes tear through the silence. Something is happening downstairs. Deep, wicked-sounding laughter echoes through the sanctum, along with familiar voices shouting, "You monster!" and "We trusted you!"

\end{DndReadAloud}

When the characters open their door, they see Malaina in a panic. She shouts for the characters to come downstairs. Read or paraphrase the following:

\begin{DndReadAloud}{}
The sanctum looks like it was hit by a tornado: there are toppled bookshelves, overturned pieces of furniture, and books and papers strewn everywhere. Scorch marks mar the once-gleaming wood and plush carpets. Only a single sofa in the sanctum's center is upright, and on it sit Alustriel and Tasha. Both are sipping potions and look dazed.

\end{DndReadAloud}

Alustriel and Tasha are in a state of shock, both because the Mordenkainen they were working with is actually the evil vampire Kas, and because they didn't catch on to his plan. Alustriel and Tasha answer questions confusedly and in short phrases.

After a few minutes, the wizards regain their poise and tell the characters all the information in the beginning of this section. Alustriel and Tasha are willing to submit to magical means of determining the truth. They never deceived the characters, and they never wanted anything more than to stop Vecna's ritual.

\subsubsection{Chasing Kas}
Before Kas fled the sanctum, he muttered a phrase both Alustriel and Tasha overheard: "Finally, it ends in Pandesmos." The women repeat the phrase to the characters. Further, Alustriel has examined the sanctum's portal and knows Kas fled to the plane of Pandemonium. They surmise, correctly, that Vecna is weaving his ritual in Pandesmos, the first layer of Pandemonium.

Neither Alustriel nor Tasha had previously met Kas, but the stories about his and Vecna's hatred for each other are well known, as is Kas's hunger for power. Alustriel correctly supposes that Kas plans to find Vecna and co-opt the lich-god's ritual to become the most powerful being in existence. This outcome is as bad as Vecna ruling the multiverse, as far as Alustriel and Tasha are concerned, since Kas is just as evil as Vecna.

Because Kas took such pains to manipulate the characters into finding the \textit{Rod of Seven Parts}, Alustriel realizes the awful truth and shares it with the characters: Kas plans to use the rod to free \textbf{Miska the Wolf-Spider}, a powerful demon lord who was trapped in Pandesmos eons ago using the rod. Kas and Miska would then march against Vecna, co-opting the lich-god's ritual to begin their reign of terror across the multiverse.

\subsubsection{The Way Forward}
Alustriel and Tasha don't know the exact location of Vecna's ritual. The characters must travel to Pandesmos and confront Kas to learn the location so they can stop the Ritual of Remaking.

During their confrontation with Kas, Alustriel and Tasha were injured. Healing supplies in the sanctum ensure the wizards aren't in real danger. They insist the characters chase after Kas immediately to prevent him from freeing Miska. Alustriel and Tasha promise to help the characters however they can.

Alustriel and Tasha correctly suspect that Vecna will have woven abjurations to protect himself from any who would stop him. Vecna's Links are the magical tie the characters need to stop the lich-god. In fact, thanks to the links tying them to Vecna, the characters are the only ones in existence who can do so. They are the multiverse's only hope.

Before the characters leave the sanctum, Alustriel suggests the characters use the \textit{Chime of Exile} to return Kas to Tovag, his Domain of Dread. Alustriel notes that the characters need to significantly weaken Kas before the magic item will work on the vampire.

\subsubsection{Malaina's Mistrust}
Once the characters learn what has happened, they notice that Malaina remains particularly solemn.

\subparagraph{Malaina's Secret}
If the characters ask Malaina what's wrong, she admits she began suspecting something wasn't right shortly after Mordenkainen arrived in the sanctum. Malaina noticed that Mordenkainen muttered to himself in an unexpectedly cruel tone when he thought no one was around. He was also angrily defensive anytime Malaina casually asked how the mage fared. Malaina assumed Mordenkainen was merely stressed, so she didn't act on her suspicions until it was too late.

Regardless of the characters' reaction to this revelation, learning it counts as a secret for the purposes of the rules in the "Power of Secrets" section in this book's introduction.

\section{Pandesmos Outlook}
If the characters decide to chase after Kas, they can step through the portal whenever they choose. They emerge in Pandesmos Outlook, entering through the door shown on map 9.1.

\subsection{Pandemonium}
The rest of this chapter and the following two chapters take place on the plane of Pandemonium.

Pandemonium is an infinite plane of overwhelming chaos, a great mass of rock riddled with tunnels and vast caverns carved by howling winds. It is cold, noisy, and mostly dark.

Wind quickly extinguishes nonmagical open flames, such as torches and campfires. Spoken conversation is possible only by yelling, and then only to a maximum distance of 10 feet.

\subsection{General Features}
The features of Pandesmos Outlook are described in the following sections.

\subsubsection{Demons}
The chaos and violence brewing in the area attracts demons. Within 10 feet of the bottom of the cliff's lowest shelf are two degloths and two vlazoks (see appendix A for stat blocks for both). The exact placement of the demons is at your discretion. If they notice the characters, the demons attack. A demon that sees a fight break out rushes to join in the combat, hoping to make an easy kill.

Demons reduced to 20 hit points or fewer might divulge key information in exchange for mercy. See the "Tracking Kas" section for how these demons might help the characters locate Kas.

\subsubsection{Doorway}
The doorway to Sigil remains open through the end of this adventure. The characters can come and go as they wish. On the Pandesmos side, the doorway is invisible in the cliff.

\subsubsection{Elevation}
The characters emerge on a shelf of a massive cliff. This shelf is 150 feet above ground level. The shelf drops 50 feet to another shelf, which is 100 feet above ground level. The entire cliff is 180 feet tall.

\subsubsection{Geysers}
In the north section of Pandesmos Outlook are five geysers that periodically spew boiling water tainted with the plane's inherent chaos.

If the characters haven't already rolled initiative, have them roll initiative when a character first moves within 30 feet of a geyser. On initiative count 20 (losing initiative ties), one of the five geysers erupts; you choose which geyser. When a geyser erupts, all creatures within 5 feet of the geyser must make a DC 18 Dexterity saving throw, taking 44 (8d10) force damage on a failed save or half as much damage on a successful one.

\subsubsection{Lighting}
Strange multicolored lights flit through the caverns of Pandesmos, casting the area in dim light at all hours.

\subsection{Tracking Kas}
Kas fled toward Carapace Ridge, where Miska is imprisoned in the nearby Ruinous Citadel and where Kas has a stronghold in a cliffside redoubt. The vampire plans to marshal his forces before marching on the Cave of Shattered Reflection and confronting Vecna.

The characters need to travel northwest from Pandesmos Outlook to reach Carapace Ridge. There are several ways the characters might learn which direction to head. If the characters have trouble determining where to go, a character who succeeds on a DC 20 Wisdom (Perception) check can hear, amid the howling wind, sounds of battle coming from the northwest.

\subsubsection{Interrogating Demons}
If the characters defeat and interrogate the demons in this area, they learn that a powerful vampire made haste recently toward the northwest. The demons stayed out of the vampire's way when they saw the vampire dispatch several glabrezus.

\subsubsection{Investigating the Battlefield}
A character who investigates the battlefield and succeeds on a DC 16 Wisdom (Survival) check spots a trail of bootprints that leads to the northwest. This is Kas's path toward Carapace Ridge.

Characters who succeed on a DC 14 Intelligence (Investigation) check find the mangled bodies of three glabrezus. The bodies and a blood trail lead northwest. Kas cut these creatures down on his way to Carapace Ridge.

\section{Next Steps}
The characters arrive at Carapace Ridge toward the end of a battle between the demonic forces of Kas and those of the demon-god Lolth, who supports Vecna's cause. The characters must decide how best to navigate the battlefield and determine where Vecna is conducting his Ritual of Remaking in the next chapter.

\chapter{The War of Pandesmos}
The characters tracked Kas to Carapace Ridge on the first layer of Pandemonium. In this chapter, the characters must find Kas and learn where Vecna is weaving his Ritual of Remaking. Along the way, the characters might become embroiled in a war raging between Kas's allies and the forces of Lolth, the Spider Queen, who is allied with Vecna.

\section{Running This Chapter}
This chapter begins where the previous chapter ended, after the characters followed Kas's trail to Carapace Ridge, a low dip in a cliff that provides access to a beach below. The characters' search for Kas takes them to several locations, from the Ruinous Citadel, where \textbf{Miska the Wolf-Spider} is in the process of breaking free from his prison; to the Hurricane Tower, where the forces of Lolth marshal; and to the cliffside redoubt, where Kas retreated to plan. Along the way, they encounter Vaeve, a \textbf{drow mage} who is loyal to neither Lolth nor Kas, and who offers the characters useful information.

The characters' choices affect where they ultimately face Kas, but regardless, the characters must weaken the vampire enough to banish him with their \textit{Chime of Exile}. While doing so, they discover Vecna's location in the Cave of Shattered Reflection.

\subsection{Character Advancement}
The characters should be 19th level when this chapter begins. The characters gain a level after they learn the location of Vecna's Ritual of Remaking.

\subsection{Power of Secrets}
The characters can learn two secrets in this chapter applicable to the Power of Secrets rules found in this book's introduction:


\begin{description}
\item[Naxa's Secret.]
The \textbf{drow mage} Naxa and her sister came to retrieve a cloak. It isn't simply a fancy magic item—it's a key piece in an important ritual they're planning. The characters can learn this secret later in the chapter in the vault of the Ruinous Citadel (area Y4).
\item[Kas's Secret.]
Vecna is weaving his Ritual of Remaking deep underground in the Cave of Shattered Reflection. The characters learn this secret when they confront Kas later in this chapter.
\end{description}

\section{Chaos Incarnate}
When Kas arrived in Pandesmos, he used the \textit{Rod of Seven Parts} to force a hole in the sealed portal that blocks Miska from exiting the demiplane where he is trapped. The hole isn't wide enough, so Miska is expanding it so he can break free.

As Miska attempts to emerge, Fiends loyal to the demon lord have flocked to the area. They've been met by the forces of Lolth, Vecna's ally in his bid to remake the multiverse. In exchange for the protection supplied by her troops while Vecna weaves his ritual, the lich-god promised the Spider Queen a place at his side when he remakes the multiverse.

When the characters arrive in Pandesmos, they see a battle unfolding on Carapace Ridge. The forces loyal to Miska and Kas are fighting those loyal to Lolth and Vecna. The battle is at a stalemate and Kas cannot yet confront Vecna and co-opt the ritual. The characters have three options to stop Vecna and Kas: confront Kas in his cliffside redoubt, stop Miska from breaking free of his prison, or face Lolth's commander at Hurricane Tower.

\section{Carapace Ridge}
This battlefield, including the ridge where the characters arrive, is presented on map 10.1. Distance is difficult to judge on Pandemonium; the physical distance between the locations on map 10.1 is flexible but shouldn't be more than a mile or two.

Read the following to present the scene:

\begin{DndReadAloud}{}
You arrive on a clifftop towering hundreds of feet over a beach, beyond which boils a sea of kaleidoscopic clouds. On the beach, an ancient stone redoubt sits half-consumed by the sea. A tower juts from atop the cliff. Both fortifications are within sight, but screaming winds carry fog that makes judging distances difficult.

Corpses of many-legged horrors lie around you. The largest corpse is a massive spider with ichor oozing from holes in its carapace.

A war between thousands of demonic combatants rages throughout the ridge. The forces organized around the tower include driders and monstrous spiders, which bear banners featuring a spider with a woman's head. The forces surging from the redoubt are enormous spiders with the heads of wolves. The spider-wolves carry banners with snarling wolf heads.

\end{DndReadAloud}

Four spyder-fiend combatants of the battle linger near the characters: a \textbf{quavilithku spyder-fiend} and three kakkuu spyder-fiends (stat blocks for both appear in appendix A). They're busily hollowing out the corpse of a massive \textbf{citadel spider} (see appendix A) as a bunker for use if the battle shifts back this way. If they spot the characters, the spyder-fiends scuttle out of the citadel spider and skitter around its corpse to attack. The spyder-fiends retreat to the carapace if they start to lose the fight.

Both sides of the battle use banners to identify units and send messages. The forces of \textbf{Camlash}, Lolth's balor general, use a stylized symbol of Lolth on their banners, while Miska's forces use banners featuring wolf heads. The troops are so numerous and powerful on both sides that even 19th-level characters can't materially affect the war by joining in the fighting.

\subsection{Vaeve's Offer}
A neutral \textbf{drow mage} named Vaeve is lurking around the battlefield, initially out of sight. Vaeve spots the fighting and realizes that the characters are her best hope for allies. She reveals herself as the battle winds down, perhaps joining the fight against the kakkuus once the characters defeat the dangerous quavilithku.

\subsubsection{Talking with Vaeve}
Vaeve hopes the characters might help her, so she is eager to talk. She identifies both sides in the fight but insists she's not loyal to either side. Vaeve shares the following key points in her conversation with the characters:


\begin{description}
\item[An Explanation.]
The battle rages between forces loyal to Lolth, who is working with Vecna, and forces loyal to Miska, who is in the process of breaking free from his prison and is allied with a vampire. Vaeve and her identical twin sister, Naxa, are adventurers from Neverwinter who came here to retrieve potent magical components for their arcane rituals. They were using the chaos of battle to steal back a magic cloak—a family heirloom, Vaeve says—recently taken to Miska's treasure vault somewhere in the stone citadel on the beach below. But Vaeve and Naxa became separated in the fighting, and Vaeve worries that Naxa is brazen enough to attempt the retrieval on her own.
\item[Assistance Needed.]
Vaeve doesn't know why the characters are here, but she assumes they're planning to cross the battlefield for some reason. Vaeve asks the characters to escort Naxa from the citadel back to Carapace Ridge. There's significant treasure in the citadel, Vaeve says, and Vaeve would be grateful if the characters helped her sister.
\item[General's Presence.]
Lolth's general, a demon named \textbf{Camlash}, directs troops from the tottering tower. Lolth's forces include giant spiders and driders. The tower looks ready to tip over because it is; only the webbing inside the tower's base keeps it from falling over.
\item[Miska's Location.]
Miska is down in the citadel by the weird ocean of churning energy. Miska was recently freed from imprisonment, but he hasn't left the citadel yet—Vaeve doesn't know why. Miska's forces are varied spyder-fiends as well as flying demons that drip deadly blood.
\item[Vampire's Presence.]
Miska is following the commands of a vampire who occasionally appears on the battlefield to turn the tide of important fights against Lolth's forces. (Vaeve doesn't know the vampire's name, but her description should make it clear to the characters that this is Kas.)
\item[Vampire's Home.]
The vampire inhabits a chamber in the cliff, another mile past the tower. The upper entrance is under siege by Camlash's forces, and the lower entrance is guarded by spyder-fiends. Vaeve points out both entrances to the cliffside redoubt, two miles distant and barely visible through the blowing fog.
\end{description}

If the characters express interest in confronting Kas, Vaeve suggests first facing one or both generals. Defeating Camlash might compel Kas to emerge to seize a final victory over Lolth's forces. Defeating Miska might bring Kas out in a rage. Doing both, Vaeve suggests, could be most effective.

Before the characters leave, Vaeve reiterates her wish that the characters help Naxa. Vaeve insists that it's the least the characters can do in exchange for the useful information Vaeve shared. She refuses to go with them into the conflict.

\subsection{Avoiding the War}
Descriptions for the Ruinous Sea, Hurricane Tower, and the Ruinous Citadel are below—but you might not need them. High-level characters have options for travel to Kas's redoubt without wading through an active battlefield. However, no options are free of danger.

\subsubsection{Teleportation}
The characters can teleport directly to the redoubt's upper or lower entrances. The characters can't teleport into the redoubt directly, since it's warded with a \textit{Forbiddance} spell.

\subsubsection{Flight}
Camlash's citadel spiders (see appendix A) function as artillery, but the skies are predominantly under Miska's control thanks to powerful flying demons pledged to his service. If the characters fly around or above the battle, three hazvongels (see appendix A) soar down from the wispy cloud cover to attack. The hazvongels fight to the death. If the characters don't fly directly toward the cliffside redoubt after this fight, another trio of hazvongels attacks from a different direction. After this second attack, the demons steer clear of the flying characters.

\subsection{The Ruinous Sea}
Pandesmos has no oceans as they're known on mortal worlds, but it contains an enormous, swirling sea of multicolored energy known as the Ruinous Sea. The Ruinous Sea's chaotic energy devours creatures and land alike. At its edges, the Ruinous Sea laps against gritty beaches beneath a steep cliff, gradually eroding the land.

Winds shriek from the Ruinous Sea, joining the gales that scour Pandesmos. Rumors claim that the chaos of Limbo leaks into Pandemonium through the Ruinous Sea; others say that the Elemental Chaos touches Pandemonium in this roiling expanse. Either way, the Ruinous Sea is inimical to all but the strange and powerful chaos krakens that inhabit it.

\subsubsection{Swimming the Ruinous Sea}
Despite their ephemeral appearance, the Ruinous Sea's roiling clouds are viscous and cloying, acting as a liquid for the purposes of movement. A creature that enters the Ruinous Sea for the first time on a turn or starts its turn there takes 70 (20d6) damage. The type of damage fluctuates with the roiling of the Ruinous Sea (choose or roll a d4): 1, acid; 2, lightning; 3, necrotic; or 4, radiant.

In addition to that danger, the Ruinous Sea is also inhabited by evil chaos krakens (a chaos kraken uses the \textbf{kraken} stat block but is immune to the Ruinous Sea's damage). A chaos kraken named Vashishax lurks in the Ruinous Sea near the beach. Vashishax considers this war among outsiders to be irrelevant, though the kraken is spying on the battle out of curiosity. Vashishax considers any intrusion into the Ruinous Sea to be a personal insult and fights any creature that spends more than a minute within it. Vashishax isn't willing to die here; if reduced to fewer than 100 hit points, the kraken flees to its lair many miles deeper in the Ruinous Sea.

\section{Hurricane Tower}
The 200-foot-tall Hurricane Tower juts from the top of the cliff above the citadel like a broken finger, threatening to fall over the edge of the 400-foot-tall cliff. Thick spider webs woven throughout the base of the tower keep it anchored to the cliff. Hurricane Tower is made of stone, but it is an unsteady structure; it's also the only easily accessible shelter within miles of the characters' initial position in Pandesmos.

Hurricane Tower is the strategic heart of Lolth's forces. The balor general \textbf{Camlash} uses the tower to organize troops and direct their efforts in the war, so she ordered her spiders to reinforce the structure with webbing. The tower remains standing unless the characters sabotage the six web anchors within it.

\subsection{Camlash}
Lolth isn't present in Pandesmos; she has given command of her legions to Camlash. The balor general specializes in leading campaigns with troops of loyal driders and mutated arachnids. For this battle, adherents of Lolth have joined Camlash's forces, despite the usual animosity between driders and people. All fight together in Lolth's name.

Camlash doesn't enjoy fighting in Pandesmos. The wind muffles shouted commands and shreds war-banners, and troops exposed to the wind for too long become slow and erratic.

Without the tower's protection, Camlash's forces are exposed to assault from Miska's flying demons and the effects of Pandemonium's winds. The tower remains stable for now, but the webs are a weak spot the characters and the opposing army could exploit.

\subsection{Reaching the Tower}
Lolth's forces surround the tower, but they are focused on the spyder-fiends below rather than on defense. Hundreds of driders, Humanoids, and demons surround the tower, but the characters don't need to fight these creatures if they create a plausible distraction, disguise, or ruse.

If the characters are careless or fail to fool Lolth's forces, they might face a squad of 1d4 driders, 1d6 giant spiders, and 1d4 assassins. Characters who defeat the squad can try again to infiltrate the tower.

\subsection{Hurricane Tower Features}
A number of features are common throughout Hurricane Tower.

\subsubsection{Ceilings}
The tower's interior floors are long gone. The ceiling in the vestibule (area X1) is 15 feet high, and the ceiling in the tower is 200 feet high.

\subsubsection{Lighting}
The tower's interior is dark. Area descriptions assume the characters have a light source or some other means of seeing in the dark.

\subsubsection{Masking Chittering}
The winds of Pandemonium aren't audible within the tower, so they don't impede hearing. However, noisy chittering from a million tiny vermin echoes throughout. The chittering masks noise from combat, so tower denizens have disadvantage on Wisdom (Perception) checks to hear anything outside the area they're in.

\subsection{Tower Locations}
The following locations are keyed to map 10.2.

\subsubsection{X1: Vestibule}
\begin{DndReadAloud}{}
This wide room's ceiling is fifteen feet high. The walls it shares with the tower have pulled away, as though the tower were eager to slide over the edge of the cliff. A fifteen-foot ladder leads up to the tower's entrance. Four enormous cables of webbing, two on each side of the entrance, keep the tower anchored in place.

On the west side of this room, webs extend over a wooden altar with misshapen candelabras on it. Tiny spiders scuttle over the altar, and a drow man and a wood elf woman stand near it. Four bloated driders lounge around the room.

\end{DndReadAloud}

A drow \textbf{assassin} named Althein and a wood elf \textbf{mage} named Vendrasha, both faithful of Lolth, pray to the Spider Queen at the altar. Of the apparent four driders, three driders are what they appear to be; the fourth is a \textbf{phisarazu spyder-fiend} (see appendix A) named Rachazz, who is masquerading as a drider using its Change Shape ability.

The three driders squad leaders are resting away from Pandemonium's winds here. Rachazz infiltrated the tower and is collecting information for Miska about Camlash's forces.

When intruders enter the vestibule, the three driders rush into melee while Althein and Vendrasha hang back to support the driders. The driders and elves fight to the death.

Rachazz doesn't participate in the fight; it feigns injury and stays out of the way. If the characters look like they're winning the battle, Rachazz tries to slip out of the tower and report back to the Ruinous Citadel about the characters' prowess. The characters might later face Rachazz in or near the citadel. If the characters struggle to defeat the drow and driders, Rachazz joins the fight against the intruders to protect its disguise.

\subparagraph{Cutting the Webs}
Four web anchors support the walls: two in the southwest and two in the southeast. The webs are tough but, except where they meet the walls, aren't sticky. Each web anchor has AC 12; 60 hit points; and immunity to bludgeoning, fire, poison, and psychic damage.

The tower creaks ominously if these four anchors are cut, but it doesn't fall unless the two anchors in area X2 are also cut.

\subparagraph{Treasure}
The eight web-crusted candelabras are covered in precious gems. They resemble yochlol demons, and each is worth 45,000 gp.

\subsubsection{X2: Tower Base}
\begin{DndReadAloud}{}
The tilting tower's interior is hollow, with broken beams showing where upper floors existed long ago. Webs cover the tower's walls, with two anchor points to the northwest and northeast.

A monstrous demon is moving scraps of cloth atop a table made of crumbling stone. The demon is covered with countless tiny, crawling spiders that appear and disappear at random.

\end{DndReadAloud}

\textbf{Camlash} (see the accompanying stat block) plans troop movements on the table when the characters arrive. Hundreds of tiny spiders constantly scuttle over and around the demon.

Camlash doesn't consider intruders much of a threat and gloats about how foolish the characters are to meddle in Lolth's affairs. After taunting the characters, Camlash fights to the death.

\subparagraph{Causing a Rout}
Lolth's forces receive regular instructions from Camlash. If the characters eliminate the general but otherwise keep their presence in Hurricane Tower secret, a squad of four driders checks on Camlash in 10 minutes, and a squad of six driders arrives a few minutes later.

Once news of Camlash's defeat becomes widely known, no one else arrives. Instead, Lolth's followers throw down their banners and flee. Their disorganized retreat spreads across Pandesmos. It takes an hour for all of Lolth's forces to leave the battlefield.

\subparagraph{Cutting the Webs}
The two web anchors have AC 12; 40 hit points; and immunity to bludgeoning, fire, poison, and psychic damage. The webs are tough, but they're sticky only where they meet the walls. Camlash has enchanted these anchors to inflict a mental backlash on anyone who tampers with them. A creature that deals damage to a web anchor takes 18 (4d8) psychic damage.

The tower creaks ominously if both anchors are cut. If the four anchors in area X1 have also been cut, the tower falls.

\subsection{The Tower Falls}
If all six web anchors are cut or destroyed, Hurricane Tower crashes 400 feet down the cliffside, leaving only its foundation behind. Once it starts falling, the characters have 1 minute to escape the tower. Characters who fail to do so take 70 (20d6) bludgeoning damage.

The tower's fall has a number of additional effects:


\begin{description}
\item[Destroyed Citadel Areas.]
If the characters cut or destroy the webbing, the falling tower smashes the creatures near the entrance to the Ruinous Citadel and obliterates most of the citadel's rooms (see the "Ruinous Citadel" section). Areas Y3–Y6 are destroyed, along with the creatures and objects in them, and the creatures in area Y1 are slain. Naxa in area Y4 manages to scamper along the ceiling into area Y2, though she hides there and won't reveal herself to the characters until the phisarazu spyder-fiends are dispatched. The \textit{Rod of Seven Parts} in area Y7 is loosened from the stones, though the characters still need to manually pull the rod free to re-imprison Miska.
\item[Forces Disperse.]
Creatures left standing on the tower's floor are in full view of Lolth's forces. The nearby troops, realizing their general is dead and their only fortification is destroyed, throw down their banners and retreat across Pandesmos.
\item[Spyder-Fiends Falter.]
The spyder-fiends can't capitalize on this sudden victory, since their raklupis battle-leaders (see appendix A) immediately devolve into infighting to seize control. The spyder-fiend battle lines break into violent packs.
\item[The Vampire Notices.]
Kas arrives at the tower's broken base within a few minutes to find who caused the collapse. If the characters are there, Kas confronts them as described in the "Fighting Kas" section later in this chapter.
\end{description}

\section{Ruinous Citadel}
Long ago, adherents of the Queen of Chaos, a powerful demon lord, built a citadel near the Ruinous Sea. The Wind Dukes of Aaqa later used what is now known as the \textit{Rod of Seven Parts} to imprison Miska within the citadel, eventually ending the terrible Dawn War between the primordials and the gods.

Since then, the Ruinous Sea has encroached on the citadel. It now stands on a beach at the edge of the Ruinous Sea. The citadel is a ruin, as the sea's chaotic energies hasten its erosion and will someday claim it entirely.

\subsection{Miska}
Kas returned the \textit{Rod of Seven Parts} to the citadel and thrust it into the sealed mystical portal that keeps Miska in his prison demiplane, creating an opening. However, the opening isn't wide enough for Miska to exit, so Miska is widening it. If he can break free, he can join the battle his forces are engaged in and revel in chaos once again. If the rod is removed or knocked aside, the seal snaps shut and binds Miska once more.

Miska isn't idle while trying to free himself. He commands the armies of spyder-fiends and other demons through his spyder-fiend battle leaders, the raklupises. Miska's other spyder-fiends (see appendix A) defend him and implement the demon lord's complex strategies.

Yet Miska's position is precarious. While he's a cunning general, he's unable to command his forces in person, and enforcing discipline is challenging; spyder-fiends, like many demons, find order and organization difficult to tolerate. Additionally, Kas is becoming increasingly frustrated with Miska's slow emergence; until Miska claws his way free and defeats Lolth's forces, Kas can't assault Vecna. And all the while, the portal to the demiplane Miska is trapped in is in danger of being swallowed by the Ruinous Sea.

\subsection{Features of the Ruinous Citadel}
The following features are common throughout the citadel.

\subsubsection{Ceilings}
Ceilings throughout the citadel are 15 feet high.

\subsubsection{Doors}
The citadel's doors are made of stone and noisily grind open and closed due to sand in the hinges. Each creature in a room beyond a door hears it open and expects visitors, unless the door's hinges are cleaned, lubricated, or silenced.

\subsubsection{Lighting}
The citadel is crumbling due to great age and proximity to the Ruinous Sea. Cracks in the ceiling and walls admit dim light throughout the citadel.

\subsubsection{Respite from the Winds}
The winds of Pandemonium are distantly audible within the citadel. They don't impede hearing within the citadel.

\subsection{Ruinous Citadel Locations}
The following locations are keyed to map 10.3. If the characters felled the Hurricane Tower, areas Y1 and Y3–Y6 are destroyed, as well as any creatures and objects in them.

\subsubsection{Y1: Drill Field}
\begin{DndReadAloud}{}
Several bloated spiders with wolf heads move back and forth on the black sand near the citadel's entrance.

\end{DndReadAloud}

Nine kakkuu spyder-fiends (see appendix A) drill on this beach while awaiting orders from the battle leaders in the citadel. They find marching in formation dull and are eager to rejoin the battle raging on the clifftop above. The kakkuus rush toward any opportunity for combat, so characters can draw them away with any ruse that promises bloodshed. A badly injured kakkuu flees the area, hoping to hide out elsewhere amid the chaos of war.

\subparagraph{Glowing Signal}
The kakkuus immediately return if they see the glyph glowing on the door to area Y2.

\subsubsection{Y2: Guard Hall}
The exterior door has an invisible glyph inscribed on it. Each time a creature other than a Fiend opens the door, the glyph glows for a split second, and the creature must make a DC 18 Wisdom saving throw, taking 21 (6d6) psychic damage on a failed save or half as much damage on a successful one. Casting \textit{Dispel Magic} on the door removes the glyph, allowing the door to be opened safely.

\begin{DndReadAloud}{}
Old carvings of monstrous, tentacled beings adorn the walls in this long room. More recently, wolf heads have been painted on each figure. Webs stretch from wall to wall and are particularly thick to the southeast, around an open stone door. Two enormous demons with crablike bodies and wolflike heads flank the door.

\end{DndReadAloud}

Two phisarazu spyder-fiends (see appendix A) flank the open door that leads to area Y3. The phisarazus don't leave this chamber, but they attack intruders on sight and fight to the death.

The spyder-fiends spun the webs in this room, making the entire room \textit{difficult terrain}. The webs don't burn, but a creature can use an action to cut away the webs in a 10-foot-square section of the room.

The old carvings depict chaos krakens from the Ruinous Sea, but a spyder-fiend recently used demonic ichor to draw wolf heads on all the krakens.

\subsubsection{Y3: Hall}
\begin{DndReadAloud}{}
This hallway is empty, and its walls are decorated with smeared drawings of wolf heads.

\end{DndReadAloud}

\subsubsection{Y4: Vault}
The door to this room bears faded markings. A character who examines the door and succeeds on a DC 18 Intelligence (Investigation) check realizes that a magic glyph was once inscribed on this door, but the glyph's energy has been carefully drained away. The door is now harmless.

\begin{DndReadAloud}{}
This long room holds piles of armor, weapons, and banners amid mounds of webbing. Tiny spiders scuttle across the mounds.

\end{DndReadAloud}

The spyder-fiends dump their spoils of war here. The room's only occupant is a neutral \textbf{drow mage} named Naxa, who is wearing a \textit{Cloak of Arachnida} she found here. The cloak allows Naxa to climb on walls and ceilings, and she is currently on the vault's ceiling.

\subparagraph{Meeting Naxa}
If the characters search the room or call out for Naxa, she reveals herself. Naxa is Vaeve's twin sister. She doesn't want to fight and is looking for a safe exit from the citadel.

As Vaeve feared, Naxa snuck into this vault to retrieve the \textit{Cloak of Arachnida} that was stolen from the sisters. The cloak has distinctive purple edging with starburst patterns. Naxa wants to leave, but even with the cloak's powers, she knows that powerful spyder-fiends can spot her. Naxa asks the characters to help her escape the citadel and return to her sister.

\subparagraph{Naxa's Secret}
When the characters talk at length with Naxa, they realize that the mage is unreasonably preoccupied with her cloak, absently touching it as she chats. If the characters ask Naxa about her behavior, she admits that the cloak isn't merely a fancy magic item with sentimental value. The cloak belonged to their great-grandfather, and Naxa and Vaeve hope to use it in a ritual to seal their home in Neverwinter against their family's enemies.

Regardless of the characters' reaction to this revelation, learning it counts as a secret for the purposes of the Power of Secrets rules. At your discretion, if the characters help Naxa return to her sister and don't know about the cloak's history, Vaeve tells them the full truth about the item.

\subparagraph{Piles of Spoils}
Amid the detritus are several bloodstained and tattered banners depicting Lolth's symbol of a spider with a woman's head. The mounds of webbing contain a few corpses of people and driders, savagely torn apart by spyder-fiends and tossed here.

If the characters want to disguise themselves to infiltrate Camlash's forces, this room is an excellent resource for drow accoutrements.

\subparagraph{Treasure}
In addition to the \textit{Cloak of Arachnida} that Naxa wears, this room contains four +1 Longsword, three +1 Dagger, a \textit{+2 Longbow}, three Potion of Poison, and a Potion of Stone Giant Strength. One banner bearing Lolth's symbol is bordered with a thick silk rope. This is a \textit{Rope of Climbing} that functions normally once unstitched from the banner. A mahogany chest the spyder-fiends took from a drow quartermaster contains 750 flawless black opals, each worth 1,000 gp.

\subsubsection{Y5: Planning Chamber}
\begin{DndReadAloud}{}
The north wall of this chamber has a three-dimensional map of the cliff that shows the tower above and the citadel below. Glowing dots are scattered across the map. Three monstrous spider demons, two human-sized and another ogre-sized, squat around the table.

\end{DndReadAloud}

A \textbf{raklupis spyder-fiend} named Jallizanx is here with two kakkuu spyder-fiends (for both stat blocks, see appendix A). The kakkuus are relaying orders to troops fighting elsewhere, and Jallizanx is deciding how to give directions the dim-witted kakkuus can be reasonably expected to convey. When they see the characters, the kakkuus rush into melee while Jallizanx hangs back to cast spells and hurl venom globes. If the kakkuus are defeated and Jallizanx is plainly losing the fight, Jallizanx teleports away and doesn't return. If the characters capture Jallizanx, the spyder-fiend tells the characters about the scrolls described in the "Treasure" section and offers this loot in exchange for its life.

\subparagraph{Reviewing the Map}
The map is a highly detailed model with magical glowing dots that represent troop locations. The map shows the terrain within several miles, including Carapace Ridge to the west and the cliffside redoubt to the east. Churning smoke on the map's lowest edge represents the Ruinous Sea. Characters who investigate the map learn the following:


\begin{description}
\item[Glowing Dots.]
The glowing dots represent groups of troops, but they aren't detailed enough to indicate specific numbers or types of foes. A careful look shows that brown and red dots are Miska's, while green and gray dots are Lolth's. A character who succeeds on a DC 18 Intelligence (Investigation) check realizes that brown and green dots indicate fresh troops, and red and gray dots indicate injured or flagging ones.
\item[Indication of Status.]
The distribution of dots shows the overall state of the war. If the characters haven't defeated either army's general, the forces are at a stalemate along the edge of the clifftop. If either general is defeated, that side is obviously at a disadvantage and faces impending loss.
\item[The Cliffside Redoubt.]
The cliffside redoubt is shown in cross-section on this model. The lower entrance, upper bunker, and central chamber are all visible. The lower entrance has a glowing red dot outside, and the upper bunker has a green dot outside. The central chamber depicts Kas's throne, which has a unique golden dot on it, representing the vampire warrior.
\item[Ruinous Citadel.]
The citadel is depicted without its roof, showing the configuration of rooms inside without any dots. A character who succeeds on a DC 15 Wisdom (Perception) check while examining this section sees faint Abyssal runes in each chamber. A character who understands Abyssal can read them as Drill Field, Guard Hall, Planning Chamber, Vault, and Miska's Prison Hall.
\item[Tower of Mystery.]
The model of Hurricane Tower is a solid box because the spyder-fiends don't have information about its layout. The webs holding the tower in place look significantly more durable in the model than in real life, since Jallizanx doesn't know how precarious the tower's anchors are.
\item[Violet Dot.]
A violet dot swims back and forth in the Ruinous Sea. Jallizanx knows there's a powerful creature out there but doesn't know what it is.
\end{description}

\subparagraph{Treasure}
A platinum scroll tube worth 2,500 gp is hidden in webbing in the southwest corner of the room, near the ceiling. It contains two Spell Scroll of \textit{Glyph of Warding}. Characters who search the room and succeed on a DC 15 Wisdom (Perception) check find the scroll tube.

\subsubsection{Y6: Food Storage}
The door to this room has an invisible glyph inscribed on it. Each time a creature other than a \textbf{raklupis spyder-fiend} opens the door, the glyph glows for a split second, and the creature must make a DC 18 Wisdom saving throw, taking 21 (6d6) psychic damage on a failed save or half as much damage on a successful one. Casting \textit{Dispel Magic} on the door removes the glyph, allowing the door to be opened safely.

\begin{DndReadAloud}{}
This room is filled with lumpy shapes bound in webs.

\end{DndReadAloud}

Dozens of driders, drow, and giant spiders are wrapped in these webs. Each is unconscious, has the incapacitated condition, and can't regain hit points until 8 hours after being freed from the web cocoons.

The spyder-fiends know that Miska will be voraciously hungry when freed, so they prepared this larder for him. Kakkuus kept sneaking in to steal food, so the raklupises changed the glyph on the door to keep other spyder-fiends away.

The webbed prisoners want to flee the citadel to rejoin Lolth's forces. They are indifferent toward characters who free them.

\subsubsection{Y7: Miska's Prison-Hall}
The door to this chamber has an invisible glyph inscribed upon it. Each time a creature other than a Fiend opens the door the glyph glows for a split second, and the creature must make a DC 18 Wisdom saving throw, taking 21 (6d6) psychic damage on a failed save or half as much damage on a successful one. Casting \textit{Dispel Magic} on the door removes the glyph, allowing the door to be opened safely.

\begin{DndReadAloud}{}
The walls and ceiling of this crumbling, ancient chamber are obscured by sheets of webbing, but the floor is mostly clear. The room's west end contains a jumble of stones stacked to create a vertical ring about five feet in diameter. A sleek staff—the \textit{Rod of Seven Parts}—is jammed between two encircling stones like a splinter, and two spider demons with wolf heads are hunkered protectively next to the circle.

Beyond the circle is a shimmering portal where a wall should be. Within the portal, a monstrous bulk of fur, chitin, and gnashing teeth presses against the stones from the other side.

The circle of stones shudders, and you realize the monstrous thing contained on the other side is expanding the circle, using the \textit{Rod of Seven Parts} like a wedge to lever the circle large enough to climb free.

\end{DndReadAloud}

A \textbf{raklupis spyder-fiend} named Uvonxu is here along with a \textbf{phisarazu spyder-fiend} bodyguard (for both stat blocks, see appendix A). Uvonxu coordinates with Miska on strategy, while the phisarazu takes notes to relay to the troops. Both immediately attack intruders. The phisarazu tries to stun enemies and then attacks in melee while Uvonxu hangs back to cast spells. These spyder-fiends fight to the death to defend Miska.

\subparagraph{Miska Intervenes}
Miska is still trapped within his prison, but the opening is sufficient for him to affect the fight. He can't move, can't be harmed, and can use only his legendary actions (see \textbf{Miska the Wolf-Spider} in appendix B), which originate from the opening. Miska uses the legendary actions to support the spyder-fiends in the fight, using his Howl to inflict psychic damage on the characters and Web to ensnare them.

\subparagraph{Demiplane Portal}
The circle of stones is the rim of a magical portal to the demiplane that has served as Miska's prison for eons. Miska is bound behind the portal and can't yet move through to Pandemonium. The stones can't be affected by mortal magic, as they are part of the prison that contains Miska. The only way to widen the portal is by using the power of the Rod of Seven Parts.

Kas jammed the \textit{Rod of Seven Parts} into the pinhole portal to Miska's prison to open it. The rod still bears some of Miska's essence from when it was used to imprison him long ago, so it is the only implement that can free him. Miska has widened the portal to 5 feet across, but he needs time to widen the portal enough to pass through it.

\subparagraph{Miska Freed}
If the characters return to Sigil without sealing Miska back in his prison, the demon lord emerges and is free to wreak havoc across the multiverse. The characters might encounter him if they return to Pandesmos, or you might decide Miska is present somewhere else in this adventure.

\subparagraph{Sealing Miska Away}
As an action, a character can try to pull the Rod of Seven Parts from the stone circle, doing so with a successful a DC 25 Athletics (Strength) check. If the Hurricane Tower collapsed on the citadel, the DC of the check is reduced to 10. On a successful check, the rod is removed and the portal snaps closed as Miska howls with rage. See the "Removing the Rod" section below for the other effects of doing so.

\subparagraph{Treasure}
The \textit{Rod of Seven Parts} is the greatest treasure here, but the spyder-fiends also left offerings to their general amid the webbing on the walls. The offerings include a Scimitar of Warning with engravings of a wolf's head on each side of the blade, a heavy platinum statue of Miska triumphantly raising his arms (worth 120,000 gp), and two huge gold fangs inset with enormous rubies (worth 35,000 gp each).

\begin{DndSidebar}[float=!b]{A Tougher Battle}
The battle in area Y7 isn't designed to pit Miska directly against the characters, as he's significantly more powerful than they are. However, if you think your players can handle a truly monumental challenge, you can use Miska's stat block in appendix B rather than the effects listed in the "Miska Intervenes" section. In this case, the portal is wide enough that Miska can reach through to use his actions, bonus actions, and legendary actions, but he can't yet move, so there's still time to seal him away by yanking the \textit{Rod of Seven Parts} free.



\end{DndSidebar}

\subsection{Removing the Rod}
Closing the portal by removing the \textit{Rod of Seven Parts} from the stones has the following consequences:


\begin{description}
\item[Spyder-Fiends Flee.]
All spyder-fiends flee the citadel, including Uvonxu (see area Y7), if she's still alive. Within a few minutes, the spyder-fiends in Pandesmos instinctively realize their general has been sealed away again. They desert their posts to engage in vicious infighting to determine who among them will lead. The spyder-fiend battle lines falter. The characters don't encounter any more raklupises in Pandesmos—the raklupises are too busy jockeying for power to risk a fight with the characters.
\item[The Vampire Notices.]
After about an hour, Kas realizes that something is wrong. He comes to Miska's prison-hall to find out what's going on, and he extends his search for the characters if Miska has been sealed or the \textit{Rod of Seven Parts} is gone. Characters who keep their presence hidden have enough time to take a short rest. If the characters do something more noticeable to alert Kas first—such as collapsing Hurricane Tower—Kas comes for a fight sooner. In either case, see the "Fighting Kas" section later in the chapter.
\end{description}

\section{Cliffside Redoubt}
Kas makes his lair inside the cliff below Carapace Ridge. This cliffside redoubt gives Kas a place to plan, organize excursions against the forces of Lolth, and brood about how long it's taking Miska to free himself.

Although the characters might encounter the vampire in his redoubt, they might also draw him out by changing the tides of war. Unless the characters draw him out by collapsing Hurricane Tower, sealing Miska away, or some other effective strategy, Kas waits in the redoubt for Miska to emerge.

\subsection{Redoubt Features}
The following features are common throughout the cliffside redoubt.

\subsubsection{Ceilings}
Ceilings within the redoubt are 12 feet high.

\subsubsection{Forbiddance}
The redoubt's interior is warded by a permanent \textit{Forbiddance} spell. The ward deals necrotic damage to Celestials and Fey who enter the redoubt.

\subsubsection{Lighting}
The redoubt's interior is dark. Area descriptions assume the characters have a light source or some other means of seeing in the dark.

\subsubsection{Respite from the Winds}
The winds of Pandemonium are distantly audible within the redoubt. They don't impede hearing within the redoubt.

\subsubsection{Ventilation Tubes}
Hidden ventilation tubes less than a half-inch in diameter draw fresh air from the cliffside into the redoubt. The holes keep out Pandemonium's whipping winds and allow Kas to come and go from within the redoubt in the mist form granted by his Change Shape bonus action. A character examining the cliffside must succeed on a DC 20 Intelligence (Investigation) check to spot the ventilation holes from the exterior; they're obvious in area F3.

\subsection{Cliffside Redoubt Locations}
The following locations are keyed to map 10.4.

\subsubsection{F1: Upper Bunker}
\begin{DndReadAloud}{}
A squat stone and metal bunker protrudes from the clifftop. Demon attackers are battering the stone walls and the thick metal door, but the door seems to be holding.

\end{DndReadAloud}

Lolth's forces batter at this bunker to reach Kas, who is lower in the redoubt. A \textbf{citadel spider} (see appendix A) attacks the structure while four driders direct its efforts. They expect retaliation from Miska's forces and are quick to attack unfamiliar creatures. If Lolth's forces were routed due to Camlash's defeat, the four driders are gone but the citadel spider remains to stubbornly bash at the bunker. It attacks anyone who comes near.

The dented metal door doesn't open with any amount of pulling or pushing, since it's anchored from the inside with metal bars driven into metal frames. The area beyond the door is included in the redoubt's \textit{Forbiddance} effect, so teleportation into the bunker fails. The door has AC 22, 75 hit points, and immunity to poison and psychic damage.

The bunker contains a narrow set of stairs descending 200 feet into Kas's chamber (area F3).

\subsubsection{F2: Lower Entrance}
\begin{DndReadAloud}{}
A cracked stone portico extends from a solid stone door set into the cliffside a few feet above the beach. Broken pillars around the portico are carved with demonic forms, blurred by scouring sand. Rubble is all that remains of the portico's collapsed roof.

\end{DndReadAloud}

A \textbf{raklupis spyder-fiend} named Kalzak and three phisarazu spyder-fiends (for both stat blocks, see appendix A) keep guard here. Kalzak uses \textit{Disguise Self} to maintain the form of a white-haired and bearded human ascetic leaning against a pillar. The phisarazus are disguised as crabs scuttling around the rubble, keeping watch while looking innocuous among the harmless crabs native to the beach.

If the characters reseal Miska in the Ruinous Citadel, Kalzak orders the phisarazus to remain on guard and then leaves to seize a higher position among the squabbling raklupises. One phisarazu stands in front of the door in its normal form while the other two continue their ruse of being mundane crabs.

Kalzak received this unfavorable posting due to his recent failed assault at Carapace Ridge. He's determined to follow Miska's order that no one enter Kas's redoubt. Kalzak engages in delaying conversation to give the phisarazus time to scuttle into position for an attack, uttering odd prophecies such as "When the winds blowing from chaos bring order, the lawmakers must falter."

The door opens with a noisy groan, revealing a small room and staircase leading 200 feet up the cliffside into the central chamber (area F3). Kas comes and goes from this entrance when he doesn't use the ventilation holes.

\subsubsection{F3: Central Chamber}
\begin{DndReadAloud}{}
This wide chamber has ascending stairs at one end and descending stairs at the other. The shrunken, bloodless corpses of four people and two driders surround a stone throne in the room's center. One wall bears a series of tiny holes, through which the wailing winds of Pandemonium are distantly audible.

\end{DndReadAloud}

Unless the characters encountered him elsewhere, \textbf{Kas the Betrayer} (see appendix B) broods here on his throne. He feasted on the corpses and drained their bodies of blood. See the "Fighting Kas" section below for information about running this confrontation.

The holes are ventilation tubes less than a half-inch wide that wind through the cliffside.

\subparagraph{Treasure}
A human corpse bears a platinum medallion of Lolth's symbol with dazzling diamonds at the end of each spider's foot (worth 3,000 gp). A drider corpse wears a bandolier woven of solidified shadow worth 30,000 gp. The bandolier holds two Potion of Supreme Healing and a \textit{Spell Scroll} of \textit{Seeming}.

\section{Fighting Kas}
The characters encounter \textbf{Kas the Betrayer} in one of three places: in area X2 of Hurricane Tower after toppling the structure, in area Y7 of the Ruinous Citadel after sealing Miska away, or in area F3 of the cliffside redoubt.

Kas spares only a few moments to belittle the characters before he draws his sword and attacks. Throughout the battle, he taunts them for being "Mordenkainen's good little puppets." If the characters had trouble recovering any piece of the \textit{Rod of Seven Parts}, Kas mocks them for those specific failures, saying they can't even dance on puppet strings properly.

Kas doesn't think he has anything to fear from the \textit{Chime of Exile}, and he's initially correct. The characters must reduce Kas to 50 hit points or fewer before they can banish him to Tovag, his Domain of Dread.

If the characters banish Kas back to Tovag, a fissure opens in the air behind him and pulls him through. Visible within the rift is a smoky, red-tinged skyline full of mountains.

\subparagraph{Kas's Secret}
As Kas is pulled back to Tovag, the characters' links to Vecna flare. The links' magic tears Kas's last great secret from his mind: Vecna is currently in the Cave of Shattered Reflection. If the characters are on the verge of killing Kas, the vampire tells the characters Vecna's location in an attempt to save his life. It's up to the characters what happens to Kas next.

Learning Vecna's location counts as a secret for the purposes of the Power of Secrets rules, and it allows the characters to move on to the next chapter of the adventure.

\section{Next Steps}
Once the characters learn the location of Vecna's ritual from Kas, they can continue to the final chapter of this adventure.

Depending on the characters' actions, the war between Lolth's forces and Miska's demons might be over. If the characters haven't defeated one or both generals, they can choose to remain in Pandesmos to do so before moving on to confront Vecna.

\chapter{Eve of Ruin}
To defeat Vecna and save the multiverse, the characters must reach the Cave of Shattered Reflection in Pandesmos and disrupt Vecna's ritual. Before that, the characters must navigate the demiplanes Vecna has already created and find a way to access the lich-god's ritual chamber.

\section{Running This Chapter}
The characters can return to the sanctum in Sigil to rest and confer with Alustriel and Tasha before they head to the Cave of Shattered Reflection. If the characters don't return to the sanctum, they find their way to the cave without trouble from their position in Pandesmos.

The majority of this chapter takes place in Vecna's Grasp, a small cave network where the characters must dismantle three demiplanar unrealities that hint at Vecna's plans to reshape the multiverse. Each demiplane contains an encounter or short exploration sequence. The characters can dismantle these demiplanes in any order, so read through all three—the "Torment of Kas," "Neverwinter's New King," and "Dead Gods" sections—before running this chapter.

Once the characters dismantle the three demiplanar unrealities, they can proceed through a new opening in Vecna's Grasp that leads to the Cave of Shattered Reflection.

\subsection{Character Advancement}
The characters must be 20th level when this chapter begins.

\subsection{Power of Secrets}
The characters can use the unspent secrets they've collected throughout this adventure in their final confrontation with Vecna. See the "Cave of Shattered Reflection" section later in this chapter for more information.

\section{Reunion in Sigil}
At the end of chapter 10, the characters learned that Vecna is performing his ritual at a site in Pandesmos called the Cave of Shattered Reflection. At some point after this revelation, the characters receive a message via a \textit{Sending} spell from Alustriel offering to provide more help once the characters are ready. If the characters don't plan to return to Sigil, encourage them to do so.

If the characters return to the sanctum in Sigil, Alustriel and Tasha convey the following information:


\begin{description}
\item[Demiplanes.]
Alustriel and Tasha know that Vecna's plans for the multiverse hinge on the potent secrets he has collected. Vecna has already used these secrets to create three demiplanes that are reachable from a small cave network in Pandesmos called Vecna's Grasp. To reach the Cave of Shattered Reflection, the characters must first destroy these demiplanes, which are the lich-god's early attempts to remake reality. The demiplanes are harbingers of what's to come if Vecna isn't stopped.
\item[Ending Unrealities.]
Each of Vecna's three demiplanes is an "unreality" constructed from a specific secret. The secret is represented by a unique item in the unreality called a "manifested secret." Alustriel and Tasha don't know the details of these manifested secrets, but they assume the items appear as important objects and bear the lich-god's symbol: a shriveled hand clutching an eye. To dismantle an unreality, a character must find its manifested secret, touch the item to focus the power of their Vecna's Link into it, and infuse the item with potent magic from an artifact such as the \textit{Rod of Seven Parts} or a high-level spell. (See the "Dismantling an Unreality" section for more information.)
\item[Teleportation Ward.]
A crystalline barrier surrounding the Cave of Shattered Reflection wards against teleportation magic. To reach the cave, the characters must first travel to Vecna's Grasp. Thanks to their supernatural connection to Vecna, the characters can teleport to Vecna's Grasp from anywhere in the multiverse.
\end{description}

\section{Vecna's Grasp}
Vecna's Grasp is a cave network in Pandesmos. Other than Vecna, only the characters can reach this cave network—and only because they carry Vecna's Links. No portal is needed to get them there; the characters need only think about entering Vecna's Grasp, and the power of their Vecna's Links will transport them directly to area E1. From Vecna's Grasp, the characters can reach Vecna's demiplanar unrealities and dismantle them, causing the large crystal at the center of Vecna's Grasp to disintegrate and leaving a pit in the floor that leads to the Cave of Crystal Reflection.

\subsection{Unrealities}
The unrealities connected to Vecna's Grasp are manifestations of Vecna's plans to remake the multiverse. These potent demiplanes exist outside time and space. Each unreality is dangerous, both in the immediate sense while the characters explore it and in the long term for the denizens of the multiverse.

The characters can reach an unreality by traveling down one of the three tunnels in Vecna's Grasp. Once the characters enter an unreality, the only way for them to leave is by finding and touching the demiplane's manifested secret, then injecting the manifested secret with enough magic to dismantle the unreality.

\subsubsection{Manifested Secrets}
Each unreality section notes the unreality's manifested secret. This is a secret that Vecna made foundational to his plan to reshape the part of the multiverse the unreality depicts. Each unreality's manifested secret is represented by an object that, although not "real" in the characters' reality, looks real in the demiplane and bears Vecna's symbol.

\subparagraph{Locating the Manifested Secret}
An unreality's manifested secret is an important item integral to the unreality. Whenever a character is in an unreality, the manifested secret's symbol of Vecna glows with purple light. The character can see this glow at any distance and through solid objects. The closer the character gets to the unreality's secret, the brighter the symbol glows. When the characters are near the manifested secret, they instinctually know how to dismantle that unreality.

\subsubsection{Dismantling an Unreality}
To dismantle an unreality, at least one character must touch the unreality's manifested secret. Meanwhile, the same or a different character must do one of the following:


\begin{itemize}
\item Expend 1 charge of the \textit{Rod of Seven Parts}
\item Expend a spell slot of 7th level or higher
\end{itemize}

As soon as the charge or spell slot is expended, the unreality disappears. The demiplane melts away, fades into nothingness, or shatters into countless fragments, making it clear it has been dismantled.

When a demiplane is dismantled, everything inside it is destroyed except for the characters, who are teleported to the mouth of the unreality's corresponding tunnel in Vecna's Grasp. Thereafter, that unreality's tunnel no longer reflects images of the unreality, and after 50 feet, the tunnel terminates at a featureless, obsidian wall.

After all three unrealities are dismantled, the crystals in area E1 shatter, clearing the way to the Cave of Shattered Reflection, which is described later in this chapter.

\subsection{Vecna's Grasp Features}
Areas in Vecna's Grasp have the following features.

\subsubsection{Ceilings}
The ceiling is 40 feet high in area E1 and 10 feet high in the tunnels leading away from that cavern.

\subsubsection{Lighting}
Areas E1 and E2 of Vecna's Grasp are unlit. Area descriptions assume the characters have a light source or some other means of seeing in the dark. Unrealities in E2 have bright light.

\subsubsection{No Teleportation}
As long as Vecna is conducting his ritual, teleportation magic doesn't function in this place or any of the connected demiplanes. There is an exception to this rule: thanks to their Vecna's Links, once the characters learn about the existence of Vecna's Grasp, each of them can use an action to transport themself to or from area E1, provided the location they're coming from or traveling to isn't one of Vecna's demiplanar unrealities or the Cave of Shattered Reflection.

\subsubsection{Walls and Floors}
Vecna's Grasp is composed of magically reinforced obsidian. Each 5-foot-square section of obsidian has AC 20, a damage threshold of 30, 60 hit points, and immunity to poison and psychic damage.

\subsection{Vecna's Grasp Locations}
These locations are keyed to map 11.1.

\subsubsection{E1: Kaleidoscopic Cavern}
Characters who use the magic of Vecna's Link to transport to Vecna's Grasp appear in this chamber. Read the following to set the scene:

\begin{DndReadAloud}{}
You appear before a thirty-foot-wide mass of translucent purple crystals embedded in the floor of a large obsidian cavern. Some of the crystal faces reflect distorted images of the cavern, while others flicker with scenes of the lich-god Vecna visiting destruction on distant worlds.

Three tunnels branch off the perimeter of the cavern. Obsidian crystal faces at the entrance to each tunnel show scenes similar to those on the central crystals.

\end{DndReadAloud}

Characters who examine the central crystals see that they seal an enormous hole in the cave floor like a cork. This hole leads down to the Cave of Shattered Reflection, but the characters can't bypass the crystals or make their descent until they explore the side tunnels (areas E2a–E2c) and destroy the demiplanes connected to them.

\subsubsection{E2: Unreality Tunnels}
Each of the three tunnels in Vecna's Grasp connects to a different demiplanar unreality. A tunnel's walls reflect vague, phantasmagoric scenes of the unreality at the tunnel's far end.

A creature can walk up to 50 feet down a tunnel and remain in Vecna's Grasp. To the creature, the tunnel appears to go on forever. After 50 feet, the creature passes an imperceptible threshold. Crossing that point causes the creature to leave the tunnel and appear in the tunnel's corresponding unreality. The link is one-way, so the creature can't leave the unreality and return to the tunnel.

\subparagraph{E2a}
Images of Kas the Betrayer flicker on the walls of this tunnel. In each reflection, Vecna torments Kas, who appears powerless against the evil god. Characters who cross the threshold appear in the unreality described in the "Torment of Kas" section.

\subparagraph{E2b}
Hazy scenes of civil unrest—peasants rioting, buildings burning, and soldiers enforcing martial law—reflect in the walls of this tunnel. A character who studies the images recognizes they take place in Neverwinter. Characters who cross the threshold appear in the unreality described in the "Neverwinter's New King" section.

\subparagraph{E2c}
This tunnel's walls flash violent scenes of Vecna slaying, dominating, and imprisoning other gods in the vast emptiness of the Astral Sea. A character proficient in the Religion skill recognizes these deities. Characters who cross the threshold appear in the unreality described in the "Dead Gods" section.

\section{Torment of Kas}
When the characters cross the threshold in area E2a, they appear in an alternate reality of Oerth, where Vecna has captured and imprisoned his archrival, \textbf{Kas the Betrayer}.

Read the following to describe the scene:

\begin{DndReadAloud}{}
Your feet sink into the mucky basin of a sprawling wasteland. The shattered remains of trees, roads, and whole mountain ranges dot the barren landscape. In the distance, a lone castle is silhouetted against the crimson horizon. Two hulking figures wander around the structure.

\end{DndReadAloud}

This unreality depicts Keoland, Kas's homeland on Oerth, as Vecna envisions it for his new multiverse. In this version, the entire nation has been transformed into an inhospitable hellscape. The ruined structure in the distance is all that remains of a castle Kas once called home and is the unreality's only noteworthy structure. The sun hangs along the horizon.

The unreality distorts distance such that characters who move toward the castle seem to get farther away from it. However, if all the characters move toward the hanging sun, the castle seems to move closer to them. Only by moving toward the sun can they reach the castle, which is a trek that takes 1 hour on foot.

If the players can't figure out how to reach the castle, have the characters make a DC 15 Intelligence (Arcana) check to recall that this unreality is a reflection of Vecna's hatred of Kas; as a vampire, Kas would loathe and fear the sun, so the sun is the key to witnessing Kas's torment.

When the characters reach the castle, read the following:

\begin{DndReadAloud}{}
Aside from a few antechambers, this once-impressive castle has been reduced to scorched rubble. Two hulking, bipedal creatures covered in spikes and smeared with blood stalk the castle's bailey.

\end{DndReadAloud}

\subsection{Castle Locations}
The following locations are keyed to map 11.2.

\subsubsection{K1: Ruined Bailey}
The hulking figures the characters saw from afar are two cadaver collectors (see appendix A). The monstrous war machines sift through the rubble of this ruined courtyard for remains. They attack creatures on sight, eager to add more corpses to their collections.

\subsubsection{K2: Statues of Vecna}
\begin{DndReadAloud}{}
Six statues of Vecna in lordly poses stand along the edges of this open-air walkway.

\end{DndReadAloud}

Each statue is a Medium object with AC 13, 20 hit points, and immunity to poison and psychic damage.

If a creature deals damage to a statue, each creature in the unreality must make a DC 15 Wisdom saving throw as psychic shock waves ripple throughout the unreality, taking 11 (2d10) psychic damage on a failed save or half as much damage on a successful one.

\subsubsection{K3: Receiving Room}
\begin{DndReadAloud}{}
An armored knight with red, glowing eyes stands on the dais at the far end of this vaulted chamber.

\end{DndReadAloud}

The armored figure is a \textbf{death knight} who stands guard in this ruined hall. It is indifferent toward creatures that enter this chamber and hostile toward creatures that try to leave area K4 or K5. As soon as it sees a creature leave area K4 or K5, the death knight moves to intercept and slay the escapees.

\subsubsection{K4: Fountain of Blood}
\begin{DndReadAloud}{}
Stagnant blood fills the basin of an old fountain at the center of this chamber. The tranquil surface seems to reflect something odd.

\end{DndReadAloud}

Vecna included this fountain of blood in his unreality to taunt Kas, who knows it's here but can't drink from it while chained in area K5. If the Kas in this unreality is destroyed, the fountain activates, causing another imagined Kas to emerge from the burbling liquid mere moments later. Kas is then supernaturally teleported into area K5, where he is shackled as described below. In Vecna's version of the multiverse, even oblivion offers his archrival no respite.

\subparagraph{Vision of Vecna}
If the characters look at the surface of the fountain's liquid, they see a glimpse of Vecna conducting his ritual in the Cave of Shattered Reflection. A character who examines the reflection sees Vecna kneeling in a crystal-lined chamber. The lich-god's eye is closed as he murmurs silent incantations. In his outstretched hand, Vecna holds a volatile, light-consuming orb of magical energy that resembles a black hole. The reflection in the fountain is illusionary; Vecna can't interact with any creature in this chamber, and vice versa.

\subsubsection{K5: Kas's Chamber}
\begin{DndReadAloud}{}
The unmistakable form of Kas sits hunched in the far corner of this dank, stone room. Each of the vampire's limbs is shackled to a ten-foot-long adamantine chain, with all four chains fastened to a large adamantine ball. Kas's attention is focused on the tattered tapestries and ruined fixtures adorning the chamber's crumbling walls. Slung on his back is a magnificent sword, the hilt of which is emblazoned with a glowing symbol of Vecna.

\end{DndReadAloud}

This is Vecna's imagined version of \textbf{Kas the Betrayer} (see appendix B). Unlike the real Kas, whom the characters encountered in the previous chapter, this imagined Kas is despondent and has no interest in fighting. Kas doesn't recognize the characters and is open to speaking with them, though he is bored by any topic of conversation that doesn't involve his own predicament. As long as he is shackled, Kas's speed is 5 feet, and he is unable to leave the room. The ball he's shackled to can't be moved, even by magical means.

This Kas doesn't know he's not real and can't comprehend the notion, but he understands that Vecna doomed him to this prison. He doesn't realize that this shattered keep was once his home on Oerth. Like the real Kas, this one bears an intense hatred toward Vecna, though in this form he is powerless against the archlich.

\subsection{Manifested Secret}
In Vecna's twisted unreality, Kas is shackled with his legendary sword, but he is forever unable to use it. The \textit{Sword of Kas} is this unreality's manifested secret.

Kas willingly lets a character take the \textit{Sword of Kas} if convinced the characters can destroy Vecna, better yet, if convinced that relinquishing his weapon will enable him to destroy Vecna someday. A character can make a DC 25 Charisma (Persuasion) check to try to convince Kas using either line of reasoning. If the characters fail to persuade him to relinquish the sword, Kas snarls and insults them, but he won't fight. If a character tells Kas that his prison is actually a shattered version of his home in Keoland, Kas howls in fury and gives over his sword immediately, asking the characters to take down Vecna.

The characters can dismantle this unreality by touching the sword and performing one of the actions detailed in the "Dismantling an Unreality" section earlier in this chapter.

\section{Neverwinter's New King}
When the characters cross the threshold in area E2b, they appear in an unreality where Vecna commands control over the city of Neverwinter in Faerûn.

In this unreality, Vecna turned Lord Neverember into a loyal death knight and iron-fisted tyrant. Neverember has turned his servants into wights and his opponents into dust.

Read the following when the characters arrive in the Neverwinter unreality:

\begin{DndReadAloud}{}
A regal figure clad in menacing spiked armor stands on the balcony of a castle ten feet above the city square. Four sets of gallows are lined up in the center of the square. The armored figure's features are shrouded in black shadow. Fitted on the figure's head is a striking golden crown, the front of which is emblazoned with a glowing symbol of Vecna.

Throughout the square, terrified citizens gawk at the king with tearful, disbelieving eyes. Beneath the balcony, twelve armored, desiccated figures stand guard and keep the crowd from advancing any closer.

\end{DndReadAloud}

Map 11.3 depicts this location.

If the characters make their presence known, Lord Neverember orders the twelve wights below the balcony to apprehend the characters and place them in the gallows. Lord Neverember doesn't join the fight unless the wights are destroyed or the characters attack him.

If a character tells the awestruck crowd that Lord Neverember is not the rightful ruler of Neverwinter and succeeds on a DC 15 Charisma (Persuasion) check, the mob breaks into a violent frenzy. In this case, the mob distracts the wights so the characters can focus their attention on Lord Neverember.

Lord Neverember (use the \textbf{death knight} stat block) is unwilling to part with his crown; he fights the characters to withhold the symbol of his rulership.

\subsection{Manifested Secret}
The secret that Lord Neverember is not the rightful ruler of Neverwinter is represented by Vecna's symbol emblazoned on Neverember's crown. The characters can dismantle this unreality by touching the crown and performing one of the actions detailed in the "Dismantling an Unreality" section earlier in this chapter. A character can touch the crown only while Lord Neverember has the incapacitated condition or after he has been destroyed. Magic can't lift the crown from Lord Neverember's head while he has at least 1 hit point.

\section{Dead Gods}
When the characters cross the threshold in area E2c, they appear in an unreality where Vecna has usurped the power of every other god in the multiverse and scattered the dead gods' bones across the Astral Sea.

Read aloud the following when the characters arrive:

\begin{DndReadAloud}{}
You float amid a vast void speckled with distant silvery lights. All around you are the floating remains of enormous stone corpses. Nearby, a monstrous behemoth resembling a horned, serpentine lobster drifts slowly through a patch of debris. The creature groans in agony, its one eye darting around in confusion. Vecna's unholy symbol glows faintly from within the behemoth's immense torso. Six spherical pests—each with a single eye and gaping maw—gnash at the behemoth's flanks.

\end{DndReadAloud}

The behemoth is a dying \textbf{astral dreadnought} that's being preyed on by six eye mongers (see appendix A for both stat blocks). The dreadnought has only 150 hit points remaining, it has no remaining uses of its Legendary Resistance trait, and its flying speed is reduced to 20 feet until its hit points are fully restored. The eye mongers focus their attacks on the astral dreadnought unless the characters stop them.

The astral dreadnought is called Arekanz. Arekanz has been consuming dead gods, but in this unreality, Vecna spitefully cursed the gods' remains to contaminate the dreadnought and induce its certain death. Arekanz is suffering the effects of the curse. It is so wracked with pain that it doesn't defend itself except to use its Donjon Visit legendary action.

A character who succeeds on a DC 15 Intelligence (Arcana) check recognizes Arekanz as an \textbf{astral dreadnought}, a creature that has a demiplane where its stomach should be. Everything the dreadnought swallows ends up there. The character also knows that if the dreadnought dies, the demiplane inside it will be destroyed, leaving everything in it to float aimlessly in the Astral Sea.

\subsection{Blue Feather of Habbakuk}
Within Arekanz's demiplanar donjon is a blue phoenix feather marked with the symbol of Vecna unreality's manifested secret. In this unreality, this single feather is all that remains of Habbakuk, a god on the world of Krynn and one of the many gods Vecna intends to destroy as part of his ritual. Any character with proficiency in the Religion skill recognizes the feather as the remains of Habbakuk. The characters must retrieve the blue phoenix feather so they can use it to dismantle this unreality.

\subsection{Arekanz's Donjon Locations}
In this unreality, before Vecna slew his rival deities and fed their remains to Arekanz, the archlich turned each rival to stone. From where creatures first appear in Arekanz's demiplanar donjon, the shattered remains of these deities stand between the party and the unreality's manifested secret to the east.

The following locations are keyed to map 11.4.

\subsubsection{G1: Crumbling Mound}
Characters who end up in Arekanz's demiplanar donjon arrive here.

\begin{DndReadAloud}{}
You stand among hills made from the half-digested bodies of countless deities. A deep valley separates two of these hills and stretches as far as can be seen in either direction. Under your feet, the ground shifts as mounds of scree continuously collapse into the valley.

\end{DndReadAloud}

The area of the hazardous scree is marked with dotted lines on the map. When any creature takes a step on the scree, the scree collapses. That creature and any other creatures within the marked area must succeed on a DC 13 Dexterity saving throw or fall down the cliff and onto the floor of the valley 100 feet below, taking 35 (10d6) bludgeoning damage.

\subsubsection{G2: Sword Bridge}
\begin{DndReadAloud}{}
The colossal, iron longsword of some legendary being was cast unceremoniously onto the ground here, creating an improvised crossing over the valley.

\end{DndReadAloud}

Three eye mongers (see appendix A) swallowed by Arekanz lurk below the sword. They attack any characters who cross the sword bridge.

\subsubsection{G3: Horror's Fissure}
\begin{DndReadAloud}{}
An wide fissure runs east to west between the sword bridge and a tall shrine.

\end{DndReadAloud}

The fissure is 30 feet deep. As the characters approach, it shudders as a \textbf{cosmic horror} (see appendix A) slithers out of it from the east to feast on the characters' minds and bodies.

\subsubsection{G4: Shrine of Habbakuk}
\begin{DndReadAloud}{}
The tall stone statue at the center of this shrine depicts a cobalt-blue biped with a humanlike body, avian features, and vestments decorated with phoenix iconography. Embedded in the statue's sternum is a blue phoenix feather emblazoned with the glowing symbol of Vecna.

\end{DndReadAloud}

The statue reflects Habbakuk, a god of animals and druids revered widely on the world of Krynn.

\subsection{Manifested Secret}
The characters can dismantle this unreality by touching the feather and performing one of the actions detailed in the "Dismantling an Unreality" section earlier in this chapter.

If the astral dreadnought dies while the characters are searching its demiplanar donjon for the manifested secret, the demiplane ceases to exist, and everything inside it appears around the dreadnought's corpse, including the characters and the feather. A character who spends 1 hour sifting through the ejected detritus finds the feather.

\section{Cave of Shattered Reflection}
Once the characters have dismantled all three demiplanes connected to Vecna's Grasp, the large crystals in area E1 shatter to dust, opening the way to Vecna's ritual chamber: a mythic cavern called the Cave of Shattered Reflection. In existence since the earliest days of the multiverse, the cavern has the ability to harness beings' life force and reveal fundamental truths to those inside. Here, Vecna twists the powers of the cave to fuel his ritual to unravel the multiverse and remake it to his specifications.

\subsection{Cave Features}
The Cave of Shattered Reflection has the following features, as shown on map 11.5.

\subsubsection{Crystal Walls}
Most of the walls are made of a translucent, purple crystal that is immune to all damage. Creatures on opposite sides of a crystal wall can see each other vaguely, provided they are both within 5 feet of the wall. Sound, including the sound made by the \textit{Chime of Exile}, can't pass through crystal walls.

\subsubsection{Diamond Doors}
Portals made of translucent, solid diamond are set into the crystal walls. Creatures on opposite sides of a diamond door can see each other vaguely, provided they are both within 5 feet of the door.

With the exception of the door in area R2, each side of each diamond door is set with a large ruby, the other with a large sapphire. (These are depicted on the map in red and blue, respectively.) Each side of the door is also etched with a sigil. This sigil matches the sigil on another door in the Cave of Shattered Reflection. Map 11.5 shows five pairs of identical sigils, labeled A1 and A2, B1 and B2, C1 and C2, D1 and D2, and E1 and E2.

\textbf{Teleporting between Sigils}. When Vecna or a creature carrying Vecna's Link touches a diamond door's gemstone, the doorway teleports the creature to the nearest unoccupied space next to the other door with a matching secret sigil.

The side of the destination doorway on which the creature appears depends on which gemstone the creature touched to teleport. If the touched gemstone was a ruby, the creature appears on the side of the destination doorway set with a ruby. If the touched gemstone was a sapphire, the creature appears on the side of the destination doorway set with a sapphire.

\subsubsection{Lighting}
The cave is bathed in eerie purple light, providing bright light.

\subsubsection{Limited Teleportation}
As long as Vecna is conducting his ritual, diamond doors are the only form of teleportation magic besides Vecna's Fell Rebuke reaction that functions in this place. The cave's crystal walls don't restrict where Vecna can teleport using Fell Rebuke, although he still must be able to see his destination.

\subsubsection{Empowered by Secrets}
In the Cave of Shattered Reflection, secrets are power. Depending on how many secrets the characters kept throughout this campaign, they might be more powerful in this cave than they would otherwise be.

Before running this part of the adventure, refer to the tracking sheet in appendix C. Tally the number of secrets the party has learned, and write the sum in the Total Secrets Learned box. Then add the number of secrets the characters traded for benefits, and write the sum in the Total Secrets Revealed box. Subtract the Total Secrets Revealed from Total Secrets Learned and write the result in the Secrets Kept Box. Consult the Power of Secrets table to determine what additional powers, if any, the characters have in the Cave of Shattered Reflection.

These powers are cumulative. For example, if the characters kept seven secrets, they start the encounter with \textit{inspiration} and also can see clearly through the cave's crystal walls.

\begin{DndTable}[header=Power of Secrets]{cX}

Secrets Kept & Additional Power\\
0–2 & —\\
3–6 & Each character starts the encounter with inspiration.\\
7–10 & Characters can see clearly through the cave's crystal walls.\\
11–14 & The first time each character enters the Cave of Shattered Reflection, the character benefits from the effects of the \textit{Haste} spell for 1 minute. The characters don't suffer from lethargy after the spell ends.\\
15–17 & Characters have advantage on any Constitution saving throws made in the cave.\\
\end{DndTable}

\subsection{Cave Locations}
These locations are keyed to map 11.5.

\subsubsection{R1: Shattered Grotto}
When the crystal mass in area E1 of Vecna's Grasp shatters in front the characters, read the following:

\begin{DndReadAloud}{}
The cave trembles as deep rents appear in the mass of purple crystals in the center of the room. With a tremendous crack, the crystals shatter, sending shards flying. The destruction of the crystals reveals a deep pit. Screaming winds and ghastly ethereal wisps of neon-green light rise from the pit.

\end{DndReadAloud}

The pit is 500 feet deep. Its sides are composed of obsidian near the top and purple crystal closer to the bottom. The howling winds and lights are expressions of Vecna's hoarded secrets. The pit is filled with wind strong enough to blow out torches. Attack rolls in the pit have disadvantage due to the wind.

Climbing down the pit's smooth walls without climbing gear or magic is perilous. A creature must succeed on a DC 20 Strength (Athletics) check to do so, falling to the bottom of the pit on a failed check.

The bottom of the pit is a flat, cylindrical chamber of smooth crystal. When the characters reach the bottom of the pit, two Undead monsters emerge from the pit's walls and attack the party. The monsters use the \textbf{death knight} stat block and represent potent secrets Vecna has hoarded.

One monster's hands and arms drip with blood, representing a king who secretly murdered his sibling. The other monster has no facial features, representing a famed individual who lied about their identity.

An opening on the east side of this area connects to area R2.

\subsubsection{R2: Threshold}
\begin{DndReadAloud}{}
A wall of translucent, purple crystal divides this obsidian cavern in two, blocking the way toward the tunnel on the other side. You feel a faint but deeply unpleasant hum coming from the tunnel.

Set in the middle of the crystal wall is a large, circular slab of lustrous white diamond. A glowing ruby is set into the diamond's center. It looks uncannily like an eyeball with a bloodred iris. An identical ruby is set into the opposite side of the diamond slab.

\end{DndReadAloud}

The first time a character comes within 5 feet of the crystal wall, a peculiar reflection of that character appears in the crystal. The reflection mirrors the character's appearance except for their face, which wears a twisted sneer. This is a \textbf{mirror shade} (see appendix A). It attacks the character once, then flees to area R3.

\subparagraph{Psychic Hum}
The tunnel east of the door emanates an unpleasant humming sensation into creatures' minds. Though painful, the hum doesn't deal damage. Once a character crosses the threshold to area R3, the hum intensifies in volume and begins to deal damage (see the "Amplified Hum" section of area R3 below).

\subparagraph{Diamond Door}
The diamond door in this hall connects the two halves of this room. Vecna or a character carrying Vecna's Link can teleport to the opposite side of the door simply by touching the ruby set into the door's diamond surface.

\subsubsection{R3: Ritual Chamber}
\begin{DndReadAloud}{}
Floor-to-ceiling walls of purple crystal divide this broad cavern and cast thousands of disorienting reflections throughout it. Large doors made of solid diamond are set into various walls. Each door has a large ruby or sapphire embedded in it.

As you enter the cave, the psychic hum intensifies, shaking you to the core. The hum is coming from somewhere in the crystal maze.

\end{DndReadAloud}

Four mirror shades (see appendix A) lurk within the crystal walls. (If the characters already destroyed the one in area R2, there are three mirror shades in this chamber instead of four.) The mirror shades are servants of Vecna tasked with slaying intruders. The mirror shades use their Mirror Movement trait to move through the crystal walls.

\subparagraph{Amplified Hum}
Any creature other than Vecna and his mirror shades that starts its turn here must succeed on a DC 20 Wisdom saving throw or take 5 (1d10) psychic damage.

\subparagraph{Door Maze}
The crystal walls divide the area into discrete chambers. To reach Vecna, the characters must use the diamond doors and figure out how to reach the central chamber (labeled "Vecna" on map 11.5).

Ten diamond doors are set into the crystal walls. Each door has a ruby on one side and a sapphire on the other, and each door is etched with a sigil, as described earlier in the "Diamond Doors" section. The doors, their gemstones, and the etched sigils—represented with labels—are marked on map 11.5.

Characters can teleport to Vecna's central chamber by touching the ruby gemstone on one of the following doors: A2, B1, C1, D1. The characters might need to use some trial and error to figure out which doors lead where. If the characters struggle, allow any character proficient in the Arcana or Investigation skill to figure out which doors lead to where the characters want to go.

\subparagraph{Confronting Vecna}
Read the following when the characters enter the central chamber where Vecna weaves his ritual:

\begin{DndReadAloud}{}
Levitating in the center of this crystal chamber is a skeletal person wearing kingly vestments: Vecna the Archlich. Vecna's desiccated head is thrown back in concentration. His wrinkled eyelids are closed, with one drooping over its empty socket. His shriveled lips are pulled back to reveal rotted teeth. Vecna spreads his arms, one with a missing hand, and murmurs in a profane-sounding gibberish while a lightless orb swirls in front of his chest. It grows bigger by the moment.

\end{DndReadAloud}

\textbf{Vecna the Archlich} (see appendix B) is weaving his Ritual of Remaking. He is too focused to pay the characters any mind unless they confront him directly. Vecna is funneling energy into the ritual in addition to the power of the secrets his cult gathered, so the lich is considerably less powerful than he normally is. Unlike the characters, Vecna can move through this area as if its doors were normal.

To achieve victory, the characters must reduce Vecna to 50 hit points or fewer. A character then must use the \textit{Chime of Exile} to target Vecna, which requires a clear line of sight to him. If a character wields the complete \textit{Rod of Seven Parts} in this encounter, that character feels the artifact's yearning to preserve the order of the multiverse and stop Vecna's ritual. Each time a character strikes Vecna with the rod, Vecna takes an extra 35 (10d6) psychic damage.

When the characters do this, read the following:

\begin{DndReadAloud}{}
The chime echoes off the crystal walls, drowning out the ambient psychic hum. Vecna, his ritual interrupted, clutches his head and screams in agony. Something in him snaps, and he falls limply to the ground. The cave becomes deathly silent. Time seems to stand still.

Then the crystal cavern shatters, overwhelming your senses. Vecna's body tumbles into an unfathomable inky void just before you also plunge into the black expanse. When you come to, you find yourself drifting among a sea of stars.

\end{DndReadAloud}

The \textit{Chime of Exile} sends Vecna hurtling toward his home world of Oerth. Moreover, the sudden cessation of Vecna's ritual causes a chain reaction that destroys the Cave of Shattered Reflection and Vecna's Grasp. Everything except Vecna is sucked through a portal to the Astral Sea, which is where the characters find themselves afterward. As if to punctuate the end of their quest, each character's Vecna's Link vanishes as their connection to the lich is severed.

\section{Conclusion}
With Vecna banished to Oerth and his ritual foiled, the characters have saved the multiverse from ruin. (If the characters fail, the multiverse becomes a terrible place indeed!)

Characters who return to the sanctum in Sigil are greeted by Alustriel and Tasha, who immediately begin planning a magnificent celebration in honor of the characters' heroics.

Though stymied, Vecna is far from undone. He might stubbornly attempt to destroy and re-create the multiverse again, or he might formulate a new scheme to bend the cosmos to his will. Either way, it will be some time before a new band of adventurers must rise to defeat the archlich once again.

\appendix
\chapter{Bestiary}
This appendix describes creatures that appear in the adventure, presenting them in alphabetical order. The \textit{introduction} of the \textit{Monster Manual} explains how to read a creature's stat block.


\begin{itemize}
\begin{multicols}{3}
\item \textbf{Blue Abishai}
\item \textbf{Green Abishai}
\item \textbf{Red Abishai}
\item \textbf{Astral Dreadnought}
\item \textbf{Black Rose Bearer}
\item \textbf{Blade Lieutenant}
\item \textbf{Blade Scout}
\item \textbf{Blazebear}
\item \textbf{Bone Roc}
\item \textbf{Borthak}
\item \textbf{Cadaver Collector}
\item \textbf{Citadel Spider}
\item \textbf{Cosmic Horror}
\item \textbf{Deadbark Dryad}
\item \textbf{Deathwolf}
\item \textbf{Degloth}
\item \textbf{Eye Monger}
\item \textbf{False Lich}
\item \textbf{Granite Juggernaut}
\item \textbf{Hazvongel}
\item \textbf{Hertilod}
\item \textbf{Inquisitor of the Tome}
\item \textbf{Adult Lunar Dragon}
\item \textbf{Mirror Shade}
\item \textbf{Moonlight Guardian}
\item \textbf{Necromancer Wizard}
\item \textbf{Night Scavver}
\item \textbf{Priest of Osybus}
\item \textbf{Relentless Impaler}
\item \textbf{Lonely Sorrowsworn}
\item \textbf{Lost Sorrowsworn}
\item \textbf{Spiderdragon}
\item \textbf{Kakkuu Spyder-Fiend}
\item \textbf{Phisarazu Spyder-Fiend}
\item \textbf{Quavilithku Spyder-Fiend}
\item \textbf{Raklupis Spyder-Fiend}
\item \textbf{Star Angler}
\item \textbf{Vlazok}
\item \textbf{Warforged Warrior}
\item \textbf{Whirling Chandelier}
\end{multicols}

\end{itemize}

\chapter{Character Dossier}
This appendix describes iconic characters the heroes might learn about during the adventure, presenting them in alphabetical order. Stat blocks are included as appropriate.

\section{Acererak}
Also known as the Devourer, Acererak is a powerful archlich feared across many worlds. He takes sadistic pleasure in killing adventurers by luring them into his trap-riddled tombs with the promise of powerful artifacts. Then he feeds on their souls to sustain his escapades across the multiverse.

\subsection{History}
Hundreds of years ago, Acererak was a mortal mage from Oerth devoted to the pursuit of power and immortality. Some say that Acererak was a pupil of Vecna, from whom he learned undeath's secrets.

Acererak travels the planes in search of powerful artifacts. When the archlich finds a particularly intriguing one, he seals it in one of his infamous tombs to tempt treasure hunters into the dungeon's depths. Acererak has filled these tombs with deadly traps, fearsome monsters, and twisted puzzles meant to torment and eventually slay would-be thieves, whose souls he then consumes.

The most famous of Acererak's tombs is the Tomb of Horrors, which has claimed many adventurers throughout its history. Another infamous tomb is the Tomb of the Nine Gods, in which Acererak sealed nine false gods he had slain there. He has additional tombs on Oerth, in Faerûn, and beyond.

Unlike many archliches, Acererak doesn't desire godhood. Nevertheless, his nefarious deeds have garnered him a substantial following. One such group of these followers founded the Bleak Academy, an institution of arcane and religious learning that extols Acererak's power.

\subsection{Personality}
Acererak delights in watching others perish. He regards himself as superior in all aspects, treating others as one would annoying insects. He views even the creatures he creates as mere tools to serve his whims. This hubris has often led him to disaster, but as a lich, Acererak continually returns to unlife.

\section{Alustriel Silverhand}
Lady \textbf{Alustriel Silverhand}, called the Shining Lady, has been an influential wizard and proponent of good across Toril for centuries. Alustriel is one of the Seven Sisters—immortal daughters of Mystra, a god of magic. The divine energy Mystra passed to Alustriel grants Alustriel incredible power over arcane magic.

Alustriel's youthful appearance as a human woman with silver hair gives no hint of her supernaturally extended life span. She typically wears long robes and wields a unicorn-headed staff, her \textit{Staff of Silverymoon}.

\subsection{Personality}
Alustriel's primary concerns are to spread kindness, reward virtue, and promote a culture of compassion throughout the multiverse. She is good at building alliances and quick to intervene when she senses a threat to the forces of good. She has traveled far and established safe houses across the planes—such as her sanctum in the city of Sigil. Alustriel doesn't seek personal glory or wealth; her style of influencing the cosmos is quiet yet steady.

\subsection{History}
Like other Chosen of Mystra, Alustriel is concerned with preserving the Weave, the primary incarnation of magic that permeates Toril. She believes that the Weave favors those who act with mercy and compassion, seeks to deliver lives of security for all, and bolsters people's efforts when they seek to right wrongs and combat evil.

Nowhere are Alustriel and her deeds better known than in the Silver Marches and its capital, Silverymoon. Alustriel ruled Silverymoon for centuries, once disguised as a wizard named Elué Dualen and then later in her true form. She helped create Silverymoon's famous Moonbridge and cofounded the Lady's College, the first school in Faerûn for mages as students rather than as apprentices in service.

Alustriel stepped down as Silverymoon's high mage more than a century ago. Her son, Methrammar Aerasumé, now leads the city and works to uphold his mother's legacy.

Alustriel has partaken in countless adventures before and after her tenure as Silverymoon's high mage. She has befriended famous adventurers such as Drizzt Do'Urden, worked with prominent organizations like the Harpers, and helped prevent or undo many kinds of evil.

\section{Kas the Betrayer}
The Betrayer, the Bloody Handed, the Destroyer—Kas has earned many epithets during his long existence. He is a vampire, legendary swordfighter, and ruthless warlord, and he is driven primarily by one thing: his hatred for Vecna.

Kas resembles a well-muscled, raven-haired, human man in his thirties, though he is far older. His fangs reveal his vampiric nature.

\subsection{Personality}
Kas is cruel, spiteful, and unrelenting in his pursuit of vengeance against Vecna. He readily lies, breaks promises, betrays allies, and taunts those who fall for his ruses. Kas has little use for those who won't validate his superiority or help advance his goals.

In this adventure, Kas fools Tasha and Alustriel, two incredibly powerful wizards. In addition to defeating Vecna, Kas desperately wants the wizards to witness and acknowledge his strength.

\section{Lolth the Spider Queen}
Lolth—the demon-god of deceit, shadows, and spiders—gains power by deceiving allies and enemies alike. Worshipers of the Spider Queen believe she can see through the eyes of all spiders and insist she is all-knowing.

\subsection{History}
In eons past, Lolth regularly plotted against her family of elven gods, seeking to kill those deities and subjugate their followers. Though her attempts mostly failed, she learned from them and grew mightier.

Her failed attempt on the life of the powerful elven god Corellon Larethian resulted in Lolth's banishment to the Abyss. There she conquered a layer known as the Demonweb, where she now makes her home. Lolth embraces her surroundings and takes the form of a demon, regularly appearing as a beautiful elven woman with the lower body of an enormous, monstrous spider. Lolth relishes her title of the Spider Queen and attempts to annihilate any who doubt her power.

Lolth became the patron god of evil monsters that dwell in Faerûn's Underdark. Some Humanoids also choose to worship her. Lolth conveys her favor in monstrous ways. One of the most accursed of the Spider Queen's gifts is transformation into a drider.

Followers of Lolth seek to foment chaos everywhere in the Spider Queen's name. The more ruthless and fanatically devoted to Lolth a worshiper is, the higher that worshiper ranks in Lolth's favor.

\subsection{Personality}
Lolth is a cruel god. She loves chaos and hurting the innocent, especially those who oppose her power-hungry ideals. The suffering of others delights her, and if that pain benefits her plans, all the better.

The Spider Queen employs duplicity and sadism against her enemies, but she also enjoys bedeviling those who claim to love her. She promises great rewards to those who follow her without question; whether she delivers on her vows depends on her whims.

The Spider Queen often allies with other gods or powerful demon lords, though she does so only for personal gain. At the first sign of any inconvenience, Lolth abandons her allies. Despite that fact, some believe her patronage is valuable.

Most recently, Lolth allied with Vecna. The archlich promised Lolth a place at his side as the second-most powerful being in existence if she aids his efforts to use the Ritual of Remaking to re-form the multiverse. Lolth has deployed armies to protect Vecna on Pandemonium, though if the tide turns against the archlich, Lolth plans to strand her forces.

\section{Lord Soth}
\textbf{Lord Soth} is the most powerful death knight on Krynn. Once a Solamnic Knight of the Order of the Rose, Soth was a paragon of virtue and justice who allowed his pride to lead him down an evil path. Soth perished during Krynn's Cataclysm, then rose from the ashes as a death knight. Soth lives in his cursed castle, Dargaard Keep.

\section{Miska the Wolf-Spider}
Miska the Wolf-Spider is a legendary demon lord and master of battlefield strategy. He has the lower body of a massive armored spider, four arms, and two enormous wolf heads that drip poison. Yet Miska's greatest strength is his cunning mind.

\subsection{History}
Ages ago, Miska led the hordes of Chaos against the forces of Law at the behest of his patron, the enigmatic Queen of Chaos. It seemed Miska's domination couldn't be stopped.

In desperation, Miska's opponents crafted an artifact to bind him in an extraplanar prison. This rod broke apart after sealing him in Pandemonium, scattering across the planes and becoming known as the \textit{Rod of Seven Parts}. The rod is the key to releasing Miska from his long imprisonment.

\section{Mordenkainen}
Mordenkainen is a human wizard from Oerth whose wanderlust and hunger for new magic take him across the multiverse. Mordenkainen is calculating and mysterious, but he isn't cruel.

Mordenkainen has been involved in several friendly rivalries, including with Tasha, another powerful spellcaster from Oerth. Many significant figures across the multiverse are aware of Mordenkainen, but few truly know him—something the vampire Kas exploits to impersonate Mordenkainen. Throughout the adventure, Kas is masquerading as Mordenkainen. The mage is unaware of the vampire warlord's actions.

\subsection{History}
Mordenkainen has told conflicting tales about his origins, but most scholars believe he accumulated impressive magical power at a young age. Known for his ingenuity, the wizard is responsible for writing several planar tomes and creating many spells that bear his name, including \textit{Mordenkainen's Faithful Hound} and \textit{Mordenkainen's Magnificent Mansion}.

As his fame on Oerth grew, Mordenkainen became the leader of a powerful group of spellcasters called the Circle of Eight. The group included several wizards who similarly pioneered new spells, including Tenser and Bigby. Eventually, the circle disbanded after it stopped an uprising organized by Vecna. That development fueled Mordenkainen's lifelong hatred of the lich.

Mordenkainen travels the multiverse while following his academic whims. The wizard owns a home called the Tower of Urm, which pops into existence in Avernus when Mordenkainen visits to study how the Nine Hells affect various schools of magic.

\subsection{Personality}
Mordenkainen is stubborn and egotistical to the point of driving even his friends to annoyance. He has few close relationships, although nearly all wizards recognize his accomplishments. He is lonely and dreams of building a community of spellcasters who support each other. Unfortunately, his ego makes that dream all but impossible.

Although Mordenkainen can be brash and off-putting, he eschews cruelty and sees no reason to cause pointless suffering. The wizard doesn't like to admit it, but he would put his own safety at risk to quell the pain of innocents.

\section{Strahd von Zarovich}
Strahd von Zarovich is the Darklord of Barovia, a Domain of Dread. Little happens there without the Darklord's knowledge, although Strahd rarely pays attention to what he considers the uninteresting dealings of lesser beings. However, once the characters arrive in Barovia, explore a place called Death House, and find a piece of the \textit{Rod of Seven Parts} there, Strahd's interest becomes piqued, and he makes an appearance. The characters encounter Strahd in his Master of Death House iteration in chapter 5.

\subsection{History}
In life, Strahd von Zarovich was a prince, a soldier, and a conqueror. His thirst for power never sated, Strahd made a pact with the Dark Powers to become immortal. Meanwhile, Strahd's evil deepened, until in a jealous rage he murdered his brother, Sergei. Sergei's betrothed, Tatyana, leapt from a tower to escape Strahd and vanished into the Mists rising around Barovia as Strahd slew everyone else in the castle. He had become a vampire, and Barovia became a Domain of Dread.

Now the Dark Powers keep Strahd trapped in his realm, tormenting him with his inability to escape for all eternity. He spends his days amusing himself as best he can, terrorizing Barovia's people and savoring the fear and worship he commands.

\section{Tasha}
Tasha's path to greatness began when she was adopted by the arch-hag Baba Yaga, who named her Natasha. Tasha went on to create various spells, including \textit{Tasha's Hideous Laughter}, and her magic-fueled ambitions brought her into contact with demons and demon lords, which she subjugated and used against her enemies. On the Material Plane, she became known as \textbf{Iggwilv the Witch Queen} and wrote the \textit{Demonomicon of Iggwilv}, the greatest of all treatises on the Abyss and its demonic inhabitants. In recent years, Tasha has sequestered herself in the Feywild, achieving incredible power and slowly turning into a Fey creature. Tasha has become Zybilna, archfey of the domain of Prismeer.

\subsection{Answering the Summons}
When Zybilna received Alustriel Silverhand's summons to combat Vecna, the archfey was sorely needed in Prismeer. As a compromise, and to honor Tasha's friendship with Alustriel, Zybilna sent a version of herself from the past to Alustriel's side. The Tasha who appears in the adventure is a powerful wizard, though she is not yet a witch queen or an archfey.

\section{Tiamat}
The terrifying dragon god Tiamat is the progenitor of chromatic dragons across the multiverse. Her earliest background is shrouded in mystery, but Tiamat's grand, five-headed silhouette has become an iconic symbol of avarice and destruction.

\subsection{History}
Tiamat is one of two primordial dragons, the other being Bahamut, the progenitor of metallic dragons. Together Tiamat and Bahamut created what became known as the First World: the first iteration of physical reality, which was subsequently sundered into the innumerable realities that now make up the Material Plane. The first inhabitant of this First World was Sardior, a ruby-red dragon with jeweled scales they made in their likeness, who worked with Bahamut and Tiamat to create metallic and chromatic dragons. The Draconic poem "Elegy for the First World" describes how hordes of creatures overran this fledgling world, waging war on the dragons and colonizing the realm in the names of their gods. Bahamut and Tiamat were defeated, and Sardior disappeared.

Sometime after the destruction of the First World, Tiamat carved out a lair in Avernus, the first layer of the Nine Hells, where she commands a fearsome legion of followers. She maintains a courteous relationship with Asmodeus, the ruler of the Nine Hells, and commands draconic devils called abishais as her agents. Each type of abishai bears the coloring of one of Tiamat's five heads.

Across many worlds, Tiamat is worshiped as a god, with clerics at her disposal. Mortals who covet wealth and power pledge their lives to the dragon queen. Tiamat can send a fragment of her power to the Material Plane, which manifests as a titanic, five-headed dragon driven by rampaging greed.

\subsection{Personality}
Tiamat embodies the vices of evil dragons. She is vengeful and covets power and wealth above all else. However, Tiamat isn't reckless in her quests to expand her hoard. The dragon queen exhibits shrewd battle tactics and beguiling charm, easily swaying mortals. Each of Tiamat's five heads has its own voice and mannerisms, but they all share the same consciousness.

The relationship between Tiamat and Bahamut is complicated. Many texts portray the two as mortal enemies, but other histories of the First World depict the two as partners; some even speculate that she still feels great love for her creations on the Material Plane and grief over the loss of Sardior.

\section{Vecna}
On countless worlds, his name evokes tales of terror and cruelty: Vecna, the Undying King. Vecna, the Whispered One. Vecna, the Lord of the Rotted Tower. But Vecna had humble beginnings on the world of Oerth, where an order of wizards used him as a bootblack and scribe. He spent the better part of his childhood secretly educating himself in his masters' arts. Once Vecna learned all he could, he massacred the wizards. He then recorded his every foul thought and dream as he started to write his \textit{Book of Vile Darkness}.

Vecna forged a kingdom on Oerth, but he grew bored with it after several centuries. He started inflicting suffering on other worlds. In this adventure, Vecna has risen to godhood on Oerth, but he seeks to become the most powerful god in existence and bend the multiverse to his will. By the time the characters are involved, Vecna's master plan is almost complete. He has funneled a significant amount of his energy into weaving his ritual. Therefore, when the characters finally confront Vecna and try to save the multiverse, he is in his archlich form and not his divine form.

\chapter{Secrets Tracker}
Use this chart to keep track of the secrets characters learn during this adventure. Since it contains spoilers, keep this sheet hidden from the players!

When the characters learn a valuable secret, check the "Secret Learned" box in the appropriate row. If the characters use the Power of Secrets rules (see the introduction) to spend a secret they've learned, or if they reveal a secret they've learned to a nonplayer character, mark a corresponding "Secret Revealed" box. The party's total "Secrets Kept"—equal to the number of Secrets Learned minus Secrets Revealed—can impact the characters' fight against Vecna in chapter 11.

Downloadable PDF


\begin{DndTable}{Xc}

Valuable Secret & Source\\
Indrina knows Lord Neverember's claim to the throne isn't legitimate. & Chapter 1\\
Sarcelle recently had a vision about a desiccated man causing something terrible to happen. & Chapter 1\\
Umberto is a historian of Vecna. & Chapter 1\\
Mordenkainen, actually Kas in disguise, is duping everyone. & Chapter 2\\
Gertrude and her friend Rockzanna knew about an imminent attack by Lolth's cultists, but they provided no warnings. & Chapter 2\\
Figaro knew about the dangers of the area where his spelljamming ship crashed but deliberately hid this information from the captain. & Chapter 3\\
Ikasa and his best friend were stranded after a pirate attack. He knows about an additional survivor of the attack. & Chapter 3\\
Mercy has never truly searched for their best friend since the two were separated after the Day of Mourning. & Chapter 4\\
Kalyth lost wealth that could have prevented the current financial strain she and her allies are under. & Chapter 4\\
Sarusanda's father, Galias, joined the evil priests of Osybus. She expected to meet and slay him in Death House. & Chapter 5\\
Gazaia hid while evil soldiers attacked her tree and looted the fifth piece of the \textit{Rod of Seven Parts}. & Chapter 6\\
Valendar led an assault against his enemies without properly planning the mission or scouting the castle. & Chapter 6\\
Marian has held a lifelong fascination with the archlich Acererak and his evil magic, and she was once tempted to study necromancy. & Chapter 7\\
Rerak resents his imprisonment in the Tomb of Wayward Souls and never wanted to enact Acererak's will there. & Chapter 7\\
Malaina was suspicious of Mordenkainen from the moment he arrived, but she said nothing to Alustriel, Tasha, or the characters. & Chapter 9\\
The cloak Naxa and her sister came to retrieve isn't just a fancy magic item, but a key piece in an important ritual they're planning. & Chapter 10\\
Kas knows that Vecna is weaving his Ritual of Remaking deep underground in the Cave of Shattered Reflection in Pandesmos. & Chapter 10\\
\end{DndTable}

\chapter{Credits}

\begin{description}
\begin{multicols}{2}
\item[Lead Designer:]
Amanda Hamon

\item[Art Director:]
Kate Irwin

\item[Designers:]
Makenzie De Armas, Ron Lundeen, Patrick Renie

\item[Rules Developers:]
Jeremy Crawford, Makenzie De Armas, Ron Lundeen, Ben Petrisor, Carl Sibley

\item[Lead Editor:]
Eytan Bernstein

\item[Editors:]
Judy Bauer, Janica Carter, Michele Carter, Adrian Ng, Christopher Perkins

\item[Lead Graphic Designer:]
Trystan Falcone

\item[Graphic Designers:]
Matt Cole, Bob Jordan

\item[Cover Illustrators:]
Hydro74, Kieran Yanner

\item[Cartographers:]
Francesca Baerald, Dyson Logos

\item[Interior Illustrators:]
Lily Abdullina, Wade Acuff, Alfven Ato, Tom Babbey, Helge C. Balzer, Luca Bancone, Mark Behm, Olivier Bernard, Zoltan Boros, Bruce Brenneise, Ekaterina Burmak, Dawn Carlos, Diana Cearley, Sidharth Chaturvedi, Jedd Chevrier, Conceptopolis, CoupleOfKooks, Kent Davis, Nikki Dawes, Simon Dominic, Olga Drebas, Evyn Fong, Caroline Gariba, Lars Grant-West, Alexandre Honoré, Ralph Horsley, Sam Keiser, Julian Kok, Katerina Ladon, Daniel Landerman, Olly Lawson, Linda Lithen, Titus Lunter, Andrew Mar, Brynn Metheney, Robson Michel, Calder Moore, Martin Mottet, Scott Murphy, David Auden Nash, Jessica Nguyen, Irina Nordsol, One Pixel Brush, Claudio Pozas, Chris Rahn, Chris Seaman, Cynthia Sheppard, Craig J Spearing, Brian Valeza, Svetlin Velinov, Lauren Walsh, Shawn Wood, Zuzanna Wuzyk, Kieran Yanner

\item[Concept Art Director:]
Josh Herman, Kate Irwin

\item[Concept Artists:]
Carlo Arellano, Ekaterina Burmak, David Kegg, One Pixel Brush, Vinod Rams, Kieran Yanner

\item[Consultants:]
Rico Corazon, Ma'at Crook, Basheer Ghouse, James Mendez, Omar Ramadan-Santiago

\item[Project Engineer:]
Cynda Callaway

\item[Imaging Technicians:]
Daniel Corona, Kevin Yee

\item[Prepress Specialist:]
Jefferson Dunlap

\item[D\&D Studio Executive Producer:]
Kyle Brink

\item[Game Architects:]
Jeremy Crawford, Christopher Perkins

\item[Design Department:]
Justice Ramin Arman, Makenzie De Armas, Dan Dillon, Amanda Hamon, Ron Lundeen, Ben Petrisor, Patrick Renie, F. Wesley Schneider, Carl Sibley, Jason Tondro, James Wyatt

\item[Studio Art Director:]
Josh Herman

\item[Art Manager:]
Rob Sather

\item[Art Department:]
Matt Cole, Trystan Falcone, Fury Galluzzi, Bree Heiss, Kate Irwin, Bob Jordan, Noor Rahman, Emi Tanji, Trish Yochum

\item[Managing Editor:]
Judy Bauer

\item[Editorial Department:]
Eytan Bernstein, Janica Carter, Adrian Ng

\item[Senior Producer:]
Dan Tovar

\item[Producers:]
Bill Benham, Siera Bruggeman, Robert Hawkey, Andy Smith

\item[Director of Studio Operations:]
Wendy Despain

\item[Director of Product Management:]
Liz Schuh

\item[Product Managers:]
Natalie Egan, Teresa Kramer, Chris Lindsay, Hilary Ross

\item Special thanks to T. Alexander Stangroom and to the playtesters whose efforts made this a better book.
\end{multicols}

\end{description}

\chapter*{Printable Statblocks}
\setlist[description]{nolistsep,listparindent=3pt,topsep=6pt,font={\bfseries\sffamily\small}}
\pagestyle{empty}

\begin{DndMonster}[float*=b,width=\textwidth + 8pt]{Adult Lunar Dragon}
\begin{multicols}{2}
\DndMonsterType{Huge dragon, neutral evil}
\DndMonsterBasics[
armor-class = {17 (natural armor)},
hit-points = {\DndDice{15d12 + 75}},
speed = {40 ft., burrow 20 ft., fly 80 ft.}
]
\DndMonsterAbilityScores[
 str = 23,
 dex = 12,
 con = 20,
 int = 10,
 wis = 13,
 cha = 15
]
\DndMonsterDetails[
saving-throws = {Con +10, Wis +6},
skills = {Perception +11, Stealth +11},
damage-immunities = {cold},
senses = {darkvision 240 ft., passive perception 21},
languages = {Draconic},
challenge = 13,
]
\DndMonsterAction{Legendary Resistance (2/Day)}
If the dragon fails a saving throw, it can choose to succeed instead.

\DndMonsterAction{Tunneler}
The dragon can burrow through solid rock at half its burrowing speed and leaves a 15-foot-diameter tunnel in its wake.

\DndMonsterAction{Unusual Nature}
The dragon doesn't require air.

\DndMonsterSection{Actions}
\DndMonsterAction{Multiattack}
The dragon makes one Bite attack and two Claw attacks.

\DndMonsterAction{Bite}
\textsl{Melee Weapon Attack:} 11 to hit, reach 10 ft., one target. \textsl{Hit:} 13 (2d6 + 6) piercing damage plus 3 (1d6) cold damage.

\DndMonsterAction{Claw}
\textsl{Melee Weapon Attack:} 11 to hit, reach 5 ft., one target. \textsl{Hit:} 13 (2d6 + 6) slashing damage.

\DndMonsterAction{Tail}
\textsl{Melee Weapon Attack:} 11 to hit, reach 15 ft., one target. \textsl{Hit:} 13 (2d6 + 6) bludgeoning damage.

\DndMonsterAction{Cold Breath (Recharge 5\textendash 6)}
The dragon exhales a blast of frost in a 60-foot cone. Each creature in the cone must make a DC 18 Constitution saving throw. On a failed save, the creature takes 36 (8d8) cold damage, and its speed is reduced to 0 until the end of its next turn. On a successful save, the creature takes half as much damage, and its speed isn't reduced.

\DndMonsterSection{Bonus Actions}
\DndMonsterAction{Phase (3/Day)}
The dragon becomes partially incorporeal for as long as it maintains concentration on the effect (as if concentration on a spell). While partially incorporeal, the dragon has resistance to bludgeoning, piercing, and slashing damage.

\DndMonsterSection{Legendary Actions}
Adult Lunar Dragon can take 2 legendary actions, choosing from the options below. Only one legendary action can be used at a time and only at the end of another creature's turn. Adult Lunar Dragon regains spent legendary actions at the start of its turn.

\begin{DndMonsterLegendaryActions}
\DndMonsterLegendaryAction{Tail Attack}{The dragon makes one Tail attack.
}
\DndMonsterLegendaryAction{Treacherous Ice}{Magical ice covers the ground in a 20-foot radius centered on a point the dragon can see within 120 feet of itself. The ice, which is difficult terrain for all creatures except lunar dragons, lasts for 10 minutes or until the dragon uses this legendary action again.
}
\end{DndMonsterLegendaryActions}
\end{multicols}
\end{DndMonster}



\begin{DndMonster}[float=t]{Alustriel Silverhand}
\DndMonsterType{Medium humanoid (human, wizard), chaotic good}
\DndMonsterBasics[
armor-class = {15 (18 with \textit{mage armor})},
hit-points = {\DndDice{32d8 + 128}},
speed = {30 ft.}
]
\DndMonsterAbilityScores[
 str = 12,
 dex = 20,
 con = 18,
 int = 24,
 wis = 23,
 cha = 22
]
\DndMonsterDetails[
saving-throws = {Con +11, Int +14, Wis +13},
skills = {Arcana +14, History +14, Insight +13, Religion +14},
damage-resistances = {radiant},
damage-immunities = {poison},
condition-immunities = {charmed, exhaustion, poisoned},
languages = {Common, Draconic, Elvish},
challenge = 21,
]
\DndMonsterAction{Ear of the Chosen}
Whenever a creature on the same plane of existence as Alustriel speaks Alustriel's name, Alustriel hears her name and the next nine words the speaker utters.

\DndMonsterAction{Legendary Resistance (3/Day)}
If Alustriel fails a saving throw, she can choose to succeed instead.

\DndMonsterAction{Special Equipment}
Alustriel carries a magic staff known as the Staff of Silverymoon. In the hands of anyone other than Alustriel, the Staff of Silverymoon is a Staff of Power.

\DndMonsterAction{Spellcasting}
Alustriel casts one of the following spells, using Intelligence as the spellcasting ability (spell save DC 22):
\begin{DndMonsterSpells}
  \DndInnateSpellLevel{Dancing Lights}
  \DndInnateSpellLevel[2]{Detect Thoughts}
  \DndInnateSpellLevel[1]{Telepathy}
\end{DndMonsterSpells}
\DndMonsterSection{Actions}
\DndMonsterAction{Multiattack}
Alustriel makes three Staff of Silverymoon attacks or two Reproving Ray attacks.

\DndMonsterAction{Staff of Silverymoon}
\textsl{Melee Weapon Attack:} 12 to hit, reach 5 ft., one target. \textsl{Hit:} 14 (2d8 + 5) bludgeoning damage plus 38 (7d10) radiant damage.

\DndMonsterAction{Reproving Ray}
\textsl{Ranged Spell Attack:} 14 to hit, range 120 ft., one target. \textsl{Hit:} 65 (9d12 + 7) force damage, and if the target is a creature, it must make a DC 22 Charisma saving throw. On a failed save, the target has the incapacitated condition until the start of Alustriel's next turn. On a successful save, the target's speed is reduced by 10 feet until the start of Alustriel's next turn.

\DndMonsterAction{Argent Blaze (Requires Silver Fire)}
Alustriel summons a 60- foot cone of silver fire. Each creature in that area must make a DC 22 Dexterity saving throw, taking 77 (14d10) radiant damage on a failed save or half as much damage on a successful one. Additionally, Alustriel or one creature of her choice within 60 feet of her then regains 10 (3d6) hit points.

\DndMonsterSection{Bonus Actions}
\DndMonsterAction{Silver Fire (2/Day)}
Brilliant silver fire harmlessly wreathes Alustriel and empowers her. The silver fire lasts for 1 hour or until she has the incapacitated condition or uses another bonus action to quench it. While wreathed in silver fire, Alustriel gains truesight within 30 feet and can use her Argent Blaze action. In addition, Alustriel is unaffected by magic that would ascertain her alignment, creature type, thoughts, or truthfulness.

\DndMonsterSection{Reactions}
\DndMonsterAction{Shining Counterspell}
Alustriel interrupts a creature she can see within 60 feet of herself that is casting a spell. If the spell is 5th level or lower, it fails and has no effect. If the spell is 6th level or higher, Alustriel makes an Intelligence check (DC 10 plus the spell's level). On a successful check, the spell fails and has no effect. Whatever the spell's level, the caster takes 11 (2d10) radiant damage if the spell fails.

\end{DndMonster}



\begin{DndMonster}[float*=b,width=\textwidth + 8pt]{Ancient Red Dragon}
\begin{multicols}{2}
\DndMonsterType{Gargantuan dragon, chaotic evil}
\DndMonsterBasics[
armor-class = {22 (natural armor)},
hit-points = {\DndDice{28d20 + 252}},
speed = {40 ft., climb 40 ft., fly 80 ft.}
]
\DndMonsterAbilityScores[
 str = 30,
 dex = 10,
 con = 29,
 int = 18,
 wis = 15,
 cha = 23
]
\DndMonsterDetails[
saving-throws = {Dex +7, Con +16, Wis +9, Cha +13},
skills = {Perception +16, Stealth +7},
damage-immunities = {fire},
senses = {blindsight 60 ft., darkvision 120 ft., passive perception 26},
languages = {Common, Draconic},
challenge = 24,
]
\DndMonsterAction{Legendary Resistance (3/Day)}
If the dragon fails a saving throw, it can choose to succeed instead.

\DndMonsterSection{Actions}
\DndMonsterAction{Multiattack}
The dragon can use its Frightful Presence. It then makes three attacks: one with its bite and two with its claws.

\DndMonsterAction{Bite}
\textsl{Melee Weapon Attack:} 17 to hit, reach 15 ft., one target. \textsl{Hit:} 21 (2d10 + 10) piercing damage plus 14 (4d6) fire damage.

\DndMonsterAction{Claw}
\textsl{Melee Weapon Attack:} 17 to hit, reach 10 ft., one target. \textsl{Hit:} 17 (2d6 + 10) slashing damage.

\DndMonsterAction{Tail}
\textsl{Melee Weapon Attack:} 17 to hit, reach 20 ft., one target. \textsl{Hit:} 19 (2d8 + 10) bludgeoning damage.

\DndMonsterAction{Frightful Presence}
Each creature of the dragon's choice that is within 120 feet of the dragon and aware of it must succeed on a DC 21 Wisdom saving throw or become frightened for 1 minute. A creature can repeat the saving throw at the end of each of its turns, ending the effect on itself on a success. If a creature's saving throw is successful or the effect ends for it, the creature is immune to the dragon's Frightful Presence for the next 24 hours.

\DndMonsterAction{Fire Breath (Recharge 5\textendash 6)}
The dragon exhales fire in a 90-foot cone. Each creature in that area must make a DC 24 Dexterity saving throw, taking 91 (26d6) fire damage on a failed save, or half as much damage on a successful one.

\DndMonsterSection{Legendary Actions}
Ancient Red Dragon can take 3 legendary actions, choosing from the options below. Only one legendary action can be used at a time and only at the end of another creature's turn. Ancient Red Dragon regains spent legendary actions at the start of its turn.

\begin{DndMonsterLegendaryActions}
\DndMonsterLegendaryAction{Detect}{The dragon makes a Wisdom (Perception) check.
}
\DndMonsterLegendaryAction{Tail Attack}{The dragon makes a tail attack.
}
\DndMonsterLegendaryAction{Wing Attack (Costs 2 Actions)}{The dragon beats its wings. Each creature within 15 feet of the dragon must succeed on a DC 25 Dexterity saving throw or take 17 (2d6 + 10) bludgeoning damage and be knocked prone. The dragon can then fly up to half its flying speed.
}
\end{DndMonsterLegendaryActions}
\end{multicols}
\end{DndMonster}



\begin{DndMonster}[float=t]{Animated Armor}
\DndMonsterType{Medium construct, unaligned}
\DndMonsterBasics[
armor-class = {18 (natural armor)},
hit-points = {\DndDice{6d8 + 6}},
speed = {25 ft.}
]
\DndMonsterAbilityScores[
 str = 14,
 dex = 11,
 con = 13,
 int = 1,
 wis = 3,
 cha = 1
]
\DndMonsterDetails[
damage-immunities = {poison, psychic},
condition-immunities = {blinded, charmed, deafened, exhaustion, frightened, paralyzed, petrified, poisoned},
senses = {blindsight 60 ft. (blind beyond this radius), passive perception 6},
challenge = 1,
]
\DndMonsterAction{Antimagic Susceptibility}
The armor is incapacitated while in the area of an \textit{antimagic field}. If targeted by \textit{dispel magic}, the armor must succeed on a Constitution saving throw against the caster's spell save DC or fall unconscious for 1 minute.

\DndMonsterAction{False Appearance}
While the armor remains motionless, it is indistinguishable from a normal suit of armor.

\DndMonsterSection{Actions}
\DndMonsterAction{Multiattack}
The armor makes two melee attacks.

\DndMonsterAction{Slam}
\textsl{Melee Weapon Attack:} 4 to hit, reach 5 ft., one target. \textsl{Hit:} 5 (1d6 + 2) bludgeoning damage.

\end{DndMonster}



\begin{DndMonster}[float=t]{Archmage}
\DndMonsterType{Medium humanoid (any race), any alignment}
\DndMonsterBasics[
armor-class = {12 (15 with \textit{mage armor})},
hit-points = {\DndDice{18d8 + 18}},
speed = {30 ft.}
]
\DndMonsterAbilityScores[
 str = 10,
 dex = 14,
 con = 12,
 int = 20,
 wis = 15,
 cha = 16
]
\DndMonsterDetails[
saving-throws = {Int +9, Wis +6},
skills = {Arcana +13, History +13},
damage-resistances = {damage from spells; bludgeoning, piercing, slashing (from stoneskin)},
languages = {any six languages},
challenge = 12,
]
\DndMonsterAction{Magic Resistance}
The archmage has advantage on saving throws against spells and other magical effects.

\DndMonsterAction{Spellcasting}
The archmage is an 18th-level spellcaster. Its spellcasting ability is Intelligence (spell save DC 17, 9 to hit with spell attacks). The archmage can cast \textit{disguise self} and \textit{invisibility} at will and has the following wizard spells prepared:
\begin{DndMonsterSpells}
  \DndInnateSpellLevel{disguise self}
  \DndMonsterSpellLevel{fire bolt}
   \DndMonsterSpellLevel[1][4]{detect magic}
   \DndMonsterSpellLevel[2][3]{detect thoughts}
   \DndMonsterSpellLevel[3][3]{counterspell}
   \DndMonsterSpellLevel[4][3]{banishment}
   \DndMonsterSpellLevel[5][3]{cone of cold}
   \DndMonsterSpellLevel[6][1]{globe of invulnerability}
   \DndMonsterSpellLevel[7][1]{teleport}
   \DndMonsterSpellLevel[8][1]{mind blank}
\end{DndMonsterSpells}
\DndMonsterSection{Actions}
\DndMonsterAction{Dagger}
\textsl{Melee or Ranged Weapon Attack:} 6 to hit, reach 5 ft. or range 20/60 ft., one target. \textsl{Hit:} 4 (1d4 + 2) piercing damage.

\end{DndMonster}



\begin{DndMonster}[float=t]{Assassin}
\DndMonsterType{Medium humanoid (any race), any non-good alignment}
\DndMonsterBasics[
armor-class = {15 (studded leather armor)},
hit-points = {\DndDice{12d8 + 24}},
speed = {30 ft.}
]
\DndMonsterAbilityScores[
 str = 11,
 dex = 16,
 con = 14,
 int = 13,
 wis = 11,
 cha = 10
]
\DndMonsterDetails[
saving-throws = {Dex +6, Int +4},
skills = {Acrobatics +6, Deception +3, Perception +3, Stealth +9},
damage-resistances = {poison},
languages = {Thieves' cant plus any two languages},
challenge = 8,
]
\DndMonsterAction{Assassinate}
During its first turn, the assassin has advantage on attack rolls against any creature that hasn't taken a turn. Any hit the assassin scores against a Surprise creature is a critical hit.

\DndMonsterAction{Evasion}
If the assassin is subjected to an effect that allows it to make a Dexterity saving throw to take only half damage, the assassin instead takes no damage if it succeeds on the saving throw, and only half damage if it fails.

\DndMonsterAction{Sneak Attack (1/Turn)}
The assassin deals an extra 14 (4d6) damage when it hits a target with a weapon attack and has advantage on the attack roll, or when the target is within 5 feet of an ally of the assassin that isn't incapacitated and the assassin doesn't have disadvantage on the attack roll.

\DndMonsterSection{Actions}
\DndMonsterAction{Multiattack}
The assassin makes two shortsword attacks.

\DndMonsterAction{Shortsword}
\textsl{Melee Weapon Attack:} 6 to hit, reach 5 ft., one target. \textsl{Hit:} 6 (1d6 + 3) piercing damage, and the target must make a DC 15 Constitution saving throw, taking 24 (7d6) poison damage on a failed save, or half as much damage on a successful one.

\DndMonsterAction{Light Crossbow}
\textsl{Ranged Weapon Attack:} 6 to hit, range 80/320 ft., one target. \textsl{Hit:} 7 (1d8 + 3) piercing damage, and the target must make a DC 15 Constitution saving throw, taking 24 (7d6) poison damage on a failed save, or half as much damage on a successful one.

\end{DndMonster}



\begin{DndMonster}[float*=b,width=\textwidth + 8pt]{Astral Dreadnought}
\begin{multicols}{2}
\DndMonsterType{Gargantuan monstrosity (titan), unaligned}
\DndMonsterBasics[
armor-class = {20 (natural armor)},
hit-points = {\DndDice{17d20 + 119}},
speed = {15 ft., fly 80 ft. \textit{(hover)}}
]
\DndMonsterAbilityScores[
 str = 28,
 dex = 7,
 con = 25,
 int = 5,
 wis = 14,
 cha = 18
]
\DndMonsterDetails[
saving-throws = {Dex +5, Wis +9},
skills = {Perception +9},
damage-resistances = {bludgeoning, piercing, slashing from nonmagical attacks},
condition-immunities = {charmed, exhaustion, frightened, paralyzed, petrified, poisoned, prone, stunned},
senses = {darkvision 120 ft., passive perception 19},
challenge = 21,
]
\DndMonsterAction{Antimagic Cone}
The dreadnought's eye creates an area of antimagic, as in the \textit{antimagic field} spell, in a 150-foot cone. At the start of each of its turns, it decides which way the cone faces. The cone doesn't function while the eye is closed or while the dreadnought is blinded.

\DndMonsterAction{Astral Entity}
The dreadnought can't leave the Astral Plane, nor can it be banished or otherwise transported out of that plane.

\DndMonsterAction{Demiplanar Donjon}
Anything the dreadnought swallows is transported to a demiplane that can be entered by no other means except a \textit{wish} spell or the dreadnought's Bite and Donjon Visit. A creature can leave the demiplane only by using magic that enables planar travel, such as the \textit{plane shift} spell. The demiplane resembles a stone cave roughly 1,000 feet in diameter with a ceiling 100 feet high. Like a stomach, it contains the remains of past meals. The dreadnought can't be harmed from within the demiplane. If the dreadnought dies, the demiplane disappears, and everything inside it appears around the dreadnought's corpse. The demiplane is otherwise indestructible.

\DndMonsterAction{Legendary Resistance (3/Day)}
If the dreadnought fails a saving throw, it can choose to succeed instead.

\DndMonsterAction{Sever Silver Cord}
If the dreadnought scores a critical hit against a creature traveling by means of the \textit{astral projection} spell, the dreadnought can cut the target's silver cord instead of dealing damage.

\DndMonsterAction{Unusual Nature}
The dreadnought doesn't require air, food, drink, or sleep.

\DndMonsterSection{Actions}
\DndMonsterAction{Multiattack}
The dreadnought makes one Bite attack and two Claw attacks.

\DndMonsterAction{Bite}
\textsl{Melee Weapon Attack:} 16 to hit, reach 10 ft., one target. \textsl{Hit:} 36 (5d10 + 9) force damage. If the target is a Huge or smaller creature and this damage reduces it to 0 hit points or it is incapacitated, the dreadnought swallows it. The swallowed target, along with everything it is wearing and carrying, appears in an unoccupied space on the floor of the Demiplanar Donjon.

\DndMonsterAction{Claw}
\textsl{Melee Weapon Attack:} 16 to hit, reach 20 ft., one target. \textsl{Hit:} 19 (3d6 + 9) force damage.

\DndMonsterSection{Legendary Actions}
Astral Dreadnought can take 3 legendary actions, choosing from the options below. Only one legendary action can be used at a time and only at the end of another creature's turn. Astral Dreadnought regains spent legendary actions at the start of its turn.

\begin{DndMonsterLegendaryActions}
\DndMonsterLegendaryAction{Claw}{The dreadnought makes one Claw attack.
}
\DndMonsterLegendaryAction{Donjon Visit (Costs 2 Actions)}{One Huge or smaller creature that the dreadnought can see within 60 feet of it must succeed on a DC 19 Charisma saving throw or be teleported to an unoccupied space on the floor of the Demiplanar Donjon. At the end of the target's next turn, it reappears in the space it left or in the nearest unoccupied space if that space is occupied.
}
\DndMonsterLegendaryAction{Psychic Projection (Costs 3 Actions)}{Each creature within 60 feet of the dreadnought must make a DC 19 Wisdom saving throw, taking 26 (4d10 + 4) psychic damage on a failed save, or half as much damage on a successful one.
}
\end{DndMonsterLegendaryActions}
\end{multicols}
\end{DndMonster}



\begin{DndMonster}[float=t]{Barbed Devil}
\DndMonsterType{Medium fiend (devil), lawful evil}
\DndMonsterBasics[
armor-class = {15 (natural armor)},
hit-points = {\DndDice{13d8 + 52}},
speed = {30 ft.}
]
\DndMonsterAbilityScores[
 str = 16,
 dex = 17,
 con = 18,
 int = 12,
 wis = 14,
 cha = 14
]
\DndMonsterDetails[
saving-throws = {Str +6, Con +7, Wis +5, Cha +5},
skills = {Deception +5, Insight +5, Perception +8},
damage-resistances = {cold; bludgeoning, piercing, slashing from nonmagical attacks that aren't silvered},
damage-immunities = {fire, poison},
condition-immunities = {poisoned},
senses = {darkvision 120 ft., passive perception 18},
languages = {Infernal, telepathy 120 ft.},
challenge = 5,
]
\DndMonsterAction{Barbed Hide}
At the start of each of its turns, the barbed devil deals 5 (1d10) piercing damage to any creature grappling it.

\DndMonsterAction{Devil's Sight}
Magical darkness doesn't impede the devil's darkvision.

\DndMonsterAction{Magic Resistance}
The devil has advantage on saving throws against spells and other magical effects.

\DndMonsterSection{Actions}
\DndMonsterAction{Multiattack}
The devil makes three melee attacks: one with its tail and two with its claws. Alternatively, it can use Hurl Flame twice.

\DndMonsterAction{Claw}
\textsl{Melee Weapon Attack:} 6 to hit, reach 5 ft., one target. \textsl{Hit:} 6 (1d6 + 3) piercing damage.

\DndMonsterAction{Tail}
\textsl{Melee Weapon Attack:} 6 to hit, reach 5 ft., one target. \textsl{Hit:} 10 (2d6 + 3) piercing damage.

\DndMonsterAction{Hurl Flame}
\textsl{Ranged Spell Attack:} 5 to hit, range 150 ft., one target. \textsl{Hit:} 10 (3d6) fire damage. If the target is a flammable object that isn't being worn or carried, it also catches fire.

\end{DndMonster}



\begin{DndMonster}[float=t]{Barlgura}
\DndMonsterType{Large fiend (demon), chaotic evil}
\DndMonsterBasics[
armor-class = {15 (natural armor)},
hit-points = {\DndDice{8d10 + 24}},
speed = {40 ft., climb 40 ft.}
]
\DndMonsterAbilityScores[
 str = 18,
 dex = 15,
 con = 16,
 int = 7,
 wis = 14,
 cha = 9
]
\DndMonsterDetails[
saving-throws = {Dex +5, Con +6},
skills = {Perception +5, Stealth +5},
damage-resistances = {cold, fire, lightning},
damage-immunities = {poison},
condition-immunities = {poisoned},
senses = {blindsight 30 ft., darkvision 120 ft., passive perception 15},
languages = {Abyssal, telepathy 120 ft.},
challenge = 5,
]
\DndMonsterAction{Reckless}
At the start of its turn, the barlgura can gain advantage on all melee weapon attack rolls it makes during that turn, but attack rolls against it have advantage until the start of its next turn.

\DndMonsterAction{Running Leap}
The barlgura's long jump is up to 40 feet and its high jump is up to 20 feet when it has a running start.

\DndMonsterAction{Innate Spellcasting}
The barlgura's spellcasting ability is Wisdom (spell save DC 13). The barlgura can innately cast the following spells, requiring no material components:
\begin{DndMonsterSpells}
  \DndInnateSpellLevel[1]{entangle}
  \DndInnateSpellLevel[2]{disguise self}
\end{DndMonsterSpells}
\DndMonsterSection{Actions}
\DndMonsterAction{Multiattack}
The barlgura makes three attacks: one with its bite and two with its fists.

\DndMonsterAction{Bite}
\textsl{Melee Weapon Attack:} 7 to hit, reach 5 ft., one target. \textsl{Hit:} 11 (2d6 + 4) piercing damage.

\DndMonsterAction{Fist}
\textsl{Melee Weapon Attack:} 7 to hit, reach 5 ft., one target. \textsl{Hit:} 9 (1d10 + 4) bludgeoning damage.

\end{DndMonster}



\begin{DndMonster}[float=t]{Bearded Devil}
\DndMonsterType{Medium fiend (devil), lawful evil}
\DndMonsterBasics[
armor-class = {13 (natural armor)},
hit-points = {\DndDice{8d8 + 16}},
speed = {30 ft.}
]
\DndMonsterAbilityScores[
 str = 16,
 dex = 15,
 con = 15,
 int = 9,
 wis = 11,
 cha = 11
]
\DndMonsterDetails[
saving-throws = {Str +5, Con +4, Wis +2},
damage-resistances = {cold; bludgeoning, piercing, slashing from nonmagical attacks that aren't silvered},
damage-immunities = {fire, poison},
condition-immunities = {poisoned},
senses = {darkvision 120 ft., passive perception 10},
languages = {Infernal, telepathy 120 ft.},
challenge = 3,
]
\DndMonsterAction{Devil's Sight}
Magical darkness doesn't impede the devil's darkvision.

\DndMonsterAction{Magic Resistance}
The devil has advantage on saving throws against spells and other magical effects.

\DndMonsterAction{Steadfast}
The devil can't be frightened while it can see an allied creature within 30 feet of it.

\DndMonsterSection{Actions}
\DndMonsterAction{Multiattack}
The devil makes two attacks: one with its beard and one with its glaive.

\DndMonsterAction{Beard}
\textsl{Melee Weapon Attack:} 5 to hit, reach 5 ft., one creature. \textsl{Hit:} 6 (1d8 + 2) piercing damage, and the target must succeed on a DC 12 Constitution saving throw or be poisoned for 1 minute. While poisoned in this way, the target can't regain hit points. The target can repeat the saving throw at the end of each of its turns, ending the effect on itself on a success.

\DndMonsterAction{Glaive}
\textsl{Melee Weapon Attack:} 5 to hit, reach 10 ft., one target. \textsl{Hit:} 8 (1d10 + 3) slashing damage. If the target is a creature other than an undead or a construct, it must succeed on a DC 12 Constitution saving throw or lose 5 (1d10) hit points at the start of each of its turns due to an infernal wound. Each time the devil hits the wounded target with this attack, the damage dealt by the wound increases by 5 (1d10). Any creature can take an action to stanch the wound with a successful DC 12 Wisdom (Medicine) check. The wound also closes if the target receives magical healing.

\end{DndMonster}



\begin{DndMonster}[float=t]{Behir}
\DndMonsterType{Huge monstrosity, neutral evil}
\DndMonsterBasics[
armor-class = {17 (natural armor)},
hit-points = {\DndDice{16d12 + 64}},
speed = {50 ft., climb 40 ft.}
]
\DndMonsterAbilityScores[
 str = 23,
 dex = 16,
 con = 18,
 int = 7,
 wis = 14,
 cha = 12
]
\DndMonsterDetails[
skills = {Perception +6, Stealth +7},
damage-immunities = {lightning},
senses = {darkvision 90 ft., passive perception 16},
languages = {Draconic},
challenge = 11,
]
\DndMonsterSection{Actions}
\DndMonsterAction{Multiattack}
The behir makes two attacks: one with its bite and one to constrict.

\DndMonsterAction{Bite}
\textsl{Melee Weapon Attack:} 10 to hit, reach 10 ft., one target. \textsl{Hit:} 22 (3d10 + 6) piercing damage.

\DndMonsterAction{Constrict}
\textsl{Melee Weapon Attack:} 10 to hit, reach 5 ft., one Large or smaller creature. \textsl{Hit:} 17 (2d10 + 6) bludgeoning damage plus 17 (2d10 + 6) slashing damage. The target is grappled (escape DC 16) if the behir isn't already constricting a creature, and the target is restrained until this grapple ends.

\DndMonsterAction{Lightning Breath (Recharge 5\textendash 6)}
The behir exhales a line of lightning that is 20 feet long and 5 feet wide. Each creature in that line must make a DC 16 Dexterity saving throw, taking 66 (12d10) lightning damage on a failed save, or half as much damage on a successful one.

\DndMonsterAction{Swallow}
The behir makes one bite attack against a Medium or smaller target it is grappling. If the attack hits, the target is also swallowed, and the grapple ends. While swallowed, the target is blinded and restrained, it has total cover against attacks and other effects outside the behir, and it takes 21 (6d6) acid damage at the start of each of the behir's turns. A behir can have only one creature swallowed at a time.

If the behir takes 30 damage or more on a single turn from the swallowed creature, the behir must succeed on a DC 14 Constitution saving throw at the end of that turn or regurgitate the creature, which falls prone in a space within 10 feet of the behir. If the behir dies, a swallowed creature is no longer restrained by it and can escape from the corpse by using 15 feet of movement, exiting prone.

\end{DndMonster}



\begin{DndMonster}[float=t]{Beholder Zombie}
\DndMonsterType{Large undead, neutral evil}
\DndMonsterBasics[
armor-class = {15 (natural armor)},
hit-points = {\DndDice{11d10 + 33}},
speed = {0 ft., fly 20 ft. \textit{(hover)}}
]
\DndMonsterAbilityScores[
 str = 10,
 dex = 8,
 con = 16,
 int = 3,
 wis = 8,
 cha = 5
]
\DndMonsterDetails[
saving-throws = {Wis +2},
damage-immunities = {poison},
condition-immunities = {poisoned, prone},
senses = {darkvision 60 ft., passive perception 9},
languages = {understands Deep Speech and Undercommon but can't speak},
challenge = 5,
]
\DndMonsterAction{Undead Fortitude}
If damage reduces the zombie to 0 hit points, it must make a Constitution saving throw with a DC of 5 + the damage taken, unless the damage is radiant or from a critical hit. On a success, the zombie drops to 1 hit point instead.

\DndMonsterSection{Actions}
\DndMonsterAction{Bite}
\textsl{Melee Weapon Attack:} 3 to hit, reach 5 ft., one target. \textsl{Hit:} 14 (4d6) piercing damage.

\DndMonsterAction{Eye Ray}
The zombie uses a random magical eye ray, choosing a target that it can see within 60 feet of it.


\begin{description}
\item[1. Paralyzing Ray.]
The targeted creature must succeed on a DC 14 Constitution saving throw or be paralyzed for 1 minute. The target can repeat the saving throw at the end of each of its turns, ending the effect on itself on a success.
\item[2. Fear Ray.]
The targeted creature must succeed on a DC 14 Wisdom saving throw or be frightened for 1 minute. The target can repeat the saving throw at the end of each of its turns, ending the effect on itself on a success.
\item[3. Enervation Ray.]
The targeted creature must make a DC 14 Constitution saving throw, taking 36 (8d8) necrotic damage on a failed save, or half as much damage on a successful one.
\item[4. Disintegration Ray.]
If the target is a creature, it must succeed on a DC 14 Dexterity saving throw or take 45 (10d8) force damage. If this damage reduces the creature to 0 hit points, its body becomes a pile of fine gray dust.

If the target is a Large or smaller nonmagical object or creation of magical force, it is disintegrated without a saving throw. If the target is a Huge or larger nonmagical object or creation of magical force, this ray disintegrates a 10-foot cube of it.

\end{description}

\end{DndMonster}



\begin{DndMonster}[float=t]{Black Rose Bearer}
\DndMonsterType{Medium undead, neutral evil}
\DndMonsterBasics[
armor-class = {11 (natural armor)},
hit-points = {\DndDice{13d8 + 52}},
speed = {20 ft.}
]
\DndMonsterAbilityScores[
 str = 17,
 dex = 6,
 con = 18,
 int = 2,
 wis = 10,
 cha = 5
]
\DndMonsterDetails[
damage-immunities = {poison},
condition-immunities = {charmed, exhaustion, frightened, poisoned},
senses = {darkvision 60 ft., passive perception 10},
languages = {understands the languages it knew in life but can't speak},
challenge = 6,
]
\DndMonsterAction{Berserk}
Whenever the bearer takes damage or makes a Strength or Dexterity saving throw, roll a d6. On a 5 or 6, the bearer goes berserk. On each of its turns while berserk, the bearer has advantage on melee attack rolls, it can Dash as a bonus action, and it must attack the nearest creature it can see. If no creature is near enough to move to and attack, the bearer attacks an object, with preference for an object smaller than itself. Once the bearer goes berserk, it remains berserk until it is destroyed or its creator gives it a pristine black rose.

\DndMonsterAction{Undead Fortitude}
If damage reduces the bearer to 0 hit points, it must make a Constitution saving throw with a DC of 5 plus the damage taken, unless the damage is radiant or from a critical hit. On a successful save, the bearer drops to 1 hit point instead.

\DndMonsterSection{Actions}
\DndMonsterAction{Multiattack}
The bearer makes two Slam attacks.

\DndMonsterAction{Slam}
\textsl{Melee Weapon Attack:} 6 to hit, reach 5 ft., one target. \textsl{Hit:} 12 (2d8 + 3) bludgeoning damage plus 11 (2d10) necrotic damage.

\end{DndMonster}



\begin{DndMonster}[float=t]{Blade Lieutenant}
\DndMonsterType{Medium construct (warforged), lawful evil}
\DndMonsterBasics[
armor-class = {16 (natural armor)},
hit-points = {\DndDice{18d8 + 72}},
speed = {30 ft.}
]
\DndMonsterAbilityScores[
 str = 18,
 dex = 13,
 con = 19,
 int = 14,
 wis = 14,
 cha = 17
]
\DndMonsterDetails[
saving-throws = {Int +6, Cha +7},
skills = {Insight +6, Intimidation +7, Perception +6},
damage-resistances = {poison},
condition-immunities = {exhaustion, poisoned},
languages = {Common},
challenge = 9,
]
\DndMonsterAction{Pack Tactics}
The lieutenant has advantage on an attack roll against a creature if at least one of the lieutenant's allies is within 5 feet of the creature and the ally doesn't have the incapacitated condition.

\DndMonsterSection{Actions}
\DndMonsterAction{Multiattack}
The lieutenant makes three Longsword or Javelin Launcher attacks.

\DndMonsterAction{Longsword}
\textsl{Melee Weapon Attack:} 8 to hit, reach 5 ft., one target. \textsl{Hit:} 17 (3d8 + 4) slashing damage, or 20 (3d10 + 4) slashing damage if used with two hands.

\DndMonsterAction{Javelin Launcher}
\textsl{Ranged Weapon Attack:} 8 to hit, range 30/120 ft., one target. \textsl{Hit:} 14 (3d6 + 4) piercing damage, and the target has the prone condition.

\DndMonsterSection{Bonus Actions}
\DndMonsterAction{Command Ally}
The lieutenant targets one ally it can see within 30 feet of itself. If the target can see or hear the lieutenant, the target can make one melee attack using its reaction, if available, and has advantage on the attack roll.

\DndMonsterAction{Rally the Troops (1/Day)}
The lieutenant ends the charmed and frightened conditions on itself and each creature of its choice that it can see within 30 feet of itself.

\DndMonsterSection{Reactions}
\DndMonsterAction{Parry}
The lieutenant adds 3 to its AC against one melee attack that would hit it. To do so, the lieutenant must see the attacker and be wielding a melee weapon.

\end{DndMonster}



\begin{DndMonster}[float=t]{Blade Scout}
\DndMonsterType{Medium construct (warforged), lawful evil}
\DndMonsterBasics[
armor-class = {18 (natural armor)},
hit-points = {\DndDice{14d8 + 42}},
speed = {30 ft.}
]
\DndMonsterAbilityScores[
 str = 14,
 dex = 20,
 con = 16,
 int = 10,
 wis = 19,
 cha = 10
]
\DndMonsterDetails[
saving-throws = {Dex +8, Wis +7},
skills = {Acrobatics +8, Perception +7, Stealth +8},
damage-resistances = {poison},
condition-immunities = {exhaustion, poisoned},
languages = {Common},
challenge = 7,
]
\DndMonsterAction{Pack Tactics}
The scout has advantage on an attack roll against a creature if at least one of the scout's allies is within 5 feet of the creature and the ally doesn't have the incapacitated condition.

\DndMonsterSection{Actions}
\DndMonsterAction{Multiattack}
The scout makes three Armblade or Bolt Launcher attacks. It can replace one of the attacks with a use of Snare Trap.

\DndMonsterAction{Armblade}
\textsl{Melee Weapon Attack:} 8 to hit, reach 5 ft., one target. \textsl{Hit:} 12 (2d6 + 5) slashing damage.

\DndMonsterAction{Bolt Launcher}
\textsl{Ranged Weapon Attack:} 8 to hit, range 80/320 ft., one target. \textsl{Hit:} 9 (1d8 + 5) piercing damage.

\DndMonsterAction{Snare Trap (1/Day)}
The scout deploys a Tiny mechanical trap on a solid surface within 5 feet of itself. The trap is hidden, requiring a successful DC 17 Intelligence (Investigation) check to find. The trap lasts for 1 minute. Whenever an enemy enters a space within 10 feet of the trap or starts its turn there, it must succeed on a DC 16 Dexterity saving throw or take 21 (6d6) piercing damage and have the prone condition. A creature makes this saving throw only once per turn.

\DndMonsterSection{Bonus Actions}
\DndMonsterAction{Dash}
The scout moves up to its speed without provoking opportunity attacks.

\end{DndMonster}



\begin{DndMonster}[float=t]{Blazebear}
\DndMonsterType{Large monstrosity, unaligned}
\DndMonsterBasics[
armor-class = {14 (natural armor)},
hit-points = {\DndDice{18d10 + 90}},
speed = {30 ft.}
]
\DndMonsterAbilityScores[
 str = 24,
 dex = 17,
 con = 21,
 int = 3,
 wis = 13,
 cha = 16
]
\DndMonsterDetails[
saving-throws = {Str +11, Cha +7},
skills = {Perception +5},
senses = {darkvision 120 ft., passive perception 15},
challenge = 12,
]
\DndMonsterAction{Magic Resistance}
The blazebear has advantage on saving throws against spells and other magical effects.

\DndMonsterAction{Mist Sight}
The blazebear can see normally through heavily obscured areas created by mist or fog, including areas created by spells such as Fog Cloud.

\DndMonsterSection{Actions}
\DndMonsterAction{Multiattack}
The blazebear makes two Bite attacks. It can replace one attack with Stunning Gaze if available.

\DndMonsterAction{Bite}
\textsl{Melee Weapon Attack:} 11 to hit, reach 5 ft., one target. \textsl{Hit:} 20 (2d12 + 7) piercing damage plus 11 (2d10) force damage.

\DndMonsterAction{Stunning Gaze (Recharge 5\textendash 6)}
The blazebear targets two creatures it can see within 120 feet of itself. Each target must succeed on a DC 15 Wisdom saving throw or have the stunned condition until the start of the blazebear's next turn.

\DndMonsterSection{Reactions}
\DndMonsterAction{Antimagic Swipe}
\textsl{Melee Weapon Attack:} 11 to hit, reach 10 ft., one creature casting a spell of 3rd level or lower. \textsl{Hit:} 22 (4d10) force damage, and the target must succeed on a DC 15 Intelligence saving throw or the spell fails and has no effect.

\end{DndMonster}



\begin{DndMonster}[float=t]{Blink Dog}
\DndMonsterType{Medium fey, lawful good}
\DndMonsterBasics[
armor-class = {13},
hit-points = {\DndDice{4d8 + 4}},
speed = {40 ft.}
]
\DndMonsterAbilityScores[
 str = 12,
 dex = 17,
 con = 12,
 int = 10,
 wis = 13,
 cha = 11
]
\DndMonsterDetails[
skills = {Perception +3, Stealth +5},
languages = {Blink Dog, understands Sylvan but can't speak it},
challenge = 1/4,
]
\DndMonsterAction{Keen Hearing and Smell}
The dog has advantage on Wisdom (Perception) checks that rely on hearing or smell.

\DndMonsterSection{Actions}
\DndMonsterAction{Bite}
\textsl{Melee Weapon Attack:} 3 to hit, reach 5 ft., one target. \textsl{Hit:} 4 (1d6 + 1) piercing damage.

\DndMonsterAction{Teleport (Recharge 4\textendash 6)}
The dog magically teleports, along with any equipment it is wearing or carrying, up to 40 feet to an unoccupied space it can see. Before or after teleporting, the dog can make one bite attack.

\end{DndMonster}



\begin{DndMonster}[float=t]{Blue Abishai}
\DndMonsterType{Medium fiend (devil, wizard), lawful evil}
\DndMonsterBasics[
armor-class = {19 (natural armor)},
hit-points = {\DndDice{27d8 + 81}},
speed = {30 ft., fly 50 ft.}
]
\DndMonsterAbilityScores[
 str = 15,
 dex = 14,
 con = 17,
 int = 22,
 wis = 23,
 cha = 18
]
\DndMonsterDetails[
saving-throws = {Int +12, Wis +12},
skills = {Arcana +12},
damage-resistances = {cold; bludgeoning, piercing, slashing from nonmagical attacks that aren't silvered},
damage-immunities = {fire, lightning, poison},
condition-immunities = {poisoned},
senses = {darkvision 120 ft., passive perception 16},
languages = {Draconic, Infernal, telepathy 120 ft.},
challenge = 17,
]
\DndMonsterAction{Devil's Sight}
Magical darkness doesn't impede the abishai's darkvision.

\DndMonsterAction{Magic Resistance}
The abishai has advantage on saving throws against spells and other magical effects.

\DndMonsterAction{Spellcasting}
The abishai casts one of the following spells, using Intelligence as the spellcasting ability (spell save DC 20):
\begin{DndMonsterSpells}
  \DndInnateSpellLevel{disguise self}
  \DndInnateSpellLevel[2]{charm person}
\end{DndMonsterSpells}
\DndMonsterSection{Actions}
\DndMonsterAction{Multiattack}
The abishai makes three Bite or Lightning Strike attacks.

\DndMonsterAction{Bite}
\textsl{Melee Weapon Attack:} 8 to hit, reach 5 ft., one target. \textsl{Hit:} 13 (2d10 + 2) piercing damage plus 14 (4d6) lightning damage.

\DndMonsterAction{Lightning Strike}
\textsl{Ranged Spell Attack:} 12 to hit, range 120 ft., one target. \textsl{Hit:} 36 (8d8) lightning damage.

\DndMonsterSection{Bonus Actions}
\DndMonsterAction{Teleport}
The abishai teleports, along with any equipment it is wearing or carrying, up to 30 feet to an unoccupied space that it can see.

\end{DndMonster}



\begin{DndMonster}[float=t]{Bone Roc}
\DndMonsterType{Huge undead, neutral evil}
\DndMonsterBasics[
armor-class = {15},
hit-points = {\DndDice{14d12 + 42}},
speed = {15 ft., fly 90 ft.}
]
\DndMonsterAbilityScores[
 str = 18,
 dex = 20,
 con = 16,
 int = 2,
 wis = 17,
 cha = 10
]
\DndMonsterDetails[
saving-throws = {Dex +8, Wis +6},
skills = {Perception +6},
damage-vulnerabilities = {bludgeoning},
damage-immunities = {poison},
condition-immunities = {exhaustion, poisoned},
senses = {darkvision 120 ft., passive perception 16},
challenge = 8,
]
\DndMonsterSection{Actions}
\DndMonsterAction{Multiattack}
The bone roc makes one Beak attack and two Talons attacks.

\DndMonsterAction{Beak}
\textsl{Melee Weapon Attack:} 8 to hit, reach 5 ft., one target. \textsl{Hit:} 14 (2d8 + 5) piercing damage.

\DndMonsterAction{Talons}
\textsl{Melee Weapon Attack:} 8 to hit, reach 5 ft., one target. \textsl{Hit:} 12 (2d6 + 5) slashing damage plus 10 (3d6) necrotic damage.

\end{DndMonster}



\begin{DndMonster}[float*=b,width=\textwidth + 8pt]{Borthak}
\begin{multicols}{2}
\DndMonsterType{Huge monstrosity, chaotic evil}
\DndMonsterBasics[
armor-class = {16 (natural armor)},
hit-points = {\DndDice{16d12 + 96}},
speed = {30 ft.}
]
\DndMonsterAbilityScores[
 str = 25,
 dex = 16,
 con = 22,
 int = 4,
 wis = 14,
 cha = 15
]
\DndMonsterDetails[
saving-throws = {Str +12, Con +11, Cha +7},
damage-resistances = {acid, cold},
senses = {darkvision 60 ft., passive perception 12},
challenge = 15,
]
\DndMonsterAction{Glacial Aura}
At the end of the borthak's turn, slippery ice covers surfaces within 10 feet of the borthak. This ice is difficult terrain. When a creature other than the borthak enters the ice's area for the first time on a turn or starts its turn there, it must succeed on a DC 16 Dexterity saving throw or have the prone condition. The ice disappears at the start of the borthak's next turn.

\DndMonsterAction{Webbed Feet}
The borthak can move across icy surfaces without needing to make an ability check. Additionally, difficult terrain composed of ice or snow doesn't cost it extra movement.

\DndMonsterSection{Actions}
\DndMonsterAction{Multiattack}
The borthak makes one Bite attack or uses Noxious Regurgitation if available, and it makes two Stomp attacks.

\DndMonsterAction{Bite}
\textsl{Melee Weapon Attack:} 12 to hit, reach 5 ft., one target. \textsl{Hit:} 16 (2d8 + 7) piercing damage plus 7 (2d6) acid damage.

\DndMonsterAction{Stomp}
\textsl{Melee Weapon Attack:} 12 to hit, reach 10 ft., one target. \textsl{Hit:} 11 (1d8 + 7) bludgeoning damage.

\DndMonsterAction{Noxious Regurgitation (Recharge 5\textendash 6)}
The borthak spews acid at a creature it can see within 120 feet of itself. The target must make a DC 21 Constitution saving throw. On a failed save, the creature takes 24 (7d6) acid damage and has the poisoned condition until the start of the borthak's next turn. On a successful save, the creature takes half as much damage only.

\DndMonsterSection{Reactions}
\DndMonsterAction{Reactive Tail}
When a creature within 10 feet of the borthak hits the borthak with an attack roll, the borthak swings its tail in retaliation. The triggering creature and any creature within 5 feet of it must succeed on a DC 16 Dexterity saving throw or take 16 (2d8 + 7) bludgeoning damage.

\DndMonsterSection{Legendary Actions}
Borthak can take 3 legendary actions, choosing from the options below. Only one legendary action can be used at a time and only at the end of another creature's turn. Borthak regains spent legendary actions at the start of its turn.

\begin{DndMonsterLegendaryActions}
\DndMonsterLegendaryAction{Move}{The borthak moves up to its speed without provoking opportunity attacks.
}
\DndMonsterLegendaryAction{Bite (Costs 2 Actions)}{The borthak makes one Bite attack.
}
\end{DndMonsterLegendaryActions}
\end{multicols}
\end{DndMonster}


\clearpage

\begin{DndMonster}[float=t]{Cadaver Collector}
\DndMonsterType{Large construct, lawful evil}
\DndMonsterBasics[
armor-class = {17 (natural armor)},
hit-points = {\DndDice{18d10 + 90}},
speed = {30 ft.}
]
\DndMonsterAbilityScores[
 str = 21,
 dex = 14,
 con = 20,
 int = 5,
 wis = 11,
 cha = 8
]
\DndMonsterDetails[
damage-immunities = {necrotic; poison; psychic; bludgeoning, piercing, slashing from nonmagical attacks that aren't adamantine},
condition-immunities = {charmed, exhaustion, frightened, paralyzed, petrified, poisoned},
senses = {darkvision 60 ft., passive perception 10},
languages = {understands all languages but can't speak},
challenge = 14,
]
\DndMonsterAction{Magic Resistance}
The collector has advantage on saving throws against spells and other magical effects.

\DndMonsterAction{Unusual Nature}
The collector doesn't require air, food, drink, or sleep.

\DndMonsterSection{Actions}
\DndMonsterAction{Multiattack}
The collector makes two Slam attacks.

\DndMonsterAction{Slam}
\textsl{Melee Weapon Attack:} 10 to hit, reach 5 ft., one target. \textsl{Hit:} 18 (3d8 + 5) bludgeoning damage plus 16 (3d10) necrotic damage.

\DndMonsterAction{Paralyzing Breath (Recharge 5\textendash 6)}
The collector releases paralyzing gas in a 30-foot cone. Each creature in that area must make a successful DC 18 Constitution saving throw or be paralyzed for 1 minute. A paralyzed creature repeats the saving throw at the end of each of its turns, ending the effect on itself with a success.

\DndMonsterSection{Bonus Actions}
\DndMonsterAction{Summon Specters (Recharges after a Short or Long Rest)}
The collector calls up the enslaved spirits of those it has slain; 1d4 specters (without Sunlight Sensitivity) arise in unoccupied spaces within 15 feet of it. The specters act right after the collector on the same initiative count and fight until they're destroyed. They disappear when the collector is destroyed.

\end{DndMonster}



\begin{DndMonster}[float=t]{Camlash}
\DndMonsterType{Huge fiend (demon), chaotic evil}
\DndMonsterBasics[
armor-class = {19 (natural armor)},
hit-points = {\DndDice{26d12 + 156}},
speed = {40 ft., fly 80 ft.}
]
\DndMonsterAbilityScores[
 str = 26,
 dex = 15,
 con = 22,
 int = 20,
 wis = 16,
 cha = 22
]
\DndMonsterDetails[
saving-throws = {Str +14, Con +12, Wis +9, Cha +12},
damage-resistances = {cold; lightning; bludgeoning, piercing, slashing from nonmagical attacks},
damage-immunities = {fire, poison},
condition-immunities = {poisoned},
senses = {truesight 120 ft., passive perception 13},
languages = {Abyssal, telepathy 120 ft.},
challenge = 19,
]
\DndMonsterAction{Death Throes}
Camlash explodes when reduced to 0 hit points, and each creature within 30 feet of Camlash must make a DC 21 Dexterity saving throw, taking 70 (20d6) fire damage on a failed save or half as much damage on a successful one. The explosion ignites flammable objects in that area that aren't being worn or carried.

\DndMonsterAction{Magic Resistance}
Camlash has advantage on saving throws against spells and other magical effects.

\DndMonsterAction{Spider Aura}
Camlash is surrounded by tiny biting spiders that magically appear and disappear from moment to moment. At the start of each of Camlash's turns, each creature within 10 feet of Camlash takes 10 (3d6) poison damage and must succeed on a DC 21 Constitution saving throw or have the paralyzed condition until the start of Camlash's next turn.

\DndMonsterSection{Actions}
\DndMonsterAction{Multiattack}
Camlash makes one Flaming Whip attack and one Lightning Blade attack. Camlash can replace one of these attacks with Teleport.

\DndMonsterAction{Flaming Whip}
\textsl{Melee Weapon Attack:} 15 to hit, reach 30 ft., one target. \textsl{Hit:} 25 (5d6 + 8) fire damage, and if the target is a creature, it must succeed on a DC 21 Strength saving throw or be pulled up to 25 feet toward Camlash.

\DndMonsterAction{Lightning Blade}
\textsl{Melee Weapon Attack:} 15 to hit, reach 10 ft., one target. \textsl{Hit:} 21 (3d8 + 8) slashing damage plus 13 (3d8) lightning damage.

\DndMonsterAction{Teleport}
Camlash magically teleports, along with any equipment she is wearing or carrying, up to 120 feet to an unoccupied space she can see.

\end{DndMonster}



\begin{DndMonster}[float=t]{Citadel Spider}
\DndMonsterType{Gargantuan monstrosity, neutral evil}
\DndMonsterBasics[
armor-class = {18 (natural armor)},
hit-points = {\DndDice{20d20 + 150}},
speed = {50 ft., climb 50 ft.}
]
\DndMonsterAbilityScores[
 str = 26,
 dex = 10,
 con = 21,
 int = 6,
 wis = 12,
 cha = 9
]
\DndMonsterDetails[
condition-immunities = {charmed, exhaustion, frightened},
senses = {blindsight 30 ft., darkvision 120 ft., passive perception 11},
challenge = 18,
]
\DndMonsterAction{Legendary Resistance (3/Day)}
If the spider fails a saving throw, it can choose to succeed instead.

\DndMonsterAction{Spider Climb}
The spider can climb difficult surfaces, including upside down on ceilings, without needing to make an ability check.

\DndMonsterAction{Web Walker}
The spider ignores movement restrictions caused by webbing.

\DndMonsterSection{Actions}
\DndMonsterAction{Multiattack}
The spider makes two Bite attacks. It can replace one of these attacks with a use of Web Bomb.

\DndMonsterAction{Bite}
\textsl{Melee Weapon Attack:} 14 to hit, reach 20 ft., one target. \textsl{Hit:} 19 (2d10 + 8) piercing damage plus 7 (2d6) poison damage.

\DndMonsterAction{Web Bomb}
\textsl{Ranged Weapon Attack:} 14 to hit, range 300 ft./600 ft., one target. \textsl{Hit:} 24 (3d10 + 8) bludgeoning damage, and the target and all creatures within 10 feet of it must succeed on a DC 19 Dexterity saving throw or take 10 (3d6) acid damage and have the restrained condition until the start of the spider's next turn.

\DndMonsterSection{Reactions}
\DndMonsterAction{Absorb Blow}
In response to being hit with an attack roll, the spider's carapace absorbs some of the blow, reducing the damage it takes by 11 (2d10).

\end{DndMonster}



\begin{DndMonster}[float=t]{Cloaker}
\DndMonsterType{Large aberration, chaotic neutral}
\DndMonsterBasics[
armor-class = {14 (natural armor)},
hit-points = {\DndDice{12d10 + 12}},
speed = {10 ft., fly 40 ft.}
]
\DndMonsterAbilityScores[
 str = 17,
 dex = 15,
 con = 12,
 int = 13,
 wis = 12,
 cha = 14
]
\DndMonsterDetails[
skills = {Stealth +5},
senses = {darkvision 60 ft., passive perception 11},
languages = {Deep Speech, Undercommon},
challenge = 8,
]
\DndMonsterAction{Damage Transfer}
While attached to a creature, the cloaker takes only half the damage dealt to it (rounded down). and that creature takes the other half.

\DndMonsterAction{False Appearance}
While the cloaker remains motionless without its underside exposed, it is indistinguishable from a dark leather cloak.

\DndMonsterAction{Light Sensitivity}
While in bright light, the cloaker has disadvantage on attack rolls and Wisdom (Perception) checks that rely on sight.

\DndMonsterSection{Actions}
\DndMonsterAction{Multiattack}
The cloaker makes two attacks: one with its bite and one with its tail.

\DndMonsterAction{Bite}
\textsl{Melee Weapon Attack:} 6 to hit, reach 5 ft., one creature. \textsl{Hit:} 10 (2d6 + 3) piercing damage, and if the target is Large or smaller, the cloaker attaches to it. If the cloaker has advantage against the target, the cloaker attaches to the target's head, and the target is blinded and unable to breathe while the cloaker is attached. While attached, the cloaker can make this attack only against the target and has advantage on the attack roll. The cloaker can detach itself by spending 5 feet of its movement. A creature, including the target, can take its action to detach the cloaker by succeeding on a DC 16 Strength check.

\DndMonsterAction{Tail}
\textsl{Melee Weapon Attack:} 6 to hit, reach 10 ft., one creature. \textsl{Hit:} 7 (1d8 + 3) slashing damage.

\DndMonsterAction{Moan}
Each creature within 60 feet of the cloaker that can hear its moan and that isn't an aberration must succeed on a DC 13 Wisdom saving throw or become frightened until the end of the cloaker's next turn. If a creature's saving throw is successful, the creature is immune to the cloaker's moan for the next 24 hours.

\DndMonsterAction{Phantasms (Recharges after a Short or Long Rest)}
The cloaker magically creates three illusory duplicates of itself if it isn't in bright light. The duplicates move with it and mimic its actions, shifting position so as to make it impossible to track which cloaker is the real one. If the cloaker is ever in an area of bright light, the duplicates disappear.

Whenever any creature targets the cloaker with an attack or a harmful spell while a duplicate remains, that creature rolls randomly to determine whether it targets the cloaker or one of the duplicates. A creature is unaffected by this magical effect if it can't see or if it relies on senses other than sight.

A duplicate has the cloaker's AC and uses its saving throws. If an attack hits a duplicate, or if a duplicate fails a saving throw against an effect that deals damage, the duplicate disappears.

\end{DndMonster}



\begin{DndMonster}[float=t]{Commoner}
\DndMonsterType{Medium humanoid (any race), any alignment}
\DndMonsterBasics[
armor-class = {10},
hit-points = {\DndDice{1d8}},
speed = {30 ft.}
]
\DndMonsterAbilityScores[
 str = 10,
 dex = 10,
 con = 10,
 int = 10,
 wis = 10,
 cha = 10
]
\DndMonsterDetails[
languages = {any one language (usually Common)},
challenge = 0,
]
\DndMonsterSection{Actions}
\DndMonsterAction{Club}
\textsl{Melee Weapon Attack:} 2 to hit, reach 5 ft., one target. \textsl{Hit:} 2 (1d4) bludgeoning damage.

\end{DndMonster}



\begin{DndMonster}[float*=b,width=\textwidth + 8pt]{Cosmic Horror}
\begin{multicols}{2}
\DndMonsterType{Gargantuan aberration, neutral evil}
\DndMonsterBasics[
armor-class = {15 (natural armor)},
hit-points = {\DndDice{16d20 + 112}},
speed = {50 ft., fly 100 ft.}
]
\DndMonsterAbilityScores[
 str = 27,
 dex = 10,
 con = 25,
 int = 24,
 wis = 15,
 cha = 24
]
\DndMonsterDetails[
saving-throws = {Int +13, Wis +8, Cha +13},
damage-immunities = {acid, poison},
condition-immunities = {charmed, frightened, poisoned},
senses = {darkvision 240 ft., passive perception 12},
languages = {Deep Speech, telepathy 240 ft.},
challenge = 18,
]
\DndMonsterAction{Legendary Resistance (3/Day)}
If the horror fails a saving throw, it can choose to succeed instead.

\DndMonsterAction{Unusual Nature}
The horror doesn't require air.

\DndMonsterSection{Actions}
\DndMonsterAction{Multiattack}
The horror makes one Bite attack and two Tentacle attacks.

\DndMonsterAction{Bite}
\textsl{Melee Weapon Attack:} 14 to hit, reach 10 ft., one target. \textsl{Hit:} 22 (4d6 + 8) piercing damage.

\DndMonsterAction{Tentacle}
\textsl{Melee Weapon Attack:} 14 to hit, reach 30 ft., one target. \textsl{Hit:} 18 (3d6 + 8) force damage, and if the target is a creature, it is grappled (escape DC 18). Until this grapple ends, the horror can't use this tentacle against other targets. The horror has 1d8 + 1 tentacles, each of which can grapple one target.

\DndMonsterAction{Psychic Whispers (Recharge 5\textendash 6)}
The horror emits dreadful whispers in a 60-foot-radius sphere centered on itself. Each creature in the sphere that isn't an Aberration must make a DC 21 Wisdom saving throw, taking 33 (6d10) psychic damage on a failed save, or half as much damage on a successful one.

\DndMonsterSection{Legendary Actions}
Cosmic Horror can take 3 legendary actions, choosing from the options below. Only one legendary action can be used at a time and only at the end of another creature's turn. Cosmic Horror regains spent legendary actions at the start of its turn.

\begin{DndMonsterLegendaryActions}
\DndMonsterLegendaryAction{Crushing Tentacle}{The horror crushes one creature it is grappling. The grappled creature must make a DC 22 Constitution saving throw, taking 18 (3d6 + 8) force damage on a failed save, or half as much damage on a successful one.
}
\DndMonsterLegendaryAction{Poison Jet (Costs 2 Actions)}{Foul gas squirts from the horror in a 30-foot line that is 5 feet wide. Each creature in the line must succeed on a DC 21 Constitution saving throw or take 14 (4d6) poison damage.
}
\DndMonsterLegendaryAction{Teleport (Costs 2 Actions)}{The horror teleports, along with any creatures it is grappling, to an unoccupied space it can see within 120 feet of itself.
}
\end{DndMonsterLegendaryActions}
\end{multicols}
\end{DndMonster}



\begin{DndMonster}[float=t]{Couatl}
\DndMonsterType{Medium celestial, lawful good}
\DndMonsterBasics[
armor-class = {19 (natural armor)},
hit-points = {\DndDice{13d8 + 39}},
speed = {30 ft., fly 90 ft.}
]
\DndMonsterAbilityScores[
 str = 16,
 dex = 20,
 con = 17,
 int = 18,
 wis = 20,
 cha = 18
]
\DndMonsterDetails[
saving-throws = {Con +5, Wis +7, Cha +6},
damage-resistances = {radiant},
damage-immunities = {psychic; bludgeoning, piercing, slashing from nonmagical attacks},
senses = {truesight 120 ft., passive perception 15},
languages = {all, telepathy 120 ft.},
challenge = 4,
]
\DndMonsterAction{Magic Weapons}
The couatl's weapon attacks are magical.

\DndMonsterAction{Shielded Mind}
The couatl is immune to scrying and to any effect that would sense its emotions, read its thoughts, or detect its location.

\DndMonsterAction{Innate Spellcasting}
The couatl's spellcasting ability is Charisma (spell save DC 14). It can innately cast the following spells, requiring only verbal components:
\begin{DndMonsterSpells}
  \DndInnateSpellLevel{detect evil and good}
  \DndInnateSpellLevel[3]{bless}
  \DndInnateSpellLevel[1]{dream}
\end{DndMonsterSpells}
\DndMonsterSection{Actions}
\DndMonsterAction{Bite}
\textsl{Melee Weapon Attack:} 8 to hit, reach 5 ft., one creature. \textsl{Hit:} 8 (1d6 + 5) piercing damage, and the target must succeed on a DC 13 Constitution saving throw or be poisoned for 24 hours. Until this poison ends, the target is unconscious. Another creature can use an action to shake the target awake.

\DndMonsterAction{Constrict}
\textsl{Melee Weapon Attack:} 6 to hit, reach 10 ft., one Medium or smaller creature. \textsl{Hit:} 10 (2d6 + 3) bludgeoning damage, and the target is grappled (escape DC 15). Until this grapple ends, the target is restrained, and the couatl can't constrict another target.

\DndMonsterAction{Change Shape}
The couatl magically polymorphs into a humanoid or beast that has a challenge rating equal to or less than its own, or back into its true form. It reverts to its true form if it dies. Any equipment it is wearing or carrying is absorbed or borne by the new form (the couatl's choice).

In a new form, the couatl retains its game statistics and ability to speak, but its AC, movement modes, Strength, Dexterity, and other actions are replaced by those of the new form, and it gains any statistics and capabilities (except class features, legendary actions, and lair actions) that the new form has but that it lacks. If the new form has a bite attack, the couatl can use its bite in that form.

\end{DndMonster}



\begin{DndMonster}[float=t]{Cult Fanatic}
\DndMonsterType{Medium humanoid (any race), any non-good alignment}
\DndMonsterBasics[
armor-class = {13 (\textit{leather armor})},
hit-points = {\DndDice{6d8 + 6}},
speed = {30 ft.}
]
\DndMonsterAbilityScores[
 str = 11,
 dex = 14,
 con = 12,
 int = 10,
 wis = 13,
 cha = 14
]
\DndMonsterDetails[
skills = {Deception +4, Persuasion +4, Religion +2},
languages = {any one language (usually Common)},
challenge = 2,
]
\DndMonsterAction{Dark Devotion}
The fanatic has advantage on saving throws against being charmed or frightened.

\DndMonsterAction{Spellcasting}
The fanatic is a 4th-level spellcaster. Its spellcasting ability is Wisdom (spell save DC 11, 3 to hit with spell attacks). The fanatic has the following cleric spells prepared:
\begin{DndMonsterSpells}
  \DndMonsterSpellLevel{light}
   \DndMonsterSpellLevel[1][4]{command}
   \DndMonsterSpellLevel[2][3]{hold person}
\end{DndMonsterSpells}
\DndMonsterSection{Actions}
\DndMonsterAction{Multiattack}
The fanatic makes two melee attacks.

\DndMonsterAction{Dagger}
\textsl{Melee or Ranged Weapon Attack:} 4 to hit, reach 5 ft. or range 20/60 ft., one creature. \textsl{Hit:} 4 (1d4 + 2) piercing damage.

\end{DndMonster}



\begin{DndMonster}[float=t]{Cultist}
\DndMonsterType{Medium humanoid (any race), any non-good alignment}
\DndMonsterBasics[
armor-class = {12 (\textit{leather armor})},
hit-points = {\DndDice{2d8}},
speed = {30 ft.}
]
\DndMonsterAbilityScores[
 str = 11,
 dex = 12,
 con = 10,
 int = 10,
 wis = 11,
 cha = 10
]
\DndMonsterDetails[
skills = {Deception +2, Religion +2},
languages = {any one language (usually Common)},
challenge = 1/8,
]
\DndMonsterAction{Dark Devotion}
The cultist has advantage on saving throws against being charmed or frightened.

\DndMonsterSection{Actions}
\DndMonsterAction{Scimitar}
\textsl{Melee Weapon Attack:} 3 to hit, reach 5 ft., one creature. \textsl{Hit:} 4 (1d6 + 1) slashing damage.

\end{DndMonster}



\begin{DndMonster}[float=t]{Cyclops}
\DndMonsterType{Huge giant, chaotic neutral}
\DndMonsterBasics[
armor-class = {14 (natural armor)},
hit-points = {\DndDice{12d12 + 60}},
speed = {30 ft.}
]
\DndMonsterAbilityScores[
 str = 22,
 dex = 11,
 con = 20,
 int = 8,
 wis = 6,
 cha = 10
]
\DndMonsterDetails[
languages = {Giant},
challenge = 6,
]
\DndMonsterAction{Poor Depth Perception}
The cyclops has disadvantage on any attack roll against a target more than 30 feet away.

\DndMonsterSection{Actions}
\DndMonsterAction{Multiattack}
The cyclops makes two greatclub attacks.

\DndMonsterAction{Greatclub}
\textsl{Melee Weapon Attack:} 9 to hit, reach 10 ft., one target. \textsl{Hit:} 19 (3d8 + 6) bludgeoning damage.

\DndMonsterAction{Rock}
\textsl{Ranged Weapon Attack:} 9 to hit, range 30/120 ft., one target. \textsl{Hit:} 28 (4d10 + 6) bludgeoning damage.

\end{DndMonster}



\begin{DndMonster}[float=t]{Deadbark Dryad}
\DndMonsterType{Medium fey, chaotic evil}
\DndMonsterBasics[
armor-class = {16 (natural armor)},
hit-points = {\DndDice{22d8 + 88}},
speed = {30 ft.}
]
\DndMonsterAbilityScores[
 str = 17,
 dex = 16,
 con = 18,
 int = 11,
 wis = 16,
 cha = 18
]
\DndMonsterDetails[
saving-throws = {Con +9, Cha +9},
skills = {Perception +8, Stealth +8},
damage-resistances = {bludgeoning, piercing, slashing from nonmagical attacks},
damage-immunities = {poison},
condition-immunities = {poisoned},
senses = {darkvision 60 ft., passive perception 18},
languages = {Elvish, Sylvan},
challenge = 13,
]
\DndMonsterAction{Bramble Walk}
Difficult terrain composed of vegetation, such as foliage or thorns, doesn't cost the dryad extra movement.

\DndMonsterAction{Magic Resistance}
The dryad has advantage on saving throws against spells and other magical effects.

\DndMonsterAction{Speak with Beasts and Plants}
The dryad can communicate with Beasts and Plants as if they shared a language.

\DndMonsterAction{Spellcasting}
The dryad casts one of the following spells, requiring no material components and using Charisma as the spellcasting ability (spell save DC 17):
\begin{DndMonsterSpells}
  \DndInnateSpellLevel{Druidcraft}
  \DndInnateSpellLevel[]{Dispel Magic}
  \DndInnateSpellLevel[2]{Pass without Trace}
\end{DndMonsterSpells}
\DndMonsterSection{Actions}
\DndMonsterAction{Multiattack}
The dryad makes two Poisonous Thorn attacks and one Sapping Vine attack.

\DndMonsterAction{Poisonous Thorn}
\textsl{Melee or Ranged Weapon Attack:} 8 to hit, reach 5 ft. or range 120 ft., one target. \textsl{Hit:} 13 (4d4 + 3) piercing damage plus 10 (3d6) poison damage. If the target is a creature, it must succeed on a DC 17 Constitution saving throw or have the poisoned condition until the start of the dryad's next turn.

\DndMonsterAction{Sapping Vine}
\textsl{Melee Weapon Attack:} 8 to hit, reach 30 ft., one target. \textsl{Hit:} The target has the grappled condition (escape DC 16). Until the grapple ends, the target has the restrained condition, and the dryad can't use the same vine on another target. A creature restrained in this way takes 13 (3d8) necrotic damage at the start of its turn.

The dryad has six vines. Each vine can be attacked (AC 20; 10 hit points; immunity to poison and psychic damage). Destroying a vine deals no damage to the dryad, but any creature grappled by that vine no longer has the grappled condition. All vines immediately wither and disappear when the dryad is reduced to 0 hit points.

\end{DndMonster}



\begin{DndMonster}[float=t]{Death Knight}
\DndMonsterType{Medium undead, chaotic evil}
\DndMonsterBasics[
armor-class = {20 (\textit{plate armor})},
hit-points = {\DndDice{19d8 + 95}},
speed = {30 ft.}
]
\DndMonsterAbilityScores[
 str = 20,
 dex = 11,
 con = 20,
 int = 12,
 wis = 16,
 cha = 18
]
\DndMonsterDetails[
saving-throws = {Dex +6, Wis +9, Cha +10},
damage-immunities = {necrotic, poison},
condition-immunities = {exhaustion, frightened, poisoned},
senses = {darkvision 120 ft., passive perception 13},
languages = {Abyssal, Common},
challenge = 17,
]
\DndMonsterAction{Magic Resistance}
The death knight has advantage on saving throws against spells and other magical effects.

\DndMonsterAction{Marshal Undead}
Unless the death knight is incapacitated, it and undead creatures of its choice within 60 feet of it have advantage on saving throws against features that turn undead.

\DndMonsterAction{Spellcasting}
The death knight is a 19th-level spellcaster. Its spellcasting ability is Charisma (spell save DC 18, 10 to hit with spell attacks). It has the following paladin spells prepared:
\begin{DndMonsterSpells}
   \DndMonsterSpellLevel[1][4]{command}
   \DndMonsterSpellLevel[2][3]{hold person}
   \DndMonsterSpellLevel[3][3]{dispel magic}
   \DndMonsterSpellLevel[4][3]{banishment}
   \DndMonsterSpellLevel[5][2]{destructive wave}
\end{DndMonsterSpells}
\DndMonsterSection{Actions}
\DndMonsterAction{Multiattack}
The death knight makes three longsword attacks.

\DndMonsterAction{Longsword}
\textsl{Melee Weapon Attack:} 11 to hit, reach 5 ft., one target. \textsl{Hit:} 9 (1d8 + 5) slashing damage, or 10 (1d10 + 5) slashing damage if used with two hands, plus 18 (4d8) necrotic damage.

\DndMonsterAction{Hellfire Orb (1/Day)}
The death knight hurls a magical ball of fire that explodes at a point it can see within 120 feet of it. Each creature in a 20-foot-radius sphere centered on that point must make a DC 18 Dexterity saving throw. The sphere spreads around corners. A creature takes 35 (10d6) fire damage and 35 (10d6) necrotic damage on a failed save, or half as much damage on a successful one.

\DndMonsterSection{Reactions}
\DndMonsterAction{Parry}
The death knight adds 6 to its AC against one melee attack that would hit it. To do so, the death knight must see the attacker and be wielding a melee weapon.

\end{DndMonster}



\begin{DndMonster}[float=t]{Death Slaad}
\DndMonsterType{Medium aberration (shapechanger), chaotic evil}
\DndMonsterBasics[
armor-class = {18 (natural armor)},
hit-points = {\DndDice{20d8 + 80}},
speed = {30 ft.}
]
\DndMonsterAbilityScores[
 str = 20,
 dex = 15,
 con = 19,
 int = 15,
 wis = 10,
 cha = 16
]
\DndMonsterDetails[
skills = {Arcana +6, Perception +8},
damage-resistances = {acid, cold, fire, lightning, thunder},
senses = {blindsight 60 ft., darkvision 60 ft., passive perception 18},
languages = {Slaad, telepathy 60 ft.},
challenge = 10,
]
\DndMonsterAction{Shapechanger}
The slaad can use its action to polymorph into a Small or Medium humanoid, or back into its true form. Its statistics, other than its size, are the same in each form. Any equipment it is wearing or carrying isn't transformed. It reverts to its true form if it dies.

\DndMonsterAction{Magic Resistance}
The slaad has advantage on saving throws against spells and other magical effects.

\DndMonsterAction{Magic Weapons}
The slaad's weapon attacks are magical.

\DndMonsterAction{Regeneration}
The slaad regains 10 hit points at the start of its turn if it has at least 1 hit point.

\DndMonsterAction{Innate Spellcasting}
The slaad's innate spellcasting ability is Charisma (spell save DC 15, 7 to hit with spell attacks). The slaad can innately cast the following spells, requiring no material components:
\begin{DndMonsterSpells}
  \DndInnateSpellLevel{detect magic}
  \DndInnateSpellLevel[2]{fear}
  \DndInnateSpellLevel[1]{cloudkill}
\end{DndMonsterSpells}
\DndMonsterSection{Actions}
\DndMonsterAction{Multiattack}
The slaad makes three attacks: one with its bite and two with its claws or greatsword.

\DndMonsterAction{Bite (Slaad Form Only)}
\textsl{Melee Weapon Attack:} 9 to hit, reach 5 ft., one target. \textsl{Hit:} 9 (1d8 + 5) piercing damage plus 7 (2d6) necrotic damage.

\DndMonsterAction{Claws (Slaad Form Only)}
\textsl{Melee Weapon Attack:} 9 to hit, reach 5 ft., one target. \textsl{Hit:} 10 (1d10 + 5) slashing damage plus 7 (2d6) necrotic damage.

\DndMonsterAction{Greatsword}
\textsl{Melee Weapon Attack:} 9 to hit, reach 5 ft., one target. \textsl{Hit:} 12 (2d6 + 5) slashing damage plus 7 (2d6) necrotic damage.

\end{DndMonster}



\begin{DndMonster}[float=t]{Deathwolf}
\DndMonsterType{Medium undead, chaotic evil}
\DndMonsterBasics[
armor-class = {15 (natural armor)},
hit-points = {\DndDice{18d8 + 72}},
speed = {40 ft.}
]
\DndMonsterAbilityScores[
 str = 20,
 dex = 16,
 con = 18,
 int = 10,
 wis = 17,
 cha = 19
]
\DndMonsterDetails[
saving-throws = {Str +10, Dex +8, Cha +9},
skills = {Perception +13, Stealth +8},
damage-immunities = {poison; bludgeoning, piercing, slashing from nonmagical attacks that aren't silvered},
condition-immunities = {charmed, exhaustion, frightened, poisoned},
senses = {darkvision 60 ft., passive perception 23},
languages = {Common},
challenge = 15,
]
\DndMonsterAction{Legendary Resistance (2/Day)}
If the deathwolf fails a saving throw, it can choose to succeed instead.

\DndMonsterAction{Moon's Grace}
When the deathwolf falls, it descends at a rate of 60 feet per round and takes no falling damage.

\DndMonsterSection{Actions}
\DndMonsterAction{Multiattack}
The deathwolf makes one Bite attack and two Claw attacks. It can replace one of these attacks with Phantom Deathwolf if available.

\DndMonsterAction{Bite}
\textsl{Melee Weapon Attack:} 10 to hit, reach 5 ft., one target. \textsl{Hit:} 14 (2d8 + 5) piercing damage plus 9 (2d8) necrotic damage. The target must succeed on a DC 16 Wisdom saving throw or have disadvantage on saving throws against the frightened condition. This curse lasts until removed by the \textit{Remove Curse} spell or other magic.

\DndMonsterAction{Claw}
\textsl{Melee Weapon Attack:} 10 to hit, reach 5 ft., one target. \textsl{Hit:} 12 (2d6 + 5) slashing damage plus 4 (1d8) necrotic damage.

\DndMonsterAction{Phantom Deathwolf (Recharge 5\textendash 6)}
The deathwolf creates a terrifying phantom of itself in the mind of a creature the deathwolf can see within 60 feet of itself. The target must succeed on a DC 17 Intelligence saving throw or have the frightened condition for 1 minute.

While the target is frightened, the phantom deals 21 (6d6) psychic damage to the target at the start of each of the target's turns. A frightened target can repeat the saving throw at the end of each of its turns, ending the effect on itself on a success.

\DndMonsterSection{Reactions}
\DndMonsterAction{Imposing Slash}
When a creature within 5 feet of the deathwolf makes an attack roll against it, the deathwolf forces the creature to succeed on a DC 17 Wisdom saving throw or have disadvantage on that roll. After the attack hits or misses, the deathwolf makes one Claw attack against the creature.

\DndMonsterAction{Phase Step}
Immediately after taking damage, the deathwolf teleports up to 30 feet to an unoccupied space it can see that is in dim light or darkness.

\end{DndMonster}



\begin{DndMonster}[float=t]{Degloth}
\DndMonsterType{Large fiend (demon), chaotic evil}
\DndMonsterBasics[
armor-class = {18 (natural armor)},
hit-points = {\DndDice{14d10 + 56}},
speed = {40 ft., climb 40 ft.}
]
\DndMonsterAbilityScores[
 str = 23,
 dex = 17,
 con = 18,
 int = 6,
 wis = 11,
 cha = 9
]
\DndMonsterDetails[
saving-throws = {Str +10, Con +8},
skills = {Athletics +10, Perception +4},
damage-resistances = {cold; fire; lightning; bludgeoning, piercing, slashing from nonmagical attacks},
damage-immunities = {poison},
condition-immunities = {charmed, exhaustion, frightened, poisoned},
senses = {darkvision 120 ft., passive perception 14},
languages = {Abyssal, telepathy 120 ft.},
challenge = 11,
]
\DndMonsterAction{Magic Resistance}
The degloth has advantage on saving throws against spells and other magical effects.

\DndMonsterSection{Actions}
\DndMonsterAction{Multiattack}
The degloth makes two Razor Fist attacks.

\DndMonsterAction{Razor Fist}
\textsl{Melee Weapon Attack:} 10 to hit, reach 10 ft., one target. \textsl{Hit:} 17 (2d10 + 6) slashing damage, and if the target is a Medium or smaller creature, the target has the grappled condition (escape DC 18). Until this grapple ends, the target has the restrained condition, and the degloth can't use this fist to grapple another target. The degloth has two fists.

\DndMonsterSection{Bonus Actions}
\DndMonsterAction{Crush}
The degloth targets one creature currently grappled by it. The target must make a DC 18 Strength saving throw, taking 15 (2d8 + 6) bludgeoning damage on a failed save or half as much damage on a successful one.

\end{DndMonster}



\begin{DndMonster}[float=t]{Deva}
\DndMonsterType{Medium celestial, lawful good}
\DndMonsterBasics[
armor-class = {17 (natural armor)},
hit-points = {\DndDice{16d8 + 64}},
speed = {30 ft., fly 90 ft.}
]
\DndMonsterAbilityScores[
 str = 18,
 dex = 18,
 con = 18,
 int = 17,
 wis = 20,
 cha = 20
]
\DndMonsterDetails[
saving-throws = {Wis +9, Cha +9},
skills = {Insight +9, Perception +9},
damage-resistances = {radiant; bludgeoning, piercing, slashing from nonmagical attacks},
condition-immunities = {charmed, exhaustion, frightened},
senses = {darkvision 120 ft., passive perception 19},
languages = {all, telepathy 120 ft.},
challenge = 10,
]
\DndMonsterAction{Angelic Weapons}
The deva's weapon attacks are magical. When the deva hits with any weapon, the weapon deals an extra 4d8 radiant damage (included in the attack).

\DndMonsterAction{Magic Resistance}
The deva has advantage on saving throws against spells and other magical effects.

\DndMonsterAction{Innate Spellcasting}
The deva's spellcasting ability is Charisma (spell save DC 17). The deva can innately cast the following spells, requiring only verbal components:
\begin{DndMonsterSpells}
  \DndInnateSpellLevel{detect evil and good}
  \DndInnateSpellLevel[1]{commune}
\end{DndMonsterSpells}
\DndMonsterSection{Actions}
\DndMonsterAction{Multiattack}
The deva makes two melee attacks.

\DndMonsterAction{Mace}
\textsl{Melee Weapon Attack:} 8 to hit, reach 5 ft., one target. \textsl{Hit:} 7 (1d6 + 4) bludgeoning damage plus 18 (4d8) radiant damage.

\DndMonsterAction{Healing Touch (3/Day)}
The deva touches another creature. The target magically regains 20 (4d8 + 2) hit points and is freed from any curse, disease, poison, blindness, or deafness.

\DndMonsterAction{Change Shape}
The deva magically polymorphs into a humanoid or beast that has a challenge rating equal to or less than its own, or back into its true form. It reverts to its true form if it dies. Any equipment it is wearing or carrying is absorbed or borne by the new form (the deva's choice).

In a new form, the deva retains its game statistics and ability to speak, but its AC, movement modes, Strength, Dexterity, and special senses are replaced by those of the new form, and it gains any statistics and capabilities (except class features, legendary actions, and lair actions) that the new form has but that it lacks.

\end{DndMonster}



\begin{DndMonster}[float=t]{Drider}
\DndMonsterType{Large monstrosity, chaotic evil}
\DndMonsterBasics[
armor-class = {19 (natural armor)},
hit-points = {\DndDice{13d10 + 52}},
speed = {30 ft., climb 30 ft.}
]
\DndMonsterAbilityScores[
 str = 16,
 dex = 16,
 con = 18,
 int = 13,
 wis = 14,
 cha = 12
]
\DndMonsterDetails[
skills = {Perception +5, Stealth +9},
senses = {darkvision 120 ft., passive perception 15},
languages = {Elvish, Undercommon},
challenge = 6,
]
\DndMonsterAction{Fey Ancestry}
The drider has advantage on saving throws against being charmed, and magic can't put the drider to sleep.

\DndMonsterAction{Spider Climb}
The drider can climb difficult surfaces, including upside down on ceilings, without needing to make an ability check.

\DndMonsterAction{Sunlight Sensitivity}
While in sunlight, the drider has disadvantage on attack rolls, as well as on Wisdom (Perception) checks that rely on sight.

\DndMonsterAction{Web Walker}
The drider ignores movement restrictions caused by webbing.

\DndMonsterAction{Innate Spellcasting}
The drider's innate spellcasting ability is Wisdom (spell save DC 13). The drider can innately cast the following spells, requiring no material components:
\begin{DndMonsterSpells}
  \DndInnateSpellLevel{dancing lights}
  \DndInnateSpellLevel[1]{darkness}
\end{DndMonsterSpells}
\DndMonsterSection{Actions}
\DndMonsterAction{Multiattack}
The drider makes three attacks, either with its longsword or its longbow. It can replace one of those attacks with a bite attack.

\DndMonsterAction{Bite}
\textsl{Melee Weapon Attack:} 6 to hit, reach 5 ft., one creature. \textsl{Hit:} 2 (1d4) piercing damage plus 9 (2d8) poison damage.

\DndMonsterAction{Longsword}
\textsl{Melee Weapon Attack:} 6 to hit, reach 5 ft., one target. \textsl{Hit:} 7 (1d8 + 3) slashing damage, or 8 (1d10 + 3) slashing damage if used with two hands.

\DndMonsterAction{Longbow}
\textsl{Ranged Weapon Attack:} 6 to hit, range 150/600 ft., one target. \textsl{Hit:} 7 (1d8 + 3) piercing damage plus 4 (1d8) poison damage.

\end{DndMonster}



\begin{DndMonster}[float=t]{Drow Mage}
\DndMonsterType{Medium humanoid (elf), neutral evil}
\DndMonsterBasics[
armor-class = {12 (15 with \textit{mage armor})},
hit-points = {\DndDice{10d8}},
speed = {30 ft.}
]
\DndMonsterAbilityScores[
 str = 9,
 dex = 14,
 con = 10,
 int = 17,
 wis = 13,
 cha = 12
]
\DndMonsterDetails[
skills = {Arcana +6, Deception +4, Perception +4, Stealth +5},
senses = {darkvision 120 ft., passive perception 14},
languages = {Elvish, Undercommon},
challenge = 7,
]
\DndMonsterAction{Fey Ancestry}
The drow has advantage on saving throws against being charmed, and magic can't put the drow to sleep.

\DndMonsterAction{Sunlight Sensitivity}
While in sunlight, the drow has disadvantage on attack rolls, as well as on Wisdom (Perception) checks that rely on sight.

\DndMonsterAction{Innate Spellcasting}
The drow's spellcasting ability is Charisma (spell save DC 12). It can innately cast the following spells, requiring no material components:
\begin{DndMonsterSpells}
  \DndInnateSpellLevel{dancing lights}
  \DndInnateSpellLevel[1]{darkness}
\end{DndMonsterSpells}
\DndMonsterAction{Spellcasting}
The drow is a 10th-level spellcaster. Its spellcasting ability is Intelligence (spell save DC 14, 6 to hit with spell attacks). The drow has the following wizard spells prepared:
\begin{DndMonsterSpells}
  \DndMonsterSpellLevel{mage hand}
   \DndMonsterSpellLevel[1][4]{mage armor}
   \DndMonsterSpellLevel[2][3]{alter self}
   \DndMonsterSpellLevel[3][3]{fly}
   \DndMonsterSpellLevel[4][3]{Evard's black tentacles}
   \DndMonsterSpellLevel[5][2]{cloudkill}
\end{DndMonsterSpells}
\DndMonsterSection{Actions}
\DndMonsterAction{Staff}
\textsl{Melee Weapon Attack:} 2 to hit, reach 5 ft., one target. \textsl{Hit:} 2 (1d6 - 1) bludgeoning damage, or 3 (1d8 - 1) bludgeoning damage if used with two hands, plus 3 (1d6) poison damage.

\DndMonsterAction{Summon Demon (1/Day)}
The drow magically summons a \textbf{quasit}, or attempts to summon a \textbf{shadow demon} with a 50\% summoning chance chance of success. The summoned demon appears in an unoccupied space within 60 feet of its summoner, acts as an ally of its summoner, and can't summon other demons. It remains for 10 minutes, until it or its summoner dies, or until its summoner dismisses it as an action.

\end{DndMonster}



\begin{DndMonster}[float=t]{Druid}
\DndMonsterType{Medium humanoid (any race), any alignment}
\DndMonsterBasics[
armor-class = {11 (16 with \textit{barkskin})},
hit-points = {\DndDice{5d8 + 5}},
speed = {30 ft.}
]
\DndMonsterAbilityScores[
 str = 10,
 dex = 12,
 con = 13,
 int = 12,
 wis = 15,
 cha = 11
]
\DndMonsterDetails[
skills = {Medicine +4, Nature +3, Perception +4},
languages = {Druidic plus any two languages},
challenge = 2,
]
\DndMonsterAction{Spellcasting}
The druid is a 4th-level spellcaster. Its spellcasting ability is Wisdom (spell save DC 12, 4 to hit with spell attacks). It has the following druid spells prepared:
\begin{DndMonsterSpells}
  \DndMonsterSpellLevel{druidcraft}
   \DndMonsterSpellLevel[1][4]{entangle}
   \DndMonsterSpellLevel[2][3]{animal messenger}
\end{DndMonsterSpells}
\DndMonsterSection{Actions}
\DndMonsterAction{Quarterstaff}
\textsl{Melee Weapon Attack:} 2 to hit (4 to hit with shillelagh), reach 5 ft., one target. \textsl{Hit:} 3 (1d6) bludgeoning damage, 4 (1d8) bludgeoning damage if wielded with two hands, or 6 (1d8 + 2) bludgeoning damage with \textit{shillelagh}.

\end{DndMonster}



\begin{DndMonster}[float=t]{Earth Elemental}
\DndMonsterType{Large elemental, neutral}
\DndMonsterBasics[
armor-class = {17 (natural armor)},
hit-points = {\DndDice{12d10 + 60}},
speed = {30 ft., burrow 30 ft.}
]
\DndMonsterAbilityScores[
 str = 20,
 dex = 8,
 con = 20,
 int = 5,
 wis = 10,
 cha = 5
]
\DndMonsterDetails[
damage-vulnerabilities = {thunder},
damage-resistances = {bludgeoning, piercing, slashing from nonmagical attacks},
damage-immunities = {poison},
condition-immunities = {exhaustion, paralyzed, petrified, poisoned, unconscious},
senses = {darkvision 60 ft., tremorsense 60 ft., passive perception 10},
languages = {Terran},
challenge = 5,
]
\DndMonsterAction{Earth Glide}
The elemental can burrow through nonmagical, unworked earth and stone. While doing so, the elemental doesn't disturb the material it moves through.

\DndMonsterAction{Siege Monster}
The elemental deals double damage to objects and structures.

\DndMonsterSection{Actions}
\DndMonsterAction{Multiattack}
The elemental makes two slam attacks.

\DndMonsterAction{Slam}
\textsl{Melee Weapon Attack:} 8 to hit, reach 10 ft., one target. \textsl{Hit:} 14 (2d8 + 5) bludgeoning damage.

\end{DndMonster}


\clearpage

\begin{DndMonster}[float=t]{Erinyes}
\DndMonsterType{Medium fiend (devil), lawful evil}
\DndMonsterBasics[
armor-class = {18 (\textit{plate armor})},
hit-points = {\DndDice{18d8 + 72}},
speed = {30 ft., fly 60 ft.}
]
\DndMonsterAbilityScores[
 str = 18,
 dex = 16,
 con = 18,
 int = 14,
 wis = 14,
 cha = 18
]
\DndMonsterDetails[
saving-throws = {Dex +7, Con +8, Wis +6, Cha +8},
damage-resistances = {cold; bludgeoning, piercing, slashing from nonmagical attacks that aren't silvered},
damage-immunities = {fire, poison},
condition-immunities = {poisoned},
senses = {truesight 120 ft., passive perception 12},
languages = {Infernal, telepathy 120 ft.},
challenge = 12,
]
\DndMonsterAction{Hellish Weapons}
The erinyes's weapon attacks are magical and deal an extra 13 (3d8) poison damage on a hit (included in the attacks).

\DndMonsterAction{Magic Resistance}
The erinyes has advantage on saving throws against spells and other magical effects.

\DndMonsterSection{Actions}
\DndMonsterAction{Multiattack}
The erinyes makes three attacks.

\DndMonsterAction{Longsword}
\textsl{Melee Weapon Attack:} 8 to hit, reach 5 ft., one target. \textsl{Hit:} 8 (1d8 + 4) slashing damage, or 9 (1d10 + 4) slashing damage if used with two hands, plus 13 (3d8) poison damage.

\DndMonsterAction{Longbow}
\textsl{Ranged Weapon Attack:} 7 to hit, range 150/600 ft., one target. \textsl{Hit:} 7 (1d8 + 3) piercing damage plus 13 (3d8) poison damage, and the target must succeed on a DC 14 Constitution saving throw or be poisoned. The poison lasts until it is removed by the \textit{lesser restoration} spell or similar magic.

\DndMonsterSection{Reactions}
\DndMonsterAction{Parry}
The erinyes adds 4 to its AC against one melee attack that would hit it. To do so, the erinyes must see the attacker and be wielding a melee weapon.

\end{DndMonster}



\begin{DndMonster}[float=t]{Ettin}
\DndMonsterType{Large giant, chaotic evil}
\DndMonsterBasics[
armor-class = {12 (natural armor)},
hit-points = {\DndDice{10d10 + 30}},
speed = {40 ft.}
]
\DndMonsterAbilityScores[
 str = 21,
 dex = 8,
 con = 17,
 int = 6,
 wis = 10,
 cha = 8
]
\DndMonsterDetails[
skills = {Perception +4},
senses = {darkvision 60 ft., passive perception 14},
languages = {Giant, Orc},
challenge = 4,
]
\DndMonsterAction{Two Heads}
The ettin has advantage on Wisdom (Perception) checks and on saving throws against being blinded, charmed, deafened, frightened, stunned, and knocked unconscious.

\DndMonsterAction{Wakeful}
When one of the ettin's heads is asleep, its other head is awake.

\DndMonsterSection{Actions}
\DndMonsterAction{Multiattack}
The ettin makes two attacks: one with its battleaxe and one with its morningstar.

\DndMonsterAction{Battleaxe}
\textsl{Melee Weapon Attack:} 7 to hit, reach 5 ft., one target. \textsl{Hit:} 14 (2d8 + 5) slashing damage.

\DndMonsterAction{Morningstar}
\textsl{Melee Weapon Attack:} 7 to hit, reach 5 ft., one target. \textsl{Hit:} 14 (2d8 + 5) piercing damage.

\end{DndMonster}



\begin{DndMonster}[float=t]{Eye Monger}
\DndMonsterType{Large aberration, lawful evil}
\DndMonsterBasics[
armor-class = {17 (natural armor)},
hit-points = {\DndDice{13d10 + 78}},
speed = {0 ft., fly 20 ft. \textit{(hover)}}
]
\DndMonsterAbilityScores[
 str = 21,
 dex = 6,
 con = 23,
 int = 7,
 wis = 13,
 cha = 7
]
\DndMonsterDetails[
senses = {darkvision 120 ft., blindsight 120 ft. while the eye monger's eye is closed, passive perception 11},
languages = {Deep Speech},
challenge = 10,
]
\DndMonsterAction{Antimagic Gullet}
Magical effects, including those produced by spells and magic items but excluding those created by artifacts or deities, are suppressed inside the eye monger's gullet. Any spell slot or charge expended by a creature in the gullet to cast a spell or activate a property of a magic item is wasted. While an effect is suppressed, it doesn't function, but the time it spends suppressed counts against its duration. No spell or magical effect that originates outside the eye monger's gullet, except one created by an artifact or a deity, can affect a creature or an object inside the gullet.

\DndMonsterAction{False Appearance}
If the eye monger is motionless and has its eye and mouth closed at the start of combat, it has advantage on its initiative roll. Moreover, if a creature hasn't observed the eye monger move or act, that creature must succeed on a DC 18 Intelligence (Investigation) check to discern that the eye monger is animate.

\DndMonsterAction{Unusual Nature}
The eye monger doesn't require air.

\DndMonsterSection{Actions}
\DndMonsterAction{Bite}
\textsl{Melee Weapon Attack:} 9 to hit, reach 5 ft., one target. \textsl{Hit:} 12 (2d6 + 5) piercing damage, and if the target is a Medium or smaller creature, it must succeed on a DC 18 Dexterity saving throw or be swallowed by the eye monger and deposited in the eye monger's gullet (see Antimagic Gullet). The eye monger can swallow one creature at a time. A swallowed creature is blinded and restrained, has total cover against attacks and other effects originating outside the eye monger, and takes 35 (10d6) acid damage at the start of each of its turns.

If the eye monger takes 25 damage or more on a single turn from a creature inside its gullet, the eye monger regurgitates the swallowed creature, which falls prone in a space within 10 feet of the eye monger. If the eye monger dies, a swallowed creature is no longer restrained by it and can escape from the corpse by using 10 feet of movement, exiting prone.

\end{DndMonster}



\begin{DndMonster}[float*=b,width=\textwidth + 8pt]{False Lich}
\begin{multicols}{2}
\DndMonsterType{Medium undead, neutral evil}
\DndMonsterBasics[
armor-class = {18 (natural armor)},
hit-points = {\DndDice{21d8 + 105}},
speed = {30 ft.}
]
\DndMonsterAbilityScores[
 str = 10,
 dex = 16,
 con = 20,
 int = 25,
 wis = 19,
 cha = 15
]
\DndMonsterDetails[
saving-throws = {Con +12, Int +14, Wis +11, Cha +9},
damage-immunities = {necrotic; poison; psychic; bludgeoning, piercing, slashing from nonmagical attacks},
condition-immunities = {charmed, exhaustion, frightened, paralyzed, poisoned, stunned},
senses = {truesight 60 ft., passive perception 14},
languages = {Abyssal, Common, Draconic, Dwarvish, Elvish, Giant, Infernal, Primordial, Undercommon},
challenge = 21,
]
\DndMonsterAction{Legendary Resistance (3/Day)}
If the false lich fails a saving throw, it can choose to succeed instead.

\DndMonsterAction{Magic Resistance}
The false lich has advantage on saving throws against spells and magical effects.

\DndMonsterAction{Spellcasting}
The false lich casts one of the following spells, requiring no material components and using Intelligence as the spellcasting ability (spell save DC 22):
\begin{DndMonsterSpells}
  \DndInnateSpellLevel{Detect Magic}
  \DndInnateSpellLevel[3]{Dispel Magic}
  \DndInnateSpellLevel[1]{Globe of Invulnerability}
\end{DndMonsterSpells}
\DndMonsterSection{Actions}
\DndMonsterAction{Multiattack}
The false lich makes two Death Rend attacks and uses Bloodcurdling Lament if available.

\DndMonsterAction{Death Rend}
\textsl{Melee Spell Attack:} 14 to hit, reach 5 ft., one target. \textsl{Hit:} 23 (3d10 + 7) necrotic damage.

\DndMonsterAction{Bloodcurdling Lament (Recharge 5\textendash 6)}
The false lich emits a hideous shriek charged with malignant energy. Each creature within 30 feet of the false lich must succeed on a DC 22 Wisdom saving throw or have the frightened condition for 1 minute. While frightened in this way, a creature also has the unconscious condition. An affected creature can repeat the saving throw at the end of each of its turns, ending the effect on itself on a success.

\DndMonsterSection{Bonus Actions}
\DndMonsterAction{Soul Siphon}
The false lich targets one creature it can see within 120 feet of itself. The target must make a DC 22 Charisma saving throw; if the target has the unconscious condition, it has disadvantage on this saving throw. The target takes 21 (6d6) force damage on a failed save or half as much damage on a successful one. The false lich then regains a number of hit points equal to the amount of force damage taken.

If this damage reduces the target to 0 hit points, the target immediately dies, its body disappears, and its soul is trapped inside one of the soul gems within the false lich's skull. After 24 hours, the gem transfers the soul to the false lich's creator.

When the false lich is reduced to 0 hit points, it is destroyed and disintegrates, leaving behind the gems. Crushing a gem releases any souls trapped within, at which point the soul's body re-forms in an unoccupied space nearest to the gem and in the same state as it was when its soul was trapped.

\DndMonsterSection{Legendary Actions}
False Lich can take 3 legendary actions, choosing from the options below. Only one legendary action can be used at a time and only at the end of another creature's turn. False Lich regains spent legendary actions at the start of its turn.

\begin{DndMonsterLegendaryActions}
\DndMonsterLegendaryAction{Spiteful Teleport}{The false lich, along with anything it is wearing or carrying, teleports to an unoccupied space it can see within 60 feet of itself. It then makes one Death Rend attack if possible.
}
\DndMonsterLegendaryAction{Cast a Spell (Costs 2 Actions)}{The false lich uses Spellcasting.
}
\end{DndMonsterLegendaryActions}
\end{multicols}
\end{DndMonster}



\begin{DndMonster}[float=t]{Flameskull}
\DndMonsterType{Tiny undead, neutral evil}
\DndMonsterBasics[
armor-class = {13},
hit-points = {\DndDice{9d4 + 18}},
speed = {0 ft., fly 40 ft. \textit{(hover)}}
]
\DndMonsterAbilityScores[
 str = 1,
 dex = 17,
 con = 14,
 int = 16,
 wis = 10,
 cha = 11
]
\DndMonsterDetails[
skills = {Arcana +5, Perception +2},
damage-resistances = {lightning, necrotic, piercing},
damage-immunities = {cold, fire, poison},
condition-immunities = {charmed, frightened, paralyzed, poisoned, prone},
senses = {darkvision 60 ft., passive perception 12},
languages = {Common},
challenge = 4,
]
\DndMonsterAction{Illumination}
The flameskull sheds either dim light in a 15-foot radius, or bright light in a 15-foot radius and dim light for an additional 15 feet. It can switch between the options as an action.

\DndMonsterAction{Magic Resistance}
The flameskull has advantage on saving throws against spells and other magical effects.

\DndMonsterAction{Rejuvenation}
If the flameskull is destroyed, it regains all its hit points in 1 hour unless holy water is sprinkled on its remains or a \textit{dispel magic} or \textit{remove curse} spell is cast on them.

\DndMonsterAction{Spellcasting}
The flameskull is a 5th-level spellcaster. Its spellcasting ability is Intelligence (spell save DC 13, 5 to hit with spell attacks). It requires no somatic or material components to cast its spells. The flameskull has the following wizard spells prepared:
\begin{DndMonsterSpells}
  \DndMonsterSpellLevel{mage hand}
   \DndMonsterSpellLevel[1][3]{magic missile}
   \DndMonsterSpellLevel[2][2]{blur}
   \DndMonsterSpellLevel[3][1]{fireball}
\end{DndMonsterSpells}
\DndMonsterSection{Actions}
\DndMonsterAction{Multiattack}
The flameskull uses Fire Ray twice.

\DndMonsterAction{Fire Ray}
\textsl{Ranged Spell Attack:} 5 to hit, range 30 ft., one target. \textsl{Hit:} 10 (3d6) fire damage.

\end{DndMonster}



\begin{DndMonster}[float=t]{Flying Sword}
\DndMonsterType{Small construct, unaligned}
\DndMonsterBasics[
armor-class = {17 (natural armor)},
hit-points = {\DndDice{5d6}},
speed = {0 ft., fly 50 ft. \textit{(hover)}}
]
\DndMonsterAbilityScores[
 str = 12,
 dex = 15,
 con = 11,
 int = 1,
 wis = 5,
 cha = 1
]
\DndMonsterDetails[
saving-throws = {Dex +4},
damage-immunities = {poison, psychic},
condition-immunities = {blinded, charmed, deafened, frightened, paralyzed, petrified, poisoned},
senses = {blindsight 60 ft. (blind beyond this radius), passive perception 7},
challenge = 1/4,
]
\DndMonsterAction{Antimagic Susceptibility}
The sword is incapacitated while in the area of an \textit{antimagic field}. If targeted by \textit{dispel magic}, the sword must succeed on a Constitution saving throw against the caster's spell save DC or fall unconscious for 1 minute.

\DndMonsterAction{False Appearance}
While the sword remains motionless and isn't flying, it is indistinguishable from a normal sword.

\DndMonsterSection{Actions}
\DndMonsterAction{Longsword}
\textsl{Melee Weapon Attack:} 3 to hit, reach 5 ft., one target. \textsl{Hit:} 5 (1d8 + 1) slashing damage.

\end{DndMonster}



\begin{DndMonster}[float=t]{Fomorian}
\DndMonsterType{Huge giant, chaotic evil}
\DndMonsterBasics[
armor-class = {14 (natural armor)},
hit-points = {\DndDice{13d12 + 65}},
speed = {30 ft.}
]
\DndMonsterAbilityScores[
 str = 23,
 dex = 10,
 con = 20,
 int = 9,
 wis = 14,
 cha = 6
]
\DndMonsterDetails[
skills = {Perception +8, Stealth +3},
senses = {darkvision 120 ft., passive perception 18},
languages = {Giant, Undercommon},
challenge = 8,
]
\DndMonsterSection{Actions}
\DndMonsterAction{Multiattack}
The fomorian attacks twice with its greatclub or makes one greatclub attack and uses Evil Eye once.

\DndMonsterAction{Greatclub}
\textsl{Melee Weapon Attack:} 9 to hit, reach 15 ft., one target. \textsl{Hit:} 19 (3d8 + 6) bludgeoning damage.

\DndMonsterAction{Evil Eye}
The fomorian magically forces a creature it can see within 60 feet of it to make a DC 14 Charisma saving throw. The creature takes 27 (6d8) psychic damage on a failed save, or half as much damage on a successful one.

\DndMonsterAction{Curse of the Evil Eye (Recharges after a Short or Long Rest)}
With a stare, the fomorian uses Evil Eye, but on a failed save, the creature is also cursed with magical deformities. While deformed, the creature has its speed halved and has disadvantage on ability checks, saving throws, and attacks based on Strength or Dexterity.

The transformed creature can repeat the saving throw whenever it finishes a long rest, ending the effect on a success.

\end{DndMonster}



\begin{DndMonster}[float=t]{Ghost}
\DndMonsterType{Medium undead, any alignment}
\DndMonsterBasics[
armor-class = {11},
hit-points = {\DndDice{10d8}},
speed = {0 ft., fly 40 ft. \textit{(hover)}}
]
\DndMonsterAbilityScores[
 str = 7,
 dex = 13,
 con = 10,
 int = 10,
 wis = 12,
 cha = 17
]
\DndMonsterDetails[
damage-resistances = {acid; fire; lightning; thunder; bludgeoning, piercing, slashing from nonmagical attacks},
damage-immunities = {cold, necrotic, poison},
condition-immunities = {charmed, exhaustion, frightened, grappled, paralyzed, petrified, poisoned, prone, restrained},
senses = {darkvision 60 ft., passive perception 11},
languages = {any languages it knew in life},
challenge = 4,
]
\DndMonsterAction{Ethereal Sight}
The ghost can see 60 feet into the Ethereal Plane when it is on the Material Plane, and vice versa.

\DndMonsterAction{Incorporeal Movement}
The ghost can move through other creatures and objects as if they were difficult terrain. It takes 5 (1d10) force damage if it ends its turn inside an object.

\DndMonsterSection{Actions}
\DndMonsterAction{Withering Touch}
\textsl{Melee Weapon Attack:} 5 to hit, reach 5 ft., one target. \textsl{Hit:} 17 (4d6 + 3) necrotic damage.

\DndMonsterAction{Etherealness}
The ghost enters the Ethereal Plane from the Material Plane, or vice versa. It is visible on the Material Plane while it is in the Border Ethereal, and vice versa, yet it can't affect or be affected by anything on the other plane.

\DndMonsterAction{Horrifying Visage}
Each non-undead creature within 60 feet of the ghost that can see it must succeed on a DC 13 Wisdom saving throw or be frightened for 1 minute. If the save fails by 5 or more, the target also ages 1d4 × 10 years. A frightened target can repeat the saving throw at the end of each of its turns, ending the frightened condition on itself on a success. If a target's saving throw is successful or the effect ends for it, the target is immune to this ghost's Horrifying Visage for the next 24 hours. The aging effect can be reversed with a  \textit{greater restoration} spell, but only within 24 hours of it occurring.

\DndMonsterAction{Possession (Recharge 6)}
One humanoid that the ghost can see within 5 feet of it must succeed on a DC 13 Charisma saving throw or be possessed by the ghost; the ghost then disappears, and the target is incapacitated and loses control of its body. The ghost now controls the body but doesn't deprive the target of awareness. The ghost can't be targeted by any attack, spell, or other effect, except ones that turn undead, and it retains its alignment, Intelligence, Wisdom, Charisma, and immunity to being charmed and frightened. It otherwise uses the possessed target's statistics, but doesn't gain access to the target's knowledge, class features, or proficiencies.

The possession lasts until the body drops to 0 hit points, the ghost ends it as a bonus action, or the ghost is turned or forced out by an effect like the \textit{dispel evil and good} spell. When the possession ends, the ghost reappears in an unoccupied space within 5 feet of the body. The target is immune to this ghost's Possession for 24 hours after succeeding on the saving throw or after the possession ends.

\end{DndMonster}



\begin{DndMonster}[float=t]{Ghoul}
\DndMonsterType{Medium undead, chaotic evil}
\DndMonsterBasics[
armor-class = {12},
hit-points = {\DndDice{5d8}},
speed = {30 ft.}
]
\DndMonsterAbilityScores[
 str = 13,
 dex = 15,
 con = 10,
 int = 7,
 wis = 10,
 cha = 6
]
\DndMonsterDetails[
damage-immunities = {poison},
condition-immunities = {charmed, exhaustion, poisoned},
senses = {darkvision 60 ft., passive perception 10},
languages = {Common},
challenge = 1,
]
\DndMonsterSection{Actions}
\DndMonsterAction{Bite}
\textsl{Melee Weapon Attack:} 2 to hit, reach 5 ft., one creature. \textsl{Hit:} 9 (2d6 + 2) piercing damage.

\DndMonsterAction{Claws}
\textsl{Melee Weapon Attack:} 4 to hit, reach 5 ft., one target. \textsl{Hit:} 7 (2d4 + 2) slashing damage. If the target is a creature other than an elf or undead, it must succeed on a DC 10 Constitution saving throw or be paralyzed for 1 minute. The target can repeat the saving throw at the end of each of its turns, ending the effect on itself on a success.

\end{DndMonster}



\begin{DndMonster}[float=t]{Giant Lizard}
\DndMonsterType{Large beast, unaligned}
\DndMonsterBasics[
armor-class = {12 (natural armor)},
hit-points = {\DndDice{3d10 + 3}},
speed = {30 ft., climb 30 ft.}
]
\DndMonsterAbilityScores[
 str = 15,
 dex = 12,
 con = 13,
 int = 2,
 wis = 10,
 cha = 5
]
\DndMonsterDetails[
senses = {darkvision 30 ft., passive perception 10},
challenge = 1/4,
]
\DndMonsterSection{Actions}
\DndMonsterAction{Bite}
\textsl{Melee Weapon Attack:} 4 to hit, reach 5 ft., one target. \textsl{Hit:} 6 (1d8 + 2) piercing damage.

\end{DndMonster}



\begin{DndMonster}[float=t]{Giant Shark}
\DndMonsterType{Huge beast, unaligned}
\DndMonsterBasics[
armor-class = {13 (natural armor)},
hit-points = {\DndDice{11d12 + 55}},
speed = {0 ft., swim 50 ft.}
]
\DndMonsterAbilityScores[
 str = 23,
 dex = 11,
 con = 21,
 int = 1,
 wis = 10,
 cha = 5
]
\DndMonsterDetails[
skills = {Perception +3},
senses = {blindsight 60 ft., passive perception 13},
challenge = 5,
]
\DndMonsterAction{Blood Frenzy}
The shark has advantage on melee attack rolls against any creature that doesn't have all its hit points.

\DndMonsterAction{Water Breathing}
The shark can breathe only underwater.

\DndMonsterSection{Actions}
\DndMonsterAction{Bite}
\textsl{Melee Weapon Attack:} 9 to hit, reach 5 ft., one target. \textsl{Hit:} 22 (3d10 + 6) piercing damage.

\end{DndMonster}



\begin{DndMonster}[float=t]{Giant Spider}
\DndMonsterType{Large beast, unaligned}
\DndMonsterBasics[
armor-class = {14 (natural armor)},
hit-points = {\DndDice{4d10 + 4}},
speed = {30 ft., climb 30 ft.}
]
\DndMonsterAbilityScores[
 str = 14,
 dex = 16,
 con = 12,
 int = 2,
 wis = 11,
 cha = 4
]
\DndMonsterDetails[
skills = {Stealth +7},
senses = {blindsight 10 ft., darkvision 60 ft., passive perception 10},
challenge = 1,
]
\DndMonsterAction{Spider Climb}
The spider can climb difficult surfaces, including upside down on ceilings, without needing to make an ability check.

\DndMonsterAction{Web Sense}
While in contact with a web, the spider knows the exact location of any other creature in contact with the same web.

\DndMonsterAction{Web Walker}
The spider ignores movement restrictions caused by webbing.

\DndMonsterSection{Actions}
\DndMonsterAction{Bite}
\textsl{Melee Weapon Attack:} 5 to hit, reach 5 ft., one creature. \textsl{Hit:} 7 (1d8 + 3) piercing damage, and the target must make a DC 11 Constitution saving throw, taking 9 (2d8) poison damage on a failed save, or half as much damage on a successful one. If the poison damage reduces the target to 0 hit points, the target is stable but poisoned for 1 hour, even after regaining hit points, and is paralyzed while poisoned in this way.

\DndMonsterAction{Web (Recharge 5\textendash 6)}
\textsl{Ranged Weapon Attack:} 5 to hit, range 30/60 ft., one creature. \textsl{Hit:} The target is restrained by webbing. As an action, the restrained target can make a DC 12 Strength check, bursting the webbing on a success. The webbing can also be attacked and destroyed (AC 10; hp 5; vulnerability to fire damage; immunity to bludgeoning, poison, and psychic damage).

\end{DndMonster}



\begin{DndMonster}[float=t]{Githyanki Knight}
\DndMonsterType{Medium humanoid (gith), lawful evil}
\DndMonsterBasics[
armor-class = {18 (\textit{plate armor})},
hit-points = {\DndDice{14d8 + 28}},
speed = {30 ft.}
]
\DndMonsterAbilityScores[
 str = 16,
 dex = 14,
 con = 15,
 int = 14,
 wis = 14,
 cha = 15
]
\DndMonsterDetails[
saving-throws = {Con +5, Int +5, Wis +5},
languages = {Gith},
challenge = 8,
]
\DndMonsterAction{Innate Spellcasting (Psionics)}
The githyanki's innate spellcasting ability is Intelligence (spell save DC 13, 5 to hit with spell attacks). It can innately cast the following spells, requiring no components:
\begin{DndMonsterSpells}
  \DndInnateSpellLevel{mage hand}
  \DndInnateSpellLevel[3]{jump}
  \DndInnateSpellLevel[1]{plane shift}
\end{DndMonsterSpells}
\DndMonsterSection{Actions}
\DndMonsterAction{Multiattack}
The githyanki makes two silver greatsword attacks.

\DndMonsterAction{Silver Greatsword}
\textsl{Melee Weapon Attack:} 9 to hit, reach 5 ft., one target. \textsl{Hit:} 13 (2d6 + 6) slashing damage plus 10 (3d6) psychic damage. This is a magic weapon attack. On a critical hit against a target in an astral body (as with the \textit{astral projection} spell), the githyanki can cut the silvery cord that tethers the target to its material body, instead of dealing damage.

\end{DndMonster}



\begin{DndMonster}[float=t]{Glabrezu}
\DndMonsterType{Large fiend (demon), chaotic evil}
\DndMonsterBasics[
armor-class = {17 (natural armor)},
hit-points = {\DndDice{15d10 + 75}},
speed = {40 ft.}
]
\DndMonsterAbilityScores[
 str = 20,
 dex = 15,
 con = 21,
 int = 19,
 wis = 17,
 cha = 16
]
\DndMonsterDetails[
saving-throws = {Str +9, Con +9, Wis +7, Cha +7},
damage-resistances = {cold; fire; lightning; bludgeoning, piercing, slashing from nonmagical attacks},
damage-immunities = {poison},
condition-immunities = {poisoned},
senses = {truesight 120 ft., passive perception 13},
languages = {Abyssal, telepathy 120 ft.},
challenge = 9,
]
\DndMonsterAction{Magic Resistance}
The glabrezu has advantage on saving throws against spells and other magical effects.

\DndMonsterAction{Innate Spellcasting}
The glabrezu's spellcasting ability is Intelligence (spell save DC 16). The glabrezu can innately cast the following spells, requiring no material components:
\begin{DndMonsterSpells}
  \DndInnateSpellLevel{darkness}
  \DndInnateSpellLevel[1]{confusion}
\end{DndMonsterSpells}
\DndMonsterSection{Actions}
\DndMonsterAction{Multiattack}
The glabrezu makes four attacks: two with its pincers and two with its fists. Alternatively, it makes two attacks with its pincers and casts one spell.

\DndMonsterAction{Pincer}
\textsl{Melee Weapon Attack:} 9 to hit, reach 10 ft., one target. \textsl{Hit:} 16 (2d10 + 5) bludgeoning damage. If the target is a Medium or smaller creature, it is grappled (escape DC 15). The glabrezu has two pincers, each of which can grapple only one target.

\DndMonsterAction{Fist}
\textsl{Melee Weapon Attack:} 9 to hit, reach 5 ft., one target. \textsl{Hit:} 7 (2d4 + 2) bludgeoning damage.

\end{DndMonster}



\begin{DndMonster}[float=t]{Glaive}
\DndMonsterType{Medium construct (warforged), chaotic evil}
\DndMonsterBasics[
armor-class = {16 (natural armor)},
hit-points = {\DndDice{22d8 + 88}},
speed = {30 ft. (50 ft. with overdrive)}
]
\DndMonsterAbilityScores[
 str = 20,
 dex = 16,
 con = 19,
 int = 11,
 wis = 16,
 cha = 9
]
\DndMonsterDetails[
saving-throws = {Str +9, Dex +7, Wis +7},
skills = {Athletics +9, Perception +7, Stealth +7, Survival +7},
damage-resistances = {poison},
condition-immunities = {charmed, exhaustion, frightened, poisoned},
languages = {Common},
challenge = 11,
]
\DndMonsterAction{Heatsink}
When Glaive takes cold damage, her Overdrive immediately recharges.

\DndMonsterAction{Pack Tactics}
Glaive has advantage on attack rolls if at least one ally is within 5 feet of the creature she's attacking and the ally doesn't have the incapacitated condition.

\DndMonsterSection{Actions}
\DndMonsterAction{Multiattack}
Glaive makes two Spiked Glaive attacks and two Serrated Bolt attacks.

\DndMonsterAction{Spiked Glaive}
\textsl{Melee Weapon Attack:} 9 to hit, reach 10 ft., one target. \textsl{Hit:} 10 (1d10 + 5) piercing or slashing damage, or 14 (1d10 + 9) piercing or slashing damage if Glaive is in overdrive.

\DndMonsterAction{Serrated Bolt}
\textsl{Ranged Weapon Attack:} 7 to hit, range 60 ft., one target. \textsl{Hit:} 13 (3d6 + 3) piercing damage. If Glaive has advantage on the attack roll, the serrated bolt lodges in the target, and the target's speed is reduced by 10 feet until the serrated bolt is removed. A target's speed can be reduced by only one serrated bolt at a time. A creature can use its action to remove a serrated bolt lodged in itself or another creature within its reach; when the bolt is removed from a creature, that creature takes 5 (2d4) slashing damage.

\DndMonsterSection{Bonus Actions}
\DndMonsterAction{Overdrive (Recharges after Finishing a Short or Long Rest)}
Glaive enters a state of overdrive that lasts for 1 minute or until she has the incapacitated condition. While in overdrive, Glaive gains the following benefits:


\begin{itemize}
\item Glaive has advantage on Strength checks and Strength saving throws.
\item When Glaive makes a melee weapon attack, she gains a +4 bonus to the damage roll.
\item Glaive's speed increases to 50 feet.
\end{itemize}

\DndMonsterSection{Reactions}
\DndMonsterAction{Self-Preservation}
In response to being hit by a weapon attack, Glaive reduces the damage by 11 (2d10).

\end{DndMonster}



\begin{DndMonster}[float=t]{Goristro}
\DndMonsterType{Huge fiend (demon), chaotic evil}
\DndMonsterBasics[
armor-class = {19 (natural armor)},
hit-points = {\DndDice{23d12 + 161}},
speed = {40 ft.}
]
\DndMonsterAbilityScores[
 str = 25,
 dex = 11,
 con = 25,
 int = 6,
 wis = 13,
 cha = 14
]
\DndMonsterDetails[
saving-throws = {Str +13, Dex +6, Con +13, Wis +7},
skills = {Perception +7},
damage-resistances = {cold; fire; lightning; bludgeoning, piercing, slashing from nonmagical attacks},
damage-immunities = {poison},
condition-immunities = {poisoned},
senses = {darkvision 120 ft., passive perception 17},
languages = {Abyssal},
challenge = 17,
]
\DndMonsterAction{Charge}
If the goristro moves at least 15 feet straight toward a target and then hits it with a gore attack on the same turn, the target takes an extra 38 (7d10) piercing damage. If the target is a creature, it must succeed on a DC 21 Strength saving throw or be pushed up to 20 feet away and knocked prone.

\DndMonsterAction{Labyrinthine Recall}
The goristro can perfectly recall any path it has traveled.

\DndMonsterAction{Magic Resistance}
The goristro has advantage on saving throws against spells and other magical effects.

\DndMonsterAction{Siege Monster}
The goristro deals double damage to objects and structures.

\DndMonsterSection{Actions}
\DndMonsterAction{Multiattack}
The goristro makes three attacks: two with its fists and one with its hoof.

\DndMonsterAction{Fist}
\textsl{Melee Weapon Attack:} 13 to hit, reach 10 ft., one target. \textsl{Hit:} 20 (3d8 + 7) bludgeoning damage.

\DndMonsterAction{Hoof}
\textsl{Melee Weapon Attack:} 13 to hit, reach 5 ft., one target. \textsl{Hit:} 23 (3d10 + 7) bludgeoning damage. If the target is a creature, it must succeed on a DC 21 Strength saving throw or be knocked prone.

\DndMonsterAction{Gore}
\textsl{Melee Weapon Attack:} 13 to hit, reach 10 ft., one target. \textsl{Hit:} 45 (7d10 + 7) piercing damage.

\end{DndMonster}



\begin{DndMonster}[float=t]{Granite Juggernaut}
\DndMonsterType{Large construct, unaligned}
\DndMonsterBasics[
armor-class = {14 (natural armor)},
hit-points = {\DndDice{15d10 + 75}},
speed = {30 ft. (in a straight line)}
]
\DndMonsterAbilityScores[
 str = 23,
 dex = 1,
 con = 20,
 int = 2,
 wis = 11,
 cha = 3
]
\DndMonsterDetails[
damage-immunities = {poison; psychic; bludgeoning, piercing, slashing from nonmagical attacks that aren't adamantine},
condition-immunities = {charmed, exhaustion, frightened, paralyzed, petrified, poisoned, prone},
senses = {blindsight 120 ft., passive perception 10},
challenge = 12,
]
\DndMonsterAction{Magic Resistance}
The juggernaut has advantage on saving throws against spells and other magical effects.

\DndMonsterAction{Siege Monster}
The juggernaut deals double damage to objects and structures.

\DndMonsterSection{Actions}
\DndMonsterAction{Slam}
\textsl{Melee Weapon Attack:} 10 to hit, reach 5 ft., one target. \textsl{Hit:} 11 (1d10 + 6) bludgeoning damage, and if the target is a Large or smaller creature, it must succeed on a DC 18 Strength saving throw or have the prone condition.

\DndMonsterSection{Bonus Actions}
\DndMonsterAction{Devastating Roll}
The juggernaut moves up to its speed. During this movement, the juggernaut can move through the spaces of creatures with the prone condition. When the juggernaut enters the space of a prone creature for the first time during this movement, the creature must make a DC 18 Dexterity saving throw, taking 55 (10d10) bludgeoning damage on a failed save or half as much damage on a successful one.

\end{DndMonster}



\begin{DndMonster}[float=t]{Green Abishai}
\DndMonsterType{Medium fiend (devil), lawful evil}
\DndMonsterBasics[
armor-class = {18 (natural armor)},
hit-points = {\DndDice{26d8 + 78}},
speed = {30 ft., fly 40 ft.}
]
\DndMonsterAbilityScores[
 str = 12,
 dex = 17,
 con = 16,
 int = 17,
 wis = 12,
 cha = 19
]
\DndMonsterDetails[
saving-throws = {Int +8, Cha +9},
skills = {Deception +9, Insight +6, Perception +6, Persuasion +9},
damage-resistances = {cold; bludgeoning, piercing, slashing from nonmagical attacks that aren't silvered},
damage-immunities = {fire, poison},
condition-immunities = {poisoned},
senses = {darkvision 120 ft., passive perception 16},
languages = {Draconic, Infernal, telepathy 120 ft.},
challenge = 15,
]
\DndMonsterAction{Devil's Sight}
Magical darkness doesn't impede the abishai's darkvision.

\DndMonsterAction{Magic Resistance}
The abishai has advantage on saving throws against spells and other magical effects.

\DndMonsterAction{Spellcasting}
The abishai casts one of the following spells, requiring no material components and using Charisma as the spellcasting ability (spell save DC 17):
\begin{DndMonsterSpells}
  \DndInnateSpellLevel{alter self}
  \DndInnateSpellLevel[3]{charm person}
  \DndInnateSpellLevel[1]{confusion}
\end{DndMonsterSpells}
\DndMonsterSection{Actions}
\DndMonsterAction{Multiattack}
The abishai makes two Fiendish Claw attacks, or it makes one Fiendish Claw attack and uses Spellcasting.

\DndMonsterAction{Fiendish Claw}
\textsl{Melee Weapon Attack:} 8 to hit, reach 5 ft., one target. \textsl{Hit:} 12 (2d8 + 3) force damage. If the target is a creature, it must succeed on a DC 16 Constitution saving throw or take 16 (3d10) poison damage and become poisoned for 1 minute. The poisoned target can repeat the saving throw at the end of each of its turns, ending the effect on itself on a success.

\end{DndMonster}



\begin{DndMonster}[float=t]{Green Slaad}
\DndMonsterType{Large aberration (shapechanger), chaotic neutral}
\DndMonsterBasics[
armor-class = {16 (natural armor)},
hit-points = {\DndDice{15d10 + 45}},
speed = {30 ft.}
]
\DndMonsterAbilityScores[
 str = 18,
 dex = 15,
 con = 16,
 int = 11,
 wis = 8,
 cha = 12
]
\DndMonsterDetails[
skills = {Arcana +3, Perception +2},
damage-resistances = {acid, cold, fire, lightning, thunder},
senses = {blindsight 30 ft., darkvision 60 ft., passive perception 12},
languages = {Slaad, telepathy 60 ft.},
challenge = 8,
]
\DndMonsterAction{Shapechanger}
The slaad can use its action to polymorph into a Small or Medium humanoid, or back into its true form. Its statistics, other than its size, are the same in each form. Any equipment it is wearing or carrying isn't transformed. It reverts to its true form if it dies.

\DndMonsterAction{Magic Resistance}
The slaad has advantage on saving throws against spells and other magical effects.

\DndMonsterAction{Regeneration}
The slaad regains 10 hit points at the start of its turn if it has at least 1 hit point.

\DndMonsterAction{Innate Spellcasting}
The slaad's innate spellcasting ability is Charisma (spell save DC 12). The slaad can innately cast the following spells, requiring no material components:
\begin{DndMonsterSpells}
  \DndInnateSpellLevel{detect magic}
  \DndInnateSpellLevel[]{fireball}
  \DndInnateSpellLevel[2]{fear}
\end{DndMonsterSpells}
\DndMonsterSection{Actions}
\DndMonsterAction{Multiattack}
The slaad makes three attacks: one with its bite and two with its claws or staff. Alternatively, it uses its Hurl Flame twice.

\DndMonsterAction{Bite (Slaad Form Only)}
\textsl{Melee Weapon Attack:} 7 to hit, reach 5 ft., one target. \textsl{Hit:} 11 (2d6 + 4) piercing damage.

\DndMonsterAction{Claw (Slaad Form Only)}
\textsl{Melee Weapon Attack:} 7 to hit, reach 5 ft., one target. \textsl{Hit:} 7 (1d6 + 4) slashing damage.

\DndMonsterAction{Staff}
\textsl{Melee Weapon Attack:} 7 to hit, reach 5 ft., one target. \textsl{Hit:} 11 (2d6 + 4) bludgeoning damage.

\DndMonsterAction{Hurl Flame}
\textsl{Ranged Spell Attack:} 4 to hit, range 60 ft., one target. \textsl{Hit:} 10 (3d6) fire damage. The fire ignites flammable objects that aren't being worn or carried.

\end{DndMonster}



\begin{DndMonster}[float=t]{Grell}
\DndMonsterType{Medium aberration, neutral evil}
\DndMonsterBasics[
armor-class = {12},
hit-points = {\DndDice{10d8 + 10}},
speed = {10 ft., fly 30 ft. \textit{(hover)}}
]
\DndMonsterAbilityScores[
 str = 15,
 dex = 14,
 con = 13,
 int = 12,
 wis = 11,
 cha = 9
]
\DndMonsterDetails[
skills = {Perception +4, Stealth +6},
damage-immunities = {lightning},
condition-immunities = {blinded, prone},
senses = {blindsight 60 ft. (blind beyond this radius), passive perception 14},
languages = {Grell},
challenge = 3,
]
\DndMonsterSection{Actions}
\DndMonsterAction{Multiattack}
The grell makes two attacks: one with its tentacles and one with its beak.

\DndMonsterAction{Tentacles}
\textsl{Melee Weapon Attack:} 4 to hit, reach 10 ft., one creature. \textsl{Hit:} 7 (1d10 + 2) piercing damage, and the target must succeed on a DC 11 Constitution saving throw or be poisoned for 1 minute. The poisoned target is paralyzed, and it can repeat the saving throw at the end of each of its turns, ending the effect on a success.

The target is also grappled (escape DC 15). If the target is Medium or smaller, it is also restrained until this grapple ends. While grappling the target, the grell has advantage on attack rolls against it and can 't use this attack against other targets. When the grell moves, any Medium or smaller target it is grappling moves with it.

\DndMonsterAction{Beak}
\textsl{Melee Weapon Attack:} 4 to hit, reach 5 ft., one target. \textsl{Hit:} 7 (2d4 + 2) piercing damage.

\end{DndMonster}


\clearpage

\begin{DndMonster}[float=t]{Grimlock}
\DndMonsterType{Medium humanoid (grimlock), neutral evil}
\DndMonsterBasics[
armor-class = {11},
hit-points = {\DndDice{2d8 + 2}},
speed = {30 ft.}
]
\DndMonsterAbilityScores[
 str = 16,
 dex = 12,
 con = 12,
 int = 9,
 wis = 8,
 cha = 6
]
\DndMonsterDetails[
skills = {Athletics +5, Perception +3, Stealth +3},
condition-immunities = {blinded},
senses = {blindsight 30 ft. or 10 ft. while deafened (blind beyond this radius), passive perception 13},
languages = {Undercommon},
challenge = 1/4,
]
\DndMonsterAction{Blind Senses}
The grimlock can't use its blindsight while deafened and unable to smell.

\DndMonsterAction{Keen Hearing and Smell}
The grimlock has advantage on Wisdom (Perception) checks that rely on hearing or smell.

\DndMonsterAction{Stone Camouflage}
The grimlock has advantage on Dexterity (Stealth) checks made to hide in rocky terrain.

\DndMonsterSection{Actions}
\DndMonsterAction{Spiked Bone Club}
\textsl{Melee Weapon Attack:} 5 to hit, reach 5 ft., one target. \textsl{Hit:} 5 (1d4 + 3) bludgeoning damage plus 2 (1d4) piercing damage.

\end{DndMonster}



\begin{DndMonster}[float=t]{Guardian Naga}
\DndMonsterType{Large monstrosity, lawful good}
\DndMonsterBasics[
armor-class = {18 (natural armor)},
hit-points = {\DndDice{15d10 + 45}},
speed = {40 ft.}
]
\DndMonsterAbilityScores[
 str = 19,
 dex = 18,
 con = 16,
 int = 16,
 wis = 19,
 cha = 18
]
\DndMonsterDetails[
saving-throws = {Dex +8, Con +7, Int +7, Wis +8, Cha +8},
damage-immunities = {poison},
condition-immunities = {charmed, poisoned},
senses = {darkvision 60 ft., passive perception 14},
languages = {Celestial, Common},
challenge = 10,
]
\DndMonsterAction{Rejuvenation}
If it dies, the naga returns to life in 1d6 days and regains all its hit points. Only a \textit{wish} spell can prevent this trait from functioning.

\DndMonsterAction{Spellcasting}
The naga is an 11th-level spellcaster. Its spellcasting ability is Wisdom (spell save DC 16, 8 to hit with spell attacks), and it needs only verbal components to cast its spells. It has the following cleric spells prepared:
\begin{DndMonsterSpells}
  \DndMonsterSpellLevel{mending}
   \DndMonsterSpellLevel[1][4]{command}
   \DndMonsterSpellLevel[2][3]{calm emotions}
   \DndMonsterSpellLevel[3][3]{bestow curse}
   \DndMonsterSpellLevel[4][3]{banishment}
   \DndMonsterSpellLevel[5][2]{flame strike}
   \DndMonsterSpellLevel[6][1]{true seeing}
\end{DndMonsterSpells}
\DndMonsterSection{Actions}
\DndMonsterAction{Bite}
\textsl{Melee Weapon Attack:} 8 to hit, reach 10 ft., one creature. \textsl{Hit:} 8 (1d8 + 4) piercing damage, and the target must make a DC 15 Constitution saving throw, taking 45 (10d8) poison damage on a failed save, or half as much damage on a successful one.

\DndMonsterAction{Spit Poison}
\textsl{Ranged Weapon Attack:} 8 to hit, range 15/30 ft., one creature. \textsl{Hit:} The target must make a DC 15 Constitution saving throw, taking 45 (10d8) poison damage on a failed save, or half as much damage on a successful one.

\end{DndMonster}



\begin{DndMonster}[float=t]{Hazvongel}
\DndMonsterType{Huge fiend (demon), chaotic evil}
\DndMonsterBasics[
armor-class = {17 (natural armor)},
hit-points = {\DndDice{25d12 + 75}},
speed = {20 ft., fly 60 ft.}
]
\DndMonsterAbilityScores[
 str = 21,
 dex = 20,
 con = 16,
 int = 12,
 wis = 15,
 cha = 11
]
\DndMonsterDetails[
saving-throws = {Con +8, Wis +7},
skills = {Perception +7},
damage-resistances = {cold; fire; lightning; bludgeoning, piercing, slashing from nonmagical attacks},
damage-immunities = {poison},
condition-immunities = {exhaustion, poisoned},
senses = {blindsight 60 ft., darkvision 120 ft., passive perception 17},
languages = {Abyssal, telepathy 120 ft.},
challenge = 14,
]
\DndMonsterAction{Magic Resistance}
The hazvongel has advantage on saving throws against spells and other magical effects.

\DndMonsterSection{Actions}
\DndMonsterAction{Multiattack}
The hazvongel makes three Talon attacks.

\DndMonsterAction{Talon}
\textsl{Melee Weapon Attack:} 10 to hit, reach 15 ft., one target. \textsl{Hit:} 18 (3d8 + 5) piercing damage.

\DndMonsterAction{Blood Barrage (Recharge 5\textendash 6)}
The hazvongel launches a spray of blood in a 90-foot cone. Each creature in that area must make a DC 18 Dexterity saving throw, taking 27 (6d8) necrotic damage on a failed save or half as much damage on a successful one.

\end{DndMonster}



\begin{DndMonster}[float=t]{Helmed Horror}
\DndMonsterType{Medium construct, unaligned}
\DndMonsterBasics[
armor-class = {20 (\textit{plate armor})},
hit-points = {\DndDice{8d8 + 24}},
speed = {30 ft., fly 30 ft.}
]
\DndMonsterAbilityScores[
 str = 18,
 dex = 13,
 con = 16,
 int = 10,
 wis = 10,
 cha = 10
]
\DndMonsterDetails[
skills = {Perception +4},
damage-resistances = {bludgeoning, piercing, slashing from nonmagical attacks that aren't adamantine},
damage-immunities = {force, necrotic, poison},
condition-immunities = {blinded, charmed, deafened, frightened, paralyzed, petrified, poisoned, stunned},
senses = {blindsight 60 ft. (blind beyond this radius), passive perception 14},
languages = {understands the languages of its creator but can't speak},
challenge = 4,
]
\DndMonsterAction{Magic Resistance}
The helmed horror has advantage on saving throws against spells and other magical effects.

\DndMonsterAction{Spell Immunity}
The helmed horror is immune to three spells chosen by its creator. Typical immunities include \textit{fireball}, \textit{heat metal}, and \textit{lightning bolt}.

\DndMonsterSection{Actions}
\DndMonsterAction{Multiattack}
The helmed horror makes two longsword attacks.

\DndMonsterAction{Longsword}
\textsl{Melee Weapon Attack:} 6 to hit, reach 5 ft., one target. \textsl{Hit:} 8 (1d8 + 4) slashing damage, or 9 (1d10 + 4) slashing damage if used with two hands.

\end{DndMonster}



\begin{DndMonster}[float*=b,width=\textwidth + 8pt]{Hertilod}
\begin{multicols}{2}
\DndMonsterType{Huge monstrosity, unaligned}
\DndMonsterBasics[
armor-class = {14},
hit-points = {\DndDice{21d12 + 105}},
speed = {50 ft., climb 50 ft.}
]
\DndMonsterAbilityScores[
 str = 23,
 dex = 18,
 con = 20,
 int = 3,
 wis = 15,
 cha = 10
]
\DndMonsterDetails[
saving-throws = {Str +12, Dex +10},
skills = {Perception +8},
damage-resistances = {necrotic; bludgeoning, piercing, slashing from nonmagical attacks},
damage-immunities = {poison},
condition-immunities = {poisoned},
senses = {blindsight 30 ft., tremorsense 60 ft., passive perception 18},
challenge = 17,
]
\DndMonsterAction{Legendary Resistances (3/Day)}
If the hertilod fails a saving throw, it can choose to succeed instead.

\DndMonsterAction{Magic Resistance}
The hertilod has advantage on saving throws against spells and other magical effects.

\DndMonsterAction{Shock Susceptibility}
If the hertilod takes lightning damage, its speed is halved until the end of its next turn, and it must succeed on a DC 15 Constitution saving throw or immediately regurgitate all swallowed creatures, each of which lands in a space within 10 feet of the hertilod and has the prone condition.

\DndMonsterAction{Spider Climb}
The hertilod can climb difficult surfaces, including upside down on ceilings, without needing to make an ability check.

\DndMonsterSection{Actions}
\DndMonsterAction{Multiattack}
The hertilod makes one Bite attack and two Claw attacks.

\DndMonsterAction{Bite}
\textsl{Melee Weapon Attack:} 12 to hit, reach 10 ft., one target. \textsl{Hit:} 15 (2d8 + 6) piercing damage plus 13 (2d12) poison damage. If the target is a Large or smaller creature, it must succeed on a DC 20 Strength saving throw or be swallowed by the hertilod. A swallowed creature has the blinded and restrained conditions, and it has total cover against attacks and other effects outside the hertilod. At the start of each of the hertilod's turns, each swallowed creature takes 13 (2d12) poison damage from the poisonous secretion in the hertilod's gullet.

The hertilod's gullet can hold up to two creatures at a time. If the hertilod takes 40 damage or more on a single turn from a swallowed creature, the hertilod must succeed on a DC 15 Constitution saving throw at the end of that turn or regurgitate all swallowed creatures, each of which lands in a space within 10 feet of the hertilod and has the prone condition. If the hertilod dies, a swallowed creature is no longer restrained and can escape from the corpse by using 10 feet of movement, exiting with the prone condition.

\DndMonsterAction{Claw}
\textsl{Melee Weapon Attack:} 12 to hit, reach 5 ft., one target. \textsl{Hit:} 17 (2d10 + 6) slashing damage.

\DndMonsterSection{Legendary Actions}
Hertilod can take 3 legendary actions, choosing from the options below. Only one legendary action can be used at a time and only at the end of another creature's turn. Hertilod regains spent legendary actions at the start of its turn.

\begin{DndMonsterLegendaryActions}
\DndMonsterLegendaryAction{Sprint}{The hertilod moves up to its speed. This movement doesn't provoke opportunity attacks.
}
\DndMonsterLegendaryAction{Feed (Costs 2 Actions)}{The hertilod drains life from the creatures in its gullet to bolster itself. Each creature in the hertilod's gullet takes 10 (3d6) necrotic damage, and the hertilod regains a number of hit points equal to the damage.
}
\end{DndMonsterLegendaryActions}
\end{multicols}
\end{DndMonster}



\begin{DndMonster}[float=t]{Hezrou}
\DndMonsterType{Large fiend (demon), chaotic evil}
\DndMonsterBasics[
armor-class = {16 (natural armor)},
hit-points = {\DndDice{13d10 + 65}},
speed = {30 ft.}
]
\DndMonsterAbilityScores[
 str = 19,
 dex = 17,
 con = 20,
 int = 5,
 wis = 12,
 cha = 13
]
\DndMonsterDetails[
saving-throws = {Str +7, Con +8, Wis +4},
damage-resistances = {cold; fire; lightning; bludgeoning, piercing, slashing from nonmagical attacks},
damage-immunities = {poison},
condition-immunities = {poisoned},
senses = {darkvision 120 ft., passive perception 11},
languages = {Abyssal, telepathy 120 ft.},
challenge = 8,
]
\DndMonsterAction{Magic Resistance}
The hezrou has advantage on saving throws against spells and other magical effects.

\DndMonsterAction{Stench}
Any creature that starts its turn within 10 feet of the hezrou must succeed on a DC 14 Constitution saving throw or be poisoned until the start of its next turn. On a successful saving throw, the creature is immune to the hezrou's stench for 24 hours.

\DndMonsterSection{Actions}
\DndMonsterAction{Multiattack}
The hezrou makes three attacks: one with its bite and two with its claws.

\DndMonsterAction{Bite}
\textsl{Melee Weapon Attack:} 7 to hit, reach 5 ft., one target. \textsl{Hit:} 15 (2d10 + 4) piercing damage.

\DndMonsterAction{Claws}
\textsl{Melee Weapon Attack:} 7 to hit, reach 5 ft., one target. \textsl{Hit:} 11 (2d6 + 4) slashing damage.

\end{DndMonster}



\begin{DndMonster}[float=t]{Horned Devil}
\DndMonsterType{Large fiend (devil), lawful evil}
\DndMonsterBasics[
armor-class = {18 (natural armor)},
hit-points = {\DndDice{17d10 + 85}},
speed = {20 ft., fly 60 ft.}
]
\DndMonsterAbilityScores[
 str = 22,
 dex = 17,
 con = 21,
 int = 12,
 wis = 16,
 cha = 17
]
\DndMonsterDetails[
saving-throws = {Str +10, Dex +7, Wis +7, Cha +7},
damage-resistances = {cold; bludgeoning, piercing, slashing from nonmagical attacks that aren't silvered},
damage-immunities = {fire, poison},
condition-immunities = {poisoned},
senses = {darkvision 120 ft., passive perception 13},
languages = {Infernal, telepathy 120 ft.},
challenge = 11,
]
\DndMonsterAction{Devil's Sight}
Magical darkness doesn't impede the devil's darkvision.

\DndMonsterAction{Magic Resistance}
The devil has advantage on saving throws against spells and other magical effects.

\DndMonsterSection{Actions}
\DndMonsterAction{Multiattack}
The devil makes three melee attacks: two with its fork and one with its tail. It can use Hurl Flame in place of any melee attack.

\DndMonsterAction{Fork}
\textsl{Melee Weapon Attack:} 10 to hit, reach 10 ft., one target. \textsl{Hit:} 15 (2d8 + 6) piercing damage.

\DndMonsterAction{Tail}
\textsl{Melee Weapon Attack:} 10 to hit, reach 10 ft., one target. \textsl{Hit:} 10 (1d8 + 6) piercing damage. If the target is a creature other than an undead or a construct, it must succeed on a DC 17 Constitution saving throw or lose 10 (3d6) hit points at the start of each of its turns due to an infernal wound. Each time the devil hits the wounded target with this attack, the damage dealt by the wound increases by 10 (3d6). Any creature can take an action to stanch the wound with a successful DC 12 Wisdom (Medicine) check. The wound also closes if the target receives magical healing.

\DndMonsterAction{Hurl Flame}
\textsl{Ranged Spell Attack:} 7 to hit, range 150 ft., one target. \textsl{Hit:} 14 (4d6) fire damage. If the target is a flammable object that isn't being worn or carried, it also catches fire.

\end{DndMonster}



\begin{DndMonster}[float=t]{Ice Devil}
\DndMonsterType{Large fiend (devil), lawful evil}
\DndMonsterBasics[
armor-class = {18 (natural armor)},
hit-points = {\DndDice{19d10 + 76}},
speed = {40 ft.}
]
\DndMonsterAbilityScores[
 str = 21,
 dex = 14,
 con = 18,
 int = 18,
 wis = 15,
 cha = 18
]
\DndMonsterDetails[
saving-throws = {Dex +7, Con +9, Wis +7, Cha +9},
damage-resistances = {bludgeoning, piercing, slashing from nonmagical attacks that aren't silvered},
damage-immunities = {fire, poison, cold},
condition-immunities = {poisoned},
senses = {blindsight 60 ft., darkvision 120 ft., passive perception 12},
languages = {Infernal, telepathy 120 ft.},
challenge = 14,
]
\DndMonsterAction{Devil's Sight}
Magical darkness doesn't impede the devil's darkvision.

\DndMonsterAction{Magic Resistance}
The devil has advantage on saving throws against spells and other magical effects.

\DndMonsterSection{Actions}
\DndMonsterAction{Multiattack}
The devil makes three attacks: one with its bite, one with its claws, and one with its tail.

\DndMonsterAction{Bite}
\textsl{Melee Weapon Attack:} 10 to hit, reach 5 ft., one target. \textsl{Hit:} 12 (2d6 + 5) piercing damage plus 10 (3d6) cold damage.

\DndMonsterAction{Claws}
\textsl{Melee Weapon Attack:} 10 to hit, reach 5 ft., one target. \textsl{Hit:} 10 (2d4 + 5) slashing damage plus 10 (3d6) cold damage.

\DndMonsterAction{Tail}
\textsl{Melee Weapon Attack:} 10 to hit, reach 10 ft., one target. \textsl{Hit:} 12 (2d6 + 5) bludgeoning damage plus 10 (3d6) cold damage.

\DndMonsterAction{Wall of Ice (Recharge 6)}
The devil magically forms an opaque wall of ice on a solid surface it can see within 60 feet of it. The wall is 1 foot thick and up to 30 feet long and 10 feet high, or it's a hemispherical dome up to 20 feet in diameter.

When the wall appears, each creature in its space is pushed out of it by the shortest route. The creature chooses which side of the wall to end up on, unless the creature is incapacitated. The creature then makes a DC 17 Dexterity saving throw, taking 35 (10d6) cold damage on a failed save, or half as much damage on a successful one.

The wall lasts for 1 minute or until the devil is incapacitated or dies. The wall can be damaged and breached; each 10-foot section has AC 5, 30 hit points, vulnerability to fire damage, and immunity to acid, cold, necrotic, poison, and psychic damage. If a section is destroyed, it leaves behind a sheet of frigid air in the space the wall occupied. Whenever a creature finishes moving through the frigid air on a turn, willingly or otherwise, the creature must make a DC 17 Constitution saving throw, taking 17 (5d6) cold damage on a failed save, or half as much damage on a successful one. The frigid air dissipates when the rest of the wall vanishes.

\end{DndMonster}



\begin{DndMonster}[float*=b,width=\textwidth + 8pt]{Iggwilv the Witch Queen}
\begin{multicols}{2}
\DndMonsterType{Medium fey (wizard), chaotic neutral}
\DndMonsterBasics[
armor-class = {19 (\textit{robe of the archmagi})},
hit-points = {\DndDice{30d8 + 120}},
speed = {30 ft.}
]
\DndMonsterAbilityScores[
 str = 10,
 dex = 18,
 con = 18,
 int = 27,
 wis = 12,
 cha = 23
]
\DndMonsterDetails[
saving-throws = {Int +14, Wis +7, Cha +12},
skills = {Arcana +20, History +14, Nature +14},
condition-immunities = {charmed, frightened},
senses = {truesight 60 ft., passive perception 11},
languages = {Abyssal, Celestial, Common, Draconic, Elvish, Infernal, Sylvan},
challenge = 20,
]
\DndMonsterAction{Boon of Immortality}
Iggwilv is immune to any effect that would age her, and she can't die from old age.

\DndMonsterAction{Legendary Resistance (3/Day)}
If Iggwilv fails a saving throw, she can choose to succeed instead.

\DndMonsterAction{Magic Resistance}
Iggwilv has advantage on saving throws against spells and other magical effects. (This trait is bestowed by her robe of the archmagi.)

\DndMonsterAction{Special Equipment}
Iggwilv wears an \textit{amulet of the planes} and a \textit{robe of the archmagi}.

\DndMonsterAction{Spellcasting}
Iggwilv casts one of the following spells, requiring no material components and using Intelligence as the spellcasting ability (spell save DC 24, 16 to hit with spell attacks):
\begin{DndMonsterSpells}
  \DndInnateSpellLevel{detect magic}
  \DndInnateSpellLevel[3]{dispel magic}
  \DndInnateSpellLevel[1]{maze}
\end{DndMonsterSpells}
\DndMonsterSection{Actions}
\DndMonsterAction{Multiattack}
Iggwilv makes two Bewitching Bolt attacks.

\DndMonsterAction{Bewitching Bolt}
\textsl{Melee or Ranged Spell Attack:} 16 to hit, reach 5 ft. or range 120 ft., one target. \textsl{Hit:} 25 (5d6 + 8) lightning damage, and if the target is a creature, it must succeed on a DC 22 Wisdom saving throw or be charmed by Iggwilv until the start of her next turn.

\DndMonsterAction{Abyssal Rift (Recharge 5\textendash 6)}
Iggwilv opens a momentary Abyssal rift within 120 feet of her. The rift is a 20-foot-radius sphere. Each creature in that area must make a DC 22 Constitution saving throw, taking 40 (9d8) necrotic damage on a failed save, or half as much damage on a successful one. In addition, there is a 50 percent chance that 3 hezrous then appear in unoccupied spaces in the sphere. They act as Iggwilv's allies, take their turns immediately after hers, and can't summon other demons. They remain until they die or until Iggwilv dismisses them as an action.

\DndMonsterSection{Bonus Actions}
\DndMonsterAction{Fey Step}
Iggwilv teleports, along with any equipment she is wearing or carrying, to an unoccupied space she can see within 30 feet of her.

\DndMonsterSection{Reactions}
\DndMonsterAction{Negate Spell (2/Day)}
When Iggwilv sees a creature within 60 feet of her casting a spell, she tries to interrupt it. If the creature is casting a spell using a spell slot of 8th level or lower, its spell fails and has no effect. If it is casting a 9th-level spell, it must succeed on a DC 22 Intelligence saving throw, or the spells fails and has no effect.

\DndMonsterSection{Legendary Actions}
Iggwilv the Witch Queen can take 3 legendary actions, choosing from the options below. Only one legendary action can be used at a time and only at the end of another creature's turn. Iggwilv the Witch Queen regains spent legendary actions at the start of its turn.

\begin{DndMonsterLegendaryActions}
\DndMonsterLegendaryAction{Witchcraft}{Iggwilv uses Spellcasting or Fey Step.
}
\DndMonsterLegendaryAction{Dark Speech (Costs 2 Actions)}{Iggwilv utters a phrase in a forbidden language and targets one or two creatures she can see within 60 feet of her. Each target must succeed on a DC 22 Wisdom saving throw or take 11 (2d10) psychic damage and be frightened of Iggwilv for 1 minute. A target can repeat the save at the end of each of its turns, ending the effect on itself on a success and thereby becoming immune to Iggwilv's Dark Speech for 24 hours.
}
\DndMonsterLegendaryAction{Fey Beguilement (Costs 3 Actions)}{Iggwilv targets one creature she can see within 60 feet of her. The target must succeed on a DC 22 Charisma saving throw or be possessed by a fey spirit. While possessed, the target must obey Iggwilv's commands. The target can repeat the saving throw at the end of each of its turns, banishing the fey spirit and ending the effect on itself on a success.
}
\end{DndMonsterLegendaryActions}
\end{multicols}
\end{DndMonster}



\begin{DndMonster}[float=t]{Imp}
\DndMonsterType{Tiny fiend (devil), lawful evil}
\DndMonsterBasics[
armor-class = {13},
hit-points = {\DndDice{3d4 + 3}},
speed = {20 ft., fly 40 ft.}
]
\DndMonsterAbilityScores[
 str = 6,
 dex = 17,
 con = 13,
 int = 11,
 wis = 12,
 cha = 14
]
\DndMonsterDetails[
skills = {Deception +4, Insight +3, Persuasion +4, Stealth +5},
damage-resistances = {cold; bludgeoning, piercing, slashing from nonmagical attacks not made with silvered weapons},
damage-immunities = {fire, poison},
condition-immunities = {poisoned},
senses = {darkvision 120 ft., passive perception 11},
languages = {Infernal, Common},
challenge = 1,
]
\DndMonsterAction{Shapechanger}
The imp can use its action to polymorph into a beast form that resembles a rat (speed 20 ft.), a raven (20 ft., fly 60 ft.), or a spider (20 ft., climb 20 ft.), or back into its true form. Its statistics are the same in each form, except for the speed changes noted. Any equipment it is wearing or carrying isn't transformed. It reverts to its true form if it dies.

\DndMonsterAction{Devil's Sight}
Magical darkness doesn't impede the imp's darkvision.

\DndMonsterAction{Magic Resistance}
The imp has advantage on saving throws against spells and other magical effects.

\DndMonsterSection{Actions}
\DndMonsterAction{Sting (Bite in Beast Form)}
\textsl{Melee Weapon Attack:} 5 to hit, reach 5 ft., one target. \textsl{Hit:} 5 (1d4 + 3) piercing damage, and the target must make a DC 11 Constitution saving throw, taking 10 (3d6) poison damage on a failed save, or half as much damage on a successful one.

\DndMonsterAction{Invisibility}
The imp magically turns invisible until it attacks, or until its concentration ends (as if concentration on a spell). Any equipment the imp wears or carries is invisible with it.

\end{DndMonster}



\begin{DndMonster}[float=t]{Inquisitor of the Tome}
\DndMonsterType{Medium humanoid, U}
\DndMonsterBasics[
armor-class = {11},
hit-points = {\DndDice{14d8 + 14}},
speed = {30 ft.}
]
\DndMonsterAbilityScores[
 str = 10,
 dex = 12,
 con = 12,
 int = 19,
 wis = 16,
 cha = 15
]
\DndMonsterDetails[
saving-throws = {Int +7, Wis +6, Cha +5},
skills = {Arcana +10, History +7, Nature +7, Religion +10},
condition-immunities = {charmed, frightened},
senses = {truesight 30 ft., passive perception 13},
languages = {any four languages, telepathy 120 ft.},
challenge = 8,
]
\DndMonsterAction{Innate Spellcasting (Psionics)}
The inquisitor casts one of the following spells, requiring no components and using Intelligence as the spellcasting ability (spell save DC 15):
\begin{DndMonsterSpells}
  \DndInnateSpellLevel{detect magic}
  \DndInnateSpellLevel[1]{Otiluke's resilient sphere}
\end{DndMonsterSpells}
\DndMonsterSection{Actions}
\DndMonsterAction{Multiattack}
The inquisitor attacks twice.

\DndMonsterAction{Force Bolt}
\textsl{Ranged Spell Attack:} 7 to hit, range 120 ft., one target. \textsl{Hit:} 22 (4d8 + 4) force damage, and if the target is a Large or smaller creature, the inquisitor can push it up to 10 feet away.

\DndMonsterAction{Silver Longsword}
\textsl{Melee Weapon Attack:} 7 to hit, reach 5 ft., one target. \textsl{Hit:} 8 (1d8 + 4) slashing damage, or 9 (1d10 + 4) if used with two hands, plus 18 (4d8) force damage.

\DndMonsterAction{Implode (Recharge 4\textendash 6)}
Each creature in a 20-foot-radius sphere centered on a point the inquisitor can see within 120 feet of it must succeed on a DC 15 Constitution saving throw or take 31 (6d8 + 4) force damage and be knocked prone and moved to the unoccupied space closest to the sphere's center. Large and smaller objects that aren't being worn or carried in the sphere automatically take the damage and are similarly moved.

\DndMonsterSection{Reactions}
\DndMonsterAction{Telekinetic Deflection}
In response to being hit by an attack roll, the inquisitor increases its AC by 4 against the attack. If this causes the attack to miss, the attacker is hit by the attack instead.

\end{DndMonster}



\begin{DndMonster}[float=t]{Invisible Stalker}
\DndMonsterType{Medium elemental, neutral}
\DndMonsterBasics[
armor-class = {14},
hit-points = {\DndDice{16d8 + 32}},
speed = {50 ft., fly 50 ft. \textit{(hover)}}
]
\DndMonsterAbilityScores[
 str = 16,
 dex = 19,
 con = 14,
 int = 10,
 wis = 15,
 cha = 11
]
\DndMonsterDetails[
skills = {Perception +8, Stealth +10},
damage-resistances = {bludgeoning, piercing, slashing from nonmagical attacks},
damage-immunities = {poison},
condition-immunities = {exhaustion, grappled, paralyzed, petrified, poisoned, prone, restrained, unconscious},
senses = {darkvision 60 ft., passive perception 18},
languages = {Auran, understands Common but doesn't speak it},
challenge = 6,
]
\DndMonsterAction{Invisibility}
The stalker is invisible.

\DndMonsterAction{Faultless Tracker}
The stalker is given a quarry by its summoner. The stalker knows the direction and distance to its quarry as long as the two of them are on the same plane of existence. The stalker also knows the location of its summoner.

\DndMonsterSection{Actions}
\DndMonsterAction{Multiattack}
The stalker makes two slam attacks.

\DndMonsterAction{Slam}
\textsl{Melee Weapon Attack:} 6 to hit, reach 5 ft., one target. \textsl{Hit:} 10 (2d6 + 3) bludgeoning damage.

\end{DndMonster}



\begin{DndMonster}[float=t]{Iron Golem}
\DndMonsterType{Large construct, unaligned}
\DndMonsterBasics[
armor-class = {20 (natural armor)},
hit-points = {\DndDice{20d10 + 100}},
speed = {30 ft.}
]
\DndMonsterAbilityScores[
 str = 24,
 dex = 9,
 con = 20,
 int = 3,
 wis = 11,
 cha = 1
]
\DndMonsterDetails[
damage-immunities = {fire; poison; psychic; bludgeoning, piercing, slashing from nonmagical attacks that aren't adamantine},
condition-immunities = {charmed, exhaustion, frightened, paralyzed, petrified, poisoned},
senses = {darkvision 120 ft., passive perception 10},
languages = {understands the languages of its creator but can't speak},
challenge = 16,
]
\DndMonsterAction{Fire Absorption}
Whenever the golem is subjected to fire damage, it takes no damage and instead regains a number of hit points equal to the fire damage dealt.

\DndMonsterAction{Immutable Form}
The golem is immune to any spell or effect that would alter its form.

\DndMonsterAction{Magic Resistance}
The golem has advantage on saving throws against spells and other magical effects.

\DndMonsterAction{Magic Weapons}
The golem's weapon attacks are magical.

\DndMonsterSection{Actions}
\DndMonsterAction{Multiattack}
The golem makes two melee attacks.

\DndMonsterAction{Slam}
\textsl{Melee Weapon Attack:} 13 to hit, reach 5 ft., one target. \textsl{Hit:} 20 (3d8 + 7) bludgeoning damage.

\DndMonsterAction{Sword}
\textsl{Melee Weapon Attack:} 13 to hit, reach 10 ft., one target. \textsl{Hit:} 23 (3d10 + 7) slashing damage.

\DndMonsterAction{Poison Breath (Recharge 5\textendash 6)}
The golem exhales poisonous gas in a 15-foot cone. Each creature in that area must make a DC 19 Constitution saving throw, taking 45 (10d8) poison damage on a failed save, or half as much damage on a successful one.

\end{DndMonster}



\begin{DndMonster}[float=t]{Kakkuu Spyder-Fiend}
\DndMonsterType{Medium fiend (demon), chaotic evil}
\DndMonsterBasics[
armor-class = {15 (natural armor)},
hit-points = {\DndDice{14d8 + 28}},
speed = {30 ft., climb 30 ft.}
]
\DndMonsterAbilityScores[
 str = 17,
 dex = 15,
 con = 14,
 int = 6,
 wis = 10,
 cha = 10
]
\DndMonsterDetails[
saving-throws = {Dex +5, Con +5, Wis +3},
skills = {Perception +3, Stealth +5},
damage-resistances = {cold, fire, lightning},
damage-immunities = {poison},
condition-immunities = {poisoned},
senses = {darkvision 60 ft., passive perception 13},
languages = {understands Abyssal but can't speak},
challenge = 5,
]
\DndMonsterAction{Magic Resistance}
The kakkuu has advantage on saving throws against spells and other magical effects.

\DndMonsterAction{Spider Climb}
The kakkuu can climb difficult surfaces, including upside down on ceilings, without needing to make an ability check.

\DndMonsterAction{Web Sense}
When in contact with a web, the kakkuu knows the exact location of any other creature in contact with the same web.

\DndMonsterAction{Web Walker}
The kakkuu ignores movement restrictions caused by webbing.

\DndMonsterSection{Actions}
\DndMonsterAction{Multiattack}
The kakkuu makes a Web Snare attack, uses Reel, and makes a Bite attack.

\DndMonsterAction{Bite}
\textsl{Melee Weapon Attack:} 6 to hit, reach 5 ft., one target. \textsl{Hit:} 14 (2d10 + 3) piercing damage plus 10 (3d6) poison damage.

\DndMonsterAction{Reel}
The kakkuu pulls each creature within 60 feet of itself that is grappled by its Web Snare up to 30 feet straight toward itself.

\DndMonsterAction{Web Snare}
\textsl{Ranged Weapon Attack:} 6 to hit, reach 30/60 ft., one Large or smaller creature. \textsl{Hit:} The target has the grappled condition (escape DC 13). While grappled, the target also has the restrained condition. A web snare grappling a creature can be attacked and destroyed (AC 10; 10 hit points; immunity to bludgeoning, poison, and psychic damage).

\end{DndMonster}



\begin{DndMonster}[float*=b,width=\textwidth + 8pt]{Kas the Betrayer}
\begin{multicols}{2}
\DndMonsterType{Medium undead (vampire), neutral evil}
\DndMonsterBasics[
armor-class = {18 (plate armor)},
hit-points = {\DndDice{30d8 + 180}},
speed = {40 ft.}
]
\DndMonsterAbilityScores[
 str = 26,
 dex = 20,
 con = 22,
 int = 24,
 wis = 19,
 cha = 26
]
\DndMonsterDetails[
saving-throws = {Con +13, Wis +11, Cha +15},
skills = {Arcana +14, Deception +22, Perception +11, Stealth +12},
damage-immunities = {necrotic; poison; bludgeoning, piercing, slashing from nonmagical attacks},
condition-immunities = {charmed, exhaustion, frightened, paralyzed, poisoned},
senses = {darkvision 120 ft., passive perception 21},
languages = {Abyssal, Common, Draconic, Infernal},
challenge = 23,
]
\DndMonsterAction{Eager Betrayer}
Kas adds 1d10 to his initiative rolls. He has advantage on attack rolls against any creature that has the frightened condition.

\DndMonsterAction{Legendary Resistance (3/Day)}
If Kas fails a saving throw, he can choose to succeed instead.

\DndMonsterAction{Regeneration}
Kas regains 20 hit points at the start of his turn if he has at least 1 hit point. If he takes radiant damage, this trait doesn't function at the start of his next turn.

\DndMonsterAction{Special Equipment}
Kas wears the Crown of Lies (see the Introduction of Vecna: Eve of Ruin).

\DndMonsterAction{Spider Climb}
Kas can climb difficult surfaces, including upside down on ceilings, without needing to make an ability check.

\DndMonsterAction{Strength of the Night}
Kas doesn't require a coffin, and he drinks blood to sow terror rather than for sustenance. If destroyed, Kas revives in 1d100 nights in an unoccupied space in Tovag, his Domain of Dread. He can be permanently destroyed only by having a stake driven through his heart and then being beheaded. The stake must be cut from a tree growing in soil from Oerth, Kas's home world.

\DndMonsterAction{Sunlight Hypersensitivity}
While in sunlight, Kas takes 20 radiant damage at the start of his turn, has disadvantage on attack rolls and ability checks, and can't use his Change Shape bonus action.

\DndMonsterSection{Actions}
\DndMonsterAction{Multiattack}
Kas makes three Vengeful Sword attacks. He can replace one of these attacks with a Bite attack.

\DndMonsterAction{Vengeful Sword}
\textsl{Melee Weapon Attack:} 15 to hit, reach 5 ft., one target. \textsl{Hit:} 20 (2d8 + 11) slashing damage. The sword scores a critical hit on a roll of 19 or 20.

\DndMonsterAction{Bite}
\textsl{Melee Weapon Attack:} 15 to hit, reach 5 ft., one creature. \textsl{Hit:} 11 (1d6 + 8) piercing damage plus 10 (3d6) necrotic damage. The target's hit point maximum is reduced by an amount equal to the necrotic damage taken, and Kas regains a number of hit points equal to that amount. The reduction lasts until the target finishes a long rest. The target dies if its hit point maximum is reduced to 0. A Humanoid slain in this way and then buried rises the following night as a vampire spawn under Kas's control.

\DndMonsterSection{Bonus Actions}
\DndMonsterAction{Change Shape}
Kas transforms into a Medium cloud of mist. While in this form, Kas has a flying speed of 20 feet, can hover, and can enter a hostile creature's space and stop there. While in mist form, Kas can pass through a space without squeezing as long as air can pass through that space, but he can't pass through water. Kas has advantage on Strength, Dexterity, and Constitution saving throws, and he is immune to all nonmagical damage except the damage he takes as part of his Vampire Weaknesses trait. While in mist form, Kas can't take any actions, speak, or manipulate objects.

\DndMonsterAction{Menacing Glare}
Kas targets one creature he can see within 60 feet of himself. The target must succeed on a DC 23 Wisdom saving throw or have the frightened condition until the start of Kas's next turn.

\DndMonsterSection{Reactions}
\DndMonsterAction{Parrying Riposte}
Kas adds 3 to his AC against one melee attack roll that would hit him. He then makes one Vengeful Sword attack against the attacker if it is within his reach. On a hit, the target takes an additional 9 (2d8) slashing damage.

\DndMonsterSection{Legendary Actions}
Kas the Betrayer can take 3 legendary actions, choosing from the options below. Only one legendary action can be used at a time and only at the end of another creature's turn. Kas the Betrayer regains spent legendary actions at the start of its turn.

\begin{DndMonsterLegendaryActions}
\DndMonsterLegendaryAction{Move}{Kas moves up to his speed without provoking opportunity attacks.
}
\DndMonsterLegendaryAction{Sword (Costs 2 Actions)}{Kas makes one Vengeful Sword attack.
}
\DndMonsterLegendaryAction{Rise, Fallen Soldier (Costs 3 Actions)}{Kas magically summons a specter. The specter appears in an unoccupied space within 30 feet of Kas, whom it obeys. The specter takes its turn immediately after Kas. It lasts for 1 hour, until Kas dies, or until Kas dismisses it as a bonus action. Kas can't have more than two specters summoned at a time.
}
\end{DndMonsterLegendaryActions}
\end{multicols}
\end{DndMonster}



\begin{DndMonster}[float*=b,width=\textwidth + 8pt]{Kraken}
\begin{multicols}{2}
\DndMonsterType{Gargantuan monstrosity (titan), chaotic evil}
\DndMonsterBasics[
armor-class = {18 (natural armor)},
hit-points = {\DndDice{27d20 + 189}},
speed = {20 ft., swim 60 ft.}
]
\DndMonsterAbilityScores[
 str = 30,
 dex = 11,
 con = 25,
 int = 22,
 wis = 18,
 cha = 20
]
\DndMonsterDetails[
saving-throws = {Str +17, Dex +7, Con +14, Int +13, Wis +11},
damage-immunities = {lightning; bludgeoning, piercing, slashing from nonmagical attacks},
condition-immunities = {frightened, paralyzed},
senses = {truesight 120 ft., passive perception 14},
languages = {Abyssal, Celestial, Infernal, Primordial, telepathy 120 ft. but can't speak},
challenge = 23,
]
\DndMonsterAction{Amphibious}
The kraken can breathe air and water.

\DndMonsterAction{Freedom of Movement}
The kraken ignores difficult terrain, and magical effects can't reduce its speed or cause it to be restrained. It can spend 5 feet of movement to escape from nonmagical restraints or being grappled.

\DndMonsterAction{Siege Monster}
The kraken deals double damage to objects and structures.

\DndMonsterSection{Actions}
\DndMonsterAction{Multiattack}
The kraken makes three tentacle attacks, each of which it can replace with one use of Fling.

\DndMonsterAction{Bite}
\textsl{Melee Weapon Attack:} 17 to hit, reach 5 ft., one target. \textsl{Hit:} 23 (3d8 + 10) piercing damage. If the target is a Large or smaller creature grappled by the kraken, that creature is swallowed, and the grapple ends. While swallowed, the creature is blinded and restrained, it has total cover against attacks and other effects outside the kraken, and it takes 42 (12d6) acid damage at the start of each of the kraken's turns. If the kraken takes 50 damage or more on a single turn from a creature inside it, the kraken must succeed on a DC 25 Constitution saving throw at the end of that turn or regurgitate all swallowed creatures, which fall prone in a space within 10 feet of the kraken. If the kraken dies, a swallowed creature is no longer restrained by it and can escape from the corpse using 15 feet of movement, exiting prone.

\DndMonsterAction{Tentacle}
\textsl{Melee Weapon Attack:} 17 to hit, reach 30 ft., one target. \textsl{Hit:} 20 (3d6 + 10) bludgeoning damage, and the target is grappled (escape DC 18). Until this grapple ends, the target is restrained. The kraken has ten tentacles, each of which can grapple one target.

\DndMonsterAction{Fling}
One Large or smaller object held or creature grappled by the kraken is thrown up to 60 feet in a random direction and knocked prone. If a thrown target strikes a solid surface, the target takes 3 (1d6) bludgeoning damage for every 10 feet it was thrown. If the target is thrown at another creature, that creature must succeed on a DC 18 Dexterity saving throw or take the same damage and be knocked prone.

\DndMonsterAction{Lightning Storm}
The kraken magically creates three bolts of lightning, each of which can strike a target the kraken can see within 120 feet of it. A target must make a DC 23 Dexterity saving throw, taking 22 (4d10) lightning damage on a failed save, or half as much damage on a successful one.

\DndMonsterSection{Legendary Actions}
Kraken can take 3 legendary actions, choosing from the options below. Only one legendary action can be used at a time and only at the end of another creature's turn. Kraken regains spent legendary actions at the start of its turn.

\begin{DndMonsterLegendaryActions}
\DndMonsterLegendaryAction{Tentacle Attack or Fling}{The kraken makes one tentacle attack or uses its Fling.
}
\DndMonsterLegendaryAction{Lightning Storm (Costs 2 Actions)}{The kraken uses Lightning Storm.
}
\DndMonsterLegendaryAction{Ink Cloud (Costs 3 Actions)}{While underwater, the kraken expels an ink cloud in a 60-foot radius. The cloud spreads around corners, and that area is heavily obscured to creatures other than the kraken. Each creature other than the kraken that ends its turn there must succeed on a DC 23 Constitution saving throw, taking 16 (3d10) poison damage on a failed save, or half as much damage on a successful one. A strong current disperses the cloud, which otherwise disappears at the end of the kraken's next turn.
}
\end{DndMonsterLegendaryActions}
\DndMonsterSection{Lair Actions}
On initiative count 20 (losing initiative ties), the kraken takes a lair action to cause one of the following magical effects:


\begin{itemize}
\item A strong current moves through the kraken's lair. Each creature within 60 feet of the kraken must succeed on a DC 23 Strength saving throw or be pushed up to 60 feet away from the kraken. On a success, the creature is pushed 10 feet away from the kraken.
\item Creatures in the water within 60 feet of the kraken have vulnerability to lightning damage until initiative count 20 on the next round.
\item The water in the kraken's lair becomes electrically charged. All creatures within 120 feet of the kraken must succeed on a DC 23 Constitution saving throw, taking 10 (3d6) lightning damage on a failed save, or half as much damage on a successful one.
\end{itemize}

\DndMonsterSection{Regional Effects}
The region containing a kraken's lair is warped by the creature's blasphemous presence, creating the following magical effects:


\begin{itemize}
\item The kraken can alter the weather at will in a 6-mile radius centered on its lair. The effect is identical to the \textit{control weather} spell.
\item Water elementals coalesce within 6 miles of the lair. These elementals can't leave the water and have Intelligence and Charisma scores of 1 (-5).
\item Aquatic creatures within 6 miles of the lair that have an Intelligence score of 2 or lower are charmed by the kraken and aggressive toward intruders in the area.
\end{itemize}

When the kraken dies, all of these regional effects fade immediately.

\end{multicols}
\end{DndMonster}



\begin{DndMonster}[float=t]{Lemure}
\DndMonsterType{Medium fiend (devil), lawful evil}
\DndMonsterBasics[
armor-class = {7},
hit-points = {\DndDice{3d8}},
speed = {15 ft.}
]
\DndMonsterAbilityScores[
 str = 10,
 dex = 5,
 con = 11,
 int = 1,
 wis = 11,
 cha = 3
]
\DndMonsterDetails[
damage-resistances = {cold},
damage-immunities = {fire, poison},
condition-immunities = {charmed, frightened, poisoned},
senses = {darkvision 120 ft., passive perception 10},
languages = {understands Infernal but can't speak},
challenge = 0,
]
\DndMonsterAction{Devil's Sight}
Magical darkness doesn't impede the lemure's darkvision.

\DndMonsterAction{Hellish Rejuvenation}
A lemure that dies in the Nine Hells comes back to life with all its hit points in 1d10 days unless it is killed by a good-aligned creature with a \textit{bless} spell cast on that creature or its remains are sprinkled with holy water.

\DndMonsterSection{Actions}
\DndMonsterAction{Fist}
\textsl{Melee Weapon Attack:} 3 to hit, reach 5 ft., one target. \textsl{Hit:} 2 (1d4) bludgeoning damage

\end{DndMonster}



\begin{DndMonster}[float=t]{Lonely Sorrowsworn}
\DndMonsterType{Medium monstrosity, neutral evil}
\DndMonsterBasics[
armor-class = {16 (natural armor)},
hit-points = {\DndDice{15d8 + 45}},
speed = {30 ft.}
]
\DndMonsterAbilityScores[
 str = 16,
 dex = 12,
 con = 17,
 int = 6,
 wis = 11,
 cha = 6
]
\DndMonsterDetails[
damage-resistances = {bludgeoning, piercing, slashing while in dim light or darkness},
senses = {darkvision 60 ft., passive perception 10},
languages = {Common},
challenge = 9,
]
\DndMonsterAction{Psychic Leech}
At the start of each of the sorrowsworn's turns, each creature within 5 feet of it must succeed on a DC 15 Wisdom saving throw or take 10 (3d6) psychic damage.

\DndMonsterAction{Thrives on Company}
The sorrowsworn has advantage on attack rolls while it is within 30 feet of at least two other creatures. It otherwise has disadvantage on attack rolls.

\DndMonsterSection{Actions}
\DndMonsterAction{Multiattack}
The sorrowsworn makes one Harpoon Arm attack, and it uses Sorrowful Embrace.

\DndMonsterAction{Harpoon Arm}
\textsl{Melee Weapon Attack:} 7 to hit, reach 60 ft., one target. \textsl{Hit:} 21 (4d8 + 3) piercing damage, and the target is grappled (escape DC 15) if it is a Large or smaller creature. The sorrowsworn has two harpoon arms and can grapple up to two creatures at once.

\DndMonsterAction{Sorrowful Embrace}
Each creature grappled by the sorrowsworn must make a DC 15 Wisdom saving throw, taking 18 (4d8) psychic damage on a failed save, or half as much damage on a successful one. In either case, the sorrowsworn pulls each of those creatures up to 30 feet straight toward it.

\end{DndMonster}



\begin{DndMonster}[float*=b,width=\textwidth + 8pt]{Lord Soth}
\begin{multicols}{2}
\DndMonsterType{Medium undead (paladin), lawful evil}
\DndMonsterBasics[
armor-class = {18 (plate armor)},
hit-points = {\DndDice{24d8 + 120}},
speed = {30 ft.}
]
\DndMonsterAbilityScores[
 str = 22,
 dex = 11,
 con = 20,
 int = 12,
 wis = 16,
 cha = 20
]
\DndMonsterDetails[
saving-throws = {Dex +6, Wis +9, Cha +11},
damage-immunities = {necrotic, poison},
condition-immunities = {exhaustion, frightened, poisoned},
senses = {darkvision 120 ft., passive perception 13},
languages = {Common, Infernal, Solamnic},
challenge = 19,
]
\DndMonsterAction{Legendary Resistance (3/Day)}
If Soth fails a saving throw, he can choose to succeed instead.

\DndMonsterAction{Magic Resistance}
Soth has advantage on saving throws against spells and other magical effects.

\DndMonsterAction{Marshal Undead}
Unless Soth is incapacitated, he and Undead creatures of his choice within 60 feet of him are immune to features that turn Undead.

\DndMonsterAction{Unusual Nature}
Soth doesn't require air, food, drink, or sleep.

\DndMonsterAction{Spellcasting}
Soth casts one of the following spells, requiring no material components and using Charisma as the spellcasting ability (spell save DC 19):
\begin{DndMonsterSpells}
  \DndInnateSpellLevel{command}
  \DndInnateSpellLevel[]{banishment}
  \DndInnateSpellLevel[2]{dispel magic}
\end{DndMonsterSpells}
\DndMonsterSection{Actions}
\DndMonsterAction{Multiattack}
Soth makes three Forsaken Brand attacks.

\DndMonsterAction{Forsaken Brand}
\textsl{Melee Weapon Attack:} 12 to hit, reach 5 ft., one target. \textsl{Hit:} 10 (1d8 + 6) bludgeoning damage plus 18 (4d8) necrotic damage, and if the target is a creature, it can't regain hit points until the start of Soth's next turn.

\DndMonsterAction{Cataclysmic Fire (1/Day)}
Soth hurls a magical ball of fire that explodes at a point he can see within 120 feet of himself. Each creature in a 20-foot-radius sphere centered on that point must make a DC 19 Dexterity saving throw. A creature takes 35 (10d6) fire damage and 35 (10d6) necrotic damage on a failed save, or half as much damage on a successful one.

Additionally, any Medium or smaller Humanoid killed by this damage, as well as every corpse of such a creature within the sphere, becomes a skeleton (see the Monster Manual) under Soth's control. The skeleton acts on Soth's initiative but immediately after his turn. Absent any other command, the skeleton tries to kill any non-Undead creature it encounters.

\DndMonsterAction{Word of Death (1/Day)}
Soth points at a creature he can see within 60 feet of himself and magically commands it to die. The target must make a DC 19 Constitution saving throw, taking 100 necrotic damage on a failed save, or half as much damage on a successful one. If this damage reduces the target to 0 hit points, the target dies.

\DndMonsterSection{Legendary Actions}
Lord Soth can take 3 legendary actions, choosing from the options below. Only one legendary action can be used at a time and only at the end of another creature's turn. Lord Soth regains spent legendary actions at the start of its turn.

\begin{DndMonsterLegendaryActions}
\DndMonsterLegendaryAction{Implacable Maneuver}{Soth moves up to his speed or commands a mount he is riding to move up to its speed. The movement from this action doesn't provoke opportunity attacks. If he or his mount moves within 5 feet of a creature during this movement, he can force the creature to make a DC 20 Strength saving throw. The creature is knocked prone unless it succeeds on the saving throw.
}
\DndMonsterLegendaryAction{Strike (Costs 2 Actions)}{Soth makes one Forsaken Brand attack.
}
\DndMonsterLegendaryAction{Cast a Spell (Costs 3 Actions)}{Soth uses Spellcasting.
}
\end{DndMonsterLegendaryActions}
\end{multicols}
\end{DndMonster}



\begin{DndMonster}[float=t]{Lost Sorrowsworn}
\DndMonsterType{Medium monstrosity, neutral evil}
\DndMonsterBasics[
armor-class = {15 (natural armor)},
hit-points = {\DndDice{12d8 + 24}},
speed = {30 ft.}
]
\DndMonsterAbilityScores[
 str = 17,
 dex = 12,
 con = 15,
 int = 6,
 wis = 7,
 cha = 5
]
\DndMonsterDetails[
skills = {Athletics +6},
damage-resistances = {bludgeoning, piercing, slashing while in dim light or darkness},
senses = {darkvision 60 ft., passive perception 8},
languages = {Common},
challenge = 7,
]
\DndMonsterSection{Actions}
\DndMonsterAction{Multiattack}
The sorrowsworn makes two Arm Spike attacks.

\DndMonsterAction{Arm Spike}
\textsl{Melee Weapon Attack:} 6 to hit, reach 10 ft., one target. \textsl{Hit:} 14 (2d10 + 3) piercing damage.

\DndMonsterAction{Embrace (Recharge 4\textendash 6)}
\textsl{Melee Weapon Attack:} 6 to hit, reach 5 ft., one target. \textsl{Hit:} 25 (4d10 + 3) piercing damage, and the target is grappled (escape DC 14) if it is a Medium or smaller creature. Until the grapple ends, the target is frightened, and it takes 27 (6d8) psychic damage at the end of each of its turns. The sorrowsworn can grapple only one creature at a time.

\DndMonsterSection{Reactions}
\DndMonsterAction{Tightening Embrace}
If the sorrowsworn takes damage, the creature grappled by Embrace takes 18 (4d8) psychic damage.

\end{DndMonster}


\clearpage

\begin{DndMonster}[float=t]{Mage}
\DndMonsterType{Medium humanoid (any race), any alignment}
\DndMonsterBasics[
armor-class = {12 (15 with \textit{mage armor})},
hit-points = {\DndDice{9d8}},
speed = {30 ft.}
]
\DndMonsterAbilityScores[
 str = 9,
 dex = 14,
 con = 11,
 int = 17,
 wis = 12,
 cha = 11
]
\DndMonsterDetails[
saving-throws = {Int +6, Wis +4},
skills = {Arcana +6, History +6},
languages = {any four languages},
challenge = 6,
]
\DndMonsterAction{Spellcasting}
The mage is a 9th-level spellcaster. Its spellcasting ability is Intelligence (spell save DC 14, 6 to hit with spell attacks). The mage has the following wizard spells prepared:
\begin{DndMonsterSpells}
  \DndMonsterSpellLevel{fire bolt}
   \DndMonsterSpellLevel[1][4]{detect magic}
   \DndMonsterSpellLevel[2][3]{misty step}
   \DndMonsterSpellLevel[3][3]{counterspell}
   \DndMonsterSpellLevel[4][3]{greater invisibility}
   \DndMonsterSpellLevel[5][1]{cone of cold}
\end{DndMonsterSpells}
\DndMonsterSection{Actions}
\DndMonsterAction{Dagger}
\textsl{Melee or Ranged Weapon Attack:} 5 to hit, reach 5 ft. or range 20/60 ft., one target. \textsl{Hit:} 4 (1d4 + 2) piercing damage.

\end{DndMonster}



\begin{DndMonster}[float=t]{Marid}
\DndMonsterType{Large elemental, chaotic neutral}
\DndMonsterBasics[
armor-class = {17 (natural armor)},
hit-points = {\DndDice{17d10 + 136}},
speed = {30 ft., fly 60 ft., swim 90 ft.}
]
\DndMonsterAbilityScores[
 str = 22,
 dex = 12,
 con = 26,
 int = 18,
 wis = 17,
 cha = 18
]
\DndMonsterDetails[
saving-throws = {Dex +5, Wis +7, Cha +8},
damage-resistances = {acid, cold, lightning},
senses = {blindsight 30 ft., darkvision 120 ft., passive perception 13},
languages = {Aquan},
challenge = 11,
]
\DndMonsterAction{Amphibious}
The marid can breathe air and water.

\DndMonsterAction{Elemental Demise}
If the marid dies, its body disintegrates into a burst of water and foam, leaving behind only equipment the marid was wearing or carrying.

\DndMonsterAction{Innate Spellcasting}
The marid's innate spellcasting ability is Charisma (spell save DC 16, 8 to hit with spell attacks). It can innately cast the following spells, requiring no material components:
\begin{DndMonsterSpells}
  \DndInnateSpellLevel{create or destroy water}
  \DndInnateSpellLevel[3]{tongues}
  \DndInnateSpellLevel[1]{conjure elemental}
\end{DndMonsterSpells}
\DndMonsterSection{Actions}
\DndMonsterAction{Multiattack}
The marid makes two trident attacks.

\DndMonsterAction{Trident}
\textsl{Melee or Ranged Weapon Attack:} 10 to hit, reach 5 ft. or range 20/60 ft., one target. \textsl{Hit:} 13 (2d6 + 6) piercing damage, or 15 (2d8 + 6) piercing damage if used with two hands to make a melee attack.

\DndMonsterAction{Water Jet}
The marid magically shoots water in a 60-foot line that is 5 feet wide. Each creature in that line must make a DC 16 Dexterity saving throw. On a failure, a target takes 21 (6d6) bludgeoning damage and, if it is Huge or smaller, is pushed up to 20 feet away from the marid and knocked prone. On a success, a target takes half the bludgeoning damage, but is neither pushed nor knocked prone.

\end{DndMonster}



\begin{DndMonster}[float=t]{Marilith}
\DndMonsterType{Large fiend (demon), chaotic evil}
\DndMonsterBasics[
armor-class = {18 (natural armor)},
hit-points = {\DndDice{18d10 + 90}},
speed = {40 ft.}
]
\DndMonsterAbilityScores[
 str = 18,
 dex = 20,
 con = 20,
 int = 18,
 wis = 16,
 cha = 20
]
\DndMonsterDetails[
saving-throws = {Str +9, Con +10, Wis +8, Cha +10},
damage-resistances = {cold; fire; lightning; bludgeoning, piercing, slashing from nonmagical attacks},
damage-immunities = {poison},
condition-immunities = {poisoned},
senses = {truesight 120 ft., passive perception 13},
languages = {Abyssal, telepathy 120 ft.},
challenge = 16,
]
\DndMonsterAction{Magic Resistance}
The marilith has advantage on saving throws against spells and other magical effects.

\DndMonsterAction{Magic Weapons}
The marilith's weapon attacks are magical.

\DndMonsterAction{Reactive}
The marilith can take one reaction on every turn in combat.

\DndMonsterSection{Actions}
\DndMonsterAction{Multiattack}
The marilith can make seven attacks: six with its longswords and one with its tail.

\DndMonsterAction{Longsword}
\textsl{Melee Weapon Attack:} 9 to hit, reach 5 ft., one target. \textsl{Hit:} 13 (2d8 + 4) slashing damage.

\DndMonsterAction{Tail}
\textsl{Melee Weapon Attack:} 9 to hit, reach 10 ft., one creature. \textsl{Hit:} 15 (2d10 + 4) bludgeoning damage. If the target is Medium or smaller, it is grappled (escape DC 19). Until this grapple ends, the target is restrained, the marilith can automatically hit the target with its tail, and the marilith can't make tail attacks against other targets.

\DndMonsterAction{Teleport}
The marilith magically teleports, along with any equipment it is wearing or carrying, up to 120 feet to an unoccupied space it can see.

\DndMonsterSection{Reactions}
\DndMonsterAction{Parry}
The marilith adds 5 to its AC against one melee attack that would hit it. To do so, the marilith must see the attacker and be wielding a melee weapon.

\end{DndMonster}



\begin{DndMonster}[float=t]{Mimic}
\DndMonsterType{Medium monstrosity (shapechanger), neutral}
\DndMonsterBasics[
armor-class = {12 (natural armor)},
hit-points = {\DndDice{9d8 + 18}},
speed = {15 ft.}
]
\DndMonsterAbilityScores[
 str = 17,
 dex = 12,
 con = 15,
 int = 5,
 wis = 13,
 cha = 8
]
\DndMonsterDetails[
skills = {Stealth +5},
damage-immunities = {acid},
condition-immunities = {prone},
senses = {darkvision 60 ft., passive perception 11},
challenge = 2,
]
\DndMonsterAction{Shapechanger}
The mimic can use its action to polymorph into an object or back into its true, amorphous form. Its statistics are the same in each form. Any equipment it is wearing or carrying isn't transformed. It reverts to its true form if it dies.

\DndMonsterAction{Adhesive (Object Form Only)}
The mimic adheres to anything that touches it. A Huge or smaller creature adhered to the mimic is also grappled by it (escape DC 13). Ability checks made to escape this grapple have disadvantage.

\DndMonsterAction{False Appearance (Object Form Only)}
While the mimic remains motionless, it is indistinguishable from an ordinary object.

\DndMonsterAction{Grappler}
The mimic has advantage on attack rolls against any creature grappled by it.

\DndMonsterSection{Actions}
\DndMonsterAction{Pseudopod}
\textsl{Melee Weapon Attack:} 5 to hit, reach 5 ft., one target. \textsl{Hit:} 7 (1d8 + 3) bludgeoning damage. If the mimic is in object form, the target is subjected to its Adhesive trait.

\DndMonsterAction{Bite}
\textsl{Melee Weapon Attack:} 5 to hit, reach 5 ft., one target. \textsl{Hit:} 7 (1d8 + 3) piercing damage plus 4 (1d8) acid damage.

\end{DndMonster}



\begin{DndMonster}[float=t]{Minotaur Skeleton}
\DndMonsterType{Large undead, lawful evil}
\DndMonsterBasics[
armor-class = {12 (natural armor)},
hit-points = {\DndDice{9d10 + 18}},
speed = {40 ft.}
]
\DndMonsterAbilityScores[
 str = 18,
 dex = 11,
 con = 15,
 int = 6,
 wis = 8,
 cha = 5
]
\DndMonsterDetails[
damage-vulnerabilities = {bludgeoning},
damage-immunities = {poison},
condition-immunities = {exhaustion, poisoned},
senses = {darkvision 60 ft., passive perception 9},
languages = {understands Abyssal but can't speak},
challenge = 2,
]
\DndMonsterAction{Charge}
If the skeleton moves at least 10 feet straight toward a target and then hits it with a gore attack on the same turn, the target takes an extra 9 (2d8) piercing damage. If the target is a creature, it must succeed on a DC 14 Strength saving throw or be pushed up to 10 feet away and knocked prone.

\DndMonsterSection{Actions}
\DndMonsterAction{Greataxe}
\textsl{Melee Weapon Attack:} 6 to hit, reach 5 ft., one target. \textsl{Hit:} 17 (2d12 + 4) slashing damage.

\DndMonsterAction{Gore}
\textsl{Melee Weapon Attack:} 6 to hit, reach 5 ft., one target. \textsl{Hit:} 13 (2d8 + 4) piercing damage.

\end{DndMonster}



\begin{DndMonster}[float=t]{Mirror Shade}
\DndMonsterType{Medium undead, chaotic evil}
\DndMonsterBasics[
armor-class = {13},
hit-points = {\DndDice{14d8 + 28}},
speed = {40 ft.}
]
\DndMonsterAbilityScores[
 str = 8,
 dex = 17,
 con = 14,
 int = 10,
 wis = 13,
 cha = 18
]
\DndMonsterDetails[
saving-throws = {Dex +7, Wis +5},
skills = {Deception +8, Stealth +7},
damage-resistances = {acid; fire; lightning; bludgeoning, piercing, slashing from nonmagical attacks},
damage-immunities = {poison, psychic, radiant},
condition-immunities = {blinded, exhaustion, frightened, grappled, paralyzed, petrified, poisoned, prone, restrained},
senses = {darkvision 60 ft., passive perception 11},
challenge = 10,
]
\DndMonsterAction{False Appearance}
If the mirror shade is within 5 feet of a reflective surface—such as a mirror, glass pane, or still water—it has advantage on its initiative roll. If a creature hasn't observed the mirror shade move or act, that creature must succeed on a DC 18 Intelligence (Investigation) check to discern that the mirror shade isn't the creature's own reflection.

\DndMonsterAction{Mirror Movement}
The mirror shade can move along the surface of reflective or translucent objects, such as mirrors, without provoking opportunity attacks. It can move through translucent objects as if they were difficult terrain.

\DndMonsterSection{Actions}
\DndMonsterAction{Multiattack}
The mirror shade makes two Phantasmal Strike attacks and uses Reflect Fear.

\DndMonsterAction{Phantasmal Strike}
\textsl{Melee Weapon Attack:} 7 to hit, reach 5 ft., one target. \textsl{Hit:} 7 (1d8 + 3) radiant damage plus 7 (2d6) psychic damage.

\DndMonsterAction{Reflect Fear}
The mirror shade targets one creature it can see within 60 feet of itself and projects an illusion of that creature's greatest fear. The target must make a DC 16 Wisdom saving throw. On a failed save, the target takes 28 (8d6) psychic damage and has the frightened condition until the start of the mirror shade's next turn. On a successful save, the target takes half as much damage only.

\DndMonsterSection{Bonus Actions}
\DndMonsterAction{Mirror Stealth}
While within 5 feet of a reflective surface, such as a mirror, the mirror shade takes the Hide action.

\end{DndMonster}



\begin{DndMonster}[float*=b,width=\textwidth + 8pt]{Miska the Wolf-Spider}
\begin{multicols}{2}
\DndMonsterType{Huge fiend (demon), chaotic evil}
\DndMonsterBasics[
armor-class = {21 (natural armor)},
hit-points = {\DndDice{38d12 + 152}},
speed = {40 ft., climb 40 ft.}
]
\DndMonsterAbilityScores[
 str = 23,
 dex = 18,
 con = 19,
 int = 18,
 wis = 21,
 cha = 22
]
\DndMonsterDetails[
saving-throws = {Dex +11, Con +11, Wis +12},
skills = {Insight +12, Perception +12, Stealth +11},
damage-resistances = {cold, fire, lightning},
damage-immunities = {poison; bludgeoning, piercing, slashing from nonmagical attacks},
condition-immunities = {poisoned},
senses = {truesight 120 ft., passive perception 21},
languages = {Abyssal, Common, telepathy 120 ft.},
challenge = 24,
]
\DndMonsterAction{Foul Ichor}
A creature that hits Miska with a melee weapon attack takes 7 (2d6) poison damage.

\DndMonsterAction{Legendary Resistance (3/Day)}
If Miska fails a saving throw, he can choose to succeed instead.

\DndMonsterAction{Magic Resistance}
Miska has advantage on saving throws against spells and other magical effects.

\DndMonsterAction{Spider Climb}
Miska can climb difficult surfaces, including upside down on ceilings, without needing to make an ability check.

\DndMonsterAction{Web Sense}
When in contact with a web, Miska knows the exact location of any other creature in contact with the same web.

\DndMonsterAction{Web Walker}
Miska ignores movement restrictions caused by webbing.

\DndMonsterAction{Spellcasting}
Miska casts one of the following spells, requiring no material components and using Charisma as the spellcasting ability (spell save DC 21):
\begin{DndMonsterSpells}
  \DndInnateSpellLevel{Disguise Self}
  \DndInnateSpellLevel[2]{Dominate Monster}
\end{DndMonsterSpells}
\DndMonsterSection{Actions}
\DndMonsterAction{Multiattack}
Miska makes one Lupine Bite attack and two Trident of Chaos attacks.

\DndMonsterAction{Lupine Bite}
\textsl{Melee Weapon Attack:} 13 to hit, reach 10 ft., one target. \textsl{Hit:} 17 (2d10 + 6) piercing damage plus 27 (6d8) poison damage. If the target is a creature, it must succeed on a DC 21 Constitution saving throw or have the poisoned condition for 1 minute. While poisoned in this way, a creature has the incapacitated condition and can't regain hit points. A poisoned creature can repeat the saving throw at the end of each of its turns, ending the effect on itself on a success.

\DndMonsterAction{Trident of Chaos}
\textsl{Melee Weapon Attack:} 13 to hit, reach 15 ft., one target. \textsl{Hit:} 13 (2d6 + 6) piercing damage plus 9 (2d8) force damage.

\DndMonsterSection{Bonus Actions}
\DndMonsterAction{Demand Loyalty}
Miska magically ends the charmed and frightened conditions on himself and on any of his allies within 120 feet of himself.

\DndMonsterSection{Legendary Actions}
Miska the Wolf-Spider can take 3 legendary actions, choosing from the options below. Only one legendary action can be used at a time and only at the end of another creature's turn. Miska the Wolf-Spider regains spent legendary actions at the start of its turn.

\begin{DndMonsterLegendaryActions}
\DndMonsterLegendaryAction{Howl}{Miska utters a bloodthirsty howl at one creature within 120 feet of himself that isn't a Fiend. The target must succeed on a DC 20 Wisdom saving throw or take 13 (2d12) psychic damage.
}
\DndMonsterLegendaryAction{Skitter}{Miska moves up to his speed without provoking opportunity attacks.
}
\DndMonsterLegendaryAction{Cast a Spell (Costs 2 Actions)}{Miska uses Spellcasting.
}
\end{DndMonsterLegendaryActions}
\end{multicols}
\end{DndMonster}



\begin{DndMonster}[float=t]{Moonlight Guardian}
\DndMonsterType{Medium construct, unaligned}
\DndMonsterBasics[
armor-class = {16 (natural armor)},
hit-points = {\DndDice{14d8 + 42}},
speed = {30 ft.}
]
\DndMonsterAbilityScores[
 str = 19,
 dex = 9,
 con = 16,
 int = 6,
 wis = 12,
 cha = 6
]
\DndMonsterDetails[
damage-resistances = {bludgeoning, piercing, slashing from nonmagical attacks},
damage-immunities = {poison, radiant},
condition-immunities = {charmed, exhaustion, frightened, paralyzed, petrified, poisoned},
senses = {truesight 120 ft., passive perception 11},
languages = {understands the languages of its creator but can't speak},
challenge = 6,
]
\DndMonsterAction{Immutable Form}
The guardian is immune to any spell or other effect that would alter its form.

\DndMonsterAction{Magic Resistance}
The guardian has advantage on saving throws against spells and other magical effects.

\DndMonsterAction{Radiant Absorption}
Whenever the guardian is subjected to radiant damage, it takes no damage and instead regains a number of hit points equal to the radiant damage.

\DndMonsterSection{Actions}
\DndMonsterAction{Multiattack}
The guardian makes two Moonlight Slam attacks.

\DndMonsterAction{Moonlight Slam}
\textsl{Melee Weapon Attack:} 7 to hit, reach 5 ft., one target. \textsl{Hit:} 8 (1d8 + 4) bludgeoning damage plus 4 (1d8) radiant damage.

\DndMonsterAction{Moonlight Blast (Recharge 5\textendash 6)}
The guardian unleashes a magical blast of moonlight in a line 60 feet long and 5 feet wide. Each creature in that area must make a DC 14 Dexterity saving throw. Creatures that aren't in their true form have disadvantage on the save. On a failed save, a creature takes 22 (5d8) radiant damage, and if it isn't in its true form, it is forced into its true form and can't change forms until the end of the guardian's next turn. On a successful save, a creature takes half as much damage only.

\end{DndMonster}



\begin{DndMonster}[float=t]{Nalfeshnee}
\DndMonsterType{Large fiend (demon), chaotic evil}
\DndMonsterBasics[
armor-class = {18 (natural armor)},
hit-points = {\DndDice{16d10 + 96}},
speed = {20 ft., fly 30 ft.}
]
\DndMonsterAbilityScores[
 str = 21,
 dex = 10,
 con = 22,
 int = 19,
 wis = 12,
 cha = 15
]
\DndMonsterDetails[
saving-throws = {Con +11, Int +9, Wis +6, Cha +7},
damage-resistances = {cold; fire; lightning; bludgeoning, piercing, slashing from nonmagical attacks},
damage-immunities = {poison},
condition-immunities = {poisoned},
senses = {truesight 120 ft., passive perception 11},
languages = {Abyssal, telepathy 120 ft.},
challenge = 13,
]
\DndMonsterAction{Magic Resistance}
The nalfeshnee has advantage on saving throws against spells and other magical effects.

\DndMonsterSection{Actions}
\DndMonsterAction{Multiattack}
The nalfeshnee uses Horror Nimbus if it can. It then makes three attacks: one with its bite and two with its claws.

\DndMonsterAction{Bite}
\textsl{Melee Weapon Attack:} 10 to hit, reach 5 ft., one target. \textsl{Hit:} 32 (5d10 + 5) piercing damage.

\DndMonsterAction{Claw}
\textsl{Melee Weapon Attack:} 10 to hit, reach 10 ft., one target. \textsl{Hit:} 15 (3d6 + 5) slashing damage.

\DndMonsterAction{Horror Nimbus (Recharge 5\textendash 6)}
The nalfeshnee magically emits scintillating, multicolored light. Each creature within 15 feet of the nalfeshnee that can see the light must succeed on a DC 15 Wisdom saving throw or be frightened for 1 minute. A creature can repeat the saving throw at the end of each of its turns, ending the effect on itself on a success. If a creature's saving throw is successful or the effect ends for it, the creature is immune to the nalfeshnee's Horror Nimbus for the next 24 hours.

\DndMonsterAction{Teleport}
The nalfeshnee magically teleports, along with any equipment it is wearing or carrying, up to 120 feet to an unoccupied space it can see.

\end{DndMonster}



\begin{DndMonster}[float=t]{Necromancer Wizard}
\DndMonsterType{Medium humanoid, any alignment}
\DndMonsterBasics[
armor-class = {12 (15 with \textit{mage armor})},
hit-points = {\DndDice{20d8 + 20}},
speed = {30 ft.}
]
\DndMonsterAbilityScores[
 str = 9,
 dex = 14,
 con = 12,
 int = 17,
 wis = 12,
 cha = 11
]
\DndMonsterDetails[
saving-throws = {Int +7, Wis +5},
skills = {Arcana +7, History +7},
damage-resistances = {necrotic},
languages = {any four languages},
challenge = 9,
]
\DndMonsterAction{Spellcasting}
The necromancer casts one of the following spells, using Intelligence as the spellcasting ability (spell save DC 15):
\begin{DndMonsterSpells}
  \DndInnateSpellLevel{dancing lights}
  \DndInnateSpellLevel[2]{bestow curse}
  \DndInnateSpellLevel[1]{circle of death}
\end{DndMonsterSpells}
\DndMonsterSection{Actions}
\DndMonsterAction{Multiattack}
The necromancer makes three Arcane Burst attacks.

\DndMonsterAction{Arcane Burst}
\textsl{Melee or Ranged Spell Attack:} 7 to hit, reach 5 ft. or range 120 ft., one target. \textsl{Hit:} 25 (4d10 + 3) necrotic damage.

\DndMonsterSection{Bonus Actions}
\DndMonsterAction{Summon Undead (1/Day)}
The necromancer magically summons five skeletons or zombies. The summoned creatures appear in unoccupied spaces within 60 feet of the necromancer, whom they obey. They take their turns immediately after the necromancer. Each lasts for 1 hour, until it or the necromancer dies, or until the necromancer dismisses it as a bonus action.

\DndMonsterSection{Reactions}
\DndMonsterAction{Grim Harvest (1/Turn)}
When the necromancer kills a creature with necrotic damage, the necromancer regains 9 (2d8) hit points. 

\end{DndMonster}



\begin{DndMonster}[float=t]{Night Scavver}
\DndMonsterType{Huge monstrosity, unaligned}
\DndMonsterBasics[
armor-class = {14 (natural armor)},
hit-points = {\DndDice{12d12 + 36}},
speed = {0 ft., fly 40 ft.}
]
\DndMonsterAbilityScores[
 str = 20,
 dex = 15,
 con = 17,
 int = 1,
 wis = 10,
 cha = 1
]
\DndMonsterDetails[
skills = {Perception +6, Stealth +8},
senses = {darkvision 120 ft., passive perception 16},
challenge = 5,
]
\DndMonsterAction{Unusual Nature}
The scavver doesn't require air.

\DndMonsterSection{Actions}
\DndMonsterAction{Bite}
\textsl{Melee Weapon Attack:} 8 to hit (with advantage if the target is a creature that is missing any hit points), reach 10 ft., one target. \textsl{Hit:} 27 (4d10 + 5) piercing damage.

\end{DndMonster}



\begin{DndMonster}[float=t]{Nightmare}
\DndMonsterType{Large fiend, neutral evil}
\DndMonsterBasics[
armor-class = {13 (natural armor)},
hit-points = {\DndDice{8d10 + 24}},
speed = {60 ft., fly 90 ft.}
]
\DndMonsterAbilityScores[
 str = 18,
 dex = 15,
 con = 16,
 int = 10,
 wis = 13,
 cha = 15
]
\DndMonsterDetails[
damage-immunities = {fire},
languages = {understands Abyssal, Common, and Infernal but can't speak },
challenge = 3,
]
\DndMonsterAction{Confer Fire Resistance}
The nightmare can grant resistance to fire damage to anyone riding it.

\DndMonsterAction{Illumination}
The nightmare sheds bright light in a 10-foot radius and dim light for an additional 10 feet.

\DndMonsterSection{Actions}
\DndMonsterAction{Hooves}
\textsl{Melee Weapon Attack:} 6 to hit, reach 5 ft., one target. \textsl{Hit:} 13 (2d8 + 4) bludgeoning damage plus 7 (2d6) fire damage.

\DndMonsterAction{Ethereal Stride}
The nightmare and up to three willing creatures within 5 feet of it magically enter the Ethereal Plane from the Material Plane, or vice versa.

\end{DndMonster}



\begin{DndMonster}[float=t]{Noble}
\DndMonsterType{Medium humanoid (any race), any alignment}
\DndMonsterBasics[
armor-class = {15 (\textit{breastplate})},
hit-points = {\DndDice{2d8}},
speed = {30 ft.}
]
\DndMonsterAbilityScores[
 str = 11,
 dex = 12,
 con = 11,
 int = 12,
 wis = 14,
 cha = 16
]
\DndMonsterDetails[
skills = {Deception +5, Insight +4, Persuasion +5},
languages = {any two languages},
challenge = 1/8,
]
\DndMonsterSection{Actions}
\DndMonsterAction{Rapier}
\textsl{Melee Weapon Attack:} 3 to hit, reach 5 ft., one target. \textsl{Hit:} 5 (1d8 + 1) piercing damage.

\DndMonsterSection{Reactions}
\DndMonsterAction{Parry}
The noble adds 2 to its AC against one melee attack that would hit it. To do so, the noble must see the attacker and be wielding a melee weapon.

\end{DndMonster}



\begin{DndMonster}[float=t]{Nothic}
\DndMonsterType{Medium aberration, neutral evil}
\DndMonsterBasics[
armor-class = {15 (natural armor)},
hit-points = {\DndDice{6d8 + 18}},
speed = {30 ft.}
]
\DndMonsterAbilityScores[
 str = 14,
 dex = 16,
 con = 16,
 int = 13,
 wis = 10,
 cha = 8
]
\DndMonsterDetails[
skills = {Arcana +3, Insight +4, Perception +2, Stealth +5},
senses = {truesight 120 ft., passive perception 12},
languages = {Undercommon},
challenge = 2,
]
\DndMonsterAction{Keen Sight}
The nothic has advantage on Wisdom (Perception) checks that rely on sight.

\DndMonsterSection{Actions}
\DndMonsterAction{Multiattack}
The nothic makes two claw attacks.

\DndMonsterAction{Claw}
\textsl{Melee Weapon Attack:} 4 to hit, reach 5 ft., one target. \textsl{Hit:} 6 (1d6 + 3) slashing damage.

\DndMonsterAction{Rotting Gaze}
The nothic targets one creature it can see within 30 feet of it. The target must succeed on a DC 12 Constitution saving throw against this magic or take 10 (3d6) necrotic damage.

\DndMonsterAction{Weird Insight}
The nothic targets one creature it can see within 30 feet of it. The target must contest its Charisma (Deception) check against the nothic's Wisdom (Insight) check. If the nothic wins, it magically learns one fact or secret about the target. The target automatically wins if it is immune to being charmed.

\end{DndMonster}



\begin{DndMonster}[float=t]{Ogre Zombie}
\DndMonsterType{Large undead, neutral evil}
\DndMonsterBasics[
armor-class = {8},
hit-points = {\DndDice{9d10 + 36}},
speed = {30 ft.}
]
\DndMonsterAbilityScores[
 str = 19,
 dex = 6,
 con = 18,
 int = 3,
 wis = 6,
 cha = 5
]
\DndMonsterDetails[
saving-throws = {Wis +0},
damage-immunities = {poison},
condition-immunities = {poisoned},
senses = {darkvision 60 ft., passive perception 8},
languages = {understands Common and Giant but can't speak},
challenge = 2,
]
\DndMonsterAction{Undead Fortitude}
If damage reduces the zombie to 0 hit points, it must make a Constitution saving throw with a DC of 5 + the damage taken, unless the damage is radiant or from a critical hit. On a success, the zombie drops to 1 hit point instead.

\DndMonsterSection{Actions}
\DndMonsterAction{Morningstar}
\textsl{Melee Weapon Attack:} 6 to hit, reach 5 ft., one target. \textsl{Hit:} 13 (2d8 + 4) bludgeoning damage.

\end{DndMonster}



\begin{DndMonster}[float=t]{Phisarazu Spyder-Fiend}
\DndMonsterType{Large fiend (demon), chaotic evil}
\DndMonsterBasics[
armor-class = {17 (natural armor)},
hit-points = {\DndDice{20d10 + 60}},
speed = {40 ft., climb 40 ft.}
]
\DndMonsterAbilityScores[
 str = 19,
 dex = 14,
 con = 17,
 int = 11,
 wis = 14,
 cha = 13
]
\DndMonsterDetails[
saving-throws = {Dex +7, Con +8, Wis +7},
skills = {Perception +7, Stealth +7},
damage-resistances = {cold; fire; lightning; bludgeoning, piercing, slashing from nonmagical attacks},
damage-immunities = {poison},
condition-immunities = {poisoned},
senses = {truesight 120 ft., passive perception 17},
languages = {Abyssal, Common, telepathy 120 ft.},
challenge = 13,
]
\DndMonsterAction{Magic Resistance}
The phisarazu has advantage on saving throws against spells and other magical effects.

\DndMonsterAction{Spider Climb}
The phisarazu can climb difficult surfaces, including upside down on ceilings, without needing to make an ability check.

\DndMonsterAction{Web Sense}
When in contact with a web, the phisarazu knows the exact location of any other creature in contact with the same web.

\DndMonsterAction{Web Walker}
The phisarazu ignores movement restrictions caused by webbing.

\DndMonsterSection{Actions}
\DndMonsterAction{Multiattack}
The phisarazu makes one Bite attack and two Claw attacks. It can replace one of these attacks with Scintillating Spray if available.

\DndMonsterAction{Bite}
\textsl{Melee Weapon Attack:} 9 to hit, reach 5 ft., one target. \textsl{Hit:} 15 (2d10 + 4) piercing damage plus 9 (2d8) poison damage, and the target has the poisoned condition until the start of the phisarazu's next turn.

\DndMonsterAction{Claw}
\textsl{Melee Weapon Attack:} 9 to hit, reach 5 ft., one target. \textsl{Hit:} 22 (4d8 + 4) slashing damage.

\DndMonsterAction{Scintillating Spray (Recharge 5\textendash 6)}
The phisarazu expels shimmering webs in a 60-foot cone. Creatures and objects in that area are outlined by the glittering webs for 1 minute, during which time they emit dim light for 10 feet and can't benefit from the invisible condition. Additionally, creatures in that area must succeed on a DC 16 Wisdom saving throw or have the stunned condition for 1 minute. A stunned creature can repeat the saving throw at the end of each of its turns, ending the effect on itself on a success.

\DndMonsterSection{Bonus Actions}
\DndMonsterAction{Change Shape}
The phisarazu transforms into a crab, drider, or giant crab, or returns to its true form. Its game statistics, except for its size, are the same in each form. Any equipment it is wearing or carrying isn't transformed.

\end{DndMonster}



\begin{DndMonster}[float=t]{Pit Fiend}
\DndMonsterType{Large fiend (devil), lawful evil}
\DndMonsterBasics[
armor-class = {19 (natural armor)},
hit-points = {\DndDice{24d10 + 168}},
speed = {30 ft., fly 60 ft.}
]
\DndMonsterAbilityScores[
 str = 26,
 dex = 14,
 con = 24,
 int = 22,
 wis = 18,
 cha = 24
]
\DndMonsterDetails[
saving-throws = {Dex +8, Con +13, Wis +10},
damage-resistances = {cold; bludgeoning, piercing, slashing from nonmagical attacks that aren't silvered},
damage-immunities = {fire, poison},
condition-immunities = {poisoned},
senses = {truesight 120 ft., passive perception 14},
languages = {Infernal, telepathy 120 ft.},
challenge = 20,
]
\DndMonsterAction{Fear Aura}
Any creature hostile to the pit fiend that starts its turn within 20 feet of the pit fiend must make a DC 21 Wisdom saving throw, unless the pit fiend is incapacitated. On a failed save, the creature is frightened until the start of its next turn. If a creature's saving throw is successful, the creature is immune to the pit fiend's Fear Aura for the next 24 hours.

\DndMonsterAction{Magic Resistance}
The pit fiend has advantage on saving throws against spells and other magical effects.

\DndMonsterAction{Magic Weapons}
The pit fiend's weapon attacks are magical.

\DndMonsterAction{Innate Spellcasting}
The pit fiend's spellcasting ability is Charisma (spell save DC 21). The pit fiend can innately cast the following spells, requiring no material components:
\begin{DndMonsterSpells}
  \DndInnateSpellLevel{detect magic}
  \DndInnateSpellLevel[3]{hold monster}
\end{DndMonsterSpells}
\DndMonsterSection{Actions}
\DndMonsterAction{Multiattack}
The pit fiend makes four attacks: one with its bite, one with its claw, one with its mace, and one with its tail.

\DndMonsterAction{Bite}
\textsl{Melee Weapon Attack:} 14 to hit, reach 5 ft., one target. \textsl{Hit:} 22 (4d6 + 8) piercing damage. The target must succeed on a DC 21 Constitution saving throw or become poisoned. While poisoned in this way, the target can't regain hit points, and it takes 21 (6d6) poison damage at the start of each of its turns. The poisoned target can repeat the saving throw at the end of each of its turns, ending the effect on itself on a success.

\DndMonsterAction{Claw}
\textsl{Melee Weapon Attack:} 14 to hit, reach 10 ft., one target. \textsl{Hit:} 17 (2d8 + 8) slashing damage.

\DndMonsterAction{Mace}
\textsl{Melee Weapon Attack:} 14 to hit, reach 10 ft., one target. \textsl{Hit:} 15 (2d6 + 8) bludgeoning damage plus 21 (6d6) fire damage.

\DndMonsterAction{Tail}
\textsl{Melee Weapon Attack:} 14 to hit, reach 10 ft., one target. \textsl{Hit:} 24 (3d10 + 8) bludgeoning damage.

\end{DndMonster}



\begin{DndMonster}[float=t]{Priest}
\DndMonsterType{Medium humanoid (any race), any alignment}
\DndMonsterBasics[
armor-class = {13 (\textit{chain shirt})},
hit-points = {\DndDice{5d8 + 5}},
speed = {30 ft.}
]
\DndMonsterAbilityScores[
 str = 10,
 dex = 10,
 con = 12,
 int = 13,
 wis = 16,
 cha = 13
]
\DndMonsterDetails[
skills = {Medicine +7, Persuasion +3, Religion +5},
languages = {any two languages},
challenge = 2,
]
\DndMonsterAction{Divine Eminence}
As a bonus action, the priest can expend a spell slot to cause its melee weapon attacks to magically deal an extra 10 (3d6) radiant damage to a target on a hit. This benefit lasts until the end of the turn. If the priest expends a spell slot of 2nd level or higher, the extra damage increases by 1d6 for each level above 1st.

\DndMonsterAction{Spellcasting}
The priest is a 5th-level spellcaster. Its spellcasting ability is Wisdom (spell save DC 13, 5 to hit with spell attacks). The priest has the following cleric spells prepared:
\begin{DndMonsterSpells}
  \DndMonsterSpellLevel{light}
   \DndMonsterSpellLevel[1][4]{cure wounds}
   \DndMonsterSpellLevel[2][3]{lesser restoration}
   \DndMonsterSpellLevel[3][2]{dispel magic}
\end{DndMonsterSpells}
\DndMonsterSection{Actions}
\DndMonsterAction{Mace}
\textsl{Melee Weapon Attack:} 2 to hit, reach 5 ft., one target. \textsl{Hit:} 3 (1d6) bludgeoning damage.

\end{DndMonster}



\begin{DndMonster}[float=t]{Priest of Osybus}
\DndMonsterType{Medium humanoid, U}
\DndMonsterBasics[
armor-class = {14 (natural armor)},
hit-points = {\DndDice{8d8 + 24}},
speed = {30 ft.}
]
\DndMonsterAbilityScores[
 str = 10,
 dex = 14,
 con = 16,
 int = 18,
 wis = 17,
 cha = 11
]
\DndMonsterDetails[
saving-throws = {Int +7, Wis +6, Cha +3},
condition-immunities = {frightened},
senses = {darkvision 120 ft., passive perception 13},
languages = {any three languages},
challenge = 6,
]
\DndMonsterAction{Tattoo of Osybus}
If the priest drops to 0 hit points, roll on the Boons of Undeath table for the boon the priest receives. The priest dies if it receives a boon it already has. If it receives a new boon, it revives at the start of its next turn with half its hit points restored, and its creature type is now Undead.

To prevent this revival, the Tattoo of Osybus on the priest's body must be destroyed. The tattoo is invulnerable while the priest has at least 1 hit point. The tattoo is otherwise an object with AC 15, and it is immune to poison and psychic damage. It has 15 hit points, but it regains all its hit points at the end of every combatant's turn.

\DndMonsterSection{Actions}
\DndMonsterAction{Multiattack}
The priest attacks twice.

\DndMonsterAction{Soul Blade}
\textsl{Melee Weapon Attack:} 5 to hit, reach 5 ft., one target. \textsl{Hit:} 7 (2d4 + 2) piercing damage, and if the target is a creature, it is paralyzed until the start of the priest's next turn. If this damage reduces a Medium or smaller creature to 0 hit points, the creature dies, and its soul is trapped in the priest's body, manifesting as a shadowy Soul Tattoo on the priest. The soul is freed if the priest dies.

\DndMonsterAction{Necrotic Bolt}
\textsl{Ranged Spell Attack:} 7 to hit, range 120 ft., one target. \textsl{Hit:} 17 (3d8 + 4) necrotic damage, and the target can't regain hit points until the start of the priest's next turn.

\DndMonsterSection{Bonus Actions}
\DndMonsterAction{Soul Tattoo (Recharge 5\textendash 6)}
The priest touches one of the Soul Tattoos on its body. The tattoo vanishes as the trapped soul manifests as a shadowy creature that appears in an unoccupied space the priest can see within 30 feet of it. The creature has the size and silhouette of its original body, but it otherwise uses the stat block of a shadow.

The shadow obeys the priest's mental commands (no action required) and takes its turn immediately after the priest. If the creature is within 5 feet of the priest, it can turn back into a tattoo as an action, reappearing on the priest's flesh and regaining all its hit points.

\end{DndMonster}



\begin{DndMonster}[float=t]{Purple Worm}
\DndMonsterType{Gargantuan monstrosity, unaligned}
\DndMonsterBasics[
armor-class = {18 (natural armor)},
hit-points = {\DndDice{15d20 + 90}},
speed = {50 ft., burrow 30 ft.}
]
\DndMonsterAbilityScores[
 str = 28,
 dex = 7,
 con = 22,
 int = 1,
 wis = 8,
 cha = 4
]
\DndMonsterDetails[
saving-throws = {Con +11, Wis +4},
senses = {blindsight 30 ft., tremorsense 60 ft., passive perception 9},
challenge = 15,
]
\DndMonsterAction{Tunneler}
The worm can burrow through solid rock at half its burrow speed and leaves a 10-foot-diameter tunnel in its wake.

\DndMonsterSection{Actions}
\DndMonsterAction{Multiattack}
The worm makes two attacks: one with its bite and one with its stinger.

\DndMonsterAction{Bite}
\textsl{Melee Weapon Attack:} 14 to hit, reach 10 ft., one target. \textsl{Hit:} 22 (3d8 + 9) piercing damage. If the target is a Large or smaller creature, it must succeed on a DC 19 Dexterity saving throw or be swallowed by the worm. A swallowed creature is blinded and restrained, it has total cover against attacks and other effects outside the worm, and it takes 21 (6d6) acid damage at the start of each of the worm's turns.

If the worm takes 30 damage or more on a single turn from a creature inside it, the worm must succeed on a DC 21 Constitution saving throw at the end of that turn or regurgitate all swallowed creatures, which fall prone in a space within 10 feet of the worm. If the worm dies, a swallowed creature is no longer restrained by it and can escape from the corpse by using 20 feet of movement, exiting prone.

\DndMonsterAction{Tail Stinger}
\textsl{Melee Weapon Attack:} 14 to hit, reach 10 ft., one creature. \textsl{Hit:} 19 (3d6 + 9) piercing damage, and the target must make a DC 19 Constitution saving throw, taking 42 (12d6) poison damage on a failed save, or half as much damage on a successful one.

\end{DndMonster}


\clearpage

\begin{DndMonster}[float=t]{Quasit}
\DndMonsterType{Tiny fiend (demon), chaotic evil}
\DndMonsterBasics[
armor-class = {13},
hit-points = {\DndDice{3d4}},
speed = {40 ft.}
]
\DndMonsterAbilityScores[
 str = 5,
 dex = 17,
 con = 10,
 int = 7,
 wis = 10,
 cha = 10
]
\DndMonsterDetails[
skills = {Stealth +5},
damage-resistances = {cold; fire; lightning; bludgeoning, piercing, slashing from nonmagical attacks},
damage-immunities = {poison},
condition-immunities = {poisoned},
senses = {darkvision 120 ft., passive perception 10},
languages = {Abyssal, Common},
challenge = 1,
]
\DndMonsterAction{Shapechanger}
The quasit can use its action to polymorph into a beast form that resembles a bat (speed 10 feet fly 40 ft.), a centipede (40 ft., climb 40 ft.), or a toad (40 ft., swim 40 ft.), or back into its true form. Its statistics are the same in each form, except for the speed changes noted. Any equipment it is wearing or carrying isn't transformed. It reverts to its true form if it dies.

\DndMonsterAction{Magic Resistance}
The quasit has advantage on saving throws against spells and other magical effects.

\DndMonsterSection{Actions}
\DndMonsterAction{Claw (Bite in Beast Form)}
\textsl{Melee Weapon Attack:} 4 to hit, reach 5 ft., one target. \textsl{Hit:} 5 (1d4 + 3) piercing damage, and the target must succeed on a DC 10 Constitution saving throw or take 5 (2d4) poison damage and become poisoned for 1 minute. The target can repeat the saving throw at the end of each of its turns, ending the effect on itself on a success.

\DndMonsterAction{Scare (1/Day)}
One creature of the quasit's choice within 20 feet of it must succeed on a DC 10 Wisdom saving throw or be frightened for 1 minute. The target can repeat the saving throw at the end of each of its turns, with disadvantage if the quasit is within line of sight, ending the effect on itself on a success.

\DndMonsterAction{Invisibility}
The quasit magically turns invisible until it attacks or uses Scare, or until its concentration ends (as if concentration on a spell). Any equipment the quasit wears or carries is invisible with it.

\end{DndMonster}



\begin{DndMonster}[float=t]{Quavilithku Spyder-Fiend}
\DndMonsterType{Large fiend (demon), chaotic evil}
\DndMonsterBasics[
armor-class = {19 (natural armor)},
hit-points = {\DndDice{27d10 + 108}},
speed = {40 ft., climb 40 ft.}
]
\DndMonsterAbilityScores[
 str = 18,
 dex = 16,
 con = 19,
 int = 17,
 wis = 14,
 cha = 12
]
\DndMonsterDetails[
saving-throws = {Dex +9, Con +10, Wis +8},
skills = {Investigation +9, Perception +8, Stealth +9},
damage-resistances = {cold, fire, lightning},
damage-immunities = {acid, poison},
condition-immunities = {poisoned},
senses = {truesight 60 ft., passive perception 18},
languages = {Abyssal, Common, telepathy 120 ft.},
challenge = 17,
]
\DndMonsterAction{Magic Resistance}
The quavilithku has advantage on saving throws against spells and other magical effects.

\DndMonsterAction{Spider Climb}
The quavilithku can climb difficult surfaces, including upside down on ceilings, without needing to make an ability check.

\DndMonsterAction{Web Sense}
When in contact with a web, the quavilithku knows the exact location of any other creature in contact with the same web.

\DndMonsterAction{Web Walker}
The quavilithku ignores movement restrictions caused by webbing.

\DndMonsterSection{Actions}
\DndMonsterAction{Multiattack}
The quavilithku makes two Bite attacks.

\DndMonsterAction{Bite}
\textsl{Melee Weapon Attack:} 10 to hit, reach 5 ft., one target. \textsl{Hit:} 15 (2d10 + 4) piercing damage plus 17 (5d6) poison damage. If the target is a creature, it must succeed on a DC 18 Constitution saving throw or have the poisoned condition for 1 minute. While poisoned in this way, a creature can't regain hit points. A poisoned creature can repeat the saving throw at the end of each of its turns, ending the effect on itself on a success.

\DndMonsterAction{Dissolving Web (Recharge 5\textendash 6)}
The quavilithku expels acid-drenched webs in a 90-foot cone. Each creature in that area must make a DC 18 Constitution saving throw, taking 44 (8d10) acid damage on a failed save or half as much damage on a successful one. Nonmagical objects in the area that aren't being worn or carried take 44 (8d10) acid damage.

\DndMonsterSection{Bonus Actions}
\DndMonsterAction{Assess Weakness}
The quavilithku sizes up a creature it can see within 40 feet of itself. Until the start of the quavilithku's next turn, it has advantage on attack rolls against the creature.

\end{DndMonster}



\begin{DndMonster}[float=t]{Raklupis Spyder-Fiend}
\DndMonsterType{Large fiend (demon), chaotic evil}
\DndMonsterBasics[
armor-class = {19 (natural armor)},
hit-points = {\DndDice{28d10 + 56}},
speed = {40 ft., climb 40 ft.}
]
\DndMonsterAbilityScores[
 str = 17,
 dex = 20,
 con = 14,
 int = 18,
 wis = 16,
 cha = 23
]
\DndMonsterDetails[
saving-throws = {Dex +11, Con +8, Wis +9},
skills = {Perception +9, Stealth +11},
damage-immunities = {cold, fire, lightning, poison},
condition-immunities = {poisoned},
senses = {truesight 120 ft., passive perception 19},
languages = {Abyssal, Common, telepathy 120 ft.},
challenge = 19,
]
\DndMonsterAction{Magic Resistance}
The raklupis has advantage on saving throws against spells and other magical effects.

\DndMonsterAction{Spider Climb}
The raklupis can climb difficult surfaces, including upside down on ceilings, without needing to make an ability check.

\DndMonsterAction{Web Sense}
When in contact with a web, the raklupis knows the exact location of any other creature in contact with the same web.

\DndMonsterAction{Web Walker}
The raklupis ignores movement restrictions caused by webbing.

\DndMonsterAction{Spellcasting}
The raklupis casts one of the following spells, requiring no material components and using Charisma as the spellcasting ability (spell save DC 20).
\begin{DndMonsterSpells}
  \DndInnateSpellLevel{Disguise Self}
  \DndInnateSpellLevel[2]{Darkness}
\end{DndMonsterSpells}
\DndMonsterSection{Actions}
\DndMonsterAction{Multiattack}
The raklupis makes a Bite attack and two Serrated Sword attacks. It can use Venom Globe in place of one of these attacks.

\DndMonsterAction{Bite}
\textsl{Melee Weapon Attack:} 11 to hit, reach 5 ft., one target. \textsl{Hit:} 16 (2d10 + 5) piercing damage plus 18 (4d8) poison damage. If the target is a creature, it must succeed on a DC 20 Constitution saving throw or have the poisoned condition for 1 minute. While poisoned in this way, a creature has the incapacitated condition and can't regain hit points. A poisoned creature can repeat the saving throw at the end of each of its turns, ending the effect on itself on a success.

\DndMonsterAction{Serrated Sword}
\textsl{Melee Weapon Attack:} 11 to hit, reach 5 ft., one target. \textsl{Hit:} 19 (4d6 + 5) slashing damage plus 18 (4d8) poison damage.

\DndMonsterAction{Venom Globe}
\textsl{Ranged Weapon Attack:} 11 to hit, range 60/180 ft., one target. \textsl{Hit:} 45 (10d8) poison damage.

\DndMonsterSection{Bonus Actions}
\DndMonsterAction{Demand Loyalty}
The raklupis magically ends the charmed and frightened conditions on itself and on any number of allies within 60 feet of itself.

\end{DndMonster}



\begin{DndMonster}[float=t]{Rakshasa}
\DndMonsterType{Medium fiend, lawful evil}
\DndMonsterBasics[
armor-class = {16 (natural armor)},
hit-points = {\DndDice{13d8 + 52}},
speed = {40 ft.}
]
\DndMonsterAbilityScores[
 str = 14,
 dex = 17,
 con = 18,
 int = 13,
 wis = 16,
 cha = 20
]
\DndMonsterDetails[
skills = {Deception +10, Insight +8},
damage-vulnerabilities = {piercing from magic weapons wielded by good creatures},
damage-immunities = {bludgeoning, piercing, slashing from nonmagical attacks},
senses = {darkvision 60 ft., passive perception 13},
languages = {Common, Infernal},
challenge = 13,
]
\DndMonsterAction{Limited Magic Immunity}
The rakshasa can't be affected or detected by spells of 6th level or lower unless it wishes to be. It has advantage on saving throws against all other spells and magical effects.

\DndMonsterAction{Innate Spellcasting}
The rakshasa's innate spellcasting ability is Charisma (spell save DC 18, 10 to hit with spell attacks). The rakshasa can innately cast the following spells, requiring no material components:
\begin{DndMonsterSpells}
  \DndInnateSpellLevel{detect thoughts}
  \DndInnateSpellLevel[3]{charm person}
  \DndInnateSpellLevel[1]{dominate person}
\end{DndMonsterSpells}
\DndMonsterSection{Actions}
\DndMonsterAction{Multiattack}
The rakshasa makes two claw attacks.

\DndMonsterAction{Claw}
\textsl{Melee Weapon Attack:} 7 to hit, reach 5 ft., one target. \textsl{Hit:} 9 (2d6 + 2) slashing damage, and the target is cursed if it is a creature. The magical curse takes effect whenever the target takes a short or long rest, filling the target's thoughts with horrible images and dreams. The cursed target gains no benefit from finishing a short or long rest. The curse lasts until it is lifted by a \textit{remove curse} spell or similar magic.

\end{DndMonster}



\begin{DndMonster}[float=t]{Raven}
\DndMonsterType{Tiny beast, unaligned}
\DndMonsterBasics[
armor-class = {12},
hit-points = {\DndDice{1d4 - 1}},
speed = {10 ft., fly 50 ft.}
]
\DndMonsterAbilityScores[
 str = 2,
 dex = 14,
 con = 8,
 int = 2,
 wis = 12,
 cha = 6
]
\DndMonsterDetails[
skills = {Perception +3},
challenge = 0,
]
\DndMonsterAction{Mimicry}
The raven can mimic simple sounds it has heard, such as a person whispering, a baby crying, or an animal chittering. A creature that hears the sounds can tell they are imitations with a successful DC 10 Wisdom (Insight) check.

\DndMonsterSection{Actions}
\DndMonsterAction{Beak}
\textsl{Melee Weapon Attack:} 4 to hit, reach 5 ft., one target. \textsl{Hit:} 1 piercing damage.

\end{DndMonster}



\begin{DndMonster}[float=t]{Red Abishai}
\DndMonsterType{Medium fiend (devil), lawful evil}
\DndMonsterBasics[
armor-class = {22 (natural armor)},
hit-points = {\DndDice{34d8 + 136}},
speed = {30 ft., fly 50 ft.}
]
\DndMonsterAbilityScores[
 str = 23,
 dex = 16,
 con = 19,
 int = 14,
 wis = 15,
 cha = 19
]
\DndMonsterDetails[
saving-throws = {Str +12, Con +10, Wis +8},
skills = {Intimidation +10, Perception +8},
damage-resistances = {cold; bludgeoning, piercing, slashing from nonmagical attacks that aren't silvered},
damage-immunities = {fire, poison},
condition-immunities = {frightened, poisoned},
senses = {darkvision 120 ft., passive perception 18},
languages = {Draconic, Infernal, telepathy 120 ft.},
challenge = 19,
]
\DndMonsterAction{Devil's Sight}
Magical darkness doesn't impede the abishai's darkvision.

\DndMonsterAction{Magic Resistance}
The abishai has advantage on saving throws against spells and other magical effects.

\DndMonsterSection{Actions}
\DndMonsterAction{Multiattack}
The abishai makes one Bite attack and one Claw attack, and it can use Frightful Presence or Incite Fanaticism.

\DndMonsterAction{Bite}
\textsl{Melee Weapon Attack:} 12 to hit, reach 5 ft., one target. \textsl{Hit:} 22 (3d10 + 6) piercing damage plus 38 (7d10) fire damage.

\DndMonsterAction{Claw}
\textsl{Melee Weapon Attack:} 12 to hit, reach 5 ft., one target. \textsl{Hit:} 17 (2d10 + 6) force damage plus 11 (2d10) fire damage.

\DndMonsterAction{Frightful Presence}
Each creature of the abishai's choice that is within 120 feet and aware of the abishai must succeed on a DC 18 Wisdom saving throw or become frightened of it for 1 minute. A creature can repeat the saving throw at the end of each of its turns, ending the effect on itself on a success. If a creature's saving throw is successful or the effect ends for it, the creature is immune to the abishai's Frightful Presence for the next 24 hours.

\DndMonsterAction{Incite Fanaticism}
The abishai chooses up to four other creatures within 60 feet of it that can see it. Until the start of the abishai's next turn, each of those creatures makes attack rolls with advantage and can't be frightened.

\DndMonsterAction{Power of the Dragon Queen}
The abishai targets one Dragon it can see within 120 feet of it. The Dragon must make a DC 18 Charisma saving throw. A chromatic dragon makes this save with disadvantage. On a successful save, the target is immune to the abishai's Power of the Dragon Queen for 1 hour. On a failed save, the target is charmed by the abishai for 1 hour. While charmed in this way, the target regards the abishai as a trusted friend to be heeded and protected. This effect ends if the abishai or its companions deal damage to the target.

\end{DndMonster}



\begin{DndMonster}[float*=b,width=\textwidth + 8pt]{Relentless Impaler}
\begin{multicols}{2}
\DndMonsterType{Large fiend, neutral evil}
\DndMonsterBasics[
armor-class = {16 (natural armor)},
hit-points = {\DndDice{16d10 + 96}},
speed = {30 ft.}
]
\DndMonsterAbilityScores[
 str = 23,
 dex = 16,
 con = 22,
 int = 12,
 wis = 15,
 cha = 18
]
\DndMonsterDetails[
saving-throws = {Str +11, Dex +8, Cha +9},
skills = {Athletics +11, Perception +7, Survival +7},
condition-immunities = {charmed, exhaustion, frightened},
senses = {darkvision 120 ft., passive perception 17},
languages = {understands all languages but can't speak},
challenge = 15,
]
\DndMonsterAction{Bloodheart Stake}
The impaler is magically bound to the ceremonial stake and the sacrificed corpse the ritual caster used to create it. If the impaler is reduced to 0 hit points, it disappears, then re-forms 1d8 hours later in the nearest unoccupied space to the stake and regains all its hit points. The impaler dies only if it is reduced to 0 hit points while either the ceremonial stake is removed from the sacrifice's corpse or the impaler is on a different plane of existence from that corpse.

\DndMonsterAction{Legendary Resistance (3/Day)}
If the impaler fails a saving throw, it can choose to succeed instead.

\DndMonsterSection{Actions}
\DndMonsterAction{Multiattack}
The impaler makes one Spike attack and two Wicked Spear attacks.

\DndMonsterAction{Spike}
\textsl{Melee Weapon Attack:} 11 to hit, reach 5 ft., one target. \textsl{Hit:} 15 (2d8 + 6) piercing damage, and the target's speed is halved until the start of the impaler's next turn.

\DndMonsterAction{Wicked Spear}
\textsl{Melee or Ranged Weapon Attack:} 11 to hit, reach 10 ft. or range 20/40 ft., one target. \textsl{Hit:} 13 (2d6 + 6) piercing damage plus 13 (3d8) necrotic damage.

\DndMonsterAction{Spike Burst (Recharge 5\textendash 6)}
Twisted, spectral spikes shoot out from the impaler's body. Each creature within 30 feet of the impaler must make a DC 19 Dexterity saving throw, taking 40 (9d8) force damage on a failed save or half as much damage on a successful one.

\DndMonsterSection{Legendary Actions}
Relentless Impaler can take 3 legendary actions, choosing from the options below. Only one legendary action can be used at a time and only at the end of another creature's turn. Relentless Impaler regains spent legendary actions at the start of its turn.

\begin{DndMonsterLegendaryActions}
\DndMonsterLegendaryAction{Speed Spike}{The impaler teleports up to 30 feet to an unoccupied space it can see. It can then make a Spike attack.
}
\DndMonsterLegendaryAction{Deepen Wounds (Costs 2 Actions)}{Each creature whose speed is currently reduced by the impaler's Spike attack takes 18 (4d8) necrotic damage.
}
\end{DndMonsterLegendaryActions}
\end{multicols}
\end{DndMonster}



\begin{DndMonster}[float=t]{Roc}
\DndMonsterType{Gargantuan monstrosity, unaligned}
\DndMonsterBasics[
armor-class = {15 (natural armor)},
hit-points = {\DndDice{16d20 + 80}},
speed = {20 ft., fly 120 ft.}
]
\DndMonsterAbilityScores[
 str = 28,
 dex = 10,
 con = 20,
 int = 3,
 wis = 10,
 cha = 9
]
\DndMonsterDetails[
saving-throws = {Dex +4, Con +9, Wis +4, Cha +3},
skills = {Perception +4},
challenge = 11,
]
\DndMonsterAction{Keen Sight}
The roc has advantage on Wisdom (Perception) checks that rely on sight.

\DndMonsterSection{Actions}
\DndMonsterAction{Multiattack}
The roc makes two attacks: one with its beak and one with its talons.

\DndMonsterAction{Beak}
\textsl{Melee Weapon Attack:} 13 to hit, reach 10 ft., one target. \textsl{Hit:} 27 (4d8 + 9) piercing damage.

\DndMonsterAction{Talons}
\textsl{Melee Weapon Attack:} 13 to hit, reach 5 ft., one target. \textsl{Hit:} 23 (4d6 + 9) slashing damage, and the target is grappled (escape DC 19). Until this grapple ends, the target is restrained, and the roc can't use its talons on another target.

\end{DndMonster}



\begin{DndMonster}[float=t]{Shambling Mound}
\DndMonsterType{Large plant, unaligned}
\DndMonsterBasics[
armor-class = {15 (natural armor)},
hit-points = {\DndDice{16d10 + 48}},
speed = {20 ft., swim 20 ft.}
]
\DndMonsterAbilityScores[
 str = 18,
 dex = 8,
 con = 16,
 int = 5,
 wis = 10,
 cha = 5
]
\DndMonsterDetails[
skills = {Stealth +2},
damage-resistances = {cold, fire},
damage-immunities = {lightning},
condition-immunities = {blinded, deafened, exhaustion},
senses = {blindsight 60 ft. (blind beyond this radius), passive perception 10},
challenge = 5,
]
\DndMonsterAction{Lightning Absorption}
Whenever the shambling mound is subjected to lightning damage, it takes no damage and regains a number of hit points equal to the lightning damage dealt.

\DndMonsterSection{Actions}
\DndMonsterAction{Multiattack}
The shambling mound makes two slam attacks. If both attacks hit a Medium or smaller target, the target is grappled (escape DC 14), and the shambling mound uses its Engulf on it.

\DndMonsterAction{Slam}
\textsl{Melee Weapon Attack:} 7 to hit, reach 5 ft., one target. \textsl{Hit:} 13 (2d8 + 4) bludgeoning damage.

\DndMonsterAction{Engulf}
The shambling mound engulfs a Medium or smaller creature grappled by it. The engulfed target is blinded, restrained, and unable to breathe, and it must succeed on a DC 14 Constitution saving throw at the start of each of the mound's turns or take 13 (2d8 + 4) bludgeoning damage. If the mound moves, the engulfed target moves with it. The mound can have only one creature engulfed at a time.

\end{DndMonster}



\begin{DndMonster}[float=t]{Shield Guardian}
\DndMonsterType{Large construct, unaligned}
\DndMonsterBasics[
armor-class = {17 (natural armor)},
hit-points = {\DndDice{15d10 + 60}},
speed = {30 ft.}
]
\DndMonsterAbilityScores[
 str = 18,
 dex = 8,
 con = 18,
 int = 7,
 wis = 10,
 cha = 3
]
\DndMonsterDetails[
damage-immunities = {poison},
condition-immunities = {charmed, exhaustion, frightened, paralyzed, poisoned},
senses = {blindsight 10 ft., darkvision 60 ft., passive perception 10},
languages = {understands commands given in any language but can't speak},
challenge = 7,
]
\DndMonsterAction{Bound}
The shield guardian is magically bound to an amulet. As long as the guardian and its amulet are on the same plane of existence, the amulet's wearer can telepathically call the guardian to travel to it, and the guardian knows the distance and direction to the amulet. If the guardian is within 60 feet of the amulet's wearer, half of any damage the wearer takes (rounded up) is transferred to the guardian.

\DndMonsterAction{Regeneration}
The shield guardian regains 10 hit points at the start of its turn if it has at least 1 hit point.

\DndMonsterAction{Spell Storing}
A spellcaster who wears the shield guardian's amulet can cause the guardian to store one spell of 4th level or lower. To do so, the wearer must cast the spell on the guardian. The spell has no effect but is stored within the guardian. When commanded to do so by the wearer or when a situation arises that was predefined by the spellcaster, the guardian casts the stored spell with any parameters set by the original caster, requiring no components. When the spell is cast or a new spell is stored, any previously stored spell is lost.

\DndMonsterSection{Actions}
\DndMonsterAction{Multiattack}
The guardian makes two fist attacks.

\DndMonsterAction{Fist}
\textsl{Melee Weapon Attack:} 7 to hit, reach 5 ft., one target. \textsl{Hit:} 11 (2d6 + 4) bludgeoning damage.

\DndMonsterSection{Reactions}
\DndMonsterAction{Shield}
When a creature makes an attack against the wearer of the guardian's amulet, the guardian grants a +2 bonus to the wearer's AC if the guardian is within 5 feet of the wearer.

\end{DndMonster}



\begin{DndMonster}[float=t]{Shrieker}
\DndMonsterType{Medium plant, unaligned}
\DndMonsterBasics[
armor-class = {5},
hit-points = {\DndDice{3d8}},
speed = {0 ft.}
]
\DndMonsterAbilityScores[
 str = 1,
 dex = 1,
 con = 10,
 int = 1,
 wis = 3,
 cha = 1
]
\DndMonsterDetails[
condition-immunities = {blinded, deafened, frightened},
senses = {blindsight 30 ft. (blind beyond this radius), passive perception 6},
challenge = 0,
]
\DndMonsterAction{False Appearance}
While the shrieker remains motionless, it is indistinguishable from an ordinary fungus.

\DndMonsterSection{Reactions}
\DndMonsterAction{Shriek}
When bright light or a creature is within 30 feet of the shrieker, it emits a shriek audible within 300 feet of it. The shrieker continues to shriek until the disturbance moves out of range and for 1d4 of the shrieker's turns afterward.

\end{DndMonster}



\begin{DndMonster}[float=t]{Skeleton}
\DndMonsterType{Medium undead, lawful evil}
\DndMonsterBasics[
armor-class = {13 (armor scraps)},
hit-points = {\DndDice{2d8 + 4}},
speed = {30 ft.}
]
\DndMonsterAbilityScores[
 str = 10,
 dex = 14,
 con = 15,
 int = 6,
 wis = 8,
 cha = 5
]
\DndMonsterDetails[
damage-vulnerabilities = {bludgeoning},
damage-immunities = {poison},
condition-immunities = {exhaustion, poisoned},
senses = {darkvision 60 ft., passive perception 9},
languages = {understands all languages it spoke in life but can't speak},
challenge = 1/4,
]
\DndMonsterSection{Actions}
\DndMonsterAction{Shortsword}
\textsl{Melee Weapon Attack:} 4 to hit, reach 5 ft., one target. \textsl{Hit:} 5 (1d6 + 2) piercing damage.

\DndMonsterAction{Shortbow}
\textsl{Ranged Weapon Attack:} 4 to hit, range 80/320 ft., one target. \textsl{Hit:} 5 (1d6 + 2) piercing damage.

\end{DndMonster}



\begin{DndMonster}[float=t]{Specter}
\DndMonsterType{Medium undead, chaotic evil}
\DndMonsterBasics[
armor-class = {12},
hit-points = {\DndDice{5d8}},
speed = {0 ft., fly 50 ft. \textit{(hover)}}
]
\DndMonsterAbilityScores[
 str = 1,
 dex = 14,
 con = 11,
 int = 10,
 wis = 10,
 cha = 11
]
\DndMonsterDetails[
damage-resistances = {acid; cold; fire; lightning; thunder; bludgeoning, piercing, slashing from nonmagical attacks},
damage-immunities = {necrotic, poison},
condition-immunities = {charmed, exhaustion, grappled, paralyzed, petrified, poisoned, prone, restrained, unconscious},
senses = {darkvision 60 ft., passive perception 10},
languages = {understands all languages it knew in life but can't speak},
challenge = 1,
]
\DndMonsterAction{Incorporeal Movement}
The specter can move through other creatures and objects as if they were difficult terrain. It takes 5 (1d10) force damage if it ends its turn inside an object.

\DndMonsterAction{Sunlight Sensitivity}
While in sunlight, the specter has disadvantage on attack rolls, as well as on Wisdom (Perception) checks that rely on sight.

\DndMonsterSection{Actions}
\DndMonsterAction{Life Drain}
\textsl{Melee Spell Attack:} 4 to hit, reach 5 ft., one creature. \textsl{Hit:} 10 (3d6) necrotic damage. The target must succeed on a DC 10 Constitution saving throw or its hit point maximum is reduced by an amount equal to the damage taken. This reduction lasts until the creature finishes a long rest. The target dies if this effect reduces its hit point maximum to 0.

\end{DndMonster}



\begin{DndMonster}[float=t]{Spiderdragon}
\DndMonsterType{Huge monstrosity, chaotic evil}
\DndMonsterBasics[
armor-class = {17 (natural armor)},
hit-points = {\DndDice{16d12 + 48}},
speed = {50 ft., climb 60 ft.}
]
\DndMonsterAbilityScores[
 str = 21,
 dex = 18,
 con = 16,
 int = 7,
 wis = 14,
 cha = 18
]
\DndMonsterDetails[
saving-throws = {Str +9, Dex +8},
skills = {Intimidation +8, Perception +6},
damage-resistances = {poison, psychic},
senses = {darkvision 90 ft., passive perception 16},
languages = {Abyssal, Draconic, Undercommon},
challenge = 11,
]
\DndMonsterAction{Magic Resistance}
The spiderdragon has advantage on saving throws against spells and other magical effects.

\DndMonsterAction{Spider Climb}
The spiderdragon can climb difficult surfaces, including upside down on ceilings, without needing to make an ability check.

\DndMonsterAction{Web Walker}
The spiderdragon ignores movement restrictions caused by webbing.

\DndMonsterSection{Actions}
\DndMonsterAction{Multiattack}
The spiderdragon makes one Bite attack and two Claw attacks.

\DndMonsterAction{Bite}
\textsl{Melee Weapon Attack:} 9 to hit, reach 10 ft., one target. \textsl{Hit:} 10 (1d10 + 5) piercing damage plus 13 (2d12) poison damage.

\DndMonsterAction{Claw}
\textsl{Melee Weapon Attack:} 9 to hit, reach 5 ft., one target. \textsl{Hit:} 12 (2d6 + 5) slashing damage.

\DndMonsterAction{Spiderling Breath (Recharge 5\textendash 6)}
The spiderdragon exhales venomous spiderlings in a 30-foot cone. Each creature in that area must make a DC 15 Dexterity saving throw, taking 33 (6d10) piercing damage and 33 (6d10) poison damage on a failed save or half as much damage on a successful one.

\DndMonsterSection{Bonus Actions}
\DndMonsterAction{Stifling Webs (Recharge 5\textendash 6)}
The spiderdragon spins a 30-foot cube of strong, sticky webbing in an area adjacent to itself. The webbing lasts for 1 minute, is difficult terrain, and lightly obscures its area. A creature that starts its turn in the webbing or enters the webbing for the first time on its turn must succeed on a DC 15 Dexterity saving throw or have the restrained condition while in the web. As an action, a creature can free itself or another creature from the web by succeeding on a DC 15 Strength check.

A 5-foot cube of the web is destroyed if it takes at least 10 acid, fire, or slashing damage on a single turn.

\end{DndMonster}



\begin{DndMonster}[float=t]{Spirit Naga}
\DndMonsterType{Large monstrosity, chaotic evil}
\DndMonsterBasics[
armor-class = {15 (natural armor)},
hit-points = {\DndDice{10d10 + 20}},
speed = {40 ft.}
]
\DndMonsterAbilityScores[
 str = 18,
 dex = 17,
 con = 14,
 int = 16,
 wis = 15,
 cha = 16
]
\DndMonsterDetails[
saving-throws = {Dex +6, Con +5, Wis +5, Cha +6},
damage-immunities = {poison},
condition-immunities = {charmed, poisoned},
senses = {darkvision 60 ft., passive perception 12},
languages = {Abyssal, Common},
challenge = 8,
]
\DndMonsterAction{Rejuvenation}
If it dies, the naga returns to life in 1d6 days and regains all its hit points. Only a \textit{wish} spell can prevent this trait from functioning.

\DndMonsterAction{Spellcasting}
The naga is a 10th-level spellcaster. Its spellcasting ability is Intelligence (spell save DC 14, 6 to hit with spell attacks), and it needs only verbal components to cast its spells. It has the following wizard spells prepared:
\begin{DndMonsterSpells}
  \DndMonsterSpellLevel{mage hand}
   \DndMonsterSpellLevel[1][4]{charm person}
   \DndMonsterSpellLevel[2][3]{detect thoughts}
   \DndMonsterSpellLevel[3][3]{lightning bolt}
   \DndMonsterSpellLevel[4][3]{blight}
   \DndMonsterSpellLevel[5][2]{dominate person}
\end{DndMonsterSpells}
\DndMonsterSection{Actions}
\DndMonsterAction{Bite}
\textsl{Melee Weapon Attack:} 7 to hit, reach 10 ft., one creature. \textsl{Hit:} 7 (1d6 + 4) piercing damage, and the target must make a DC 13 Constitution saving throw, taking 31 (7d8) poison damage on a failed save, or half as much damage on a successful one.

\end{DndMonster}



\begin{DndMonster}[float=t]{Star Angler}
\DndMonsterType{Large monstrosity, unaligned}
\DndMonsterBasics[
armor-class = {15 (natural armor)},
hit-points = {\DndDice{14d10 + 42}},
speed = {0 ft., fly 40 ft.}
]
\DndMonsterAbilityScores[
 str = 21,
 dex = 15,
 con = 17,
 int = 3,
 wis = 14,
 cha = 6
]
\DndMonsterDetails[
skills = {Perception +5, Stealth +8},
senses = {blindsight 120 ft. (can't see beyond this radius), passive perception 15},
challenge = 8,
]
\DndMonsterAction{Avoidance}
If the star angler is subjected to an effect that allows it to make a saving throw to take only half damage, it instead takes no damage if it succeeds on the saving throw and only half damage if it fails.

\DndMonsterAction{Illumination}
The star angler's lure sheds bright light in a 30- foot radius and dim light for an additional 30 feet.

\DndMonsterSection{Actions}
\DndMonsterAction{Multiattack}
The star angler makes three Bite attacks.

\DndMonsterAction{Bite}
\textsl{Melee Weapon Attack:} 8 to hit, reach 5 ft., one target. \textsl{Hit:} 16 (2d10 + 5) piercing damage.

\DndMonsterSection{Bonus Actions}
\DndMonsterAction{Lure Charm}
The star angler's lure flares with enchanting starlight, targeting one creature the star angler can see within 120 feet of itself. The target must succeed on a DC 13 Wisdom saving throw or have the charmed condition until the start of the star angler's next turn. While charmed in this way, the target has the incapacitated condition and must use its movement on its turn to move directly toward the star angler; a charmed target doesn't avoid opportunity attacks, but it does avoid damaging terrain. A target can be charmed by only one star angler at a time.

\end{DndMonster}



\begin{DndMonster}[float=t]{Stone Golem}
\DndMonsterType{Large construct, unaligned}
\DndMonsterBasics[
armor-class = {17 (natural armor)},
hit-points = {\DndDice{17d10 + 85}},
speed = {30 ft.}
]
\DndMonsterAbilityScores[
 str = 22,
 dex = 9,
 con = 20,
 int = 3,
 wis = 11,
 cha = 1
]
\DndMonsterDetails[
damage-immunities = {poison; psychic; bludgeoning, piercing, slashing from nonmagical attacks that aren't adamantine},
condition-immunities = {charmed, exhaustion, frightened, paralyzed, petrified, poisoned},
senses = {darkvision 120 ft., passive perception 10},
languages = {understands the languages of its creator but can't speak},
challenge = 10,
]
\DndMonsterAction{Immutable Form}
The golem is immune to any spell or effect that would alter its form.

\DndMonsterAction{Magic Resistance}
The golem has advantage on saving throws against spells and other magical effects.

\DndMonsterAction{Magic Weapons}
The golem's weapon attacks are magical.

\DndMonsterSection{Actions}
\DndMonsterAction{Multiattack}
The golem makes two slam attacks.

\DndMonsterAction{Slam}
\textsl{Melee Weapon Attack:} 10 to hit, reach 5 ft., one target. \textsl{Hit:} 19 (3d8 + 6) bludgeoning damage.

\DndMonsterAction{Slow (Recharge 5\textendash 6)}
The golem targets one or more creatures it can see within 10 feet of it. Each target must make a DC 17 Wisdom saving throw against this magic. On a failed save, a target can't use reactions, its speed is halved, and it can't make more than one attack on its turn. In addition, the target can take either an action or a bonus action on its turn, not both. These effects last for 1 minute. A target can repeat the saving throw at the end of each of its turns, ending the effect on itself on a success.

\end{DndMonster}



\begin{DndMonster}[float*=b,width=\textwidth + 8pt]{Strahd, Master of Death House}
\begin{multicols}{2}
\DndMonsterType{Medium undead (vampire, wizard), lawful evil}
\DndMonsterBasics[
armor-class = {16 (natural armor)},
hit-points = {\DndDice{16d8 + 64}},
speed = {30 ft.}
]
\DndMonsterAbilityScores[
 str = 18,
 dex = 18,
 con = 18,
 int = 20,
 wis = 15,
 cha = 18
]
\DndMonsterDetails[
saving-throws = {Dex +9, Wis +7, Cha +9},
skills = {Arcana +15, Perception +12, Religion +10, Stealth +14},
damage-resistances = {necrotic; bludgeoning, piercing, slashing from nonmagical attacks},
senses = {darkvision 120 ft., passive perception 22},
languages = {Abyssal, Common, Draconic, Elvish, Giant, Infernal},
challenge = 15,
]
\DndMonsterAction{Legendary Resistance (3/Day)}
If Strahd fails a saving throw, he can choose to succeed instead.

\DndMonsterAction{Master of the House}
When Strahd is reduced to 0 hit points, he dissolves into mist and immediately teleports to his lair in Castle Ravenloft. After 1d4 hours, Strahd re-forms in a random unoccupied space within his lair, regaining all his hit points.

\DndMonsterAction{Regeneration}
Strahd regains 20 hit points at the start of his turn if he has at least 1 hit point. If he takes radiant damage, this trait doesn't function at the start of his next turn.

\DndMonsterAction{Spider Climb}
Strahd can climb difficult surfaces, including upside down on ceilings, without needing to make an ability check.

\DndMonsterAction{Vampire Weaknesses}
Strahd has the following flaws:

\DndMonsterAction{Harmed by Running Water}
While in running water, Strahd takes 20 acid damage if he ends his turn there, and he can't use his Change Shape.

\DndMonsterAction{Sunlight Hypersensitivity}
While in sunlight, Strahd takes 20 radiant damage at the start of his turn, has disadvantage on attack rolls and ability checks, and can't use his Change Shape bonus action.

\DndMonsterAction{Spellcasting}
Strahd casts one of the following spells, using Intelligence as the spellcasting ability (spell save DC 18):
\begin{DndMonsterSpells}
  \DndInnateSpellLevel{Detect Thoughts}
  \DndInnateSpellLevel[2]{Animate Dead}
  \DndInnateSpellLevel[1]{Greater Invisibility}
\end{DndMonsterSpells}
\DndMonsterSection{Actions}
\DndMonsterAction{Multiattack}
Strahd makes two Death Strike attacks. He can replace one of these attacks with Blighted Fire if available.

\DndMonsterAction{Death Strike}
\textsl{Melee Weapon Attack:} 9 to hit, reach 5 ft., one target. \textsl{Hit:} 8 (1d8 + 4) slashing damage plus 14 (4d6) necrotic damage. If the target is a creature, Strahd can forgo dealing slashing damage; the target then has the grappled condition (escape DC 18) instead. Strahd can grapple only one creature at a time.

\DndMonsterAction{Blighted Fire (Recharge 5\textendash 6)}
Strahd summons shadowy, necrotic fire that fills a 20-foot-radius sphere centered on a point he can see within 90 feet of himself. Each creature in that area must make a DC 18 Dexterity saving throw, taking 14 (4d6) fire damage plus 14 (4d6) necrotic damage on a failed save or half as much damage on a successful one.

\DndMonsterAction{Charm}
Strahd targets one Humanoid he can see within 30 feet of himself. The target must succeed on a DC 17 Wisdom saving throw or have the charmed condition. The charmed target regards Strahd as a trusted friend to be heeded and protected. The target isn't under Strahd's control, but it takes Strahd's requests and actions in the most favorable way.

Each time Strahd or his companions deal damage to the target, it can repeat the saving throw, ending the effect on itself on a success. Otherwise, the effect lasts 24 hours or until Strahd is reduced to 0 hit points, is on a different plane of existence than the target, or uses a bonus action to end the effect.

\DndMonsterSection{Bonus Actions}
\DndMonsterAction{Bite}
\textsl{Melee Weapon Attack:} 9 to hit, reach 5 ft., one creature that has the charmed or grappled condition. \textsl{Hit:} 7 (1d6 + 4) piercing damage plus 10 (3d6) necrotic damage. The target's hit point maximum is reduced by an amount equal to the necrotic damage taken, and Strahd regains a number of hit points equal to that amount. The reduction lasts until the target finishes a long rest. The target dies if its hit point maximum is reduced to 0. A Humanoid slain in this way and then buried rises the following night as a vampire spawn under Strahd's control.

\DndMonsterAction{Change Shape}
Strahd transforms into a new form or back into his true form. Anything he is wearing transforms with him, but nothing he is carrying does. He reverts to his true form if he dies. When transforming into a new form, Strahd chooses one of the following options:

\DndMonsterAction{Beast Form}
Strahd transforms into a Tiny bat (flying speed 30 ft.) or a Medium wolf (speed 40 ft.). While in that form, he can't speak, and he retains his game statistics other than his size and speed.

\DndMonsterAction{Mist Form}
Strahd transforms into a Medium cloud of mist. While in this form, Strahd has a flying speed of 20 feet, can hover, and can enter a hostile creature's space and stop there. While in mist form, Strahd can pass through a space without squeezing as long as air can pass through that space, but he can't pass through water. Strahd has advantage on Strength, Dexterity, and Constitution saving throws, and he is immune to all nonmagical damage except the damage he takes as part of his Vampire Weaknesses trait. While in mist form, Strahd can't take any actions, speak, or manipulate objects.

\DndMonsterSection{Legendary Actions}
Strahd, Master of Death House can take 3 legendary actions, choosing from the options below. Only one legendary action can be used at a time and only at the end of another creature's turn. Strahd, Master of Death House regains spent legendary actions at the start of its turn.

\begin{DndMonsterLegendaryActions}
\DndMonsterLegendaryAction{Cunning Escape}{Strahd moves up to his speed without provoking opportunity attacks.
}
\DndMonsterLegendaryAction{Strike (Costs 2 Actions)}{Strahd makes one Death Strike attack.
}
\end{DndMonsterLegendaryActions}
\end{multicols}
\end{DndMonster}



\begin{DndMonster}[float=t]{Swarm of Poisonous Snakes}
\DndMonsterType{Medium beast, unaligned}
\DndMonsterBasics[
armor-class = {14},
hit-points = {\DndDice{8d8}},
speed = {30 ft., swim 30 ft.}
]
\DndMonsterAbilityScores[
 str = 8,
 dex = 18,
 con = 11,
 int = 1,
 wis = 10,
 cha = 3
]
\DndMonsterDetails[
damage-resistances = {bludgeoning, piercing, slashing},
condition-immunities = {charmed, frightened, grappled, paralyzed, petrified, prone, restrained, stunned},
senses = {blindsight 10 ft., passive perception 10},
challenge = 2,
]
\DndMonsterAction{Swarm}
The swarm can occupy another creature's space and vice versa, and the swarm can move through any opening large enough for a Tiny snake. The swarm can't regain hit points or gain temporary hit points.

\DndMonsterSection{Actions}
\DndMonsterAction{Bites}
\textsl{Melee Weapon Attack:} 6 to hit, reach 0 ft., one creature in the swarm's space. \textsl{Hit:} 7 (2d6) piercing damage, or 3 (1d6) piercing damage if the swarm has half of its hit points or fewer. The target must make a DC 10 Constitution saving throw, taking 14 (4d6) poison damage on a failed save, or half as much damage on a successful one.

\end{DndMonster}



\begin{DndMonster}[float=t]{Tasha the Witch}
\DndMonsterType{Medium humanoid (human, wizard), chaotic neutral}
\DndMonsterBasics[
armor-class = {19 (\textit{robe of the archmagi})},
hit-points = {\DndDice{28d8 + 84}},
speed = {30 ft.}
]
\DndMonsterAbilityScores[
 str = 10,
 dex = 18,
 con = 17,
 int = 23,
 wis = 12,
 cha = 22
]
\DndMonsterDetails[
saving-throws = {Int +12, Wis +7, Cha +12},
skills = {Arcana +18, History +12, Persuasion +12},
condition-immunities = {charmed, frightened},
languages = {Abyssal, Celestial, Common, Draconic, Elvish, Infernal, Sylvan},
challenge = 19,
]
\DndMonsterAction{Legendary Resistance (3/Day)}
If Tasha fails a saving throw, she can choose to succeed instead.

\DndMonsterAction{Magic Resistance}
Tasha has advantage on saving throws against spells and other magical effects. (This trait is bestowed by her Robe of the Archmagi.)

\DndMonsterAction{Special Equipment}
Tasha wears a Robe of the Archmagi.

\DndMonsterAction{Spellcasting}
Tasha casts one of the following spells, using Intelligence as the spellcasting ability (spell save DC 22, 14 to hit with spell attacks):
\begin{DndMonsterSpells}
  \DndInnateSpellLevel{Detect Magic}
  \DndInnateSpellLevel[]{Polymorph}
  \DndInnateSpellLevel[1]{Maze}
\end{DndMonsterSpells}
\DndMonsterSection{Actions}
\DndMonsterAction{Multiattack}
Tasha makes two Caustic Blast attacks and uses Psychic Whip once.

\DndMonsterAction{Caustic Blast}
\textsl{Melee or Ranged Spell Attack:} 14 to hit, reach 5 ft. or range 120 ft., one target. \textsl{Hit:} 21 (6d4 + 6) acid damage.

\DndMonsterAction{Psychic Whip}
Tasha psychically lashes out at one creature she can see within 90 feet of herself. The target must make a DC 20 Intelligence saving throw. On a failed save, the target takes 21 (6d6) psychic damage and has the stunned condition until the start of Tasha's next turn. On a successful save, the target takes half as much damage only.

\DndMonsterSection{Bonus Actions}
\DndMonsterAction{Abyssal Visage (2/Day)}
For 1 minute, Tasha gains a flying speed of 30 feet, is immune to poison damage and the poisoned condition, and has advantage on attack rolls against any creature that doesn't have all its hit points. These benefits end early if Tasha has the incapacitated condition or if she uses another bonus action to dismiss them.

\DndMonsterSection{Reactions}
\DndMonsterAction{Arcane Rebuff}
Immediately after Tasha takes damage, she unleashes arcane energy in a 10-foot-radius sphere centered on herself. All other creatures in that area must make a DC 20 Dexterity saving throw, taking 19 (3d12) lightning damage on a failed save or half as much damage on a successful one. Tasha then teleports, along with any equipment she is wearing or carrying, to an unoccupied space she can see within 60 feet of herself.

\end{DndMonster}


\clearpage

\begin{DndMonster}[float=t]{Treant}
\DndMonsterType{Huge plant, chaotic good}
\DndMonsterBasics[
armor-class = {16 (natural armor)},
hit-points = {\DndDice{12d12 + 60}},
speed = {30 ft.}
]
\DndMonsterAbilityScores[
 str = 23,
 dex = 8,
 con = 21,
 int = 12,
 wis = 16,
 cha = 12
]
\DndMonsterDetails[
damage-vulnerabilities = {fire},
damage-resistances = {bludgeoning, piercing},
languages = {Common, Druidic, Elvish, Sylvan},
challenge = 9,
]
\DndMonsterAction{False Appearance}
While the treant remains motionless, it is indistinguishable from a normal tree.

\DndMonsterAction{Siege Monster}
The treant deals double damage to objects and structures.

\DndMonsterSection{Actions}
\DndMonsterAction{Multiattack}
The treant makes two slam attacks.

\DndMonsterAction{Slam}
\textsl{Melee Weapon Attack:} 10 to hit, reach 5 ft., one target. \textsl{Hit:} 16 (3d6 + 6) bludgeoning damage.

\DndMonsterAction{Rock}
\textsl{Ranged Weapon Attack:} 10 to hit, range 60/180 ft., one target. \textsl{Hit:} 28 (4d10 + 6) bludgeoning damage.

\DndMonsterAction{Animate Trees (1/Day)}
The treant magically animates one or two trees it can see within 60 feet of it. These trees have the same statistics as a \textbf{treant}, except they have Intelligence and Charisma scores of 1, they can't speak, and they have only the Slam action option. An animated tree acts as an ally of the treant. The tree remains animate for 1 day or until it dies; until the treant dies or is more than 120 feet from the tree; or until the treant takes a bonus action to turn it back into an inanimate tree. The tree then takes root if possible.

\end{DndMonster}



\begin{DndMonster}[float*=b,width=\textwidth + 8pt]{Unicorn}
\begin{multicols}{2}
\DndMonsterType{Large celestial, lawful good}
\DndMonsterBasics[
armor-class = {12},
hit-points = {\DndDice{9d10 + 18}},
speed = {50 ft.}
]
\DndMonsterAbilityScores[
 str = 18,
 dex = 14,
 con = 15,
 int = 11,
 wis = 17,
 cha = 16
]
\DndMonsterDetails[
damage-immunities = {poison},
condition-immunities = {charmed, paralyzed, poisoned},
senses = {darkvision 60 ft., passive perception 13},
languages = {Celestial, Elvish, Sylvan, telepathy 60 ft.},
challenge = 5,
]
\DndMonsterAction{Charge}
If the unicorn moves at least 20 feet straight toward a target and then hits it with a horn attack on the same turn, the target takes an extra 9 (2d8) piercing damage. If the target is a creature, it must succeed on a DC 15 Strength saving throw or be knocked prone.

\DndMonsterAction{Magic Resistance}
The unicorn has advantage on saving throws against spells and other magical effects.

\DndMonsterAction{Magic Weapons}
The unicorn's weapon attacks are magical.

\DndMonsterAction{Innate Spellcasting}
The unicorn's innate spellcasting ability is Charisma (spell save DC 14). The unicorn can innately cast the following spells, requiring no components:
\begin{DndMonsterSpells}
  \DndInnateSpellLevel{detect evil and good}
  \DndInnateSpellLevel[1]{calm emotions}
\end{DndMonsterSpells}
\DndMonsterSection{Actions}
\DndMonsterAction{Multiattack}
The unicorn makes two attacks: one with its hooves and one with its horn.

\DndMonsterAction{Hooves}
\textsl{Melee Weapon Attack:} 7 to hit, reach 5 ft., one target. \textsl{Hit:} 11 (2d6 + 4) bludgeoning damage.

\DndMonsterAction{Horn}
\textsl{Melee Weapon Attack:} 7 to hit, reach 5 ft., one target. \textsl{Hit:} 8 (1d8 + 4) piercing damage.

\DndMonsterAction{Healing Touch (3/Day)}
The unicorn touches another creature with its horn. The target magically regains 11 (2d8 + 2) hit points. In addition, the touch removes all diseases and neutralizes all poisons afflicting the target.

\DndMonsterAction{Teleport (1/Day)}
The unicorn magically teleports itself and up to three willing creatures it can see within 5 feet of it, along with any equipment they are wearing or carrying, to a location the unicorn is familiar with, up to 1 mile away.

\DndMonsterSection{Legendary Actions}
Unicorn can take 3 legendary actions, choosing from the options below. Only one legendary action can be used at a time and only at the end of another creature's turn. Unicorn regains spent legendary actions at the start of its turn.

\begin{DndMonsterLegendaryActions}
\DndMonsterLegendaryAction{Hooves}{The unicorn makes one attack with its hooves.
}
\DndMonsterLegendaryAction{Shimmering Shield (Costs 2 Actions)}{The unicorn creates a shimmering, magical field around itself or another creature it can see within 60 feet of it. The target gains a +2 bonus to AC until the end of the unicorn's next turn.
}
\DndMonsterLegendaryAction{Heal Self (Costs 3 Actions)}{The unicorn magically regains 11 (2d8 + 2) hit points.
}
\end{DndMonsterLegendaryActions}
\DndMonsterSection{Regional Effects}
Transformed by the creature's celestial presence, the domain of a unicorn might include any of the following magical effects:


\begin{itemize}
\item Open flames of a non magical nature are extinguished within the unicorn's domain. Torches and campfires refuse to burn, but closed lanterns are unaffected.
\item Creatures native to the unicorn's domain have an easier time hiding; they have advantage on all Dexterity (Stealth) checks made to hide.
\item When a good-aligned creature casts a spell or uses a magical effect that causes another good-aligned creature to regain hit points, the target regains the maximum number of hit points possible for the spell or effect.
\item Curses affecting any good-aligned creature are suppressed.
\end{itemize}

If the unicorn dies, these effects end immediately.

\end{multicols}
\end{DndMonster}



\begin{DndMonster}[float*=b,width=\textwidth + 8pt]{Vampire}
\begin{multicols}{2}
\DndMonsterType{Medium undead (shapechanger), lawful evil}
\DndMonsterBasics[
armor-class = {16 (natural armor)},
hit-points = {\DndDice{17d8 + 68}},
speed = {30 ft.}
]
\DndMonsterAbilityScores[
 str = 18,
 dex = 18,
 con = 18,
 int = 17,
 wis = 15,
 cha = 18
]
\DndMonsterDetails[
saving-throws = {Dex +9, Wis +7, Cha +9},
skills = {Perception +7, Stealth +9},
damage-resistances = {necrotic; bludgeoning, piercing, slashing from nonmagical attacks},
senses = {darkvision 120 ft., passive perception 17},
languages = {the languages it knew in life},
challenge = 13,
]
\DndMonsterAction{Shapechanger}
If the vampire isn't in sunlight or running water, it can use its action to polymorph into a Tiny bat or a Medium cloud of mist, or back into its true form.

While in bat form, the vampire can't speak, its walking speed is 5 feet, and it has a flying speed of 30 feet. Its statistics, other than its size and speed, are unchanged. Anything it is wearing transforms with it, but nothing it is carrying does. It reverts to its true form if it dies.

While in mist form, the vampire can't take any actions, speak, or manipulate objects. It is weightless, has a flying speed of 20 feet, can hover, and can enter a hostile creature's space and stop there. In addition, if air can pass through a space, the mist can do so without squeezing, and it can't pass through water. It has advantage on Strength, Dexterity, and Constitution saving throws, and it is immune to all nonmagical damage, except the damage it takes from sunlight.

\DndMonsterAction{Legendary Resistance (3/Day)}
If the vampire fails a saving throw, it can choose to succeed instead.

\DndMonsterAction{Misty Escape}
When it drops to 0 hit points outside its resting place, the vampire transforms into a cloud of mist (as in the Shapechanger trait) instead of falling unconscious, provided that it isn't in sunlight or running water. If it can't transform, it is destroyed.

While it has 0 hit points in mist form, it can't revert to its vampire form, and it must reach its resting place within 2 hours or be destroyed. Once in its resting place, it reverts to its vampire form. It is then paralyzed until it regains at least 1 hit point. After spending 1 hour in its resting place with 0 hit points, it regains 1 hit point.

\DndMonsterAction{Regeneration}
The vampire regains 20 hit points at the start of its turn if it has at least 1 hit point and isn't in sunlight or running water. If the vampire takes radiant damage or damage from holy water, this trait doesn't function at the start of the vampire's next turn.

\DndMonsterAction{Spider Climb}
The vampire can climb difficult surfaces, including upside down on ceilings, without needing to make an ability check.

\DndMonsterAction{Vampire Weaknesses}
The vampire has the following flaws:

\textit{Forbiddance.} The vampire can't enter a residence without an invitation from one of the occupants.

\textit{Harmed by Running Water.} The vampire takes 20 acid damage if it ends its turn in running water.

\textit{Stake to the Heart.} If a piercing weapon made of wood is driven into the vampire's heart while the vampire is incapacitated in its resting place, the vampire is paralyzed until the stake is removed.

\textit{Sunlight Hypersensitivity.} The vampire takes 20 radiant damage when it starts its turn in sunlight. While in sunlight, it has disadvantage on attack rolls and ability checks.

\DndMonsterSection{Actions}
\DndMonsterAction{Multiattack (Vampire Form Only)}
The vampire makes two attacks, only one of which can be a bite attack.

\DndMonsterAction{Unarmed Strike (Vampire Form Only)}
\textsl{Melee Weapon Attack:} 9 to hit, reach 5 ft., one creature. \textsl{Hit:} 8 (1d8 + 4) bludgeoning damage. Instead of dealing damage, the vampire can grapple the target (escape DC 18).

\DndMonsterAction{Bite (Bat or Vampire Form Only)}
\textsl{Melee Weapon Attack:} 9 to hit, reach 5 ft., one willing creature, or a creature that is grappled by the vampire, incapacitated, or restrained. \textsl{Hit:} 7 (1d6 + 4) piercing damage plus 10 (3d6) necrotic damage. The target's hit point maximum is reduced by an amount equal to the necrotic damage taken, and the vampire regains hit points equal to that amount. The reduction lasts until the target finishes a long rest. The target dies if this effect reduces its hit point maximum to 0. A humanoid slain in this way and then buried in the ground rises the following night as a \textbf{vampire spawn} under the vampire's control.

\DndMonsterAction{Charm}
The vampire targets one humanoid it can see within 30 feet of it. If the target can see the vampire, the target must succeed on a DC 17 Wisdom saving throw against this magic or be charmed by the vampire. The charmed target regards the vampire as a trusted friend to be heeded and protected. Although the target isn't under the vampire's control, it takes the vampire's requests or actions in the most favorable way it can, and it is a willing target for the vampire's bite attack.

Each time the vampire or the vampire's companions do anything harmful to the target, it can repeat the saving throw, ending the effect on itself on a success. Otherwise, the effect lasts 24 hours or until the vampire is destroyed, is on a different plane of existence than the target, or takes a bonus action to end the effect.

\DndMonsterAction{Children of the Night (1/Day)}
The vampire magically calls 2d4 swarms of bats or rats, provided that the sun isn't up. While outdoors, the vampire can call 3d6 wolves instead. The called creatures arrive in 1d4 rounds, acting as allies of the vampire and obeying its spoken commands. The beasts remain for 1 hour, until the vampire dies, or until the vampire dismisses them as a bonus action.

\DndMonsterSection{Legendary Actions}
Vampire can take 3 legendary actions, choosing from the options below. Only one legendary action can be used at a time and only at the end of another creature's turn. Vampire regains spent legendary actions at the start of its turn.

\begin{DndMonsterLegendaryActions}
\DndMonsterLegendaryAction{Move}{The vampire moves up to its speed without provoking opportunity attacks.
}
\DndMonsterLegendaryAction{Unarmed Strike}{The vampire makes one unarmed strike.
}
\DndMonsterLegendaryAction{Bite (Costs 2 Actions)}{The vampire makes one bite attack.
}
\end{DndMonsterLegendaryActions}
\DndMonsterSection{Regional Effects}
The region surrounding a vampire's lair is warped by the creature's unnatural presence, creating any of the following effects:


\begin{itemize}
\item There's a noticeable increase in the populations of bats, rats, and wolves in the region.
\item Plants within 500 feet of the lair wither, and their stems and branches become twisted and thorny.
\item Shadows cast within 500 feet of the lair seem abnormally gaunt and sometimes move as though alive.
\item A creeping fog clings to the ground within 500 feet of the vampire's lair. The fog occasionally takes eerie forms, such as grasping claws and writhing serpents.
\end{itemize}

If the vampire is destroyed, these effects end after 2d6 days.

\end{multicols}
\end{DndMonster}



\begin{DndMonster}[float=t]{Vampire Spawn}
\DndMonsterType{Medium undead, neutral evil}
\DndMonsterBasics[
armor-class = {15 (natural armor)},
hit-points = {\DndDice{11d8 + 33}},
speed = {30 ft.}
]
\DndMonsterAbilityScores[
 str = 16,
 dex = 16,
 con = 16,
 int = 11,
 wis = 10,
 cha = 12
]
\DndMonsterDetails[
saving-throws = {Dex +6, Wis +3},
skills = {Perception +3, Stealth +6},
damage-resistances = {necrotic; bludgeoning, piercing, slashing from nonmagical attacks},
senses = {darkvision 60 ft., passive perception 13},
languages = {the languages it knew in life},
challenge = 5,
]
\DndMonsterAction{Regeneration}
The vampire regains 10 hit points at the start of its turn if it has at least 1 hit point and isn't in sunlight or running water. If the vampire takes radiant damage or damage from holy water, this trait doesn't function at the start of the vampire's next turn.

\DndMonsterAction{Spider Climb}
The vampire can climb difficult surfaces, including upside down on ceilings, without needing to make an ability check.

\DndMonsterAction{Vampire Weaknesses}
The vampire has the following flaws:

\textit{Forbiddance.} The vampire can't enter a residence without an invitation from one of the occupants.

\textit{Harmed by Running Water.} The vampire takes 20 acid damage when it ends its turn in running water.

\textit{Stake to the Heart.} The vampire is destroyed if a piercing weapon made of wood is driven into its heart while it is incapacitated in its resting place.

\textit{Sunlight Hypersensitivity.} The vampire takes 20 radiant damage when it starts its turn in sunlight. While in sunlight, it has disadvantage on attack rolls and ability checks.

\DndMonsterSection{Actions}
\DndMonsterAction{Multiattack}
The vampire makes two attacks, only one of which can be a bite attack.

\DndMonsterAction{Bite}
\textsl{Melee Weapon Attack:} 6 to hit, reach 5 ft., one willing creature, or a creature that is grappled by the vampire, incapacitated, or restrained. \textsl{Hit:} 6 (1d6 + 3) piercing damage plus 7 (2d6) necrotic damage. The target's hit point maximum is reduced by an amount equal to the necrotic damage taken, and the vampire regains hit points equal to that amount. The reduction lasts until the target finishes a long rest. The target dies if this effect reduces its hit point maximum to 0.

\DndMonsterAction{Claws}
\textsl{Melee Weapon Attack:} 6 to hit, reach 5 ft., one creature. \textsl{Hit:} 8 (2d4 + 3) slashing damage. Instead of dealing damage, the vampire can grapple the target (escape DC 13).

\end{DndMonster}



\begin{DndMonster}[float=t]{Vecna the Archlich}
\DndMonsterType{Medium undead (wizard), lawful evil}
\DndMonsterBasics[
armor-class = {18 (natural armor)},
hit-points = {\DndDice{32d8 + 128}},
speed = {30 ft.}
]
\DndMonsterAbilityScores[
 str = 14,
 dex = 16,
 con = 18,
 int = 22,
 wis = 24,
 cha = 16
]
\DndMonsterDetails[
saving-throws = {Con +12, Int +14, Wis +15},
skills = {Arcana +22, History +14, Insight +15, Perception +15},
damage-resistances = {cold, lightning, necrotic},
damage-immunities = {poison; bludgeoning, piercing, slashing from nonmagical attacks},
condition-immunities = {charmed, exhaustion, frightened, paralyzed, poisoned, stunned},
senses = {truesight 120 ft., passive perception 25},
languages = {Common, Draconic, Elvish, Infernal},
challenge = 26,
]
\DndMonsterAction{Legendary Resistance (5/Day)}
If Vecna fails a saving throw, he can choose to succeed instead.

\DndMonsterAction{Special Equipment}
Vecna carries a magic dagger named Afterthought. In the hands of anyone other than Vecna, Afterthought is a +2 Dagger.

\DndMonsterAction{Undying}
If Vecna is slain, his soul refuses to accept its fate and lives on as a disembodied spirit that fashions a new body for itself after 1d100 years. Vecna's new body appears within 100 miles of where he was slain. When the new body is complete, Vecna regains all his hit points and becomes active again.

\DndMonsterAction{Spellcasting}
Vecna casts one of the following spells, requiring no material components and using Intelligence as the spellcasting ability (spell save DC 22):
\begin{DndMonsterSpells}
  \DndInnateSpellLevel{Animate Dead}
  \DndInnateSpellLevel[2]{Dimension Door}
  \DndInnateSpellLevel[1]{Dominate Monster}
\end{DndMonsterSpells}
\DndMonsterSection{Actions}
\DndMonsterAction{Multiattack}
Vecna uses Flight of the Damned (if available), Rotten Fate, or Spellcasting. He then makes two attacks with Afterthought.

\DndMonsterAction{Afterthought}
\textsl{Melee Weapon Attack:} 13 to hit, reach 5 ft., one target. \textsl{Hit:} 7 (1d4 + 5) piercing damage plus 9 (2d8) necrotic damage. If the target is a creature, it is afflicted by entropic magic, taking 9 (2d8) necrotic damage at the start of each of its turns. Immediately after taking this damage on its turn, the target must make a DC 20 Constitution saving throw, ending the effect on itself on a success. Until it succeeds on this save, the afflicted target can't regain hit points.

\DndMonsterAction{Flight of the Damned (Recharge 5\textendash 6)}
Vecna conjures a torrent of flying, spectral entities that fill a 120-foot cone and pass through all creatures in that area before dissipating. Each creature in that area must make a DC 22 Constitution saving throw. On a failed save, the creature takes 36 (8d8) necrotic damage and has the frightened condition for 1 minute. On a successful save, the creature takes half as much damage only. A frightened creature can repeat the saving throw at the end of each of its turns, ending the effect on itself on a success.

\DndMonsterAction{Rotten Fate}
Vecna causes necrotic magic to engulf one creature he can see within 120 feet of himself. The target must make a DC 22 Constitution saving throw, taking 96 (8d8 + 60) necrotic damage on a failed save or half as much damage on a successful one. A Humanoid killed by this magic rises as a zombie at the start of Vecna's next turn and acts immediately after Vecna in the initiative order. The zombie is under Vecna's control.

\DndMonsterSection{Bonus Actions}
\DndMonsterAction{Vile Teleport}
Vecna teleports, along with any equipment he is wearing or carrying, up to 30 feet to an unoccupied space he can see. He can cause each creature of his choice within 15 feet of his destination space to take 10 (3d6) psychic damage. If at least one creature takes this damage, Vecna regains 80 hit points.

\DndMonsterSection{Reactions}
\DndMonsterAction{Dread Counterspell}
Vecna utters a dread word to interrupt a creature he can see that is casting a spell. If the spell is 4th level or lower, it fails and has no effect. If the spell is 5th level or higher, Vecna makes an Intelligence check (DC 10 plus the spell's level). On a successful check, the spell fails and has no effect. Whatever the spell's level, the caster takes 10 (3d6) psychic damage if the spell fails.

\DndMonsterAction{Fell Rebuke}
In response to being hit by an attack, Vecna utters a fell word, dealing 10 (3d6) necrotic damage to the attacker, and Vecna teleports, along with any equipment he is wearing or carrying, up to 30 feet to an unoccupied space he can see.

\end{DndMonster}



\begin{DndMonster}[float=t]{Veteran}
\DndMonsterType{Medium humanoid (any race), any alignment}
\DndMonsterBasics[
armor-class = {17 (\textit{splint armor})},
hit-points = {\DndDice{9d8 + 18}},
speed = {30 ft.}
]
\DndMonsterAbilityScores[
 str = 16,
 dex = 13,
 con = 14,
 int = 10,
 wis = 11,
 cha = 10
]
\DndMonsterDetails[
skills = {Athletics +5, Perception +2},
languages = {any one language (usually Common)},
challenge = 3,
]
\DndMonsterSection{Actions}
\DndMonsterAction{Multiattack}
The veteran makes two longsword attacks. If it has a shortsword drawn, it can also make a shortsword attack.

\DndMonsterAction{Longsword}
\textsl{Melee Weapon Attack:} 5 to hit, reach 5 ft., one target. \textsl{Hit:} 7 (1d8 + 3) slashing damage, or 8 (1d10 + 3) slashing damage if used with two hands.

\DndMonsterAction{Shortsword}
\textsl{Melee Weapon Attack:} 5 to hit, reach 5 ft., one target. \textsl{Hit:} 6 (1d6 + 3) piercing damage.

\DndMonsterAction{Heavy Crossbow}
\textsl{Ranged Weapon Attack:} 3 to hit, range 100/400 ft., one target. \textsl{Hit:} 6 (1d10 + 1) piercing damage.

\end{DndMonster}



\begin{DndMonster}[float=t]{Vlazok}
\DndMonsterType{Large fiend (demon), chaotic evil}
\DndMonsterBasics[
armor-class = {17 (natural armor)},
hit-points = {\DndDice{16d10 + 48}},
speed = {30 ft., climb 30 ft.}
]
\DndMonsterAbilityScores[
 str = 21,
 dex = 18,
 con = 16,
 int = 6,
 wis = 9,
 cha = 9
]
\DndMonsterDetails[
saving-throws = {Str +9, Con +7},
skills = {Perception +3},
damage-resistances = {cold, fire, lightning},
damage-immunities = {poison},
condition-immunities = {charmed, exhaustion, frightened, poisoned},
senses = {darkvision 120 ft., passive perception 13},
languages = {Abyssal, telepathy 120 ft.},
challenge = 11,
]
\DndMonsterAction{All-Around Vision}
The vlazok can't be surprised.

\DndMonsterAction{Blood Frenzy}
The vlazok has advantage on melee attack rolls against any creature that doesn't have all its hit points.

\DndMonsterAction{Magic Resistance}
The vlazok has advantage on saving throws against spells and other magical effects.

\DndMonsterAction{Spider Climb}
The vlazok can climb difficult surfaces, including upside down on ceilings, without needing to make an ability check.

\DndMonsterSection{Actions}
\DndMonsterAction{Multiattack}
The vlazok makes two Gore attacks.

\DndMonsterAction{Gore}
\textsl{Melee Weapon Attack:} 9 to hit, reach 5 ft., one target. \textsl{Hit:} 23 (4d8 + 5) piercing damage, and if the target is a Large or smaller creature, it has the prone condition.

\DndMonsterSection{Bonus Actions}
\DndMonsterAction{Stomp}
\textsl{Melee Weapon Attack:} 9 to hit, reach 5 ft., one prone creature. \textsl{Hit:} 27 (4d10 + 5) bludgeoning damage.

\end{DndMonster}



\begin{DndMonster}[float=t]{Vrock}
\DndMonsterType{Large fiend (demon), chaotic evil}
\DndMonsterBasics[
armor-class = {15 (natural armor)},
hit-points = {\DndDice{11d10 + 44}},
speed = {40 ft., fly 60 ft.}
]
\DndMonsterAbilityScores[
 str = 17,
 dex = 15,
 con = 18,
 int = 8,
 wis = 13,
 cha = 8
]
\DndMonsterDetails[
saving-throws = {Dex +5, Wis +4, Cha +2},
damage-resistances = {cold; fire; lightning; bludgeoning, piercing, slashing from nonmagical attacks},
damage-immunities = {poison},
condition-immunities = {poisoned},
senses = {darkvision 120 ft., passive perception 11},
languages = {Abyssal, telepathy 120 ft.},
challenge = 6,
]
\DndMonsterAction{Magic Resistance}
The vrock has advantage on saving throws against spells and other magical effects.

\DndMonsterSection{Actions}
\DndMonsterAction{Multiattack}
The vrock makes two attacks: one with its beak and one with its talons.

\DndMonsterAction{Beak}
\textsl{Melee Weapon Attack:} 6 to hit, reach 5 ft., one target. \textsl{Hit:} 10 (2d6 + 3) piercing damage.

\DndMonsterAction{Talons}
\textsl{Melee Weapon Attack:} 6 to hit, reach 5 ft., one target. \textsl{Hit:} 14 (2d10 + 3) slashing damage.

\DndMonsterAction{Spores (Recharge 6)}
A 15-foot-radius cloud of toxic spores extends out from the vrock. The spores spread around corners. Each creature in that area must succeed on a DC 14 Constitution saving throw or become poisoned. While poisoned in this way, a target takes 5 (1d10) poison damage at the start of each of its turns. A target can repeat the saving throw at the end of each of its turns, ending the effect on itself on a success. Emptying a vial of holy water on the target also ends the effect on it.

\DndMonsterAction{Stunning Screech (1/Day)}
The vrock emits a horrific screech. Each creature within 20 feet of it that can hear it and that isn't a demon must succeed on a DC 14 Constitution saving throw or be stunned until the end of the vrock's next turn.

\end{DndMonster}



\begin{DndMonster}[float=t]{Warforged Warrior}
\DndMonsterType{Medium construct, any alignment}
\DndMonsterBasics[
armor-class = {16 (natural armor)},
hit-points = {\DndDice{4d8 + 12}},
speed = {30 ft.}
]
\DndMonsterAbilityScores[
 str = 16,
 dex = 12,
 con = 16,
 int = 10,
 wis = 14,
 cha = 11
]
\DndMonsterDetails[
skills = {Athletics +5, Perception +4, Survival +4},
damage-resistances = {poison},
condition-immunities = {poisoned},
languages = {Common},
challenge = 1,
]
\DndMonsterSection{Actions}
\DndMonsterAction{Multiattack}
The warforged makes two Armblade attacks.

\DndMonsterAction{Armblade}
\textsl{Melee Weapon Attack:} 5 to hit, reach 5 ft., one target. \textsl{Hit:} 6 (1d6 + 3) slashing damage.

\DndMonsterAction{Javelin}
\textsl{Melee or Ranged Weapon Attack:} 5 to hit, reach 5 ft. or range 30/120 ft., one target. \textsl{Hit:} 6 (1d6 + 3) piercing damage.

\DndMonsterSection{Reactions}
\DndMonsterAction{Protection}
When an attacker the warforged can see makes an attack roll against a creature within 5 feet of the warforged, the warforged can impose disadvantage on the attack roll.

\end{DndMonster}



\begin{DndMonster}[float=t]{Water Elemental}
\DndMonsterType{Large elemental, neutral}
\DndMonsterBasics[
armor-class = {14 (natural armor)},
hit-points = {\DndDice{12d10 + 48}},
speed = {30 ft., swim 90 ft.}
]
\DndMonsterAbilityScores[
 str = 18,
 dex = 14,
 con = 18,
 int = 5,
 wis = 10,
 cha = 8
]
\DndMonsterDetails[
damage-resistances = {acid; bludgeoning, piercing, slashing from nonmagical attacks},
damage-immunities = {poison},
condition-immunities = {exhaustion, grappled, paralyzed, petrified, poisoned, prone, restrained, unconscious},
senses = {darkvision 60 ft., passive perception 10},
languages = {Aquan},
challenge = 5,
]
\DndMonsterAction{Water Form}
The elemental can enter a hostile creature's space and stop there. It can move through a space as narrow as 1 inch wide without squeezing.

\DndMonsterAction{Freeze}
If the elemental takes cold damage, it partially freezes; its speed is reduced by 20 feet until the end of its next turn.

\DndMonsterSection{Actions}
\DndMonsterAction{Multiattack}
The elemental makes two slam attacks.

\DndMonsterAction{Slam}
\textsl{Melee Weapon Attack:} 7 to hit, reach 5 ft., one target. \textsl{Hit:} 13 (2d8 + 4) bludgeoning damage.

\DndMonsterAction{Whelm (Recharge 4\textendash 6)}
Each creature in the elemental's space must make a DC 15 Strength saving throw. On a failure, a target takes 13 (2d8 + 4) bludgeoning damage. If it is Large or smaller, it is also grappled (escape DC 14). Until this grapple ends, the target is restrained and unable to breathe unless it can breathe water. If the saving throw is successful, the target is pushed out of the elemental's space.

The elemental can grapple one Large creature or up to two Medium or smaller creatures at one time. At the start of each of the elemental's turns, each target grappled by it takes 13 (2d8 + 4) bludgeoning damage. A creature within 5 feet of the elemental can pull a creature or object out of it by taking an action to make a DC 14 Strength check and succeeding.

\end{DndMonster}



\begin{DndMonster}[float=t]{Water Weird}
\DndMonsterType{Large elemental, neutral}
\DndMonsterBasics[
armor-class = {13},
hit-points = {\DndDice{9d10 + 9}},
speed = {0 ft., swim 60 ft.}
]
\DndMonsterAbilityScores[
 str = 17,
 dex = 16,
 con = 13,
 int = 11,
 wis = 10,
 cha = 10
]
\DndMonsterDetails[
damage-resistances = {fire; bludgeoning, piercing, slashing from nonmagical attacks},
damage-immunities = {poison},
condition-immunities = {exhaustion, grappled, paralyzed, poisoned, restrained, prone, unconscious},
senses = {blindsight 30 ft., passive perception 10},
languages = {understands Aquan but doesn't speak},
challenge = 3,
]
\DndMonsterAction{Invisible in Water}
The water weird is invisible while fully immersed in water.

\DndMonsterAction{Water Bound}
The water weird dies if it leaves the water to which it is bound or if that water is destroyed.

\DndMonsterSection{Actions}
\DndMonsterAction{Constrict}
\textsl{Melee Weapon Attack:} 5 to hit, reach 10 ft., one creature. \textsl{Hit:} 13 (3d6 + 3) bludgeoning damage. If the target is Medium or smaller, it is grappled (escape DC 13) and pulled 5 feet toward the water weird. Until this grapple ends, the target is restrained, the water weird tries to drown it, and the water weird can't constrict another target.

\end{DndMonster}



\begin{DndMonster}[float=t]{Werewolf}
\DndMonsterType{Medium humanoid (human, shapechanger), chaotic evil}
\DndMonsterBasics[
armor-class = {11 in humanoid form (12 (natural armor) in wolf or hybrid form)},
hit-points = {\DndDice{9d8 + 18}},
speed = {30 ft. (40 ft. in wolf form)}
]
\DndMonsterAbilityScores[
 str = 15,
 dex = 13,
 con = 14,
 int = 10,
 wis = 11,
 cha = 10
]
\DndMonsterDetails[
skills = {Perception +4, Stealth +3},
damage-immunities = {bludgeoning, piercing, slashing from nonmagical attacks that aren't silvered},
languages = {Common (can't speak in wolf form)},
challenge = 3,
]
\DndMonsterAction{Shapechanger}
The werewolf can use its action to polymorph into a wolf-humanoid hybrid or into a wolf, or back into its true form, which is humanoid. Its statistics, other than its AC, are the same in each form. Any equipment it is wearing or carrying isn't transformed. It reverts to its true form if it dies.

\DndMonsterAction{Keen Hearing and Smell}
The werewolf has advantage on Wisdom (Perception) checks that rely on hearing or smell.

\DndMonsterSection{Actions}
\DndMonsterAction{Multiattack (Humanoid or Hybrid Form Only)}
The werewolf makes two attacks: two with its spear (humanoid form) or one with its bite and one with its claws (hybrid form).

\DndMonsterAction{Bite (Wolf or Hybrid Form Only)}
\textsl{Melee Weapon Attack:} 4 to hit, reach 5 ft., one target. \textsl{Hit:} 6 (1d8 + 2) piercing damage. If the target is a humanoid, it must succeed on a DC 12 Constitution saving throw or be cursed with werewolf lycanthropy.

\DndMonsterAction{Claws (Hybrid Form Only)}
\textsl{Melee Weapon Attack:} 4 to hit, reach 5 ft., one creature. \textsl{Hit:} 7 (2d4 + 2) slashing damage.

\DndMonsterAction{Spear (Humanoid Form Only)}
\textsl{Melee or Ranged Weapon Attack:} 4 to hit, reach 5 ft. or range 20/60 ft., one creature. \textsl{Hit:} 5 (1d6 + 2) piercing damage, or 6 (1d8 + 2) piercing damage if used with two hands to make a melee attack.

\end{DndMonster}



\begin{DndMonster}[float=t]{Whirling Chandelier}
\DndMonsterType{Large construct, unaligned}
\DndMonsterBasics[
armor-class = {13 (natural armor)},
hit-points = {\DndDice{14d10 + 28}},
speed = {30 ft., fly 30 ft. \textit{(hover)}}
]
\DndMonsterAbilityScores[
 str = 18,
 dex = 15,
 con = 15,
 int = 3,
 wis = 5,
 cha = 1
]
\DndMonsterDetails[
damage-resistances = {fire},
damage-immunities = {poison, psychic},
condition-immunities = {blinded, charmed, deafened, exhaustion, frightened, paralyzed, petrified, poisoned},
senses = {blindsight 60 ft. (blind beyond this radius), passive perception 7},
languages = {understands Common but can't speak},
challenge = 7,
]
\DndMonsterAction{False Appearance}
If the chandelier is motionless at the start of combat, it has advantage on its initiative roll. Moreover, if a creature hasn't observed the chandelier move or act, that creature must succeed on a DC 18 Intelligence (Investigation) check to discern that the chandelier is animate.

\DndMonsterAction{Fiery Aura}
Any creature that starts its turn within 5 feet of the chandelier takes 7 (2d6) fire damage.

\DndMonsterSection{Actions}
\DndMonsterAction{Multiattack}
The chandelier makes three Chain attacks, three Lamp attacks, or a combination thereof.

\DndMonsterAction{Chain}
\textsl{Melee Weapon Attack:} 7 to hit, reach 15 ft., one target. \textsl{Hit:} 13 (2d8 + 4) bludgeoning damage, and the target must succeed on a DC 15 Strength saving throw or be pulled into an unoccupied space within 5 feet of the chandelier.

\DndMonsterAction{Lamp}
\textsl{Melee Weapon Attack:} 7 to hit, reach 5 ft., one target. \textsl{Hit:} 9 (2d4 + 4) bludgeoning damage plus 13 (3d8) fire damage.

\DndMonsterAction{Blazing Vortex (Recharge 5\textendash 6)}
Each creature within 20 feet of the chandelier and not behind total cover must succeed on a DC 14 Constitution saving throw or take 36 (8d8) fire damage and have the blinded condition until the start of the chandelier's next turn.

\end{DndMonster}



\begin{DndMonster}[float=t]{Wight}
\DndMonsterType{Medium undead, neutral evil}
\DndMonsterBasics[
armor-class = {14 (studded leather armor)},
hit-points = {\DndDice{6d8 + 18}},
speed = {30 ft.}
]
\DndMonsterAbilityScores[
 str = 15,
 dex = 14,
 con = 16,
 int = 10,
 wis = 13,
 cha = 15
]
\DndMonsterDetails[
skills = {Perception +3, Stealth +4},
damage-resistances = {necrotic; bludgeoning, piercing, slashing from nonmagical attacks that aren't silvered},
damage-immunities = {poison},
condition-immunities = {exhaustion, poisoned},
senses = {darkvision 60 ft., passive perception 13},
languages = {the languages it knew in life},
challenge = 3,
]
\DndMonsterAction{Sunlight Sensitivity}
While in sunlight, the wight has disadvantage on attack rolls, as well as on Wisdom (Perception) checks that rely on sight.

\DndMonsterSection{Actions}
\DndMonsterAction{Multiattack}
The wight makes two longsword attacks or two longbow attacks. It can use its Life Drain in place of one longsword attack.

\DndMonsterAction{Life Drain}
\textsl{Melee Weapon Attack:} 4 to hit, reach 5 ft., one creature. \textsl{Hit:} 5 (1d6 + 2) necrotic damage. The target must succeed on a DC 13 Constitution saving throw or its hit point maximum is reduced by an amount equal to the damage taken. This reduction lasts until the target finishes a long rest. The target dies if this effect reduces its hit point maximum to 0.

A humanoid slain by this attack rises 24 hours later as a \textbf{zombie} under the wight's control, unless the humanoid is restored to life or its body is destroyed. The wight can have no more than twelve zombies under its control at one time.

\DndMonsterAction{Longsword}
\textsl{Melee Weapon Attack:} 4 to hit, reach 5 ft., one target. \textsl{Hit:} 6 (1d8 + 2) slashing damage, or 7 (1d10 + 2) slashing damage if used with two hands.

\DndMonsterAction{Longbow}
\textsl{Ranged Weapon Attack:} 4 to hit, range 150/600 ft., one target. \textsl{Hit:} 6 (1d8 + 2) piercing damage.

\end{DndMonster}



\begin{DndMonster}[float=t]{Will-o'-Wisp}
\DndMonsterType{Tiny undead, chaotic evil}
\DndMonsterBasics[
armor-class = {19},
hit-points = {\DndDice{9d4}},
speed = {0 ft., fly 50 ft. \textit{(hover)}}
]
\DndMonsterAbilityScores[
 str = 1,
 dex = 28,
 con = 10,
 int = 13,
 wis = 14,
 cha = 11
]
\DndMonsterDetails[
damage-resistances = {acid; cold; fire; necrotic; thunder; bludgeoning, piercing, slashing from nonmagical attacks},
damage-immunities = {lightning, poison},
condition-immunities = {exhaustion, grappled, paralyzed, poisoned, prone, restrained, unconscious},
senses = {darkvision 120 ft., passive perception 12},
languages = {the languages it knew in life},
challenge = 2,
]
\DndMonsterAction{Consume Life}
As a bonus action, the will-o'-wisp can target one creature it can see within 5 feet of it that has 0 hit points and is still alive. The target must succeed on a DC 10 Constitution saving throw against this magic or die. If the target dies, the will-o'-wisp regains 10 (3d6) hit points.

\DndMonsterAction{Ephemeral}
The will-o'-wisp can't wear or carry anything.

\DndMonsterAction{Incorporeal Movement}
The will-o'-wisp can move through other creatures and objects as if they were difficult terrain. It takes 5 (1d10) force damage if it ends its turn inside an object.

\DndMonsterAction{Variable Illumination}
The will-o'-wisp sheds bright light in a 5 to 20-foot radius and dim light for an additional number of ft. equal to the chosen radius. The will-o'-wisp can alter the radius as a bonus action.

\DndMonsterSection{Actions}
\DndMonsterAction{Shock}
\textsl{Melee Spell Attack:} 4 to hit, reach 5 ft., one creature. \textsl{Hit:} 9 (2d8) lightning damage.

\DndMonsterAction{Invisibility}
The will-o'-wisp and its light magically become invisible until it attacks or uses its Consume Life, or until its concentration ends (as if concentration on a spell).

\end{DndMonster}



\begin{DndMonster}[float*=b,width=\textwidth + 8pt]{Windfall}
\begin{multicols}{2}
\DndMonsterType{Medium humanoid (bard, tiefling), chaotic evil}
\DndMonsterBasics[
armor-class = {19 (\textit{studded leather armor})},
hit-points = {\DndDice{34d8 + 170}},
speed = {30 ft., fly 30 ft.}
]
\DndMonsterAbilityScores[
 str = 14,
 dex = 24,
 con = 20,
 int = 22,
 wis = 18,
 cha = 26
]
\DndMonsterDetails[
saving-throws = {Str +9, Dex +14, Wis +11, Cha +15},
skills = {Arcana +13, Deception +22, Insight +18, Perception +11, Performance +22, Persuasion +22, Sleight Of Hand +14},
damage-resistances = {acid, cold, fire, lightning, thunder},
condition-immunities = {charmed, frightened},
senses = {darkvision 60 ft., passive perception 21},
languages = {Common, Draconic, Infernal},
challenge = 23,
]
\DndMonsterAction{Dazzling Visage}
A brilliant array of chromatic colors emanates from Windfall, causing attack rolls against her to have disadvantage. This trait ceases to function while Windfall has the incapacitated condition or has a speed of 0.

\DndMonsterAction{Legendary Resistance (3/Day)}
If Windfall fails a saving throw, she can choose to succeed instead.

\DndMonsterAction{Special Equipment}
Windfall wears an iridescent magic coat that was tailored specifically for her and imbued with Tiamat's power. When she dies, the coat functions as a Robe of Scintillating Colors.

\DndMonsterAction{Spellcasting}
Windfall casts one of the following spells, requiring no material components and using Charisma as the spellcasting ability (spell save DC 23):
\begin{DndMonsterSpells}
  \DndInnateSpellLevel{Detect Magic}
  \DndInnateSpellLevel[]{Hold Monster}
  \DndInnateSpellLevel[3]{Shatter}
  \DndInnateSpellLevel[2]{Hypnotic Pattern}
\end{DndMonsterSpells}
\DndMonsterSection{Actions}
\DndMonsterAction{Multiattack}
Windfall makes two Chromatic Rapier attacks and uses Dragon's Fury once.

\DndMonsterAction{Chromatic Rapier}
\textsl{Melee Weapon Attack:} 14 to hit, reach 5 ft., one target. \textsl{Hit:} 11 (1d8 + 7) piercing damage plus 21 (6d6) acid, cold, fire, lightning, or poison damage (Windfall's choice).

\DndMonsterAction{Dragon's Fury}
Windfall targets one creature she can see within 60 feet of herself and unleashes a burst of magical ire. The target must make a DC 23 Wisdom saving throw. On a failed save, the target takes 36 (8d8) psychic damage and has the frightened condition until the start of Windfall's next turn. On a successful save, the target takes half as much damage only.

\DndMonsterSection{Bonus Actions}
\DndMonsterAction{Stunning Scintillation (Recharge 5\textendash 6)}
Windfall emits an overwhelming array of colors from her coat. Each creature within 30 feet of Windfall that can see her must succeed on a DC 23 Constitution saving throw or have the stunned condition until the start of Windfall's next turn.

\DndMonsterSection{Legendary Actions}
Windfall can take 3 legendary actions, choosing from the options below. Only one legendary action can be used at a time and only at the end of another creature's turn. Windfall regains spent legendary actions at the start of its turn.

\begin{DndMonsterLegendaryActions}
\DndMonsterLegendaryAction{Deft Dance}{Windfall moves up to her speed without provoking opportunity attacks.
}
\DndMonsterLegendaryAction{Dragon's Flare}{Windfall flares with multicolored flames and targets a creature she can see within 30 feet of herself. The target must make a DC 23 Dexterity saving throw. On a failed save, the target takes 26 (4d12) damage of a type chosen by Windfall: acid, cold, fire, lightning, or poison. On a successful save, the target takes half as much damage.
}
\DndMonsterLegendaryAction{Cast a Spell (Costs 2 Actions)}{Windfall uses Spellcasting.
}
\end{DndMonsterLegendaryActions}
\end{multicols}
\end{DndMonster}



\begin{DndMonster}[float=t]{Wraith}
\DndMonsterType{Medium undead, neutral evil}
\DndMonsterBasics[
armor-class = {13},
hit-points = {\DndDice{9d8 + 27}},
speed = {0 ft., fly 60 ft. \textit{(hover)}}
]
\DndMonsterAbilityScores[
 str = 6,
 dex = 16,
 con = 16,
 int = 12,
 wis = 14,
 cha = 15
]
\DndMonsterDetails[
damage-resistances = {acid; cold; fire; lightning; thunder; bludgeoning, piercing, slashing from nonmagical attacks that aren't silvered},
damage-immunities = {necrotic, poison},
condition-immunities = {charmed, exhaustion, grappled, paralyzed, petrified, poisoned, prone, restrained},
senses = {darkvision 60 ft., passive perception 12},
languages = {the languages it knew in life},
challenge = 5,
]
\DndMonsterAction{Incorporeal Movement}
The wraith can move through other creatures and objects as if they were difficult terrain. It takes 5 (1d10) force damage if it ends its turn inside an object.

\DndMonsterAction{Sunlight Sensitivity}
While in sunlight, the wraith has disadvantage on attack rolls, as well as on Wisdom (Perception) checks that rely on sight.

\DndMonsterSection{Actions}
\DndMonsterAction{Life Drain}
\textsl{Melee Weapon Attack:} 6 to hit, reach 5 ft., one creature. \textsl{Hit:} 21 (4d8 + 3) necrotic damage. The target must succeed on a DC 14 Constitution saving throw or its hit point maximum is reduced by an amount equal to the damage taken. This reduction lasts until the target finishes a long rest. The target dies if this effect reduces its hit point maximum to 0.

\DndMonsterAction{Create Specter}
The wraith targets a humanoid within 10 feet of it that has been dead for no longer than 1 minute and died violently. The target's spirit rises as a \textbf{specter} in the space of its corpse or in the nearest unoccupied space. The \textbf{specter} is under the wraith's control. The wraith can have no more than seven specters under its control at one time.

\end{DndMonster}



\begin{DndMonster}[float=t]{Yochlol}
\DndMonsterType{Medium fiend (demon, shapechanger), chaotic evil}
\DndMonsterBasics[
armor-class = {15 (natural armor)},
hit-points = {\DndDice{16d8 + 64}},
speed = {30 ft., climb 30 ft.}
]
\DndMonsterAbilityScores[
 str = 15,
 dex = 14,
 con = 18,
 int = 13,
 wis = 15,
 cha = 15
]
\DndMonsterDetails[
saving-throws = {Dex +6, Int +5, Wis +6, Cha +6},
skills = {Deception +10, Insight +6},
damage-resistances = {cold; fire; lightning; bludgeoning, piercing, slashing from nonmagical attacks},
damage-immunities = {poison},
condition-immunities = {poisoned},
senses = {darkvision 120 ft., passive perception 12},
languages = {Abyssal, Elvish, Undercommon},
challenge = 10,
]
\DndMonsterAction{Shapechanger}
The yochlol can use its action to polymorph into a form that resembles a female drow or giant spider, or back into its true form. Its statistics are the same in each form. Any equipment it is wearing or carrying isn't transformed. It reverts to its true form if it dies.

\DndMonsterAction{Magic Resistance}
The yochlol has advantage on saving throws against spells and other magical effects.

\DndMonsterAction{Spider Climb}
The yochlol can climb difficult surfaces, including upside down on ceilings, without needing to make an ability check.

\DndMonsterAction{Web Walker}
The yochlol ignores movement restrictions caused by webbing.

\DndMonsterAction{Innate Spellcasting}
The yochlol's spellcasting ability is Charisma (spell save DC 14). The yochlol can innately cast the following spells, requiring no material components:
\begin{DndMonsterSpells}
  \DndInnateSpellLevel{detect thoughts}
  \DndInnateSpellLevel[]{dominate person}
\end{DndMonsterSpells}
\DndMonsterSection{Actions}
\DndMonsterAction{Multiattack}
The yochlol makes two melee attacks.

\DndMonsterAction{Slam (Bite in Spider Form)}
\textsl{Melee Weapon Attack:} 6 to hit, reach 5 ft. (10 feet in demon form), one target. \textsl{Hit:} 5 (1d6 + 2) bludgeoning (piercing in spider form) damage plus 21 (6d6) poison damage.

\DndMonsterAction{Mist Form}
The yochlol transforms into toxic mist or reverts to its true form. Any equipment it is wearing or carrying is also transformed. It reverts to its true form if it dies.

While in mist form, the yochlol is incapacitated and can't speak. It has a flying speed of 30 feet, can hover, and can pass through any space that isn't airtight. It has advantage on Strength, Dexterity, and Constitution saving throws, and it is immune to nonmagical damage.

While in mist form, the yochlol can enter a creature's space and stop there. Each time that creature starts its turn with the yochlol in its space, the creature must succeed on a DC 14 Constitution saving throw or be poisoned until the start of its next turn. While poisoned in this way, the target is incapacitated.

\end{DndMonster}



\begin{DndMonster}[float=t]{Zombie}
\DndMonsterType{Medium undead, neutral evil}
\DndMonsterBasics[
armor-class = {8},
hit-points = {\DndDice{3d8 + 9}},
speed = {20 ft.}
]
\DndMonsterAbilityScores[
 str = 13,
 dex = 6,
 con = 16,
 int = 3,
 wis = 6,
 cha = 5
]
\DndMonsterDetails[
saving-throws = {Wis +0},
damage-immunities = {poison},
condition-immunities = {poisoned},
senses = {darkvision 60 ft., passive perception 8},
languages = {understands all languages it spoke in life but can't speak},
challenge = 1/4,
]
\DndMonsterAction{Undead Fortitude}
If damage reduces the zombie to 0 hit points, it must make a Constitution saving throw with a DC of 5 + the damage taken, unless the damage is radiant or from a critical hit. On a success, the zombie drops to 1 hit point instead.

\DndMonsterSection{Actions}
\DndMonsterAction{Slam}
\textsl{Melee Weapon Attack:} 3 to hit, reach 5 ft., one target. \textsl{Hit:} 4 (1d6 + 1) bludgeoning damage.

\end{DndMonster}


\end{document}
